% Generated mechanically from frozen input p05/ko/more-morphisms.tex.
% Do not edit this generated file; edit the bound locale source instead.

% OK, start here.
%

\setcounter{chapter}{36}
\chapter{사상 심화}



\phantomsection
\label{more-morphisms-section-phantom}


\section{들어가며}
\label{more-morphisms-section-introduction}

\noindent
이 장에서는 스킴의 사상이 갖는 성질을 계속 연구한다.
기본 참고문헌은 \cite{EGA}이다.






\section{비후화}
\label{more-morphisms-section-thickenings}

\noindent
다음 용어들은 완전히 표준적이지 않을 수 있지만 편리하다.

\begin{definition}
\label{more-morphisms-definition-thickening}
비후화(thickening).
\begin{enumerate}
\item 스킴 $X$가 $X'$의 닫힌 부분스킴이고 두 스킴의 바탕 위상공간이
같으면, 스킴 $X'$을 스킴 $X$의 {\it 비후화}라 한다.
\item $X$가 $X'$의 닫힌 부분스킴이고 $X$를 정의하는 준연접 아이디얼층
$\mathcal{I} \subset \mathcal{O}_{X'}$의 제곱이 영이면, 스킴 $X'$을
스킴 $X$의 {\it 1차 비후화}라 한다.
\item $X$가 $X'$의 닫힌 부분스킴이고 $X$를 정의하는 준연접 아이디얼층
$\mathcal{I} \subset \mathcal{O}_{X'}$이 멱영, 즉
$\mathcal{I}^{n + 1} = 0$을 만족하는 정수 $n \geq 0$이 존재하면,
스킴 $X'$을 스킴 $X$의 {\it 유한 차수 비후화}라 한다.
% P05-EN-CH37-0001: frozen source repeats X as its own ambient scheme.
\item $X$가 $X$의 닫힌 부분스킴이고 $\mathcal{I}^{n + 1} = 0$이며,
여기서 $\mathcal{I} \subset \mathcal{O}_{X'}$가 $X$를 정의하는 준연접
아이디얼층이면, 스킴 $X'$을 스킴 $X$의 {\it $n$차 비후화}라 한다.
\item 두 비후화 $X \subset X'$과 $Y \subset Y'$이 주어졌다고 하자.
{\it 비후화의 사상}은 $f'(X) \subset Y$, 즉 $f'|_X$가 닫힌 부분스킴
$Y$를 통해 분해되는 사상 $f' : X' \to Y'$이다. 이 상황에서
$f = f'|_X : X \to Y$로 놓고,
$(f, f') : (X \subset X') \to (Y \subset Y')$를 비후화의 사상이라 한다.
\item $S$를 스킴이라 하자. 마찬가지로 {$S$ 위의 비후화}와
{{$S$ 위의 비후화의 사상}}을 정의한다. 이는 위의 스킴
$X, X', Y, Y'$이 $S$ 위의 스킴이고 사상
$X \to X'$, $Y \to Y'$, $f' : X' \to Y'$이 $S$ 위의 사상이라는 뜻이다.
\end{enumerate}
\end{definition}

\noindent
$i_X : X \to X'$를 비후화라 하자. 핵
$\mathcal{I} = \Ker(i_X^\sharp)$의 모든 국소 절단은 국소적으로 멱영이다.
이로부터 $\mathcal{O}_{X'}$의 구조층을 어느 정도 제어할 수 있지만,
기술적으로는 유한 차수 비후화 $X \subset X'$의 부류를 다루기가 훨씬 쉽다.
% P05-EN-CH37-0002: frozen source says "order thinking".
구체적으로 $X'$이 $n$차 비후화이면 여과
$$
0 = \mathcal{I}^{n + 1} \subset
\mathcal{I}^n \subset
\mathcal{I}^{n - 1} \subset \ldots \subset
\mathcal{I} \subset \mathcal{O}_{X'}
$$
가 존재하며, $X'$은 닫힌 부분공간들의 여과
$$
X = X_1 \subset X_2 \subset \ldots \subset X_n \subset X_{n + 1} = X'
$$
를 갖는다. 각 쌍 $X_i \subset X_{i + 1}$은 $S$ 위의 1차 비후화이다.
간단한 귀납 논증을 사용하면 1차 비후화에 대해 증명된 많은 결과를
유한 차수 비후화에 관한 결과로 다시 서술할 수 있다.

\medskip\noindent
% P05-EN-CH37-0003: frozen source has singular/plural disagreement.
1차 비후화는 다음과 같이 기술된다(Modules, Lemma
\ref{modules-lemma-double-structure-gives-derivation}을 보라).

\begin{lemma}
\label{more-morphisms-lemma-first-order-thickening}
% P05-EN-CH37-0004: frozen source uses f^{-1}O_S without binding f.
$X$를 기저 $S$ 위의 스킴이라 하자. $X$ 위의 층들의 짧은 완전열
$$
0 \to \mathcal{I} \to \mathcal{A} \to \mathcal{O}_X \to 0
$$
을 생각하자. 여기서 $\mathcal{A}$는 $f^{-1}\mathcal{O}_S$-대수층이고,
$\mathcal{A} \to \mathcal{O}_X$는 $f^{-1}\mathcal{O}_S$-대수층의
전사이며, $\mathcal{I}$는 그 핵이다. 다음을 가정하자.
\begin{enumerate}
\item $\mathcal{I}$는 $\mathcal{A}$ 안에서 제곱이 영인 아이디얼이다.
\item $\mathcal{I}$는 $\mathcal{O}_X$-가군으로서 준연접이다.
\end{enumerate}
그러면 $X' = (X, \mathcal{A})$는 스킴이고 $X \to X'$는 $S$ 위의
1차 비후화이다. 또한 $S$ 위의 모든 1차 비후화는 이 꼴이다.
\end{lemma}

\begin{proof}
$X'$이 국소 환 달린 공간임은 분명하다. $U = \Spec(B)$를 $X$의 아핀
열린집합이라 하고 $A = \Gamma(U, \mathcal{A})$로 놓자.
$H^1(U, \mathcal{I}) = 0$이므로(Cohomology of Schemes, Lemma
\ref{coherent-lemma-quasi-coherent-affine-cohomology-zero})
사상 $A \to B$는 전사이다. 가정에 의해 핵 $I = \mathcal{I}(U)$는
환 $A$ 안에서 제곱이 영인 아이디얼이다. Schemes, Lemma
\ref{schemes-lemma-morphism-into-affine}에 의해 국소 환 달린 공간의
표준 사상
$$
(U, \mathcal{A}|_U) \longrightarrow \Spec(A)
$$
이 존재하며, 이는 사상 $B \to \Gamma(U, \mathcal{A})$에서 온다.
% P05-EN-CH37-0005: frozen source has B in this map; no reverse map was bound.
이 사상은 가환도식
$$
\xymatrix{
(U, \mathcal{O}_X|_U) \ar[d] \ar[r] & \Spec(B) \ar[d] \\
(U, \mathcal{A}|_U) \ar[r] & \Spec(A)
}
$$
에 들어맞으므로 바탕 위상공간 위의 위상동형이다. 따라서 이 사상이
동형사상임을 보이려면 국소환 위에 동형사상을 유도함을 확인하면 충분하다.
소 아이디얼 $\mathfrak p \subset A$에 대응하는 $u \in U$에 대해
짧은 완전열들의 가환도식
$$
\xymatrix{
0 \ar[r] &
I_{\mathfrak p} \ar[r] \ar[d] &
A_{\mathfrak p} \ar[r] \ar[d] &
B_{\mathfrak p} \ar[r] \ar[d] & 0 \\
0 \ar[r] &
\mathcal{I}_u \ar[r] &
\mathcal{A}_u \ar[r] &
\mathcal{O}_{X, u} \ar[r] & 0.
}
$$
을 얻는다. $\mathcal{I}$와 $\mathcal{O}_X$가 준연접층이므로 왼쪽과
오른쪽 수직 화살표는 동형사상이다. 따라서 가운데 사상도 동형사상이다.
그러므로 $X' = (X, \mathcal{A})$의 모든 점은 아핀 근방을 가지며,
$X'$은 원하는 대로 스킴이다.
\end{proof}

\begin{lemma}
\label{more-morphisms-lemma-thickening-affine-scheme}
\begin{slogan}
아핀성은 비후화에 영향을 받지 않는다
\end{slogan}
\begin{reference}
유한 차수 비후화의 경우는 \cite[Proposition 5.1.9]{EGA1}이다.
\end{reference}
아핀 스킴의 모든 비후화는 아핀이다.
\end{lemma}

\begin{proof}
이는 Limits, Proposition \ref{limits-proposition-affine}의 특수한 경우이다.
\end{proof}

\begin{proof}[유한 차수 비후화에 대한 증명]
$X \subset X'$이 유한 차수 비후화이고 $X$가 아핀이라고 하자. 세르의
판정법을 사용하여 $X'$이 아핀임을 증명할 수 있다. 더 정확히는
Cohomology of Schemes, Lemma
\ref{coherent-lemma-quasi-compact-h1-zero-covering}을 사용하겠다.
$\mathcal{F}$를 준연접 $\mathcal{O}_{X'}$-가군이라 하자.
$H^1(X', \mathcal{F}) = 0$임을 보이면 충분하다. 주어진 닫힌 몰입을
$i : X \to X'$로 나타내고
$\mathcal{I} = \Ker(i^\sharp : \mathcal{O}_{X'} \to i_*\mathcal{O}_X)$.
로 나타내자. Definition \ref{more-morphisms-definition-thickening} 뒤의 유한 차수
비후화에 관한 논의에 의해 $n \geq 0$과 준연접 부분가군들의 여과
$$
0 = \mathcal{F}_{n + 1} \subset \mathcal{F}_n \subset
\mathcal{F}_{n - 1} \subset \ldots \subset
\mathcal{F}_0 = \mathcal{F}
$$
가 존재하며, $\mathcal{F}_a/\mathcal{F}_{a + 1}$은 $\mathcal{I}$에 의해
소멸한다. 실제로 $\mathcal{F}_a = \mathcal{I}^a\mathcal{F}$로 취할 수 있다.
그러면 어떤 준연접 $\mathcal{O}_X$-가군 $\mathcal{G}_a$에 대해
$\mathcal{F}_a/\mathcal{F}_{a + 1} = i_*\mathcal{G}_a$이다. Morphisms,
Lemma \ref{morphisms-lemma-i-star-equivalence}를 보라. 따라서
$$
H^1(X', \mathcal{F}_a/\mathcal{F}_{a + 1}) =
H^1(X', i_*\mathcal{G}_a) = H^1(X, \mathcal{G}_a) = 0
$$
을 얻는다. 두 번째 등식은 Cohomology of Schemes, Lemma
\ref{coherent-lemma-relative-affine-cohomology}에서, 마지막 등식은
Cohomology of Schemes, Lemma
\ref{coherent-lemma-quasi-coherent-affine-cohomology-zero}에서 따른다.
따라서 $\mathcal{F}$는 모든 연속 몫의 1차 코호몰로지가 소멸하는 유한
여과를 가지며, 간단한 귀납 논증으로 $H^1(X', \mathcal{F}) = 0$이다.
\end{proof}

\begin{lemma}
\label{more-morphisms-lemma-base-change-thickening}
$S \subset S'$을 스킴의 비후화라 하자. 사상 $X' \to S'$이 주어졌다고 하고
$X = S \times_{S'} X'$로 놓자. 그러면
$(X \subset X') \to (S \subset S')$은 비후화의 사상이다.
$S \subset S'$이 1차 비후화(각각 유한 차수 비후화)이면
$X \subset X'$도 1차 비후화(각각 유한 차수 비후화)이다.
\end{lemma}

\begin{proof}
생략한다.
\end{proof}

\begin{lemma}
\label{more-morphisms-lemma-composition-thickening}
\begin{slogan}
비후화의 합성은 비후화이다
\end{slogan}
$S \subset S'$과 $S' \subset S''$이 비후화이면 $S \subset S''$도
비후화이다.
\end{lemma}

\begin{proof}
생략한다.
\end{proof}

\begin{lemma}
\label{more-morphisms-lemma-descending-property-thickening}
비후화라는 성질은 fpqc-국소적이다. 1차 비후화에 대해서도 마찬가지이다.
\end{lemma}

\begin{proof}
이 명제의 뜻은 다음과 같다. $X \to X'$을 스킴의 사상이라 하고,
$\{g_i : X'_i \to X'\}$를 모든 $i$에 대해 기저변환
$X_i \to X'_i$가 비후화인 fpqc 덮개라 하자. 그러면 $X \to X'$은
비후화이다. 사상들 $g_i$가 함께 전사이므로 $X \to X'$이 전사임을
알 수 있다. Descent, Lemma
\ref{descent-lemma-descending-property-closed-immersion}에 의해
$X \to X'$은 닫힌 몰입이다. 따라서 $X \to X'$은 비후화이다.
1차 비후화의 경우에 대한 증명은 생략한다.
\end{proof}









\section{비후화의 사상}
\label{more-morphisms-section-morphisms-thickenings}

\noindent
$(f, f') : (X \subset X') \to (Y \subset Y')$가 스킴의 비후화들 사이의
사상이면, 흔히 사상 $f$의 성질이 $f'$에도 계승된다. 여기에는 여러
변형이 있다.

\begin{lemma}
\label{more-morphisms-lemma-thicken-property-morphisms}
$(f, f') : (X \subset X') \to (S \subset S')$를 비후화들의 사상이라
하자. 그러면
\begin{enumerate}
\item $f$가 아핀 사상일 필요충분조건은 $f'$이 아핀 사상인 것이다.
\item $f$가 전사 사상일 필요충분조건은 $f'$이 전사 사상인 것이다.
% P05-EN-CH37-0006
\item $f$가 준콤팩트일 필요충분조건은 $f'$이 준콤팩트인 것이다.
\item $f$가 보편 닫힌 사상일 필요충분조건은 $f'$이 보편 닫힌 사상인 것이다.
\item $f$가 정수적일 필요충분조건은 $f'$이 정수적인 것이다.
\item $f$가 (준)분리일 필요충분조건은 $f'$이 (준)분리인 것이다.
\item $f$가 보편 단사일 필요충분조건은 $f'$이 보편 단사인 것이다.
\item $f$가 보편 열린 사상일 필요충분조건은 $f'$이 보편 열린 사상인 것이다.
\item $f$가 준아핀일 필요충분조건은 $f'$이 준아핀인 것이다.
% P05-EN-CH37-0007
\item 여기에 더 추가할 것.
\end{enumerate}
\end{lemma}

\begin{proof}
$S \to S'$과 $X \to X'$은 보편 위상동형사상임에 유의하자(예를 들어
Morphisms, Lemma \ref{morphisms-lemma-reduction-universal-homeomorphism} 참조).
이로부터 (2), (3), (4), (7), (8)이 즉시 따른다. (1)은 $S$와 $S'$의
아핀 열린집합들 사이 및 $X$와 $X'$의 아핀 열린집합들 사이에 일대일
대응이 있음을 알려 주는 Lemma \ref{more-morphisms-lemma-thickening-affine-scheme}에서
따른다. (9)는 Limits, Lemma \ref{limits-lemma-thickening-quasi-affine}와
방금 언급한 $S$와 $S'$의 아핀 열린집합에 관한 사실에서 따른다.
(5)는 Morphisms, Lemma
\ref{morphisms-lemma-integral-universally-closed}에 의해 (1)과 (4)에서
따른다. 마지막으로 다음 사상이
% P05-EN-CH37-0008
$$
S \times_X S = S \times_{X'} S \to S \times_{X'} S' \to S' \times_{X'} S'
$$
비후화임에 유의하자(두 화살표는 Lemma
\ref{more-morphisms-lemma-base-change-thickening}에 의해 비후화이다). 따라서 비후화의
사상 $(S \subset S') \to (S \times_X S \to S' \times_{X'} S')$에
(3)과 (4)를 적용하면 (6)을 얻는다.
\end{proof}

\begin{lemma}
\label{more-morphisms-lemma-thicken-property-relatively-ample}
$(f, f') : (X \subset X') \to (S \subset S')$를 비후화들의 사상이라
하자. $\mathcal{L}'$을 $X'$ 위의 가역층이라 하고, $\mathcal{L}$로
$X$에 대한 그 제한을 나타내자. 그러면 $\mathcal{L}'$이 $f'$-풍부일
필요충분조건은 $\mathcal{L}$이 $f$-풍부인 것이다.
\end{lemma}

\begin{proof}
상대적으로 풍부하다는 것은 밑의 각 아핀 열린집합에 대한 조건임을
상기하자. Morphisms, Definition
\ref{morphisms-definition-relatively-ample}을 보라. Lemma
\ref{more-morphisms-lemma-thickening-affine-scheme}에 의해 $S$와 $S'$의 아핀
열린집합들 사이에는 일대일 대응이 있다. 따라서 $S$와 $S'$이 아핀이라고
가정할 수 있고, $\mathcal{L}'$이 풍부일 필요충분조건은
$\mathcal{L}$이 풍부인 것임을 증명하는 문제로 환원된다. 이는 Limits,
Lemma \ref{limits-lemma-ample-on-reduction}이다.
\end{proof}

\begin{lemma}
\label{more-morphisms-lemma-thicken-property-morphisms-cartesian}
$(f, f') : (X \subset X') \to (S \subset S')$를
$X = S \times_{S'} X'$을 만족하는 비후화들의 사상이라 하자.
$S \subset S'$이 유한 차수 비후화이면 다음이 성립한다.
\begin{enumerate}
\item $f$가 닫힌 몰입일 필요충분조건은 $f'$이 닫힌 몰입인 것이다.
\item $f$가 국소 유한형일 필요충분조건은 $f'$이 국소 유한형인 것이다.
\item $f$가 국소 준유한일 필요충분조건은 $f'$이 국소 준유한인 것이다.
\item $f$가 상대 차원 $d$인 국소 유한형일 필요충분조건은 $f'$이 상대
차원 $d$인 국소 유한형인 것이다.
\item $\Omega_{X/S} = 0$일 필요충분조건은 $\Omega_{X'/S'} = 0$인 것이다.
\item $f$가 비분기일 필요충분조건은 $f'$이 비분기인 것이다.
\item $f$가 고유일 필요충분조건은 $f'$이 고유인 것이다.
\item $f$가 유한일 필요충분조건은 $f'$이 유한인 것이다.
\item $f$가 모노사상일 필요충분조건은 $f'$이 모노사상인 것이다.
\item $f$가 몰입일 필요충분조건은 $f'$이 몰입인 것이다.
\item 여기에 더 추가할 것.
\end{enumerate}
\end{lemma}

\begin{proof}
보조정리에 열거된 성질 $\mathcal{P}$는 모두 밑변환에 안정적이다. 따라서
$f'$이 성질 $\mathcal{P}$를 가지면 $f$도 그러하다. Schemes, Lemmas
\ref{schemes-lemma-base-change-immersion} 및
\ref{schemes-lemma-base-change-monomorphism}, 그리고 Morphisms, Lemmas
\ref{morphisms-lemma-base-change-finite-type},
\ref{morphisms-lemma-base-change-quasi-finite},
\ref{morphisms-lemma-base-change-relative-dimension-d},
\ref{morphisms-lemma-base-change-differentials},
\ref{morphisms-lemma-base-change-unramified},
\ref{morphisms-lemma-base-change-proper},
\ref{morphisms-lemma-base-change-finite}을 보라.

\medskip\noindent
그러므로 각 경우에서 실질적인 방향은 $f$가 그 성질을 가진다고 가정하고
$f'$도 그 성질을 가짐을 도출하는 것이다. 비후화의 차수에 대한 귀납법에
의해 $S \subset S'$이 1차 비후화라고 가정할 수 있다. Definition
\ref{more-morphisms-definition-thickening} 바로 뒤의 논의를 보라.

\medskip\noindent
대부분의 증명에서는 아핀 경우로의 환원을 사용한다. $U' \subset S'$을
아핀 열린집합이라 하고, $V' \subset X'$을 $U'$ 위에 놓인 아핀
열린집합이라 하자. $U' = \Spec(A')$이라 하고, 닫힌 부분스킴
$U' \cap S$를 정의하는 아이디얼을 $I \subset A'$로 나타내자.
% P05-EN-CH37-0009
$V' = \Spec(B')$이라 하자. 그러면
$V' \cap X = \Spec(B'/IB')$이다. $A = A'/I$, $B = B'/IB'$로 놓으면
다음 가환 도식을 얻는다.
$$
\xymatrix{
0 \ar[r] &
IB' \ar[r] &
B' \ar[r] &
B \ar[r] & 0 \\
0 \ar[r] &
IA' \ar[r] \ar[u] &
A' \ar[r] \ar[u] &
A \ar[r] \ar[u] & 0
}
$$
두 행은 완전하고 $I^2 = 0$이다.

\medskip\noindent
(1)을 대수학으로 옮기면 다음과 같다. $A \to B$가 전사이면
$A' \to B'$도 전사이다. 이는 Nakayama 보조정리(Algebra, Lemma
\ref{algebra-lemma-NAK})에서 따른다.

\medskip\noindent
(2)를 대수학으로 옮기면 다음과 같다. $A \to B$가 유한형 환 사상이면
$A' \to B'$도 유한형 환 사상이다. 이는
$A[x_1, \ldots, x_n] \to B$가 전사이도록 택한 사상
$A'[x_1, \ldots, x_n] \to B'$에 Nakayama 보조정리(Algebra, Lemma
\ref{algebra-lemma-NAK})를 적용하여 얻는다.

\medskip\noindent
(3)의 증명. (2), 그리고 국소 유한형인 사상의 준유한성은 올에서 확인할
수 있다는 사실에서 따른다. Morphisms, Lemma
\ref{morphisms-lemma-quasi-finite-at-point-characterize}를 보라.

\medskip\noindent
(4)의 증명. (2), 그리고 ``상대 차원 $d$이다''라는 추가 성질은 올에서
확인할 수 있다는 사실에서 따른다(정의에 의해 그렇다. Morphisms,
Definition \ref{morphisms-definition-relative-dimension-d} 참조).
% P05-EN-CH37-0010

\medskip\noindent
(5)를 대수학으로 옮기면 다음과 같다. $\Omega_{B/A} = 0$이면
$\Omega_{B'/A'} = 0$이다. Algebra, Lemma
\ref{algebra-lemma-differentials-base-change}에 의해
$0 = \Omega_{B/A} = \Omega_{B'/A'}/I\Omega_{B'/A'}$이다. 따라서
Nakayama 보조정리(Algebra, Lemma \ref{algebra-lemma-NAK})에 의해
$\Omega_{B'/A'} = 0$이다.

\medskip\noindent
(6)을 대수학으로 옮기면 다음과 같다. $A \to B$가 비분기 사상이면
$A' \to B'$도 비분기이다. $A \to B$가 유한형이므로 위의 (2)에 의해
$A' \to B'$도 유한형이다. $A \to B$가 비분기이므로
$\Omega_{B/A} = 0$이고, (5)에 의해 $\Omega_{B'/A'} = 0$이다.
따라서 $A' \to B'$는 비분기이다.
% P05-EN-CH37-0011

\medskip\noindent
(7)의 증명. (2), Lemma \ref{more-morphisms-lemma-thicken-property-morphisms}의 결과,
그리고 고유가 준콤팩트 $+$ 분리 $+$ 국소 유한형 $+$ 보편 닫힌과
같다는 사실을 결합하면 된다.

\medskip\noindent
(8)의 증명. (2)와 Lemma \ref{more-morphisms-lemma-thicken-property-morphisms}의 결과를
결합하고, 유한이 정수적 $+$ 국소 유한형과 같다는 사실(Morphisms,
Lemma \ref{morphisms-lemma-finite-integral})을 사용하면 된다.

\medskip\noindent
(9)의 증명. $f$가 모노사상이므로 $X = X \times_S X$이다. 지금까지
증명한 결과들을 비후화의 사상
$(X \subset X') \to (X \times_S X \subset X' \times_{S'} X')$에
적용할 수 있다. (1)에 의해 $X' \to X' \times_{S'} X'$은 닫힌
몰입이다. 사실 이는 1차 비후화이다. 닫힌 몰입
$X' \to X' \times_{S'} X'$을 정의하는 아이디얼은 $S'$ 안에서
$S$를 잘라내는 아이디얼
$\mathcal{I} \subset \mathcal{O}_{S'}$의 역상에 포함되기 때문이다.
실제로 $X = X \times_S X = (X' \times_{S'} X') \times_{S'} S$는
$X'$에 포함된다. 따라서 Morphisms, Lemma
\ref{morphisms-lemma-differentials-diagonal}에 의해
$\Omega_{X'/S'} = 0$임을 보이면 충분한데, 이는 (5)와 $X/S$에 대한
대응하는 명제에서 따른다.

\medskip\noindent
(10)의 증명. $f : X \to S$가 몰입이면, $U \to S$는 열린 몰입이고
$X \to U$는 닫힌 몰입인 $X \to U \to S$로 인수분해된다.
$U' \subset S'$을 그 바탕 위상공간이 $U$와 같은 열린 부분스킴이라
하자. 그러면 $X' \to S'$은 $U'$을 통하여 인수분해되고, (1)에 의해
$X' \to U'$이 닫힌 몰입이라고 결론짓는다. 이것으로 증명이 끝난다.
\end{proof}

\noindent
다음 보조정리는 앞의 보조정리의 변형이다. 여기서는 관련된 비후화들이
유한 차수라고 가정하여 국소 유한형이라는 성질을 $f$에서 $f'$으로
옮기는 대신, $f$와 $f'$이 각각 국소 유한형이라는 것을 가정한다.

\begin{lemma}
\label{more-morphisms-lemma-properties-that-extend-over-thickenings}
$(f, f') : (X \subset X') \to (Y \to Y')$를 비후화들의 사상이라
하자. $f$와 $f'$이 국소 유한형이고 $X = Y \times_{Y'} X'$이라고
가정하자. 그러면
% P05-EN-CH37-0012
\begin{enumerate}
\item $f$가 국소 준유한일 필요충분조건은 $f'$이 국소 준유한인 것이다.
\item $f$가 유한일 필요충분조건은 $f'$이 유한인 것이다.
\item $f$가 닫힌 몰입일 필요충분조건은 $f'$이 닫힌 몰입인 것이다.
\item $\Omega_{X/Y} = 0$일 필요충분조건은 $\Omega_{X'/Y'} = 0$인 것이다.
\item $f$가 비분기일 필요충분조건은 $f'$이 비분기인 것이다.
\item $f$가 모노사상일 필요충분조건은 $f'$이 모노사상인 것이다.
\item $f$가 몰입일 필요충분조건은 $f'$이 몰입인 것이다.
\item $f$가 고유일 필요충분조건은 $f'$이 고유인 것이다.
% P05-EN-CH37-0013
\item 여기에 더 추가할 것.
\end{enumerate}
\end{lemma}

\begin{proof}
보조정리에 열거된 성질 $\mathcal{P}$는 모두 밑변환에 안정적이다. 따라서
$f'$이 성질 $\mathcal{P}$를 가지면 $f$도 그러하다. Schemes, Lemmas
\ref{schemes-lemma-base-change-immersion} 및
\ref{schemes-lemma-base-change-monomorphism}, 그리고 Morphisms, Lemmas
\ref{morphisms-lemma-base-change-quasi-finite},
\ref{morphisms-lemma-base-change-relative-dimension-d},
\ref{morphisms-lemma-base-change-differentials},
\ref{morphisms-lemma-base-change-unramified},
\ref{morphisms-lemma-base-change-proper},
\ref{morphisms-lemma-base-change-finite}을 보라. 따라서 각 경우에는
$f$가 원하는 성질을 가지면 $f'$도 그러함만 증명하면 된다.

\medskip\noindent
사상이 국소 준유한일 필요충분조건은 그것이 국소 유한형이고 그
스킴론적 올들이 이산공간인 것이다. Morphisms, Lemma
\ref{morphisms-lemma-locally-quasi-finite-fibres}를 보라. 비후화로
넘어갈 때 바탕 위상공간은 변하지 않으므로, $f'$이 국소 준유한일
필요충분조건은 $f$가 국소 준유한인 것임을 알 수 있다. 이것으로 (1)이
증명된다.

\medskip\noindent
(2)는 Lemma \ref{more-morphisms-lemma-thicken-property-morphisms}의 (5)와, 유한
사상은 정확히 국소 유한형인 정수적 사상이라는 사실(Morphisms, Lemma
\ref{morphisms-lemma-finite-integral})에서 따른다.

\medskip\noindent
(3)은 Morphisms, Lemma
\ref{morphisms-lemma-check-closed-infinitesimally}에서 즉시 따른다.

\medskip\noindent
(4)는 다음 대수학 명제에서 따른다. $A$를 환이라 하고
$I \subset A$를 국소 멱영 아이디얼이라 하자. $B$를 유한형
$A$-대수라 하자. $\Omega_{(B/IB)/(A/I)} = 0$이면
$\Omega_{B/A} = 0$이다. 실제로 이 가정은
$I\Omega_{B/A} = 0$임을 뜻한다. Algebra, Lemma
\ref{algebra-lemma-differentials-base-change}를 보라.
% P05-EN-CH37-0014
한편 $\Omega_{B/A}$는 유한 $B$-가군이다. Algebra, Lemma
\ref{algebra-lemma-differentials-finitely-generated}를 보라. 따라서
$IB$가 $B$의 Jacobson 근기에 포함된다는 사실과 Nakayama 보조정리
(Algebra, Lemma \ref{algebra-lemma-NAK})로부터
$\Omega_{B/A}$의 소멸이 따른다.

\medskip\noindent
(5)는 (4)와 Morphisms, Lemma
\ref{morphisms-lemma-unramified-omega-zero}에서 즉시 따른다.

\medskip\noindent
(6)의 증명. $f$가 모노사상이므로 $X = X \times_Y X$이다. 지금까지
증명한 결과를 비후화의 사상
$(X \subset X') \to (X \times_Y X \subset X' \times_{Y'} X')$에
적용할 수 있다. (3)에 의해
$\Delta_{X'/Y'} : X' \to X' \times_{Y'} X'$은 닫힌 몰입이다. 실제로
$\Delta_{X'/Y'}$는 바탕 집합들 위에서 전단사이므로
$\Delta_{X'/Y'}$는 비후화이다. 한편 Morphisms, Lemma
\ref{morphisms-lemma-diagonal-morphism-finite-type}에 의해
$\Delta_{X'/Y'}$는 국소 유한표현이다. 달리 말해,
$\Delta_{X'/Y'}(X')$는 유한형 준연접 아이디얼 층
$\mathcal{J} \subset \mathcal{O}_{X' \times_{Y'} X'}$에 의해 잘려
나온다. (5)에 의해 $\Omega_{X'/Y'} = 0$이므로, Morphisms, Lemma
\ref{morphisms-lemma-differentials-diagonal}에 의해
$X' \to X' \times_{Y'} X'$의 conormal 층이 영이다. 달리 말해
$\mathcal{J}/\mathcal{J}^2 = 0$이다. 예를 들어 Algebra, Lemma
\ref{algebra-lemma-ideal-is-squared-union-connected}에 의해 이로부터
$\Delta_{X'/Y'}$가 동형임이 따른다.

\medskip\noindent
(7)의 증명. $f : X \to Y$가 몰입이면, $V \to Y$는 열린 몰입이고
$X \to V$는 닫힌 몰입인 $X \to V \to Y$로 인수분해된다.
$V' \subset Y'$을 그 바탕 위상공간이 $V$와 같은 열린 부분스킴이라
하자. 그러면 $X' \to V'$은 $V'$을 통하여 인수분해되고, (3)에 의해
$X' \to V'$이 닫힌 몰입이라고 결론짓는다.
% P05-EN-CH37-0015

\medskip\noindent
(8)은 Lemma \ref{more-morphisms-lemma-thicken-property-morphisms}와, 고유 사상은
국소 유한형이면서 준콤팩트이고 보편 닫힌이며 분리인 사상이라는 정의에서
따른다.
\end{proof}








\section{비후화의 Picard 군}
\label{more-morphisms-section-picard-group-thickening}

\noindent
비후화의 Picard 군에 관한 몇 가지 내용을 다룬다.

\begin{lemma}
\label{more-morphisms-lemma-picard-group-first-order-thickening}
$X \subset X'$을 아이디얼 층 $\mathcal{I}$를 갖는 1차 비후화라 하자.
그러면 다음과 같은 아벨 군의 표준 완전열이 존재한다.
$$
\xymatrix{
0 \ar[r] &
H^0(X, \mathcal{I}) \ar[r] &
H^0(X', \mathcal{O}_{X'}^*) \ar[r] &
H^0(X, \mathcal{O}^*_X) \ar `r[d] `d[l] `l[llld] `d[dll] [dll] \\
& H^1(X, \mathcal{I}) \ar[r] &
\Pic(X') \ar[r] &
\Pic(X) \ar `r[d] `d[l] `l[llld] `d[dll] [dll] \\
& H^2(X, \mathcal{I}) \ar[r] & \ldots \ar[r] & \ldots
}
$$
\end{lemma}

\begin{proof}
이는 다음 아벨 군 층의 짧은 완전열에 결부된 긴 코호몰로지 완전열이다.
$$
0 \to \mathcal{I} \to \mathcal{O}_{X'}^* \to \mathcal{O}_X^* \to 0
$$
여기서 첫째 사상은 $\mathcal{I}$의 국소 절단 $f$를
$\mathcal{O}_{X'}$의 가역 절단 $1 + f$로 보낸다. 또한 환 달린 공간의
Picard 군을 가역함수 층의 첫째 코호몰로지 군과 동일시하는 사실을
사용한다. Cohomology, Lemma \ref{cohomology-lemma-h1-invertible}을 보라.
\end{proof}

\begin{lemma}
\label{more-morphisms-lemma-torsion-pic-thickening}
$X \subset X'$을 비후화라 하자. $n$을 $\mathcal{O}_X$에서 가역인
정수라 하자. 그러면 사상 $\Pic(X')[n] \to \Pic(X)[n]$은 전단사이다.
\end{lemma}

\begin{proof}[유한 차수 비후화의 경우에 대한 증명]
Definition \ref{more-morphisms-definition-thickening} 뒤에서 설명한 일반 원리에 의해
1차 비후화의 경우로 환원된다. 그러면 Lemma
\ref{more-morphisms-lemma-picard-group-first-order-thickening}을 사용하여
$H^1(X, \mathcal{I})[n]$, $H^1(X, \mathcal{I})/n$ 및
$H^2(X, \mathcal{I})[n]$이 영임을 보이면 충분함을 알 수 있다.
% P05-EN-CH37-0016
$\mathcal{I}$는 $\mathcal{O}_X$-가군이고 그 위에서 $n$배 사상이
동형이므로 이 소멸이 따른다.
\end{proof}

\begin{proof}[일반적인 경우의 증명]
$\mathcal{I} \subset \mathcal{O}_{X'}$을 $X$를 잘라내는 준연접
아이디얼 층이라 하자. 그러면 다음 아벨 군의 짧은 완전열이 있다.
$$
0 \to (1 + \mathcal{I})^* \to \mathcal{O}_{X'}^* \to \mathcal{O}_X^* \to 0
$$
Lemma \ref{more-morphisms-lemma-picard-group-first-order-thickening}의 명제에서
$H^i(X, \mathcal{I})$를 $H^i(X, (1 + \mathcal{I})^*)$로 바꾼 것과
같은 긴 코호몰로지 완전열을 얻는다. 따라서 $n$제곱 사상이
$(1 + \mathcal{I})^* \to (1 + \mathcal{I})^*$의 동형임을 보이면
충분하다. 아핀 열린집합 위의 절단을 취하면 이는 Algebra, Lemma
\ref{algebra-lemma-lift-nth-roots}에서 따른다.
\end{proof}





\section{무한소 근방}
\label{more-morphisms-section-first-order-infinitesimal-neighbourhood}

\noindent
유한 차수 비후화를 자연스럽게 구성하는 방법은 다음과 같다.
$i : Z \to X$를 스킴의 몰입이라 하자.
% P05-EN-CH37-0017
$i$가 $Z$를 닫힌 부분스킴 $Z \subset U$와 동일시하도록 하는 열린
부분스킴 $U \subset X$를 택하자. $\mathcal{I} \subset \mathcal{O}_U$를
$U$ 안의 $Z$를 정의하는 준연접 아이디얼 층이라 하자. $n \geq 1$에
대해 준연접 아이디얼 층 $\mathcal{I}^{n + 1}$로 정의되는 닫힌 부분스킴
$Z_n \subset U$를 생각할 수 있다.

\begin{definition}
\label{more-morphisms-definition-first-order-infinitesimal-neighbourhood}
$i : Z \to X$를 스킴의 몰입이라 하자.
\begin{enumerate}
\item $X$ 안의 $Z$의 {\it 1차 무한소 근방}은 위에서 설명한 $X$ 위의
1차 비후화 $Z \subset Z_1$이다.
\item $X$ 안의 $Z$의 {\it $n$차 무한소 근방}은 위에서 설명한 $X$ 위의
$n$차 비후화 $Z \subset Z_n$이다.
\end{enumerate}
\end{definition}

\noindent
이 비후화들은 다음 보편 성질을 갖는다(따라서 위의 구성이 열린집합
$U$의 선택에 의존할지 모른다는 우려가 사라진다).

\begin{lemma}
\label{more-morphisms-lemma-first-order-infinitesimal-neighbourhood}
$i : Z \to X$를 스킴의 몰입이라 하자.
\begin{enumerate}
\item $X$ 안의 $Z$의 1차 무한소 근방 $Z'$은 다음 보편 성질을 갖는다.
임의의 가환 도식
$$
\xymatrix{
Z \ar[d]_i & T \ar[l]^a \ar[d] \\
X & T' \ar[l]_b
}
$$
이 주어지고 $T \subset T'$이 $X$ 위의 1차 비후화이면, $X$ 위의
비후화의 사상
$(a', a) : (T \subset T') \to (Z \subset Z')$가 유일하게 존재한다.
% P05-EN-CH37-0018
\item $n \geq 1$에 대해 $X$ 안의 $Z$의 $n$차 무한소 근방 $Z_n$은
다음 보편 성질을 갖는다. 임의의 가환 도식
$$
\xymatrix{
Z \ar[d]_i & T \ar[l]^a \ar[d] \\
X & T' \ar[l]_b
}
$$
이 주어지고 $T \subset T'$이 $X$ 위의 $n$차 비후화이면, $X$ 위의
비후화의 사상
$(a', a) : (T \subset T') \to (Z \subset Z_n)$가 유일하게 존재한다.
% P05-EN-CH37-0018
\end{enumerate}
\end{lemma}

\begin{proof}
(1)만 증명한다. $U \subset X$를 $Z'$의 구성에 사용한 열린집합, 즉
준연접 아이디얼 층 $\mathcal{I}$에 의해 잘려 나오는 $U$의 닫힌
부분스킴과 $Z$가 동일시되도록 하는 열린집합이라 하자.
$|T| = |T'|$이므로 $b(T') \subset U$임을 알 수 있다. 따라서 $b$를
$U$로 가는 사상으로 생각할 수 있다.
$\mathcal{J} \subset \mathcal{O}_{T'}$를 $T$를 잘라내는 아이디얼이라
하자. 위 도식에 의해 $b(T) \subset Z$이므로
$b^\sharp(b^{-1}\mathcal{I}) \subset \mathcal{J}$이다. $T'$이 $T$의
1차 비후화이므로 $\mathcal{J}^2 = 0$이고, 따라서
$b^\sharp(b^{-1}(\mathcal{I}^2)) = 0$이다. Schemes, Lemma
\ref{schemes-lemma-characterize-closed-subspace}에 의해 이로부터 $b$가
$Z'$을 통하여 인수분해됨이 따른다. 이 인수분해를
$a' : T' \to Z'$로 나타내면 나머지는 명백하다.
% P05-EN-CH37-0019
\end{proof}

\begin{lemma}
\label{more-morphisms-lemma-infinitesimal-neighbourhood-conormal}
$i : Z \to X$를 스킴의 몰입이라 하고, $Z \subset Z'$을 $X$ 안의
$Z$의 1차 무한소 근방이라 하자. 그러면 도식
$$
\xymatrix{
Z \ar[r] \ar[d] & Z' \ar[d] \\
Z \ar[r] & X
}
$$
은 Morphisms, Lemma \ref{morphisms-lemma-conormal-functorial}에 의해
conormal 층의 사상 $\mathcal{C}_{Z/X} \to \mathcal{C}_{Z/Z'}$를
유도한다. 이 사상은 동형이다.
\end{lemma}

\begin{proof}
이는 위의 $Z'$의 구성에서 명백하다.
\end{proof}
















\section{형식적으로 비분기인 사상}
\label{more-morphisms-section-formally-unramified}

\noindent
환 사상 $R \to A$가 {\it 형식적으로 비분기}라 함은, 가환인 실선 도식
$$
\xymatrix{
A \ar[r] \ar@{-->}[rd] & B/I \\
R \ar[r] \ar[u] & B \ar[u]
}
$$
마다, 여기서 $I \subset B$는 제곱이 영인 아이디얼이며, 도식을 가환으로
만드는 점선 화살표가 많아야 하나 존재하는 것임을 상기하자(Algebra,
Definition \ref{algebra-definition-formally-unramified} 참조). 이는 스킴의
사상에 대한 다음 유사 개념의 동기가 된다.

\begin{definition}
\label{more-morphisms-definition-formally-unramified}
$f : X \to S$를 스킴의 사상이라 하자. 임의의 가환인 실선 도식
$$
\xymatrix{
X \ar[d]_f & T \ar[d]^i \ar[l] \\
S & T' \ar[l] \ar@{-->}[lu]
}
$$
이 주어지고 $T \subset T'$이 $S$ 위의 아핀 스킴들의 1차 비후화일 때,
도식을 가환으로 만드는 점선 화살표가 많아야 하나 존재하면 $f$가
{\it 형식적으로 비분기}라고 한다.
\end{definition}

\noindent
먼저 몇 가지 형식적인 보조정리, 즉 대응하는 도식을 그려서 증명할 수 있는
보조정리들을 증명한다.

\begin{lemma}
\label{more-morphisms-lemma-formally-unramified-not-affine}
$f : X \to S$가 형식적으로 비분기인 사상이라고 하자. 임의의 가환인
실선 도식
$$
\xymatrix{
X \ar[d]_f & T \ar[d]^i \ar[l] \\
S & T' \ar[l] \ar@{-->}[lu]
}
$$
이 주어지고 $T \subset T'$이 $S$ 위의 스킴들의 1차 비후화이면, 도식을
가환으로 만드는 점선 화살표가 많아야 하나 존재한다. 달리 말해 Definition
\ref{more-morphisms-definition-formally-unramified}에서 $T$가 아핀이라는 조건을 버려도
된다.
\end{lemma}

\begin{proof}
사상은 아핀 열린집합들에 대한 제한들로 결정되기 때문이다.
\end{proof}

\begin{lemma}
\label{more-morphisms-lemma-composition-formally-unramified}
형식적으로 비분기인 사상들의 합성은 형식적으로 비분기이다.
\end{lemma}

\begin{proof}
이는 형식적이다.
\end{proof}

\begin{lemma}
\label{more-morphisms-lemma-base-change-formally-unramified}
형식적으로 비분기인 사상의 밑변환은 형식적으로 비분기이다.
\end{lemma}

\begin{proof}
이는 형식적이다.
\end{proof}

\begin{lemma}
\label{more-morphisms-lemma-formally-unramified-on-opens}
$f : X \to S$를 스킴의 사상이라 하자. $U \subset X$와
$V \subset S$를 $f(U) \subset V$인 열린집합들이라 하자. $f$가
형식적으로 비분기이면 $f|_U : U \to V$도 형식적으로 비분기이다.
\end{lemma}

\begin{proof}
Definition \ref{more-morphisms-definition-formally-unramified}에서와 같은 실선 도식
$$
\xymatrix{
U \ar[d]_{f|_U} & T \ar[d]^i \ar[l]^a \\
V & T' \ar[l] \ar@{-->}[lu]
}
$$
을 생각하자. $f$가 형식적으로 분기이면 $a'|_T = a$를 만족하는
$S$-사상 $a' : T' \to X$가 많아야 하나 존재한다. 따라서 그러한
$U$로 가는 사상도 많아야 하나 존재함이 명백하다.
% P05-EN-CH37-0020
\end{proof}

\begin{lemma}
\label{more-morphisms-lemma-affine-formally-unramified}
$f : X \to S$를 스킴의 사상이라 하고, $X$와 $S$가 아핀이라고
가정하자. 그러면 $f$가 형식적으로 비분기일 필요충분조건은 환 사상
$\mathcal{O}_S(S) \to \mathcal{O}_X(X)$가 형식적으로 비분기인 것이다.
\end{lemma}

\begin{proof}
환의 범주와 아핀 스킴의 범주 사이의 동치에 의해 정의들(Definition
\ref{more-morphisms-definition-formally-unramified} 및 Algebra, Definition
\ref{algebra-definition-formally-unramified})에서 즉시 따른다. Schemes,
Lemma \ref{schemes-lemma-category-affine-schemes}를 보라.
\end{proof}

\noindent
미분층을 이용한 다음 특징짓기가 있다.

\begin{lemma}
\label{more-morphisms-lemma-formally-unramified-differentials}
$f : X \to S$를 스킴의 사상이라 하자. $f$가 형식적으로 비분기일
필요충분조건은 $\Omega_{X/S} = 0$인 것이다.
\end{lemma}

\begin{proof}
Morphisms, Lemma \ref{morphisms-lemma-differentials-affine}의 증명에
나오는 논증의 일부를 상기한다. $W \subset X \times_S X$를
$\Delta : X \to X \times_S X$가 $W$로의 닫힌 몰입을 유도하도록 하는
열린집합이라 하자. $\mathcal{J} \subset \mathcal{O}_W$를 이 닫힌
몰입의 아이디얼 층이라 하자. $X' \subset W$를 준연접 아이디얼 층
$\mathcal{J}^2$으로 정의되는 닫힌 부분스킴이라 하자. 두 사영
$X \times_S X \to X$가 유도하는 두 사상
$p_1, p_2 : X' \to X$를 생각하자. $p_1$과 $p_2$는
$\Delta : X \to X'$와 합성하면 일치하고, $X \to X'$은 제곱이 영인
아이디얼로 정의되는 닫힌 몰입임에 유의하자.
% P05-EN-CH37-0021
더욱이 다음 짧은 완전열이 있고
$$
0 \to \mathcal{J}/\mathcal{J}^2 \to \mathcal{O}_{X'} \to \mathcal{O}_X \to 0
$$
$\Omega_{X/S} = \mathcal{J}/\mathcal{J}^2$이다. 또한
$\mathcal{J}/\mathcal{J}^2$는 $\mathcal{O}_X$의 국소 절단 $f$에 대한
국소 절단 $p_1^\sharp(f) - p_2^\sharp(f)$들로 생성된다.

\medskip\noindent
$f : X \to S$가 형식적으로 비분기라고 가정하자. 가정에 의해 이는
임의의 아핀 열린집합 $T' \subset X'$에 제한하면 $p_1 = p_2$라는
뜻이다. 따라서 $p_1 = p_2$이다. 위에서 말한 것으로부터
$\Omega_{X/S} = \mathcal{J}/\mathcal{J}^2 = 0$이라고 결론짓는다.

\medskip\noindent
거꾸로 $\Omega_{X/S} = 0$이라고 가정하자. 그러면 $X' = X$이다.
Definition \ref{more-morphisms-definition-formally-unramified}의 도식에서 점선 화살표로
들어맞는 임의의 두 사상 $f'_1, f'_2 : T' \to X$를 택하자. 이들은
사상 $(f'_1, f'_2) : T' \to X \times_S X$를 준다.
$f'_1|_T = f'_2|_T$이고 $|T| =|T'|$이므로 $(f'_1, f'_2)$에 의한
$T'$의 상은 위에서 택한 열린집합 $W$에 포함된다. 또한
$(f'_1, f'_2)(T) \subset \Delta(X)$이고 $T$는 $T'$ 안에서 제곱이
영인 아이디얼로 정의되므로, $(f'_1, f'_2)$는 $X'$을 통하여
인수분해된다. $X' = X$이므로 원하는 대로 $f_1' = f'_2$라고
결론짓는다.
\end{proof}

\begin{lemma}
\label{more-morphisms-lemma-formally-unramified}
$f : X \to Y$를 스킴의 사상이라 하자. 다음은 서로 동치이다.
\begin{enumerate}
\item $f$는 형식적으로 비분기이다.
\item 모든 $x \in X$에 대해 $f(U) \subset V$이고
$f|_U : U \to V$가 형식적으로 비분기인 열린집합
$x \in U \subset X$와 $f(x) \in V \subset Y$가 존재한다.
\item $f(U) \subset V$인 모든 아핀 열린집합 쌍 $U \subset X$와
$V \subset Y$에 대해 환 사상
$\mathcal{O}_Y(V) \to \mathcal{O}_X(U)$가 형식적으로 비분기이다.
\item 아핀 열린 덮개 $Y = \bigcup V_j$가 존재하고, 각 $j$에 대해
아핀 열린 덮개 $f^{-1}(V_j) = \bigcup U_{ji}$가 존재하여 모든 $j$와
$i$에 대해 환 사상 $\mathcal{O}_Y(V) \to \mathcal{O}_X(U)$가
형식적으로 비분기이다.
% P05-EN-CH37-0022
\end{enumerate}
\end{lemma}

\begin{proof}
Lemmas \ref{more-morphisms-lemma-formally-unramified-on-opens}와
\ref{more-morphisms-lemma-affine-formally-unramified}를 결합하면
(1) $\Rightarrow$ (3)을 얻는다. (3) $\Rightarrow$ (4)는 즉시 따른다.
Lemma \ref{more-morphisms-lemma-affine-formally-unramified}에 의해
(4) $\Rightarrow$ (2)이다. (2)가 성립하면 Lemma
\ref{more-morphisms-lemma-formally-unramified-differentials}에 의해 모든 점의 근방에서
$\Omega_{X/S}$가 소멸하고(여기서는 Morphisms, Lemma
\ref{morphisms-lemma-differentials-restrict-open}도 사용한다), 다시
Lemma \ref{more-morphisms-lemma-formally-unramified-differentials}에 의해 $f$가
형식적으로 비분기임을 알 수 있다.
% P05-EN-CH37-0023
\end{proof}

\begin{lemma}
\label{more-morphisms-lemma-unramified-formally-unramified}
\begin{slogan}
비분기 사상은 형식적으로 비분기이면서 국소 유한형인 사상과 같다.
% P05-EN-CH37-0024
\end{slogan}
$f : X \to S$를 스킴의 사상이라 하자. 다음은 서로 동치이다.
\begin{enumerate}
\item 사상 $f$는 비분기이다(각각 G-비분기이다).
\item 사상 $f$는 국소 유한형(각각 국소 유한표현)이고 형식적으로
비분기이다.
\end{enumerate}
\end{lemma}

\begin{proof}
Lemma \ref{more-morphisms-lemma-formally-unramified-differentials}와 Morphisms, Lemma
\ref{morphisms-lemma-unramified-omega-zero}을 사용하라.
\end{proof}









\section{보편 1차 비후화}
\label{more-morphisms-section-universal-thickening}

\noindent
$h : Z \to X$를 스킴의 사상이라 하자. $X$ 위의 $Z$의
{\it 보편 1차 비후화}란 다음 성질을 갖는 $X$ 위의 1차 비후화
$Z \subset Z'$이다. $X$ 위의 임의의 1차 비후화 $T \subset T'$과
가환인 실선 도식
$$
\xymatrix{
& Z \ar[ld] & & T \ar[rd] \ar[ll]^a \\
Z' \ar[rrd] & & & & T' \ar@{..>}[llll]_{a'} \ar[lld]^b \\
 & & X
}
$$
이 주어지면 도식을 가환으로 만드는 점선 화살표가 유일하게 존재한다.
이 상황에서 $(a, a') : (T \subset T') \to (Z \subset Z')$는 $X$ 위의
비후화의 사상임에 유의하자. 따라서 보편 1차 비후화가 존재하면 그것은
유일한 동형을 제외하고 유일하다. 일반적으로 보편 1차 비후화는 존재하지
않지만, $h$가 형식적으로 비분기이면 존재한다.

\begin{lemma}
\label{more-morphisms-lemma-universal-thickening}
$h : Z \to X$를 형식적으로 비분기인 스킴의 사상이라 하자. $X$ 위의
$Z$의 보편 1차 비후화 $Z \subset Z'$이 존재한다.
\end{lemma}

\begin{proof}
이 증명에서는 $Z \subset Z'$이 보조정리의 조건을 만족하면 그것을
$X$ 위의 $Z$의 보편 1차 비후화라고 하겠다. 먼저 Algebra, Lemma
\ref{algebra-lemma-universal-thickening}인 아핀 경우에서 시작하여 붙이기로
$X$ 위의 보편 1차 비후화 $Z \subset Z'$을 구성할 것이다. 몇 가지
일반적인 언급으로 시작한다.

\medskip\noindent
$X$ 위의 $Z$의 보편 1차 비후화가 존재하면, 그것은 유일한 동형을
제외하고 유일하다. 또한 $V \subset Z$와 $U \subset X$를
$h(V) \subset U$인 열린 부분스킴들이라 하자. $Z \subset Z'$을
$X$ 위의 $Z$의 보편 1차 비후화라 하고, $V' \subset Z'$을
$V = Z \cap V'$인 열린 부분스킴이라 하자. 그러면 $V \subset V'$이
$U$ 위의 $V$의 보편 1차 비후화라고 주장한다. 실제로 임의의 도식
$$
\xymatrix{
V \ar[d]_h & T \ar[l]^a \ar[d] \\
U & T' \ar[l]_b
}
$$
이 주어지고 $T \subset T'$이 $U$ 위의 1차 비후화라고 하자. $Z'$의
보편 성질에 의해
$(a, a') : (T \subset T') \to (Z \subset Z')$를 얻는다. 그런데 바탕
위상공간들의 등식 $|T| = |T'|$이 있으므로 $a'(T') \subset V'$임을
알 수 있다. 따라서 $(a, a')$를 $U$ 위의 비후화의 사상
$(a, a') : (T \subset T') \to (V \subset V')$로 생각할 수 있다.
유일성도 명백하다. 완전히 같은 방식으로, $h(Z) \subset U$이고
$Z \subset Z'$이 $U$ 위의 보편 1차 비후화이면 $Z \subset Z'$이
$X$ 위의 보편 1차 비후화임을 증명한다.

\medskip\noindent
아핀 조각들을 붙이기 전에 $Z$와 $X$가 아핀일 때 보조정리가 성립함을
보이자. $X = \Spec(R)$이고 $Z = \Spec(S)$라 하자. Algebra, Lemma
\ref{algebra-lemma-universal-thickening}에 의해 다음과 같은 도식에 대해
보조정리의 보편 성질을 갖는 $X$ 위의 1차 비후화 $Z \subset Z'$이
존재한다.
$$
\xymatrix{
Z \ar[d]_h & T \ar[l]^a \ar[d] \\
X & T' \ar[l]_b
}
$$
여기서 $T, T'$은 아핀이다. 일반적인 도식이 주어지면 아핀 열린 덮개
$T' = \bigcup T'_i$를 택할 수 있고, $a'_i|_{T_i} = a|_{T_i}$를
만족하는 $X$ 위의 사상 $a'_i : T'_i \to Z'$를 얻는다. 유일성에 의해
$a'_i$와 $a'_j$는 $T'_i \cap T'_j$의 임의의 아핀 열린집합 위에서
일치한다. 따라서 사상들 $a'_i$는 원하는 $X$ 위의 전역 사상
$a' : T' \to Z'$로 붙는다. 그러므로 $X$와 $Z$가 아핀이면 보조정리가
성립한다.

\medskip\noindent
각 $Z_i$가 $X$의 어떤 아핀 열린집합 $U_i$로 가도록 하는 아핀 열린
덮개 $Z = \bigcup Z_i$를 택하자. Lemma
\ref{more-morphisms-lemma-formally-unramified-on-opens}에 의해 사상 $Z_i \to U_i$는
형식적으로 비분기이다. 따라서 아핀 경우에 의해 $U_i$ 위의 보편 1차
비후화 $Z_i \subset Z_i'$을 얻는다. 위의 일반적인 언급에 의해
$Z_i \subset Z_i'$은 $X$ 위의 $Z_i$의 보편 1차 비후화이기도 하다.
$Z'_{i, j} \subset Z'_i$를
$Z_i \cap Z_j = Z'_{i, j} \cap Z_i$인 열린 부분스킴이라 하자. 일반적인
언급에 의해 $Z'_{i, j}$와 $Z'_{j, i}$는 모두 $X$ 위의
$Z_i \cap Z_j$의 보편 1차 비후화이다. 따라서 첫째 일반적인 언급에
의해 $Z_i \cap Z_j$ 위에서 항등사상을 유도하는 표준 동형
$\varphi_{ij} : Z'_{i, j} \to Z'_{j, i}$가 존재한다. 이 사상들이
Schemes, Section \ref{schemes-section-glueing-schemes}의 코사이클 조건을
만족한다고 주장한다. (검증은 생략한다. 힌트: 위에서 말한 것에 의해
유일한 동형을 제외하고 결정되는 $Z_i \cap Z_j \cap Z_k$의 보편 1차
비후화가 $Z'_{i, j} \cap Z'_{i, k}$라는 사실을 사용하라.) 따라서
Schemes, Section \ref{schemes-section-glueing-schemes}의 결과를 사용하여
$Z_i = Z'_i \cap Z$인 열린 부분스킴 $Z'_i \subset Z'$이 $X$ 위의
$Z_i$의 보편 1차 비후화라는 성질을 갖는 $X$ 위의 1차 비후화
$Z \subset Z'$을 얻는다.
% P05-EN-CH37-0025

\medskip\noindent
이로부터 형식적으로 $Z'$이 $X$ 위의 $Z$의 보편 1차 비후화임이 따른다.
실제로 다음 도식에서 $a(T)$가 어떤 $Z_i$에 포함될 때에는 보편 성질이
성립한다.
$$
\xymatrix{
Z \ar[d]_h & T \ar[l]^a \ar[d] \\
X & T' \ar[l]_b
}
$$
일반적인 도식이 주어지면 $a(T_i) \subset Z_i$가 되도록 하는 열린 덮개
$T' = \bigcup T'_i$를 택할 수 있다. $a'_i|_{T_i} = a|_{T_i}$를
만족하는 $X$ 위의 사상 $a'_i : T'_i \to Z'$를 얻는다.
$a'_i|_{T'_i \cap T'_j}$와 $a'_j|_{T'_i \cap T'_j}$는 모두 $X$ 위의
$Z_i \cap Z_j$의 보편 1차 비후화 $Z'_i \cap Z'_j$로 사상하는 문제의
해이므로, $a'_i$와 $a'_j$는 $T'_i \cap T'_j$ 위에서 반드시 일치한다.
따라서 사상들 $a'_i$는 원하는 $X$ 위의 전역 사상
$a' : T' \to Z'$로 붙는다. 이것으로 증명이 끝난다.
\end{proof}

\begin{definition}
\label{more-morphisms-definition-universal-thickening}
$h : Z \to X$를 형식적으로 비분기인 스킴의 사상이라 하자.
\begin{enumerate}
\item $X$ 위의 $Z$의 {\it 보편 1차 비후화}는 Lemma
\ref{more-morphisms-lemma-universal-thickening}에서 구성한 비후화 $Z \subset Z'$이다.
\item {$X$ 위의 $Z$의 conormal 층}은 $X$ 위의 보편 1차 비후화
$Z'$ 안의 $Z$의 conormal 층이다.
\end{enumerate}
이 상황에서 conormal 층을 흔히 $\mathcal{C}_{Z/X}$로 나타낸다.
\end{definition}

\noindent
따라서 $Z$ 위에 다음 층들의 짧은 완전열이 있음을 알 수 있다.
$$
0 \to \mathcal{C}_{Z/X} \to \mathcal{O}_{Z'} \to \mathcal{O}_Z \to 0
$$
다음 보조정리는 이 정의가 $Z \to X$가 몰입일 때의 정의와 충돌하지
않음을 증명한다.

\begin{lemma}
\label{more-morphisms-lemma-immersion-universal-thickening}
$i : Z \to X$를 스킴의 몰입이라 하자. 그러면
\begin{enumerate}
\item $i$는 형식적으로 비분기이다.
\item $X$ 위의 $Z$의 보편 1차 비후화는 Definition
\ref{more-morphisms-definition-first-order-infinitesimal-neighbourhood}의 $X$ 안의
$Z$의 1차 무한소 근방이다.
\item Morphisms, Definition \ref{morphisms-definition-conormal-sheaf}의
의미에서 $i$의 conormal 층은 Definition
\ref{more-morphisms-definition-universal-thickening}의 의미에서 $i$의 conormal 층과
일치한다.
\end{enumerate}
\end{lemma}

\begin{proof}
Morphisms, Lemmas \ref{morphisms-lemma-open-immersion-unramified}와
\ref{morphisms-lemma-closed-immersion-unramified}에 의해 몰입은
비분기이고, 따라서 Lemma \ref{more-morphisms-lemma-unramified-formally-unramified}에
의해 형식적으로 비분기이다. 나머지 주장들은 Lemmas
\ref{more-morphisms-lemma-first-order-infinitesimal-neighbourhood}와
\ref{more-morphisms-lemma-infinitesimal-neighbourhood-conormal} 및 정의들을 결합하면
따른다.
\end{proof}

\begin{lemma}
\label{more-morphisms-lemma-universal-thickening-unramified}
$Z \to X$를 형식적으로 비분기인 스킴의 사상이라 하자. 그러면 보편
1차 비후화 $Z'$은 $X$ 위에서 형식적으로 비분기이다.
\end{lemma}

\begin{proof}
두 가지 증명이 있다. 첫째 증명은 아핀 국소적으로 작업하고 Algebra,
Lemma \ref{algebra-lemma-differentials-universal-thickening}을 적용하여
$\Omega_{Z'/X} = 0$임을 보이는 것이다. 그러면 Lemma
\ref{more-morphisms-lemma-formally-unramified-differentials}에서 원하는 결론이 따른다.
둘째 증명은 다음의 직접적인 논증이다.

\medskip\noindent
$T \subset T'$을 1차 비후화라 하자. 가환 도식
$$
\xymatrix{
Z' \ar[d] & T \ar[l]^c \ar[d] \\
X & T' \ar[l] \ar[lu]^{a, b}
}
$$
을 생각하고, 이 도식에 들어맞는 두 사상 $a, b : T' \to Z'$를
생각하자. $T_0 = c^{-1}(Z) \subset T$ 및
$T'_a = a^{-1}(Z)$로 놓자(스킴론적으로). $Z'$이 $Z$의 1차
비후화이므로 $T'$은 $T'_a$의 1차 비후화이다. 또한 $c = a|_T$이므로
$T_0 = T \cap T'_a$이다(스킴론적으로). $T'$이 $T$의 1차
비후화이므로 $T'_a$는 $T_0$의 1차 비후화이다. 이제
$a|_{T'_a}$와 $b|_{T'_a}$는 $X$ 위에서 $T'_a$에서 $Z'$으로 가는
사상들이며, $T_0$ 위에서는 $Z$로 가는 사상으로서 일치한다. 따라서
$Z'$의 보편 성질에 의해 $a|_{T'_a} = b|_{T'_a}$이다. 그러므로 $a$와
$b$는 $T'_a$의 1차 비후화 $T'$에서 출발하는 사상들이고, 그 제한들은
$T'_a$ 위에서 $Z$로 가는 사상으로서 일치한다.
% P05-EN-CH37-0026
$Z'$의 보편 성질을 한 번 더 사용하면 $a = b$라고 결론짓는다. 달리
말해 형식적으로 비분기인 사상의 정의 성질이 원하는 대로 $Z' \to X$에
대해 성립한다.
\end{proof}

\begin{lemma}
\label{more-morphisms-lemma-universal-thickening-functorial}
스킴의 가환 도식
$$
\xymatrix{
Z \ar[r]_h \ar[d]_f & X \ar[d]^g \\
W \ar[r]^{h'} & Y
}
$$
을 생각하고 $h$와 $h'$이 형식적으로 비분기라고 하자.
$Z \subset Z'$을 $X$ 위의 $Z$의 보편 1차 비후화라 하고,
$W \subset W'$을 $Y$ 위의 $W$의 보편 1차 비후화라 하자. 다음 가환
도식에 들어맞는 $Y$ 위의 비후화의 표준 사상
$(f, f') : (Z, Z') \to (W, W')$가 존재한다.
% P05-EN-CH37-0027
$$
\xymatrix{
& & & Z' \ar[ld] \ar[d]^{f'} \\
Z \ar[rr] \ar[d]_f \ar[rrru] & & X \ar[d] & W' \ar[ld] \\
W \ar[rrru]|!{[rr];[rruu]}\hole \ar[rr] & & Y
}
$$
특히 비후화의 사상 $(f, f')$는 conormal 층의 사상
$f^*\mathcal{C}_{W/Y} \to \mathcal{C}_{Z/X}$를 유도한다.
\end{lemma}

\begin{proof}
첫째 주장은 $W'$의 보편 성질에서 명백하다. 유도된 conormal 층의 사상은
Morphisms, Lemma \ref{morphisms-lemma-conormal-functorial}의 사상을
$(Z \subset Z') \to (W \subset W')$에 적용한 것이다.
\end{proof}

\begin{lemma}
\label{more-morphisms-lemma-universal-thickening-fibre-product}
다음 도식이
$$
\xymatrix{
Z \ar[r]_h \ar[d]_f & X \ar[d]^g \\
W \ar[r]^{h'} & Y
}
$$
스킴의 범주에서 올곱 도식이고 $h'$이 형식적으로 비분기라고 하자. 그러면
$h$는 형식적으로 비분기이고, $W \subset W'$이 $Y$ 위의 $W$의 보편
1차 비후화이면
$Z = X \times_Y W \subset X \times_Y W'$은 $X$ 위의 $Z$의 보편
1차 비후화이다. 특히 Lemma
\ref{more-morphisms-lemma-universal-thickening-functorial}의 표준 사상
$f^*\mathcal{C}_{W/Y} \to \mathcal{C}_{Z/X}$는 전사이다.
\end{lemma}

\begin{proof}
Lemma \ref{more-morphisms-lemma-base-change-formally-unramified}에 의해 사상 $h$는
형식적으로 비분기이다. $X \times_Y W'$이 1차 비후화임은 명백하다.
$W'$이 보편 성질을 가지므로(올곱의 사상 성질에 의해) 이것도 보편
성질을 가짐을 곧바로 확인할 수 있다. 이로부터 conormal 층의 사상이
전사임이 따르는 이유는 Morphisms, Lemma
\ref{morphisms-lemma-conormal-functorial-flat}을 보라.
\end{proof}

\begin{lemma}
\label{more-morphisms-lemma-universal-thickening-fibre-product-flat}
다음 도식이
$$
\xymatrix{
Z \ar[r]_h \ar[d]_f & X \ar[d]^g \\
W \ar[r]^{h'} & Y
}
$$
스킴의 범주에서 올곱 도식이고, $h'$은 형식적으로 비분기이며 $g$는
평탄하다고 하자. 이 경우 대응하는 보편 1차 비후화들의 사상
$Z' \to W'$은 평탄하고,
$f^*\mathcal{C}_{W/Y} \to \mathcal{C}_{Z/X}$는 동형이다.
\end{lemma}

\begin{proof}
평탄성은 밑변환에 의해 보존된다. Morphisms, Lemma
\ref{morphisms-lemma-base-change-flat}을 보라. 따라서 첫째 명제는 Lemma
\ref{more-morphisms-lemma-universal-thickening-fibre-product}의 $W'$의 서술에서 따른다.
% P05-EN-CH37-0028
$X \times_Y W'$이 1차 비후화임은 명백하다. $W'$이 보편 성질을
가지므로(올곱의 사상 성질에 의해) 이것도 보편 성질을 가짐을 곧바로
확인할 수 있다. 이로부터 conormal 층의 사상이 동형임이 따르는 이유는
Morphisms, Lemma \ref{morphisms-lemma-conormal-functorial-flat}을 보라.
\end{proof}

\begin{lemma}
\label{more-morphisms-lemma-universal-thickening-localize}
보편 1차 비후화를 취하는 것은 열린부분을 취하는 것과 가환한다. 더
정확히, $h : Z \to X$를 형식적으로 비분기인 스킴의 사상이라 하자.
$V \subset Z$, $U \subset X$를 $h(V) \subset U$인 열린집합들이라
하고, $Z'$을 $X$ 위의 $Z$의 보편 1차 비후화라 하자. 그러면
$h|_V : V \to U$는 형식적으로 비분기이고, $U$ 위의 $V$의 보편 1차
비후화는 $V = Z \cap V'$인 열린 부분스킴 $V' \subset Z'$이다. 특히
$\mathcal{C}_{Z/X}|_V = \mathcal{C}_{V/U}$이다.
\end{lemma}

\begin{proof}
첫째 명제는 Lemma \ref{more-morphisms-lemma-formally-unramified-on-opens}이다. 보편
비후화의 호환성은 Lemma \ref{more-morphisms-lemma-universal-thickening}의 증명에서,
또는 Algebra, Lemma \ref{algebra-lemma-universal-thickening-localize}에서,
또는 Lemma \ref{more-morphisms-lemma-universal-thickening-fibre-product-flat}에서
도출할 수 있다.
\end{proof}

\begin{lemma}
\label{more-morphisms-lemma-differentials-universally-unramified}
$h : Z \to X$를 $S$ 위에서 형식적으로 비분기인 스킴의 사상이라 하자.
$Z \subset Z'$을 구조 사상 $h' : Z' \to X$를 갖는 $X$ 위의 $Z$의
보편 1차 비후화라 하자. 표준 사상
$$
c_{h'} : (h')^*\Omega_{X/S} \longrightarrow \Omega_{Z'/S}
$$
은 동형 $h^*\Omega_{X/S} \to \Omega_{Z'/S} \otimes \mathcal{O}_Z$를
유도한다.
\end{lemma}

\begin{proof}
사상 $c_{h'}$은 Morphisms, Lemma
\ref{morphisms-lemma-functoriality-differentials}에서 정의한 사상이다.
$i : Z \to Z'$이 주어진 닫힌 몰입이면 $i^*c_{h'}$은 사상
$h^*\Omega_{X/S} \to \Omega_{Z'/S} \otimes \mathcal{O}_Z$이다. 국소화에
의해 이것이 동형인지 확인하는 문제는 아핀 경우로 환원된다. Lemma
\ref{more-morphisms-lemma-universal-thickening-localize}와 Morphisms, Lemma
\ref{morphisms-lemma-differentials-restrict-open}을 보라. 이 경우 결과는
Algebra, Lemma \ref{algebra-lemma-differentials-universal-thickening}이다.
\end{proof}

\begin{lemma}
\label{more-morphisms-lemma-universally-unramified-differentials-sequence}
$h : Z \to X$를 $S$ 위에서 형식적으로 비분기인 스킴의 사상이라 하자.
다음 표준 완전열이 있다.
$$
\mathcal{C}_{Z/X} \to h^*\Omega_{X/S} \to \Omega_{Z/S} \to 0.
$$
첫째 화살표는 $\text{d}_{Z'/S}$에 의해 유도되며, 여기서 $Z'$은 $X$
위의 $Z$의 보편 1차 근방이다.
\end{lemma}

\begin{proof}
다음 표준 완전열이 있음을 안다.
$$
\mathcal{C}_{Z/Z'} \to
\Omega_{Z'/S} \otimes \mathcal{O}_Z \to
\Omega_{Z/S} \to 0.
$$
Morphisms, Lemma \ref{morphisms-lemma-differentials-relative-immersion}을
보라. 따라서 Lemma \ref{more-morphisms-lemma-differentials-universally-unramified}를
적용하면 결과가 따른다.
\end{proof}

\begin{lemma}
\label{more-morphisms-lemma-two-unramified-morphisms}
다음 스킴의 도식이
$$
\xymatrix{
Z \ar[r]_i \ar[rd]_j & X \ar[d] \\
& Y
}
$$
가환이고 $i$와 $j$가 형식적으로 비분기라고 하자. 그러면 다음 표준
완전열이 있다.
$$
\mathcal{C}_{Z/Y} \to
\mathcal{C}_{Z/X} \to
i^*\Omega_{X/Y} \to 0
$$
여기서 첫째 화살표는 Lemma
\ref{more-morphisms-lemma-universal-thickening-functorial}에서, 둘째 화살표는 Lemma
\ref{more-morphisms-lemma-universally-unramified-differentials-sequence}에서 온다.
\end{lemma}

\begin{proof}
$X$ 위의 $Z$의 보편 1차 비후화를 $Z \to Z'$로 나타내고, $Y$ 위의
$Z$의 보편 1차 비후화를 $Z \to Z''$로 나타내자.
% P05-EN-CH37-0029
Lemma \ref{more-morphisms-lemma-universally-unramified-differentials-sequence}에 의해
다음 가환 도식을 주는 표준 사상 $Z' \to Z''$가 있다.
% P05-EN-CH37-0030
$$
\xymatrix{
Z \ar[r]_{i'} \ar[rd]_{j'} & Z' \ar[r] \ar[d] & X \ar[d] \\
& Z'' \ar[r] & Y
}
$$
왼쪽 삼각형에 Morphisms, Lemma
\ref{morphisms-lemma-two-immersions}를 적용하여 완전열
$$
\mathcal{C}_{Z/Z''} \to
\mathcal{C}_{Z/Z'} \to
(i')^*\Omega_{Z'/Z''} \to 0
$$
을 얻는다. $Z''$은 $Y$ 위에서 형식적으로 비분기이므로(Lemma
\ref{more-morphisms-lemma-universal-thickening-unramified} 참조), Lemma
\ref{more-morphisms-lemma-formally-unramified-differentials}와 Morphisms, Lemma
\ref{morphisms-lemma-triangle-differentials}를 결합하여
$\Omega_{Z'/Z''} = \Omega_{Z/Y}$를 얻는다.
% P05-EN-CH37-0031
그런 다음 Lemma \ref{more-morphisms-lemma-differentials-universally-unramified}에 의해
$(i')^*\Omega_{Z'/Y} = i^*\Omega_{X/Y}$를 얻는다.
\end{proof}

\begin{lemma}
\label{more-morphisms-lemma-transitivity-conormal}
$Z \to Y \to X$를 형식적으로 비분기인 스킴의 사상들이라 하자.
\begin{enumerate}
\item $Z \subset Z'$이 $X$ 위의 $Z$의 보편 1차 비후화이고
$Y \subset Y'$이 $X$ 위의 $Y$의 보편 1차 비후화이면, 사상
$Z' \to Y'$가 존재하고 $Y \times_{Y'} Z'$은 $Y$ 위의 $Z$의 보편 1차
비후화이다.
\item 다음 표준 완전열이 있다.
$$
i^*\mathcal{C}_{Y/X} \to
\mathcal{C}_{Z/X} \to
\mathcal{C}_{Z/Y} \to 0
$$
여기서 사상들은 Lemma \ref{more-morphisms-lemma-universal-thickening-functorial}에서
오고, $i : Z \to Y$는 첫째 사상이다.
\end{enumerate}
\end{lemma}

\begin{proof}
(1)의 사상 $h : Z' \to Y'$은 Lemma
\ref{more-morphisms-lemma-universal-thickening-functorial}에서 온다.
$Y \times_{Y'} Z'$이 $Y$ 위의 $Z$의 보편 1차 비후화라는 주장은
$Z'$과 $Y'$의 보편 성질에서 명백하다. Morphisms, Lemma
\ref{morphisms-lemma-transitivity-conormal}에 의해 완전열
$$
(i')^*\mathcal{C}_{Y \times_{Y'} Z'/Z'} \to
\mathcal{C}_{Z/Z'} \to
\mathcal{C}_{Z/Y \times_{Y'} Z'} \to 0
$$
이 있고, 여기서 $i' : Z \to Y \times_{Y'} Z'$은 주어진 사상이다.
Morphisms, Lemma \ref{morphisms-lemma-conormal-functorial-flat}에 의해
전사
$h^*\mathcal{C}_{Y/Y'} \to \mathcal{C}_{Y \times_{Y'} Z'/Z'}$가
존재한다. 이를 등식들
$\mathcal{C}_{Y/Y'} = \mathcal{C}_{Y/X}$,
$\mathcal{C}_{Z/Z'} = \mathcal{C}_{Z/X}$, 그리고
$\mathcal{C}_{Z/Y \times_{Y'} Z'} = \mathcal{C}_{Z/Y}$
과 결합하면 보조정리가 증명된다.
\end{proof}











\section{형식적으로 에탈인 사상}
\label{more-morphisms-section-formally-etale}

\noindent
환 사상 $R \to A$가 {\it 형식적으로 에탈}이라 함은, 가환인 실선 도식
$$
\xymatrix{
A \ar[r] \ar@{-->}[rd] & B/I \\
R \ar[r] \ar[u] & B \ar[u]
}
$$
마다, 여기서 $I \subset B$는 제곱이 영인 아이디얼이며, 도식을 가환으로
만드는 점선 화살표가 정확히 하나 존재하는 것임을 상기하자(Algebra,
Definition \ref{algebra-definition-formally-etale} 참조). 이는 스킴의
사상에 대한 다음 유사 개념의 동기가 된다.

\begin{definition}
\label{more-morphisms-definition-formally-etale}
$f : X \to S$를 스킴의 사상이라 하자. 임의의 가환인 실선 도식
$$
\xymatrix{
X \ar[d]_f & T \ar[d]^i \ar[l] \\
S & T' \ar[l] \ar@{-->}[lu]
}
$$
이 주어지고 $T \subset T'$이 $S$ 위의 아핀 스킴들의 1차 비후화일 때,
도식을 가환으로 만드는 점선 화살표가 정확히 하나 존재하면 $f$가
{\it 형식적으로 에탈}이라고 한다.
\end{definition}

\noindent
형식적으로 에탈인 사상이 형식적으로 비분기임은 명백하다. 따라서
$f : X \to S$가 형식적으로 에탈이면 $\Omega_{X/S}$는 영이다. Lemma
\ref{more-morphisms-lemma-formally-unramified-differentials}를 보라.

\begin{lemma}
\label{more-morphisms-lemma-formally-etale-not-affine}
$f : X \to S$가 형식적으로 에탈인 사상이라고 하자. 임의의 가환인
실선 도식
$$
\xymatrix{
X \ar[d]_f & T \ar[d]^i \ar[l] \\
S & T' \ar[l] \ar@{-->}[lu]
}
$$
이 주어지고 $T \subset T'$이 $S$ 위의 스킴들의 1차 비후화이면, 도식을
가환으로 만드는 점선 화살표가 정확히 하나 존재한다. 달리 말해 Definition
\ref{more-morphisms-definition-formally-etale}에서 $T$가 아핀이라는 조건을 버려도 된다.
\end{lemma}

\begin{proof}
$T' = \bigcup T'_i$를 아핀 열린 덮개라 하고 $T_i = T \cap T'_i$로
놓자. 그러면 도식에 들어맞는 사상 $a'_i : T'_i \to X$를 얻는다.
유일성에 의해 $a'_i$와 $a'_j$는 $T'_i \cap T'_j$의 임의의 아핀 열린
부분스킴 위에서 일치한다. 따라서 $a'_i$와 $a'_j$는
$T'_i \cap T'_j$ 위에서 일치한다. 그러므로 사상들 $a'_i$는 전역 사상
$a' : T' \to X$로 붙는다. $a'$의 유일성은 Lemma
\ref{more-morphisms-lemma-formally-unramified-not-affine}에서 보았다.
\end{proof}

\begin{lemma}
\label{more-morphisms-lemma-composition-formally-etale}
형식적으로 에탈인 사상들의 합성은 형식적으로 에탈이다.
\end{lemma}

\begin{proof}
이는 형식적이다.
\end{proof}

\begin{lemma}
\label{more-morphisms-lemma-base-change-formally-etale}
형식적으로 에탈인 사상의 밑변환은 형식적으로 에탈이다.
\end{lemma}

\begin{proof}
이는 형식적이다.
\end{proof}

\begin{lemma}
\label{more-morphisms-lemma-formally-etale-on-opens}
$f : X \to S$를 스킴의 사상이라 하자. $U \subset X$와
$V \subset S$를 $f(U) \subset V$인 열린 부분스킴들이라 하자. $f$가
형식적으로 에탈이면 $f|_U : U \to V$도 형식적으로 에탈이다.
\end{lemma}

\begin{proof}
Definition \ref{more-morphisms-definition-formally-etale}에서와 같은 실선 도식
$$
\xymatrix{
U \ar[d]_{f|_U} & T \ar[d]^i \ar[l]^a \\
V & T' \ar[l] \ar@{-->}[lu]
}
$$
을 생각하자. $f$가 형식적으로 분기이면 $a'|_T = a$를 만족하는
$S$-사상 $a' : T' \to X$가 정확히 하나 존재한다.
% P05-EN-CH37-0032
$|T'| = |T|$이므로 $a'(T') \subset U$라고 결론짓고, 이로써 $T'$에서
$U$로 가는 유일한 사상을 얻는다.
\end{proof}

\begin{lemma}
\label{more-morphisms-lemma-characterize-formally-etale}
$f : X \to S$를 스킴의 사상이라 하자. 다음은 서로 동치이다.
\begin{enumerate}
\item $f$는 형식적으로 에탈이다.
\item $f$는 형식적으로 비분기이고 $S$ 위의 $X$의 보편 1차 비후화는
$X$와 같다.
\item $f$는 형식적으로 비분기이고 $\mathcal{C}_{X/S} = 0$이다.
\item $\Omega_{X/S} = 0$이고 $\mathcal{C}_{X/S} = 0$이다.
\end{enumerate}
\end{lemma}

\begin{proof}
실제로 마지막 주장이 뜻을 갖는 것은 $\Omega_{X/S} = 0$으로부터 Lemma
\ref{more-morphisms-lemma-formally-unramified-differentials}와 Definition
\ref{more-morphisms-definition-universal-thickening}을 통해 $\mathcal{C}_{X/S}$가
정의되기 때문이다. 이로부터 (3)과 (4)가 동치임도 명백해진다.
% P05-EN-CH37-0033

\medskip\noindent
가정 (1), (2), (3) 가운데 어느 것이든 $f$가 형식적으로 비분기임을
함의한다. 따라서 $f$가 형식적으로 비분기라고 가정할 수 있다.
% P05-EN-CH37-0034
(1), (2), (3)의 동치는 $S$ 위의 $X$의 보편 1차 비후화 $X'$의 보편
성질과 $X = X' \Leftrightarrow \mathcal{C}_{X/S} = 0$이라는 사실에서
따른다. 왜냐하면 결국 정의에 의해
$\mathcal{C}_{X/S} = \mathcal{C}_{X/X'}$는 $X'$ 안의 $X$의 아이디얼
층이기 때문이다.
% P05-EN-CH37-0035
\end{proof}

\begin{lemma}
\label{more-morphisms-lemma-unramified-flat-formally-etale}
비분기이고 평탄한 사상은 형식적으로 에탈이다.
\end{lemma}

\begin{proof}
$X \to S$가 비분기이고 평탄하다고 하자. 그러면
$\Delta : X \to X \times_S X$는 열린 몰입이다. Morphisms, Lemma
\ref{morphisms-lemma-diagonal-unramified-morphism}을 보라.
$\mathcal{C}_{X/S}$가 영임을 보여야 한다. 두 사영
$p, q : X \times_S X \to X$를 생각하자. $f$가 형식적으로
비분기이므로(Lemma \ref{more-morphisms-lemma-unramified-formally-unramified} 참조),
$q$는 형식적으로 비분기이다(Lemma
\ref{more-morphisms-lemma-base-change-formally-unramified} 참조). $f$가 평탄하므로
$p$는 평탄하다. Morphisms, Lemma \ref{morphisms-lemma-base-change-flat}을
보라. 따라서 Lemma \ref{more-morphisms-lemma-universal-thickening-fibre-product-flat}에
의해 $p^*\mathcal{C}_{X/S} = \mathcal{C}_q$이며, 여기서
$\mathcal{C}_q$는 형식적으로 비분기인 사상
$q : X \times_S X \to X$의 conormal 층을 나타낸다. 그런데
$\Delta(X) \subset X \times_S X$는 $q$에 의해 $X$로 동형사상되는 열린
부분스킴이다. 따라서 Lemma \ref{more-morphisms-lemma-universal-thickening-localize}에
의해 $\mathcal{C}_q|_{\Delta(X)} = \mathcal{C}_{X/X} = 0$이다. 달리
말해 항등사상에 의한 $\mathcal{C}_{X/S}$의 $X$로의 역상은 영이다. 즉
$\mathcal{C}_{X/S} = 0$이다.
\end{proof}

\begin{lemma}
\label{more-morphisms-lemma-affine-formally-etale}
$f : X \to S$를 스킴의 사상이라 하고, $X$와 $S$가 아핀이라고
가정하자. 그러면 $f$가 형식적으로 에탈일 필요충분조건은 환 사상
$\mathcal{O}_S(S) \to \mathcal{O}_X(X)$가 형식적으로 에탈인 것이다.
\end{lemma}

\begin{proof}
환의 범주와 아핀 스킴의 범주 사이의 동치에 의해 정의들(Definition
\ref{more-morphisms-definition-formally-etale} 및 Algebra, Definition
\ref{algebra-definition-formally-etale})에서 즉시 따른다. Schemes, Lemma
\ref{schemes-lemma-category-affine-schemes}를 보라.
\end{proof}

\begin{lemma}
\label{more-morphisms-lemma-formally-etale}
$f : X \to Y$를 스킴의 사상이라 하자. 다음은 서로 동치이다.
\begin{enumerate}
\item $f$는 형식적으로 에탈이다.
\item 모든 $x \in X$에 대해 $f(U) \subset V$이고
$f|_U : U \to V$가 형식적으로 에탈인 열린집합 $x \in U \subset X$와
$f(x) \in V \subset Y$가 존재한다.
\item $f(U) \subset V$인 모든 아핀 열린집합 쌍 $U \subset X$와
$V \subset Y$에 대해 환 사상
$\mathcal{O}_Y(V) \to \mathcal{O}_X(U)$가 형식적으로 에탈이다.
\item 아핀 열린 덮개 $Y = \bigcup V_j$가 존재하고, 각 $j$에 대해
아핀 열린 덮개 $f^{-1}(V_j) = \bigcup U_{ji}$가 존재하여 모든 $j$와
$i$에 대해 환 사상 $\mathcal{O}_Y(V) \to \mathcal{O}_X(U)$가
형식적으로 에탈이다.
% P05-EN-CH37-0036
\end{enumerate}
\end{lemma}

\begin{proof}
Lemmas \ref{more-morphisms-lemma-formally-etale-on-opens}와
\ref{more-morphisms-lemma-affine-formally-etale}를 결합하면 (1) $\Rightarrow$ (3)을
얻는다. (3) $\Rightarrow$ (4)는 즉시 따른다. Lemma
\ref{more-morphisms-lemma-affine-formally-etale}에 의해 (4) $\Rightarrow$ (2)이다.
(2)가 성립하면 Lemma \ref{more-morphisms-lemma-characterize-formally-etale}에 의해 모든
점의 근방에서 $\Omega_{X/S}$와 $\mathcal{C}_{X/S}$가 소멸하고(여기서는
Morphisms, Lemma \ref{morphisms-lemma-differentials-restrict-open}과 Lemma
\ref{more-morphisms-lemma-universal-thickening-localize}도 사용한다), 다시 Lemma
\ref{more-morphisms-lemma-characterize-formally-etale}에 의해 $f$가 형식적으로
비분기임을 알 수 있다.
% P05-EN-CH37-0037 P05-EN-CH37-0038
\end{proof}

\begin{lemma}
\label{more-morphisms-lemma-etale-formally-etale}
$f : X \to S$를 스킴의 사상이라 하자. 다음은 서로 동치이다.
\begin{enumerate}
\item 사상 $f$는 에탈이다.
\item 사상 $f$는 국소 유한표현이고 형식적으로 에탈이다.
\end{enumerate}
\end{lemma}

\begin{proof}
$f$가 에탈이라고 가정하자. 에탈 사상은 국소 유한표현이고 평탄하며
비분기이다. Morphisms, Section \ref{morphisms-section-etale}을 보라.
따라서 $f$는 국소 유한표현이고 형식적으로 에탈이다. Lemma
\ref{more-morphisms-lemma-unramified-flat-formally-etale}를 보라.

\medskip\noindent
거꾸로 $f$가 국소 유한표현이고 형식적으로 에탈이라고 가정하자. 에탈인
것은 $X$와 $S$ 위의 Zariski 위상에서 국소적이다. Morphisms, Lemma
\ref{morphisms-lemma-etale-characterize}를 보라. Lemma
\ref{more-morphisms-lemma-formally-etale-on-opens}에 의해 $X$를 다음과 같은 아핀
열린집합 $U$들로 덮을 수 있다. 각 열린집합은 어떤 아핀 열린집합 $V$로
가고, $U \to V$는 형식적으로 에탈이다(또한 유한표현이다. Morphisms,
Lemma \ref{morphisms-lemma-locally-finite-presentation-characterize} 참조).
Lemma \ref{more-morphisms-lemma-affine-formally-etale}에 의해 환 사상
$\mathcal{O}(V) \to \mathcal{O}(U)$는 형식적으로 에탈이다(또한
유한표현이다). Algebra, Lemma \ref{algebra-lemma-formally-etale-etale}에
의해 결론을 얻는다. (형식적으로 매끄러운 사상을 논의할 때 이 함의의
다른 증명을 제시할 것이다.)
\end{proof}













\section{사상의 무한소 변형}
\label{more-morphisms-section-action-by-derivations}

\noindent
이 절에서는 미분을 사용하여 사상을 무한소적으로 움직이는 방법을 설명한다.
이 절 전체에서 $X$의 비후화 $X'$ 위의 층을 $X$ 위의 층으로 볼 수 있다는
사실을 사용한다.

\begin{lemma}
\label{more-morphisms-lemma-difference-derivation}
$S$를 스킴이라 하자.
$X \subset X'$와 $Y \subset Y'$를 $S$ 위의 두 1차 비후화라 하자.
$(a, a'), (b, b') : (X \subset X') \to (Y \subset Y')$를
$S$ 위의 비후화의 두 사상이라 하자. 다음을 가정하자.
\begin{enumerate}
\item $a = b$이고,
\item 두 사상 $a^*\mathcal{C}_{Y/Y'} \to \mathcal{C}_{X/X'}$
(Morphisms, Lemma \ref{morphisms-lemma-conormal-functorial})가 같다.
\end{enumerate}
그러면 사상 $(a')^\sharp - (b')^\sharp$는 다음과 같이 인수분해된다.
$$
\mathcal{O}_{Y'} \to \mathcal{O}_Y \xrightarrow{D}
a_*\mathcal{C}_{X/X'} \to a_*\mathcal{O}_{X'}
$$
여기서 $D$는 $\mathcal{O}_S$-미분이다.
\end{lemma}

\begin{proof}
$Y$ 위에서 작업하는 대신 $X$ 위에서 작업한다. 그러면 당김 함자
$a^{-1}$가 완전하다는 이점이 있다. (1)과 (2)를 사용하면 각 행이
완전한 다음 가환 그림을 얻는다.
$$
\xymatrix{
0 \ar[r] &
\mathcal{C}_{X/X'} \ar[r] &
\mathcal{O}_{X'} \ar[r] &
\mathcal{O}_X \ar[r] & 0 \\
0 \ar[r] &
a^{-1}\mathcal{C}_{Y/Y'} \ar[r] \ar[u] &
a^{-1}\mathcal{O}_{Y'}
\ar[r] \ar@<1ex>[u]^{(a')^\sharp} \ar@<-1ex>[u]_{(b')^\sharp} &
a^{-1}\mathcal{O}_Y \ar[r] \ar[u] & 0
}
$$
이제 이러한 상황에서 $\mathcal{O}_S$-대수 사상 $(a')^\sharp$와
$(b')^\sharp$의 차는 $a^{-1}\mathcal{O}_Y$에서
$\mathcal{C}_{X/X'}$로 가는 $\mathcal{O}_S$-미분이라는 일반적인 사실이 있다.
함자 $a^{-1}$와 $a_*$의 수반성에 의해 이는 $\mathcal{O}_Y$에서
$a_*\mathcal{C}_{X/X'}$로 가는 $\mathcal{O}_S$-미분과 같은 것이다.
일부 세부사항은 생략한다.
\end{proof}

\noindent
위 보조정리의 상황에서 $D$를 다음과 같이 쓸 수 있음에 유의하자.
\begin{equation}
\label{more-morphisms-equation-D}
% P05-EN-CH37-0039
D = \text{d}_{Y/S} \circ \theta
\end{equation}
여기서 $\theta$는 $\mathcal{O}_Y$-선형 사상
$\theta : \Omega_{Y/S} \to a_*\mathcal{C}_{X/X'}$이다.
물론 다시 수반성을 사용하면 $\theta$를 $\mathcal{O}_X$-선형 사상
$\theta : a^*\Omega_{Y/S} \to \mathcal{C}_{X/X'}$로 볼 수 있다.

\begin{lemma}
\label{more-morphisms-lemma-action-by-derivations}
$S$를 스킴이라 하자.
$(a, a') : (X \subset X') \to (Y \subset Y')$를
$S$ 위의 1차 비후화의 사상이라 하자.
$$
\theta : a^*\Omega_{Y/S} \to \mathcal{C}_{X/X'}
$$
를 $\mathcal{O}_X$-선형 사상이라 하자. 그러면 Lemma
\ref{more-morphisms-lemma-difference-derivation}의 (1)과 (2)를 만족하고 미분 $D$와
$\theta$가 Equation (\ref{more-morphisms-equation-D})로 관계되는 쌍의 사상
$(b, b') : (X \subset X') \to (Y \subset Y')$가 유일하게 존재한다.
\end{lemma}

\begin{proof}
$b = a$로 두고 $(b')^\sharp$를 다음 사상으로 정의하면 된다.
$$
(a')^\sharp + D : a^{-1}\mathcal{O}_{Y'} \to \mathcal{O}_{X'}
$$
여기서 $D$는 Equation (\ref{more-morphisms-equation-D})의 미분이다. $(b')^\sharp$가
$\mathcal{O}_S$-대수 층의 사상이고 Lemma
\ref{more-morphisms-lemma-difference-derivation}의 (1)과 (2)가 성립한다는 확인은 생략한다.
Equation (\ref{more-morphisms-equation-D})는 구성에 의해 성립한다.
\end{proof}

\begin{remark}
\label{more-morphisms-remark-action-by-derivations}
가정과 표기는 Lemma \ref{more-morphisms-lemma-action-by-derivations}와 같다고 하자.
국소 단면 $\theta$의 $a'$에 대한 작용을 때때로 $\theta \cdot a'$로 나타낸다.
이는 $(a')^\sharp$와 $(\theta \cdot a')^\sharp$의 차가 Equation
(\ref{more-morphisms-equation-D})에 의해 $\theta$와 관계되는 미분 $D$라는 뜻일 뿐이다.
\end{remark}

\begin{lemma}
\label{more-morphisms-lemma-sheaf}
$S$를 스킴이라 하자.
$X \subset X'$와 $Y \subset Y'$를 $S$ 위의 1차 비후화라 하자.
사상 $a : X \to Y$와 $\mathcal{O}_X$-가군의 사상
$A : a^*\mathcal{C}_{Y/Y'} \to \mathcal{C}_{X/X'}$가 주어졌다고 하자.
열린 부분스킴 $U' \subset X'$에 대하여 다음을 만족하는 사상
$a' : U' \to Y'$를 생각하자.
\begin{enumerate}
\item $a'$는 $S$ 위의 사상이고,
\item $a'|_U = a|_U$이며,
\item 유도된 사상
$a^*\mathcal{C}_{Y/Y'}|_U \to \mathcal{C}_{X/X'}|_U$
는 $A$를 $U$에 제한한 것이다.
\end{enumerate}
여기서 $U = X \cap U'$이다. 그러면 규칙
\begin{equation}
\label{more-morphisms-equation-sheaf}
U' \mapsto
\{a' : U' \to Y'\text{이며 (1), (2), (3)을 만족하는 것}\}
\end{equation}
은 $X'$ 위의 집합의 층을 정의한다.
\end{lemma}

\begin{proof}
% P05-EN-CH37-0040
보조정리의 규칙을 $\mathcal{F}$라 하자.
$V' \subset U' \subset X'$에 대한 $\mathcal{F}$의 제한 사상
$\mathcal{F}(U') \to \mathcal{F}(V')$는 바로 제한 사상
$a' \mapsto a'|_{V'}$이다. 사상은 국소적으로 정의되므로 이 정의 아래
$\mathcal{F}$가 층이라는 것은 명백하다.
\end{proof}

\noindent
다음 보조정리에서는 $X$ 위의 층과 $X$의 임의의 비후화 위의 층을 동일시한다.

\begin{lemma}
\label{more-morphisms-lemma-action-sheaf}
표기와 가정은 Lemma \ref{more-morphisms-lemma-sheaf}와 같다고 하자.
층
$$
\SheafHom_{\mathcal{O}_X}(a^*\Omega_{Y/S}, \mathcal{C}_{X/X'})
$$
은 층 (\ref{more-morphisms-equation-sheaf})에 작용한다. 더욱이 층
(\ref{more-morphisms-equation-sheaf})가 단면을 갖는 임의의 열린 부분스킴
$U' \subset X'$ 위에서 이 작용은 단순추이적이다.
\end{lemma}

\begin{proof}
이는 Lemmas \ref{more-morphisms-lemma-difference-derivation},
\ref{more-morphisms-lemma-action-by-derivations}, 그리고 \ref{more-morphisms-lemma-sheaf}를
결합한 것이다.
\end{proof}

\begin{remark}
\label{more-morphisms-remark-special-case}
Lemmas \ref{more-morphisms-lemma-difference-derivation},
\ref{more-morphisms-lemma-action-by-derivations},
\ref{more-morphisms-lemma-sheaf}, 그리고
\ref{more-morphisms-lemma-action-sheaf}의 한 특수한 경우는 $Y = Y'$인 경우이다.
이 경우 사상 $A$는 항상 영이다. Lemma \ref{more-morphisms-lemma-sheaf}의 층은
단순히 다음 규칙으로 주어진다.
$$
U' \mapsto
\{a' : U' \to Y\text{는 }S\text{ 위이고 } a'|_U = a|_U\text{를 만족한다}\}
$$
그리고 층
$\SheafHom_{\mathcal{O}_X}(a^*\Omega_{Y/S}, \mathcal{C}_{X/X'})$가
이에 작용한다.
\end{remark}

\begin{remark}
\label{more-morphisms-remark-another-special-case}
Lemmas \ref{more-morphisms-lemma-difference-derivation},
\ref{more-morphisms-lemma-action-by-derivations},
\ref{more-morphisms-lemma-sheaf}, 그리고
\ref{more-morphisms-lemma-action-sheaf}의 또 다른 특수한 경우는 $S$ 자체가 비후화
$Z \subset Z' = S$이고 $Y = Z \times_{Z'} Y'$인 경우이다. 그림은 다음과 같다.
$$
\xymatrix{
(X \subset X') \ar@{..>}[rr]_{(a, ?)} \ar[rd]_{(g, g')} & &
(Y \subset Y') \ar[ld]^{(h, h')} \\
& (Z \subset Z')
}
$$
이 경우 사상 $A : a^*\mathcal{C}_{Y/Y'} \to \mathcal{C}_{X/X'}$는
$a$에 의해 결정된다. 실제로 사상
$h^*\mathcal{C}_{Z/Z'} \to \mathcal{C}_{Y/Y'}$는 전사이다
($Y = Z \times_{Z'} Y'$라고 가정했기 때문이다). 따라서 당김
$g^*\mathcal{C}_{Z/Z'} = a^*h^*\mathcal{C}_{Z/Z'} \to
a^*\mathcal{C}_{Y/Y'}$는 전사이고, 합성
$g^*\mathcal{C}_{Z/Z'} \to a^*\mathcal{C}_{Y/Y'} \to \mathcal{C}_{X/X'}$
은 $g'$에 의해 유도되는 표준 사상이어야 한다. 그러므로 Lemma
\ref{more-morphisms-lemma-sheaf}의 층은 단순히 다음 규칙으로 주어진다.
$$
U' \mapsto
\{a' : U' \to Y'\text{는 }Z'\text{ 위이고 } a'|_U = a|_U\text{를 만족한다}\}
$$
그리고 층
$\SheafHom_{\mathcal{O}_X}(a^*\Omega_{Y/Z}, \mathcal{C}_{X/X'})$가
이에 작용한다.
\end{remark}

\begin{lemma}
\label{more-morphisms-lemma-omega-deformation}
$S$를 스킴이라 하자. $X \subset X'$를 $S$ 위의 1차 비후화라 하고,
$Y$를 $S$ 위의 스킴이라 하자. $a = a'|_X = b'|_X$를 만족하는
$S$ 위의 두 사상 $a', b' : X' \to Y$를 택하자. 이로부터 가환 그림
$$
\xymatrix{
X \ar[r] \ar[d]_a & X' \ar[d]^{(b', a')} \\
Y \ar[r]^-{\Delta_{Y/S}} & Y \times_S Y
}
$$
을 얻는다. 가로 화살표들은 각각 conormal 층
$\mathcal{C}_{X/X'}$와 $\Omega_{Y/S}$를 갖는 몰입이므로,
Morphisms, Lemma \ref{morphisms-lemma-conormal-functorial}에 의해 사상
$\theta : a^*\Omega_{Y/S} \to \mathcal{C}_{X/X'}$를 얻는다.
이 $\theta$와 Lemma \ref{more-morphisms-lemma-difference-derivation}의 미분 $D$는
Equation (\ref{more-morphisms-equation-D})로 관계된다.
\end{lemma}

\begin{proof}
생략한다. 힌트: 이 등식은 아핀 열린 부분스킴 위에서 확인할 수 있으며,
그곳에서는 다음 계산으로부터 따른다. $f$가 $\mathcal{O}_Y$의 국소 단면이면
$1 \otimes f - f \otimes 1$은 $\text{d}_{Y/S}(f)$에 대응하는
$\mathcal{C}_{Y/(Y \times_S Y)}$의 국소 단면이다. 이는
$\mathcal{C}_{X/X'}$의 국소 단면
$(a')^\sharp(f) - (b')^\sharp(f) = D(f)$로 보내진다. 다시 말해
$\theta(\text{d}_{Y/S}(f)) = D(f)$이다.
\end{proof}

\noindent
나중에 사용하기 위해 Lemma \ref{more-morphisms-lemma-action-by-derivations}의 구성이
에탈 국소화와 양립함을 대략적으로 말해 주는 결과가 필요하다.

\begin{lemma}
\label{more-morphisms-lemma-sheaf-differentials-etale-localization}
다음 가환 그림을 생각하자.
$$
\xymatrix{
X_1 \ar[d] & X_2 \ar[l]^f \ar[d] \\
S_1 & S_2 \ar[l]
}
$$
$X_2 \to X_1$과 $S_2 \to S_1$은 에탈이라고 하자. 그러면
Morphisms, Lemma \ref{morphisms-lemma-functoriality-differentials}의 사상
$c_f : f^*\Omega_{X_1/S_1} \to \Omega_{X_2/S_2}$는 동형이다.
\end{lemma}

\begin{proof}
에탈 사상 $U \to V$는 $\Omega_{U/V} = 0$을 만족하는 매끄러운 사상임을
상기하자. 이를 사용하면 Morphisms, Lemma
\ref{morphisms-lemma-triangle-differentials}로부터
$\Omega_{X_2/S_2} = \Omega_{X_2/S_1}$을 얻고, Morphisms, Lemma
\ref{morphisms-lemma-triangle-differentials-smooth}로부터 사상
$f^*\Omega_{X_1/S_1} \to \Omega_{X_2/S_1}$이 동형임을 얻는다
($f$를 $S_1$ 위의 사상으로 본다). 따라서 보조정리가 따른다.
\end{proof}

\begin{lemma}
\label{more-morphisms-lemma-action-by-derivations-etale-localization}
1차 비후화의 가환 그림
$$
\vcenter{
\xymatrix{
(T_2 \subset T_2') \ar[d]_{(h, h')} \ar[rr]_{(a_2, a_2')} & &
(X_2 \subset X_2') \ar[d]^{(f, f')} \\
(T_1 \subset T_1') \ar[rr]^{(a_1, a_1')} & &
(X_1 \subset X_1')
}
}
\quad
\begin{matrix}
\text{그리고 스킴의} \\
\text{가환 그림}
\end{matrix}
\quad
\vcenter{
\xymatrix{
X_2' \ar[r] \ar[d] & S_2 \ar[d] \\
X_1' \ar[r] & S_1
}
}
$$
을 생각하고, $X_2 \to X_1$과 $S_2 \to S_1$은 에탈이라고 하자.
임의의 $\mathcal{O}_{T_1}$-선형 사상
$\theta_1 : a_1^*\Omega_{X_1/S_1} \to \mathcal{C}_{T_1/T'_1}$에 대해
$\theta_2$를 다음 합성으로 정의하자.
$$
\xymatrix{
a_2^*\Omega_{X_2/S_2} \ar@{=}[r] &
h^*a_1^*\Omega_{X_1/S_1} \ar[r]^-{h^*\theta_1} &
h^*\mathcal{C}_{T_1/T'_1} \ar[r] &
\mathcal{C}_{T_2/T'_2}
}
$$
(등호는 증명에서 설명한다). 그러면 그림
$$
\xymatrix{
T_2' \ar[rr]_{\theta_2 \cdot a_2'} \ar[d] & & X'_2 \ar[d] \\
T_1' \ar[rr]^{\theta_1 \cdot a_1'} & & X'_1
}
$$
은 가환이다. 여기서 작용 $\theta_2 \cdot a_2'$와
$\theta_1 \cdot a_1'$는 Remark \ref{more-morphisms-remark-action-by-derivations}의 작용이다.
\end{lemma}

\begin{proof}
등호는 Lemma \ref{more-morphisms-lemma-sheaf-differentials-etale-localization}의 동일시
$f^*\Omega_{X_1/S_1} = \Omega_{X_2/S_2}$에서 온다. 즉, 이를 사용하면
$f \circ a_2 = a_1 \circ h$이므로
$a_2^*\Omega_{X_2/S_2} = a_2^*f^*\Omega_{X_1/S_1} =
h^*a_1^*\Omega_{X_1/S_1}$이다. 이를 확인했으므로 그림의 가환성은
아핀 열린 부분스킴 위에서 검사할 수 있다. 따라서 처음의 큰 그림에 나타나는
스킴들이 아핀이라고 가정해도 된다. 그러면 다음 가환 그림들을 얻는다.
$$
\vcenter{
\xymatrix{
(B'_2, I_2) & & (A'_2, J_2) \ar[ll]^{a_2'} \\
(B'_1, I_1) \ar[u]^{h'} & & (A'_1, J_1) \ar[ll]_{a_1'} \ar[u]_{f'}
}
}
\quad\text{그리고}\quad
\vcenter{
\xymatrix{
A'_2 & & R_2 \ar[ll] \\
A'_1 \ar[u] & & R_1 \ar[ll] \ar[u]
}
}
$$
이 표기는 $I_1, I_2, J_1, J_2$가 제곱영 아이디얼이고 쌍의 사상은
아이디얼을 아이디얼로 보내는 환 사상임을 뜻한다.
$A_1 = A'_1/J_1$, $A_2 = A'_2/J_2$, $B_1 = B'_1/I_1$, 그리고
$B_2 = B'_2/I_2$로 두자. 사상
$$
A_2 \otimes_{A_1} \Omega_{A_1/R_1} \longrightarrow \Omega_{A_2/R_2}
$$
이 동형이라고 주어져 있다. 그러면
$\theta_1 : B_1 \otimes_{A_1} \Omega_{A_1/R_1} \to I_1$은
$B_1$-선형이다. 이로부터 $R_1$-미분
$D_1 = \theta_1 \circ \text{d}_{A_1/R_1} : A_1 \to I_1$을 얻는다.
마찬가지로
$\theta_2 : B_2 \otimes_{A_2} \Omega_{A_2/R_2} \to I_2$로부터
$R_2$-미분
$D_2 = \theta_2 \circ \text{d}_{A_2/R_2} : A_2 \to I_2$를 얻는다.
$\theta_2$의 구성은 $\theta_1$과 $\theta_2$ 사이의 다음 양립성을 함의한다.
모든 $x \in A_1$에 대해
$$
h'(D_1(x)) = D_2(f'(x))
$$
$I_2$의 원소로서 등식이 성립한다. $D_1$은 사상
% P05-EN-CH37-0041
$A'_1 \to A_1 \xrightarrow{D_1} I_1 \to B_1$을 사용하여
$A'_1 \to B'_1$인 사상으로 볼 수 있고, 마찬가지로 $D_2$를
$A'_2 \to B'_2$인 사상으로 볼 수 있다. 그러면 위 등식은
$x \in A'_1$에 대해 성립한다. Lemma
\ref{more-morphisms-lemma-action-by-derivations}와 Remark
\ref{more-morphisms-remark-action-by-derivations}에서의 작용의 구성에 의해
$\theta_1 \cdot a_1'$는 환 사상 $a_1' + D_1 : A'_1 \to B'_1$에 대응하고,
$\theta_2 \cdot a_2'$는 환 사상 $a_2' + D_2 : A'_2 \to B'_2$에 대응한다.
위 등식으로부터
$h' \circ (a_1' + D_1) = (a_2' + D_2) \circ f'$
를 얻는다. 이것이 원하는 결과이다.
\end{proof}

\begin{remark}
\label{more-morphisms-remark-tiny-improvement}
Lemma \ref{more-morphisms-lemma-action-by-derivations-etale-localization}는 다음과 같이
개선할 수 있다. Lemma \ref{more-morphisms-lemma-action-by-derivations-etale-localization}와
같은 가환 그림들이 있지만 $X_2 \to X_1$과 $S_2 \to S_1$이
에탈이라고 가정하지는 않자. 다음으로
$\theta_1 : a_1^*\Omega_{X_1/S_1} \to \mathcal{C}_{T_1/T'_1}$
과
$\theta_2 : a_2^*\Omega_{X_2/S_2} \to \mathcal{C}_{T_2/T'_2}$
가 있어서
$$
\xymatrix{
f_*\mathcal{O}_{X_2} \ar[rr]_{f_*D_2} & &
f_*a_{2, *}\mathcal{C}_{T_2/T_2'} \\
\mathcal{O}_{X_1} \ar[rr]^{D_1} \ar[u]^{f^\sharp} & &
a_{1, *}\mathcal{C}_{T_1/T_1'} \ar[u]_{\text{다음에 의해 유도됨: }(h')^\sharp}
}
$$
가환이라고 하자. 여기서 $D_i$는 Equation (\ref{more-morphisms-equation-D})에서와 같이
$\theta_i$에 대응한다. 그러면 Lemma
\ref{more-morphisms-lemma-action-by-derivations-etale-localization}의 결론이 성립한다.
$X_2 \to X_1$과 $S_2 \to S_1$이 모두 에탈이라는 조건의 중요성은
$\theta_1$로부터 $\theta_2$를 구성할 수 있게 한다는 데 있다.
\end{remark}








\section{스킴의 무한소 변형}
\label{more-morphisms-section-deform}

\noindent
다음의 간단한 보조정리는 사상의 무한소 변형이 평탄한지 검사하는 데
흔히 편리한 도구가 된다.

\begin{lemma}
\label{more-morphisms-lemma-deform}
$(f, f') : (X \subset X') \to (S \subset S')$를 1차 비후화의 사상이라 하자.
$f$가 평탄하다고 가정하자. 그러면 다음은 동치이다.
\begin{enumerate}
\item $f'$가 평탄하고 $X = S \times_{S'} X'$이다.
\item 표준 사상 $f^*\mathcal{C}_{S/S'} \to \mathcal{C}_{X/X'}$가 동형이다.
\end{enumerate}
\end{lemma}

\begin{proof}
문제는 $X'$ 위에서 국소적이므로 $X, X', S, S'$가 아핀 스킴이라고
가정해도 된다. $S' = \Spec(A')$, $X' = \Spec(B')$,
$S = \Spec(A)$, $X = \Spec(B)$라 쓰자. 여기서 어떤 제곱영 아이디얼에 대해
$A = A'/I$이고 $B = B'/J$이다. 그러면 다음 가환 그림을 얻는다.
$$
\xymatrix{
0 \ar[r] &
J \ar[r] &
B' \ar[r] &
B \ar[r] & 0 \\
0 \ar[r] &
I \ar[r] \ar[u] &
A' \ar[r] \ar[u] &
A \ar[r] \ar[u] & 0
}
$$
각 행은 완전하다. 보조정리의 표준 사상은 다음 사상이다.
$$
I \otimes_A B = I \otimes_{A'} B' \longrightarrow J.
$$
$f$가 평탄하다는 가정은 $A \to B$가 평탄하다는 뜻이다.

\medskip\noindent
(1)을 가정하자. 그러면 $A' \to B'$는 평탄하고 $J = IB'$이다.
평탄성으로부터 $\text{Tor}_1^{A'}(B', A) = 0$을 얻는다
(Algebra, Lemma \ref{algebra-lemma-characterize-flat} 참조).
이는 $I \otimes_{A'} B' \to B'$가 단사라는 뜻이다
(Algebra, Remark \ref{algebra-remark-Tor-ring-mod-ideal} 참조).
따라서 $I \otimes_A B \to J$는 동형이다.

\medskip\noindent
(2)를 가정하자. 그러면 $J = IB'$이고 따라서 $X = S \times_{S'} X'$이다.
더욱이 앞 단락의 함의를 거꾸로 적용하면
$\text{Tor}_1^{A'}(B', A'/I) = 0$을 얻는다. 그러므로 Algebra, Lemma
\ref{algebra-lemma-what-does-it-mean}에 의해 $B'$는 $A'$ 위에서 평탄하다.
\end{proof}

\noindent
다음 보조정리는 ``crit\`ere de platitude par fibres''의 ``멱영'' 판이다.
Section \ref{more-morphisms-section-criterion-flat-fibres}를 보라.

\begin{lemma}
\label{more-morphisms-lemma-flatness-morphism-thickenings}
비후화의 가환 그림
$$
\xymatrix{
(X \subset X') \ar[rr]_{(f, f')} \ar[rd] & & (Y \subset Y') \ar[ld] \\
& (S \subset S')
}
$$
을 생각하자. 다음을 가정하자.
\begin{enumerate}
\item $X'$는 $S'$ 위에서 평탄하다.
\item $f$는 평탄하다.
\item $S \subset S'$는 유한 차수 비후화이다.
\item $X = S \times_{S'} X'$이고 $Y = S \times_{S'} Y'$이다.
\end{enumerate}
그러면 $f'$는 평탄하고, $f'$의 상에 있는 모든 점에서 $Y'$는
$S'$ 위에서 평탄하다.
\end{lemma}

\begin{proof}
Algebra, Lemma \ref{algebra-lemma-criterion-flatness-fibre-nilpotent}의
직접적인 귀결이다.
\end{proof}

\noindent
스킴의 사상이 갖는 많은 성질은 평탄 변형 아래에서 보존된다.

\begin{lemma}
\label{more-morphisms-lemma-deform-property}
비후화의 가환 그림
$$
\xymatrix{
(X \subset X') \ar[rr]_{(f, f')} \ar[rd] & & (Y \subset Y') \ar[ld] \\
& (S \subset S')
}
$$
을 생각하자. $S \subset S'$는 유한 차수 비후화이고, $X'$는 $S'$ 위에서
평탄하며, $X = S \times_{S'} X'$이고 $Y = S \times_{S'} Y'$라고 가정하자.
그러면 다음이 성립한다.
\begin{enumerate}
\item $f$가 평탄일 필요충분조건은 $f'$가 평탄인 것이다.
\label{more-morphisms-item-flat}
\item $f$가 동형일 필요충분조건은 $f'$가 동형인 것이다.
\label{more-morphisms-item-isomorphism}
\item $f$가 열린 몰입일 필요충분조건은 $f'$가 열린 몰입인 것이다.
\label{more-morphisms-item-open-immersion}
\item $f$가 준콤팩트일 필요충분조건은 $f'$가 준콤팩트인 것이다.
\label{more-morphisms-item-quasi-compact}
\item $f$가 보편적으로 닫혔을 필요충분조건은 $f'$가 보편적으로 닫힌 것이다.
\label{more-morphisms-item-universally-closed}
\item $f$가 (준)분리일 필요충분조건은 $f'$가 (준)분리인 것이다.
\label{more-morphisms-item-separated}
\item $f$가 모노사상일 필요충분조건은 $f'$가 모노사상인 것이다.
\label{more-morphisms-item-monomorphism}
\item $f$가 전사일 필요충분조건은 $f'$가 전사인 것이다.
\label{more-morphisms-item-surjective}
\item $f$가 보편 단사일 필요충분조건은 $f'$가 보편 단사인 것이다.
\label{more-morphisms-item-universally-injective}
\item $f$가 아핀일 필요충분조건은 $f'$가 아핀인 것이다.
\label{more-morphisms-item-affine}
\item
\label{more-morphisms-item-finite-type}
$f$가 국소 유한형일 필요충분조건은 $f'$가 국소 유한형인 것이다.
\item $f$가 국소 준유한일 필요충분조건은 $f'$가 국소 준유한인 것이다.
\label{more-morphisms-item-quasi-finite}
\item
\label{more-morphisms-item-finite-presentation}
$f$가 국소 유한표현일 필요충분조건은 $f'$가 국소 유한표현인 것이다.
\item
\label{more-morphisms-item-relative-dimension-d}
$f$가 상대차원 $d$인 국소 유한형일 필요충분조건은 $f'$가
상대차원 $d$인 국소 유한형인 것이다.
\item $f$가 보편적으로 열렸을 필요충분조건은 $f'$가 보편적으로 열린 것이다.
\label{more-morphisms-item-universally-open}
\item $f$가 신토믹일 필요충분조건은 $f'$가 신토믹인 것이다.
\label{more-morphisms-item-syntomic}
\item $f$가 매끄러울 필요충분조건은 $f'$가 매끄러운 것이다.
\label{more-morphisms-item-smooth}
\item $f$가 비분기일 필요충분조건은 $f'$가 비분기인 것이다.
\label{more-morphisms-item-unramified}
\item $f$가 에탈일 필요충분조건은 $f'$가 에탈인 것이다.
\label{more-morphisms-item-etale}
\item $f$가 고유일 필요충분조건은 $f'$가 고유인 것이다.
\label{more-morphisms-item-proper}
\item $f$가 정수일 필요충분조건은 $f'$가 정수인 것이다.
\label{more-morphisms-item-integral}
\item $f$가 유한일 필요충분조건은 $f'$가 유한인 것이다.
\label{more-morphisms-item-finite}
\item
\label{more-morphisms-item-finite-locally-free}
$f$가 (계수 $d$의) 유한 국소 자유일 필요충분조건은 $f'$가
(계수 $d$의) 유한 국소 자유인 것이고,
% P05-EN-CH37-0042
\item 여기에 더 추가할 것.
\end{enumerate}
\end{lemma}

\begin{proof}
$X$와 $Y$에 대한 가정은 $f$가 $X \to X'$에 의한 $f'$의 밑변환이라는
뜻이다. 위 (1) -- (23)에 열거한 성질 $\mathcal{P}$는 모두 밑변환 아래에서
안정적이므로, $f'$가 성질 $\mathcal{P}$를 가지면 $f$도 그 성질을 갖는다. 다음을 보라.
Schemes, Lemmas \ref{schemes-lemma-base-change-immersion},
\ref{schemes-lemma-quasi-compact-preserved-base-change},
\ref{schemes-lemma-separated-permanence}, 그리고
\ref{schemes-lemma-base-change-monomorphism}
그리고
Morphisms, Lemmas
\ref{morphisms-lemma-base-change-surjective},
\ref{morphisms-lemma-base-change-universally-injective},
\ref{morphisms-lemma-base-change-affine},
\ref{morphisms-lemma-base-change-finite-type},
\ref{morphisms-lemma-base-change-quasi-finite},
\ref{morphisms-lemma-base-change-finite-presentation},
\ref{morphisms-lemma-base-change-relative-dimension-d},
\ref{morphisms-lemma-base-change-syntomic},
\ref{morphisms-lemma-base-change-smooth},
\ref{morphisms-lemma-base-change-unramified},
\ref{morphisms-lemma-base-change-etale},
\ref{morphisms-lemma-base-change-proper},
\ref{morphisms-lemma-base-change-finite}, 그리고
\ref{morphisms-lemma-base-change-finite-locally-free}.

\medskip\noindent
따라서 각 경우에 흥미로운 방향은 $f$가 그 성질을 갖는다고 가정하고
$f'$도 그 성질을 가짐을 보이는 방향이다. 비후화의 차수에 관한 귀납법으로
$S \subset S'$가 1차 비후화라고 가정해도 된다. Definition
\ref{more-morphisms-definition-thickening} 바로 뒤의 논의를 보라. 아래의 논증에서 더 언급하지
않고 사용할 몇 가지 일반적인 주의를 먼저 둔다.
(I) $W' \subset S'$를 아핀 열린 부분스킴이라 하고, $U' \subset X'$와
$V' \subset Y'$를 $W'$ 위에 놓인 아핀 열린 부분스킴으로서
$f'(U') \subset V'$를 만족한다고 하자. $W' = \Spec(R')$라 하고,
% P05-EN-CH37-0043
$I \subset R'$를 닫힌 부분스킴 $W' \cap S$를 정의하는 아이디얼이라 하자.
$U' = \Spec(B')$와 $V' = \Spec(A')$라 쓰자. 그러면 가환 그림
$$
\xymatrix{
0 \ar[r] &
IB' \ar[r] &
B' \ar[r] &
B \ar[r] & 0 \\
0 \ar[r] &
IA' \ar[r] \ar[u] &
A' \ar[r] \ar[u] &
A \ar[r] \ar[u] & 0
}
$$
을 얻고 각 행은 완전하다. 또한 $IB' \cong I \otimes_R B$이다.
Lemma \ref{more-morphisms-lemma-deform}의 증명을 보라.
(II) 사상 $X \to X'$와 $Y \to Y'$는 보편 위상동형이다. 따라서 임의의
밑변환 뒤에도 사상 $f$와 $f'$의 위상은 동일하다.
(III) $f$가 평탄하면 $f'$도 평탄하고, $f'$의 상에 있는 모든 점에서
$Y' \to S'$가 평탄하다. Lemma
\ref{more-morphisms-lemma-flatness-morphism-thickenings}를 보라.

\medskip\noindent
(\ref{more-morphisms-item-flat})에 관하여. 이는 일반 주의 (III)이다.

\medskip\noindent
(\ref{more-morphisms-item-isomorphism})에 관하여. $f$가 동형이라고 가정하자.
(III)에 의해 $Y' \to S'$는 평탄하다. 아핀 열린 부분스킴
$V' \subset Y'$를 택하고 $U' = (f')^{-1}(V')$로 두자. 그러면
$V = Y \cap V'$는 아핀이고, 따라서
$V \cong f^{-1}(V) = U = Y \times_{Y'} U'$도 아핀이다. Lemma
\ref{more-morphisms-lemma-thickening-affine-scheme}에 의해 $U'$는 아핀이다. 그러므로
일반 주의 (I)의 그림을 얻으며, $R' \to A'$가 평탄하므로 또한
$IA \cong I \otimes_R A$이다. 이제
$IB' \cong I \otimes_R B \cong I \otimes_R A \cong IA'$
이고 $A \cong B$이다. 위 그림의 각 행이 완전하므로
$A' \cong B'$, 즉 $U' \cong V'$임을 알 수 있다. 따라서 $f'$는 동형이다.

\medskip\noindent
(\ref{more-morphisms-item-open-immersion})에 관하여. $f$가 열린 몰입이라고 가정하자.
그러면 $f$는 $X$와 어떤 열린 부분스킴 $V \subset Y$ 사이의 동형이다.
$V' \subset Y'$를 바탕 위상공간이 $V$인 열린 부분스킴이라 하자.
그러면 $f'$는 $X'$에서 $V'$로 가는 사상이고
(\ref{more-morphisms-item-isomorphism})에 의해 동형이다. 따라서 $f'$는 열린 몰입이다.

\medskip\noindent
(\ref{more-morphisms-item-quasi-compact})에 관하여. 주의 (II)에서 바로 따른다. 더 일반적인
명제는 Lemma \ref{more-morphisms-lemma-thicken-property-morphisms}도 보라.

\medskip\noindent
(\ref{more-morphisms-item-universally-closed})에 관하여. 주의 (II)에서 바로 따른다.
더 일반적인 명제는 Lemma \ref{more-morphisms-lemma-thicken-property-morphisms}도 보라.

\medskip\noindent
(\ref{more-morphisms-item-separated})에 관하여. 다음에 유의하자.
$X \times_Y X = Y \times_{Y'} (X' \times_{Y'} X')$이므로
$X' \times_{Y'} X'$는 $X \times_Y X$의 비후화이다. 따라서 사상
$\Delta_{X/Y}$와 $\Delta_{X'/Y'}$의 위상이 일치하므로 결론을 얻는다.
더 일반적인 명제는 Lemma \ref{more-morphisms-lemma-thicken-property-morphisms}도 보라.

\medskip\noindent
(\ref{more-morphisms-item-monomorphism})에 관하여. $f$가 모노사상이라고 가정하자.
대각 사상 $\Delta_{X'/Y'} : X' \to X' \times_{Y'} X'$를 생각하자.
$S \to S'$에 의한 $\Delta_{X'/Y'}$의 밑변환은 $\Delta_{X/Y}$이고,
이는 가정에 의해 동형이다. (\ref{more-morphisms-item-isomorphism})에 의해
$\Delta_{X'/Y'}$도 동형이다.

\medskip\noindent
(\ref{more-morphisms-item-surjective})에 관하여. 이는 명백하다. 더 일반적인 명제는
Lemma \ref{more-morphisms-lemma-thicken-property-morphisms}도 보라.

\medskip\noindent
(\ref{more-morphisms-item-universally-injective})에 관하여. 주의 (II)에서 바로 따른다.
더 일반적인 명제는 Lemma \ref{more-morphisms-lemma-thicken-property-morphisms}도 보라.

\medskip\noindent
(\ref{more-morphisms-item-affine})에 관하여. $f$가 아핀이라고 가정하자. 아핀 열린
부분스킴 $V' \subset Y'$를 택하고 $U' = (f')^{-1}(V')$로 두자.
그러면 $V = Y \cap V'$가 아핀이므로 $U = Y \times_{Y'} U'$도 아핀이다.
Lemma \ref{more-morphisms-lemma-thickening-affine-scheme}에 의해 $U'$는 아핀이다.
따라서 $f'$는 아핀이다. 더 일반적인 명제는 Lemma
\ref{more-morphisms-lemma-thicken-property-morphisms}도 보라.

\medskip\noindent
(\ref{more-morphisms-item-finite-type})에 관하여. 주의 (I)에 의해 $A \to B$가 유한형일 때
$A' \to B'$도 유한형임을 보이는 문제로 귀착된다. $x_1, \ldots, x_n \in B'$의
$B$에서의 상들이 $B$를 $A$-대수로서 생성한다고 하자. 그러면
% P05-EN-CH37-0044
$A'[x_1, \ldots, x_n] \to B$는 전사이다. 실제로
$A'[x_1, \ldots, x_n] \to B$와
$I \otimes_R A[x_1, \ldots, x_n] \to I \otimes_R B$가 모두 전사이기 때문이다.
더 일반적인 명제는 Lemma
\ref{more-morphisms-lemma-thicken-property-morphisms-cartesian}도 보라.

\medskip\noindent
(\ref{more-morphisms-item-quasi-finite})에 관하여. (\ref{more-morphisms-item-finite-type}) 및 유한형 사상의
준유한성은 올 위에서 검사할 수 있다는 사실로부터 따른다. Morphisms, Lemma
\ref{morphisms-lemma-quasi-finite-at-point-characterize}를 보라. 더 일반적인
명제는 Lemma \ref{more-morphisms-lemma-thicken-property-morphisms-cartesian}도 보라.

\medskip\noindent
(\ref{more-morphisms-item-finite-presentation})에 관하여. 주의 (I)에 의해 $A \to B$가
유한표현일 때 $A' \to B'$도 유한표현임을 보이는 문제로 귀착된다.
(\ref{more-morphisms-item-finite-type})에 의해 어떤 아이디얼 $K'$에 대해
$B' = A'[x_1, \ldots, x_n]/K'$라고 가정해도 된다. 짧은 완전열
$$
0 \to K' \to A'[x_1, \ldots, x_n] \to B' \to 0
$$
을 얻는다. $B'$가 $R'$ 위에서 평탄하므로 $K' \otimes_{R'} R$은 전사
$A[x_1, \ldots, x_n] \to B$의 핵이다. $A \to B$에 대한 가정에 의해
$A[x_1, \ldots, x_n]$에서의 상들이 이 핵을 생성하는 유한 개의 원소
$f'_1, \ldots, f'_m \in K'$가 존재한다. $I$가 멱영이므로 Nakayama의
보조정리에 의해 $f'_1, \ldots, f'_m$이 $K'$를 생성한다. Algebra, Lemma
\ref{algebra-lemma-NAK}를 보라.

\medskip\noindent
(\ref{more-morphisms-item-relative-dimension-d})에 관하여. (\ref{more-morphisms-item-finite-type})와
일반 주의 (II)로부터 따른다. 더 일반적인 명제는 Lemma
\ref{more-morphisms-lemma-thicken-property-morphisms-cartesian}도 보라.

\medskip\noindent
(\ref{more-morphisms-item-universally-open})에 관하여. 일반 주의 (II)에서 바로 따른다.
더 일반적인 명제는 Lemma \ref{more-morphisms-lemma-thicken-property-morphisms}도 보라.

\medskip\noindent
(\ref{more-morphisms-item-syntomic})에 관하여. $f$가 신토믹이라고 가정하자.
(\ref{more-morphisms-item-finite-presentation})에 의해 $f'$는 국소 유한표현이고,
일반 주의 (III)에 의해 $f'$는 평탄하며, $f'$의 올들은 $f$의 올들이다.
따라서 Morphisms, Lemma \ref{morphisms-lemma-syntomic-flat-fibres}에 의해
$f'$는 신토믹이다.

\medskip\noindent
(\ref{more-morphisms-item-smooth})에 관하여. $f$가 매끄럽다고 가정하자.
(\ref{more-morphisms-item-finite-presentation})에 의해 $f'$는 국소 유한표현이고,
일반 주의 (III)에 의해 $f'$는 평탄하며, $f'$의 올들은 $f$의 올들이다.
따라서 Morphisms, Lemma \ref{morphisms-lemma-smooth-flat-smooth-fibres}에 의해
$f'$는 매끄럽다.

\medskip\noindent
(\ref{more-morphisms-item-unramified})에 관하여. $f$가 비분기라고 가정하자.
(\ref{more-morphisms-item-finite-type})에 의해 $f'$는 국소 유한형이고 $f'$의 올들은
$f$의 올들이다. 따라서 Morphisms, Lemma
\ref{morphisms-lemma-unramified-etale-fibres}에 의해 $f'$는 비분기이다.
더 일반적인 명제는 Lemma
\ref{more-morphisms-lemma-thicken-property-morphisms-cartesian}도 보라.

\medskip\noindent
(\ref{more-morphisms-item-etale})에 관하여. $f$가 에탈이라고 가정하자.
(\ref{more-morphisms-item-finite-presentation})에 의해 $f'$는 국소 유한표현이고,
일반 주의 (III)에 의해 $f'$는 평탄하며, $f'$의 올들은 $f$의 올들이다.
따라서 Morphisms, Lemma \ref{morphisms-lemma-etale-flat-etale-fibres}에 의해
$f'$는 에탈이다.

\medskip\noindent
(\ref{more-morphisms-item-proper})에 관하여. (\ref{more-morphisms-item-separated}),
(\ref{more-morphisms-item-finite-type}), (\ref{more-morphisms-item-quasi-compact}), 그리고
(\ref{more-morphisms-item-universally-closed})를 결합하면 따른다. 더 일반적인 명제는
Lemma \ref{more-morphisms-lemma-thicken-property-morphisms-cartesian}도 보라.

\medskip\noindent
(\ref{more-morphisms-item-integral})에 관하여. (\ref{more-morphisms-item-universally-closed})와
(\ref{more-morphisms-item-affine})를 Morphisms, Lemma
\ref{morphisms-lemma-integral-universally-closed}와 결합한다. 더 일반적인
명제는 Lemma \ref{more-morphisms-lemma-thicken-property-morphisms}도 보라.

\medskip\noindent
(\ref{more-morphisms-item-finite})에 관하여. (\ref{more-morphisms-item-integral})과
(\ref{more-morphisms-item-finite-type})를 Morphisms, Lemma
\ref{morphisms-lemma-finite-integral}과 결합한다. 더 일반적인 명제는
Lemma \ref{more-morphisms-lemma-thicken-property-morphisms-cartesian}도 보라.

\medskip\noindent
(\ref{more-morphisms-item-finite-locally-free})에 관하여. $f$가 유한 국소 자유라고
가정하자. (\ref{more-morphisms-item-finite})에 의해 $f'$는 유한이고, 일반 주의 (III)에
의해 $f'$는 평탄하며, (\ref{more-morphisms-item-finite-presentation})에 의해 $f'$는
국소 유한표현이다. 따라서 Morphisms, Lemma
\ref{morphisms-lemma-finite-flat}에 의해 $f'$는 유한 국소 자유이다.
\end{proof}

\noindent
다음 보조정리는 ``crit\`ere de platitude par fibres''의 ``국소 멱영'' 판이다.
Section \ref{more-morphisms-section-criterion-flat-fibres}를 보라.

\begin{lemma}
\label{more-morphisms-lemma-flatness-morphism-thickenings-fp-over-ft}
비후화의 가환 그림
$$
\xymatrix{
(X \subset X') \ar[rr]_{(f, f')} \ar[rd] & & (Y \subset Y') \ar[ld] \\
& (S \subset S')
}
$$
을 생각하자. 다음을 가정하자.
\begin{enumerate}
\item $Y' \to S'$는 국소 유한형이다.
\item $X' \to S'$는 평탄하고 국소 유한표현이다.
\item $f$는 평탄하다.
\item $X = S \times_{S'} X'$이고 $Y = S \times_{S'} Y'$이다.
\end{enumerate}
그러면 $f'$는 평탄하고, $f'$의 상에 있는 모든 $y' \in Y'$에 대해
국소환 $\mathcal{O}_{Y', y'}$는 $\mathcal{O}_{S', s'}$ 위에서 평탄하고
본질적으로 유한표현이다.
\end{lemma}

\begin{proof}
Algebra, Lemma
\ref{algebra-lemma-criterion-flatness-fibre-locally-nilpotent}의
직접적인 귀결이다.
\end{proof}

\noindent
스킴의 사상이 갖는 많은 성질은 위 보조정리와 같은 평탄 변형 아래에서 보존된다.

\begin{lemma}
\label{more-morphisms-lemma-deform-property-fp-over-ft}
비후화의 가환 그림
$$
\xymatrix{
(X \subset X') \ar[rr]_{(f, f')} \ar[rd] & & (Y \subset Y') \ar[ld] \\
& (S \subset S')
}
$$
을 생각하자. $Y' \to S'$는 국소 유한형이고, $X' \to S'$는 평탄하고
국소 유한표현이며, $X = S \times_{S'} X'$이고
$Y = S \times_{S'} Y'$라고 가정하자. 그러면 다음이 성립한다.
\begin{enumerate}
\item $f$가 평탄일 필요충분조건은 $f'$가 평탄인 것이다.
\label{more-morphisms-item-flat-fp-over-ft}
\item $f$가 동형일 필요충분조건은 $f'$가 동형인 것이다.
\label{more-morphisms-item-isomorphism-fp-over-ft}
\item $f$가 열린 몰입일 필요충분조건은 $f'$가 열린 몰입인 것이다.
\label{more-morphisms-item-open-immersion-fp-over-ft}
\item $f$가 준콤팩트일 필요충분조건은 $f'$가 준콤팩트인 것이다.
\label{more-morphisms-item-quasi-compact-fp-over-ft}
\item $f$가 보편적으로 닫혔을 필요충분조건은 $f'$가 보편적으로 닫힌 것이다.
\label{more-morphisms-item-universally-closed-fp-over-ft}
\item $f$가 (준)분리일 필요충분조건은 $f'$가 (준)분리인 것이다.
\label{more-morphisms-item-separated-fp-over-ft}
\item $f$가 모노사상일 필요충분조건은 $f'$가 모노사상인 것이다.
\label{more-morphisms-item-monomorphism-fp-over-ft}
\item $f$가 전사일 필요충분조건은 $f'$가 전사인 것이다.
\label{more-morphisms-item-surjective-fp-over-ft}
\item $f$가 보편 단사일 필요충분조건은 $f'$가 보편 단사인 것이다.
\label{more-morphisms-item-universally-injective-fp-over-ft}
\item $f$가 아핀일 필요충분조건은 $f'$가 아핀인 것이다.
\label{more-morphisms-item-affine-fp-over-ft}
\item $f$가 국소 준유한일 필요충분조건은 $f'$가 국소 준유한인 것이다.
\label{more-morphisms-item-quasi-finite-fp-over-ft}
\item
\label{more-morphisms-item-relative-dimension-d-fp-over-ft}
$f$가 상대차원 $d$인 국소 유한형일 필요충분조건은 $f'$가
상대차원 $d$인 국소 유한형인 것이다.
\item $f$가 보편적으로 열렸을 필요충분조건은 $f'$가 보편적으로 열린 것이다.
\label{more-morphisms-item-universally-open-fp-over-ft}
\item $f$가 신토믹일 필요충분조건은 $f'$가 신토믹인 것이다.
\label{more-morphisms-item-syntomic-fp-over-ft}
\item $f$가 매끄러울 필요충분조건은 $f'$가 매끄러운 것이다.
\label{more-morphisms-item-smooth-fp-over-ft}
\item $f$가 비분기일 필요충분조건은 $f'$가 비분기인 것이다.
\label{more-morphisms-item-unramified-fp-over-ft}
\item $f$가 에탈일 필요충분조건은 $f'$가 에탈인 것이다.
\label{more-morphisms-item-etale-fp-over-ft}
\item $f$가 고유일 필요충분조건은 $f'$가 고유인 것이다.
\label{more-morphisms-item-proper-fp-over-ft}
\item $f$가 유한일 필요충분조건은 $f'$가 유한인 것이다.
\label{more-morphisms-item-finite-fp-over-ft}
\item
\label{more-morphisms-item-finite-locally-free-fp-over-ft}
$f$가 (계수 $d$의) 유한 국소 자유일 필요충분조건은 $f'$가
(계수 $d$의) 유한 국소 자유인 것이고,
% P05-EN-CH37-0045
\item 여기에 더 추가할 것.
\end{enumerate}
\end{lemma}

\begin{proof}
$X$와 $Y$에 대한 가정은 $f$가 $X \to X'$에 의한 $f'$의 밑변환이라는
뜻이다. 위 (1) -- (20)에 열거한 성질 $\mathcal{P}$는 모두 밑변환 아래에서
안정적이므로, $f'$가 성질 $\mathcal{P}$를 가지면 $f$도 그 성질을 갖는다.
다음을 보라.
Schemes, Lemmas \ref{schemes-lemma-base-change-immersion},
\ref{schemes-lemma-quasi-compact-preserved-base-change},
\ref{schemes-lemma-separated-permanence}, 그리고
\ref{schemes-lemma-base-change-monomorphism}
그리고
Morphisms, Lemmas
\ref{morphisms-lemma-base-change-surjective},
\ref{morphisms-lemma-base-change-universally-injective},
\ref{morphisms-lemma-base-change-affine},
\ref{morphisms-lemma-base-change-quasi-finite},
\ref{morphisms-lemma-base-change-relative-dimension-d},
\ref{morphisms-lemma-base-change-syntomic},
\ref{morphisms-lemma-base-change-smooth},
\ref{morphisms-lemma-base-change-unramified},
\ref{morphisms-lemma-base-change-etale},
\ref{morphisms-lemma-base-change-proper},
\ref{morphisms-lemma-base-change-finite}, 그리고
\ref{morphisms-lemma-base-change-finite-locally-free}.

\medskip\noindent
따라서 각 경우에 흥미로운 방향은 $f$가 그 성질을 갖는다고 가정하고
$f'$도 그 성질을 가짐을 보이는 방향이다. 아래의 논증에서 더 언급하지
않고 사용할 몇 가지 일반적인 주의를 먼저 둔다.
(I) $W' \subset S'$를 아핀 열린 부분스킴이라 하고, $U' \subset X'$와
$V' \subset Y'$를 $W'$ 위에 놓인 아핀 열린 부분스킴으로서
$f'(U') \subset V'$를 만족한다고 하자. $W' = \Spec(R')$라 하고,
% P05-EN-CH37-0046
$I \subset R'$를 닫힌 부분스킴 $W' \cap S$를 정의하는 아이디얼이라 하자.
$U' = \Spec(B')$와 $V' = \Spec(A')$라 쓰자. 그러면 가환 그림
$$
\xymatrix{
0 \ar[r] &
IB' \ar[r] &
B' \ar[r] &
B \ar[r] & 0 \\
0 \ar[r] &
IA' \ar[r] \ar[u] &
A' \ar[r] \ar[u] &
A \ar[r] \ar[u] & 0
}
$$
을 얻고 각 행은 완전하다.
(II) 사상 $X \to X'$와 $Y \to Y'$는 보편 위상동형이다. 따라서 임의의
밑변환 뒤에도 사상 $f$와 $f'$의 위상은 동일하다.
(III) $f$가 평탄하면 $f'$도 평탄하고, $f'$의 상에 있는 모든 점에서
$Y' \to S'$가 평탄하다. 다음을 보라.
% P05-EN-CH37-0047
Lemma \ref{more-morphisms-lemma-flatness-morphism-thickenings}.

\medskip\noindent
(\ref{more-morphisms-item-flat-fp-over-ft})에 관하여. 이는 일반 주의 (III)이다.

\medskip\noindent
(\ref{more-morphisms-item-isomorphism-fp-over-ft})에 관하여. $f$가 동형이라고 가정하자.
아핀 열린 부분스킴 $V' \subset Y'$를 택하고 $U' = (f')^{-1}(V')$로 두자.
그러면 $V = Y \cap V'$가 아핀이므로
$V \cong f^{-1}(V) = U = Y \times_{Y'} U'$도 아핀이다. Lemma
\ref{more-morphisms-lemma-thickening-affine-scheme}에 의해 $U'$는 아핀이다. 따라서 일반
주의 (I)의 그림을 얻는다. Algebra, Lemma
\ref{algebra-lemma-isomorphism-modulo-locally-nilpotent}에 의해
$A' \to B'$는 동형, 즉 $U' \cong V'$이다. 따라서 $f'$는 동형이다.

\medskip\noindent
(\ref{more-morphisms-item-open-immersion-fp-over-ft})에 관하여. $f$가 열린 몰입이라고
가정하자. 그러면 $f$는 $X$와 어떤 열린 부분스킴 $V \subset Y$ 사이의
동형이다. $V' \subset Y'$를 바탕 위상공간이 $V$인 열린 부분스킴이라 하자.
그러면 $f'$는 $X'$에서 $V'$로 가는 사상이고
(\ref{more-morphisms-item-isomorphism-fp-over-ft})에 의해 동형이다. 따라서 $f'$는 열린 몰입이다.

\medskip\noindent
(\ref{more-morphisms-item-quasi-compact-fp-over-ft})에 관하여. 주의 (II)에서 바로 따른다.
더 일반적인 명제는 Lemma \ref{more-morphisms-lemma-thicken-property-morphisms}도 보라.

\medskip\noindent
(\ref{more-morphisms-item-universally-closed-fp-over-ft})에 관하여. 주의 (II)에서 바로
따른다. 더 일반적인 명제는 Lemma
\ref{more-morphisms-lemma-thicken-property-morphisms}도 보라.

\medskip\noindent
(\ref{more-morphisms-item-separated-fp-over-ft})에 관하여. 다음에 유의하자.
$X \times_Y X = Y \times_{Y'} (X' \times_{Y'} X')$이므로
$X' \times_{Y'} X'$는 $X \times_Y X$의 비후화이다. 따라서 사상
$\Delta_{X/Y}$와 $\Delta_{X'/Y'}$의 위상이 일치하므로 결론을 얻는다.
더 일반적인 명제는 Lemma \ref{more-morphisms-lemma-thicken-property-morphisms}도 보라.

\medskip\noindent
(\ref{more-morphisms-item-monomorphism-fp-over-ft})에 관하여. $f$가 모노사상이라고
가정하자. 대각 사상 $\Delta_{X'/Y'} : X' \to X' \times_{Y'} X'$를
생각하자. $X' \times_{Y'} X' \to S'$는 국소 유한형임에 유의하자.
$S \to S'$에 의한 $\Delta_{X'/Y'}$의 밑변환은 $\Delta_{X/Y}$이고,
이는 가정에 의해 동형이다. (\ref{more-morphisms-item-isomorphism-fp-over-ft})에 의해
$\Delta_{X'/Y'}$도 동형이다.

\medskip\noindent
(\ref{more-morphisms-item-surjective-fp-over-ft})에 관하여. 이는 명백하다. 더 일반적인
명제는 Lemma \ref{more-morphisms-lemma-thicken-property-morphisms}도 보라.

\medskip\noindent
(\ref{more-morphisms-item-universally-injective-fp-over-ft})에 관하여. 주의 (II)에서 바로
따른다. 더 일반적인 명제는 Lemma
\ref{more-morphisms-lemma-thicken-property-morphisms}도 보라.

\medskip\noindent
(\ref{more-morphisms-item-affine-fp-over-ft})에 관하여. $f$가 아핀이라고 가정하자.
아핀 열린 부분스킴 $V' \subset Y'$를 택하고 $U' = (f')^{-1}(V')$로 두자.
그러면 $V = Y \cap V'$가 아핀이므로 $U = Y \times_{Y'} U'$도 아핀이다.
Lemma \ref{more-morphisms-lemma-thickening-affine-scheme}에 의해 $U'$는 아핀이다.
따라서 $f'$는 아핀이다. 더 일반적인 명제는 Lemma
\ref{more-morphisms-lemma-thicken-property-morphisms}도 보라.

\medskip\noindent
(\ref{more-morphisms-item-quasi-finite-fp-over-ft})에 관하여. Morphisms, Lemma
\ref{morphisms-lemma-permanence-finite-type}에 의해 $f'$가 국소 유한형이라는
사실과, 유한형 사상의 준유한성은 올 위에서 검사할 수 있다는 사실로부터
따른다. Morphisms, Lemma
\ref{morphisms-lemma-quasi-finite-at-point-characterize}를 보라.

\medskip\noindent
(\ref{more-morphisms-item-relative-dimension-d-fp-over-ft})에 관하여. 일반 주의 (II)와
$f'$가 국소 유한형이라는 사실(Morphisms, Lemma
\ref{morphisms-lemma-permanence-finite-type})로부터 따른다.

\medskip\noindent
(\ref{more-morphisms-item-universally-open-fp-over-ft})에 관하여. 일반 주의 (II)에서 바로
따른다. 더 일반적인 명제는 Lemma
\ref{more-morphisms-lemma-thicken-property-morphisms}도 보라.

\medskip\noindent
(\ref{more-morphisms-item-syntomic-fp-over-ft})에 관하여. $f$가 신토믹이라고 가정하자.
Morphisms, Lemma \ref{morphisms-lemma-finite-presentation-permanence}에 의해
$f'$는 국소 유한표현이다. 일반 주의 (III)에 의해 $f'$는 평탄하다.
$f'$의 올들은 $f$의 올들이다. 따라서 Morphisms, Lemma
\ref{morphisms-lemma-syntomic-flat-fibres}에 의해 $f'$는 신토믹이다.

\medskip\noindent
(\ref{more-morphisms-item-smooth-fp-over-ft})에 관하여. $f$가 매끄럽다고 가정하자.
Morphisms, Lemma \ref{morphisms-lemma-finite-presentation-permanence}에 의해
$f'$는 국소 유한표현이다. 일반 주의 (III)에 의해 $f'$는 평탄하다.
$f'$의 올들은 $f$의 올들이다. 따라서 Morphisms, Lemma
\ref{morphisms-lemma-smooth-flat-smooth-fibres}에 의해 $f'$는 매끄럽다.

\medskip\noindent
(\ref{more-morphisms-item-unramified-fp-over-ft})에 관하여. $f$가 비분기라고 가정하자.
Morphisms, Lemma \ref{morphisms-lemma-permanence-finite-type}에 의해
$f'$는 국소 유한형이다. $f'$의 올들은 $f$의 올들이다. 따라서 Morphisms,
Lemma \ref{morphisms-lemma-unramified-etale-fibres}에 의해 $f'$는 비분기이다.

\medskip\noindent
(\ref{more-morphisms-item-etale-fp-over-ft})에 관하여. $f$가 에탈이라고 가정하자.
Morphisms, Lemma \ref{morphisms-lemma-finite-presentation-permanence}에 의해
$f'$는 국소 유한표현이다. 일반 주의 (III)에 의해 $f'$는 평탄하다.
$f'$의 올들은 $f$의 올들이다. 따라서 Morphisms, Lemma
\ref{morphisms-lemma-etale-flat-etale-fibres}에 의해 $f'$는 에탈이다.

\medskip\noindent
(\ref{more-morphisms-item-proper-fp-over-ft})에 관하여.
(\ref{more-morphisms-item-separated-fp-over-ft}), $f$가 국소 유한형이라는 사실
(Morphisms, Lemma \ref{morphisms-lemma-permanence-finite-type}),
(\ref{more-morphisms-item-quasi-compact-fp-over-ft}), 그리고
(\ref{more-morphisms-item-universally-closed-fp-over-ft})를 결합하면 따른다.

\medskip\noindent
(\ref{more-morphisms-item-finite-fp-over-ft})에 관하여.
(\ref{more-morphisms-item-universally-closed-fp-over-ft}),
(\ref{more-morphisms-item-affine-fp-over-ft}), Morphisms, Lemma
\ref{morphisms-lemma-integral-universally-closed}, $f$가 국소 유한형이라는
사실(Morphisms, Lemma \ref{morphisms-lemma-permanence-finite-type}), 그리고
Morphisms, Lemma \ref{morphisms-lemma-finite-integral}을 결합한다.

\medskip\noindent
(\ref{more-morphisms-item-finite-locally-free-fp-over-ft})에 관하여. $f$가 유한 국소
자유라고 가정하자. (\ref{more-morphisms-item-finite-fp-over-ft})에 의해 $f'$는 유한이다.
일반 주의 (III)에 의해 $f'$는 평탄하다. Morphisms, Lemma
\ref{morphisms-lemma-finite-presentation-permanence}에 의해 $f'$는
국소 유한표현이다. 따라서 Morphisms, Lemma
\ref{morphisms-lemma-finite-flat}에 의해 $f'$는 유한 국소 자유이다.
\end{proof}

\begin{lemma}[사영 스킴의 변형]
\label{more-morphisms-lemma-deform-projective}
$f : X \to S$를 고유이고 평탄하며 유한표현인 스킴의 사상이라 하자.
$\mathcal{L}$이 $f$-풍부라고 하자. $S$가 준콤팩트라고 가정하자.
다음을 만족하는 $d_0 \geq 0$이 존재한다. 모든 데카르트 그림
$$
\vcenter{
\xymatrix{
X \ar[r]_{i'} \ar[d]_f & X' \ar[d]^{f'} \\
S \ar[r]^i & S'
}
}
\quad\text{그리고}\quad
\begin{matrix}
\text{가역 }\mathcal{O}_{X'}\text{-가군}\\
\mathcal{L}'\text{이며 }\mathcal{L} \cong (i')^*\mathcal{L}'
\end{matrix}
$$
에 대해, 여기서 $S \subset S'$는 비후화이고 $f'$는 고유이고 평탄하며
유한표현이라고 하자. 그러면 다음이 성립한다.
\begin{enumerate}
\item $R^p(f')_*(\mathcal{L}')^{\otimes d} = 0$
는 모든 $p > 0$과 $d \geq d_0$에 대해 성립한다.
\item $\mathcal{A}'_d = (f')_*(\mathcal{L}')^{\otimes d}$
는 $d \geq d_0$에 대해 유한 국소 자유이다.
\item $\mathcal{A}' =
\mathcal{O}_{S'} \oplus \bigoplus_{d \geq d_0} \mathcal{A}'_d$
는 유한표현인 준연접 $\mathcal{O}_{S'}$-대수이다.
\item 표준 동형
$r' : X' \to \underline{\text{Proj}}_{S'}(\mathcal{A}')$이 존재한다.
\item 표준 동형
$\theta' : (r')^*\mathcal{O}_{\underline{\text{Proj}}_{S'}(\mathcal{A}')}(1)
\to \mathcal{L}'$이 존재한다.
\end{enumerate}
$\mathcal{A}'$, $r'$, $\theta'$의 구성은 자료
$(X', S', i, i', f', \mathcal{L}')$에 관하여 함자적이다.
\end{lemma}

\begin{proof}
먼저 사상 $r'$와 $\theta'$를 설명한다. $\mathcal{L}'$은 $f'$-풍부임에
유의하자. Lemma \ref{more-morphisms-lemma-thicken-property-relatively-ample}를 보라.
준연접 등급 $\mathcal{O}_{S'}$-대수의 표준 사상
$\mathcal{A}' \to \bigoplus_{d \geq 0} (f')_*(\mathcal{L}')^{\otimes d}$
이 존재하며 차수 $\geq d_0$에서는 동형이다. 따라서 이는 구조층의 Serre
꼬임과 양립하는 상대 Proj 위의 동형을 유도한다. Constructions, Lemma
\ref{constructions-lemma-eventual-iso-graded-rings-map-relative-proj}를 보라.
그러므로 Morphisms, Lemma
\ref{morphisms-lemma-characterize-relatively-ample}에 의해 사상 $r'$를 얻고
(이 보조정리는 다시 Constructions, Lemma
\ref{constructions-lemma-invertible-map-into-relative-proj}의 구성을 사용한다),
Morphisms, Lemma \ref{morphisms-lemma-proper-ample-is-proj}에 의해 이는
동형이다. 사상 $\theta'$는 Constructions, Lemma
\ref{constructions-lemma-invertible-map-into-relative-proj}에서 얻는다.
Properties, Lemma \ref{properties-lemma-ample-gcd-is-one}에 의해 $\theta'$가
동형임을 알 수 있다(여기서는 $S'$의 아핀 열린 부분스킴 위의 사상 $r'$가
Properties, Lemma \ref{properties-lemma-map-into-proj}의 사상과 같다는
사실도 사용한다. 이는 Morphisms, Lemma
\ref{morphisms-lemma-proper-ample-is-proj}의 증명에서 설명되어 있다).

\medskip\noindent
(1)과 (2)에 서술된 소멸과 국소 자유성을 가정하면, 구성의 함자성은 다음과
같이 볼 수 있다. $h : T \to S'$가 스킴의 사상이라고 하자. $f'$의 밑변환을
$f_T : X'_T \to T$라 하고,
% P05-EN-CH37-0048
$\mathcal{L}$의 $X'_T$로의 당김을 $\mathcal{L}_T$라 하자. 코호몰로지와
밑변환에 의해(예를 들어 Derived Categories of Schemes, Lemma
\ref{perfect-lemma-compare-base-change}의 정식화를 사용하면) $T$ 위에서도
대응하는 소멸이 성립하고, 더욱이
$h^*\mathcal{A}'_d = f_{T, *}\mathcal{L}_T^{\otimes d}$
이다(따라서 밀어앞선 것들의 국소 자유성과 대응하는 등급
$\mathcal{O}_T$-대수 $\mathcal{A}_T$의 유한 생성성도 성립한다).
그러므로 사상
% P05-EN-CH37-0049
$r_T : X_T \to
\underline{\text{Proj}}_T(\bigoplus f_{T, *}\mathcal{L}_T^{\otimes d})$
은 단순히 $r'$의 $T$로의 밑변환이고, $\theta'$의 당김은 사상
$\theta_T$이다.

\medskip\noindent
이상을 모두 확인했으므로 (1), (2), (3)을 증명하면 충분하다.
$R^pf_*\mathcal{L}^{\otimes d} = 0$을 모든 $d \geq d_0$과 $p > 0$에 대해
만족하는 $d_0$를 택하자.
Cohomology of Schemes, Lemma
\ref{coherent-lemma-coherent-proper-ample}를 보라. 이 $d_0$가 원하는
것이라고 주장한다.

\medskip\noindent
코호몰로지와 밑변환(Derived Categories of Schemes, Lemma
\ref{perfect-lemma-flat-proper-perfect-direct-image-general})에 의해
$E'_d = Rf'_*(\mathcal{L}')^{\otimes d}$는 $D(\mathcal{O}_{S'})$의 완전
대상이고 그 구성은 임의의 밑변환과 가환한다. 특히
$E_d = Li^*E'_d = Rf_*\mathcal{L}^{\otimes d}$이다. Derived Categories of
Schemes, Lemma \ref{perfect-lemma-vanishing-implies-locally-free}에 의해
$d \geq d_0$일 때 복합체 $E_d$는 코호몰로지 차수 $0$에 놓인 유한 국소
자유 $\mathcal{O}_S$-가군 $f_*\mathcal{L}^{\otimes d}$와 동형이다. 그러면
Derived Categories of Schemes, Lemma
\ref{perfect-lemma-open-where-cohomology-in-degree-i-rank-r}에 의해 $E'_d$는
코호몰로지 차수 $0$에 놓인 유한 국소 자유 가군과 동형이다. 물론 이는
$E'_d = \mathcal{A}'_d[0]$이고,
$R^pf'_*(\mathcal{L}')^{\otimes d} = 0$이 $p > 0$에 대해 성립하며,
$\mathcal{A}'_d$가 유한 국소 자유라는 뜻이다. 이로써 (1)과 (2)를 증명했다.

\medskip\noindent
마지막으로 보여야 할 것은 $\mathcal{A}'$가 $\mathcal{O}_{S'}$-대수의 층으로서
유한표현이라는 것이다(이 개념은 Properties, Section
\ref{properties-section-extending-quasi-coherent-sheaves}에서 도입했다).
$U' = \Spec(R') \subset S'$를 아핀 열린 부분스킴이라 하자. 그러면
$A' = \mathcal{A}'(U')$는 등급 $R'$-대수이고 그 등급 성분들은 유한 사영
$R'$-가군이다. $A'$가 유한표현 $R'$-대수임을 보여야 한다. 이를 Noetherian
경우로 환원하여 증명하겠다. 즉, 밑변환이
$(X'_{U'}, \mathcal{L}'|_{X'_{U'}})$인 유한형 $\mathbf{Z}$-부분대수
$R'_0 \subset R'$와 $R'_0$ 위의 쌍\footnote{$X' \to S'$와
$\mathcal{L}'$이 갖는 것과 같은 성질을 갖는다는 뜻이다. 즉,
$X'_0 \to \Spec(R'_0)$는 평탄하고 고유이며 $\mathcal{L}'_0$은 풍부하다.}
$(X'_0, \mathcal{L}'_0)$를 찾을 수 있다. 다음을 보라.
Limits, Lemmas
\ref{limits-lemma-descend-modules-finite-presentation},
\ref{limits-lemma-descend-invertible-modules},
\ref{limits-lemma-eventually-proper},
\ref{limits-lemma-descend-flat-finite-presentation}, 그리고
\ref{limits-lemma-limit-ample}.
Cohomology of Schemes, Lemma \ref{coherent-lemma-coherent-proper-ample}에 의해
$A'_0 = \bigoplus_{d \geq 0} H^0(X'_0, (\mathcal{L}'_0)^{\otimes d})$
는 유한 생성 등급 $R'_0$-대수이고, 다음을 만족하는 $d'_0$이 존재한다.
$H^p(X'_0, (\mathcal{L}'_0)^{\otimes d}) = 0$, $p > 0$이고 $d \geq d'_0$일 때.
위의 논증을
$X'_0 \to \Spec(R'_0)$와 $\mathcal{L}'_0$에 적용하면 $(A'_0)_d$가 유한
사영 $R'_0$-가군이고 다음 등식이 성립함을 알 수 있다.
$$
A'_d = \mathcal{A}'_d(U') =
H^0(X'_{U'}, (\mathcal{L}')^{\otimes d}|_{X'_{U'}}) =
H^0(X'_0, (\mathcal{L}'_0)^{\otimes d}) \otimes_{R'_0} R' =
(A'_0)_d \otimes_{R'_0} R'
$$
이는 $d \geq d'_0$에 대해 성립한다. 이제 논증에 한 가지 작은 문제가 있다.
$d'_0$를 $d_0$와 같게 택할 수 있는지 알 수 없다는 점이다\footnote{실제로는
극한 논증을 더 수행하여 이 경우로 환원할 수 있다.}. 이를 피하기 위해 다음
일련의 논증으로 증명을 마무리한다.
\begin{enumerate}
\item[(a)] 대수
$B = R'_0 \oplus \bigoplus_{d \geq \max(d_0, d'_0)} (A'_0)_d$
는 유한형 $R'_0$-대수이다. $B \subset A'_0$에 Artin--Tate 보조정리를
적용하라. Algebra, Lemma \ref{algebra-lemma-Artin-Tate}를 보라.
\item[(b)] $R'_0$가 Noetherian이므로 $B$는 유한표현 $R'_0$-대수이다.
\item[(c)] 텐서곱의 오른쪽 완전성에 의해
$B \otimes_{R'_0} R'$는 유한표현 $R'$-대수이다.
\item[(d)] 위에 표시한 등식들에 의해 이는 정확히
$C = R' \oplus \bigoplus_{d \geq \max(d_0, d'_0)} A'_d$
가 유한표현 $R'$-대수라는 뜻이다.
% P05-EN-CH37-0050
\item[(e)] 몫 $A'/C$는 유한 사영 $R'$-가군 $A'_d$,
$d_0 \leq d \leq \max(d_0, d'_0)$의 직합이므로 유한표현 $R'$-가군이다.
\item[(f)] Algebra, Lemma
\ref{algebra-lemma-finitely-presented-over-subring}에 의해 몫 $A'/C$는
유한표현 $C$-가군이다.
\item[(g)] 따라서 Algebra, Lemma \ref{algebra-lemma-extension}에 의해
$A'$는 유한표현 $C$-가군이다.
\item[(h)] Algebra, Lemma \ref{algebra-lemma-finite-finite-type}에 의해
이는 $A'$가 유한표현 $C$-대수임을 함의한다.
\item[(i)] 마지막으로 Algebra, Lemma
\ref{algebra-lemma-compose-finite-type}을 $R' \to C \to A'$에 적용하면
$A'$가 유한표현 $R'$-대수임을 얻는다.
\end{enumerate}
이로써 증명이 끝난다.
\end{proof}









\section{형식적으로 매끄러운 사상}
\label{more-morphisms-section-formally-smooth}

\noindent
매끄러움의 미분 판정법(예를 들어 Morphisms, Lemma
\ref{morphisms-lemma-smooth-at-point})에 대한 Michael Artin의 입장은 그것들이
(실제로는) 기본적으로 쓸모없다는 것이다. 이 절에서는 형식적으로 매끄러운
사상 $X \to S$의 개념을 도입한다. 이러한 사상은 $T$가 아핀일 때 $X$의
$T$-값 점들이 $T$의 무한소 비후화로 올라간다는 성질로 특징지어진다.
주요 결과는 형식적으로 매끄럽고 국소 유한표현인 사상은 매끄럽다는 것이다.
Lemma \ref{more-morphisms-lemma-smooth-formally-smooth}를 보라. 이 판정법은 위에서 언급한
미분 판정법보다 흔히 사용하기 쉽다는 것이 밝혀진다.

\medskip\noindent
환 사상 $R \to A$가 {\it 형식적으로 매끄럽다}는 것은 다음과 같은 모든
실선으로 이루어진 가환 그림에 대해
$$
\xymatrix{
A \ar[r] \ar@{-->}[rd] & B/I \\
R \ar[r] \ar[u] & B \ar[u]
}
$$
$I \subset B$가 제곱영 아이디얼일 때 그림을 가환으로 만드는 점선 화살표가
존재한다는 뜻임을 상기하자(Algebra, Definition
\ref{algebra-definition-formally-smooth} 참조). 이는 스킴의 사상에 대한
다음 유사 개념의 동기가 된다.

\begin{definition}
\label{more-morphisms-definition-formally-smooth}
$f : X \to S$를 스킴의 사상이라 하자. 다음과 같은 임의의 실선 가환 그림에 대해
$$
\xymatrix{
X \ar[d]_f & T \ar[d]^i \ar[l] \\
S & T' \ar[l] \ar@{-->}[lu]
}
$$
$T \subset T'$가 $S$ 위의 아핀 스킴의 1차 비후화일 때 그림을 가환으로 만드는
점선 화살표가 존재하면 $f$가 {\it 형식적으로 매끄럽다}고 한다.
\end{definition}

\noindent
형식적으로 비분기인 사상과 형식적으로 에탈인 사상의 경우에는 $T'$가
아핀이라는 조건을 없앨 수 있었다. Lemmas
\ref{more-morphisms-lemma-formally-unramified-not-affine}와
\ref{more-morphisms-lemma-formally-etale-not-affine}를 보라. 형식적으로 매끄러운 사상의
경우에는 더 이상 그렇지 않다. 사실 점선 화살표를 $T'$ 위에서 Zariski
국소적으로 채울 수 있어야 한다는 조건이 조금 더 자연스러울 것이다.
Lemma \ref{more-morphisms-lemma-formally-smooth}의 증명을 분석하면 이것이 현재의 정의와
동치임을 알 수 있다. 특히 형식적으로 매끄러운 것은 시작점 위에서 Zariski
국소적이다(사실 시작점 위에서 매끄러운 위상에 관하여 국소적이다.
% P05-EN-CH37-0051
여기에 미래의 참고문헌을 삽입할 것).

\begin{lemma}
\label{more-morphisms-lemma-composition-formally-smooth}
형식적으로 매끄러운 사상들의 합성은 형식적으로 매끄럽다.
\end{lemma}

\begin{proof}
생략한다.
\end{proof}

\begin{lemma}
\label{more-morphisms-lemma-base-change-formally-smooth}
형식적으로 매끄러운 사상의 밑변환은 형식적으로 매끄럽다.
\end{lemma}

\begin{proof}
생략한다. 다만 대수적 판은 Algebra, Lemma
\ref{algebra-lemma-base-change-fs}를 보라.
\end{proof}

\begin{lemma}
\label{more-morphisms-lemma-formally-etale-unramified-smooth}
$f : X \to S$를 스킴의 사상이라 하자. 그러면 $f$가 형식적으로 에탈일
필요충분조건은 $f$가 형식적으로 매끄럽고 형식적으로 비분기인 것이다.
\end{lemma}

\begin{proof}
생략한다.
\end{proof}

\begin{lemma}
\label{more-morphisms-lemma-formally-smooth-on-opens}
$f : X \to S$를 스킴의 사상이라 하자. $U \subset X$와 $V \subset S$를
$f(U) \subset V$를 만족하는 열린 부분스킴이라 하자. $f$가 형식적으로
매끄러우면 $f|_U : U \to V$도 형식적으로 매끄럽다.
\end{lemma}

\begin{proof}
Definition \ref{more-morphisms-definition-formally-smooth}에서와 같은 실선 그림
$$
\xymatrix{
U \ar[d]_{f|_U} & T \ar[d]^i \ar[l]^a \\
V & T' \ar[l] \ar@{-->}[lu]
}
$$
을 생각하자. $f$가 형식적으로 매끄러우면 $a'|_T = a$를 만족하는
$S$-사상 $a' : T' \to X$가 존재한다. $T$와 $T'$의 바탕 집합은 같으므로
$a'$가 $U$로 가는 사상임을 알 수 있다(Schemes, Section
\ref{schemes-section-open-immersion} 참조). 또한 이는 분명히 $V$-사상이다.
따라서 원하는 위 점선 화살표를 얻는다.
% P05-EN-CH37-0052
\end{proof}

\begin{lemma}
\label{more-morphisms-lemma-affine-formally-smooth}
$f : X \to S$를 스킴의 사상이라 하고 $X$와 $S$가 아핀이라고 가정하자.
그러면 $f$가 형식적으로 매끄러울 필요충분조건은
$\mathcal{O}_S(S) \to \mathcal{O}_X(X)$가 형식적으로 매끄러운 환 사상인 것이다.
\end{lemma}

\begin{proof}
환과 아핀 스킴의 범주의 동치에 의해 이는 정의들(Definition
\ref{more-morphisms-definition-formally-smooth}와 Algebra, Definition
\ref{algebra-definition-formally-smooth})에서 바로 따른다. Schemes, Lemma
\ref{schemes-lemma-category-affine-schemes}를 보라.
\end{proof}

\noindent
다음 보조정리가 이 절의 주요 결과이다. 이는 함수자적 관점의 승리라고 할 수
있다. 실제로 이 결과를 Limits, Proposition
\ref{limits-proposition-characterize-locally-finite-presentation}과 결합하면
사상 $f : X \to S$가 매끄러운지를 함수자
$h_X : \Sch/S \to \textit{Sets}$의 ``단순한'' 성질들로 판별할 수 있다.

\begin{lemma}[무한소 올림 판정법]
\label{more-morphisms-lemma-smooth-formally-smooth}
$f : X \to S$를 스킴의 사상이라 하자. 다음은 동치이다.
\begin{enumerate}
\item 사상 $f$는 매끄럽다.
\item 사상 $f$는 국소 유한표현이고 형식적으로 매끄럽다.
\end{enumerate}
\end{lemma}

\begin{proof}
$f : X \to S$가 국소 유한표현이고 형식적으로 매끄럽다고 가정하자.
$f(U) \subset V$를 만족하는 한 쌍의 아핀 열린집합
$\Spec(A) = U \subset X$와 $\Spec(R) = V \subset S$를 생각하자.
Lemma \ref{more-morphisms-lemma-formally-smooth-on-opens}에 의해 $U \to V$는 형식적으로
매끄럽다. Lemma \ref{more-morphisms-lemma-affine-formally-smooth}에 의해 $R \to A$는
형식적으로 매끄럽다. Morphisms, Lemma
\ref{morphisms-lemma-locally-finite-presentation-characterize}에 의해
$R \to A$는 유한표현이다. Algebra, Proposition
\ref{algebra-proposition-smooth-formally-smooth}에 의해 $R \to A$는 매끄럽다.
따라서 매끄러운 사상의 정의에 의해 $X \to S$는 매끄럽다.

\medskip\noindent
역으로 $f : X \to S$가 매끄럽다고 가정하자. Definition
\ref{more-morphisms-definition-formally-smooth}에서와 같은 실선 가환 그림
$$
\xymatrix{
X \ar[d]_f & T \ar[d]^i \ar[l]^a \\
S & T' \ar[l] \ar@{-->}[lu]
}
$$
을 생각하자. 점선 화살표가 존재함을 보여 $f$가 형식적으로 매끄러움을
증명할 것이다.

\medskip\noindent
Remark \ref{more-morphisms-remark-special-case}에서 논의한 특별한 경우에 Lemma
\ref{more-morphisms-lemma-sheaf}가 주는 $T'$ 위의 집합의 층을 $\mathcal{F}$라 하자. 또한
$$
\mathcal{H} =
\SheafHom_{\mathcal{O}_T}(a^*\Omega_{X/S}, \mathcal{C}_{T/T'})
$$
를 Lemma \ref{more-morphisms-lemma-action-sheaf}에서와 같은 작용
$\mathcal{H} \times \mathcal{F} \to \mathcal{F}$을 갖는
$\mathcal{O}_T$-가군의 층이라 하자. 목표는 단지
$\mathcal{F}(T) \not = \emptyset$임을 보이는 것이다. 달리 말해
$\mathcal{F}$가 $T$ 위의 자명한 $\mathcal{H}$-토서임을 보이고자 한다
(Cohomology, Section \ref{cohomology-section-h1-torsors} 참조).
두 단계가 있다. (I) $\mathcal{F}$가 토서임을 보이려면 모든 $t \in T$에
대해 $\mathcal{F}_t \not = \emptyset$임을 보여야 한다(Cohomology,
Definition \ref{cohomology-definition-torsor} 참조). (II) $\mathcal{F}$가
자명한 토서임을 보이려면 $H^1(T, \mathcal{H}) = 0$임을 보이면 충분하다
(Cohomology, Lemma \ref{cohomology-lemma-torsors-h1} 참조. $\mathcal{H}$를
아벨 층으로 보아도, $\mathcal{O}_T$-가군으로 보아도 그 코호몰로지를 사용할
수 있다. Cohomology, Lemma \ref{cohomology-lemma-modules-abelian} 참조).

\medskip\noindent
먼저 (I)을 증명하자. 각 $t \in T$에 대해 $a(U)$가 아핀 열린집합
$\Spec(A) = W \subset X$에 포함되고 그 열린집합이 아핀 열린집합
$\Spec(R) = V \subset S$로 사상되도록 하는 $t$의 아핀 열린 근방
$U \subset T$를 택할 수 있다. Morphisms, Lemma
\ref{morphisms-lemma-smooth-characterize}에 의해 환 사상 $R \to A$는
매끄럽다. 따라서 Algebra, Proposition
\ref{algebra-proposition-smooth-formally-smooth}에 의해 환 사상
$R \to A$는 형식적으로 매끄럽다. 그러면 Lemma
\ref{more-morphisms-lemma-affine-formally-smooth}에 의해 $W \to V$는 형식적으로
매끄럽다. 따라서 $a|_U : U \to W$를 $V$-사상
$a' : U' \to W \subset X$로 올릴 수 있고, 이는
$\mathcal{F}(U) \not = \emptyset$임을 보인다.

\medskip\noindent
끝으로 (II)를 증명하자. Morphisms, Lemma
\ref{morphisms-lemma-finite-presentation-differentials}에 의해
$\Omega_{X/S}$는 유한표현이다(사실 Morphisms, Lemma
\ref{morphisms-lemma-smooth-omega-finite-locally-free}에 의해 유한 국소
자유이기도 하다). 따라서 $a^*\Omega_{X/S}$는 유한표현이다(Modules,
Lemma \ref{modules-lemma-pullback-finite-presentation} 참조). 그러므로 층
$\mathcal{H} =
\SheafHom_{\mathcal{O}_T}(a^*\Omega_{X/S}, \mathcal{C}_{T/T'})$
는 Schemes, Section \ref{schemes-section-quasi-coherent}의 논의에 의해
준연접이다. 따라서 Cohomology of Schemes, Lemma
\ref{coherent-lemma-quasi-coherent-affine-cohomology-zero}에 의해 원하는 대로
$H^1(T, \mathcal{H}) = 0$이다.
\end{proof}

\noindent
국소 사영 준연접 가군은 Properties, Section
\ref{properties-section-locally-projective}에서 정의한다.

\begin{lemma}
\label{more-morphisms-lemma-formally-smooth-sheaf-differentials}
$f : X \to Y$를 형식적으로 매끄러운 스킴의 사상이라 하자.
그러면 $\Omega_{X/Y}$는 $X$ 위에서 국소 사영이다.
\end{lemma}

\begin{proof}
$f(U) \subset V$를 만족하는 아핀 열린집합 $U \subset X$와
$V \subset Y$를 택하자. Lemma \ref{more-morphisms-lemma-formally-smooth-on-opens}에 의해
$f|_U : U \to V$는 형식적으로 매끄럽다. 따라서 Lemma
\ref{more-morphisms-lemma-affine-formally-smooth}에 의해
$\Gamma(V, \mathcal{O}_V) \to \Gamma(U, \mathcal{O}_U)$는 형식적으로
매끄러운 환 사상이다. 그러므로 Algebra, Lemma
\ref{algebra-lemma-characterize-formally-smooth-again}에 의해
$\Gamma(U, \mathcal{O}_U)$-가군
$\Omega_{\Gamma(U, \mathcal{O}_U)/\Gamma(V, \mathcal{O}_V)}$
은 사영이다. 따라서 $\Omega_{U/V}$는 국소 사영이다(Properties, Section
\ref{properties-section-locally-projective} 참조).
\end{proof}

\begin{lemma}
\label{more-morphisms-lemma-h1-is-zero}
$T$를 아핀 스킴이라 하자. $\mathcal{F}$와 $\mathcal{G}$를 준연접
$\mathcal{O}_T$-가군이라 하고
$\mathcal{H} = \SheafHom_{\mathcal{O}_T}(\mathcal{F}, \mathcal{G})$라 하자.
$\mathcal{F}$가 국소 사영이면 $H^1(T, \mathcal{H}) = 0$이다.
\end{lemma}

\begin{proof}
스킴 위의 국소 사영 층의 정의(Properties, Definition
\ref{properties-definition-locally-projective} 참조)에 의해
$\mathcal{F}$는 자유 $\mathcal{O}_T$-가군의 직합인자이다. 따라서
$\mathcal{F} = \bigoplus_{i \in I} \mathcal{O}_T$가 자유 가군이라고
가정해도 된다. 이 경우 $\mathcal{H} = \prod_{i \in I} \mathcal{G}$는
준연접 가군들의 곱이다. 아핀 스킴 위의 준연접층의 코호몰로지는 영이므로
(Cohomology of Schemes, Lemma
\ref{coherent-lemma-quasi-coherent-affine-cohomology-zero} 참조), Cohomology,
Lemma \ref{cohomology-lemma-cohomology-products}에 의해 $H^1 = 0$이다.
\end{proof}

\begin{lemma}
\label{more-morphisms-lemma-formally-smooth}
$f : X \to Y$를 스킴의 사상이라 하자. 다음은 동치이다.
\begin{enumerate}
\item $f$는 형식적으로 매끄럽다.
\item 모든 $x \in X$에 대해 $f(U) \subset V$이고
$f|_U : U \to V$가 형식적으로 매끄럽도록 하는 열린집합
$x \in U \subset X$와 $f(x) \in V \subset Y$가 존재한다.
\item $f(U) \subset V$를 만족하는 모든 아핀 열린집합의 쌍
$U \subset X$, $V \subset Y$에 대해 환 사상
$\mathcal{O}_Y(V) \to \mathcal{O}_X(U)$는 형식적으로 매끄럽다.
% P05-EN-CH37-0053
\item 아핀 열린 덮개 $Y = \bigcup V_j$와 각 $j$에 대한 아핀 열린 덮개
$f^{-1}(V_j) = \bigcup U_{ji}$가 존재하여, 모든 $j$와 $i$에 대해
$\mathcal{O}_Y(V) \to \mathcal{O}_X(U)$가 형식적으로 매끄러운 환
사상이다.
\end{enumerate}
\end{lemma}

\begin{proof}
함의 (1) $\Rightarrow$ (2), (1) $\Rightarrow$ (3), (2) $\Rightarrow$ (4)는
Lemma \ref{more-morphisms-lemma-formally-smooth-on-opens}에서 따른다.
% P05-EN-CH37-0054
함의 (3) $\Rightarrow$ (4)는 즉시 성립한다.

\medskip\noindent
(4)를 가정하자. $f$가 형식적으로 매끄럽다는 증명은 Lemma
\ref{more-morphisms-lemma-smooth-formally-smooth}의 증명 후반부와 같다. Definition
\ref{more-morphisms-definition-formally-smooth}에서와 같은 실선 가환 그림
$$
\xymatrix{
X \ar[d]_f & T \ar[d]^i \ar[l]^a \\
Y & T' \ar[l] \ar@{-->}[lu]
}
$$
을 생각하자. 점선 화살표가 존재함을 보여 $f$가 형식적으로 매끄러움을
증명할 것이다. Remark \ref{more-morphisms-remark-special-case}에서 논의한 특별한 경우에
Lemma \ref{more-morphisms-lemma-sheaf}가 주는 $T'$ 위의 집합의 층을 $\mathcal{F}$라 하자.
또한
$$
\mathcal{H} =
\SheafHom_{\mathcal{O}_T}(a^*\Omega_{X/Y}, \mathcal{C}_{T/T'})
$$
를 Lemma \ref{more-morphisms-lemma-action-sheaf}에서와 같은 작용
$\mathcal{H} \times \mathcal{F} \to \mathcal{F}$을 갖는 $T$ 위의
$\mathcal{O}_T$-가군의 층이라 하자. 작용
$\mathcal{H} \times \mathcal{F} \to \mathcal{F}$은 $\mathcal{F}$를 유사
$\mathcal{H}$-토서로 만든다(Cohomology, Definition
\ref{cohomology-definition-torsor} 참조). 목표는 $\mathcal{F}$가 자명한
$\mathcal{H}$-토서임을 보이는 것이다. 두 단계가 있다. (I) $\mathcal{F}$가
토서임을 보이려면 $\mathcal{F}$가 국소적으로 단면을 가짐을 보여야 한다.
(II) $\mathcal{F}$가 자명한 토서임을 보이려면
$H^1(T, \mathcal{H}) = 0$임을 보이면 충분하다(Cohomology, Lemma
\ref{cohomology-lemma-torsors-h1} 참조).

\medskip\noindent
먼저 (I)을 증명하자. 각 $t \in T$에 대해 어떤 $i, j$에 대하여
$a(W)$가 $U_{ji}$에 포함되도록 하는 $t$의 아핀 열린 근방
$W \subset T$를 택할 수 있다. 대응하는 열린 부분스킴을
$W' \subset T'$이라 하자. 가정 (4)에 의해
$a|_W : W \to U_{ji}$를 $V_j$-사상 $a' : W' \to U_{ji}$로 올릴 수 있고,
이는 $\mathcal{F}(W')$가 공집합이 아님을 보인다.

\medskip\noindent
끝으로 (II)를 증명하자. Lemma
\ref{more-morphisms-lemma-formally-smooth-sheaf-differentials}에 의해
% P05-EN-CH37-0055
$\Omega_{U_{ji}/V_j}$는 국소 사영이다. 따라서 Properties, Lemma
\ref{properties-lemma-locally-projective}에 의해 $\Omega_{X/Y}$는 국소
사영이다. 그러므로 Properties, Lemma
\ref{properties-lemma-locally-projective-pullback}에 의해
$a^*\Omega_{X/Y}$는 국소 사영이다. 따라서
$$
H^1(T, \mathcal{H}) =
H^1(T, \SheafHom_{\mathcal{O}_T}(a^*\Omega_{X/Y}, \mathcal{C}_{T/T'})) = 0
$$
이고, 이는 Lemma \ref{more-morphisms-lemma-h1-is-zero}에 의한 원하는 결론이다.
\end{proof}

\begin{lemma}
\label{more-morphisms-lemma-triangle-differentials-formally-smooth}
$f : X \to Y$와 $g : Y \to S$를 스킴의 사상이라 하자. $f$가
형식적으로 매끄럽다고 가정하자. 그러면
$$
0 \to f^*\Omega_{Y/S} \to \Omega_{X/S} \to \Omega_{X/Y} \to 0
$$
은 짧은 완전열이다(Morphisms, Lemma
\ref{morphisms-lemma-triangle-differentials} 참조).
\end{lemma}

\begin{proof}
이 보조정리의 대수적 판은 다음과 같다. 환 사상 $A \to B \to C$가
주어지고 $B \to C$가 형식적으로 매끄러우면 열
$$
0 \to C \otimes_B \Omega_{B/A} \to \Omega_{C/A} \to \Omega_{C/B} \to 0
$$
은 완전하다(Algebra, Lemma
\ref{algebra-lemma-exact-sequence-differentials}). 이것이 Algebra, Lemma
\ref{algebra-lemma-ses-formally-smooth}이다.
\end{proof}

\begin{lemma}
\label{more-morphisms-lemma-differentials-formally-unramified-formally-smooth}
$h : Z \to X$를 $S$ 위의 형식적으로 비분기인 스킴의 사상이라 하자.
$Z$가 $S$ 위에서 형식적으로 매끄럽다고 가정하자. 그러면 Lemma
\ref{more-morphisms-lemma-universally-unramified-differentials-sequence}의 표준 완전열
$$
0 \to \mathcal{C}_{Z/X} \to h^*\Omega_{X/S} \to \Omega_{Z/S} \to 0
$$
은 짧은 완전열이다.
\end{lemma}

\begin{proof}
$Z \to Z'$을 $X$ 위의 $Z$의 보편 1차 비후화라 하자. Lemma
\ref{more-morphisms-lemma-universally-unramified-differentials-sequence}의 증명에서 우리의
열은 다음 열과 동일시됨을 알 수 있다.
$$
\mathcal{C}_{Z/Z'} \to \Omega_{Z'/S} \otimes \mathcal{O}_Z \to
\Omega_{Z/S} \to 0.
$$
$Z \to S$가 형식적으로 매끄러우므로 $Z'$ 위에서 국소적으로 포함 사상
$Z \to Z'$의 $S$ 위 왼쪽 역사상 $Z' \to Z$를 찾을 수 있다. 따라서 이
열은 국소적으로 분할된다(Morphisms, Lemma
\ref{morphisms-lemma-differentials-relative-immersion-section} 참조).
\end{proof}

\begin{lemma}
\label{more-morphisms-lemma-two-unramified-morphisms-formally-smooth}
다음 가환 그림을 생각하자.
$$
\xymatrix{
Z \ar[r]_i \ar[rd]_j & X \ar[d]^f \\
& Y
}
$$
$i$와 $j$가 형식적으로 비분기이고 $f$가 형식적으로 매끄럽다고 하자.
그러면 Lemma \ref{more-morphisms-lemma-two-unramified-morphisms}의 표준 완전열
$$
0 \to
\mathcal{C}_{Z/Y} \to
\mathcal{C}_{Z/X} \to
i^*\Omega_{X/Y} \to 0
$$
은 완전하고 국소적으로 분할된다.
\end{lemma}

\begin{proof}
% P05-EN-CH37-0056
$Z \to Z'$을 $X$ 위의 $Z$의 보편 1차 비후화라 하고,
$Z \to Z''$을 $Y$ 위의 $Z$의 보편 1차 비후화라 하자.
% P05-EN-CH37-0057
Lemma \ref{more-morphisms-lemma-universally-unramified-differentials-sequence}에 의해 표준
사상 $Z' \to Z''$이 존재하여 다음 가환 그림을 얻는다.
$$
\xymatrix{
Z \ar[r]_{i'} \ar[rd]_{j'} & Z' \ar[r]_a \ar[d]^k & X \ar[d]^f \\
& Z'' \ar[r]^b & Y
}
$$
Lemma \ref{more-morphisms-lemma-two-unramified-morphisms}의 증명에서 위의 열을 다음 열과
동일시했다.
$$
\mathcal{C}_{Z/Z''} \to
\mathcal{C}_{Z/Z'} \to
(i')^*\Omega_{Z'/Z''} \to 0
$$
$U'' \subset Z''$을 아핀 열린집합이라 하자. 대응하는 아핀 열린 부분스킴을
$U \subset Z$와 $U' \subset Z'$이라 하자. $f$가 형식적으로 매끄러우므로
$U$ 위에서 $i$와 일치하고 $f \circ h$가 $b|_{U''}$와 같은 사상
$h : U'' \to X$가 존재한다. $Z'$이 보편 1차 비후화이므로
% P05-EN-CH37-0058
$g = a \circ h$를 만족하는 유일한 사상 $g : U'' \to Z'$을 얻는다.
$Z''$의 보편 성질에 의해 $k \circ g$는 포함 사상 $U'' \to Z''$이다.
따라서 $g$는 $k$의 왼쪽 역사상이다. 그림으로 쓰면
$$
\xymatrix{
U \ar[d] \ar[r] & Z' \ar[d]^k \\
U'' \ar[r] \ar[ru]^g & Z''
}
$$
이다. 따라서 $g$는 사상
$\mathcal{C}_{Z/Z'}|_U \to \mathcal{C}_{Z/Z''}|_U$를 유도하며, 이는
$U$ 위의 사상 $\mathcal{C}_{Z/Z''} \to \mathcal{C}_{Z/Z'}$의 왼쪽
역사상이다.
\end{proof}



































\section{Noetherian 밑 위의 매끄러움}
\label{more-morphisms-section-smooth-Noetherian}

\noindent
밑이 Noetherian이면 형식적 매끄러움의 정식화에서 더 약한 조건으로도
충분하다는 것이 밝혀진다. 어떤 의미에서 다음 보조정리들은 변형 이론의
출발점이다.

\begin{lemma}
\label{more-morphisms-lemma-lifting-along-artinian-at-point}
$f : X \to S$를 스킴의 사상이라 하고 $x \in X$라 하자.
$S$가 국소 Noetherian이고 $f$가 국소 유한형이라고 가정하자.
% P05-EN-CH37-0059
다음은 동치이다.
\begin{enumerate}
\item $f$는 $x$에서 매끄럽다.
\item 다음과 같은 모든 실선 가환 그림에 대해
$$
\xymatrix{
X \ar[d]_f & \Spec(B) \ar[d]^i \ar[l]^-\alpha \\
S & \Spec(B') \ar[l]_-{\beta} \ar@{-->}[lu]
}
$$
여기서 $B' \to B$는 국소환의 전사이고 $\Ker(B' \to B)$는 제곱영이며,
$\alpha$는 $\Spec(B)$의 닫힌점을 $x$로 보낸다. 이때 그림을 가환으로
만드는 점선 화살표가 존재한다.
\item (2)와 같되 $B' \to B$가 작은 확대 전체를 움직인다(Algebra,
Definition \ref{algebra-definition-small-extension} 참조).
\item (2)와 같되 $B' \to B$는, $\mathfrak m \subset B$가 극대 아이디얼일 때
$\alpha$가 동형사상 $\kappa(x) \to \kappa(\mathfrak m)$을 유도하는 작은
확대 전체를 움직인다.
\end{enumerate}
\end{lemma}

\begin{proof}
$f(x)$의 아핀 근방 $V \subset S$와 $f(U) \subset V$를 만족하는 $x$의
아핀 근방 $U \subset X$를 택하자. (2)에서와 같은 임의의 ``시험'' 그림에
대해 사상 $\alpha$는 $\Spec(B)$를 $U$ 안으로 보내고, 사상 $\beta$는
$\Spec(B')$를 $V$ 안으로 보낸다(Schemes, Section
\ref{schemes-section-points} 참조). 따라서 이 보조정리는 아핀 스킴 사이의
사상 $f|_U : U \to V$의 경우로 환원된다. (실제로 $V$는 Noetherian이고
$f|_U$는 유한형이다. Properties, Lemma
\ref{properties-lemma-locally-Noetherian}과 Morphisms, Lemma
\ref{morphisms-lemma-locally-finite-type-characterize}를 보라.) 이 아핀의
경우 보조정리는 Algebra, Lemma
\ref{algebra-lemma-smooth-test-artinian}와 동일하다.
\end{proof}

\noindent
때로는 유한형인 점을 ``중심으로 하는'' 작은 확대에 대해서만 올림 판정법을
확인하면 충분하다는 사실이 유용하다(Morphisms, Section
\ref{morphisms-section-points-finite-type} 참조).

\begin{lemma}
\label{more-morphisms-lemma-lifting-along-artinian}
$f : X \to S$를 스킴의 사상이라 하자. $S$가 국소 Noetherian이고 $f$가
국소 유한형이라고 가정하자.
% P05-EN-CH37-0059
다음은 동치이다.
\begin{enumerate}
\item $f$는 매끄럽다.
\item 다음과 같은 모든 실선 가환 그림에 대해
$$
\xymatrix{
X \ar[d]_f & \Spec(B) \ar[d]^i \ar[l]^-\alpha \\
S & \Spec(B') \ar[l]_-{\beta} \ar@{-->}[lu]
}
$$
여기서 $B' \to B$는 Artin 국소환의 작은 확대이고 $\beta$는 유한형이다
(!). 이때 그림을 가환으로 만드는 점선 화살표가 존재한다.
% P05-EN-CH37-0060
\end{enumerate}
\end{lemma}

\begin{proof}
$f$가 매끄러우면 무한소 올림 판정법(Lemma
\ref{more-morphisms-lemma-smooth-formally-smooth})에 의해 $f$는 형식적으로 매끄럽고
(2)가 성립한다.

\medskip\noindent
(2)를 가정하자. $f$가 매끄럽지 않은 점 $x \in X$들의 집합은 $X$의 닫힌
부분집합 $T$를 이룬다. Morphisms, Section
\ref{morphisms-section-points-finite-type}의 논의에 의해
$T \not = \emptyset$이면 사상
$$
\Spec(\kappa(x)) \to X \to S
$$
이 유한형이 되는 점 $x \in T \subset X$가 존재한다(즉 $X$의 어떤 아핀
열린집합에서 닫혀 있는 $T$의 임의의 점 $x$를 택한다). Morphisms, Lemma
\ref{morphisms-lemma-artinian-finite-type}에 의해, 잉여체가 $\kappa(x)$인
임의의 Artin 국소환 $B'$에 대해 모든 사상
$\beta : \Spec(B') \to S$는 유한형이다. 따라서 Lemma
\ref{more-morphisms-lemma-lifting-along-artinian-at-point}의 (4)에서 사용하는 모든 그림은
현재 보조정리의 (2)에서와 같은 그림에 대응한다. 그러므로
$X \to S$는 $x$에서 매끄러우며, 이는 모순이다.
% P05-EN-CH37-0061
\end{proof}

\noindent
다음은 유용한 응용이다.

\begin{lemma}
\label{more-morphisms-lemma-check-smoothness-on-infinitesimal-nbhds}
$f : X \to S$를 국소 Noetherian 스킴 사이의 유한형 사상이라 하자.
$Z \subset S$를 닫힌 부분스킴이라 하고, 그 $n$차 무한소 근방을
$Z_n \subset S$라 하자. $X_n = Z_n \times_S X$라 두자.
\begin{enumerate}
\item 모든 $n$에 대해 $X_n \to Z_n$가 매끄러우면 $f$는
$f^{-1}(Z)$의 모든 점에서 매끄럽다.
\item 모든 $n$에 대해 $X_n \to Z_n$가 에탈이면 $f$는
$f^{-1}(Z)$의 모든 점에서 에탈이다.
\end{enumerate}
\end{lemma}

\begin{proof}
모든 $n$에 대해 $X_n \to Z_n$가 매끄럽다고 가정하자. $x \in X$를
$Z$의 한 점 위에 놓인 점이라 하자. Lemma
\ref{more-morphisms-lemma-lifting-along-artinian-at-point}의 (3)에서와 같은 작은 확대
$B' \to B$와 사상 $\alpha$, $\beta$가 주어졌다고 하자. $B'$의 극대
아이디얼은 멱영이다($B'$가 Artin 환이므로). 따라서 적절한 $n$에 대해
사상 $\beta$는 $Z_n$을 경유하고 $\alpha$는 $X_n$을 경유한다. 이제
$X_n \to Z_n$의 올림 성질을 적용하면 그림에서 원하는 점선 화살표를 얻는다.
이로써 (1)을 증명했다. (2)는 (1)과, 사상이 에탈일 필요충분조건은 상대차원
$0$인 매끄러운 사상이라는 사실에서 따른다.
\end{proof}

\begin{lemma}
\label{more-morphisms-lemma-check-flatness-on-infinitesimal-nbhds}
$f : X \to S$를 국소 Noetherian 스킴 사이의 사상이라 하자.
$Z \subset S$를 닫힌 부분스킴이라 하고, 그 $n$차 무한소 근방을
$Z_n \subset S$라 하자. $X_n = Z_n \times_S X$라 두자.
모든 $n$에 대해 $X_n \to Z_n$가 평탄하면 $f$는 $f^{-1}(Z)$의 모든
점에서 평탄하다.
\end{lemma}

\begin{proof}
이는 Algebra, Lemma \ref{algebra-lemma-flat-module-powers}를 스킴의 언어로
옮긴 것이다.
\end{proof}








\section{소박한 여접 복합체}
\label{more-morphisms-section-netherlander}

\noindent
이 절은 Modules, Section \ref{modules-section-netherlander}의 연속이며,
그 절은 다시 Algebra, Section \ref{algebra-section-netherlander}의 논의를
이어간다.

\begin{definition}
\label{more-morphisms-definition-netherlander}
$f : X \to Y$를 스킴의 사상이라 하자. {\it $f$의 소박한 여접 복합체}는
Modules, Definition
\ref{modules-definition-cotangent-complex-morphism-ringed-topoi}에서 정의한
복합체이다. 기호는 $\NL_f$ 또는 $\NL_{X/Y}$를 쓴다.
\end{definition}

\begin{lemma}
\label{more-morphisms-lemma-NL-affine}
$f : X \to Y$를 스킴의 사상이라 하자. 아핀 열린집합
$\Spec(A) = U \subset X$와
% P05-EN-CH37-0062
$\Spec(R) = V \subset S$가 $f(U) \subset V$를 만족한다고 하자.
복합체의 표준 사상
$$
\widetilde{\NL_{A/R}} \longrightarrow \NL_{X/Y}|_U
$$
이 존재하며, 이는 $D(\mathcal{O}_U)$에서 동형사상이다.
\end{lemma}

\begin{proof}
Modules, Section \ref{modules-section-netherlander}의 $\NL_{X/Y}$ 구성에서
복합체의 표준 사상
$\NL_{\mathcal{O}_X(U)/f^{-1}\mathcal{O}_Y(U)} \to \NL_{X/Y}(U)$
이 있음을 알 수 있다. 이는 $A = \mathcal{O}_X(U)$-가군의 사상이며 더
제한하는 것과 양립한다. 표준 사상 $R \to f^{-1}\mathcal{O}_Y(U)$를 사용하면
복합체의 표준 사상
$\NL_{A/R} \to \NL_{\mathcal{O}_X(U)/f^{-1}\mathcal{O}_Y(U)}$
을 얻으며, 이는 $A$-가군 복합체의 사상이다. $\widetilde{\ }$ 함수자의
보편 성질(Schemes, Lemma \ref{schemes-lemma-compare-constructions} 참조)을
사용하면 보조정리의 명제에 있는 사상을 얻는다. 이 사상이 코호몰로지층에서
동형사상인지는 줄기에서 동형사상을 유도하는지를 확인하여 검사할 수 있다.
이는 Algebra, Lemmas \ref{algebra-lemma-NL-localize-bottom}와
\ref{algebra-lemma-localize-NL}, 그리고 Modules, Lemma
\ref{modules-lemma-stalk-NL}에서 따른다. 여기서는 점
$\mathfrak p \in \Spec(A)$에서 $\mathcal{O}_X$와
$f^{-1}\mathcal{O}_Y$의 줄기가 각각 $A_\mathfrak p$와 $R_\mathfrak q$이고
$\mathfrak q = R \cap \mathfrak p$라는 기술도 사용한다. 이에 대한 참조는
Schemes, Lemma \ref{schemes-lemma-spec-sheaves}와 Sheaves, Lemma
\ref{sheaves-lemma-stalk-pullback}이다.
\end{proof}

\begin{lemma}
\label{more-morphisms-lemma-netherlander-quasi-coherent}
$f : X \to Y$를 스킴의 사상이라 하자. 복합체 $\NL_{X/Y}$의 코호몰로지층은
준연접이고 차수 $-1$, $0$ 밖에서는 영이며, 차수 $0$에서는
$\Omega_{X/Y}$와 같다.
\end{lemma}

\begin{proof}
Modules, Section \ref{modules-section-netherlander}의 소박한 여접 복합체
구성에 의해 $\NL_{X/Y}$는 차수 $-1$, $0$에 놓인 복합체이고 차수 $0$의
코호몰로지는 $\Omega_{X/Y}$이다. 미분의 층은 준연접이다(Morphisms,
Lemma \ref{morphisms-lemma-differentials-diagonal}). 증명을 끝내려면
$H^{-1}(\NL_{X/Y})$가 준연접임을 보이면 충분하다. 이는 Lemma
\ref{more-morphisms-lemma-NL-affine}를 사용하여 아핀 위에서 확인하면 따른다.
\end{proof}

\begin{lemma}
\label{more-morphisms-lemma-netherlander-fp}
$f : X \to Y$를 스킴의 사상이라 하자. $f$가 국소 유한표현이면
$\NL_{X/Y}$는 $X$ 위에서 국소적으로 다음 복합체와 준동형이다.
$$
\ldots \to 0 \to \mathcal{F}^{-1} \to \mathcal{F}^0 \to 0 \to \ldots
$$
이는 준연접 $\mathcal{O}_X$-가군의 복합체이며, $\mathcal{F}^0$는
유한표현이고 $\mathcal{F}^{-1}$은 유한형이다.
\end{lemma}

\begin{proof}
Lemma \ref{more-morphisms-lemma-NL-affine}에 의해 $R \to A$가 유한표현 환 사상일 때
$\NL_{A/R}$가 이 꼴임을 보이면 충분하다. $I$가 유한생성이고
$A = R[x_1, \ldots, x_n]/I$라고 쓰자. 그러면 $I/I^2$은 유한
$A$-가군이고 $\NL_{A/R}$는
$$
\ldots \to 0 \to I/I^2 \to
\bigoplus\nolimits_{i = 1, \ldots, n} A\text{d}x_i \to 0 \to \ldots
$$
와 준동형이다. 이는 Algebra, Section \ref{algebra-section-netherlander},
특히 Algebra, Lemma \ref{algebra-lemma-NL-homotopy}에 의한다.
\end{proof}

\begin{lemma}
\label{more-morphisms-lemma-NL-formally-smooth}
$f : X \to Y$를 스킴의 사상이라 하자. 다음은 동치이다.
\begin{enumerate}
\item $f$는 형식적으로 매끄럽다.
\item $H^{-1}(\NL_{X/Y}) = 0$이고
$H^0(\NL_{X/Y}) = \Omega_{X/Y}$는 국소 사영이다.
\end{enumerate}
\end{lemma}

\begin{proof}
이는 Algebra, Proposition
\ref{algebra-proposition-characterize-formally-smooth}과 Lemma
\ref{more-morphisms-lemma-formally-smooth}에서 따른다.
\end{proof}

\begin{lemma}
\label{more-morphisms-lemma-NL-formally-etale}
$f : X \to Y$를 스킴의 사상이라 하자. 다음은 동치이다.
\begin{enumerate}
\item $f$는 형식적으로 에탈이다.
\item $H^{-1}(\NL_{X/Y}) = H^0(\NL_{X/Y}) = 0$이다.
\end{enumerate}
\end{lemma}

\begin{proof}
형식적으로 에탈인 사상은 형식적으로 매끄러우므로 Lemma
\ref{more-morphisms-lemma-NL-formally-smooth}에 의해 $H^{-1}(\NL_{X/Y}) = 0$이다.
한편 Lemma \ref{more-morphisms-lemma-characterize-formally-etale}에 의해
$\Omega_{X/Y} = 0$이다. 역으로 (2)가 성립하면 Lemma
\ref{more-morphisms-lemma-NL-formally-smooth}에 의해 $f$는 형식적으로 매끄럽고, Lemma
\ref{more-morphisms-lemma-formally-unramified-differentials}에 의해 형식적으로 비분기이다.
따라서 Lemma
% P05-EN-CH37-0063
\ref{more-morphisms-lemma-formally-etale-unramified-smooth}에 의해 형식적으로 에탈이다.
\end{proof}

\begin{lemma}
\label{more-morphisms-lemma-NL-smooth}
$f : X \to Y$를 스킴의 사상이라 하자. 다음은 동치이다.
\begin{enumerate}
\item $f$는 매끄럽다.
\item $f$는 국소 유한표현이고 $H^{-1}(\NL_{X/Y}) = 0$이며
$H^0(\NL_{X/Y}) = \Omega_{X/Y}$는 유한 국소 자유이다.
\end{enumerate}
\end{lemma}

\begin{proof}
이는 매끄러운 환 준동형의 정의(Algebra, Definition
\ref{algebra-definition-smooth}), Lemma \ref{more-morphisms-lemma-NL-affine}, 그리고
스킴의 매끄러운 사상의 정의(Morphisms, Definition
\ref{morphisms-definition-smooth})에서 따른다. 환 위의 가군에 대해서 유한
국소 자유와 유한 사영이 같다는 사실도 사용한다(Algebra, Lemma
\ref{algebra-lemma-finite-projective}).
\end{proof}

\begin{lemma}
\label{more-morphisms-lemma-NL-etale}
$f : X \to Y$를 스킴의 사상이라 하자. 다음은 동치이다.
\begin{enumerate}
\item $f$는 에탈이다.
\item $f$는 국소 유한표현이고
$H^{-1}(\NL_{X/Y}) = H^0(\NL_{X/Y}) = 0$이다.
\end{enumerate}
\end{lemma}

\begin{proof}
이는 에탈 환 준동형의 정의(Algebra, Definition
\ref{algebra-definition-etale}), Lemma \ref{more-morphisms-lemma-NL-affine}, 그리고 스킴의
에탈 사상의 정의(Morphisms, Definition
\ref{morphisms-definition-etale})에서 따른다.
\end{proof}


\begin{lemma}
\label{more-morphisms-lemma-NL-immersion}
$i : Z \to X$를 스킴의 몰입이라 하자. 그러면 $\NL_{Z/X}$는
$D(\mathcal{O}_Z)$에서 $\mathcal{C}_{Z/X}[1]$과 동형이다. 여기서
$\mathcal{C}_{Z/X}$는 $X$ 안의 $Z$의 conormal 층이다.
\end{lemma}

\begin{proof}
이는 Algebra, Lemma \ref{algebra-lemma-NL-surjection}, Morphisms, Lemma
\ref{morphisms-lemma-affine-conormal}, 그리고 Lemma
\ref{more-morphisms-lemma-NL-affine}에서 따른다.
\end{proof}

\begin{lemma}
\label{more-morphisms-lemma-exact-sequence-NL}
$f : X \to Y$와 $g : Y \to Z$를 스킴의 사상이라 하자. 코호몰로지층의
표준 6항 완전열
$$
H^{-1}(f^*\NL_{Y/Z}) \to
H^{-1}(\NL_{X/Z}) \to
H^{-1}(\NL_{X/Y}) \to
f^*\Omega_{Y/Z} \to \Omega_{X/Z} \to \Omega_{X/Y} \to 0
$$
이 존재한다.
\end{lemma}

\begin{proof}
Modules, Lemma \ref{modules-lemma-exact-sequence-NL-ringed-topoi}의 특별한
경우이다.
\end{proof}

\begin{lemma}
\label{more-morphisms-lemma-get-triangle-NL}
$f : X \to Y$와 $Y \to Z$를 스킴의 사상이라 하자. $X \to Y$가
완전교차 사상이라고 가정하자. 그러면 표준 구별 삼각형
$$
f^*\NL_{Y/Z} \to \NL_{X/Z} \to \NL_{X/Y} \to f^*\NL_{Y/Z}[1]
$$
이 $D(\mathcal{O}_X)$에 존재하며, 이는 Lemma
\ref{more-morphisms-lemma-exact-sequence-NL}의 $6$항 완전열을 복원한다.
\end{lemma}

\begin{proof}
Modules, Lemma \ref{modules-lemma-exact-sequence-NL-ringed-topoi}의 표준 사상
$$
f^*\NL_{Y/Z} \to \text{Cone}(\NL_{X/Z} \to \NL_{X/Y})[-1]
$$
이 $D(\mathcal{O}_X)$에서 동형사상임을 보이면 충분하다. 이를 보이려면
$6$항 열의 왼쪽에 영이 있음을, 즉
$H^{-1}(f^*\NL_{Y/Z}) \to H^{-1}(\NL_{X/Z})$가 단사임을 보이면 충분하다.
이는 아핀 국소적으로 More on Algebra, Lemma
\ref{more-algebra-lemma-transitive-lci-at-end}의 대응하는 대수적 결과에서
따른다. 대수로 옮길 때에는 Lemma \ref{more-morphisms-lemma-NL-affine}를 사용한다.
\end{proof}

\begin{lemma}
\label{more-morphisms-lemma-smooth-etale-permanence}
$X \to Y \to Z$를 스킴의 사상이라 하자. $X \to Z$가 매끄럽고
$Y \to Z$가 에탈이라고 가정하자. 그러면 $X \to Y$는 매끄럽다.
% P05-EN-CH37-0064
\end{lemma}

\begin{proof}
Morphisms, Lemma \ref{morphisms-lemma-finite-presentation-permanence}에 의해
사상 $X \to Y$는 국소 유한표현이다. Lemma \ref{more-morphisms-lemma-NL-smooth}에 의해
$H^{-1}(\NL_{X/Z}) = 0$이고 가군 $\Omega_{X/Z}$는 유한 국소 자유이다.
Lemma \ref{more-morphisms-lemma-NL-etale}에 의해
$H^{-1}(\NL_{Y/Z}) = H^0(\NL_{Y/Z}) = 0$이다. Lemma
\ref{more-morphisms-lemma-exact-sequence-NL}에 의해 $H^{-1}(\NL_{X/Y}) = 0$이고
$\Omega_{X/Y} \cong \Omega_{X/Z}$는 유한 국소 자유이다. Lemma
\ref{more-morphisms-lemma-NL-smooth}에 의해 사상 $X \to Y$는 매끄럽다.
\end{proof}

\begin{lemma}
\label{more-morphisms-lemma-get-NL}
$f : X \to Y$를 스킴의 사상이라 하자. $i$는 몰입이고
$g : P \to Y$는 형식적으로 매끄러우며(예를 들어 매끄러우며),
$f = g \circ i$로 분해된다고 하자. 그러면 표준 동형사상
$$
\NL_{X/Y} \cong \left(\mathcal{C}_{X/P} \to i^*\Omega_{P/Y}\right)
$$
이 $D(\mathcal{O}_X)$에 존재한다. 여기서 conormal 층
$\mathcal{C}_{X/P}$는 차수 $-1$에 놓인다.
\end{lemma}

\begin{proof}
(괄호 안의 명제는 Lemma \ref{more-morphisms-lemma-smooth-formally-smooth}를 보라.) Lemmas
\ref{more-morphisms-lemma-NL-immersion}와 \ref{more-morphisms-lemma-NL-formally-smooth}에 의해
$\NL_{X/P} = \mathcal{C}_{X/P}[1]$이고 $\NL_{P/Y} = \Omega_{P/Y}$이며,
$\Omega_{P/Y}$는 국소 사영이다. 따라서
$i^*\NL_{P/Y} \to i^*\Omega_{P/Y}$도 준동형이다. 작은 세부사항은
생략한다. 그 이유는 $i^*\NL_{P/Y}$가
$\tau_{\geq -1}Li^*\NL_{P/Y}$와 같기 때문이다(More on Algebra, Lemma
\ref{more-algebra-lemma-tensor-NL} 참조). 따라서 Modules, Lemma
\ref{modules-lemma-exact-sequence-NL-ringed-topoi}의 표준 사상
$$
i^*\NL_{P/Y} \to \text{Cone}(\NL_{X/Y} \to \NL_{X/P})[-1]
$$
은 $D(\mathcal{O}_X)$에서 동형사상이다. 실제로 위에서 말했듯이 코호몰로지군
$H^{-1}(i^*\NL_{P/Y})$은 영이다. 달리 말해 구별 삼각형
$$
i^*\NL_{P/Y} \to \NL_{X/Y} \to \NL_{X/P} \to i^*\NL_{P/Y}[1]
$$
을 얻는다. 분명히 이는 $\NL_{X/Y}$가 사상
$\NL_{X/P}[-1] \to i^*\NL_{P/Y}$의 원뿔이라는 뜻이다. 위에서 유도범주에
있는 이 대상들의 코호몰로지층을 계산했으므로 이는 보조정리의 명제와
동치이다.
\end{proof}

\begin{lemma}
\label{more-morphisms-lemma-base-change-NL}
스킴의 데카르트 도식
$$
\xymatrix{
X' \ar[r]_{g'} \ar[d] & X \ar[d] \\
Y' \ar[r] & Y
}
$$
을 생각하자. 표준 사상 $(g')^*\NL_{X/Y} \to \NL_{X'/Y'}$은 $H^0$에서
동형사상을, $H^{-1}$에서 전사를 유도한다.
\end{lemma}

\begin{proof}
이를 대수로 옮기면 More on Algebra, Lemma
\ref{more-algebra-lemma-base-change-NL}이다. 옮길 때에는 Lemma
\ref{more-morphisms-lemma-NL-affine}를 사용한다.
\end{proof}

\begin{lemma}
\label{more-morphisms-lemma-flat-base-change-NL}
스킴의 데카르트 도식
$$
\xymatrix{
X' \ar[d] \ar[r]_{g'} & X \ar[d] \\
Y' \ar[r] & Y
}
$$
을 생각하자. $Y' \to Y$가 평탄하면 표준 사상
$(g')^*\NL_{X/Y} \to \NL_{X'/Y'}$은 준동형이다.
\end{lemma}

\begin{proof}
이는 Lemma \ref{more-morphisms-lemma-NL-affine}와 Algebra, Lemma
\ref{algebra-lemma-change-base-NL}에서 따른다.
\end{proof}

\begin{lemma}
\label{more-morphisms-lemma-base-change-NL-flat}
스킴의 데카르트 도식
$$
\xymatrix{
X' \ar[r]_{g'} \ar[d] & X \ar[d] \\
Y' \ar[r] & Y
}
$$
을 생각하자. $X \to Y$가 평탄하면 표준 사상
$(g')^*\NL_{X/Y} \to \NL_{X'/Y'}$은 준동형이다. 또한 $\NL_{X/Y}$가
$[-1, 0]$에서 Tor 진폭을 가지면
$L(g')^*\NL_{X/Y} \to \NL_{X'/Y'}$도 준동형이다.
\end{lemma}

\begin{proof}
이를 대수로 옮기면 More on Algebra, Lemma
\ref{more-algebra-lemma-base-change-NL-flat}이다. 옮길 때에는 Lemma
\ref{more-morphisms-lemma-NL-affine}와 Derived Categories of Schemes, Lemmas
\ref{perfect-lemma-affine-compare-bounded} 및
\ref{perfect-lemma-tor-dimension-affine}를 사용한다.
\end{proof}











\section{스킴의 범주에서의 밀어내기, I}
\label{more-morphisms-section-pushouts}

\noindent
이 절에서는 $X \to Y$가 아핀이고 $X \to X'$이 비후화인 경우에
$Y \leftarrow X \rightarrow X'$의 밀어내기를 구성한다. 이는 실제로 우리에게
중요한 경우이므로 자세히 논의할 가치가 있다. Section
\ref{more-morphisms-section-pushouts-II}에서는 더 흥미롭고 더 어려운 경우를 논의한다.
임의의 범주에서 밀어내기에 관한 일반적인 논의는 Categories, Section
\ref{categories-section-pushouts}를 보라.

\begin{lemma}
\label{more-morphisms-lemma-basic-example-pushout}
$A' \to A$를 환의 전사라 하고 $B \to A$를 환 사상이라 하자.
$B' = B \times_A A'$을 환의 섬유곱이라 하자.
$S = \Spec(A)$, $S' = \Spec(A')$, $T = \Spec(B)$,
$T' = \Spec(B')$라 두자. 그러면
$$
\vcenter{
\xymatrix{
S \ar[r]_i \ar[d]_f & S' \ar[d]^{f'} \\
T \ar[r]^{i'} & T'
}
}
\quad\text{다음에 대응하는}\quad
\vcenter{
\xymatrix{
A & A' \ar[l] \\
B \ar[u] & B' \ar[l] \ar[u]
}
}
$$
은 스킴의 밀어내기이다.
\end{lemma}

\begin{proof}
More on Algebra, Lemma \ref{more-algebra-lemma-points-of-fibre-product}에 의해
위상공간으로서 $T' = T \amalg_S S'$이다. 즉 이 도식은 위상공간의 범주에서
밀어내기이다. 다음 사상을 생각하자.
$$
((i')^\sharp, (f')^\sharp) :
\mathcal{O}_{T'}
\longrightarrow
i'_*\mathcal{O}_T \times_{g_*\mathcal{O}_S} f'_*\mathcal{O}_{S'}
$$
여기서 $g = i' \circ f = f' \circ i$이다. 이 사상이 환의 층의 동형사상이라고
주장한다. 실제로 양변을 준연접 $\mathcal{O}_{T'}$-가군으로 볼 수 있고
(우변에는 Schemes,
% P05-EN-CH37-0065
Lemma \ref{schemes-lemma-push-forward-quasi-coherent}를 사용한다), 이 사상은
$\mathcal{O}_{T'}$-선형이다. 따라서 전역절단에서 이 사상이 동형사상임을
보이면 충분하다(Schemes, Lemma
\ref{schemes-lemma-equivalence-quasi-coherent}). 전역절단에서는 보조정리의
명제에 있는 동일시 $B' \to B \times_A A'$를 되찾는다(이것이 우리가
$B'$을 택한 방식이다).

\medskip\noindent
$X$를 스킴이라 하자. $m' \circ i = n \circ f$를 만족하는 스킴의 사상
$m' : S' \to X$와 $n : T \to X$가 주어졌다고 하자(이 공통 사상을
$m$이라 하자). 위에서 보았듯이 $T' = T \amalg_S S'$이므로 $m'$ 및 $n$과
양립하는 유일한 위상공간의 사상 $n' : T' \to X$를 얻는다. 앞 문단의
$\mathcal{O}_{T'}$에 대한 기술에 의해 환의 층의 유일한 준동형
% P05-EN-CH37-0066
$$
(n')^\sharp :
\mathcal{O}_X
\longrightarrow
(n')_*\mathcal{O}_{T'} =
m'_*\mathcal{O}_T \times_{m_*\mathcal{O}_T} n_*\mathcal{O}_S
$$
을 얻으며, 이는 $(m')^\sharp$와 $n^\sharp$로 주어진다. 따라서
$(n', (n')^\sharp)$는 $m'$ 및 $n$과 양립하는 환 달린 공간의 유일한 사상
$T' \to X$이다. 증명을 끝내려면 $n'$가 스킴의 사상, 즉 국소환 달린
공간의 사상임을 보이면 충분하다.

\medskip\noindent
$t' \in T'$의 상을 $x \in X$라 하자.
$\mathcal{O}_{X, x} \to \mathcal{O}_{T', t'}$가 국소임을 보여야 한다.
$t' \not \in T$이면 $t'$은 유일한 점 $s' \in S'$의 상이고
$\mathcal{O}_{T', t'} = \mathcal{O}_{S', s'}$이다. 실제로 $B' \to A'$은
동형사상 $\Ker(B' \to B) = \Ker(A' \to A)$를 유도하므로
$S' \setminus S \to T' \setminus T$는 스킴의 동형사상이다. $t'$이
$t \in T$의 상이면 합성
$\mathcal{O}_{X, x} \to \mathcal{O}_{T', t'} \to \mathcal{O}_{T, t}$가
국소임을 알고 있으므로 이 경우에도 국소임을 얻는다.
% P05-EN-CH37-0067
\end{proof}

\begin{lemma}
\label{more-morphisms-lemma-pushout-fpqc-local}
$\mathcal{I} \to (\Sch/S)_{fppf}$, $i \mapsto X_i$를 스킴의 도식이라 하자.
$(W, X_i \to W)$를 스킴의 범주에서 이 도식의 쌍대원뿔이라 하자
(Categories, Remark \ref{categories-remark-cones-and-cocones}). 다음을
만족하는 스킴의 fpqc 덮개 $\{W_a \to W\}_{a \in A}$가 존재한다고 하자.
\begin{enumerate}
\item 모든 $a \in A$에 대해 스킴의 범주에서
$W_a = \colim X_i \times_W W_a$
이다.
\item 모든 $a, b \in A$에 대해 스킴의 범주에서
$W_a \times_W W_b = \colim X_i \times_W W_a \times_W W_b$
이다.
\end{enumerate}
그러면 스킴의 범주에서 $W = \colim X_i$이다.
\end{lemma}

\begin{proof}
실제로 스킴 $T$에 대해 사상 $W \to T$는 겹침
$W_a \times_W W_b$에서 일치하는 사상들의 모음
$W_a \to T$, $a \in A$와 같은 것이다. Descent, Lemma
\ref{descent-lemma-fpqc-universal-effective-epimorphisms}를 보라.
% P05-EN-CH37-0068
\end{proof}

\begin{lemma}
\label{more-morphisms-lemma-pushout-along-thickening}
$X \to X'$을 스킴의 비후화라 하고 $X \to Y$를 스킴의 아핀 사상이라 하자.
그러면 스킴의 범주에서 밀어내기
$$
\xymatrix{
X \ar[r] \ar[d]_f
&
X' \ar[d]^{f'}
\\
Y \ar[r]
&
Y'
}
$$
가 존재한다. 더욱이 $Y \subset Y'$은 비후화이고
$X = Y \times_{Y'} X'$이며, $|Y| = |Y'|$ 위의 층으로서
$$
\mathcal{O}_{Y'} = \mathcal{O}_Y \times_{f_*\mathcal{O}_X} f'_*\mathcal{O}_{X'}
$$
이다.
\end{lemma}

\begin{proof}
먼저 $Y'$을 환 달린 공간으로 구성한다. 위상공간으로는 $Y' = Y$라 둔다.
% P05-EN-CH37-0069
$f$와 같은 위상공간의 사상을 $f' : X' \to Y'$라 하자. 구조층
$\mathcal{O}_{Y'}$로는 보조정리의 등식 우변을 취한다. $Y'$이 스킴임을
보이려면 모든 점이 아핀 근방을 가짐을 보여야 한다. 층의 섬유곱을 형성하는
것은 열린집합으로 제한하는 것과 가환하므로 $Y$가 아핀이라고 가정해도 된다.
그러면 $f$가 아핀이므로 $X$는 아핀이고, $X'$도 아핀이다(Lemma
\ref{more-morphisms-lemma-thickening-affine-scheme} 참조). $Y \leftarrow X \rightarrow X'$이
$B \rightarrow A \leftarrow A'$에 대응한다고 하자.
$B' = B \times_A A'$라 두자. 이는 $\mathcal{O}_{Y'}$의 전역절단이다.
$A' \to A$가 국소 멱영인 핵을 갖는 전사이므로 $B' \to B$도 국소 멱영인
핵을 갖는 전사이다. 따라서 위상공간으로서
$\Spec(B') = \Spec(B)$이다. $Y' = \Spec(B')$라고 주장한다. 이를 보이기
위해 상이 $g \in B$인 $g' \in B'$에 대해
$\mathcal{O}_{Y'}(D(g)) = B'_{g'}$임을 보이겠다. More on Algebra, Lemma
\ref{more-algebra-lemma-diagram-localize}에 의해
$$
(B')_{g'} = B_g \times_{A_h} A'_{h'}
$$
이다. 여기서 $h \in A$, $h' \in A'$은 $g'$의 상이다. $B_g$, 각각
$A_h$, 각각 $A'_{h'}$는 $\mathcal{O}_Y(D(g))$, 각각
$f_*\mathcal{O}_X(D(g))$, 각각 $f'_*\mathcal{O}_{X'}(D(g))$와 같으므로
주장이 따른다.

\medskip\noindent
이제 $Y'$이 밀어내기임을 보이면 된다. 위의 논의에 의해 스킴 $Y'$에는
아핀 열린 덮개 $Y' = \bigcup W'_i$가 있어서 대응하는 열린집합
$U'_i \subset X'$, $W_i \subset Y$, 그리고
$U_i \subset X$은 아핀이다. 더욱이 $A'_i$, $B_i$, $A_i$가 각각 $U'_i$, $W_i$, $U_i$에
대응하는 환이면 $W'_i$는 $B_i \times_{A_i} A'_i$에 대응한다. 따라서
Lemmas \ref{more-morphisms-lemma-basic-example-pushout}와
\ref{more-morphisms-lemma-pushout-fpqc-local}를 적용하면 우리의 구성이 스킴의 범주에서
밀어내기라는 결론을 얻는다.
\end{proof}

\noindent
다음 보조정리에서는 Categories, Example
\ref{categories-example-2-fibre-product-categories}에서 정의한
범주들의 올곱을 사용한다.

\begin{lemma}
\label{more-morphisms-lemma-equivalence-categories-schemes-over-pushout}
$X \to X'$를 스킴의 비후화라 하고 $X \to Y$를 스킴의 아핀 사상이라 하자.
$Y' = Y \amalg_X X'$가 밀어내기라 하자(Lemma
\ref{more-morphisms-lemma-pushout-along-thickening} 참조). 밑변환은 함자
$$
F : (\Sch/Y') \longrightarrow (\Sch/Y) \times_{(\Sch/X)} (\Sch/X')
$$
를 주는데, 이는 $V' \longmapsto (V' \times_{Y'} Y, V' \times_{Y'} X', 1)$로
주어지고 왼쪽 수반
$$
G : (\Sch/Y) \times_{(\Sch/X)} (\Sch/X') \longrightarrow (\Sch/Y')
$$
을 가진다. 이 수반은 삼중항 $(V, U', \varphi)$를 밀어내기
$V \amalg_{(V \times_Y X)} U'$로 보낸다. 끝으로 $F \circ G$는 항등함자와
동형이다.
\end{lemma}

\begin{proof}
$(V, U', \varphi)$를 올곱 범주의 대상으로 두자.
$U = U' \times_{X'} X$라 놓는다. $U \to U'$가 비후화임에 유의하자.
$\varphi : V \times_Y X \to U' \times_{X'} X = U$가 동형이므로,
$U$를 올곱 $X \times_Y V$와 동일시하는 $X \to Y$ 위의 사상
$U \to V$가 있다. 특히 $U \to V$는 아핀이다. Morphisms, Lemma
\ref{morphisms-lemma-base-change-affine}를 보라.
따라서 Lemma \ref{more-morphisms-lemma-pushout-along-thickening}을 적용하여 밀어내기
$V' = V \amalg_U U'$를 얻을 수 있다.
% P05-EN-CH37-0070
$V'$가 밀어내기라는 사실과, $Y'$로 가는 사상으로 보았을 때 $U$ 위에서
일치하는 사상 $V \to Y$ 및 $U' \to X'$가 주어졌다는 사실로부터 얻는
사상을 $V' \to Y'$로 나타내자. $G(V, U', \varphi) = V'$로 놓으면
함자 $G$를 얻는다.

\medskip\noindent
$G$가 $F$의 왼쪽 수반임을 보이자. $Z$를 $Y'$ 위의 스킴이라 하자.
다음을 보이면 된다.
$$
\Mor(V', Z) = \Mor((V, U', \varphi), F(Z))
$$
여기서 사상들의 집합은 각각 해당 범주에서 취한다.
% P05-EN-CH37-0071
$g' : V' \to Z$를 사상이라 하자. $g'$와 사상 $V \to V'$, 각각
$U' \to V'$의 합성을 $\tilde g$, 각각 $\tilde f'$로 나타내자.
% P05-EN-CH37-0072
$\tilde g$, 각각 $\tilde f'$를 $Y \to Y'$, 각각 $X' \to Y'$로 밑변환하면
사상 $g : V \to Z \times_{Y'} Y$, 각각
$f' : U' \to Z \times_{Y'} X'$를 얻는다. 그러면 $(g, f')$는 위 등식의
우변의 원소이다(세부사항은 생략한다).
반대로 $(g, f') : (V, U', \varphi) \to F(Z)$가 우변의 원소라고 하자.
% P05-EN-CH37-0073
$g$, 각각 $f$를 $Z \times_{Y'} X' \to Z$, 각각
$Z \times_{Y'} Y \to Z$와 합성하여 $\tilde g : V \to Z$, 각각
$\tilde f' : U' \to Z$를 생각할 수 있다.
% P05-EN-CH37-0074
그러면 $\tilde g$와 $\tilde f'$는 $U$에서 $Z$로 가는 사상으로서 일치한다.
밀어내기의 보편 성질에 의해 사상 $g' : V' \to Z$, 즉 좌변의 원소를
얻는다. 이 구성들이 서로 역이라는 확인은 생략한다.

\medskip\noindent
$F \circ G$가 항등함자와 동형임을 보이려면 수반 사상
$(V, U', \varphi) \to F(G(V, U', \varphi))$이 동형임을 보이면 된다.
이를 위해 아핀 국소적으로 논할 수 있다. $X = \Spec(A)$,
$X' = \Spec(A')$, 그리고 $Y = \Spec(B)$라 하자. 그러면 $A' \to A$와
$B \to A$는 More on Algebra, Lemma
\ref{more-algebra-lemma-module-over-fibre-product}에서와 같은 환 사상이고,
$B' = B \times_A A'$인 $Y' = \Spec(B')$를 얻는다. 다음으로
$V = \Spec(D)$, $U' = \Spec(C')$이고 $\varphi$가 $A$-대수 동형
$D \otimes_B A \to C' \otimes_{A'} A = C'/IC'$로 주어진다고 하자.
$D' = D \times_{C'/IC'} C'$라 놓는다. 이 경우 보여야 할 주장은
$D' \otimes_{B'} B \cong D$ 및 $D' \otimes_{B'} A' \cong C'$이다.
이는 More on Algebra, Lemma
\ref{more-algebra-lemma-module-over-fibre-product}의 특수한 경우이다.
\end{proof}

\begin{lemma}
\label{more-morphisms-lemma-scheme-over-pushout-flat-modules}
$X \to X'$를 스킴의 비후화라 하고 $X \to Y$를 스킴의 아핀 사상이라 하자.
$Y' = Y \amalg_X X'$가 밀어내기라 하자(Lemma
\ref{more-morphisms-lemma-pushout-along-thickening} 참조). $V' \to Y'$를 스킴의 사상이라
하자. 다음과 같이 놓는다.
$V = Y \times_{Y'} V'$, $U' = X' \times_{Y'} V'$, 그리고 $U = X \times_{Y'} V'$.
다음 두 범주는 서로 동치이다.
\begin{enumerate}
\item $Y'$ 위에서 평탄한 준연접 $\mathcal{O}_{V'}$-가군들의 범주,
\item 다음 조건들을 만족하는 삼중항 $(\mathcal{G}, \mathcal{F}', \varphi)$의 범주:
\begin{enumerate}
\item $\mathcal{G}$는 $Y$ 위에서 평탄한 준연접 $\mathcal{O}_V$-가군이다.
\item $\mathcal{F}'$는 $X'$ 위에서 평탄한 준연접 $\mathcal{O}_{U'}$-가군이다.
\item $\varphi : (U \to V)^*\mathcal{G} \to (U \to U')^*\mathcal{F}'$는
$\mathcal{O}_U$-가군의 동형이다.
\end{enumerate}
\end{enumerate}
이 동치는 $\mathcal{G}'$를
$((V \to V')^*\mathcal{G}', (U' \to V')^*\mathcal{G}', can)$으로 보낸다.
$\mathcal{G}'$가 삼중항 $(\mathcal{G}, \mathcal{F}', \varphi)$에 대응한다고
하자. 그러면 다음이 성립한다.
% P05-EN-CH37-0075
\begin{enumerate}
\item[(a)] $\mathcal{G}'$가 유한형 $\mathcal{O}_{V'}$-가군일 필요충분조건은
$\mathcal{G}$와 $\mathcal{F}'$가 각각 유한형 $\mathcal{O}_Y$-가군과
$\mathcal{O}_{U'}$-가군인 것이다.
\item[(b)] $V' \to Y'$가 국소 유한 표시이면, $\mathcal{G}'$가 유한 표시
$\mathcal{O}_{V'}$-가군일 필요충분조건은 $\mathcal{G}$와 $\mathcal{F}'$가
각각 유한 표시 $\mathcal{O}_Y$-가군과 $\mathcal{O}_{U'}$-가군인 것이다.
\end{enumerate}
\end{lemma}

\begin{proof}
준역함자는 삼중항 $(\mathcal{G}, \mathcal{F}', \varphi)$에 올곱
$$
(V \to V')_*\mathcal{G}
\times_{(U \to V')_*\mathcal{F}}
(U' \to V')_*\mathcal{F}'
$$
을 대응시킨다. 여기서 $\mathcal{F} = (U \to U')^*\mathcal{F}'$이다.
실제로 아핀 위에서는 More on Algebra, Lemma
\ref{more-algebra-lemma-relative-flat-module-over-fibre-product}의 동치를
되찾기 때문에 이 구성은 잘 작동한다. 일부 세부사항은 생략한다.

\medskip\noindent
(a)와 (b)는 More on Algebra, Lemmas
\ref{more-algebra-lemma-relative-finite-module-over-fibre-product} 및
\ref{more-algebra-lemma-relative-finitely-presented-module-over-fibre-product}에서
따른다.
\end{proof}

\begin{lemma}
\label{more-morphisms-lemma-equivalence-categories-schemes-over-pushout-flat}
Lemma \ref{more-morphisms-lemma-equivalence-categories-schemes-over-pushout}의 상황을
생각하자. 어떤 삼중항 $(V, U', \varphi)$에 대해
$V' = G(V, U', \varphi)$이면 다음이 성립한다.
\begin{enumerate}
\item $V' \to Y'$가 국소 유한형일 필요충분조건은 $V \to Y$와
$U' \to X'$가 국소 유한형인 것이다.
\item $V' \to Y'$가 평탄할 필요충분조건은 $V \to Y$와 $U' \to X'$가
평탄한 것이다.
\item $V' \to Y'$가 평탄하고 국소 유한 표시일 필요충분조건은
$V \to Y$와 $U' \to X'$가 평탄하고 국소 유한 표시인 것이다.
\item $V' \to Y'$가 매끄러울 필요충분조건은 $V \to Y$와
$U' \to X'$가 매끄러운 것이다.
\item $V' \to Y'$가 \'etale일 필요충분조건은 $V \to Y$와 $U' \to X'$가
\'etale인 것이다.
% P05-EN-CH37-0076
\item 필요에 따라 여기에 더 추가한다.
\end{enumerate}
$W'$가 $Y'$ 위에서 평탄하면 수반 사상 $G(F(W')) \to W'$는 동형이다.
따라서 $F$와 $G$는 $Y'$ 위에서 평탄한 스킴들의 범주와 $V \to Y$ 및
$U' \to X'$가 평탄한 삼중항 $(V, U', \varphi)$의 범주 사이의 서로
준역인 함자를 정의한다.
\end{lemma}

\begin{proof}
아핀 조각들 위에서 보면 이 보조정리의 주장들은 More on Algebra, Lemma
\ref{more-algebra-lemma-properties-algebras-over-fibre-product}의 대응하는
주장들과 동치이다.
\end{proof}







\section{평탄 자취의 열림성}
\label{more-morphisms-section-open-flat}

\noindent
이 결과를 증명하려면 다소 수고가 필요하고, (어쩌면) 독립된 절을
할애할 가치가 있다. 그 결과는 다음과 같다.

\begin{theorem}
\label{more-morphisms-theorem-openness-flatness}
\begin{reference}
\cite[IV Theorem 11.3.1]{EGA}
\end{reference}
$S$를 스킴이라 하자. $f : X \to S$를 국소 유한 표시인 사상이라 하자.
$\mathcal{F}$를 국소 유한 표시인 준연접 $\mathcal{O}_X$-가군이라 하자.
그러면
$$
U = \{x \in X \mid \mathcal{F}\text{가 }S\text{ 위에서 }x\text{에서 평탄하다}\}
$$
는 $X$에서 열린집합이다.
\end{theorem}

\begin{proof}
열림성은 $X$ 위에서 국소적으로 검사할 수 있으므로 $f$가 아핀 스킴들의
사상이라고 가정할 수 있다. 이 경우 정리는 정확히 Algebra, Theorem
\ref{algebra-theorem-openness-flatness}이다.
\end{proof}

\begin{lemma}
\label{more-morphisms-lemma-flat-locus-base-change}
$S$를 스킴이라 하고, 다음을 스킴들의 데카르트 그림이라 하자.
$$
\xymatrix{
X' \ar[r]_{g'} \ar[d]_{f'} & X \ar[d]^f \\
S' \ar[r]^g & S
}
$$
$\mathcal{F}$를 준연접 $\mathcal{O}_X$-가군이라 하자. $x' \in X'$의
상들이 $x = g'(x')$ 및 $s' = f'(x')$라고 하자.
\begin{enumerate}
\item $\mathcal{F}$가 $x$에서 $S$ 위에서 평탄하면,
$(g')^*\mathcal{F}$는 $x'$에서 $S'$ 위에서 평탄하다.
\item $g$가 $s'$에서 평탄하고 $(g')^*\mathcal{F}$가 $x'$에서 $S'$ 위에서
평탄하면, $\mathcal{F}$는 $x$에서 $S$ 위에서 평탄하다.
\end{enumerate}
특히 $g$가 평탄하고 $f$가 국소 유한 표시이며 $\mathcal{F}$가 국소 유한
표시이면, Theorem \ref{more-morphisms-theorem-openness-flatness}의 열린부분집합의 구성은
밑변환과 가환한다.
\end{lemma}

\begin{proof}
국소환들의 가환 그림
$$
\xymatrix{
\mathcal{O}_{X', x'} & \mathcal{O}_{X, x} \ar[l] \\
\mathcal{O}_{S', s'} \ar[u] & \mathcal{O}_{S, s} \ar[l] \ar[u]
}
$$
을 생각하자. $\mathcal{O}_{X', x'}$는
$\mathcal{O}_{X, x} \otimes_{\mathcal{O}_{S, s}} \mathcal{O}_{S', s'}$의
국소화이고, $((g')^*\mathcal{F})_{x'}$는
$\mathcal{F}_x \otimes_{\mathcal{O}_{X, x}} \mathcal{O}_{X', x'}$와 같음에
유의하자. 따라서 보조정리는 Algebra, Lemma
\ref{algebra-lemma-base-change-flat-up-down}에서 따른다.
\end{proof}






\section{올에 의한 평탄성 판정법}
\label{more-morphisms-section-criterion-flat-fibres}

\noindent
스킴들의 가환 그림(왼쪽 그림)
$$
\xymatrix{
X \ar[rr]_f \ar[dr] & & Y \ar[dl] \\
& S
}
\quad
\xymatrix{
X_s \ar[rr]_{f_s} \ar[rd] & & Y_s \ar[dl] \\
& \Spec(\kappa(s))
}
$$
과 준연접 $\mathcal{O}_X$-가군 $\mathcal{F}$를 생각하자.
점 $x \in X$가 $s \in S$ 위에 놓이고 상이 $y = f(x)$라고 하자.
$\mathcal{F}$는 $Y$ 위에서 $x$에서 평탄한가라는 질문을 생각한다.
$\mathcal{F}$가 $S$ 위에서 $x$에서 평탄하면, 아래 정리는 이 질문이
$\mathcal{F}$의 올에 대한 제한
$$
\mathcal{F}_s = (X_s \to X)^*\mathcal{F}
$$
이 $Y_s$ 위에서 $x$에서 평탄한가라는 질문과 밀접하게 관련됨을 말해 준다.
아래에는 ``뇌터'' 버전과 ``유한 표시'' 버전이 있으며, 앞에서는
``멱영'' 버전을 다루었다. Lemma
\ref{more-morphisms-lemma-flatness-morphism-thickenings}를 보라.

\begin{theorem}
\label{more-morphisms-theorem-criterion-flatness-fibre-Noetherian}
$S$를 스킴이라 하자. $f : X \to Y$를 $S$ 위 스킴들의 사상이라 하고,
$\mathcal{F}$를 준연접 $\mathcal{O}_X$-가군이라 하자. $x \in X$라 하자.
$y = f(x)$라 놓고, $s \in S$를 $x$의 $S$에서의 상이라 하자.
% P05-EN-CH37-0077
$S$, $X$, $Y$가 국소 뇌터이고 $\mathcal{F}$가 연접이며
$\mathcal{F}_x \not = 0$이라고 가정하자. 그러면 다음은 서로 동치이다.
\begin{enumerate}
\item $\mathcal{F}$는 $S$ 위에서 $x$에서 평탄하고,
$\mathcal{F}_s$는 $Y_s$ 위에서 $x$에서 평탄하다.
\item $Y$는 $S$ 위에서 $y$에서 평탄하고, $\mathcal{F}$는
$Y$ 위에서 $x$에서 평탄하다.
\end{enumerate}
\end{theorem}

\begin{proof}
환 사상들
$$
\mathcal{O}_{S, s} \longrightarrow
\mathcal{O}_{Y, y} \longrightarrow \mathcal{O}_{X, x}
$$
과 가군 $\mathcal{F}_x$를 생각하자. $\mathcal{F}_s$의 $x$에서의 줄기는
가군 $\mathcal{F}_x/\mathfrak m_s \mathcal{F}_x$이고, $Y_s$의 $y$에서의
국소환은 $\mathcal{O}_{Y, y}/\mathfrak m_s \mathcal{O}_{Y, y}$이다.
따라서 함의 (1) $\Rightarrow$ (2)는 Algebra, Lemma
\ref{algebra-lemma-criterion-flatness-fibre-Noetherian}이다.
(2)가 성립하면 첫 번째 환 사상은 충실 평탄이고 $\mathcal{F}_x$는
$\mathcal{O}_{Y, y}$ 위에서 평탄하므로, Algebra, Lemma
\ref{algebra-lemma-composition-flat}에 의해 $\mathcal{F}_x$가
$\mathcal{O}_{S, s}$ 위에서 평탄함을 알 수 있다. 더욱이
$\mathcal{F}_x/\mathfrak m_s \mathcal{F}_x$는 평탄 가군 $\mathcal{F}_x$를
$\mathcal{O}_{Y, y} \to \mathcal{O}_{Y, y}/\mathfrak m_s \mathcal{O}_{Y, y}$로
밑변환한 것이므로, Algebra, Lemma
\ref{algebra-lemma-flat-base-change}에 의해 평탄하다.
\end{proof}

\noindent
다음은 비뇌터 버전이다.

\begin{theorem}
\label{more-morphisms-theorem-criterion-flatness-fibre}
$S$를 스킴이라 하자. $f : X \to Y$를 $S$ 위 스킴들의 사상이라 하고,
$\mathcal{F}$를 준연접 $\mathcal{O}_X$-가군이라 하자. 다음을 가정하자.
\begin{enumerate}
\item $X$는 $S$ 위에서 국소 유한 표시이다.
% P05-EN-CH37-0078
\item $\mathcal{F}$는 유한 표시 $\mathcal{O}_X$-가군이다.
\item $Y$는 $S$ 위에서 국소 유한형이다.
\end{enumerate}
$x \in X$라 하자. $y = f(x)$라 놓고, $s \in S$를 $x$의 $S$에서의 상이라
하자. $\mathcal{F}_x \not = 0$이면 다음은 서로 동치이다.
\begin{enumerate}
\item $\mathcal{F}$는 $S$ 위에서 $x$에서 평탄하고,
$\mathcal{F}_s$는 $Y_s$ 위에서 $x$에서 평탄하다.
\item $Y$는 $S$ 위에서 $y$에서 평탄하고, $\mathcal{F}$는
$Y$ 위에서 $x$에서 평탄하다.
\end{enumerate}
더욱이 (1)과 (2)가 성립하는 점 $x$들의 집합은
$\text{Supp}(\mathcal{F})$에서 열린집합이다.
\end{theorem}

\begin{proof}
환 사상들
$$
\mathcal{O}_{S, s} \longrightarrow
\mathcal{O}_{Y, y} \longrightarrow \mathcal{O}_{X, x}
$$
과 가군 $\mathcal{F}_x$를 생각하자. $\mathcal{F}_s$의 $x$에서의 줄기는
가군 $\mathcal{F}_x/\mathfrak m_s \mathcal{F}_x$이고, $Y_s$의 $y$에서의
국소환은 $\mathcal{O}_{Y, y}/\mathfrak m_s \mathcal{O}_{Y, y}$이다.
따라서 함의 (1) $\Rightarrow$ (2)는 Algebra, Lemma
\ref{algebra-lemma-criterion-flatness-fibre-fp-over-ft}이다.
(2)가 성립하면 첫 번째 환 사상은 충실 평탄이고 $\mathcal{F}_x$는
$\mathcal{O}_{Y, y}$ 위에서 평탄하므로, Algebra, Lemma
\ref{algebra-lemma-composition-flat}에 의해 $\mathcal{F}_x$가
$\mathcal{O}_{S, s}$ 위에서 평탄함을 알 수 있다. 더욱이
$\mathcal{F}_x/\mathfrak m_s \mathcal{F}_x$는 평탄 가군 $\mathcal{F}_x$를
$\mathcal{O}_{Y, y} \to \mathcal{O}_{Y, y}/\mathfrak m_s \mathcal{O}_{Y, y}$로
밑변환한 것이므로, Algebra, Lemma
\ref{algebra-lemma-flat-base-change}에 의해 평탄하다.

\medskip\noindent
Morphisms, Lemma \ref{morphisms-lemma-finite-presentation-permanence}에
의해 사상 $f$는 국소 유한 표시이다. 집합
\begin{equation}
\label{more-morphisms-equation-open}
U = \{x \in X \mid \mathcal{F} \text{는 }x\text{에서 }
Y\text{와 }S\text{ 모두의 위에서 평탄하다}\}.
\end{equation}
이 집합은 Theorem \ref{more-morphisms-theorem-openness-flatness}에 의해 $X$에서
열린집합이다. $x \in U$이면 $\mathcal{F}_s$는 $x$에서 $Y_s$ 위에서
평탄한데, 이는 평탄 가군을 사상 $Y_s \to Y$ 아래에서 밑변환한 것이다.
Morphisms, Lemma \ref{morphisms-lemma-base-change-module-flat}를 보라.
따라서 $U \cap \text{Supp}(\mathcal{F})$의 모든 점에서 조건 (1)이
성립한다. 한편 $x \in \text{Supp}(\mathcal{F})$가 (1)과 (2)를 만족하면
$x \in U$임은 명백하다. 따라서 우리가 찾는 열린집합은
$U \cap \text{Supp}(\mathcal{F})$이다.
\end{proof}

\noindent
이 정리들은 흔히 다음과 같은 단순화된 형태로 사용된다. 여기서는 대역적
진술만 제시하지만, 물론 점별 버전들도 있다.

\begin{lemma}
\label{more-morphisms-lemma-morphism-between-flat-Noetherian}
$S$를 스킴이라 하고 $f : X \to Y$를 $S$ 위 스킴들의 사상이라 하자.
다음을 가정하자.
\begin{enumerate}
\item $S$, $X$, $Y$는 국소 뇌터이다.
\item $X$는 $S$ 위에서 평탄하다.
\item 모든 $s \in S$에 대해 사상 $f_s : X_s \to Y_s$는 평탄하다.
\end{enumerate}
그러면 $f$는 평탄하다. $f$가 전사이기도 하면 $Y$는 $S$ 위에서 평탄하다.
\end{lemma}

\begin{proof}
이는 Theorem \ref{more-morphisms-theorem-criterion-flatness-fibre-Noetherian}의 특수한
경우이다.
\end{proof}

\begin{lemma}
\label{more-morphisms-lemma-morphism-between-flat}
$S$를 스킴이라 하고 $f : X \to Y$를 $S$ 위 스킴들의 사상이라 하자.
다음을 가정하자.
\begin{enumerate}
\item $X$는 $S$ 위에서 국소 유한 표시이다.
\item $X$는 $S$ 위에서 평탄하다.
\item 모든 $s \in S$에 대해 사상 $f_s : X_s \to Y_s$는 평탄하다.
\item $Y$는 $S$ 위에서 국소 유한형이다.
\end{enumerate}
그러면 $f$는 평탄하다. $f$가 전사이기도 하면 $Y$는 $S$ 위에서 평탄하다.
\end{lemma}

\begin{proof}
이는 Theorem \ref{more-morphisms-theorem-criterion-flatness-fibre}의 특수한 경우이다.
\end{proof}

\begin{lemma}
\label{more-morphisms-lemma-base-change-criterion-flatness-fibre}
$S$를 스킴이라 하고 $f : X \to Y$를 $S$ 위 스킴들의 사상이라 하자.
$\mathcal{F}$를 준연접 $\mathcal{O}_X$-가군이라 하자. 다음을 가정하자.
\begin{enumerate}
\item $X$는 $S$ 위에서 국소 유한 표시이다.
% P05-EN-CH37-0079
\item $\mathcal{F}$는 유한 표시 $\mathcal{O}_X$-가군이다.
\item $\mathcal{F}$는 $S$ 위에서 평탄하다.
\item $Y$는 $S$ 위에서 국소 유한형이다.
\end{enumerate}
그러면 집합
% P05-EN-CH37-0081
$$
U = \{x \in X \mid \mathcal{F}\text{는 }x\text{에서 }Y\text{ 위에서 평탄하다}\}.
$$
은 $X$에서 열려 있고 그 구성은 임의의 밑변환과 가환한다. 즉,
$S' \to S$가 스킴의 사상이고 $U'$가 $X' = X \times_S S'$의 점들 가운데
$\mathcal{F}' = \mathcal{F} \times_S S'$가 $Y' = Y \times_S S'$ 위에서
평탄한 점들의 집합이면, $U' = U \times_S S'$이다.
% P05-EN-CH37-0080
\end{lemma}

\begin{proof}
Morphisms, Lemma \ref{morphisms-lemma-finite-presentation-permanence}에
의해 사상 $f$는 국소 유한 표시이다. 따라서 Theorem
\ref{more-morphisms-theorem-openness-flatness}에 의해 $U$는 열려 있다.
$\mathcal{F}$가 $S$ 위에서 평탄하다고 가정했으므로, Theorem
\ref{more-morphisms-theorem-criterion-flatness-fibre}에서 다음이 따른다.
% P05-EN-CH37-0082
$$
U = \{x \in X \mid \mathcal{F}_s\text{는 }x\text{에서 }Y_s\text{ 위에서 평탄하다}\}.
$$
여기서 $s$는 언제나 $x$의 $S$에서의 상을 뜻한다. (이 기술은
$\mathcal{F}_x = 0$일 때도 자명하게 성립한다.) 더욱이 보조정리의 가정들은
사상 $f' : X' \to Y'$와 층 $\mathcal{F}'$에 대해서도 계속 성립한다.
따라서 $U'$도 비슷하게 기술된다. 다시 말해 $s \in S$로 가는
$s' \in S'$가 주어졌을 때,
$$
\{x' \in X'_{s'} \mid
\mathcal{F}'_{s'}\text{는 }x'\text{에서 }Y'_{s'}\text{ 위에서 평탄하다}\}
$$
가 $X_s$에서 대응하는 자취의 역상임을 보이면 충분하다. 이는 Lemma
\ref{more-morphisms-lemma-flat-locus-base-change}에 의해 참이다. 실제로 데카르트 그림
$$
\xymatrix{
X'_{s'} \ar[d] \ar[r] & X_s \ar[d] \\
Y'_{s'} \ar[r] & Y_s
}
$$
에서 수평 사상들은 평탄 사상
$\Spec(\kappa(s')) \to \Spec(\kappa(s))$에 의한 밑변환이므로 평탄하다.
\end{proof}

\begin{lemma}
\label{more-morphisms-lemma-base-change-flatness-fibres}
$S$를 스킴이라 하고 $f : X \to Y$를 $S$ 위 스킴들의 사상이라 하자.
다음을 가정하자.
\begin{enumerate}
\item $X$는 $S$ 위에서 국소 유한 표시이다.
\item $X$는 $S$ 위에서 평탄하다.
\item $Y$는 $S$ 위에서 국소 유한형이다.
\end{enumerate}
그러면 집합
% P05-EN-CH37-0083
$$
U = \{x \in X \mid X\text{는 }x\text{에서 }Y\text{ 위에서 평탄하다}\}.
$$
은 $X$에서 열려 있고 그 구성은 임의의 밑변환과 가환한다.
\end{lemma}

\begin{proof}
이는 Lemma \ref{more-morphisms-lemma-base-change-criterion-flatness-fibre}의 특수한 경우이다.
\end{proof}

\noindent
다음 보조정리는 Algebra, Lemma
\ref{algebra-lemma-free-fibre-flat-free}의 변형이다.
$(\mathcal{F}_s)_x$가 평탄 $\mathcal{O}_{X_s, x}$-가군이라는 가정은
$(\mathcal{F}_s)_x$가 자유 $\mathcal{O}_{X_s, x}$-가군이라는 뜻임에
유의하자. $x \in X_s$가 $X_s$의 기약 성분의 일반점이고 $X_s$가
축약이면 이는 언제나 성립한다(실제로 이 경우 $\mathcal{O}_{X_s, x}$는
체이다. Algebra, Lemma \ref{algebra-lemma-minimal-prime-reduced-ring} 참조).

\begin{lemma}
\label{more-morphisms-lemma-flat-and-free-at-point-fibre}
$f : X \to S$를 유한 표시인 스킴의 사상이라 하고, $\mathcal{F}$를
유한 표시 $\mathcal{O}_X$-가군이라 하자. $x \in X$의 상이 $s \in S$라고
하자. $\mathcal{F}$가 $x$에서 $S$ 위에서 평탄하고 $(\mathcal{F}_s)_x$가
평탄 $\mathcal{O}_{X_s, x}$-가군이면, $\mathcal{F}$는 $x$의 한 근방에서
유한 자유이다.
\end{lemma}

\begin{proof}
$\mathcal{F}_x \otimes \kappa(x)$가 영이면 나카야마 보조정리(Algebra,
Lemma \ref{algebra-lemma-NAK})에 의해 $\mathcal{F}_x = 0$이고, 따라서
$\mathcal{F}$는 $x$의 한 근방에서 영이다(Modules, Lemma
\ref{modules-lemma-finite-type-stalk-zero}). 그러므로 보조정리가 성립한다.
따라서 $\mathcal{F}_x \otimes \kappa(x)$가 영이 아니라고 가정할 수 있고,
$f = \text{id} : X \to X$에 Theorem
\ref{more-morphisms-theorem-criterion-flatness-fibre}를 적용할 수 있다. 그 결과
$\mathcal{F}_x$는 $\mathcal{O}_{X, x}$ 위에서 평탄하다. 따라서
$\mathcal{F}_x$는 자유이다. 예를 들어 Algebra, Lemma
\ref{algebra-lemma-finite-flat-local}를 보라. 열린 근방
$x \in U \subset X$와 단면들 $s_1, \ldots, s_r \in \mathcal{F}(U)$를
택하되, 이들이 $\mathcal{F}_x$에서 기저로 가게 하자. 대응하는 사상
$\psi : \mathcal{O}_U^{\oplus r} \to \mathcal{F}|_U$는 $U$를 줄인 뒤
전사이다(Modules, Lemma \ref{modules-lemma-finite-type-stalk-zero}).
그러면 $\Ker(\psi)$는 유한형이고(Modules, Lemma
\ref{modules-lemma-kernel-surjection-finite-free-onto-finite-presentation}
참조), $\Ker(\psi)_x = 0$이다. 따라서 $U$를 한 번 더 줄이면 $\psi$는
동형이다.
\end{proof}

\begin{lemma}
\label{more-morphisms-lemma-finite-free-open}
$f : X \to S$를 국소 유한 표시인 스킴의 사상이라 하자.
$\mathcal{F}$를 유한 표시 $\mathcal{O}_X$-가군이며 $S$ 위에서 평탄하다고 하자.
그러면 집합
$$
\{x \in X : \mathcal{F}\text{는 }x\text{의 한 근방에서 자유이다}\}
$$
은 $X$에서 열려 있고 그 구성은 임의의 밑변환 $S' \to S$와 가환한다.
\end{lemma}

\begin{proof}
열림성은 자명하다. $x \in X$의 상이 $s \in S$라고 하자. Lemma
\ref{more-morphisms-lemma-flat-and-free-at-point-fibre}에 의해 $x$가 이 집합에 속할
필요충분조건은 $\mathcal{F}|_{X_s}$가 $x$에서 $X_s$ 위에서 평탄한 것이다.
이는 또한 $\mathcal{F}$가 $x$에서 $X$ 위에서 평탄한 것과 동치임이
명백하다(이 진술은 $\mathcal{F}_x$의 자유성으로부터 따르고,
$\mathcal{F}|_{X_s}$가 $x$에서 $X_s$ 위에서 평탄함을 함의하기 때문이다).
따라서 밑변환 진술은 Lemma
\ref{more-morphisms-lemma-base-change-criterion-flatness-fibre}를 $\text{id} : X \to X$에
$S$ 위에서 적용하여 따른다.
\end{proof}






\section{매끄러운 스킴 사이의 닫힌 몰입}
\label{more-morphisms-section-closed-immersions-smooth}

\noindent
다른 곳에는 잘 들어맞지 않는 몇 가지 결과를 모은다.

\begin{lemma}
\label{more-morphisms-lemma-etale-local-structure}
$S$를 스킴이라 하자. $Y \to X$를 $S$ 위에서 매끄러운 스킴들의 닫힌
몰입이라 하자. 모든 $y \in Y$에 대해 정수 $0 \leq m, n$과 가환 그림
$$
\xymatrix{
Y \ar[d] &
V \ar[l] \ar[d] \ar[r] &
\mathbf{A}^m_S
\ar[d]^{(a_1, \ldots, a_m) \mapsto (a_1, \ldots, a_m, 0 \ldots, 0)} \\
X &
U \ar[l] \ar[r]^-\pi &
\mathbf{A}^{m + n}_S
}
$$
이 존재하여 $U \subset X$는 열려 있고, $V = Y \cap U$이며,
$\pi$는 \'etale이고, $V = \pi^{-1}(\mathbf{A}^m_S)$이며, $y \in V$이다.
\end{lemma}

\begin{proof}
이 문제는 $X$ 위에서 국소적이므로 $X$를 $y$의 열린 근방으로 바꿀 수
있다. Divisors, Lemma
\ref{divisors-lemma-immersion-smooth-into-smooth-regular-immersion}에 의해
$Y \to X$가 정칙 몰입이므로, $X = \Spec(A)$가 아핀이고
$Y = V(f_1, \ldots, f_n)$가 되게 하는 정칙열
$f_1, \ldots, f_n \in A$가 존재한다고 가정할 수 있다. $X$를 (따라서
$Y$도) 더 줄이면 \'etale 사상 $Y \to \mathbf{A}^m_S$가 존재한다고
가정할 수 있다. Morphisms, Lemma
\ref{morphisms-lemma-smooth-etale-over-affine-space}를 보라.
$\overline{g}_1, \ldots, \overline{g}_m$을 $\mathcal{O}_Y(Y)$에 속하는 이
\'etale 사상의 좌표함수들이라 하자. 이 함수들의 올림
$g_1, \ldots, g_m \in A$를 택하고 사상
$$
(g_1, \ldots, g_m, f_1, \ldots, f_n) :
X
\longrightarrow
\mathbf{A}^{m + n}_S
$$
을 $S$ 위에서 생각하자. 이는 $S$ 위에서 국소 유한 표시인 스킴들의 사상이므로
국소 유한 표시이다(Morphisms, Lemma
\ref{morphisms-lemma-finite-presentation-permanence}). 이 사상을
$\mathbf{A}^m_S \subset \mathbf{A}^{m + n}_S$에 제한하면 구성에 의해
\'etale이다. 따라서 $X \to \mathbf{A}^{m + n}_S$가 $y$에서 \'etale임을
보이려면 $X \to \mathbf{A}^{m + n}_S$가 $y$에서 평탄함을 보이면 충분하다.
Morphisms, Lemma \ref{morphisms-lemma-etale-at-point}를 보라.
$s \in S$를 $y$의 상이라 하자. $X_s \to \mathbf{A}^{m + n}_s$가
$y$에서 평탄함을 확인하면 충분하다. Theorem
\ref{more-morphisms-theorem-criterion-flatness-fibre}를 보라.
$z \in \mathbf{A}^{m + n}_s$를 $y$의 상이라 하자. 국소환 사상
$$
\mathcal{O}_{\mathbf{A}^{m + n}_s, z}
\longrightarrow
\mathcal{O}_{X_s, y}
$$
은 Algebra, Lemma \ref{algebra-lemma-CM-over-regular-flat}에 의해 평탄하다.
실제로 체 위에서 매끄러운 스킴은 정칙이고 정칙환은 코언--매콜리이다.
Varieties, Lemma \ref{varieties-lemma-smooth-regular} 및 Algebra, Lemma
\ref{algebra-lemma-regular-ring-CM}를 보라. 따라서 정의역과 공역은 모두
정칙 국소환이다(따라서 CM이다). 정의역과 공역의 차원은 같다. 실제로
$\dim(\mathcal{O}_{Y_s, y}) = \dim(\mathcal{O}_{\mathbf{A}^m_s, z})$
가 More on Algebra, Lemma
\ref{more-algebra-lemma-dimension-etale-extension}에 의해 성립하고,
$\dim(\mathcal{O}_{\mathbf{A}^{m + n}_s, z}) = n +
\dim(\mathcal{O}_{\mathbf{A}^m_s, z})$이며,
$\dim(\mathcal{O}_{X_s, y}) = n + \dim(\mathcal{O}_{Y_s, y})$
이다. 마지막 등식은 $\mathcal{O}_{Y_s, y}$가 $\mathcal{O}_{X_s, y}$를
길이 $n$의 정칙열 $f_1, \ldots, f_n$로 나눈 몫이기 때문이다(Divisors,
Remark \ref{divisors-remark-relative-regular-immersion-elements} 참조).
끝으로 $Y_s \to \mathbf{A}^m_s$가 $y$에서 \'etale이므로 위에 표시된
화살표의 올환은 $\kappa(z)$ 위에서 유한이다. 이것으로 증명이 끝난다.
\end{proof}

\begin{remark}
\label{more-morphisms-remark-relative-blowup}
환 $R$을 고정하고 $S = \Spec(R)$라 놓자. 정수 $0 \leq m$과
$1 \leq n$을 고정하자. 닫힌 몰입
% P05-EN-CH37-0084
$$
Z = \mathbf{A}^m_S \longrightarrow \mathbf{A}^{m + n}_S = X,\quad
(a_1, \ldots, a_m) \mapsto (a_1, \ldots, a_m, 0, \ldots 0).
$$
을 생각하자. 닫힌 부분스킴 $Z$를 중심으로 한 $X$의 블로업을 $X'$이라
하자. 다음과 같이 쓰자.
$$
X = \Spec(A)
\quad\text{with}\quad
A = R[x_1, \ldots, x_m, y_1, \ldots, y_n]
$$
그러면 $X'$은 아이디얼 $(y_1, \ldots, y_n)$에 관한 $A$의 리스 대수의
Proj이다. 이 리스 대수는 $B = A[T_1, \ldots, T_n]/(y_iT_j - y_jT_i)$와
같다. 세부사항은 생략한다. 따라서 $X' = \text{Proj}(B)$는 다음 아핀
열린집합들로 덮이므로 $S$ 위에서 매끄럽다.
% P05-EN-CH37-0085
\begin{align*}
D_+(T_i)
& =
\Spec(B_{(T_i)}) \\
& =
\Spec(A[t_1, \ldots, \hat t_i, \ldots t_n]/(y_j - y_i t_j)) \\
& =
\Spec(R[x_1, \ldots, x_m, y_i, t_1, \ldots, \hat t_i, \ldots, t_n])
\end{align*}
이들은 $\mathbf{A}^{n + m}_S$와 동형이다. 이 좌표도표에서 예외 인자는
$y_i = 0$으로 잘려 나오므로 예외 인자도 $S$ 위에서 매끄럽다.
\end{remark}

\begin{lemma}
\label{more-morphisms-lemma-blowup}
$S$를 스킴이라 하자. $Z \to X$를 $S$ 위에서 매끄러운 스킴들의 닫힌
몰입이라 하자. $b : X' \to X$를 $Z$의 블로업이라 하고 예외 인자를
$E \subset X'$라 하자. 그러면 $X'$과 $E$는 $S$ 위에서 매끄럽다.
사상 $p : E \to Z$는 사영공간 다발
$$
\mathbf{P}(\mathcal{I}/\mathcal{I}^2) \longrightarrow Z
$$
과 표준적으로 동형이다. 여기서 $\mathcal{I} \subset \mathcal{O}_X$는
$Z$의 아이디얼 층이다. 사영공간 다발 구조에서 오는 상대
$\mathcal{O}_E(1)$은 $\mathcal{O}_{X'}(-E)$를 $E$로 제한한 것과 동형이다.
\end{lemma}

\begin{proof}
Divisors, Lemma
\ref{divisors-lemma-immersion-smooth-into-smooth-regular-immersion}에 의해
% P05-EN-CH37-0086
몰입 $Z \to X$는 정칙 몰입이다. 따라서 아이디얼 층 $\mathcal{I}$는
유한형이고, 따라서 $b$는 상대 풍부 가역층
$\mathcal{O}_{X'}(1) = \mathcal{O}_{X'}(-E)$을 가진 사영 사상이다.
Divisors, Lemmas
\ref{divisors-lemma-blowing-up-gives-effective-Cartier-divisor} 및
\ref{divisors-lemma-blowing-up-projective}를 보라. 표준 사상
$\mathcal{I} \to b_*\mathcal{O}_{X'}(1)$은 닫힌 몰입
$$
X' \longrightarrow
\mathbf{P}\left(\bigoplus\nolimits_{n \geq 0}
\text{Sym}^n_{\mathcal{O}_X}(\mathcal{I})\right)
$$
을 주는데, 이는 바로 블로업의 구성에서 따른다. 이 사상을 $E$로
제한하면 표준 사상
$$
E \longrightarrow
\mathbf{P}\left(\bigoplus\nolimits_{n \geq 0}
\text{Sym}^n_{\mathcal{O}_Z}(\mathcal{I}/\mathcal{I}^2)\right)
$$
을 $Z$ 위에서 얻는다. $\mathcal{I}/\mathcal{I}^2$는 유한 국소 자유이다.
% P05-EN-CH37-0087
이 표준 사상이 동형이면 보조정리의 마지막 부분이 성립한다. 이제 이 모든
사실을 말했으므로 문제는 $X$ 위에서 \'etale 국소적이다. 실제로
Divisors, Lemma \ref{divisors-lemma-flat-base-change-blowing-up}에 의해
블로업은 평탄 밑변환과 가환하고, 전사 \'etale 사상과 앞합성한 뒤
매끄러움을 검사할 수 있다. 따라서 Lemma
\ref{more-morphisms-lemma-etale-local-structure}에 주어진 $S$ 위 스킴의 닫힌 몰입의
\'etale 국소 구조에 의해, Remark \ref{more-morphisms-remark-relative-blowup}에서 논한
경우로 환원된다.
\end{proof}








\section{평탄 가군과 상대 동반점}
\label{more-morphisms-section-flat-relative-assassin}

\noindent
이 절에서는 기하학적 상황에서 평탄 가군의 받침이 (어떤 의미에서)
밑공간 위에서 등차원임을 증명한다. 뇌터인 경우에 대해서는
\cite[IV Proposition 12.1.1.5]{EGA}를 보라. 먼저 두 개의 보조정리를
증명한다.

\begin{lemma}
\label{more-morphisms-lemma-mod-injective-valuation-ring}
$A$를 값매김환이라 하자.
% P05-EN-CH37-0088
$A \to B$를 본질적 유한형인 국소환의 국소 준동형이라 하자.
$u : N \to M$을 유한 $B$-가군의 사상이라 하자. $M$이 $A$ 위에서
평탄하고 $\overline{u} : N/\mathfrak m_A N \to M/\mathfrak m_A M$이
단사라고 가정하자. 그러면 $u$는 단사이고 $M/u(N)$은 $A$ 위에서
평탄하다.
\end{lemma}

\begin{proof}
이 보조정리를 Algebra, Lemma \ref{algebra-lemma-mod-injective-general}에서
도출하겠다($M$과 $N$의 역할을 바꾸었음에 유의하자). 이 환원을 위해
More on Algebra, Lemma
\ref{more-algebra-lemma-valuation-ring-flat-essentially-finite-type}를 사용하여
유한 표시인 경우로 환원한다.

\medskip\noindent
가정에 의해 $B$를 다항식 대수 $P = A[x_1, \ldots, x_n]$을
소아이디얼 $\mathfrak q$에서 국소화한 것의 몫으로 쓸 수 있다. 그러면
$u : N \to M$을 유한 $P_\mathfrak q$-가군의 사상으로 생각할 수 있다.
따라서 $B$가 $A$ 위에서 본질적 유한 표시라고 가정할 수 있고 그렇게
가정한다.

\medskip\noindent
다음으로 $B$-가군 $N$은 유한이지만 유한 표시가 아닐 수도 있다.
$N = \colim N_\lambda$를 전이 사상이 전사인 유한 표시 $B$-가군들의
여과 쌍극한으로 쓰자. 예를 들어 표시
$0 \to K \to B^{\oplus r} \to N \to 0$을 택하고, $K$를 유한 부분가군들
$K_\lambda$의 합집합으로 쓴 뒤,
$N_\lambda = \Coker(K_\lambda \to B^{\oplus r})$라 놓는다.
가군 $N/\mathfrak m_A N$은 뇌터환 $B/\mathfrak m_A B$ 위에서 유한
표시이다. 따라서 충분히 큰 $\lambda$에 대해
$N_\lambda/\mathfrak m_A N_\lambda = N/\mathfrak m_A N$이다. 이제 합성
$u_\lambda : N_\lambda \to M$에 대해 보조정리를 보일 수 있으면
$N_\lambda = N$이고 결과가 $u$에 대해 성립한다고 결론지을 수 있다.
따라서 $N$이 $B$ 위에서 유한 표시라고 가정할 수 있고 그렇게 가정한다.

\medskip\noindent
More on Algebra, Lemma
\ref{more-algebra-lemma-valuation-ring-flat-essentially-finite-type}에 의해
가군 $M$은 $B$ 위에서 유한 표시이다. 따라서 Algebra, Lemma
\ref{algebra-lemma-mod-injective-general}의 모든 가정이 성립하므로 결론을
얻는다.
\end{proof}

\begin{lemma}
\label{more-morphisms-lemma-make-base-change}
\begin{reference}
이는 \cite[IV Proposition 12.1.1.5]{EGA}의 증명에서 찾을 수 있다.
\end{reference}
$f : X \to S$를 스킴의 사상이라 하자. $y \in X$를 상이 $t \in S$인
점이라 하자.
% P05-EN-CH37-0089
$Y \subset X$를 $X$의 정수적 닫힌 부분스킴으로 본 $\{y\}$의 폐포라
하자. $s \in S$라 하고, $x \in Y_s$를 $Y_s$의 한 기약 성분의 일반점이라
하자. 다음 성질들을 가진 데카르트 그림
$$
\xymatrix{
X' \ar[r]_{g'} \ar[d]_{f'} & X \ar[d]^f \\
S' \ar[r]^g & S
}
$$
이 존재한다.
\begin{enumerate}
\item $S'$은 일반점이 $t'$이고 닫힌점이 $s'$인 값매김환의 스펙트럼이다.
\item $g(t') = t$이고 $g(s') = s$이다.
\item $(S' \times_S Y)_{t'} = Y_t \times_t t'$의 한 기약 성분의 일반점이며
$g'(y') = y$를 만족하는 점 $y' \in X'_{t'}$가 존재한다.
\item $Y' \subset X'$를 $X'$의 정수적 닫힌 부분스킴으로 본 $\{y'\}$의
폐포라 하면, $Y'_{s'}$의 한 기약 성분의 일반점이며 $g'(x') = x$를
만족하는 점 $x' \in Y'_{s'}$가 존재한다.
\end{enumerate}
\end{lemma}

\begin{proof}
값매김환 $R$을 택하고, 일반점이 $t'$이고 닫힌점이 $s'$인
$S' = \Spec(R)$라 놓은 뒤, $h(t') = y$ 및 $h(s') = x$를 만족하는 사상
$h : S' \to X$를 택한다. Schemes, Lemma
\ref{schemes-lemma-points-specialize}를 보라. $g(t') = t$ 및 $g(s') = s$가
되도록 $g = f \circ h$라 놓는다. 밑변환
$$
\xymatrix{
X' \ar[r]_{g'} \ar[d] & X \ar[d] \\
S' \ar@/^1em/[u]^\sigma \ar[r]^-g & S
}
$$
을 생각하자. $h = g' \circ \sigma$를 만족하는 이 밑변환의 단면
$\sigma$를 얻는다.

\medskip\noindent
물론 $h$가 $Y$를 통하여 인수분해되므로 $\sigma$는 $Y$의 밑변환
$S' \times_S Y$를 통하여 인수분해된다. $y' \in X'_{t'} \subset X'$를 올
$$
(S' \times_S Y)_{t'} = Y_t \times_t t'
$$
에서 점 $\sigma(t')$를 포함하는 한 기약 성분의 일반점, 즉
$y' \leadsto \sigma(t')$를 만족하는 점이라 하자.
% P05-EN-CH37-0090
$g'(y') \in Y_t$이고 $g(y') \leadsto g(\sigma(t')) = y$이므로,
$y$가 올 $Y_t$의 일반점이라는 사실에서 $g'(y') = y$를 얻는다.
% P05-EN-CH37-0091
$Y' \subset X'$를 $\{y'\}$의 $X'$에서의 폐포에 주는 정수적 닫힌
부분스킴 구조라 하자. $\sigma(t') \in Y'$이므로 $\sigma$는 $Y'$를
통하여 인수분해된다. $x' \in Y'_{s'}$를 $\sigma(s')$를 포함하는
$Y'_{s'}$의 한 기약 성분의 일반점으로 택하자. 즉,
$x' \leadsto \sigma(s')$를 얻고, 따라서
$g'(x') \leadsto g'(\sigma(s')) = x$이다. 다시 말해 가정에 의해 $x$가
$Y_s$의 한 기약 성분의 일반점이고 $g'(Y') \subset Y$이므로
$g'(x') = x$라고 결론짓는다.
\end{proof}

\begin{lemma}
\label{more-morphisms-lemma-associated-point-specializes}
$f : X \to S$를 국소 유한형인 스킴의 사상이라 하자. $\mathcal{F}$를
준연접 유한형 $\mathcal{O}_X$-가군이라 하자. 상이 $t \in S$인
$y \in \text{Ass}_{X/S}(\mathcal{F})$라 하자.
% P05-EN-CH37-0092
$Y \subset X$를 $\{y\}$의 $X$에서의 폐포에 주는 정수적 닫힌 부분스킴
구조라 하자. $s \in S$라 하고, $x \in Y_s$를 $Y_s$의 한 기약 성분의
일반점이라 하자. $\mathcal{F}$가 $S$ 위에서 $x$에서 평탄하면,
$x \in \text{Ass}_{X/S}(\mathcal{F})$이고
$\dim_x(Y_s) = \dim(Y_t)$이다.
\end{lemma}

\begin{proof}
Lemma \ref{more-morphisms-lemma-make-base-change}에서와 같은 그림을 택하자.
$\mathcal{F}' = (g')^*\mathcal{F}$라 놓는다. Divisors, Lemma
\ref{divisors-lemma-base-change-relative-assassin}에 의해
$y' \in \text{Ass}_{X'/S'}(\mathcal{F}')$이다. $y'$의 선택에 의해 또한
$\dim(Y'_{t'}) = \dim(Y_t)$임을 알 수 있다. 예를 들어 Algebra, Lemma
\ref{algebra-lemma-inequalities-under-field-extension}을 보라. Algebra,
Lemma \ref{algebra-lemma-finite-type-domain-over-valuation-ring-dim-fibres}에
의해 $Y'_{s'}$는 차원이 $\dim(Y_t)$와 같은 등차원 스킴이다.
$\mathcal{F}$가 $x$에서 $S$ 위에서 평탄하므로 $\mathcal{F}'$는
$x'$에서 $S'$ 위에서 평탄하다. Morphisms, Lemma
\ref{morphisms-lemma-base-change-module-flat}를 보라.

\medskip\noindent
% P05-EN-CH37-0093
$x' \in \text{Ass}_{X'/S}(\mathcal{F}')$임을 보일 수 있다고 하자. 그러면
Divisors, Lemma \ref{divisors-lemma-base-change-relative-assassin}에 의해
$x \in \text{Ass}_{X/S}(\mathcal{F})$이고, $x'$를 포함하는 $Y'_{s'}$의
기약 성분 $C'$은 $x$를 포함하는 기약 성분 $C \subset Y_s$에 대한
$C \times_s s'$의 한 기약 성분이다. 따라서
$\dim(C) = \dim(C') = \dim(Y_t)$이고(위를 보라) 증명이 완성된다.
이로써 다음 단락에서 논할 경우로 환원된다.

\medskip\noindent
$A$가 값매김환인 $S = \Spec(A)$라고 가정하고, $t$와 $s$를 $S$의
일반점과 닫힌점이라 하자. 모순을 얻기 위해
$x \not \in \text{Ass}_{X/S}(\mathcal{F})$라고 가정하겠다. 다시 말해,
$\mathcal{F}_s$가 $\mathcal{F}$를 $X_s$로 당긴 것일 때
$x \not \in \text{Ass}_{X_s}(\mathcal{F}_s)$라고 가정한다. 환 사상
$$
A \longrightarrow \mathcal{O}_{X, x} = B
$$
과 $B = \mathcal{O}_{X, x}$ 위의 가군 $N = \mathcal{F}_x$를 생각하자.
그러면 $B/\mathfrak m_A B = \mathcal{O}_{X_s, x}$이고
$N/\mathfrak m_A N$은 점 $x$에서 $\mathcal{F}_s$의 줄기이다.
% P05-EN-CH37-0094
$\mathfrak q \subset B$를 점 $y$에 대응하는 소아이디얼이라 하자.
Schemes, Lemma \ref{schemes-lemma-specialize-points}를 보라. $x$가 $Y_s$의
일반점이므로 $\mathfrak q + \mathfrak m_A B$의 근기는 $\mathfrak m_B$이다.
그러면 $\text{Ass}_{B/\mathfrak m_A B}(N/\mathfrak m_A N)$은 소아이디얼들의
유한집합인데(Algebra, Lemma \ref{algebra-lemma-finite-ass}),
$x \not \in \text{Ass}_{X/S}(\mathcal{F})$이므로 $B/\mathfrak m_A B$의
극대아이디얼을 포함하지 않는다.
% P05-EN-CH37-0095
따라서 $B/\mathfrak m_A B$에서 $\mathfrak q$의 상은 이 소아이디얼들 중
어느 것에도 포함되지 않는다. 그러므로 소아이디얼 회피(Algebra, Lemma
\ref{algebra-lemma-silly})에 의해, $B/\mathfrak m_A B$에서의 상이
$N/\mathfrak m_A N$ 위의 영인자가 아닌 원소 $g \in \mathfrak q$를 찾을
수 있다(이는 Algebra, Lemma \ref{algebra-lemma-ass-zero-divisors}의
영인자에 대한 기술을 사용한다).
% P05-EN-CH37-0096
$N = \mathcal{F}_x$는 Lemma \ref{more-morphisms-lemma-mod-injective-valuation-ring}에 의해
$A$-평탄이므로 다음 사상은 단사이다.
$$
g : N \longrightarrow N
$$
특히 $K = \text{Frac}(A)$가 $A$의 분수체이면 다음 사상이 단사임을 알 수
있다.
$$
g : N \otimes_A K \longrightarrow N \otimes_A K
$$
$\mathfrak q$가 $B \otimes_A K$의 소아이디얼에 대응함에 유의하자.
% P05-EN-CH37-0097
$\mathcal{F}_t$를 $\mathcal{F}$의 일반올 $X_t$로의 제한이라 하자.
$(B \otimes_A K)_{\mathfrak q} = \mathcal{O}_{X_t, y}$이고,
$(N \otimes_A K)_\mathfrak q$는 $y$에서 $\mathcal{F}_t$의 줄기이다.
따라서 $g \in \mathfrak m_y \subset \mathcal{O}_{X_t, y}$가 줄기
$(\mathcal{F}_t)_y$ 위의 영인자가 아님을 얻는다. 이는
$y \in \text{Ass}_{X/S}(\mathcal{F})$라는 가정과 모순이다.
\end{proof}

\begin{lemma}
\label{more-morphisms-lemma-relative-dimension-support-flat}
$f : X \to S$를 국소 유한형인 스킴의 사상이라 하자. $\mathcal{F}$를
$S$ 위에서 평탄한 유한형 준연접 $\mathcal{O}_X$-가군이라 하자.
$S$가 일반점 $\eta$를 가진 기약 스킴이라고 가정하자.
$\dim(\text{Supp}(\mathcal{F}_\eta)) \leq r$이면 모든 $s \in S$에 대해
$\dim(\text{Supp}(\mathcal{F}_s)) \leq r$이다.
\end{lemma}

\begin{proof}
$x \in \text{Supp}(\mathcal{F}_s)$를 $\text{Supp}(\mathcal{F}_s)$의 한 기약
성분의 일반점이라 하자. Algebra, Lemma
\ref{algebra-lemma-going-down-flat-module}에 의해 $f(y) = \eta$를 만족하는
$\text{Supp}(\mathcal{F})$에서의 특수화 $y \leadsto x$를 찾을 수 있다.
물론 $y$가 $\text{Supp}(\mathcal{F}_\eta)$의 한 기약 성분의 일반점이라고
가정할 수 있다. Lemma \ref{more-morphisms-lemma-associated-point-specializes}로부터
$\overline{\{x\}}$의 차원이 $r$ 이하라고 결론짓는다.
\end{proof}

\begin{lemma}
\label{more-morphisms-lemma-flat-associated-equidimensional}
$f : X \to S$를 국소 유한형인 스킴의 사상이라 하자. $\mathcal{F}$를
유한형 준연접 $\mathcal{O}_X$-가군이라 하자.
$y \in \text{Ass}_{X/S}(\mathcal{F})$라 하자.
% P05-EN-CH37-0098
$Y \subset X$를 $\{y\}$의 $X$에서의 폐포에 주는 정수적 닫힌 부분스킴
구조라 하고, $T \subset S$를 정수적 닫힌 부분스킴으로 본 $\{f(y)\}$의
폐포라 하자. 가환 그림
$$
\xymatrix{
Y \ar[r] \ar[d] & X \ar[d] \\
T \ar[r] & S
}
$$
을 얻으며, 여기서 $Y \to T$는 우세하다. $\mathcal{F}$가 $Y \to T$의
올의 기약 성분들의 모든 일반점에서 $S$ 위에서 평탄하다고 가정하자
(예를 들어 $\mathcal{F}$가 $S$ 위에서 평탄하면 그러하다). 그러면
\begin{enumerate}
\item $s \in S$이고 $x \in Y_s$가 $Y_s$의 한 기약 성분의 일반점이면
$x \in \text{Ass}_{X/S}(\mathcal{F})$이다.
\item $Y \to T$의 상대차원이 $d$가 되게 하는 정수 $d \geq 0$이 존재한다.
Morphisms, Definition \ref{morphisms-definition-relative-dimension-d}를 보라.
\end{enumerate}
\end{lemma}

\begin{proof}
% P05-EN-CH37-0099
이는 Lemma \ref{more-morphisms-lemma-associated-point-specializes}의 점별 버전으로부터
즉시 따른다. 국소 대수적 스킴 $Y_s$의 차원을 계산하려면 일반점들의
근방을 살펴보면 충분함에 유의하자. Varieties, Section
\ref{varieties-section-algebraic-schemes}을 보라.
\end{proof}

\begin{remark}
\label{more-morphisms-remark-flat-equidimensional}
위의 내용, 특히 Lemma \ref{more-morphisms-lemma-flat-associated-equidimensional}를 이용하여
스킴의 사상 $f : X \to S$가 Morphisms, Definition
\ref{morphisms-definition-relative-dimension-d}에서 정의한 상대차원 $d$를
가진다고 결론지을 수 있는 몇 가지 경우를 제시한다. 예를 들어 다음이
성립하면 그러하다.
\begin{enumerate}
\item $X$는 일반점 $\xi$를 가진 정수적 스킴이다.
\item $\kappa(f(\xi))$ 위의 $\kappa(\xi)$의 초월차수는 $d$이다.
\item $f$는 국소 유한형이다.
\item $S$ 위에서 평탄하고 $\text{Supp}(\mathcal{F}) = X$를 만족하는
유한형 준연접 $\mathcal{O}_X$-가군 $\mathcal{F}$가 존재한다.
\end{enumerate}
다음과 같은 또 다른 가정들도 충분하다.
\begin{enumerate}
\item $S$는 일반점 $\eta$를 가진 기약 스킴이다.
\item $X_\eta$는 $X$에서 조밀하다.
\item $X_\eta$의 모든 기약 성분의 차원은 $d$이다.
\item $f$는 국소 유한형이다.
\item $S$ 위에서 평탄하고 $\text{Supp}(\mathcal{F}) = X$를 만족하는
유한형 준연접 $\mathcal{O}_X$-가군 $\mathcal{F}$가 존재한다.
\end{enumerate}
물론 $\mathcal{F}$에 대한 평탄성 조건을 완화하여 $\mathcal{F}$가
여차원 $0$에서 $S$ 위에서 평탄하다는 것, 즉 $\mathcal{F}$가 모든 올의
모든 일반점에서 $S$ 위에서 평탄하다는 것만 요구할 수도 있다.
% P05-EN-CH37-0100
이 결과들이 필요해지면
여기에 주의 깊게 진술하고 증명할 것이다.
\end{remark}











\section{정규화 다시 보기}
\label{more-morphisms-section-normalization}

\noindent
정규화는 매끄러운 밑변환과 가환한다.

\begin{lemma}
\label{more-morphisms-lemma-integral-closure-smooth-pullback}
$f : Y \to X$를 스킴의 매끄러운 사상이라 하자. $\mathcal{A}$를 준연접
$\mathcal{O}_X$-대수층이라 하자. $f^*\mathcal{A}$에서 $\mathcal{O}_Y$의
정수적 폐포는 $f^*\mathcal{A}'$와 같다. 여기서
$\mathcal{A}' \subset \mathcal{A}$은 $\mathcal{A}$에서 $\mathcal{O}_X$의
정수적 폐포이다.
\end{lemma}

\begin{proof}
이는 Algebra, Lemma \ref{algebra-lemma-integral-closure-commutes-smooth}를
스킴의 언어로 옮긴 것이다. 세부사항은 생략한다.
\end{proof}

\begin{lemma}[정규화는 매끄러운 밑변환과 가환한다]
\label{more-morphisms-lemma-normalization-smooth-localization}
다음을 스킴의 범주에서의 올곱 그림이라 하자.
$$
\xymatrix{
Y_2 \ar[r] \ar[d]_{f_2} & Y_1 \ar[d]^{f_1} \\
X_2 \ar[r]^\varphi & X_1
}
$$
$f_1$이 준콤팩트이고 준분리이며, $\varphi$가 매끄럽다고 가정하자.
$Y_i \to X_i' \to X_i$를 $Y_i$에서 $X_i$의 정규화라 하자. 그러면
$X_2' \cong X_2 \times_{X_1} X_1'$이다.
\end{lemma}

\begin{proof}
인수분해 $Y_1 \to X_1' \to X_1$을 $X_2$로 밑변환하면 인수분해
$Y_2 \to X_2 \times_{X_1} X_1' \to X_2$가 되고,
$X_2 \times_{X_1} X_1' \to X_2$는 정수적이다(Morphisms, Lemma
\ref{morphisms-lemma-base-change-finite}). 따라서 Morphisms, Lemma
\ref{morphisms-lemma-characterize-normalization}의 보편 성질에 의해 사상
$h : X_2' \to X_2 \times_{X_1} X_1'$를 얻는다. $X_2'$은
$f_{2, *}\mathcal{O}_{Y_2}$에서 $\mathcal{O}_{X_2}$의 정수적 폐포의 상대
스펙트럼임에 유의하자. $\mathcal{A}' \subset f_{1, *}\mathcal{O}_{Y_1}$가
$\mathcal{O}_{X_1}$의 정수적 폐포를 나타내면,
$X_2 \times_{X_1} X_1'$은 $\varphi^*\mathcal{A}'$의 상대 스펙트럼이다.
Constructions, Lemma \ref{constructions-lemma-spec-properties}를 보라.
Cohomology of Schemes, Lemma
\ref{coherent-lemma-flat-base-change-cohomology}에 의해
$f_{2, *}\mathcal{O}_{Y_2} = \varphi^*f_{1, *}\mathcal{O}_{Y_1}$임을 안다.
따라서 결과는 Lemma \ref{more-morphisms-lemma-integral-closure-smooth-pullback}에서 따른다.
\end{proof}

\begin{lemma}[정규화와 매끄러운 사상]
\label{more-morphisms-lemma-normalization-and-smooth}
$X \to Y$를 스킴의 매끄러운 사상이라 하자. $Y$의 모든 준콤팩트
열린집합이 유한 개의 기약 성분을 가진다고 가정하자. 그러면 $X$에
대해서도 마찬가지이고, $X^\nu$, $Y^\nu$가 각각 $X$, $Y$의 정규화일 때
$X$ 위의 유일한 동형 $X^\nu = X \times_Y Y^\nu$가 존재한다.
\end{lemma}

\begin{proof}
Descent, Lemma \ref{descent-lemma-locally-finite-nr-irred-local-fppf}에 의해
$X$의 모든 준콤팩트 열린집합은 유한 개의 기약 성분을 가진다. 축약
스킴 위에서 매끄러운 스킴은 축약이므로
$X_{red} = X \times_Y Y_{red}$임에 유의하자. Descent, Lemma
\ref{descent-lemma-reduced-local-smooth}을 보라. 따라서 $X$와 $Y$가
축약이라고 가정할 수 있다(정의에 의해 스킴의 정규화는 그 축약의
정규화와 같기 때문이다). 다음으로 Descent, Lemma
\ref{descent-lemma-normal-local-smooth}에 의해
$X' = X \times_Y Y^\nu$가 정규 스킴임에 유의하자. 사상
$X' \to Y^\nu$는 매끄럽고(따라서 평탄하고), 그러므로 $X'$의 기약
성분의 일반점들은 $Y^\nu$의 기약 성분의 일반점들 위에 놓인다.
$Y^\nu \to Y$가 쌍유리이므로 $X' \to X$도 쌍유리라고 결론짓는다
($X' \to Y^\nu$가 $Y$의 일반점들 위의 올에서 동형을 유도하기 때문이다).
따라서 인수분해 $X^\nu \to X' \to X$가 존재한다. Morphisms, Lemma
\ref{morphisms-lemma-normalization-normal}를 보라. $X'$이 정규이고 $X$
위에서 정수적이므로 이 인수분해는 동형이다.
\end{proof}

\begin{lemma}[정규화와 헨젤화]
\label{more-morphisms-lemma-normalization-and-henselization}
$X$를 국소 뇌터 스킴이라 하고 $\nu : X^\nu \to X$를 정규화 사상이라
하자. 그러면 임의의 점 $x \in X$에 대해 밑변환
$$
X^\nu \times_X \Spec(\mathcal{O}_{X, x}^h) \to \Spec(\mathcal{O}_{X, x}^h),
\quad\text{resp.}\quad
X^\nu \times_X \Spec(\mathcal{O}_{X, x}^{sh}) \to \Spec(\mathcal{O}_{X, x}^{sh})
$$
은 $\Spec(\mathcal{O}_{X, x}^h)$, 각각
$\Spec(\mathcal{O}_{X, x}^{sh})$의 정규화이다.
\end{lemma}

\begin{proof}
$\eta_1, \ldots, \eta_r$를 $x$를 지나는 $X$의 기약 성분들의 일반점들이라
하자. 정규화를 $\Spec(\mathcal{O}_{X, x})$로 밑변환한 것은
$\prod \kappa(\eta_i)$에서 $\mathcal{O}_{X, x}$의 정수적 폐포의
스펙트럼이다. 이는 Morphisms, Definition
\ref{morphisms-definition-normalization}과 Morphisms, Lemma
\ref{morphisms-lemma-integral-closure}에서 $X$의 정규화를 구성한 방식에서
따른다. Morphisms, Lemma
\ref{morphisms-lemma-description-normalization}의 정규화 기술을 사용해도
된다. 따라서 다음 대수 문제로 환원된다. $A$를 뇌터 국소환이라 하자.
그러면 헨젤화 $A^h$와 강한 헨젤화 $A^{sh}$도 뇌터임을 상기하자(More on
Algebra, Lemma \ref{more-algebra-lemma-henselization-noetherian}).
$\mathfrak p_1, \ldots, \mathfrak p_r$을 그 극소 소아이디얼들이라 하자.
$A'$을 $\prod \kappa(\mathfrak p_i)$에서 $A$의 정수적 폐포라 하자.
문제는 $A' \otimes_A A^h$, 각각 $A' \otimes_A A^{sh}$가 뇌터 국소환
$A^h$, 각각 $A^{sh}$로부터 같은 방식으로 구성됨을 보이는 것이다.

\medskip\noindent
$A^h$, 각각 $A^{sh}$는 \'etale $A$-대수들의 쌍극한이므로,
% P05-EN-CH37-0101
$A$와 $A^{sh}$의 극소 소아이디얼들은 정확히 $A$의 극소 소아이디얼들
위에 놓이는 $A^h$, 각각 $A^{sh}$의 소아이디얼들임을 알 수 있다
(하강정리에 의한다. Algebra, Lemmas
\ref{algebra-lemma-flat-going-down} 및
\ref{algebra-lemma-minimal-prime-image-minimal-prime} 참조). 따라서 More on
Algebra, Lemma \ref{more-algebra-lemma-fibres-henselization}에 의해
$A^h \otimes_A \prod \kappa(\mathfrak p_i)$,
각각
$A^{sh} \otimes_A \prod \kappa(\mathfrak p_i)$
가 $A^h$, 각각 $A^{sh}$의 극소 소아이디얼에서의 잉여체들의 곱임을
알 수 있다. 과환에서 정수적 폐포를 취하는 것은 \'etale 밑변환과
가환한다. Algebra, Lemma
\ref{algebra-lemma-integral-closure-commutes-etale}를 보라.
% P05-EN-CH37-0102
$A^h$와 $A^{sh}$를 \'etale $A$-대수들의 극한으로 쓰면 $A^h$와
$A^{sh}$로의 밑변환에 대해서도 같은 것이 참임을 알 수 있다(더 일반적인
Algebra, Lemma \ref{algebra-lemma-integral-closure-commutes-colim-smooth}를
사용해도 된다).
\end{proof}








\section{정규 사상}
\label{more-morphisms-section-normal}

\noindent
Deligne와 Mumford의 논문 \cite{DM}에는 정규 사상의 개념이 언급된다.
이는 일련의 사상 유형들 가운데 하나일 뿐이다\footnote{
다른 유형은 coprof $\leq k$, Cohen-Macaulay, $(S_k)$,
정칙, $(R_k)$, 축약이다. \cite[IV Definition 6.8.1.]{EGA}을 참조하라.
Gorenstein 사상은
Duality for Schemes, Section \ref{duality-section-gorenstein}에서 정의한다.}
이 유형들은 모두 비슷하게 정의할 수 있다.
% P05-EN-CH37-0103: frozen source lines 5509-5516 expose a rolling future-work note.
앞으로 필요에 따라
각 유형을 별도의 절에서 추가할 것이다.

\begin{definition}
\label{more-morphisms-definition-normal}
$f : X \to Y$를 스킴의 사상이라 하자.
모든 올 $X_y$가 국소 뇌터 스킴이라고 가정하자.
\begin{enumerate}
\item $x \in X$라 하고 $y = f(x)$라 하자. $f$가 $x$에서 평탄하고
스킴 $X_y$가 $\kappa(y)$ 위에서 $x$에서 기하학적으로 정규이면
$f$가 {\it $x$에서 정규이다}라고 한다(Varieties, Definition
\ref{varieties-definition-geometrically-normal} 참조).
\item $f$가 $X$의 모든 점에서 정규이면 $f$를 {\it 정규 사상}이라 한다.
\end{enumerate}
\end{definition}

\noindent
따라서 사상 $X \to Y$가 정규라는 조건은 모든 올이 정규인
국소 뇌터 스킴이라고만 요구하는 것보다 강하다.

\begin{lemma}
\label{more-morphisms-lemma-normal}
$f : X \to Y$를 스킴의 사상이라 하자.
$f$의 모든 올이 국소 뇌터라고 가정하자.
다음은 동치이다.
\begin{enumerate}
\item $f$는 정규이다.
\item $f$는 평탄하고 그 올들은 기하학적으로 정규인 스킴이다.
\end{enumerate}
\end{lemma}

\begin{proof}
이는 정의에서 바로 따른다.
\end{proof}

\begin{lemma}
\label{more-morphisms-lemma-smooth-normal}
매끄러운 사상은 정규이다.
\end{lemma}

\begin{proof}
$f : X \to Y$를 매끄러운 사상이라 하자.
Morphisms, Lemma \ref{morphisms-lemma-smooth-locally-finite-presentation}에
따라 $f$는 국소 유한 표시이므로 올 $X_y$는 체 위의 국소 유한형이고,
따라서 국소 뇌터이다. 또한 Morphisms, Lemma
\ref{morphisms-lemma-smooth-flat}에 따라 $f$는 평탄하다.
끝으로 올 $X_y$는 체 위에서 매끄럽고(Morphisms, Lemma
\ref{morphisms-lemma-base-change-smooth}), 따라서 Varieties, Lemma
\ref{varieties-lemma-smooth-geometrically-normal}에 의해 기하학적으로 정규이다.
그러므로 Lemma \ref{more-morphisms-lemma-normal}에 따라 $f$는 정규이다.
\end{proof}

\noindent
이 개념이 매끄러운 위상에 관하여 정의역과 공역에서 국소적임을
보이고자 한다. 먼저 국소 뇌터 올을 갖는 성질을 다룬다.

\begin{lemma}
\label{more-morphisms-lemma-locally-Noetherian-fibres-fppf-local-source-and-target}
성질 $\mathcal{P}(f)=$``$f$의 올들은 국소 뇌터이다''는 정의역과
공역의 fppf 위상에서 국소적이다.
\end{lemma}

\begin{proof}
$f : X \to Y$를 스킴의 사상이라 하자.
$\{\varphi_i : Y_i \to Y\}_{i \in I}$를 $Y$의 fppf 덮개라 하자.
% P05-EN-CH37-0104: frozen source line 5586 omits “by/as” in “Denote ... the base change”.
$f_i : X_i \to Y_i$를 $\varphi_i$에 의한 $f$의 밑변환이라 쓰자.
$i \in I$이고 $y_i \in Y_i$가 한 점이라고 하자.
$y = \varphi_i(y_i)$라 놓자. 다음에 유의하라.
$$
X_{i, y_i} = \Spec(\kappa(y_i)) \times_{\Spec(\kappa(y))} X_y.
$$
더욱이 $\varphi_i$는 유한 표시이므로 체 확대
$\kappa(y_i)/\kappa(y)$는 유한 생성이다.
따라서 이 상황에서 $X_y$가 국소 뇌터일 필요충분조건은
$X_{i, y_i}$가 국소 뇌터인 것이다(Varieties, Lemma
\ref{varieties-lemma-locally-Noetherian-base-change} 참조).
이 사실은 공역에서의 국소성을 함의한다.

\medskip\noindent
$\{X_i \to X\}$를 $X$의 fppf 덮개라 하자.
$y \in Y$라 하자. 이 경우 $\{X_{i, y} \to X_y\}$는 올의
fppf 덮개이다. 따라서 정의역에서의 국소성은 Descent, Lemma
\ref{descent-lemma-Noetherian-local-fppf}에서 따른다.
\end{proof}

\begin{lemma}
\label{more-morphisms-lemma-normal-fppf-local-source-and-target}
성질
$\mathcal{P}(f)=$``$f$의 올들은 국소 뇌터이고 $f$는 정규이다''는
공역의 fppf 위상에서 국소적이고 정의역의 매끄러운 위상에서 국소적이다.
\end{lemma}

\begin{proof}
다음과 같이 쓸 수 있다.
$\mathcal{P}(f) =
\mathcal{P}_1(f) \wedge \mathcal{P}_2(f) \wedge \mathcal{P}_3(f)$
여기서
$\mathcal{P}_1(f)=$``$f$의 올들은 국소 뇌터이다'',
$\mathcal{P}_2(f)=$``$f$는 평탄하다'', 그리고
$\mathcal{P}_3(f)=$``$f$의 올들은 기하학적으로 정규이다''이다.
Lemma \ref{more-morphisms-lemma-locally-Noetherian-fibres-fppf-local-source-and-target}와
Descent, Lemmas \ref{descent-lemma-descending-property-flat},
\ref{descent-lemma-flat-fpqc-local-source}에서 이미
$\mathcal{P}_1$과 $\mathcal{P}_2$가 정의역과 공역의 fppf 위상에서
국소적임을 보았다. 따라서 $\mathcal{P}_3$를 다루면 된다.

\medskip\noindent
$f : X \to Y$를 스킴의 사상이라 하자.
$\{\varphi_i : Y_i \to Y\}_{i \in I}$를 $Y$의 fpqc 덮개라 하자.
% P05-EN-CH37-0105: frozen source line 5632 omits “by/as” in “Denote ... the base change”.
$f_i : X_i \to Y_i$를 $\varphi_i$에 의한 $f$의 밑변환이라 쓰자.
$i \in I$이고 $y_i \in Y_i$가 한 점이라고 하자.
$y = \varphi_i(y_i)$라 놓자. 다음에 유의하라.
$$
X_{i, y_i} = \Spec(\kappa(y_i)) \times_{\Spec(\kappa(y))} X_y.
$$
따라서 이 상황에서 $X_y$가 기하학적으로 정규일 필요충분조건은
$X_{i, y_i}$가 기하학적으로 정규인 것이다(Varieties, Lemma
\ref{varieties-lemma-geometrically-normal-upstairs} 참조).
이 사실은 $\mathcal{P}_3$가 공역에서 fpqc 국소적임을 함의한다.

\medskip\noindent
$\{X_i \to X\}$를 $X$의 매끄러운 덮개라 하자.
$y \in Y$라 하자. 이 경우 $\{X_{i, y} \to X_y\}$는 올의
매끄러운 덮개이다. 따라서 정의역의 매끄러운 위상에 관한
$\mathcal{P}_3$의 국소성은 Descent, Lemma
\ref{descent-lemma-normal-local-smooth}에서 따른다.
위의 결과들을 결합하면 보조정리가 따른다.
\end{proof}



\section{정칙 사상}
\label{more-morphisms-section-regular}

\noindent
Section \ref{more-morphisms-section-normal}과 비교하라. 이 개념의 대수적 판본은
More on Algebra, Section \ref{more-algebra-section-regular}에서 논의한다.

\begin{definition}
\label{more-morphisms-definition-regular}
$f : X \to Y$를 스킴의 사상이라 하자.
모든 올 $X_y$가 국소 뇌터 스킴이라고 가정하자.
\begin{enumerate}
\item $x \in X$라 하고 $y = f(x)$라 하자. $f$가 $x$에서 평탄하고
스킴 $X_y$가 $\kappa(y)$ 위에서 $x$에서 기하학적으로 정칙이면
$f$가 {\it $x$에서 정칙이다}라고 한다(Varieties, Definition
\ref{varieties-definition-geometrically-regular} 참조).
\item $f$가 $X$의 모든 점에서 정칙이면 $f$를 {\it 정칙 사상}이라 한다.
\end{enumerate}
\end{definition}

\noindent
사상 $X \to Y$가 정칙이라는 조건은 모든 올이 정칙인
국소 뇌터 스킴이라고만 요구하는 것보다 강하다.

\begin{lemma}
\label{more-morphisms-lemma-regular}
$f : X \to Y$를 스킴의 사상이라 하자.
$f$의 모든 올이 국소 뇌터라고 가정하자.
다음은 동치이다.
\begin{enumerate}
\item $f$는 정칙이다.
\item $f$는 평탄하고 그 올들은 기하학적으로 정칙인 스킴이다.
\item $f(U) \subset V$를 만족하는 모든 아핀 열린집합의 쌍
$U \subset X$, $V \subset Y$에 대하여 환 준동형
$\mathcal{O}_Y(V) \to \mathcal{O}_X(U)$가 정칙이다.
\item 열린 덮개 $Y = \bigcup_{j \in J} V_j$와 열린 덮개들
$f^{-1}(V_j) = \bigcup_{i \in I_j} U_i$가 존재하여 각 사상
$U_i \to V_j$가 정칙이다.
\item 아핀 열린 덮개 $Y = \bigcup_{j \in J} V_j$와 아핀 열린 덮개들
$f^{-1}(V_j) = \bigcup_{i \in I_j} U_i$가 존재하여 환 준동형
$\mathcal{O}_Y(V_j) \to \mathcal{O}_X(U_i)$가 정칙이다.
\end{enumerate}
\end{lemma}

\begin{proof}
(1)과 (2)의 동치는 정의에서 바로 따른다.
$y = f(x)$인 $x \in X$를 잡자. 정의에 따라 $f$가 $x$에서 평탄일
필요충분조건은 $\mathcal{O}_{Y, y} \to \mathcal{O}_{X, x}$가
평탄한 환 준동형인 것이고, $X_y$가 $\kappa(y)$ 위에서 $x$에서
기하학적으로 정칙일 필요충분조건은
$\mathcal{O}_{X_y, x} = \mathcal{O}_{X, x}/\mathfrak m_y\mathcal{O}_{X, x}$
가 $\kappa(y)$ 위의 기하학적으로 정칙인 대수인 것이다. 따라서
% P05-EN-CH37-0106: frozen source line 5709 capitalizes “Whether” mid-sentence.
$f$가 $x$에서 정칙인지 여부는 국소환의 국소 준동형
$\mathcal{O}_{Y, y} \to \mathcal{O}_{X, x}$에만 의존한다.
그러므로 (1)과 (4)의 동치는 분명하다.

\medskip\noindent
환 준동형 $A \to B$가 정칙일 필요충분조건은 그것이 평탄하고,
모든 소 아이디얼 $\mathfrak p \subset A$에 대하여 올환
$B \otimes_A \kappa(\mathfrak p)$가 뇌터이면서 기하학적으로 정칙인
것임을 상기하라(More on Algebra, Definition
\ref{more-algebra-definition-regular}). Varieties, Lemma
\ref{varieties-lemma-geometrically-regular}에 따라 이는
$\Spec(B \otimes_A \kappa(\mathfrak p))$가 $\kappa(\mathfrak p)$ 위에서
기하학적으로 정칙인 스킴인 것과 동치이다. 따라서 (2)는
(3)을 함의한다. (3)이 (5)를 함의함은 분명하다.
끝으로 (5)를 가정하자. 그러면 $f$는 평탄하다(Morphisms, Lemma
\ref{morphisms-lemma-flat-characterize} 참조). 더욱이 $y \in Y$이면
어떤 $j$에 대하여 $y \in V_j$이고,
$X_y = \bigcup_{i \in I_j} U_{i, y}$이다. Varieties, Lemma
\ref{varieties-lemma-geometrically-regular}에 따라 각 $U_{i, y}$는
$\kappa(y)$ 위에서 기하학적으로 정칙이다. Varieties, Lemma
\ref{varieties-lemma-geometrically-regular}를 다시 적용하면
$X_y$가 기하학적으로 정칙임을 알 수 있다. 따라서 (2)가 성립하고
보조정리의 증명이 끝난다.
\end{proof}

\begin{lemma}
\label{more-morphisms-lemma-smooth-regular}
매끄러운 사상은 정칙이다.
\end{lemma}

\begin{proof}
$f : X \to Y$를 매끄러운 사상이라 하자.
Morphisms, Lemma \ref{morphisms-lemma-smooth-locally-finite-presentation}에
따라 $f$는 국소 유한 표시이므로 올 $X_y$는 체 위의 국소 유한형이고,
따라서 국소 뇌터이다. 또한 Morphisms, Lemma
\ref{morphisms-lemma-smooth-flat}에 따라 $f$는 평탄하다.
끝으로 올 $X_y$는 체 위에서 매끄럽고(Morphisms, Lemma
\ref{morphisms-lemma-base-change-smooth}), 따라서 Varieties, Lemma
\ref{varieties-lemma-smooth-geometrically-normal}에 의해 기하학적으로 정칙이다.
그러므로 Lemma \ref{more-morphisms-lemma-regular}에 따라 $f$는 정칙이다.
\end{proof}

\begin{lemma}
\label{more-morphisms-lemma-regular-fppf-local-source-and-target}
성질 $\mathcal{P}(f)=$``$f$의 올들은
국소 뇌터이고 $f$는 정칙이다''는
공역의 fppf 위상에서 국소적이고
정의역의 매끄러운 위상에서 국소적이다.
\end{lemma}

\begin{proof}
다음과 같이 쓸 수 있다.
$\mathcal{P}(f) =
\mathcal{P}_1(f) \wedge \mathcal{P}_2(f) \wedge \mathcal{P}_3(f)$
여기서
$\mathcal{P}_1(f)=$``$f$의 올들은 국소 뇌터이다'',
$\mathcal{P}_2(f)=$``$f$는 평탄하다'', 그리고
$\mathcal{P}_3(f)=$``$f$의 올들은 기하학적으로 정칙이다''이다.
Lemma \ref{more-morphisms-lemma-locally-Noetherian-fibres-fppf-local-source-and-target}와
Descent, Lemmas \ref{descent-lemma-descending-property-flat},
\ref{descent-lemma-flat-fpqc-local-source}에서 이미
$\mathcal{P}_1$과 $\mathcal{P}_2$가 정의역과 공역의 fppf 위상에서
국소적임을 보았다. 따라서 $\mathcal{P}_3$를 다루면 된다.

\medskip\noindent
$f : X \to Y$를 스킴의 사상이라 하자.
$\{\varphi_i : Y_i \to Y\}_{i \in I}$를 $Y$의 fpqc 덮개라 하자.
% P05-EN-CH37-0107: frozen source line 5781 omits “by/as” in “Denote ... the base change”.
$f_i : X_i \to Y_i$를 $\varphi_i$에 의한 $f$의 밑변환이라 쓰자.
$i \in I$이고 $y_i \in Y_i$가 한 점이라고 하자.
$y = \varphi_i(y_i)$라 놓자. 다음에 유의하라.
$$
X_{i, y_i} = \Spec(\kappa(y_i)) \times_{\Spec(\kappa(y))} X_y.
$$
따라서 이 상황에서 $X_y$가 기하학적으로 정칙일 필요충분조건은
$X_{i, y_i}$가 기하학적으로 정칙인 것이다(Varieties, Lemma
\ref{varieties-lemma-geometrically-regular-upstairs} 참조).
이 사실은 $\mathcal{P}_3$가 공역에서 fpqc 국소적임을 함의한다.

\medskip\noindent
$\{X_i \to X\}$를 $X$의 매끄러운 덮개라 하자.
$y \in Y$라 하자. 이 경우 $\{X_{i, y} \to X_y\}$는 올의
매끄러운 덮개이다. 따라서 정의역의 매끄러운 위상에 관한
$\mathcal{P}_3$의 국소성은 Descent, Lemma
\ref{descent-lemma-regular-local-smooth}에서 따른다.
위의 결과들을 결합하면 보조정리가 따른다.
\end{proof}




\section{코언--매콜리 사상}
\label{more-morphisms-section-CM}

\noindent
Section \ref{more-morphisms-section-normal}과 비교하라.
Algebra, Section \ref{algebra-section-geometrically-CM}과
Varieties, Section \ref{varieties-section-CM}에서 지적했듯이
``기하학적으로 코언--매콜리''인 것은 단순히 코언--매콜리인 것과 같다.

\begin{definition}
\label{more-morphisms-definition-CM}
$f : X \to Y$를 스킴의 사상이라 하자.
모든 올 $X_y$가 국소 뇌터 스킴이라고 가정하자.
\begin{enumerate}
\item $x \in X$라 하고 $y = f(x)$라 하자. $f$가 $x$에서 평탄하고
스킴 $X_y$의 $x$에서의 국소환이 코언--매콜리이면
$f$가 {\it $x$에서 코언--매콜리이다}라고 한다.
\item $f$가 $X$의 모든 점에서 코언--매콜리이면
$f$를 {\it 코언--매콜리 사상}이라 한다.
\end{enumerate}
\end{definition}

\noindent
이를 다음과 같이 다시 쓸 수 있다.

\begin{lemma}
\label{more-morphisms-lemma-CM}
$f : X \to Y$를 스킴의 사상이라 하자.
$f$의 모든 올이 국소 뇌터라고 가정하자.
다음은 동치이다.
\begin{enumerate}
\item $f$는 코언--매콜리이다.
\item $f$는 평탄하고 그 올들은 코언--매콜리 스킴이다.
\end{enumerate}
\end{lemma}

\begin{proof}
이는 정의에서 바로 따른다.
\end{proof}

\begin{lemma}
\label{more-morphisms-lemma-CM-dimension}
$f : X \to Y$를 국소 뇌터 스킴의 국소 유한형 코언--매콜리 사상이라 하자.
$Y$에서의 상이 $y$인 $X$의 모든 점 $x$에 대하여
$$
\dim_x(X) = \dim_y(Y) + \dim_x(X_y),
$$
이다. 여기서 $X_y$는 $y$ 위의 올을 나타낸다.
\end{lemma}

\begin{proof}
$X$를 $x$의 한 열린 근방으로 바꾸면 자연수 $d$가 존재하여
$X \to Y$의 모든 올이 각 점에서 차원 $d$를 갖는다(Morphisms, Lemma
\ref{morphisms-lemma-flat-finite-presentation-CM-fibres-relative-dimension} 참조).
그러면 $f$는 평탄하고 국소 유한형이며 상대 차원이 $d$이다.
따라서 결과는 Morphisms, Lemma
\ref{morphisms-lemma-rel-dimension-dimension}에서 따른다.
\end{proof}

\begin{lemma}
\label{more-morphisms-lemma-composition-CM}
$f : X \to Y$와 $g : Y \to Z$를 스킴의 사상이라 하자.
$f$, $g$, $g \circ f$의 올들이 국소 뇌터라고 가정하자.
각각 $y \in Y$, $z \in Z$로 가는 $x \in X$를 잡자.
\begin{enumerate}
\item $f$가 $x$에서 코언--매콜리이고 $g$가 $f(x)$에서
코언--매콜리이면 $g \circ f$는 $x$에서 코언--매콜리이다.
\item $f$와 $g$가 코언--매콜리이면 $g \circ f$는 코언--매콜리이다.
\item $g \circ f$가 $x$에서 코언--매콜리이고 $f$가 $x$에서 평탄이면,
$f$는 $x$에서 코언--매콜리이고 $g$는 $f(x)$에서 코언--매콜리이다.
\item $g \circ f$가 코언--매콜리이고 $f$가 평탄이면,
$f$는 코언--매콜리이고 $g$는 $f$의 상에 속하는 모든 점에서
코언--매콜리이다.
\end{enumerate}
\end{lemma}

\begin{proof}
다음 뇌터 국소환의 준동형을 생각하자.
$$
\mathcal{O}_{Y_z, y} \to \mathcal{O}_{X_z, x}
$$
그 올은 다음과 같다.
$$
\mathcal{O}_{X_z, x}/\mathfrak m_{Y_z, y}\mathcal{O}_{X_z, x} =
\mathcal{O}_{X_y, x}
$$
% P05-EN-CH37-0108: frozen source line 5895 contains the extraneous word “this”.
따라서 보조정리는 Algebra, Lemma \ref{algebra-lemma-CM-goes-up}에서 따른다.
\end{proof}

\begin{lemma}
\label{more-morphisms-lemma-flat-morphism-from-CM-scheme}
$f : X \to Y$를 국소 뇌터 스킴의 평탄 사상이라 하자.
$X$가 코언--매콜리이면 $f$는 코언--매콜리이고, 모든 $x \in X$에
대하여 $\mathcal{O}_{Y, f(x)}$가 코언--매콜리이다.
\end{lemma}

\begin{proof}
대수의 언어로 옮기면 이는 Algebra, Lemma
\ref{algebra-lemma-CM-goes-up}에서 따른다.
\end{proof}

\begin{lemma}
\label{more-morphisms-lemma-base-change-CM}
$f : X \to Y$를 스킴의 사상이라 하자.
모든 올 $X_y$가 국소 뇌터 스킴이라고 가정하자.
$Y' \to Y$가 국소 유한형이라고 하자. $f' : X' \to Y'$를
$f$의 밑변환이라 하자.
$x' \in X'$를 그 상이 $x \in X$인 점이라 하자.
\begin{enumerate}
\item $f$가 $x$에서 코언--매콜리이면
$f' : X' \to Y'$는 $x'$에서 코언--매콜리이다.
\item $f$가 $x$에서 평탄이고 $f'$가 $x'$에서 코언--매콜리이면,
$f$는 $x$에서 코언--매콜리이다.
\item $Y' \to Y$가 $f'(x')$에서 평탄이고 $f'$가 $x'$에서
코언--매콜리이면, $f$는 $x$에서 코언--매콜리이다.
\end{enumerate}
\end{lemma}

\begin{proof}
$Y' \to Y$에 대한 가정은 $y \in Y$로 가는 $y' \in Y'$에 대하여
체 확대 $\kappa(y')/\kappa(y)$가 유한 생성임을 함의한다.
따라서 모든 올 $X'_{y'} = (X_y)_{\kappa(y')}$도 국소 뇌터이다
(Varieties, Lemma \ref{varieties-lemma-locally-Noetherian-base-change} 참조).
그러므로 보조정리의 진술은 의미가 있다. $y' = f'(x')$와 $y = f(x)$라
놓자. 그러면 다음 국소환의 가환 그림을 얻는다.
$$
\xymatrix{
\mathcal{O}_{X', x'} & \mathcal{O}_{X, x} \ar[l] \\
\mathcal{O}_{Y', y'} \ar[u] & \mathcal{O}_{Y, y} \ar[l] \ar[u]
}
$$
여기서 왼쪽 위 꼭짓점은 오른쪽 위 꼭짓점과 왼쪽 아래 꼭짓점의
오른쪽 아래 꼭짓점 위 텐서곱을 국소화한 것이다.

\medskip\noindent
$f$가 $x$에서 코언--매콜리라고 가정하자.
$\mathcal{O}_{Y, y} \to \mathcal{O}_{X, x}$의 평탄성은
$\mathcal{O}_{Y', y'} \to \mathcal{O}_{X', x'}$의 평탄성을 함의한다
(Algebra, Lemma \ref{algebra-lemma-base-change-flat-up-down} 참조).
$\mathcal{O}_{X, x}/\mathfrak m_y\mathcal{O}_{X, x}$가 코언--매콜리라는
사실은 $\mathcal{O}_{X', x'}/\mathfrak m_{y'}\mathcal{O}_{X', x'}$가
코언--매콜리임을 함의한다(Varieties, Lemma
\ref{varieties-lemma-CM-base-change} 참조). 따라서 $f'$는
$x'$에서 코언--매콜리이다.

\medskip\noindent
$f$가 $x$에서 평탄이고 $f'$가 $x'$에서 코언--매콜리라고 가정하자.
$\mathcal{O}_{X', x'}/\mathfrak m_{y'}\mathcal{O}_{X', x'}$가
코언--매콜리라는 사실은
$\mathcal{O}_{X, x}/\mathfrak m_y\mathcal{O}_{X, x}$가
코언--매콜리임을 함의한다(Varieties, Lemma
\ref{varieties-lemma-CM-base-change} 참조).
따라서 $f$는 $x$에서 코언--매콜리이다.

\medskip\noindent
$Y' \to Y$가 $y'$에서 평탄이고 $f'$가 $x'$에서 코언--매콜리라고
가정하자. $\mathcal{O}_{Y', y'} \to \mathcal{O}_{X', x'}$와
$\mathcal{O}_{Y, y} \to \mathcal{O}_{Y', y'}$의 평탄성은
$\mathcal{O}_{Y, y} \to \mathcal{O}_{X, x}$의 평탄성을 함의한다
(Algebra, Lemma \ref{algebra-lemma-base-change-flat-up-down} 참조).
$\mathcal{O}_{X', x'}/\mathfrak m_{y'}\mathcal{O}_{X', x'}$가
코언--매콜리라는 사실은
$\mathcal{O}_{X, x}/\mathfrak m_y\mathcal{O}_{X, x}$가
코언--매콜리임을 함의한다(Varieties, Lemma
\ref{varieties-lemma-CM-base-change} 참조). 따라서 $f$는
$x$에서 코언--매콜리이다.
\end{proof}

\begin{lemma}
\label{more-morphisms-lemma-flat-finite-presentation-CM-open}
\begin{reference}
\cite[IV Corollary 12.1.7(iii)]{EGA}
\end{reference}
$f : X \to S$를 평탄한 국소 유한 표시 스킴 사상이라 하자. 다음과 놓자.
$$
W = \{x \in X \mid f\text{ is Cohen-Macaulay at }x\}
$$
그러면 다음이 성립한다.
\begin{enumerate}
\item $W = \{x \in X \mid \mathcal{O}_{X_{f(x)}, x}\text{ is Cohen-Macaulay}\}$이다.
\item $W$는 $X$에서 열린집합이다.
\item $W$는 $X \to S$의 모든 올에서 조밀하다.
\item $W$의 형성은 $f$의 임의의 밑변환과 가환한다.
임의의 사상 $g : S' \to S$에 대하여 $f$의 밑변환
$f' : X' \to S'$와 사영 $g' : X' \to X$를 생각하자. 그러면
% P05-EN-CH37-0109: frozen source line 5999 redundantly names W' on both sides.
사상 $f'$에 대응하는 집합 $W'$은 $W' = (g')^{-1}(W)$와 같다.
\end{enumerate}
\end{lemma}

\begin{proof}
$f$는 국소 뇌터 올을 갖는 평탄 사상이므로 (1)의 등식은 정의에 따라
성립한다. (2)와 (3)은 Algebra, Lemma
\ref{algebra-lemma-generic-CM-flat-finite-presentation}에서 따른다.
(4)는 Algebra, Lemma \ref{algebra-lemma-CM-locus-commutes-base-change}
또는 Varieties, Lemma \ref{varieties-lemma-CM-base-change}에서 따른다.
\end{proof}

\begin{lemma}
\label{more-morphisms-lemma-flat-finite-presentation-characterize-CM}
$f : X \to S$를 평탄한 국소 유한 표시 스킴 사상이라 하자.
$s \in S$로 가는 $x \in X$를 잡자.
$d = \dim_x(X_s)$라 놓자.
다음은 동치이다.
\begin{enumerate}
\item $f$는 $x$에서 코언--매콜리이다.
\item $x$의 열린 근방 $U \subset X$와 $S$ 위의 국소 유사유한 사상
$U \to \mathbf{A}^d_S$가 존재하여 이 사상이 $x$에서 평탄하다.
\item $x$의 열린 근방 $U \subset X$와 $S$ 위의 국소 유사유한 평탄 사상
$U \to \mathbf{A}^d_S$가 존재한다.
\item $x$의 열린 근방 $U \subset X$에서 정의된 임의의 $S$-사상
$g : U \to \mathbf{A}^d_S$에 대하여 다음이 성립한다.
$g$가 $x$에서 유사유한이다 $\Rightarrow$ $g$가 $x$에서 평탄이다.
\end{enumerate}
\end{lemma}

\begin{proof}
평탄성의 열림성에 따라 (2)와 (3)은 동치이다(Theorem
\ref{more-morphisms-theorem-openness-flatness} 참조).

\medskip\noindent
$x \in U$인 아핀 열린집합 $U = \Spec(A) \subset X$와
$f(U) \subset V$인 아핀 열린집합 $V = \Spec(R) \subset S$를 택하자.
그러면 $R \to A$는 유한 표시인 평탄 환 준동형이다.
$\mathfrak p \subset A$를 $x$에 대응하는 소 아이디얼이라 하자.
$A$를 한 주 열린 국소화로 바꾸면 유사유한 준동형
$R[x_1, \ldots, x_d]  \to A$가 존재한다고 가정할 수 있다
(Algebra, Lemma \ref{algebra-lemma-quasi-finite-over-polynomial-algebra} 참조).
따라서 $x$의 열린 근방 $U \subset X$와 국소적으로\footnote{$S$가
준분리이면 $g$는 유사유한이다.} 유사유한인 사상
$g : U \to \mathbf{A}^d_S$로 이루어진 쌍 $(U, g)$가 적어도 하나 존재한다.

\medskip\noindent
주장: $R \to A$가 평탄하고 유한 표시이며, $\mathfrak p \subset A$가
소 아이디얼이고, $\varphi : R[x_1, \ldots, x_d] \to A$가
$\mathfrak p$에서 유사유한이라고 하자. 그러면
$\Spec(\varphi)$가 $\mathfrak p$에서 평탄일 필요충분조건은
$\Spec(A) \to \Spec(R)$가 $\mathfrak p$에서 코언--매콜리인 것이다.
실제로 Theorem \ref{more-morphisms-theorem-criterion-flatness-fibre}에 따라
평탄성은 올에서 확인할 수 있다. 코언--매콜리라는 성질도 마찬가지이다
($A$가 이미 $R$ 위에서 평탄이라고 가정했기 때문이다).
따라서 주장은 Algebra, Lemma \ref{algebra-lemma-where-CM}에서 따른다.

\medskip\noindent
이 주장은 (1)과 (4)가 동치임을 보인다. 또한 둘째 문단에서 적절한
$(U, g)$를 구성했다는 사실과 결합하면, 이 주장은 (1)과 (2)가
동치임도 보인다.
\end{proof}

\begin{lemma}
\label{more-morphisms-lemma-flat-finite-presentation-CM-pieces}
$f : X \to S$를 평탄한 국소 유한 표시 스킴 사상이라 하자.
$d \geq 0$에 대하여 다음 성질을 갖는 열린집합 $U_d \subset X$가 존재한다.
\begin{enumerate}
\item $W = \bigcup_{d \geq 0} U_d$는 $f$의 모든 올에서 조밀하다.
\item $U_d \to S$의 상대 차원은 $d$이다(Morphisms, Definition
\ref{morphisms-definition-relative-dimension-d} 참조).
\end{enumerate}
\end{lemma}

\begin{proof}
이는 Lemma \ref{more-morphisms-lemma-flat-finite-presentation-CM-open}과
Morphisms, Lemma
\ref{morphisms-lemma-flat-finite-presentation-CM-fibres-relative-dimension}을
결합하면 따른다.
\end{proof}

\begin{lemma}
\label{more-morphisms-lemma-flat-finite-presentation-specialization-dimension}
$f : X \to S$를 평탄한 국소 유한 표시 스킴 사상이라 하자.
$x' \leadsto x$가 $X$의 점들의 특수화이고 $S$에서의 상이
$s' \leadsto s$라고 하자. $x$가 $X_s$의 한 기약 성분의 일반점이면
$\dim_{x'}(X_{s'}) = \dim_x(X_s)$이다.
\end{lemma}

\begin{proof}
% P05-EN-CH37-0110: frozen proof lines 6096-6099 stop after locating x in U_d.
점 $x$는 어떤 $d$에 대하여 $U_d$에 속한다. 여기서 $U_d$는
Lemma \ref{more-morphisms-lemma-flat-finite-presentation-CM-pieces}에서와 같다.
\end{proof}

\begin{lemma}
\label{more-morphisms-lemma-CM-local-source-and-target}
성질
$\mathcal{P}(f)=$``$f$의 올들은 국소 뇌터이고 $f$는
코언--매콜리이다''는 공역의 fppf 위상에서 국소적이고
정의역의 신토믹 위상에서 국소적이다.
\end{lemma}

\begin{proof}
다음과 같이 쓸 수 있다.
$\mathcal{P}(f) =
\mathcal{P}_1(f) \wedge \mathcal{P}_2(f)$
여기서
$\mathcal{P}_1(f)=$``$f$는 평탄하다'', 그리고
$\mathcal{P}_2(f)=$``$f$의 올들은 국소 뇌터이고
코언--매콜리이다''이다.
Descent, Lemmas \ref{descent-lemma-descending-property-flat},
\ref{descent-lemma-flat-fpqc-local-source}에 따라
$\mathcal{P}_1$은 정의역과 공역의 fppf 위상에서 국소적이다.
따라서 $\mathcal{P}_2$를 다루면 된다.

\medskip\noindent
$f : X \to Y$를 스킴의 사상이라 하자.
$\{\varphi_i : Y_i \to Y\}_{i \in I}$를 $Y$의 fppf 덮개라 하자.
% P05-EN-CH37-0111: frozen source line 6126 omits “by/as” in “Denote ... the base change”.
$f_i : X_i \to Y_i$를 $\varphi_i$에 의한 $f$의 밑변환이라 쓰자.
$i \in I$이고 $y_i \in Y_i$가 한 점이라고 하자.
$y = \varphi_i(y_i)$라 놓자. 다음에 유의하라.
$$
X_{i, y_i} = \Spec(\kappa(y_i)) \times_{\Spec(\kappa(y))} X_y.
$$
% P05-EN-CH37-0112: frozen source line 6132 continues the sentence after a period.
또한 $\kappa(y_i)/\kappa(y)$는 유한 생성 체 확대이다.
따라서 $X_y$가 국소 뇌터이면 $X_{i, y_i}$도 국소 뇌터이다
(Varieties, Lemma \ref{varieties-lemma-locally-Noetherian-base-change} 참조).
또한 $X_y$가 코언--매콜리이기까지 하면 $X_{i, y_i}$도
코언--매콜리이다(Varieties, Lemma \ref{varieties-lemma-CM-base-change} 참조).
따라서 $\mathcal{P}_2$는 공역에서 fppf 국소적이다.

\medskip\noindent
$\{X_i \to X\}$를 $X$의 신토믹 덮개라 하자.
$y \in Y$라 하자. 이 경우 $\{X_{i, y} \to X_y\}$는 올의
신토믹 덮개이다. 따라서 정의역의 신토믹 위상에 관한
$\mathcal{P}_2$의 국소성은 Descent, Lemma
\ref{descent-lemma-CM-local-syntomic}에서 따른다.
위의 결과들을 결합하면 보조정리가 따른다.
\end{proof}






\section{코언--매콜리 사상의 절단}
\label{more-morphisms-section-slicing-CM}

\noindent
이 절의 결과들은 결국 fppf 위상이
``유한 표시, 평탄, 유사유한'' 위상과 같다는 명제로 이어진다.
다음 보조정리는 Divisors, Lemma
\ref{divisors-lemma-fibre-Cartier}와 매우 밀접하다.

\begin{lemma}
\label{more-morphisms-lemma-slice-given-element}
$f : X \to S$를 스킴의 사상이라 하자.
$x \in X$를 그 상이 $s \in S$인 점이라 하자.
$h \in \mathfrak m_x \subset \mathcal{O}_{X, x}$라 하자.
다음을 가정하자.
\begin{enumerate}
\item $f$는 국소 유한 표시이다.
\item $f$는 $x$에서 평탄이다.
\item $h$의 다음 환에서의 상 $\overline{h}$는 영인자가 아닌 원소이다.
$\mathcal{O}_{X_s, x} = \mathcal{O}_{X, x}/\mathfrak m_s\mathcal{O}_{X, x}$
\end{enumerate}
그러면 $x$의 아핀 열린 근방 $U \subset X$가 존재하여,
$h$는 $h \in \Gamma(U, \mathcal{O}_U)$에서 오고,
$D = V(h)$는 $U$의 유효 Cartier 인자이며 $x \in D$이고,
$D \to S$는 평탄한 국소 유한 표시 사상이다.
\end{lemma}

\begin{proof}
뇌터인 경우로 환원하여 이를 증명한다.
평탄성의 열림성(Theorem \ref{more-morphisms-theorem-openness-flatness} 참조)에 따라
$X$를 $x$의 한 열린 근방으로 바꾸어 $X \to S$가 평탄이라고
가정할 수 있다. 또한 $X$와 $S$가 아핀이라고 가정할 수 있다.
$X$를 조금 더 줄여, 주어진 $h$로 가는
$h \in \Gamma(X, \mathcal{O}_X)$가 존재한다고 가정할 수 있다.

\medskip\noindent
$S = \Spec(A)$라 쓰고, $A = \colim_i A_i$를 유한형
$\mathbf{Z}$-대수들의 유향 쌍극한으로 쓸 수 있다.
% P05-EN-CH37-0113: frozen source lines 6200-6202 repeat the same Limits citation.
그러면 Algebra, Lemma \ref{algebra-lemma-flat-finite-presentation-limit-flat}
또는 Limits, Lemmas \ref{limits-lemma-descend-finite-presentation},
\ref{limits-lemma-descend-affine-finite-presentation},
\ref{limits-lemma-descend-finite-presentation}에 따라 다음 데카르트 그림을
찾을 수 있다.
$$
\xymatrix{
X \ar[r] \ar[d]_f & X_0 \ar[d]^{f_0} \\
S \ar[r] & S_0
}
$$
여기서 $f_0$는 평탄하고 유한 표시이며, $X_0$는 아핀이고,
$S_0$는 아핀 뇌터이다. $x_0 \in X_0$, 각각 $s_0 \in S_0$를
$x$, 각각 $s$의 상이라 하자. 또한
$h_0 \in \Gamma(X_0, \mathcal{O}_{X_0})$가 존재하여 $X$에서 $h$로
제한된다고 가정할 수 있다(위의 대수 참고문헌을 사용했다면 이는
분명하고, 극한 장의 참고문헌을 사용했다면 $h$를 사상
$X \to \mathbf{A}^1_S$로 생각하여 Limits, Lemma
\ref{limits-lemma-descend-finite-presentation}를 적용하면 된다).
$\mathcal{O}_{X_s, x}$는
$\mathcal{O}_{(X_0)_{s_0}, x_0} \otimes_{\kappa(s_0)} \kappa(s)$의
국소화이므로
$\mathcal{O}_{(X_0)_{s_0}, x_0} \to \mathcal{O}_{X_s, x}$는
평탄한 국소환 준동형이고, 특히 충실하게 평탄이다. 따라서 상
$\overline{h}_0 \in \mathcal{O}_{(X_0)_{s_0}, x_0}$는
$\mathfrak m_{(X_0)_{s_0}, x_0}$에 속하고 영인자가 아닌 원소이다.
$X_0$를 $x_0$의 한 주 열린 근방으로 바꾸면 원소 $h_0$가
$B_0 = \Gamma(X_0, \mathcal{O}_{X_0})$에서 영인자가 아니고
$B_0/h_0B_0$가 $A_0 = \Gamma(S_0, \mathcal{O}_{S_0})$ 위에서
평탄이라고 주장한다. 그렇다면
$$
0 \to B_0 \xrightarrow{h_0} B_0 \to B_0/h_0B_0 \to 0
$$
은 평탄한 $A_0$-가군의 짧은 완전열이다. 따라서 $A$와 텐서곱한
뒤에도 완전하며(Algebra, Lemma \ref{algebra-lemma-flat-tor-zero}),
보조정리가 따른다.

\medskip\noindent
위 주장을 증명하면 된다. 이에 대응하는 대수 명제는 다음과 같다
(여기서는 아래첨자 ${}_0$을 생략한다).
$A \to B$를 뇌터 환들의 평탄한 유한형 환 준동형이라 하자.
$\mathfrak q \subset B$를 $\mathfrak p \subset A$ 위에 놓인 소 아이디얼이라 하자.
$h \in \mathfrak q$의 $B_{\mathfrak q}/\mathfrak p B_{\mathfrak q}$에서의
상이 영인자가 아닌 원소라고 가정하자.
% P05-EN-CH37-0114: frozen source line 6244 says “after possible replacing”.
목표는 어떤 $g \in B$, $g \not \in \mathfrak q$에 대하여 $B$를 $B_g$로
바꾸면 $h$가 영인자가 아닌 원소가 되고 $B/hB$가 $A$ 위에서
평탄임을 보이는 것이다. Algebra, Lemma
\ref{algebra-lemma-grothendieck}에 따라 $h$는 $B_{\mathfrak q}$에서
영인자가 아닌 원소이고 $B_{\mathfrak q}/hB_{\mathfrak q}$는
$A$ 위에서 평탄이다. 평탄성의 열림성(Algebra, Theorem
\ref{algebra-theorem-openness-flatness} 또는 Theorem
\ref{more-morphisms-theorem-openness-flatness} 참조)에 따라 어떤 $g \in B$,
$g \not \in \mathfrak q$에 대하여 $B$를 $B_g$로 바꾸면
$B/hB$가 $A$ 위에서 평탄이다. 끝으로
$I = \{b \in B \mid hb = 0\}$를 $h$의 소멸자라 하자.
$h$는 $B_{\mathfrak q}$에서 영인자가 아닌 원소이므로
$IB_{\mathfrak q} = 0$이다. 또한 $B$는 뇌터이므로 $I$는 유한 생성이다.
따라서 $IB_g = 0$을 만족하는 $g \in B$, $g \not \in \mathfrak q$가 존재한다.
$B$를 $B_g$로 바꾸면 $h$가 영인자가 아닌 원소임을 알 수 있다.
\end{proof}

\begin{lemma}
\label{more-morphisms-lemma-slice-given-elements}
$f : X \to S$를 스킴의 사상이라 하자.
$x \in X$를 그 상이 $s \in S$인 점이라 하자.
$h_1, \ldots, h_r \in \mathcal{O}_{X, x}$라 하자.
다음을 가정하자.
\begin{enumerate}
\item $f$는 국소 유한 표시이다.
\item $f$는 $x$에서 평탄이다.
\item $h_1, \ldots, h_r$의 다음 환에서의 상들은 정칙열을 이룬다.
$\mathcal{O}_{X_s, x} = \mathcal{O}_{X, x}/\mathfrak m_s\mathcal{O}_{X, x}$
\end{enumerate}
그러면 $x$의 아핀 열린 근방 $U \subset X$가 존재하여,
$h_1, \ldots, h_r$는
$h_1, \ldots, h_r \in \Gamma(U, \mathcal{O}_U)$에서 오고,
$Z = V(h_1, \ldots, h_r) \to U$는 $x \in Z$를 만족하는 정칙 몰입이며,
$Z \to S$는 평탄한 국소 유한 표시 사상이다. 더욱이 $S$ 위의
임의의 스킴 $S'$에 대하여 밑변환 $Z_{S'} \to U_{S'}$는 정칙 몰입이다.
\end{lemma}

\begin{proof}
(정칙열에 대한 우리의 규약은 각 $i$에 대하여
$h_i \in \mathfrak m_x$임을 함의한다.) $r = 1$인 경우는
Lemma \ref{more-morphisms-lemma-slice-given-element}과 Divisors, Lemma
\ref{divisors-lemma-relative-Cartier}를 결합하여 $V(h_1)$가
밑변환 뒤에도 유효 Cartier 인자로 남음을 확인하면 따른다.
$r > 1$인 경우는 $r$에 대한 직접적인 귀납법으로 따른다
($r = 1$인 결과를 정확히 $r$번 적용하며, 자세한 내용은 생략한다).

\medskip\noindent
보조정리를 증명하는 다른 방법은 Divisors, Section
\ref{divisors-section-relative-regular-immersion}의 내용을 사용하는 것이다.
먼저 평탄성의 열림성(Theorem \ref{more-morphisms-theorem-openness-flatness} 참조)에 따라
$X$를 $x$의 한 열린 근방으로 바꾸어 $X \to S$가 평탄이라고
가정할 수 있다. 또한 $X$와 $S$가 아핀이라고 가정할 수 있다.
$X$를 조금 줄여 $h_1, \ldots, h_r \in \Gamma(X, \mathcal{O}_X)$라고
가정할 수 있다. $Z = V(h_1, \ldots, h_r)$라 놓자.
$X_s$는 대수적 $\kappa(s)$-스킴이므로 뇌터 스킴이고
(Varieties, Section \ref{varieties-section-algebraic-schemes} 참조),
$X_s$의 위상은 $X$의 위상에서 유도된다(Schemes, Lemma
\ref{schemes-lemma-fibre-topological} 참조). 따라서 $X$를 조금 더
줄여 $Z_s \subset X_s$가 $r$개 원소 $h_i|_{X_s}$로 정의되는
정칙 몰입이라고 가정할 수 있다(Divisors, Lemma
\ref{divisors-lemma-Noetherian-scheme-regular-ideal}과 그 증명 참조).
또한 다음 때문에 $r = \dim_x(X_s) - \dim_x(Z_s)$임이 분명하다.
\begin{align*}
\dim_x(X_s) & = \dim(\mathcal{O}_{X_s, x}) +
\text{trdeg}_{\kappa(s)}(\kappa(x)), \\
\dim_x(Z_s) & = \dim(\mathcal{O}_{Z_s, x}) +
\text{trdeg}_{\kappa(s)}(\kappa(x)), \\
\dim(\mathcal{O}_{X_s, x}) & = \dim(\mathcal{O}_{Z_s, x}) + r
\end{align*}
처음 두 등식은 Algebra, Lemma
\ref{algebra-lemma-dimension-at-a-point-finite-type-field}에서 따르고,
% P05-EN-CH37-0115: frozen source line 6325 calls the remaining third equality “the second”.
둘째 등식은 Algebra, Lemma \ref{algebra-lemma-one-equation}을
$r$번 적용하면 따른다. 따라서 Divisors, Lemma
\ref{divisors-lemma-fibre-quasi-regular}의 (3)을 적용하면
(Zariski 의미에서 $X$를 줄인 뒤) 사상 $Z \to X$가 정칙 몰입이고,
Divisors, Lemma
\ref{divisors-lemma-relative-regular-immersion-flat-in-neighbourhood}을
적용할 수 있음을 알 수 있다(이 보조정리가 평탄성과 밑변환에 관한
진술을 준다).
\end{proof}

\begin{lemma}
\label{more-morphisms-lemma-slice-once}
$f : X \to S$를 스킴의 사상이라 하자.
$x \in X$를 그 상이 $s \in S$인 점이라 하자.
다음을 가정하자.
\begin{enumerate}
\item $f$는 국소 유한 표시이다.
\item $f$는 $x$에서 평탄이다.
\item $\mathcal{O}_{X_s, x}$의 $\text{depth} \geq 1$이다.
\end{enumerate}
그러면 $x$의 아핀 열린 근방 $U \subset X$와 $x$를 포함하는
유효 Cartier 인자 $D \subset U$가 존재하여
$D \to S$는 평탄하고 유한 표시이다.
\end{lemma}

\begin{proof}
$\mathcal{O}_{X_s, x}$에서 영인자가 아닌 원소로 가는 임의의
$h \in \mathfrak m_x \subset \mathcal{O}_{X, x}$를 택하고
Lemma \ref{more-morphisms-lemma-slice-given-element}을 적용하라.
\end{proof}

\begin{lemma}
\label{more-morphisms-lemma-slice-CM}
\begin{reference}
\cite[IV Proposition 17.16.1]{EGA}
\end{reference}
$f : X \to S$를 스킴의 사상이라 하자.
$x \in X$를 그 상이 $s \in S$인 점이라 하자.
다음을 가정하자.
\begin{enumerate}
\item $f$는 국소 유한 표시이다.
\item $f$는 $x$에서 코언--매콜리이다.
\item $x$는 $X_s$의 닫힌점이다.
\end{enumerate}
그러면 $x$를 포함하는 정칙 몰입 $Z \to X$가 존재하여 다음이 성립한다.
\begin{enumerate}
\item[(a)] $Z \to S$는 평탄하고 국소 유한 표시이다.
\item[(b)] $Z \to S$는 국소 유사유한이다.
\item[(c)] 집합론적으로 $Z_s = \{x\}$이다.
\end{enumerate}
\end{lemma}

\begin{proof}
$S$를 $s$의 한 아핀 열린 근방으로 바꾸자.
아핀 $S$에 대하여 $d = \dim_x(X_s)$에 대한 귀납법으로 보조정리를 증명한다.

\medskip\noindent
$d = 0$인 경우. 이 경우 $Z$를 $x$의 한 열린 근방으로 택할 수
있음을 보인다(열린 몰입은 정칙 몰입임에 유의하라).
실제로 $d = 0$이면 $X \to S$는 $x$에서 유사유한이다
(Morphisms, Lemma
\ref{morphisms-lemma-locally-quasi-finite-rel-dimension-0} 참조).
따라서 $U \to S$가 유사유한인 아핀 열린 근방 $U \subset X$가 존재한다
(Morphisms, Lemma \ref{morphisms-lemma-quasi-finite-points-open} 참조).
그러므로 $X$를 $U$로 바꾸면 올 $X_s$는 유한 이산 집합이다.
% P05-EN-CH37-0116: frozen source line 6393 says an affine open neighbourhood of X, not x.
따라서 $X$를 $X$의 더 작은 아핀 열린 근방으로 바꾸면
$f^{-1}(\{s\}) = \{x\}$임을 알 수 있다($X_s$의 위상은 $X$의
위상에서 유도되기 때문이다. Schemes, Lemma
\ref{schemes-lemma-fibre-topological} 참조).
이는 이 경우의 보조정리를 증명한다.

\medskip\noindent
이제 $d > 0$이라 가정하자. $x$는 그 올의 닫힌점이므로
체 확대 $\kappa(x)/\kappa(s)$는 유한이다(Hilbert 영점정리에 의하며,
Morphisms, Lemma
\ref{morphisms-lemma-closed-point-fibre-locally-finite-type} 참조).
따라서 다음을 얻는다.
$$
\text{depth}(\mathcal{O}_{X_s, x}) = \dim(\mathcal{O}_{X_s, x}) = d > 0
$$
첫째 등식은 $\mathcal{O}_{X_s, x}$가 코언--매콜리이므로 성립하고,
둘째 등식은 Morphisms, Lemma
\ref{morphisms-lemma-dimension-fibre-at-a-point}에서 따른다.
따라서 Lemma \ref{more-morphisms-lemma-slice-once}을 적용하여 다음 그림을 얻는다.
$$
\xymatrix{
D \ar[r] \ar[rrd] & U \ar[r] \ar[rd] & X \ar[d] \\
& & S
}
$$
여기서 $x \in D$이다. 어떤 영인자가 아닌 원소 $\overline{h}$에 대하여
$\mathcal{O}_{D_s, x} = \mathcal{O}_{X_s, x}/(\overline{h})$임에 유의하라
(Divisors, Lemma \ref{divisors-lemma-relative-Cartier} 참조).
따라서 $\mathcal{O}_{D_s, x}$는 코언--매콜리이고 그 차원은
$\mathcal{O}_{X_s, x}$의 차원보다 하나 작다(예를 들어 Algebra, Lemma
\ref{algebra-lemma-reformulate-CM} 참조). 그러므로 사상
$D \to S$는 평탄하고 국소 유한 표시이며 $x$에서 코언--매콜리이고,
$\dim_x(D_s) = \dim_x(X_s) - 1 = d - 1$이다. 귀납 가정에 따라
(a), (b), (c)의 성질을 갖는 정칙 몰입 $Z \to D$를 찾을 수 있다.
$Z \to D \to U$는 둘 다 정칙 몰입이므로 Divisors, Lemma
\ref{divisors-lemma-composition-regular-immersion}에 따라
$Z \to U$도 정칙 몰입이다. 이로써 증명이 끝난다.
\end{proof}

\begin{lemma}
\label{more-morphisms-lemma-qf-fp-flat-neighbourhood-dominates-fppf}
$f : X \to S$를 국소 유한 표시인 평탄 스킴 사상이라 하자.
$s \in S$를 $f$의 상에 속하는 점이라 하자.
그러면 다음 가환 그림이 존재한다.
$$
\xymatrix{
S' \ar[rr] \ar[rd]_g & & X \ar[ld]^f \\
& S
}
$$
여기서 $g : S' \to S$는 평탄하고 국소 유한 표시이며 국소 유사유한이고,
$s \in g(S')$이다.
\end{lemma}

\begin{proof}
가정에 따라 올 $X_s$는 공집합이 아니다. 따라서 $f$가
코언--매콜리인 닫힌점 $x \in X_s$가 존재한다(Lemma
\ref{more-morphisms-lemma-flat-finite-presentation-CM-open} 참조).
Lemma \ref{more-morphisms-lemma-slice-CM}을 적용하고 $S' = Z$라 놓자.
\end{proof}

\noindent
다음 보조정리는 fppf 위상의 층이
``유사유한, 평탄, 유한 표시'' 위상의 층과 같은 것임을 보인다.

\begin{lemma}
\label{more-morphisms-lemma-qf-fp-flat-dominates-fppf}
$S$를 스킴이라 하자. $\mathcal{U} = \{S_i \to S\}_{i \in I}$를
$S$의 fppf 덮개라 하자(Topologies, Definition
\ref{topologies-definition-fppf-covering} 참조).
그러면 $\mathcal{U}$를 세분하고(Sites, Definition
\ref{sites-definition-morphism-coverings} 참조) 각 $T_j \to S$가
국소 유사유한인 fppf 덮개
$\mathcal{V} = \{T_j \to S\}_{j \in J}$가 존재한다.
\end{lemma}

\begin{proof}
모든 $s \in S$에 대하여 $s$가 $S_i \to S$의 상에 속하는
$i \in I$가 존재한다. Lemma
\ref{more-morphisms-lemma-qf-fp-flat-neighbourhood-dominates-fppf}에 따라
$s \in g_s(T_s)$를 만족하며 평탄하고 국소 유한 표시이며 국소 유사유한이고
$g_s$가 $S_i \to S$를 통해 분해되는 사상 $g_s : T_s \to S$를 찾을 수 있다.
따라서 $\{T_s \to S\}$는 $\mathcal{U}$를 세분하는, 구하고자 한
$S$의 덮개이다.
\end{proof}



\section{일반올}
\label{more-morphisms-section-generic}

\noindent
일반올과 근방의 올 사이의 관계에 관한 몇 가지 결과를 제시한다.

\begin{lemma}
\label{more-morphisms-lemma-empty-generic-fibre}
$f : X \to Y$를 스킴의 유한형 사상이라 하자. $Y$가 일반점
$\eta$를 가진 기약 스킴이라고 가정하자. $X_\eta = \emptyset$이면
$X_V = V \times_Y X = \emptyset$을 만족하는 비어 있지 않은
열린집합 $V \subset Y$가 존재한다.
\end{lemma}

\begin{proof}
더 일반적인 Morphisms,
Lemma \ref{morphisms-lemma-quasi-compact-generic-point-not-in-image}에서
즉시 따른다.
\end{proof}

\begin{lemma}
\label{more-morphisms-lemma-nonempty-generic-fibre}
$f : X \to Y$를 스킴의 유한형 사상이라 하자. $Y$가 일반점
$\eta$를 가진 기약 스킴이라고 가정하자. $X_\eta \not = \emptyset$이면
$X_V = V \times_Y X \to V$가 전사인 비어 있지 않은 열린집합
$V \subset Y$가 존재한다.
\end{lemma}

\begin{proof}
아핀 열린집합들을 택하면 Algebra,
Lemma \ref{algebra-lemma-characterize-image-finite-type}에서 따른다.
(물론 일반 평탄성에서도 따른다.)
\end{proof}

\begin{lemma}
\label{more-morphisms-lemma-nowhere-dense-generic-fibre}
$f : X \to Y$를 스킴의 유한형 사상이라 하자. $Y$가 일반점
$\eta$를 가진 기약 스킴이라고 가정하자.
$Z \subset X$가 닫힌 부분집합이고 $Z_\eta$가 $X_\eta$에서
어디에서도 조밀하지 않다면, 모든 $y \in V$에 대하여 $Z_y$가
$X_y$에서 어디에서도 조밀하지 않게 되는 비어 있지 않은 열린집합
$V \subset Y$가 존재한다.
\end{lemma}

\begin{proof}
$Y' \subset Y$를 $Y$의 축약이라 하자.
$X' = Y' \times_Y X$ 및 $Z' = Y' \times_Y Z$로 놓자.
Morphisms, Lemma
\ref{morphisms-lemma-reduction-universal-homeomorphism}에 의해
$Y' \to Y$가 보편 위상동형이므로 $Z' \subset X' \to Y'$에 대하여
보조정리를 증명하면 충분하다.
따라서 $Y$가 정수적이라고 가정할 수 있다. Properties, Lemma
\ref{properties-lemma-characterize-integral}을 참조하라.
Morphisms, Proposition
\ref{morphisms-proposition-generic-flatness}에 의해
$X_V \to V$와 $Z_V \to V$가 평탄하고 유한 표시인 비어 있지 않은
아핀 열린집합 $V \subset Y$가 존재한다. 이 $V$가 조건을 만족한다고
주장한다.
$y \in V$를 택하자. $Z_y$의 내부가 비어 있지 않다면, $Z_y$는
$X_y$의 한 기약 성분의 일반점 $\xi$를 포함한다.
$\eta \leadsto f(\xi)$임에 유의하자. $Z_V \to V$가 평탄하므로
$f(\xi') = \eta$를 만족하는 특수화 $\xi' \leadsto \xi$,
$\xi' \in Z$를 택할 수 있다. Morphisms, Lemma
\ref{morphisms-lemma-generalizations-lift-flat}을 참조하라.
Lemma \ref{more-morphisms-lemma-flat-finite-presentation-specialization-dimension}에
의해
$$
\dim_{\xi'}(Z_\eta) = \dim_{\xi}(Z_y) = \dim_{\xi}(X_y) = \dim_{\xi'}(X_\eta).
$$
따라서 $\xi'$을 지나는 $Z_\eta$의 어떤 기약 성분은 차원이
$\dim_{\xi'}(X_\eta)$이다. 이는 $Z_\eta$가 $X_\eta$에서 어디에서도
조밀하지 않다는 가정에 모순이므로 증명이 끝난다.
\end{proof}

\begin{lemma}
\label{more-morphisms-lemma-scheme-theoretically-dense-generic-fibre}
$f : X \to Y$를 스킴의 유한형 사상이라 하자.
$Y$가 일반점 $\eta$를 가진 기약 스킴이라고 가정하자.
$U_\eta$가 $X_\eta$에서 스킴론적으로 조밀하도록 $U \subset X$를
열린 부분스킴이라 하자. 그러면 모든 $y \in V$에 대하여 $U_y$가
$X_y$에서 스킴론적으로 조밀하게 되는 비어 있지 않은 열린집합
$V \subset Y$가 존재한다.
\end{lemma}

\begin{proof}
$Y' \subset Y$를 $Y$의 축약이라 하자.
$X' = Y' \times_Y X$ 및 $U' = Y' \times_Y U$로 놓자.
$Y' \to Y$가 점들의 전단사를 유도하고 $U' \to U$와 $X' \to X$가
스킴론적 올들의 동형을 유도하므로, $Y$를 $Y'$으로, $X$를 $X'$으로
바꾸어도 된다. 따라서 $Y$가 정수적이라고 가정할 수 있다. Properties,
Lemma \ref{properties-lemma-characterize-integral}을 참조하라.
또한 $Y$를 비어 있지 않은 아핀 열린집합으로 바꾸어도 된다. 다시 말해
$A$가 분수체 $K$를 가진 정역이고 $Y = \Spec(A)$라고 가정할 수 있다.

\medskip\noindent
$f$가 유한형이므로 $X$는 준콤팩트이다.
어떤 아핀 열린집합 $X_i$들에 대하여
$X = X_1 \cup \ldots \cup X_n$이라 쓰자. Morphisms, Definition
\ref{morphisms-definition-scheme-theoretically-dense}에 의해
$U_i = X_i \cap U$는 $X_i$의 열린 부분스킴이고 $U_{i, \eta}$는
$X_{i, \eta}$에서 스킴론적으로 조밀하다.
따라서 각 쌍 $(X_i, U_i)$에 대하여 결과를 증명하면 충분하다.
다시 말해 $X$가 아핀이라고 가정할 수 있다.

\medskip\noindent
$X = \Spec(B)$라 쓰자. $B_K$는 유한형 $K$-대수이므로 뇌터임에
유의하자. 따라서 $U_\eta$는 준콤팩트이다. 그러므로
$D(g_j) \subset U$이고
$U_\eta = D(g_1)_\eta \cup \ldots \cup D(g_m)_\eta$가 되도록 하는
유한 개의 $g_1, \ldots, g_m \in B$를 찾을 수 있다.
$U_\eta$가 $X_\eta$에서 스킴론적으로 조밀하다는 것은
$B_K \to \bigoplus_j (B_K)_{g_j}$가 단사라는 뜻이다. Morphisms,
Example \ref{morphisms-example-scheme-theoretic-closure}를 참조하라.
Algebra, Lemma \ref{algebra-lemma-when-injective-covering}에 의해
이는 다음 사상의 단사성과 동치이다.
$B_K \to \bigoplus\nolimits_{j = 1, \ldots, m} B_K$,
$b \mapsto (g_1b, \ldots, g_mb)$. 이 사상을 $A$ 위에서 취했을 때의
여핵을 $M$이라 하자. 즉, 다음 완전열이 있다.
$$
0 \to I \to B \xrightarrow{(g_1, \ldots, g_m)}
\bigoplus\nolimits_{j = 1, \ldots, m} B \to M \to 0
$$
어떤 영이 아닌 $h$에 대하여 $A$를 $A_h$로 바꾸면, $B$가 평탄한
유한 표시 $A$-대수이고 $M$이 $A$ 위에서 평탄이라고 가정할 수 있다.
Algebra, Lemma \ref{algebra-lemma-generic-flatness}를 참조하라.
$B$가 $A$ 위에서 평탄하므로 $B$는 $A$-가군으로서 꼬임이 없다.
More on Algebra, Lemma \ref{more-algebra-lemma-flat-torsion-free}를
참조하라. 따라서 $B \subset B_K$이다. 가정에 의해 $I_K = 0$이고,
이는 $I = 0$을 함의한다($I \subset B \subset B_K$가 $I_K$의
부분집합이기 때문이다). 그러므로 이제 다음 짧은 완전열이 있다.
$$
0 \to B \xrightarrow{(g_1, \ldots, g_m)}
\bigoplus\nolimits_{j = 1, \ldots, m} B \to M \to 0
$$
여기서 $M$은 $A$ 위에서 평탄하다. 따라서 $\kappa$가 체일 때의 모든
준동형 $A \to \kappa$에 대하여 다음 짧은 완전열을 얻는다.
$$
0 \to B \otimes_A \kappa \xrightarrow{(g_1 \otimes 1, \ldots, g_m \otimes 1)}
\bigoplus\nolimits_{j = 1, \ldots, m} B \otimes_A \kappa \to
M \otimes_A \kappa \to 0
$$
Algebra, Lemma \ref{algebra-lemma-flat-tor-zero}를 참조하라.
위 논증을 거꾸로 적용하면 이는 $\bigcup D(g_j \otimes 1)$이
$\Spec(B \otimes_A \kappa)$에서 스킴론적으로 조밀하다는 뜻이다.
$\bigcup D(g_j \otimes 1) = \bigcup D(g_j)_\kappa \subset U_\kappa$이므로
$U_\kappa$가 $X_\kappa$에서 스킴론적으로 조밀함을 얻는다. 이것이
증명하려던 바이다.
\end{proof}

\noindent
스킴의 사상 $f : X \to Y$와 점 $y \in Y$가 주어졌다고 하자.
올 $X_y$는 유도된 위상을 가진 $X$의 부분집합
$f^{-1}(\{y\})$와 위상동형임을 상기하자. Schemes, Lemma
\ref{schemes-lemma-fibre-topological}을 참조하라.
닫힌 부분집합 $T(y) \subset X_y$가 주어졌다고 하자.
$T$를 $X$에서 $T(y)$의 폐포라 하고, $T$에 유도된 축약 스킴 구조를
주자. 그러면 $T$는 집합론적으로 $T_y = T(y)$를 만족하는 $X$의
닫힌 부분스킴이다. 사실 $T$는 이 성질을 가진 $X$의 가장 작은 닫힌
부분스킴이다. 따라서 원한다면 $X_y$의 닫힌 부분집합을 $T_y$로
표기해도 “무방하다”. 다음 보조정리에서는 이를 $f$의 일반올에
적용한다.

\begin{lemma}
\label{more-morphisms-lemma-cover-generic-fibre-neighbourhood}
$f : X \to Y$를 스킴의 유한형 사상이라 하자. $Y$가 일반점
$\eta$를 가진 기약 스킴이라고 가정하자.
$X_\eta = Z_{1, \eta} \cup \ldots \cup Z_{n, \eta}$를 $X_\eta$의
닫힌 부분집합들로 이루어진 일반올의 덮개라 하자.
$Z_i$를 $X$에서 $Z_{i, \eta}$의 폐포라 하자(위 논의를 참조하라).
그러면 모든 $y \in V$에 대하여
$X_y = Z_{1, y} \cup \ldots \cup Z_{n, y}$를 만족하는 비어 있지
않은 열린집합 $V \subset Y$가 존재한다.
\end{lemma}

\begin{proof}
$Y$가 뇌터이면
$U = X \setminus (Z_1 \cup \ldots \cup Z_n)$은 $Y$ 위에서 유한형이고,
Lemma \ref{more-morphisms-lemma-empty-generic-fibre}를 직접 적용하여 비어 있지 않은
열린집합 $V \subset Y$에 대하여 $U_V = \emptyset$을 얻는다.
일반적인 경우에는 다음과 같이 논증한다. 문제는 위상적이므로 $Y$를
그 축약으로 바꾸어도 된다. 따라서 $Y$는 정수적이다. Properties,
Lemma \ref{properties-lemma-characterize-integral}을 참조하라.
$Y$를 축소한 뒤 $X \to Y$가 평탄이라고 가정할 수 있다. Morphisms,
Proposition \ref{morphisms-proposition-generic-flatness}를 참조하라.
이 경우 Morphisms, Lemma
\ref{morphisms-lemma-generalizations-lift-flat}에 의해 $X_y$의 모든 점
$x$는 어떤 점 $x' \in X_\eta$의 특수화이다. $Z_i$들이 $X$에서 닫혀
있고 일반올을 덮으므로, 원하는 대로 $y \in Y$에 대하여
$X_y = \bigcup Z_{i, y}$를 얻는다.
\end{proof}

\noindent
다음 보조정리는 정의역이 축약인 사상의 일반올도 축약임을 말한다.

\begin{lemma}
\label{more-morphisms-lemma-reduction-generic-fibre}
$f : X \to Y$를 스킴의 사상이라 하자. $\eta \in Y$를 $Y$의 한 기약
성분의 일반점이라 하자. 그러면
$(X_\eta)_{red} = (X_{red})_\eta$.
\end{lemma}

\begin{proof}
$\eta$의 아핀 근방 $\Spec(A) \subset Y$를 택하자.
사상 $f$에 의해 $\Spec(A)$로 가는 아핀 열린집합
$\Spec(B) \subset X$를 택하자. $\mathfrak p \subset A$를 $\eta$에
대응하는 극소 소 아이디얼이라 하자. $B_{red}$를 $B$를 멱영근기
$\sqrt{(0)}$로 나눈 몫이라 하자. 보조정리의 대수적 내용은
$C = B_{red} \otimes_A \kappa(\mathfrak p)$가 축약이라는 것이다.
% P05-EN-CH37-0117
$I \subset A$를 멱영근기라 표기하여 $A_{red} = A/I$로 놓자.
$\mathfrak p_{red} = \mathfrak p/I$라 표기하자. 이는 $A_{red}$의
극소 소 아이디얼이고
$\kappa(\mathfrak p) = \kappa(\mathfrak p_{red})$를 만족한다.
$A \to B_{red}$와 $A \to \kappa(\mathfrak p)$가 모두
$A \to A_{red}$를 통해 분해되므로
$C = B_{red} \otimes_{A_{red}} \kappa(\mathfrak p_{red})$.
이제 Algebra, Lemma
\ref{algebra-lemma-minimal-prime-reduced-ring}에 의해
$\kappa(\mathfrak p_{red}) = (A_{red})_{\mathfrak p_{red}}$는
국소화이다. 따라서 $C$는 $B_{red}$의 국소화이고
(Algebra, Lemma \ref{algebra-lemma-tensor-localization}) 그러므로
축약이다.
\end{proof}

\begin{lemma}
\label{more-morphisms-lemma-make-generic-fibre-geometrically-reduced}
$f : X \to Y$를 스킴의 사상이라 하자.
$Y$가 기약이고 $f$가 유한형이라고 가정하자.
다음 조건들을 만족하는 도식이 존재한다.
$$
\xymatrix{
X' \ar[d]_{f'} \ar[r]_{g'} & X_V \ar[r] \ar[d] & X \ar[d]^f \\
Y' \ar[r]^g & V \ar[r] & Y
}
$$
여기서
\begin{enumerate}
\item $V$는 $Y$의 비어 있지 않은 열린집합이고,
\item $X_V = V \times_Y X$,
\item $g : Y' \to V$는 유한 보편 위상동형이며,
\item $X' = (Y' \times_Y X)_{red} = (Y' \times_V X_V)_{red}$,
\item $g'$은 유한 보편 위상동형이고,
\item $Y'$은 정수적 아핀 스킴이며,
\item $f'$은 평탄하고 유한 표시이고,
\item $f'$의 일반올은 기하학적으로 축약이다.
\end{enumerate}
\end{lemma}

\begin{proof}
$V = \Spec(A)$를 $Y$의 비어 있지 않은 아핀 열린집합이라 하자.
가정에 의해 $A$의 멱영근기 $\mathfrak p = \sqrt{(0)}$는 소
아이디얼이다. $K = \kappa(\mathfrak p)$라 하자.
$K$의 표수가 양수이면 그 표수를 $p$라 하고, 표수가 영이면
$1$이라 하자. Varieties, Lemma
\ref{varieties-lemma-finite-extension-geometrically-reduced}에 의해
$(X_{K'})_{red}$가 $K'$ 위에서 기하학적으로 축약이 되도록 하는
유한 순수 비분리 체 확대 $K'/K$가 존재한다.
$K$ 위에서 $K'$을 생성하고 각 $x_i$의 어떤 $p$-거듭제곱이
$A/\mathfrak p$에 속하도록 하는 원소
$x_1, \ldots, x_n \in K'$을 택하자.
$A' \subset K'$을 $x_1, \ldots, x_n$으로 생성되는 $K'$의 유한
$A$-부분대수라 하자. $A'$은 분수체가 $K'$인 정역임에 유의하자.
Algebra, Lemma \ref{algebra-lemma-p-ring-map}에 의해
$A \to A'$은 스펙트럼 위에서 보편 위상동형을 유도한다.
$Y' = \Spec(A')$ 및 $X' = (Y' \times_Y X)_{red}$로 놓자.
Lemma \ref{more-morphisms-lemma-reduction-generic-fibre}에 의해 $X' \to Y'$의
일반올은 $(X_{K'})_{red}$이고, 이는 구성상 기하학적으로 축약이다.
$X' \to X_V$는 축약 사상 $X' \to Y' \times_Y X$
(Morphisms, Lemma
\ref{morphisms-lemma-reduction-universal-homeomorphism} 참조)과
$g$의 밑변환의 합성이므로 유한 보편 위상동형임에 유의하자.
이제 보조정리의 성질들은 (7)을 제외하면 모두 성립한다.
$Y'$을, 따라서 $V$를 축소하면 (7)도 성립하게 할 수 있다. Morphisms,
Proposition \ref{morphisms-proposition-generic-flatness}를 참조하라.
\end{proof}

\begin{lemma}
\label{more-morphisms-lemma-make-components-generic-fibre-geometrically-irreducible}
$f : X \to Y$를 스킴의 사상이라 하자.
$Y$가 기약이고 $f$가 유한형이라고 가정하자.
다음 조건들을 만족하는 도식이 존재한다.
$$
\xymatrix{
X' \ar[d]_{f'} \ar[r]_{g'} & X_V \ar[r] \ar[d] & X \ar[d]^f \\
Y' \ar[r]^g & V \ar[r] & Y
}
$$
여기서
\begin{enumerate}
\item $V$는 $Y$의 비어 있지 않은 열린집합이고,
\item $X_V = V \times_Y X$,
\item $g : Y' \to V$는 전사 유한 에탈 사상이며,
\item $X' = Y' \times_Y X = Y' \times_V X_V$,
\item $g'$은 전사 유한 에탈 사상이고,
\item $Y'$은 기약 아핀 스킴이며,
\item $f'$의 일반올의 모든 기약 성분은 기하학적으로 기약이다.
\end{enumerate}
\end{lemma}

\begin{proof}
$V = \Spec(A)$를 $Y$의 비어 있지 않은 아핀 열린집합이라 하자.
가정에 의해 $A$의 멱영근기 $\mathfrak p = \sqrt{(0)}$는 소
아이디얼이다. $K = \kappa(\mathfrak p)$라 하자. Varieties, Lemma
\ref{varieties-lemma-finite-extension-geometrically-irreducible-components}에
의해 $X_{K'}$의 모든 기약 성분이 $K'$ 위에서 기하학적으로 기약이
되도록 하는 유한 분리 체 확대 $K'/K$가 존재한다.
$K$ 위에서 $K'$을 생성하는 원소 $\alpha \in K'$을 택하자. Fields,
Lemma \ref{fields-lemma-primitive-element}를 참조하라.
% P05-EN-CH37-0118
$P(T) \in K[T]$를 $K$ 위에서 $\alpha$의 최소다항식이라 하자.
어떤 $f \in A$, $f \not \in \mathfrak p$에 대하여 $\alpha$를
$f \alpha$로 바꾼 뒤, 사상 $A[T] \to K[T]$ 아래에서
$P(T) \in K[T]$로 가는 일계수 다항식
$T^d + a_1T^{d - 1} + \ldots + a_d \in A[T]$가 존재한다고
가정할 수 있다. $A' = A[T]/(P)$로 놓자. 그러면 $A \to A'$은
유한 자유 환 사상이고, $\mathfrak p$ 위에 놓이는 유일한 소
아이디얼 $\mathfrak q$가 존재하여
$K = \kappa(\mathfrak p) \subset \kappa(\mathfrak q) = K'$
는 유한 분리 확대이고, $\mathfrak pA'_{\mathfrak q}$는
$A'_{\mathfrak q}$의 극대 아이디얼이다.
따라서 $g : Y' = \Spec(A') \to V = \Spec(A)$는 $\mathfrak q$에서
에탈이다. Algebra, Lemma
\ref{algebra-lemma-characterize-etale}을 참조하라.
이는 $g|_W : W \to \Spec(A)$가 에탈이 되는 열린집합
$W \subset \Spec(A')$가 존재한다는 뜻이다.
$g$가 유한이고 $\mathfrak q$가 $\mathfrak p$ 위에 놓이는 유일한
점이므로 $Z = g(Y' \setminus W)$는 $\mathfrak p$를 포함하지 않는
$V$의 닫힌 부분집합이다. 따라서 $V$를 $Z$와 만나지 않는 $V$의
주 아핀 열린집합으로 바꾼 뒤 $g$가 유한 에탈임을 얻는다.
\end{proof}







\section{상대 동반점}
\label{more-morphisms-section-assassin}

\begin{lemma}
\label{more-morphisms-lemma-relative-assassin-in-neighbourhood}
$f : X \to S$를 스킴의 사상이라 하자.
$\mathcal{F}$를 준연접 $\mathcal{O}_X$-가군이라 하자.
$\xi \in \text{Ass}_{X/S}(\mathcal{F})$라 하고
$Z = \overline{\{\xi\}} \subset X$로 놓자.
$f$가 국소 유한형이고 $\mathcal{F}$가 유한형
$\mathcal{O}_X$-가군이면, 모든 $s \in f(V)$에 대하여 $V_s$의
일반점들이 $\text{Ass}_{X/S}(\mathcal{F})$에 속하도록 하는 비어 있지
않은 열린집합 $V \subset Z$가 존재한다.
\end{lemma}

\begin{proof}
$S$를 $f(\xi)$의 아핀 열린 근방으로, $X$를 $\xi$의 아핀 열린
근방으로 바꾸어도 된다. 따라서 $S = \Spec(A)$,
$X = \Spec(B)$이고 $f$가 유한형 환 사상 $A \to B$로 주어진다고
가정할 수 있다. Morphisms, Lemma
\ref{morphisms-lemma-locally-finite-type-characterize}를 참조하라.
더욱이 어떤 유한 $B$-가군 $M$에 대하여
$\mathcal{F} = \widetilde{M}$이라 쓸 수 있다. Properties, Lemma
\ref{properties-lemma-finite-type-module}을 참조하라.
$\mathfrak q \subset B$를 $\xi$에 대응하는 소아이디얼이라 하고,
$\mathfrak p \subset A$를 이에 대응하는 $A$의 소아이디얼이라 하자.
가정에 의해
$\mathfrak q \in \text{Ass}_B(M \otimes_A \kappa(\mathfrak p))$이다.
Algebra, Remark \ref{algebra-remark-bourbaki} 및
Divisors, Lemma \ref{divisors-lemma-associated-affine-open}을 참조하라.
이 표기에서 $Z = V(\mathfrak q) \subset \Spec(B)$이다.
특히 $f(Z) \subset V(\mathfrak p)$이다. 따라서 $A$, $B$, $M$을 각각
$A/\mathfrak pA$, $B/\mathfrak pB$, $M/\mathfrak pM$으로 바꾼 뒤
보조정리를 증명하면 충분함은 명백하다. 다시 말해 $A$가 분수체
$K$를 가진 정역이고 $\mathfrak q \subset B$가
$M \otimes_A K$의 동반 소아이디얼이라고 가정할 수 있다.

\medskip\noindent
이제 일반 평탄성을 사용할 수 있다. 즉, Algebra, Lemma
\ref{algebra-lemma-generic-flatness}에 의해 $M_g$가 $A_g$-가군으로서
평탄하게 되는 영이 아닌 $g \in A$가 존재한다. $A$를 $A_g$로
바꾸면 $M$이 $A$-가군으로서 평탄이라고 가정할 수 있다.

\medskip\noindent
이 경우 Algebra, Lemma \ref{algebra-lemma-post-bourbaki}에 의해
$\mathfrak q$는 $M$의 동반 소아이디얼이기도 하다.
따라서 단사 $B$-가군 사상 $B/\mathfrak q \to M$을 얻는다.
$Q$를 그 여핵이라 하면 다음 짧은 완전열을 얻는다.
$$
0 \to B/\mathfrak q \to M \to Q \to 0
$$
이는 유한 $B$-가군들의 완전열이다. 일반 평탄성 Algebra, Lemma
\ref{algebra-lemma-generic-flatness}를 이번에는 $B$-가군 $Q$에 한 번 더
적용하면, $Q$가 평탄한 $A$-가군이라고 가정할 수 있다. 특히 위 짧은
완전열이 보편 단사라고 가정할 수 있다. Algebra, Lemma
\ref{algebra-lemma-flat-tor-zero}를 참조하라. 이 상황에서는
$(B/\mathfrak q) \otimes_A \kappa(\mathfrak p')
\subset M \otimes_A \kappa(\mathfrak p')$
가 $A$의 모든 소아이디얼 $\mathfrak p'$에 대하여 성립한다.
보조정리는 다음 사실에서 따른다. 즉,
$(B/\mathfrak q) \otimes_A \kappa(\mathfrak p')$
의 받침의 극소 소아이디얼 $\mathfrak q'$는 Divisors, Lemma
\ref{divisors-lemma-minimal-support-in-ass}에 의해
$(B/\mathfrak q) \otimes_A \kappa(\mathfrak p')$의 동반
소아이디얼이다.
\end{proof}

\begin{lemma}
\label{more-morphisms-lemma-bad-case}
$f : X \to Y$를 스킴의 사상이라 하자. $\mathcal{F}$를 준연접
$\mathcal{O}_X$-가군이라 하자. $U \subset X$를 열린 부분스킴이라
하자. 다음을 가정하자.
\begin{enumerate}
\item $f$는 유한형이다.
\item $\mathcal{F}$는 유한형이다.
\item $Y$는 일반점 $\eta$를 가진 기약 스킴이다.
\item $\text{Ass}_{X_\eta}(\mathcal{F}_\eta)$는 $U_\eta$에 포함되지 않는다.
\end{enumerate}
그러면 모든 $y \in V$에 대하여 집합
$\text{Ass}_{X_y}(\mathcal{F}_y)$가 $U_y$에 포함되지 않도록 하는
비어 있지 않은 열린 부분스킴 $V \subset Y$가 존재한다.
\end{lemma}

\begin{proof}
% P05-EN-CH37-0119
$Z \subset X$를 $\mathcal{F}$의 스킴론적 받침이라 하자. Morphisms,
Definition \ref{morphisms-definition-scheme-theoretic-support}를 참조하라.
그러면 $Z_\eta$는 $\mathcal{F}_\eta$의 스킴론적 받침이다
(Morphisms, Lemma \ref{morphisms-lemma-flat-pullback-support}).
따라서 Divisors, Lemma \ref{divisors-lemma-minimal-support-in-ass}에
의해 $Z_\eta$의 기약 성분들의 일반점들은
$\text{Ass}_{X_\eta}(\mathcal{F}_\eta)$에 포함된다.
따라서 $Z_\eta \cap U_\eta = \emptyset$임을 알 수 있다.
그러므로 $T = Z \setminus U$는 $T_\eta = \emptyset$을 만족하는
$Z$의 닫힌 부분집합이다. $T$에 유도된 축약 스킴 구조를 주면
$T \to Y$는 유한형 사상이다. Lemma
\ref{more-morphisms-lemma-empty-generic-fibre}에 의해 $T_V = \emptyset$을 만족하는
비어 있지 않은 열린집합 $V \subset Y$가 존재한다. 이 $V$가 조건을
만족한다.
\end{proof}

\begin{lemma}
\label{more-morphisms-lemma-good-case}
$f : X \to Y$를 스킴의 사상이라 하자. $\mathcal{F}$를 준연접
$\mathcal{O}_X$-가군이라 하자. $U \subset X$를 열린 부분스킴이라
하자. 다음을 가정하자.
\begin{enumerate}
\item $f$는 유한형이다.
\item $\mathcal{F}$는 유한형이다.
\item $Y$는 일반점 $\eta$를 가진 기약 스킴이다.
\item $\text{Ass}_{X_\eta}(\mathcal{F}_\eta) \subset U_\eta$.
\end{enumerate}
그러면 모든 $y \in V$에 대하여
$\text{Ass}_{X_y}(\mathcal{F}_y) \subset U_y$를 만족하는 비어 있지
않은 열린 부분스킴 $V \subset Y$가 존재한다.
\end{lemma}

\begin{proof}
(이 증명은 Lemma
\ref{more-morphisms-lemma-scheme-theoretically-dense-generic-fibre}의 증명과 같다.
독자에게 먼저 그 증명을 읽기를 권한다.)
명제가 올에 관한 것이므로 $Y$를 그 축약으로 바꾸어도 됨은 명백하다.
따라서 $Y$가 정수적이라고 가정할 수 있다. Properties, Lemma
\ref{properties-lemma-characterize-integral}을 참조하라.
또한 $Y = \Spec(A)$가 아핀이라고 가정할 수 있다. 그러면 $A$는
분수체 $K$를 가진 정역이다.

\medskip\noindent
$f$가 유한형이므로 $X$는 준콤팩트이다. 어떤 아핀 열린집합
$X_i$들에 대하여 $X = X_1 \cup \ldots \cup X_n$이라 쓰고
$\mathcal{F}_i = \mathcal{F}|_{X_i}$로 놓자. 가정에 의해
$U_i = X_i \cap U$의 일반올은
$\text{Ass}_{X_{i, \eta}}(\mathcal{F}_{i, \eta})$를 포함한다.
따라서 각 삼중항 $(X_i, \mathcal{F}_i, U_i)$에 대하여 결과를
증명하면 충분하다. 다시 말해 $X$가 아핀이라고 가정할 수 있다.

\medskip\noindent
$X = \Spec(B)$라 쓰자. $\mathcal{F} = \widetilde{N}$을 만족하는 유한
$B$-가군을 $N$이라 하자. $B_K$는 유한형 $K$-대수이므로 뇌터임에
유의하자. 따라서 $U_\eta$는 준콤팩트이다. 그러므로
$D(g_j) \subset U$이고
$U_\eta = D(g_1)_\eta \cup \ldots \cup D(g_m)_\eta$가 되도록 하는
유한 개의 $g_1, \ldots, g_m \in B$를 찾을 수 있다.
$\text{Ass}_{X_\eta}(\mathcal{F}_\eta) \subset U_\eta$이므로
$N_K \to \bigoplus_j (N_K)_{g_j}$는 단사이다. Algebra, Lemma
\ref{algebra-lemma-when-injective-covering}에 의해 이는 다음 사상의
단사성과 동치이다.
$N_K \to \bigoplus\nolimits_{j = 1, \ldots, m} N_K$,
$n \mapsto (g_1n, \ldots, g_mn)$. 이 사상을 $A$ 위에서 취했을 때의
핵과 여핵을 각각 $I$와 $M$이라 하자. 즉, 다음 완전열이 있다.
$$
0 \to I \to N \xrightarrow{(g_1, \ldots, g_m)}
\bigoplus\nolimits_{j = 1, \ldots, m} N \to M \to 0
$$
어떤 영이 아닌 $h$에 대하여 $A$를 $A_h$로 바꾸면, $B$가 평탄한
유한 표시 $A$-대수이고 $M$과 $N$이 모두 $A$ 위에서 평탄이라고
가정할 수 있다. Algebra, Lemma
\ref{algebra-lemma-generic-flatness}를 참조하라. $N$이 $A$ 위에서
평탄하므로 $N$은 $A$-가군으로서 꼬임이 없다. More on Algebra,
Lemma \ref{more-algebra-lemma-flat-torsion-free}를 참조하라.
따라서 $N \subset N_K$이다. 구성에 의해 $I_K = 0$이고, 이는
$I = 0$을 함의한다($I \subset N \subset N_K$가 $I_K$의 부분집합이기
때문이다). 그러므로 이제 다음 짧은 완전열이 있다.
$$
0 \to N \xrightarrow{(g_1, \ldots, g_m)}
\bigoplus\nolimits_{j = 1, \ldots, m} N \to M \to 0
$$
여기서 $M$은 $A$ 위에서 평탄하다. 따라서 $\kappa$가 체일 때의 모든
준동형 $A \to \kappa$에 대하여 다음 짧은 완전열을 얻는다.
$$
0 \to N \otimes_A \kappa \xrightarrow{(g_1 \otimes 1, \ldots, g_m \otimes 1)}
\bigoplus\nolimits_{j = 1, \ldots, m} N \otimes_A \kappa \to
M \otimes_A \kappa \to 0
$$
Algebra, Lemma \ref{algebra-lemma-flat-tor-zero}를 참조하라.
위 논증을 거꾸로 적용하면 이는 $\bigcup D(g_j \otimes 1)$이
$\text{Ass}_{B \otimes_A \kappa}(N \otimes_A \kappa)$를 포함한다는
뜻이다. $\bigcup D(g_j \otimes 1) =
\bigcup D(g_j)_\kappa \subset U_\kappa$이므로 $U_\kappa$가
$\text{Ass}_{X \otimes \kappa}(\mathcal{F} \otimes \kappa)$를
포함함을 얻는다. 이것이 증명하려던 바이다.
\end{proof}

\begin{lemma}
\label{more-morphisms-lemma-base-change-assassin-in-U}
$f : X \to S$를 국소 유한형인 사상이라 하자.
$\mathcal{F}$를 유한형 준연접 $\mathcal{O}_X$-가군이라 하자.
$U \subset X$를 열린 부분스킴이라 하자. $g : S' \to S$를 스킴의
사상이라 하고, $f' : X' = X_{S'} \to S'$을 $f$의 밑변환이라 하며,
$g' : X' \to X$를 사영이라 하자.
$\mathcal{F}' = (g')^*\mathcal{F}$ 및 $U' = (g')^{-1}(U)$로 놓자.
끝으로 $s' \in S'$의 상이 $s = g(s')$라고 하자. 이 경우
$$
\text{Ass}_{X_s}(\mathcal{F}_s) \subset U_s
\Leftrightarrow
\text{Ass}_{X'_{s'}}(\mathcal{F}'_{s'}) \subset U'_{s'}.
$$
\end{lemma}

\begin{proof}
Divisors, Lemma \ref{divisors-lemma-base-change-relative-assassin}에서
즉시 따른다. Divisors, Remark
\ref{divisors-remark-base-change-relative-assassin}도 참조하라.
\end{proof}

\begin{lemma}
\label{more-morphisms-lemma-relative-assassin-constructible}
$f : X \to Y$를 유한 표시인 사상이라 하자.
$\mathcal{F}$를 유한 표시인 준연접 $\mathcal{O}_X$-가군이라 하자.
$U \to Y$가 준콤팩트가 되도록 $U \subset X$를 열린 부분스킴이라
하자. 그러면 집합
$$
E = \{y \in Y \mid \text{Ass}_{X_y}(\mathcal{F}_y) \subset U_y\}
$$
는 $Y$에서 국소 구성가능이다.
\end{lemma}

\begin{proof}
$y \in Y$라 하자. $E \cap V$가 $V$에서 구성가능이 되도록 하는
$Y$에서의 $y$의 열린 근방 $V$가 존재함을 보여야 한다. 따라서
$Y$가 아핀이라고 가정할 수 있다. $Y = \Spec(A)$라 쓰자. 유한형
$\mathbf{Z}$-대수들의 유향 극한
$A = \colim A_i$로 쓰자. Limits, Lemma
\ref{limits-lemma-descend-finite-presentation}에 의해, $Y$로 밑변환하면
$f$를 복원하는 유한 표시인 사상
$f_i : X_i \to \Spec(A_i)$와 지표 $i$를 찾을 수 있다.
$i$를 필요하면 증가시켜, $X$로의 당김이 $\mathcal{F}$와 동형인
유한 표시 준연접 $\mathcal{O}_{X_i}$-가군 $\mathcal{F}_i$가
존재한다고 가정할 수 있다. Limits, Lemma
\ref{limits-lemma-descend-modules-finite-presentation}을 참조하라.
$i$를 필요하면 한 번 더 증가시켜, $X$에서의 역상이 $U$인 열린
부분스킴 $U_i \subset X_i$가 존재한다고 가정할 수 있다. Limits,
Lemma \ref{limits-lemma-descend-opens}를 참조하라.
Lemma \ref{more-morphisms-lemma-base-change-assassin-in-U}에 의해 $f_i$에 대하여
보조정리를 증명하면 충분하다. 따라서 $Y$가 뇌터환의 스펙트럼인
경우로 환원된다.

\medskip\noindent
$Y$가 뇌터 스킴인 경우 $E$가 구성가능임을 증명하기 위해 Topology,
Lemma \ref{topology-lemma-characterize-constructible-Noetherian}의
판정법을 사용하겠다. 이를 위해 $Z \subset Y$를 기약 닫힌
부분스킴이라 하자. $E \cap Z$가 비어 있지 않은 열린 부분집합을
포함하거나 $Z$에서 조밀하지 않음을 보여야 한다. 이는 $Z$ 위에서의
밑변환 $(X, \mathcal{F}, U) \times_Y Z$에 Lemmas
\ref{more-morphisms-lemma-bad-case} 및 \ref{more-morphisms-lemma-good-case}를 적용하면 따른다.
\end{proof}











\section{축약 올}
\label{more-morphisms-section-reduced}

\begin{lemma}
\label{more-morphisms-lemma-nonreduced-in-neighbourhood}
$f : X \to Y$를 스킴의 사상이라 하자. $Y$가 일반점 $\eta$를 가진
기약 스킴이고 $f$가 유한형이라고 가정하자. $X_\eta$가 축약이
아니면, 모든 $y \in V$에 대하여 올 $X_y$가 축약이 아니도록 하는
비어 있지 않은 열린집합 $V \subset Y$가 존재한다.
\end{lemma}

\begin{proof}
$Y' \subset Y$를 $Y$의 축약이라 하자. $X' \to Y'$을 $f$의
밑변환이라 하자. $Y' \to Y$가 점들의 전단사를 유도하고
$X' \to X$가 올들을 식별함에 유의하자.
% P05-EN-CH37-0120
따라서 $Y'$이 축약, 즉 정수적이라고 가정할 수 있다. Properties,
Lemma \ref{properties-lemma-characterize-integral}을 참조하라.
또한 $Y$를 아핀 열린집합으로 바꾸어도 된다. 따라서 $A$가 정역이고
$Y = \Spec(A)$라고 가정할 수 있다.
% P05-EN-CH37-0121
$K$를 $A$의 분수체라 하자.
아핀 열린집합 $\Spec(B) = U \subset X$와 영이 아닌 멱영 절단
$h_\eta \in \Gamma(U_\eta, \mathcal{O}_{U_\eta}) = B_K$를 택하자.
$Y$를 축소한 뒤 $h$가
$h \in \Gamma(U, \mathcal{O}_U) = B$에서 온다고 가정할 수 있다.
$Y$를 조금 더 축소하면 $h$가 멱영이라고 가정할 수 있다.
$I = \{b \in B \mid hb = 0\}$를 $h$의 소멸자라 하자.
그러면 $C = B/I$는 일반올 $(B/I)_K$가 영이 아닌 유한형 $A$-대수이다
($h_\eta \not = 0$이므로). $A \to C$와 $A \to B/hB$에 일반 평탄성을
적용하자. Algebra, Lemma \ref{algebra-lemma-generic-flatness}를
참조하라. 그러면 $C_g$가 자유 $A_g$-가군이고 $(B/hB)_g$가 평탄한
$A_g$-가군이 되도록 하는 $g \in A$, $g \not = 0$을 얻는다.
$Y$를 $D(g) \subset Y$로 바꾸자. 이제 다음 짧은 완전열이 있다.
$$
0 \to C \to B \to B/hB \to 0.
$$
여기서 $B/hB$는 $A$ 위에서 평탄하고 $C$는 영이 아닌 자유
$A$-가군이다. 따라서 체로 가는 임의의 준동형 $A \to \kappa$에
대하여 환 $C \otimes_A \kappa$는 영이 아니고 다음 열은 완전하다.
$$
0 \to C \otimes_A \kappa \to B \otimes_A \kappa \to
B/hB \otimes_A \kappa \to 0
$$
Algebra, Lemma \ref{algebra-lemma-flat-tor-zero}를 참조하라.
텐서곱의 오른쪽 완전성에 의해
$B/hB \otimes_A \kappa = (B \otimes_A \kappa) / h(B \otimes_A \kappa)$
임에 유의하자. 따라서 $h$에 의한 곱셈은
$B \otimes_A \kappa$ 위에서 영이 아니라고 결론짓는다.
이는 모든 점 $y \in Y$에 대하여 원소 $h$의 $U_y$로의 제한이
영이 아님을 명백히 뜻하고, 따라서 $X_y$는 축약이 아니다.
\end{proof}

\begin{lemma}
\label{more-morphisms-lemma-base-change-fibres-geometrically-reduced}
$f : X \to Y$를 스킴의 사상이라 하자.
$g : Y' \to Y$를 임의의 사상이라 하고,
$f' : X' \to Y'$을 $f$의 밑변환이라 하자. 그러면
\begin{align*}
\{y' \in Y' \mid X'_{y'}\text{는 기하학적으로 축약이다}\} \\
= g^{-1}(\{y \in Y \mid X_y\text{는 기하학적으로 축약이다}\}).
\end{align*}
\end{lemma}

\begin{proof}
이는 상이 $y \in Y$인 $y' \in Y'$에 대하여 올
$X'_{y'} = X_y \times_y y'$가 $\kappa(y')$ 위에서 기하학적으로
축약일 필요충분조건이 $X_y$가 $\kappa(y)$ 위에서 기하학적으로
축약인 것이라는 명제로 귀착된다. 이는 Varieties, Lemma
\ref{varieties-lemma-geometrically-reduced-upstairs}에서 따른다.
\end{proof}

\begin{lemma}
\label{more-morphisms-lemma-not-geometrically-reduced-in-neighbourhood}
$f : X \to Y$를 스킴의 사상이라 하자. $Y$가 일반점 $\eta$를 가진
기약 스킴이고 $f$가 유한형이라고 가정하자. $X_\eta$가 기하학적으로
축약이 아니면, 모든 $y \in V$에 대하여 올 $X_y$가 기하학적으로
축약이 아니도록 하는 비어 있지 않은 열린집합 $V \subset Y$가 존재한다.
\end{lemma}

\begin{proof}
Lemma \ref{more-morphisms-lemma-make-generic-fibre-geometrically-reduced}를 적용하여
$$
\xymatrix{
X' \ar[d]_{f'} \ar[r]_{g'} & X_V \ar[r] \ar[d] & X \ar[d]^f \\
Y' \ar[r]^g & V \ar[r] & Y
}
$$
그 보조정리에서 언급한 모든 성질을 가진 도식을 얻는다.
$\eta'$을 $Y'$의 일반점이라 하자.
사상 $X' \to X_{Y'}$(이는 축약 사상이다)과 그로부터 얻는 일반올의
사상 $X'_{\eta'} \to X_{\eta'}$을 생각하자.
$X'_{\eta'}$은 기하학적으로 축약이지만 $X_\eta$는 그렇지 않으므로
이 사상은 동형일 수 없다. Varieties, Lemma
\ref{varieties-lemma-geometrically-reduced-upstairs}를 참조하라.
따라서 $X_{\eta'}$은 축약이 아니다. 그러므로 Lemma
\ref{more-morphisms-lemma-nonreduced-in-neighbourhood}에 의해, 비어 있지 않은
열린집합 $V' \subset Y'$의 모든 점 $y' \in V'$에서
$X_{Y'} \to Y'$의 올은 축약이 아니다. $g : Y' \to V$가
위상동형이므로 Lemma
\ref{more-morphisms-lemma-base-change-fibres-geometrically-reduced}에 의해
$g(V')$가 찾던 열린집합이다.
\end{proof}

\begin{lemma}
\label{more-morphisms-lemma-geometrically-reduced-generic-fibre}
$f : X \to Y$를 스킴의 사상이라 하자. 다음을 가정하자.
\begin{enumerate}
\item $Y$는 일반점 $\eta$를 가진 기약 스킴이다.
\item $X_\eta$는 기하학적으로 축약이다.
\item $f$는 유한형이다.
\end{enumerate}
그러면 $X_V \to V$의 올들이 기하학적으로 축약이 되도록 하는 비어
있지 않은 열린 부분스킴 $V \subset Y$가 존재한다.
\end{lemma}

\begin{proof}
$Y' \subset Y$를 $Y$의 축약이라 하자. $X' \to Y'$을 $f$의
밑변환이라 하자. $Y' \to Y$가 점들의 전단사를 유도하고
$X' \to X$가 올들을 식별함에 유의하자.
% P05-EN-CH37-0120
따라서 $Y'$이 축약, 즉 정수적이라고 가정할 수 있다. Properties,
Lemma \ref{properties-lemma-characterize-integral}을 참조하라.
또한 $Y$를 아핀 열린집합으로 바꾸어도 된다. 따라서 $A$가 정역이고
$Y = \Spec(A)$라고 가정할 수 있다.
% P05-EN-CH37-0121
$K$를 $A$의 분수체라 하자. $Y$를 조금 축소하면 $X \to Y$가
평탄하고 유한 표시라고 가정할 수도 있다. Morphisms, Proposition
\ref{morphisms-proposition-generic-flatness}를 참조하라.

\medskip\noindent
$X_\eta$가 기하학적으로 축약이므로 $V \to \Spec(K)$가 매끄러운
열린 조밀 부분집합 $V \subset X_\eta$가 존재한다. Varieties, Lemma
\ref{varieties-lemma-geometrically-reduced-dense-smooth-open}을 참조하라.
$U \subset X$를 $f$가 매끄러운 점들의 집합이라 하자. Morphisms,
Lemma \ref{morphisms-lemma-set-points-where-fibres-smooth}에 의해
$V \subset U_\eta$임을 알 수 있다. 따라서 $U$의 일반올은 $X$의
일반올에서 조밀하다. $X_\eta$가 축약이므로 $U_\eta$는 $X_\eta$에서
스킴론적으로 조밀하다. Morphisms, Lemma
\ref{morphisms-lemma-reduced-scheme-theoretically-dense}를 참조하라.
$U \to Y$가 매끄러우므로 $U \to Y$의 모든 올은 기하학적으로
축약임에 유의하자. 따라서 $Y$를 축소한 뒤 모든 $y \in Y$에 대하여
스킴 $U_y$가 $X_y$에서 스킴론적으로 조밀함을 보이면 충분하다.
Morphisms, Lemma \ref{morphisms-lemma-reduced-subscheme-closure}를
참조하라. 이는 Lemma
\ref{more-morphisms-lemma-scheme-theoretically-dense-generic-fibre}에서 따른다.
\end{proof}

\begin{lemma}
\label{more-morphisms-lemma-geometrically-reduced-constructible}
$f : X \to Y$를 준콤팩트이고 국소 유한 표시인 사상이라 하자.
그러면 집합
$$
E = \{y \in Y \mid X_y\text{는 기하학적으로 축약이다}\}
$$
는 $Y$에서 국소 구성가능이다.
\end{lemma}

\begin{proof}
$y \in Y$라 하자. $E \cap V$가 $V$에서 구성가능이 되도록 하는
$Y$에서의 $y$의 열린 근방 $V$가 존재함을 보여야 한다. 따라서
$Y$가 아핀이라고 가정할 수 있다. 그러면 $X$는 준콤팩트이다.
유한 아핀 열린 덮개 $X = U_1 \cup \ldots \cup U_n$을 택하자.
그러면 $y$에서 $U_i \to Y$의 올들은 $y$에서 $X \to Y$의 올의
아핀 열린 덮개를 이룬다. 따라서 $X$도 아핀이라고 가정할 수 있다.
$Y = \Spec(A)$라 쓰자. 유한형 $\mathbf{Z}$-대수들의 유향 극한
$A = \colim A_i$로 쓰자. Limits, Lemma
\ref{limits-lemma-descend-finite-presentation}에 의해, $Y$로
밑변환하면 $f$를 복원하는 유한 표시인 사상
$f_i : X_i \to \Spec(A_i)$와 지표 $i$를 찾을 수 있다.
Lemma \ref{more-morphisms-lemma-base-change-fibres-geometrically-reduced}에 의해
$f_i$에 대하여 보조정리를 증명하면 충분하다. 따라서 $Y$가 뇌터환의
스펙트럼인 경우로 환원된다.

\medskip\noindent
$Y$가 뇌터 스킴인 경우 $E$가 구성가능임을 증명하기 위해 Topology,
Lemma \ref{topology-lemma-characterize-constructible-Noetherian}의
판정법을 사용하겠다. 이를 위해 $Z \subset Y$를 일반점 $\xi$를 가진
기약 닫힌 부분스킴이라 하자. $E \cap Z$가 비어 있지 않은 열린
부분집합을 포함하거나 $Z$에서 조밀하지 않음을 보여야 한다.
$X_\xi$가 기하학적으로 축약이면, 사상 $X_Z \to Z$에 적용한
Lemma \ref{more-morphisms-lemma-geometrically-reduced-generic-fibre}에 의해 비어
있지 않은 열린집합 $V \subset Z$의 모든 올 $X_y$가 기하학적으로
축약이다. $X_\xi$가 기하학적으로 축약이 아니면, 사상
$X_Z \to Z$에 적용한
Lemma \ref{more-morphisms-lemma-not-geometrically-reduced-in-neighbourhood}에 의해
비어 있지 않은 열린집합 $V \subset Z$의 모든 올 $X_y$가
기하학적으로 축약이 아니다. 이로써 증명이 끝난다.
\end{proof}

\begin{lemma}
\label{more-morphisms-lemma-proper-flat-over-dvr-reduced-fibre}
$R$이 이산 값매김환일 때 $f : X \to \Spec(R)$를 고유 사상이라 하자.
$X$의 모든 기약 성분이 $\Spec(R)$ 위에서 우세하다고 가정하고
(예를 들어 $f$가 평탄이면 그러하다), 특수올이 축약이라고 가정하자.
그러면 $X$와 일반올 $X_\eta$는 모두 축약이다.
\end{lemma}

\begin{proof}
(평탄 사상이 기약 성분에 관한 가정을 만족한다는 사실은 예를 들어
Divisors, Lemma \ref{divisors-lemma-bourbaki} 또는 평탄 환 사상에 대한
내려가기에서 따른다.)
특수올 $X_s$가 축약이라고 가정하자. $x \in X$를 임의의 점이라 하고
$\mathcal{O}_{X, x}$가 축약임을 보이자. 그러면 $X$와 $X_\eta$가
축약임이 증명된다. $x'$이 특수올에 속하는 특수화
$x \leadsto x'$을 택하자. 고유 사상은 닫힌 사상이므로 이러한
특수화가 존재한다. 국소환 $A = \mathcal{O}_{X, x'}$를 생각하자.
그러면 $\mathcal{O}_{X, x}$는 $A$의 국소화이므로 $A$가 축약임을
보이면 충분하다. $\pi \in R$을 균등화원이라 하자.
$a \in A$가 영이 아니면 $a = \pi^n a'$ 및 $a' \not \in \pi A$를
만족하는 $n \geq 0$과 원소 $a' \in A$가 존재한다.
이는 Krull 교차 정리
(Algebra, Lemma \ref{algebra-lemma-intersect-powers-ideal-module-zero})에서
따른다.
% P05-EN-CH37-0122
$a$가 멱영이면 $a$는 $A$의 모든 극소 소아이디얼에 포함되고, 따라서
$a'$도 그러하다. $X$의 기약 성분에 관한 가정에 의해 $\pi$가 $A$의
어떤 극소 소아이디얼에도 포함되지 않기 때문이다. 따라서 $a'$는 멱영이다.
하지만 $a'$는 축약환 $A/\pi A = \mathcal{O}_{X_s, x'}$의 영이 아닌
원소로 간다. $A$가 축약이 아니라면 이는 모순이다. 이것이 보이려던 바이다.
\end{proof}

\begin{lemma}
\label{more-morphisms-lemma-geometrically-reduced-open}
$f : X \to Y$를 평탄한 유한 표시 고유 사상이라 하자.
그러면 집합 $\{y \in Y \mid X_y\text{는 기하학적으로 축약이다}\}$는
$Y$에서 열려 있다.
\end{lemma}

\begin{proof}
$Y$가 아핀이라고 가정할 수 있다. 그러면 $Y$는 $\mathbf{Z}$ 위의
유한형 아핀 스킴들의 공여과 극한이다.
따라서 $Y_0$가 유한형 $\mathbf{Z}$-대수의 스펙트럼이고
$X_0 \to Y_0$가 평탄하고 고유일 때, $X \to Y$가 $X_0 \to Y_0$의
밑변환이라고 가정할 수 있다. Limits, Lemma
\ref{limits-lemma-descend-finite-presentation},
\ref{limits-lemma-descend-flat-finite-presentation} 및
\ref{limits-lemma-eventually-proper}를 참조하라.
올이 기하학적으로 축약인 점들의 집합의 형성은 밑변환과 교환하므로
(Lemma \ref{more-morphisms-lemma-base-change-fibres-geometrically-reduced}),
밑스킴이 뇌터라고 가정할 수 있다.

\medskip\noindent
$Y$가 뇌터라고 가정하자. Lemma
\ref{more-morphisms-lemma-geometrically-reduced-constructible}에 의해 이 집합은
구성가능이다. 따라서 일반화에 대해 안정함을 보이면 충분하다
(Topology, Lemma \ref{topology-lemma-characterize-closed-Noetherian}).
Properties, Lemma
\ref{properties-lemma-locally-Noetherian-specialization-dvr}에 의해
$R$이 이산 값매김환이고 $Y = \Spec(R)$이며 닫힌올 $X_y$가
기하학적으로 축약인 경우로 환원된다. 보여야 할 것은 일반올
$X_\eta$가 기하학적으로 축약이라는 것이다.

\medskip\noindent
그렇지 않다면 $X_L$이 축약이 아니도록 하는 $R$의 분수체의 유한
확대 $L$이 존재한다. Varieties, Lemma
\ref{varieties-lemma-geometrically-reduced}를 참조하라.
분수체가 $L$이고 $R$을 지배하는 이산 값매김환 $R' \subset L$이
존재한다. Algebra, Lemma
\ref{algebra-lemma-integral-closure-Dedekind}를 참조하라.
$R$을 $R'$으로 바꾸면
Lemma \ref{more-morphisms-lemma-proper-flat-over-dvr-reduced-fibre}로 환원된다.
\end{proof}







\section{올의 기약 성분}
\label{more-morphisms-section-irreducible}

\begin{lemma}
\label{more-morphisms-lemma-irreducible-components-in-neighbourhood}
$f : X \to Y$를 스킴의 사상이라 하자. $Y$가 일반점 $\eta$를 가진
기약 스킴이고 $f$가 유한형이라고 가정하자. $X_\eta$의 기약 성분이
$n$개이면, 모든 $y \in V$에 대하여 올 $X_y$의 기약 성분이 적어도
$n$개가 되도록 하는 비어 있지 않은 열린집합 $V \subset Y$가 존재한다.
\end{lemma}

\begin{proof}
문제가 순전히 위상적이므로 $X$와 $Y$를 각각 그 축약으로 바꾸어도 된다.
특히 그러면 $Y$가 정수적이다. Properties, Lemma
\ref{properties-lemma-characterize-integral}을 참조하라.
$X_\eta = X_{1, \eta} \cup \ldots \cup X_{n, \eta}$를 $X_\eta$의
기약 성분 분해라 하자. $X_i \subset X$를 일반올이 $X_{i, \eta}$인
축약 닫힌 부분스킴이라 하자. $Z_{i, j} = X_i \cap X_j$는 $X_i$의
닫힌 부분집합이고 그 일반올 $Z_{i, j, \eta}$는 $X_{i, \eta}$에서
어디에서도 조밀하지 않음에 유의하자. 따라서 $Y$를 축소하면 모든
$y \in Y$에 대하여 $Z_{i, j, y}$가 $X_{i, y}$에서 어디에서도
조밀하지 않다고 가정할 수 있다. Lemma
\ref{more-morphisms-lemma-nowhere-dense-generic-fibre}를 참조하라.
$Y$를 더 축소하면 $y \in Y$에 대하여 $X_y = \bigcup X_{i, y}$라고
가정할 수 있다. Lemma \ref{more-morphisms-lemma-cover-generic-fibre-neighbourhood}를
참조하라. 더욱이 $Y$를 축소하면 각 $X_i \to Y$가 평탄하고 유한
표시라고 가정할 수 있다. Morphisms, Proposition
\ref{morphisms-proposition-generic-flatness}를 참조하라.
사상 $X_i \to Y$는 열린 사상이다. Morphisms, Lemma
\ref{morphisms-lemma-fppf-open}을 참조하라.
따라서 각 $i$에 대하여 $f(X_i)$에 포함되는 $\eta$의 열린 근방 $V$가
존재한다. 각 $y \in V$에 대하여 스킴 $X_{i, y}$들은 $X_y$의 비어 있지
않은 닫힌 부분집합이고, $X_y = \bigcup X_{i, y}$이며, 교집합
$Z_{i, j, y} = X_{i, y} \cap X_{j, y}$는 $X_{i, y}$에서 조밀하지 않다.
이는 $X_y$의 기약 성분이 적어도 $n$개임을 명백히 함의한다.
\end{proof}

\begin{lemma}
\label{more-morphisms-lemma-base-change-fibres-geometrically-irreducible}
$f : X \to Y$를 스킴의 사상이라 하자.
$g : Y' \to Y$를 임의의 사상이라 하고,
$f' : X' \to Y'$을 $f$의 밑변환이라 하자. 그러면
\begin{align*}
\{y' \in Y' \mid X'_{y'}\text{는 기하학적으로 기약이다}\} \\
= g^{-1}(\{y \in Y \mid X_y\text{는 기하학적으로 기약이다}\}).
\end{align*}
\end{lemma}

\begin{proof}
이는 상이 $y \in Y$인 $y' \in Y'$에 대하여 올
$X'_{y'} = X_y \times_y y'$가 $\kappa(y')$ 위에서 기하학적으로
기약일 필요충분조건이 $X_y$가 $\kappa(y)$ 위에서 기하학적으로
기약인 것이라는 명제로 귀착된다. 이는 Varieties,
Lemma \ref{varieties-lemma-geometrically-irreducible-check-after-extension}에서
따른다.
\end{proof}

\begin{lemma}
\label{more-morphisms-lemma-base-change-fibres-nr-geometrically-irreducible-components}
$f : X \to Y$를 스킴의 사상이라 하자. 함수
$$
n_{X/Y} : Y \to \{0, 1, 2, 3, \ldots, \infty\}
$$
를 $y \in Y$에 $(X_y)_K$의 기약 성분의 수를 대응시키는 함수라 하자.
여기서 $K$는 $\kappa(y)$의 분리적으로 닫힌 확대체이다.
이는 잘 정의되며, $g : Y' \to Y$가 사상이면
$$
n_{X'/Y'} = n_{X/Y} \circ g
$$
이다. 여기서 $X' \to Y'$은 $f$의 밑변환이다.
\end{lemma}

\begin{proof}
$y' \in Y'$의 상이 $y \in Y$라고 하자.
$K \supset \kappa(y)$와 $K' \supset \kappa(y')$가 분리적으로 닫힌
확대체라고 하자. 그러면 체들의 가환 도식
$$
\xymatrix{
K \ar[r] & K'' & K' \ar[l] \\
\kappa(y) \ar[u] \ar[rr] & & \kappa(y') \ar[u]
}
$$
을 택할 수 있다. 스킴의 사상
$$
\xymatrix{
(X'_{y'})_{K'} &
(X'_{y'})_{K''} = (X_y)_{K''} \ar[l] \ar[r] &
(X_y)_K
}
$$
들이 기약 성분 사이의 전단사를 유도하므로 결과가 따른다.
Varieties,
Lemma \ref{varieties-lemma-separably-closed-field-irreducible-components}를
참조하라.
\end{proof}

\begin{lemma}
\label{more-morphisms-lemma-irreducible-polynomial-over-domain}
$A$를 분수체 $K$를 가진 정역이라 하자.
$P \in A[x_1, \ldots, x_n]$이라 하자.
% P05-EN-CH37-0123
$\overline{K}$를 $K$의 대수적 폐포라 하자.
$P$가 $\overline{K}[x_1, \ldots, x_n]$에서 기약이라고 가정하자.
그러면 체로 가는 모든 준동형 $\varphi : A_f \to \kappa$에 대하여
$P^\varphi \in \kappa[x_1, \ldots, x_n]$가 기약이 되도록 하는
$f \in A$가 존재한다.
\end{lemma}

\begin{proof}
$A$ 위의 $A[x_1, \ldots, x_n]$의 자기동형 $\Psi$ 중에서
$\Psi(P) = ax_n^d +$ ($x_n$에 관한 저차항들)이고 $a \not = 0$을
만족하는 것이 존재한다. Algebra, Lemma
\ref{algebra-lemma-helper-polynomial}을 참조하라.
$P$를 $\Psi(P)$로 바꾸고 $A$를 $A_a$로 바꾸어도 된다.
따라서 $P$가 $x_n$에 관한 차수가 $d > 0$인 일계수 다항식이라고
가정할 수 있다. $i = 1, \ldots, n - 1$에 대하여 $d_i$를
$P$의 $x_i$에 관한 차수라 하자. 그러면 $\kappa$가 체일 때의 모든
준동형 $\varphi : A \to \kappa$에 대하여 $P^\varphi$가 $x_n$에
관한 차수 $d$의 일계수 다항식이고 $x_i$에 관한 차수는 $\leq d_i$임에
유의하자. 따라서 $P^\varphi$가 가약이면
$$
P^\varphi = Q_1 Q_2
$$
로 쓸 수 있다.
% P05-EN-CH37-0124
여기서 $Q_1, Q_2$는 $x_n$에 관한 차수가 $e_1, e_2 \geq 0$인
일계수 다항식이고, $e_1 + e_2 = d$이며,
$i = 1, \ldots, n - 1$에 대하여 $x_i$에 관한 차수는 $\leq d_i$이다.
다시 말해 다음과 같이 쓸 수 있다.
\begin{equation}
\label{more-morphisms-equation-factors}
Q_j = x_n^{e_j} + \sum\nolimits_{0 \leq l < e_j}
\left( \sum\nolimits_{L \in \mathcal{L}} a_{j, l, L} x^L \right) x_n^l
\end{equation}
여기서 합은 $0 \leq l_i \leq d_i$를 만족하는
$L = (l_1, \ldots, l_{n - 1})$ 형태의 다중지표 $L$들의 집합
$\mathcal{L}$ 위에서 취한다.
$e_1 + e_2 = d$를 만족하는 임의의 $e_1, e_2 \geq 0$에 대하여
다음 $A$-대수를 생각하자.
$$
B_{e_1, e_2} =
A[\{a_{1, l, L}\}_{0 \leq l < e_1, L \in \mathcal{L}},
\{a_{2, l, L}\}_{0 \leq l < e_2, L \in \mathcal{L}}]/(\text{관계식들})
$$
% P05-EN-CH37-0129
여기서 $(\text{관계식들})$은 다음 다항식의 계수들로 생성되는 아이디얼이다.
$$
P - Q_1Q_2 \in
A[\{a_{1, l, L}\}_{0 \leq l < e_1, L \in \mathcal{L}},
\{a_{2, l, L}\}_{0 \leq l < e_2, L \in \mathcal{L}}][x_1, \ldots, x_n]
$$
$Q_1$과 $Q_2$는 (\ref{more-morphisms-equation-factors})에서와 같이 정의한다.
이제 $P$가 $\overline{K}$ 위에서 기약이라는 가정은 $A$-대수
준동형 $B_{e_1, e_2} \to \overline{K}$가 존재하지 않음을 함의한다.
Hilbert 영점 정리, 즉 Algebra, Theorem
\ref{algebra-theorem-nullstellensatz}에 의해 이는
$B_{e_1, e_2} \otimes_A K = 0$이라는 뜻이다.
$B_{e_1, e_2}$는 유한 생성 $A$-대수이므로,
$(B_{e_1, e_2})_{f_{e_1, e_2}} = 0$을 만족하는
$f_{e_1, e_2} \in A$를 찾을 수 있다. 구성에 의해 이는
$\varphi : A_{f_{e_1, e_2}} \to \kappa$가 체로 가는 준동형이면
$P^\varphi$가 $Q_1$의 $x_n$에 관한 차수가 $e_1$이고 $Q_2$의
$x_n$에 관한 차수가 $e_2$인 인수분해 $P^\varphi = Q_1 Q_2$를
갖지 않는다는 뜻이다. 따라서
$f = \prod_{e_1, e_2 \geq 0, e_1 + e_2 = d} f_{e_1, e_2}$로 놓으면
증명이 끝난다.
\end{proof}

\begin{lemma}
\label{more-morphisms-lemma-geom-irreducible-generic-fibre}
$f : X \to Y$를 스킴의 사상이라 하자. 다음을 가정하자.
\begin{enumerate}
\item $Y$는 일반점 $\eta$를 가진 기약 스킴이다.
\item $X_\eta$는 기하학적으로 기약이다.
\item $f$는 유한형이다.
\end{enumerate}
그러면 $X_V \to V$의 올들이 기하학적으로 기약이 되도록 하는 비어
있지 않은 열린 부분스킴 $V \subset Y$가 존재한다.
\end{lemma}

\begin{proof}[Lemma \ref{more-morphisms-lemma-geom-irreducible-generic-fibre}의 첫 번째 증명]
보조정리의 두 가지 증명을 제시한다. 이들은 본질적으로 동등하며,
두 번째 증명은 더 자족적이지만 조금 더 길다.
Lemma \ref{more-morphisms-lemma-make-generic-fibre-geometrically-reduced}에서와 같은
도식을 택하자.
$$
\xymatrix{
X' \ar[d]_{f'} \ar[r]_{g'} & X_V \ar[r] \ar[d] & X \ar[d]^f \\
Y' \ar[r]^g & V \ar[r] & Y
}
$$
$f'$의 일반올은 $f$의 일반올의 축약임에 유의하자
(Lemma \ref{more-morphisms-lemma-reduction-generic-fibre} 참조).
따라서 이는 기하학적으로 기약이다.
사상 $f'$에 대하여 보조정리가 성립한다고 하자. 그러면 $V$를 축소한
뒤 $f'$의 모든 올이 기하학적으로 기약이다.
$X' = (Y' \times_V X_V)_{red}$이므로 $Y' \times_V X_V$의 모든 올도
기하학적으로 기약이다. 따라서 Lemma
\ref{more-morphisms-lemma-base-change-fibres-geometrically-irreducible}에 의해
$X_V \to V$의 모든 올이 기하학적으로 기약이고 증명이 끝난다.
이렇게 하여 일반올이 기하학적으로 기약일 뿐 아니라 기하학적으로
축약이며, $A$가 정역이고 $Y = \Spec(A)$라고 가정할 수 있다.

\medskip\noindent
$x \in X_\eta$를 일반점이라 하자. $X_\eta$가 기하학적으로 기약이고
축약이므로 $L = \kappa(x)$는 $K = \kappa(\eta)$의 유한 생성
확대체이며 기하학적으로 축약이고 기하학적으로 기약이다.
Varieties, Lemmas \ref{varieties-lemma-geometrically-reduced-at-point}
및 \ref{varieties-lemma-geometrically-irreducible-function-field}를
참조하라. 특히 체 확대 $L/K$는 분리이다. Algebra, Lemma
\ref{algebra-lemma-characterize-separable-field-extensions}를 참조하라.
따라서 $K$ 위에서 $L$을 생성하고 $x_1, \ldots, x_r$가 $K$ 위에서
$L$의 초월기저가 되도록 하는 $x_1, \ldots, x_{r + 1} \in L$을
찾을 수 있다. Algebra, Lemma
\ref{algebra-lemma-generating-finitely-generated-separable-field-extensions}를
참조하라. $P \in K(x_1, \ldots, x_r)[T]$를 $x_{r + 1}$의
최소다항식이라 하자. 분모를 없애면 $P$의 계수들이
$A[x_1, \ldots, x_r]$에 속한다고 가정할 수 있다.
$L$이 $K$ 위에서 기하학적으로 축약이고 기하학적으로 기약이므로,
다항식 $P$는 $\overline{K}[x_1, \ldots, x_r, T]$에서 기약임에
유의하자. 여기서 $\overline{K}$는 $K$의 대수적 폐포이다.
% P05-EN-CH37-0125
다음과 놓자.
$$
B' = A[x_1, \ldots, x_{r + 1}]/(P(x_{r + 1}))
$$
또한 $X' = \Spec(B')$로 놓자. 구성에 의해 $B'$의 분수체는
$K$의 확대체로서 $L = \kappa(x)$와 동형이다.
% P05-EN-CH37-0130
따라서 열린집합 $U \subset X$와 $U' \subset X'$ 및
$Y$-동형 $U \to U'$이 존재한다. Morphisms, Lemma
\ref{morphisms-lemma-common-open}을 참조하라.
이를 도식으로 나타내면 다음과 같다.
$$
\xymatrix{
X \ar[rd] &
U \ar[l] \ar@{=}[r] \ar[d] &
U' \ar[r] \ar[d] &
X' \ar[ld] \ar@{=}[r] & \Spec(B') \\
& Y \ar@{=}[r] & Y &
}
$$
$U_\eta \subset X_\eta$와 $U'_\eta \subset X'_\eta$는 조밀한
열린집합임에 유의하자. 따라서 Lemma
\ref{more-morphisms-lemma-nowhere-dense-generic-fibre}를 적용하여 $Y$를 축소하면
모든 $y \in Y$에 대하여 $U_y$가 $X_y$에서 조밀하고 $U'_y$가
$X'_y$에서 조밀함을 얻는다. 그러므로 $X' \to Y$에 대하여
보조정리를 증명하면 충분하고, 이는
Lemma \ref{more-morphisms-lemma-irreducible-polynomial-over-domain}의 내용이다.
\end{proof}

\begin{proof}[Lemma \ref{more-morphisms-lemma-geom-irreducible-generic-fibre}의 두 번째 증명]
$Y' \subset Y$를 $Y$의 축약이라 하고, $X' \to X$를 $X$의 축약이라
하자. $X' \to X  \to Y$는 $Y'$을 통해 분해됨에 유의하자.
Schemes, Lemma \ref{schemes-lemma-map-into-reduction}을 참조하라.
Morphisms, Lemma \ref{morphisms-lemma-reduction-universal-homeomorphism}에
의해 $Y' \to Y$와 $X' \to X$가 보편 위상동형이므로,
$X' \to Y'$에 대하여 보조정리를 증명하면 충분하다.
따라서 $X$와 $Y$가 축약이라고 가정할 수 있다. 특히 $Y$는 정수적이다.
Properties, Lemma \ref{properties-lemma-characterize-integral}을 참조하라.
따라서 Morphisms, Proposition
\ref{morphisms-proposition-generic-flatness}에 의해 $X_V \to V$가
평탄하고 유한 표시인 비어 있지 않은 아핀 열린집합 $V \subset Y$가
존재한다. $Y$를 $V$로 바꾸면 (1), (2), (3)에 더하여 $Y$가 정수적
아핀 스킴이고, $X$가 축약이며, $f$가 평탄하고 유한 표시라고
가정할 수 있다. 특히 $f$는 보편 열린 사상이다. Morphisms, Lemma
\ref{morphisms-lemma-fppf-open}을 참조하라.

\medskip\noindent
비어 있지 않은 아핀 열린집합 $U \subset X$를 택하자.
그러면 $U \to Y$는 평탄하고 유한 표시이며 일반올이 기하학적으로
기약이다. 여집합 $X_\eta \setminus U_\eta$는 어디에서도 조밀하지 않다.
따라서 $Y$를 축소하면 모든 $y \in Y$에 대하여 $U_y \subset X_y$가
열려 있고 조밀하다고 가정할 수 있다. Lemma
\ref{more-morphisms-lemma-nowhere-dense-generic-fibre}를 참조하라.
그러므로 $X$를 $U$로 바꾸어도 되고, $Y$가 정수적 아핀 스킴이며
$X$가 축약 아핀 스킴이고 $Y$ 위에서 평탄하고 유한 표시이며
일반올 $X_\eta$가 기하학적으로 기약인 경우로 환원된다.

\medskip\noindent
$X = \Spec(B)$ 및 $Y = \Spec(A)$라 쓰자. 그러면 $A$는 정역이고,
$B$는 축약이며, $A \to B$는 평탄한 유한 표시 사상이고, $B_K$는
$A$의 분수체 $K$ 위에서 기하학적으로 기약이다.
특히 $B_K$는 정역이다. $L$을 $B_K$의 분수체라 하자.
$B$가 유한 표시 $A$-대수이므로 $L$은 $K$의 유한 생성 확대체임에
유의하자. $(L \otimes_K K')_{red}$가 분리 생성 확대체가 되도록 하는
유한 순수 비분리 확대 $K'/K$를 택하자. Algebra, Lemma
\ref{algebra-lemma-make-separable}를 참조하라.
% P05-EN-CH37-0126
$K$ 위에서 $K'$을 생성하고, 어떤 소수의 거듭제곱 $q_i$에 대하여
$x_i^{q_i} \in A$를 만족하는 $x_1, \ldots, x_n \in K'$을 택하자
($x_i$의 존재 증명은 생략한다). $A'$을 $x_1, \ldots, x_n$으로
생성되는 $K'$의 $A$-부분대수라 하자. 그러면 $A'$은 분수체가
$K'$인 유한 $A$-부분대수 $A' \subset K'$이다.
$\Spec(A') \to \Spec(A)$는 보편 위상동형임에 유의하자. Algebra,
Lemma \ref{algebra-lemma-p-ring-map}을 참조하라.
따라서 $\Spec(A')$로 밑변환한 뒤 결과를 증명하면 충분하다.
$A$를 $A'$으로, $B$를 $(B \otimes_A A')_{red}$로 바꾸어 $L$이
$K$의 분리 생성 확대체인 상황에 도달하고자 한다.
물론 $(B \otimes_A A')_{red}$가 더 이상 $A'$ 위에서 평탄하지 않거나
유한 표시가 아닐 수도 있다. 그러나 적당한 $f \in A'$에 대하여
$A'$을 $A'_f$로 바꾸면 이를 해결할 수 있다. Algebra, Lemma
\ref{algebra-lemma-generic-flatness}를 참조하라.

\medskip\noindent
이제 $A$는 정역이고, $B$는 축약이며, $A \to B$는 평탄하고 유한
표시이고, $B_K$는 그 분수체 $L$이 $A$의 분수체 $K$의 분리 생성
확대체가 되는 정역이다. Algebra, Lemma
\ref{algebra-lemma-generating-finitely-generated-separable-field-extensions}에
의해 $L = K(x_1, \ldots, x_{r + 1})$이라 쓸 수 있다.
여기서 $x_1, \ldots, x_r$는 $K$ 위에서 대수적으로 독립이고
$x_{r + 1}$은 $K(x_1, \ldots, x_r)$ 위에서 분리이다.
분모를 없애면 $K(x_1, \ldots, x_r)$ 위에서 $x_{r + 1}$의
최소다항식 $P \in K(x_1, \ldots, x_r)[T]$의 계수가
$A[x_1, \ldots, x_r]$에 속한다고 가정할 수 있다.
$L/K$가 분리이고 $L$이 $K$ 위에서 기하학적으로 기약이므로, 다항식
$P$는 $K$의 대수적 폐포 $\overline{K}$ 위에서 기약임에 유의하자.
% P05-EN-CH37-0125
다음과 놓자.
$$
B' = A[x_1, \ldots, x_{r + 1}]/(P(x_{r + 1})).
$$
구성에 의해 $B$와 $B'$의 분수체는 $K$의 확대체로서 동형이다.
따라서 적당한 $h \in B$와 $h' \in B'$에 대하여 $A$-대수의 동형
$B_h \cong B'_{h'}$가 존재한다. Morphisms, Lemma
\ref{morphisms-lemma-common-open}을 참조하라. 다시 말해 $X$와
$X' = \Spec(B')$는 공통 아핀 열린집합 $U$를 가진다.
이를 도식으로 나타내면 다음과 같다.
$$
\xymatrix{
X = \Spec(B) \ar[rd] &
U \ar[l] \ar[r] \ar[d] &
\Spec(B') = X' \ar[ld] \\
& Y = \Spec(A) &
}
$$
$Y$를 한 번 더 축소하면(Lemma
\ref{more-morphisms-lemma-nowhere-dense-generic-fibre}를 $X$ 안의 $Z = X \setminus U$와
$X'$ 안의 $Z' = X' \setminus U$에 적용하여), 모든 $y \in Y$에 대하여
$U_y$가 $X_y$에서 조밀하고 $U_y$가 $X'_y$에서 조밀함을 얻는다.
따라서 $X' \to Y$에 대하여 보조정리를 증명하면 충분하고, 이는
Lemma \ref{more-morphisms-lemma-irreducible-polynomial-over-domain}의 내용이다.
\end{proof}

\begin{lemma}
\label{more-morphisms-lemma-nr-geom-irreducible-components-good}
$f : X \to Y$를 스킴의 사상이라 하자.
$n_{X/Y}$를 Lemma
\ref{more-morphisms-lemma-base-change-fibres-nr-geometrically-irreducible-components}에서
도입한, $f$의 올들의 기하학적 기약 성분 수를 세는 $Y$ 위의 함수라 하자.
% P05-EN-CH37-0128
$f$가 유한형이라고 가정하자.
$y \in Y$를 점이라 하자. 그러면 $n_{X/Y}|_V$가 상수가 되도록 하는
비어 있지 않은 열린집합 $V \subset \overline{\{y\}}$가 존재한다.
\end{lemma}

\begin{proof}
$Z$를 $\overline{\{y\}}$에 유도된 축약 스킴 구조를 준 것이라 하자.
$f_Z : X_Z \to Z$를 $f$의 밑변환이라 하자.
$f_Z$와 $Z$의 일반점에 대하여 보조정리를 증명하면 충분함은 명백하다.
따라서 $Y$가 정수적 스킴이라고 가정할 수 있다. Properties, Lemma
\ref{properties-lemma-characterize-integral}을 참조하라.
이 경우 우리의 목표는 $n_{X/Y}|_V$가 상수가 되도록 하는 비어 있지
않은 열린집합 $V \subset Y$를 만드는 것이다.

\medskip\noindent
Lemma \ref{more-morphisms-lemma-make-components-generic-fibre-geometrically-irreducible}를
$f : X \to Y$에 적용하여 $g : Y' \to V \subset Y$를 얻는다.
$g : Y' \to V$는 전사 유한 에탈 사상이므로 특히 열린 사상이다
(Morphisms, Lemma \ref{morphisms-lemma-etale-open} 참조).
% P05-EN-CH37-0127
따라서 $n_{X'/Y'}|_{V'}$가 상수가 되도록 하는 열린집합
$V' \subset Y'$가 존재함을 증명하면 충분하다. Lemma
\ref{more-morphisms-lemma-base-change-fibres-nr-geometrically-irreducible-components}를
참조하라. 이로써 일반올 $X_\eta$의 모든 기약 성분이 $\kappa(\eta)$
위에서 기하학적으로 기약이라고 가정할 수 있다.

\medskip\noindent
이제
$X_\eta = X_{1, \eta} \bigcup \ldots \bigcup X_{n, \eta}$
를 일반올의 (기하학적) 기약 성분 분해라 하자.
특히 $n_{X/Y}(\eta) = n$이다.
$X_i$를 $X$에서 $X_{i, \eta}$의 폐포라 하자.
$Y$를 축소하면 $X = \bigcup X_i$라고 가정할 수 있다. Lemma
\ref{more-morphisms-lemma-cover-generic-fibre-neighbourhood}를 참조하라.
$Y$를 더 축소하면 $f$의 각 올의 기약 성분 수는 적어도 $n$개이다.
Lemma \ref{more-morphisms-lemma-irreducible-components-in-neighbourhood}를 참조하라.
따라서 모든 $y \in Y$에 대하여 $n_{X/Y}(y) \geq n$이다.
$Y$를 더 축소하면 각 $i$와 모든 $y \in Y$에 대하여 $X_{i, y}$가
기하학적으로 기약임을 얻는다. Lemma
\ref{more-morphisms-lemma-geom-irreducible-generic-fibre}를 참조하라.
$X_y = \bigcup X_{i, y}$이므로 $n_{X/Y}(y) \leq n$이고
증명이 끝난다.
\end{proof}

\begin{lemma}
\label{more-morphisms-lemma-nr-geom-irreducible-components-constructible}
$f : X \to Y$를 스킴의 사상이라 하자.
$n_{X/Y}$를 Lemma
\ref{more-morphisms-lemma-base-change-fibres-nr-geometrically-irreducible-components}에서
도입한, $f$의 올들의 기하학적 기약 성분 수를 세는 $Y$ 위의 함수라 하자.
% P05-EN-CH37-0128
$f$가 유한 표시라고 가정하자. 그러면 $n_{X/Y}$의 등위집합
$$
E_n = \{y \in Y \mid n_{X/Y}(y) = n\}
$$
은 $Y$에서 국소 구성가능이다.
\end{lemma}

\begin{proof}
$n$을 고정하자. $y \in Y$라 하자. $E_n \cap V$가 $V$에서 구성가능이
되도록 하는 $Y$에서의 $y$의 열린 근방 $V$가 존재함을 보여야 한다.
따라서 $Y$가 아핀이라고 가정할 수 있다. $Y = \Spec(A)$라 쓰고
유한형 $\mathbf{Z}$-대수들의 유향 극한 $A = \colim A_i$로 쓰자.
Limits, Lemma \ref{limits-lemma-descend-finite-presentation}에 의해
$Y$로 밑변환하면 $f$를 복원하는 유한 표시 사상
$f_i : X_i \to \Spec(A_i)$와 지표 $i$를 찾을 수 있다.
Lemma \ref{more-morphisms-lemma-base-change-fibres-nr-geometrically-irreducible-components}에
의해 $f_i$에 대하여 보조정리를 증명하면 충분하다.
따라서 $Y$가 뇌터환의 스펙트럼인 경우로 환원된다.

\medskip\noindent
$Y$가 뇌터 스킴인 경우 $E_n$이 구성가능임을 증명하기 위해 Topology,
Lemma \ref{topology-lemma-characterize-constructible-Noetherian}의
판정법을 사용하겠다. 이를 위해 $Z \subset Y$를 기약 닫힌
부분스킴이라 하자. $E_n \cap Z$가 비어 있지 않은 열린 부분집합을
포함하거나 $Z$에서 조밀하지 않음을 보여야 한다. $\xi \in Z$를
일반점이라 하자. 그러면 Lemma
\ref{more-morphisms-lemma-nr-geom-irreducible-components-good}에 의해 $n_{X/Y}$는
$Z$에서의 $\xi$의 근방에서 상수이다. 이는 원하는 결론을 명백히 함의한다.
\end{proof}













\section{올의 연결 성분}
\label{more-morphisms-section-connected}

\begin{lemma}
\label{more-morphisms-lemma-connected-components-in-neighbourhood}
$f : X \to Y$를 스킴의 사상이라 하자. $Y$가 일반점 $\eta$를 가진
기약 스킴이고 $f$가 유한형이라고 가정하자. $X_\eta$의 연결 성분이
$n$개이면, 모든 $y \in V$에 대하여 올 $X_y$의 연결 성분이 적어도
$n$개가 되도록 하는 비어 있지 않은 열린집합 $V \subset Y$가 존재한다.
\end{lemma}

\begin{proof}
문제가 순전히 위상적이므로 $X$와 $Y$를 각각 그 축약으로 바꾸어도 된다.
특히 그러면 $Y$가 정수적이다. Properties, Lemma
\ref{properties-lemma-characterize-integral}을 참조하라.
$X_\eta = X_{1, \eta} \cup \ldots \cup X_{n, \eta}$를 $X_\eta$의
연결 성분 분해라 하자. $X_i \subset X$를 일반올이 $X_{i, \eta}$인
축약 닫힌 부분스킴이라 하자. $Z_{i, j} = X_i \cap X_j$는 $X$의
닫힌 부분집합이고 그 일반올 $Z_{i, j, \eta}$는 비어 있음에 유의하자.
따라서 $Y$를 축소하면 $Z_{i, j} = \emptyset$이라고 가정할 수 있다.
Lemma \ref{more-morphisms-lemma-empty-generic-fibre}를 참조하라.
$Y$를 더 축소하면 $y \in Y$에 대하여 $X_y = \bigcup X_{i, y}$라고
가정할 수 있다. Lemma \ref{more-morphisms-lemma-cover-generic-fibre-neighbourhood}를
참조하라. 더욱이 $Y$를 축소하면 각 $X_i \to Y$가 평탄하고 유한
표시라고 가정할 수 있다. Morphisms, Proposition
\ref{morphisms-proposition-generic-flatness}를 참조하라.
사상 $X_i \to Y$는 열린 사상이다. Morphisms, Lemma
\ref{morphisms-lemma-fppf-open}을 참조하라.
따라서 각 $i$에 대하여 $f(X_i)$에 포함되는 $\eta$의 열린 근방 $V$가
존재한다. 각 $y \in V$에 대하여 스킴 $X_{i, y}$들은 $X_y$의
비어 있지 않은 닫힌 부분집합이고, $X_y = \bigcup X_{i, y}$이며,
교집합 $Z_{i, j, y} = X_{i, y} \cap X_{j, y}$는 비어 있다!
이는 $X_y$의 연결 성분이 적어도 $n$개임을 명백히 함의한다.
\end{proof}

\begin{lemma}
\label{more-morphisms-lemma-base-change-fibres-geometrically-connected}
$f : X \to Y$를 스킴의 사상이라 하자.
$g : Y' \to Y$를 임의의 사상이라 하고,
$f' : X' \to Y'$을 $f$의 밑변환이라 하자. 그러면
\begin{align*}
\{y' \in Y' \mid X'_{y'}\text{는 기하학적으로 연결이다}\} \\
= g^{-1}(\{y \in Y \mid X_y\text{는 기하학적으로 연결이다}\}).
\end{align*}
\end{lemma}

\begin{proof}
이는 상이 $y \in Y$인 $y' \in Y'$에 대하여 올
$X'_{y'} = X_y \times_y y'$가 $\kappa(y')$ 위에서 기하학적으로
연결일 필요충분조건이 $X_y$가 $\kappa(y)$ 위에서 기하학적으로
연결인 것이라는 명제로 귀착된다. 이는 Varieties,
Lemma \ref{varieties-lemma-geometrically-connected-check-after-extension}에서
따른다.
\end{proof}

\begin{lemma}
\label{more-morphisms-lemma-base-change-fibres-nr-geometrically-connected-components}
$f : X \to Y$를 스킴의 사상이라 하자. 함수
$$
n_{X/Y} : Y \to \{0, 1, 2, 3, \ldots, \infty\}
$$
를 $y \in Y$에 $(X_y)_K$의 연결 성분의 수를 대응시키는 함수라 하자.
여기서 $K$는 $\kappa(y)$의 분리적으로 닫힌 확대체이다.
이는 잘 정의되며, $g : Y' \to Y$가 사상이면
$$
n_{X'/Y'} = n_{X/Y} \circ g
$$
이다. 여기서 $X' \to Y'$은 $f$의 밑변환이다.
\end{lemma}

\begin{proof}
$y' \in Y'$의 상이 $y \in Y$라고 하자.
$K \supset \kappa(y)$와 $K' \supset \kappa(y')$가 분리적으로 닫힌
확대체라고 하자. 그러면 체들의 가환 도식
$$
\xymatrix{
K \ar[r] & K'' & K' \ar[l] \\
\kappa(y) \ar[u] \ar[rr] & & \kappa(y') \ar[u]
}
$$
을 택할 수 있다. 스킴의 사상
$$
\xymatrix{
(X'_{y'})_{K'} &
(X'_{y'})_{K''} = (X_y)_{K''} \ar[l] \ar[r] &
(X_y)_K
}
$$
들이 연결 성분 사이의 전단사를 유도하므로 결과가 따른다.
Varieties,
Lemma \ref{varieties-lemma-separably-closed-field-connected-components}를
참조하라.
\end{proof}

\begin{lemma}
\label{more-morphisms-lemma-geometrically-connected-generic-fibre}
$f : X \to Y$를 스킴의 사상이라 하자. 다음을 가정하자.
\begin{enumerate}
\item $Y$는 일반점 $\eta$를 가진 기약 스킴이다.
\item $X_\eta$는 기하학적으로 연결이다.
\item $f$는 유한형이다.
\end{enumerate}
그러면 $X_V \to V$의 올들이 기하학적으로 연결이 되도록 하는 비어
있지 않은 열린 부분스킴 $V \subset Y$가 존재한다.
\end{lemma}

\begin{proof}
Lemma \ref{more-morphisms-lemma-make-components-generic-fibre-geometrically-irreducible}에서와
같은 도식을 택하자.
$$
\xymatrix{
X' \ar[d]_{f'} \ar[r]_{g'} & X_V \ar[r] \ar[d] & X \ar[d]^f \\
Y' \ar[r]^g & V \ar[r] & Y
}
$$
$f'$의 일반올은 기하학적으로 연결임에 유의하자
(예를 들어 Lemma
\ref{more-morphisms-lemma-base-change-fibres-nr-geometrically-connected-components}에 의해).
사상 $f'$에 대하여 보조정리가 성립한다고 하자. 이는 $W$ 위에서
$X' \to Y'$의 모든 올이 기하학적으로 연결이 되도록 하는 비어 있지
않은 열린집합 $W \subset Y'$가 존재한다는 뜻이다.
Morphisms, Lemma \ref{morphisms-lemma-etale-open}에 의해 $g$는
열린 사상이므로, 비어 있지 않은 열린집합 $V = g(W)$의 점들에서
$f$의 모든 올은 기하학적으로 연결이다. Lemma
\ref{more-morphisms-lemma-base-change-fibres-nr-geometrically-connected-components}를
참조하라. 이렇게 하여 일반올 $X_\eta$의 기약 성분들이 기하학적으로
기약이라고 가정할 수 있다.

\medskip\noindent
$Y'$을 $Y$의 축약이라 하고 $X' = Y' \times_Y X$로 놓자.
그러면 사상 $X' \to Y'$에 대하여 보조정리를 증명하면 충분하다
(예를 들어 Lemma
\ref{more-morphisms-lemma-base-change-fibres-nr-geometrically-connected-components}를
다시 적용하면 된다). $X' \to Y'$의 일반올은 $X \to Y$의 일반올과
같으므로, $Y$가 기약이고 축약이며(즉 정수적이다. Properties,
Lemma \ref{properties-lemma-characterize-integral} 참조) 일반올
$X_\eta$의 기약 성분들은 기하학적으로 기약이라고 가정할 수 있다.

\medskip\noindent
이제
$X_\eta = X_{1, \eta} \bigcup \ldots \bigcup X_{n, \eta}$
를 일반올의 (기하학적) 기약 성분 분해라 하자.
$X_i$를 $X$에서 $X_{i, \eta}$의 폐포라 하자.
$Y$를 축소하면 $X = \bigcup X_i$라고 가정할 수 있다. Lemma
\ref{more-morphisms-lemma-cover-generic-fibre-neighbourhood}를 참조하라.
$Z_{i, j} = X_i \cap X_j$라 하자. 다음과 같이 놓자.
$$
\{1, \ldots, n\} \times \{1, \ldots, n\} = I \amalg J
$$
여기서 $Z_{i, j, \eta} = \emptyset$이면 $(i, j) \in I$이고,
$Z_{i, j, \eta} \not = \emptyset$이면 $(i, j) \in J$이다.
$Y$를 축소하면 모든 $(i, j) \in I$에 대하여
$Z_{i, j} = \emptyset$이라고 가정할 수 있다. Lemma
\ref{more-morphisms-lemma-empty-generic-fibre}를 참조하라.
$Y$를 축소하면 각 $i$와 모든 $y \in Y$에 대하여 $X_{i, y}$가
기하학적으로 기약임을 얻는다. Lemma
\ref{more-morphisms-lemma-geom-irreducible-generic-fibre}를 참조하라.
$Y$를 더 축소하면 모든 $(i, j) \in J$에 대하여 각
$Z_{i, j} \to Y$가 평탄하고 유한 표시인 상황을 얻는다. Morphisms,
Proposition \ref{morphisms-proposition-generic-flatness}를 참조하라.
이는 $f(Z_{i, j}) \subset Y$가 열려 있다는 뜻이다. Morphisms, Lemma
\ref{morphisms-lemma-fppf-open}을 참조하라. 다음 집합이 조건을
만족한다고 주장한다.
$$
V  = \bigcap\nolimits_{(i, j) \in J} f(Z_{i, j})
$$
즉, 각 $y \in V$에 대하여 $X_y$가 기하학적으로 연결이라고 주장한다.
실제로 $X_\eta$가 연결이라는 사실은 $J$의 쌍들로 생성되는 동치관계가
오직 하나의 동치류를 가짐을 함의한다. 이제 $y \in V$이고
$K \supset \kappa(y)$가 분리적으로 닫힌 확대체이면, $(X_y)_K$의 기약
성분들은 올 $(X_{i, y})_K$이다. 더욱이 구성과 $y \in V$에 의해,
$(X_{i, y})_K$가 $(X_{j, y})_K$와 만날 필요충분조건은
$(i, j) \in J$인 것이다. 따라서 동치류에 관한 위 관찰로부터
$(X_y)_K$가 연결임을 얻고 증명이 끝난다.
\end{proof}

\begin{lemma}
\label{more-morphisms-lemma-nr-geom-connected-components-good}
$f : X \to Y$를 스킴의 사상이라 하자.
$n_{X/Y}$를 Lemma
\ref{more-morphisms-lemma-base-change-fibres-nr-geometrically-connected-components}에서
도입한, $f$의 올들의 기하학적 연결 성분 수를 세는 $Y$ 위의 함수라 하자.
% P05-EN-CH37-0133
$f$가 유한형이라고 가정하자.
$y \in Y$를 점이라 하자. 그러면 $n_{X/Y}|_V$가 상수가 되도록 하는
비어 있지 않은 열린집합 $V \subset \overline{\{y\}}$가 존재한다.
\end{lemma}

\begin{proof}
$Z$를 $\overline{\{y\}}$에 유도된 축약 스킴 구조를 준 것이라 하자.
$f_Z : X_Z \to Z$를 $f$의 밑변환이라 하자.
$f_Z$와 $Z$의 일반점에 대하여 보조정리를 증명하면 충분함은 명백하다.
따라서 $Y$가 정수적 스킴이라고 가정할 수 있다. Properties, Lemma
\ref{properties-lemma-characterize-integral}을 참조하라.
이 경우 우리의 목표는 $n_{X/Y}|_V$가 상수가 되도록 하는 비어 있지
않은 열린집합 $V \subset Y$를 만드는 것이다.

\medskip\noindent
Lemma \ref{more-morphisms-lemma-make-components-generic-fibre-geometrically-irreducible}를
$f : X \to Y$에 적용하여 $g : Y' \to V \subset Y$를 얻는다.
$g : Y' \to V$는 전사 유한 에탈 사상이므로 특히 열린 사상이다
(Morphisms, Lemma \ref{morphisms-lemma-etale-open} 참조).
% P05-EN-CH37-0131
따라서 $n_{X'/Y'}|_{V'}$가 상수가 되도록 하는 열린집합
$V' \subset Y'$가 존재함을 증명하면 충분하다.
% P05-EN-CH37-0132
Lemma \ref{more-morphisms-lemma-base-change-fibres-nr-geometrically-irreducible-components}를
참조하라. 이로써 일반올 $X_\eta$의 모든 기약 성분이 $\kappa(\eta)$
위에서 기하학적으로 기약이라고 가정할 수 있다.
Varieties, Lemma
\ref{varieties-lemma-irreducible-components-geometrically-irreducible}에
의해 $X_\eta$의 연결 성분들도 기하학적으로 연결이다.

\medskip\noindent
이제
$X_\eta = X_{1, \eta} \bigcup \ldots \bigcup X_{n, \eta}$
를 일반올의 (기하학적) 연결 성분 분해라 하자.
특히 $n_{X/Y}(\eta) = n$이다.
$X_i$를 $X$에서 $X_{i, \eta}$의 폐포라 하자.
$Y$를 축소하면 $X = \bigcup X_i$라고 가정할 수 있다. Lemma
\ref{more-morphisms-lemma-cover-generic-fibre-neighbourhood}를 참조하라.
$Y$를 더 축소하면 $f$의 각 올의 연결 성분 수는 적어도 $n$개이다.
Lemma \ref{more-morphisms-lemma-connected-components-in-neighbourhood}를 참조하라.
따라서 모든 $y \in Y$에 대하여 $n_{X/Y}(y) \geq n$이다.
$Y$를 더 축소하면 각 $i$와 모든 $y \in Y$에 대하여 $X_{i, y}$가
기하학적으로 연결임을 얻는다. Lemma
\ref{more-morphisms-lemma-geometrically-connected-generic-fibre}를 참조하라.
$X_y = \bigcup X_{i, y}$이므로 $n_{X/Y}(y) \leq n$이고
증명이 끝난다.
\end{proof}

\begin{lemma}
\label{more-morphisms-lemma-nr-geom-connected-components-constructible}
$f : X \to Y$를 스킴의 사상이라 하자.
% P05-EN-CH37-0134
$n_{X/Y}$를 Lemma
\ref{more-morphisms-lemma-base-change-fibres-nr-geometrically-connected-components}에서
도입한, $f$의 올들의 기하학적 연결 성분 수를 세는 $Y$ 위의 함수라 하자.
% P05-EN-CH37-0133
$f$가 유한 표시라고 가정하자. 그러면 $n_{X/Y}$의 등위집합
$$
E_n = \{y \in Y \mid n_{X/Y}(y) = n\}
$$
은 $Y$에서 국소 구성가능이다.
\end{lemma}

\begin{proof}
$n$을 고정하자. $y \in Y$라 하자. $E_n \cap V$가 $V$에서 구성가능이
되도록 하는 $Y$에서의 $y$의 열린 근방 $V$가 존재함을 보여야 한다.
따라서 $Y$가 아핀이라고 가정할 수 있다. $Y = \Spec(A)$라 쓰고
유한형 $\mathbf{Z}$-대수들의 유향 극한 $A = \colim A_i$로 쓰자.
Limits, Lemma \ref{limits-lemma-descend-finite-presentation}에 의해
$Y$로 밑변환하면 $f$를 복원하는 유한 표시 사상
$f_i : X_i \to \Spec(A_i)$와 지표 $i$를 찾을 수 있다.
Lemma \ref{more-morphisms-lemma-base-change-fibres-nr-geometrically-connected-components}에
의해 $f_i$에 대하여 보조정리를 증명하면 충분하다.
따라서 $Y$가 뇌터환의 스펙트럼인 경우로 환원된다.

\medskip\noindent
$Y$가 뇌터 스킴인 경우 $E_n$이 구성가능임을 증명하기 위해 Topology,
Lemma \ref{topology-lemma-characterize-constructible-Noetherian}의
판정법을 사용하겠다. 이를 위해 $Z \subset Y$를 기약 닫힌
부분스킴이라 하자. $E_n \cap Z$가 비어 있지 않은 열린 부분집합을
포함하거나 $Z$에서 조밀하지 않음을 보여야 한다. $\xi \in Z$를
일반점이라 하자. 그러면 Lemma
\ref{more-morphisms-lemma-nr-geom-connected-components-good}에 의해 $n_{X/Y}$는
$Z$에서의 $\xi$의 근방에서 상수이다. 이는 원하는 결론을 명백히 함의한다.
\end{proof}

\begin{lemma}
\label{more-morphisms-lemma-connected-flat-over-dvr}
\begin{slogan}
연결이 아닌 스킴의 평탄 퇴화는 연결이 아니거나 축약이 아니다.
\end{slogan}
$f : X \to S$를 스킴의 사상이라 하자. 다음을 가정하자.
\begin{enumerate}
\item $S$는 이산 값매김환의 스펙트럼이다.
\item $f$는 평탄이다.
\item $X$는 연결이다.
\item 닫힌올 $X_s$는 축약이다.
\end{enumerate}
그러면 일반올 $X_\eta$는 연결이다.
\end{lemma}

\begin{proof}
$S = \Spec(R)$라 쓰고 $\pi \in R$을 균등화원이라 하자.
모순을 얻기 위해 $X_\eta$가 연결이 아니라고 가정하자.
이는 비자명한 멱등원 $e \in \Gamma(X_\eta, \mathcal{O}_{X_\eta})$가
존재한다는 뜻이다. $U = \Spec(A)$를 $X$의 임의의 아핀 열린집합이라
하자. $A$가 $R$ 위에서 평탄하므로 $\pi$는 $A$ 위의 영인자가 아닌
원소임에 유의하자. 예를 들어 More on Algebra, Lemma
\ref{more-algebra-lemma-flat-torsion-free}를 참조하라.
그러면 $e|_{U_\eta}$는 원소 $e \in A[1/\pi]$에 대응한다.
$e = z/\pi^n$을 만족하고 $n \geq 0$이 최소가 되도록 하는
원소 $z \in A$를 택하자. $z^2 = \pi^nz$임에 유의하자.
이는 $n > 0$이면 $z \bmod \pi A$가 멱영이라는 뜻이다.
가정에 의해 $A/\pi A$는 축약이므로 $n$의 최소성에 의해 $n = 0$이다.
따라서 $e \in A$라고 결론짓는다! 다시 말해
$e \in \Gamma(X, \mathcal{O}_X)$이다. $X$가 연결이므로 $e$는 자명한
멱등원이고, 이는 모순이다.
\end{proof}







\section{단면과 만나는 연결 성분}
\label{more-morphisms-section-connected-components}

\noindent
이 절의 결과들은 특히 군 스킴 $G \to S$와 그 항등 단면
$e : S \to G$에 적용할 수 있다.

\begin{situation}
\label{more-morphisms-situation-connected-along-section}
% P05-EN-CH37-0135
여기서 $f : X \to Y$를 스킴의 사상이라 하고,
$s : Y \to X$를 $f$의 단면이라 하자.
% P05-EN-CH37-0136
각 $y \in Y$에 대하여 $s(y)$를 포함하는 $X_y$의 연결 성분을
$X^0_y$로 표기한다. 끝으로 $X^0 = \bigcup_{y \in Y} X^0_y$로 놓는다.
\end{situation}

\begin{lemma}
\label{more-morphisms-lemma-base-change-connected-along-section}
$f : X \to Y$, $s : Y \to X$가 Situation
\ref{more-morphisms-situation-connected-along-section}에서와 같다고 하자.
임의의 사상 $g : Y' \to Y$에 대하여 밑변환 도식
$$
\xymatrix{
X' \ar[r]_{g'} \ar[d]^{f'} & X \ar[d]_f \\
Y' \ar@/^1pc/[u]^{s'} \ar[r]^g & Y \ar@/_1pc/[u]_s
}
$$
을 생각하면 $(X')^0 \subset X'$을 얻는다.
그러면 $(X')^0 = (g')^{-1}(X^0)$이다.
\end{lemma}

\begin{proof}
$y' \in Y'$의 상이 $y \in Y$라고 하자. $X^0_y$를 $X_y$의 닫힌
부분스킴으로 생각할 수 있다. 예를 들어 Morphisms,
Definition \ref{morphisms-definition-scheme-structure-connected-component}을
참조하라. $s(y) \in X^0_y$이므로 Varieties, Lemma
\ref{varieties-lemma-geometrically-connected-if-connected-and-point}에 의해
$X_y^0$는 $\kappa(y)$ 위에서 기하학적으로 연결인 스킴이다.
따라서 $X_y^0 \times_y y' \to X'_{y'}$는 $s'(y')$를 포함하는
연결 닫힌 부분스킴이다. 그러므로
$X_y^0 \times_y y' \subset (X'_{y'})^0$이다.
반대 포함관계 $X_y^0 \times_y y' \supset (X'_{y'})^0$는 명백하다.
$(X'_{y'})^0$의 $X_y$에서의 상은 $s(y)$를 포함하는 $X_y$의
연결 부분집합이기 때문이다.
\end{proof}

\begin{lemma}
\label{more-morphisms-lemma-connected-along-section-good}
$f : X \to Y$, $s : Y \to X$가 Situation
\ref{more-morphisms-situation-connected-along-section}에서와 같다고 하자.
% P05-EN-CH37-0137
$f$가 유한형이라고 가정하자. $y \in Y$를 점이라 하자.
그러면 밑변환 $X_V$에서 $X^0$의 역상이 $X_V$에서 열려 있고 닫혀
있도록 하는 비어 있지 않은 열린집합 $V \subset \overline{\{y\}}$가
존재한다.
\end{lemma}

\begin{proof}
$Z \subset Y$를 $\overline{\{y\}}$에 유도된 축약 닫힌 부분스킴
구조를 준 것이라 하자. $f_Z : X_Z \to Z$와 $s_Z : Z \to X_Z$를
각각 $f$와 $s$의 밑변환이라 하자.
Lemma \ref{more-morphisms-lemma-base-change-connected-along-section}에 의해
$(X_Z)^0 = (X^0)_Z$이다.
% P05-EN-CH37-0138
따라서 사상 $X_Z \to Z$와 $Z$의 일반점으로 가는 점 $x \in X_Z$에
대하여 보조정리를 증명하면 충분하다. 다시 말해 $Y$가 일반점
$\eta$를 가진 정수적 스킴인 경우로 문제를 환원하였다
(Properties, Lemma \ref{properties-lemma-characterize-integral} 참조).
우리의 목표는 $Y$를 축소하면 부분집합 $X^0$가 $X$의 열리고 닫힌
부분집합이 됨을 보이는 것이다.

\medskip\noindent
스킴 $X_\eta$는 체 위에서 유한형이므로 뇌터임에 유의하자.
따라서 그 연결 성분들은 닫혀 있을 뿐 아니라 열려 있다.
그러므로 $X_\eta$의 어떤 열리고 닫힌 부분집합 $T_\eta$에 대하여
$X_\eta = X_\eta^0 \amalg T_\eta$라 쓸 수 있다.
다음으로 $T \subset X$를 $T_\eta$의 폐포라 하고,
$X^{00} \subset X$를 $X_\eta^0$의 폐포라 하자.
$T_\eta$와 $X^0_\eta$는 각각 $T$와 $X^{00}$의 일반올임에 유의하자.
Lemma \ref{more-morphisms-lemma-cover-generic-fibre-neighbourhood} 앞의 논의를
참조하라. 더욱이 그 보조정리에 의해 $Y$를 축소하면 집합론적으로
$X = X^{00} \cup T$라고 가정할 수 있다.
$(T \cap X^{00})_\eta = T_\eta \cap X^0_\eta = \emptyset$임에 유의하자.
따라서 $Y$를 축소하면 $T \cap X^{00} = \emptyset$이라고 가정할 수
있다. Lemma \ref{more-morphisms-lemma-empty-generic-fibre}를 참조하라.
특히 $X^{00}$는 $X$에서 열려 있다. $X^0_\eta$는 연결이고 유리점
$s(\eta)$를 가지므로 기하학적으로 연결임에 유의하자. Varieties,
Lemma \ref{varieties-lemma-geometrically-connected-if-connected-and-point}를
참조하라. 따라서 $Y$를 축소하면 $X^{00} \to Y$의 모든 올이
기하학적으로 연결이라고 가정할 수 있다. Lemma
\ref{more-morphisms-lemma-geometrically-connected-generic-fibre}를 참조하라.
이제 올 $X^{00}_y$는 $s(y)$를 포함하는 $X_y$의 열리고 닫힌 연결
부분집합이다. 따라서 $X^0 = X^{00}$이고 증명이 끝난다.
\end{proof}

\begin{lemma}
\label{more-morphisms-lemma-connected-along-section-locally-constructible}
$f : X \to Y$, $s : Y \to X$가 Situation
\ref{more-morphisms-situation-connected-along-section}에서와 같다고 하자.
$f$가 유한 표시이면 $X^0$는 $X$에서 국소 구성가능이다.
\end{lemma}

\begin{proof}
$x \in X$라 하자. $X^0 \cap U$가 $U$에서 구성가능이 되도록 하는
$x$의 열린 근방 $U$가 존재함을 보여야 한다.
따라서 $Y$가 아핀인 경우로 환원된다.
$Y = \Spec(A)$라 쓰고 유한형 $\mathbf{Z}$-대수들의 유향 극한
$A = \colim A_i$로 쓰자. Limits, Lemma
\ref{limits-lemma-descend-finite-presentation}에 의해, 지표 $i$와
단면 $s_i : \Spec(A_i) \to X_i$를 갖춘 유한 표시 사상
$f_i : X_i \to \Spec(A_i)$를 찾을 수 있으며, 이를 $Y$로 밑변환하면
$f$와 그 단면 $s$를 복원한다.
Lemma \ref{more-morphisms-lemma-base-change-connected-along-section}에 의해
$f_i, s_i$에 대하여 보조정리를 증명하면 충분하다.
따라서 $Y$가 뇌터환의 스펙트럼인 경우로 환원된다.

\medskip\noindent
$Y$가 뇌터 아핀 스킴이라고 가정하자. $f$가 유한 표시, 즉 유한형이므로
$X$도 뇌터 스킴이다. Morphisms, Lemma
\ref{morphisms-lemma-finite-type-noetherian}을 참조하라.
이 경우 보조정리를 증명하려면 모든 기약 닫힌 부분집합 $Z \subset X$에
대하여 교집합 $Z \cap X^0$가 $Z$의 비어 있지 않은 열린집합을 포함하거나
$Z$에서 조밀하지 않음을 보이면 충분하다. Topology, Lemma
\ref{topology-lemma-characterize-constructible-Noetherian}을 참조하라.
$x \in Z$를 일반점이라 하고 $y = f(x)$라 하자.
Lemma \ref{more-morphisms-lemma-connected-along-section-good}에 의해
$X^0 \cap X_V$가 $X_V$에서 열려 있고 닫혀 있도록 하는 비어 있지 않은
열린 부분집합 $V \subset \overline{\{y\}}$가 존재한다.
$f(Z) \subset \overline{\{y\}}$이고 $f(x) = y \in V$이므로
$W = f^{-1}(V) \cap Z$는 $Z$의 비어 있지 않은 열린 부분집합이다.
따라서 $X^0 \cap W$는 $W$에서 열려 있고 닫혀 있다.
$W$가 기약이므로 $X^0 \cap W$는 비어 있거나 $W$와 같다.
이로써 보조정리가 증명된다.
\end{proof}

\begin{lemma}
\label{more-morphisms-lemma-connected-along-section-open-neighbourhood}
$f : X \to Y$, $s : Y \to X$가 Situation
\ref{more-morphisms-situation-connected-along-section}에서와 같다고 하자.
$y \in Y$를 점이라 하자. 다음을 가정하자.
\begin{enumerate}
\item $f$는 유한 표시이고 평탄이다.
\item 올 $X_y$는 기하학적으로 축약이다.
\end{enumerate}
그러면 $X^0$는 $X$에서 $X^0_y$의 근방이다.
\end{lemma}

\begin{proof}
$Y$를 $y$의 아핀 열린 근방으로 바꾸어도 된다.
$Y = \Spec(A)$라 쓰고 유한형 $\mathbf{Z}$-대수들의 유향 극한
$A = \colim A_i$로 쓰자. Limits, Lemma
\ref{limits-lemma-descend-finite-presentation}에 의해, 지표 $i$와
단면 $s_i : \Spec(A_i) \to X_i$를 갖춘 유한 표시 사상
$f_i : X_i \to \Spec(A_i)$를 찾을 수 있으며, 이를 $Y$로 밑변환하면
$f$와 그 단면 $s$를 복원한다.
필요하면 $i$를 증가시켜 $f_i$가 평탄이라고 가정할 수도 있다.
Limits, Lemma \ref{limits-lemma-descend-flat-finite-presentation}을 참조하라.
$y_i$를 $Y_i$에서 $y$의 상이라 하자.
$X_y = (X_{i, y_i}) \times_{y_i} y$임에 유의하자.
따라서 $X_{i, y_i}$는 기하학적으로 축약이다. Varieties, Lemma
\ref{varieties-lemma-geometrically-reduced-upstairs}를 참조하라.
Lemma \ref{more-morphisms-lemma-base-change-connected-along-section}에 의해
자료 $f_i, s_i, y_i \in Y_i$에 대하여 보조정리를 증명하면 충분하다.
따라서 $Y$가 뇌터환의 스펙트럼인 경우로 환원된다.

\medskip\noindent
$Y$가 뇌터환의 스펙트럼이라고 가정하자.
$f$가 유한 표시, 즉 유한형이므로 $X$도 뇌터 스킴이다.
Morphisms, Lemma \ref{morphisms-lemma-finite-type-noetherian}을 참조하라.
$x \in X^0$를 $y$ 위에 놓이는 점이라 하자. Topology, Lemma
\ref{topology-lemma-constructible-neighbourhood-Noetherian}에 의해,
$x$를 지나는 임의의 기약 닫힌 부분집합 $Z \subset X$에 대하여
교집합 $X^0 \cap Z$가 $Z$에서 조밀함을 증명하면 충분하다.
특히 일반점 $x' \in Z$가 $X^0$에 속함을 증명하면 충분하다.
Properties, Lemma
\ref{properties-lemma-locally-Noetherian-specialization-dvr}에 의해,
특수점을 $x$로, 일반점을 $x'$으로 보내는 사상 $\Spec(R) \to X$와
이산 값매김환 $R$을 찾을 수 있다.
합성 $\Spec(R) \to X \to Y$를 통해 $\Spec(R)$을 $Y$ 위의 스킴으로
생각하겠다. Lemma \ref{more-morphisms-lemma-base-change-connected-along-section}에
의해 $(X_R)^0$는 $X^0$의 역상이다.
구성에 의해 구조 사상 $X_R \to \Spec(R)$의 두 번째 단면
$t : \Spec(R) \to X_R$이 존재하며($s$의 밑변환 $s_R$ 외에),
$t(\eta_R)$는 $x'$으로 가는 $X_R$의 점이고 $t(0_R)$는 $x$로 가는
$X_R$의 점이다. $t(0_R)$가 $(X_R)^0$에 속하고
$t(\eta_R) \leadsto t(0_R)$임에 유의하자.
따라서 이로부터 $t(\eta_R) \in (X_R)^0$가 따름을 증명하면 충분하다.
그러므로 $Y$가 이산 값매김환의 스펙트럼이고 $y$가 그 닫힌점인
경우에 보조정리를 증명하면 충분하다.

\medskip\noindent
$Y$가 이산 값매김환의 스펙트럼이고 $y$가 그 닫힌점이라고 가정하자.
우리의 목표는 $X^0$가 $X^0_y$의 근방임을 증명하는 것이다.
$X_y$는 기약 성분이 유한 개이므로 $X^0_y$는 $X_y$에서 열려 있고
닫혀 있음에 유의하자. 따라서 여집합 $C = X_y \setminus X_y^0$는
$X$에서 닫혀 있다. 그러므로 $U = X \setminus C$는 $X^0_y$의 열린
근방이고 $U^0 = X^0$이다. 따라서 사상 $U \to Y$에 대하여 결과를
증명하면 충분하다. 다시 말해 $X_y$가 연결이라고 가정할 수 있다.
$X$가 연결이 아니며, $X = X_1 \amalg \ldots \amalg X_n$이 연결
성분 분해라고 하자. 그러면 $s(Y)$는 $X_i$ 중 하나에 완전히 포함된다.
$s(Y) \subset X_1$이라고 하자. 그러면 $X^0 \subset X_1$이다.
따라서 $X$를 $X_1$으로 바꾸어 $X$가 연결이라고 가정할 수 있다.
이제 Lemma \ref{more-morphisms-lemma-connected-flat-over-dvr}에 의해 $X_\eta$가
연결이므로, $X^0 = X$이고 증명이 끝난다.
\end{proof}

\begin{lemma}
\label{more-morphisms-lemma-connected-along-section-open}
$f : X \to Y$, $s : Y \to X$가 Situation
\ref{more-morphisms-situation-connected-along-section}에서와 같다고 하자.
다음을 가정하자.
\begin{enumerate}
\item $f$는 유한 표시이고 평탄이다.
\item $f$의 모든 올은 기하학적으로 축약이다.
\end{enumerate}
그러면 $X^0$는 $X$에서 열려 있다.
\end{lemma}

\begin{proof}
Lemma \ref{more-morphisms-lemma-connected-along-section-open-neighbourhood}에서
즉시 따른다.
\end{proof}






\section{올의 차원}
\label{more-morphisms-section-dimension}

\begin{lemma}
\label{more-morphisms-lemma-dimension-in-neighbourhood}
$f : X \to Y$를 스킴의 사상이라 하자.
% P05-EN-CH37-0139
$Y$가 일반점 $\eta$를 가진 기약 스킴이고 $f$가 유한형이라고 가정하자.
$X_\eta$의 차원이 $n$이면, 모든 $y \in V$에 대하여 올 $X_y$의 차원이
$n$이 되도록 하는 비어 있지 않은 열린집합 $V \subset Y$가 존재한다.
\end{lemma}

\begin{proof}
$Z = \{x \in X \mid \dim_x(X_{f(x)}) > n \}$이라 하자.
Morphisms, Lemma \ref{morphisms-lemma-openness-bounded-dimension-fibres}
에 의해 이는 $X$의 닫힌 부분집합이다. 가정에 의해 $Z_\eta = \emptyset$이다.
따라서 Lemma \ref{more-morphisms-lemma-empty-generic-fibre}에 의해
$Y$를 축소하여 $Z = \emptyset$이라고 가정할 수 있다.
$Z' = \{x \in X \mid \dim_x(X_{f(x)}) > n - 1 \} =
\{x \in X \mid \dim_x(X_{f(x)}) = n\}$이라 하자. 앞에서와 같이 이는
$X$의 닫힌 부분집합이다. 가정에 의해 $Z'_\eta \not = \emptyset$이다.
따라서 $Y$를 축소하면 $Z' \to Y$가 전사라고 가정할 수 있다.
Lemma \ref{more-morphisms-lemma-nonempty-generic-fibre}를 참조하라.
이로써 원하는 결론을 얻는다.
\end{proof}

\begin{lemma}
\label{more-morphisms-lemma-base-change-dimension-fibres}
$f : X \to Y$를 유한형 사상이라 하자. 함수
% P05-EN-CH37-0141
$$
n_{X/Y} : Y \to \{0, 1, 2, 3, \ldots, \infty\}
$$
를 $y \in Y$에 $X_y$의 차원을 대응시키는 함수라 하자.
$g : Y' \to Y$가 사상이면
$$
n_{X'/Y'} = n_{X/Y} \circ g
$$
이다. 여기서 $X' \to Y'$은 $f$의 밑변환이다.
\end{lemma}

\begin{proof}
Morphisms, Lemma \ref{morphisms-lemma-dimension-fibre-after-base-change}에서
따른다.
\end{proof}

\begin{lemma}
\label{more-morphisms-lemma-dimension-fibres-constructible}
$f : X \to Y$를 스킴의 사상이라 하자.
$n_{X/Y}$를 Lemma \ref{more-morphisms-lemma-base-change-dimension-fibres}에서
도입한, $f$의 올들의 차원을 주는 $Y$ 위의 함수라 하자.
% P05-EN-CH37-0140
$f$가 유한 표시라고 가정하자. 그러면 $n_{X/Y}$의 등위집합
$$
E_n = \{y \in Y \mid n_{X/Y}(y) = n\}
$$
은 $Y$에서 국소 구성가능이다.
\end{lemma}

\begin{proof}
$n$을 고정하자. $y \in Y$라 하자. $E_n \cap V$가 $V$에서 구성가능이
되도록 하는 $Y$에서의 $y$의 열린 근방 $V$가 존재함을 보여야 한다.
따라서 $Y$가 아핀이라고 가정할 수 있다. $Y = \Spec(A)$라 쓰고
유한형 $\mathbf{Z}$-대수들의 유향 극한 $A = \colim A_i$로 쓰자.
Limits, Lemma \ref{limits-lemma-descend-finite-presentation}에 의해
$Y$로 밑변환하면 $f$를 복원하는 유한 표시 사상
$f_i : X_i \to \Spec(A_i)$와 지표 $i$를 찾을 수 있다.
Lemma \ref{more-morphisms-lemma-base-change-dimension-fibres}에 의해
$f_i$에 대하여 보조정리를 증명하면 충분하다. 따라서
$Y$가 뇌터환의 스펙트럼인 경우로 환원된다.

\medskip\noindent
$Y$가 뇌터 스킴인 경우 $E_n$이 구성가능임을 증명하기 위해
Topology, Lemma \ref{topology-lemma-characterize-constructible-Noetherian}
의 판정법을 사용하겠다. 이를 위해 $Z \subset Y$를 기약 닫힌
부분스킴이라 하자. $E_n \cap Z$가 비어 있지 않은 열린 부분집합을
포함하거나 $Z$에서 조밀하지 않음을 보여야 한다. $\xi \in Z$를
일반점이라 하자. 그러면 Lemma \ref{more-morphisms-lemma-dimension-in-neighbourhood}
에 의해 $n_{X/Y}$는 $Z$에서의 $\xi$의 근방에서 상수이다.
이는 원하는 결론을 함의한다.
\end{proof}

\begin{lemma}
\label{more-morphisms-lemma-dimension-fibres-flat}
$f : X \to Y$를 유한 표시인 스킴의 평탄 사상이라 하자.
$n_{X/Y}$를 Lemma \ref{more-morphisms-lemma-base-change-dimension-fibres}에서
도입한, $f$의 올들의 차원을 주는 $Y$ 위의 함수라 하자.
그러면 $n_{X/Y}$는 하반연속이다.
\end{lemma}

\begin{proof}
$W \subset X$, $W = \coprod_{d \geq 0} U_d$를
Lemmas \ref{more-morphisms-lemma-flat-finite-presentation-CM-open}와
\ref{more-morphisms-lemma-flat-finite-presentation-CM-pieces}에서 구성한 열린집합이라 하자.
$y \in Y$를 점이라 하자. $n_{X/Y}(y) = \dim(X_y) = n$이면
$y$는 $U_n \to Y$의 상에 속한다.
Morphisms, Lemma \ref{morphisms-lemma-fppf-open}에 의해
$f(U_n)$은 $Y$에서 열려 있다.
따라서 $n_{X/Y}$가 $\geq n$인 $y$의 열린 근방이 존재한다.
\end{proof}

\begin{lemma}
\label{more-morphisms-lemma-dimension-fibres-proper}
$f : X \to Y$를 스킴의 고유 사상이라 하자.
$n_{X/Y}$를 Lemma \ref{more-morphisms-lemma-base-change-dimension-fibres}에서
도입한, $f$의 올들의 차원을 주는 $Y$ 위의 함수라 하자.
그러면 $n_{X/Y}$는 상반연속이다.
\end{lemma}

\begin{proof}
$Z_d = \{x \in X \mid \dim_x(X_{f(x)}) > d\}$라 하자.
그러면 Morphisms, Lemma
\ref{morphisms-lemma-openness-bounded-dimension-fibres}에 의해
$Z_d$는 $X$의 닫힌 부분집합이다.
$f$가 고유이므로 $f(Z_d)$는 닫혀 있다.
$y \in f(Z_d) \Leftrightarrow n_{X/Y}(y) > d$이므로
보조정리가 성립한다.
\end{proof}

\begin{lemma}
\label{more-morphisms-lemma-dimension-fibres-proper-flat}
$f : X \to Y$를 유한 표시인 스킴의 고유 평탄 사상이라 하자.
$n_{X/Y}$를 Lemma \ref{more-morphisms-lemma-base-change-dimension-fibres}에서
도입한, $f$의 올들의 차원을 주는 $Y$ 위의 함수라 하자.
그러면 $n_{X/Y}$는 국소 상수이다.
\end{lemma}

\begin{proof}
Lemmas \ref{more-morphisms-lemma-dimension-fibres-flat}와
\ref{more-morphisms-lemma-dimension-fibres-proper}에서 즉시 따른다.
\end{proof}

















\section{약한 상대 뇌터 정규화}
\label{more-morphisms-section-relative-noether-normalization}

\noindent
이 절의 목표는 Lemma \ref{more-morphisms-lemma-weak-relative-noether-normalization}을
증명하는 것이다.

\begin{lemma}
\label{more-morphisms-lemma-silly}
$R$을 환이라 하자. $\mathfrak p_1, \ldots, \mathfrak p_r$를
$i \not = j$이면 $\mathfrak p_i \not \subset \mathfrak p_j$인
$R$의 소아이디얼들이라 하자. $k_i \subset \kappa(\mathfrak p_i)$를
확대 $\kappa(\mathfrak p_i)/k_i$가 대수적이지 않게 하는 부분체들이라 하자.
$J \subset R$을 어느 $\mathfrak p_i$에도 포함되지 않는 아이디얼이라 하자.
그러면 $i = 1, \ldots, r$에 대하여 $\kappa(\mathfrak p_i)$에서
$x$의 상이 $k_i$ 위에서 초월적이 되도록 하는 원소 $x \in J$가 존재한다.
\end{lemma}

\begin{proof}
아이디얼
$J_i = J \mathfrak p_1 \ldots \hat{\mathfrak p}_i \ldots \mathfrak p_r$는
$\mathfrak p_i$에 포함되지 않는다. Algebra, Lemma
\ref{algebra-lemma-product-ideals-in-prime}를 참조하라.
% P05-EN-CH37-0142
따라서 $\kappa(\mathfrak p_i) = \text{Frac}(B/\mathfrak p_i)$의
모든 원소 $\xi$는 $a, b \in J_i$이고 $b \not \in \mathfrak p_i$인
$\xi = a/b$의 꼴로 쓸 수 있다.
$\xi$를 $k_i$ 위에서 초월적으로 택하면, $a$ 또는 $b$ 중 하나는
$k_i$ 위에서 초월적인 $\kappa(\mathfrak p_i)$의 원소로 간다.
따라서 모든 $i = 1, \ldots, r$에 대하여
$k_i$ 위에서 초월적인 $\kappa(\mathfrak p_i)$의 원소로 가는 원소
$x_i \in J_i = J \mathfrak p_1 \ldots \hat{\mathfrak p}_i \ldots \mathfrak p_r$
를 찾을 수 있다. 그러면 $x = x_1 + \ldots + x_r$가 조건을 만족한다.
\end{proof}

\begin{lemma}
\label{more-morphisms-lemma-sum-is-ok}
$R \to S$를 유한형 환 준동형이라 하자. $d \geq 0$이라 하고
$a, b \in S$라 하자. $x$를 $a$로 보내는 $R$-대수 준동형
$R[x] \to S$로 주어지는 사상
$$
f_a : \Spec(S) \longrightarrow \mathbf{A}^1_R
$$
의 올들의 차원이 $\leq d$라고 가정하자. 그러면 $n \geq n_0$일 때,
$x$를 $a^n + b$로 보내는 $R$-대수 준동형 $R[x] \to S$로 주어지는 사상
$$
f_{a^n + b} : \Spec(S) \longrightarrow \mathbf{A}^1_R
$$
의 올들의 차원이 $\leq d$가 되도록 하는 $n_0$가 존재한다.
\end{lemma}

\begin{proof}
이 문단에서는 $R \to S$가 유한 표시인 경우로 환원한다.
% P05-EN-CH37-0143
즉, 어떤 아이디얼 $J \subset R[x_1, \ldots, x_n]$에 대하여
$S = R[A, B, x_1, \ldots, x_n]/J$라 쓰자. 여기서 $A$와 $B$는
$S$의 $a$와 $b$로 간다. 그러면 $J$는 그 유한 생성 아이디얼들
$J_\lambda \subset J$의 합집합이다.
$S_\lambda = R[A, B, x_1, \ldots, x_n]/J_\lambda$로 놓고,
% P05-EN-CH37-0144
$A$와 $B$의 상을 $a_\lambda, b_\lambda \in S_\lambda$로 표기하자.
그러면 어떤 $\lambda$에 대하여
$$
f_{a_\lambda} : \Spec(S_\lambda) \longrightarrow \mathbf{A}^1_R
$$
의 올들의 차원이 $\leq d$이다. Limits, Lemma
\ref{limits-lemma-limit-dimension}을 참조하라.
그러한 $\lambda$를 고정하자. $R \to S_\lambda$, $a_\lambda$, $b_\lambda$에
대하여 조건을 만족하는 $n_0$를 찾을 수 있다면, 이 $n_0$는 $R \to S$에
대해서도 조건을 만족한다. 실제로
$f_{a_\lambda^n + b_\lambda} : \Spec(S_\lambda) \to \mathbf{A}^1_R$의
올들은 $f_{a^n + b} : \Spec(S) \to \mathbf{A}^1_R$의 올들을 포함한다.
이로써 다음 문단에서 다루는 경우로 환원된다.

\medskip\noindent
$R \to S$가 유한 표시라고 가정하자. 이 문단에서는 $R$이
$\mathbf{Z}$ 위에서 유한형인 경우로 환원한다.
Algebra, Lemma \ref{algebra-lemma-limit-module-finite-presentation}에 의해,
유향집합 $\Lambda$와 $\Lambda$ 위의 환 준동형들의 계 $R_\lambda \to S_\lambda$를
찾을 수 있다. 그 여극한은 $R \to S$이고, $\mu \geq \lambda$에 대하여
$S_\mu = S_\lambda \otimes_{R_\lambda} R_\mu$이며,
각 $R_\lambda$와 $S_\lambda$는 $\mathbf{Z}$ 위에서 유한형이다.
$\lambda_0 \in \Lambda$와 $a, b \in S$로 가는 원소들
$a_{\lambda_0}, b_{\lambda_0} \in S_{\lambda_0}$를 택하자.
$\lambda \geq \lambda_0$에 대하여 $a_{\lambda_0}, b_{\lambda_0}$의
상을 $a_\lambda, b_\lambda \in S_\lambda$로 표기하자.
그러면 충분히 큰 $\lambda \geq \lambda_0$에 대하여
$$
f_{a_\lambda} : 
\Spec(S_\lambda) \longrightarrow \mathbf{A}^1_{R_\lambda}
$$
의 올들의 차원이 $\leq d$이다. Limits, Lemma
\ref{limits-lemma-descend-dimension-d}를 참조하라.
그러한 $\lambda$를 고정하자.
$R_\lambda \to S_\lambda$, $a_\lambda$, $b_\lambda$에 대하여
조건을 만족하는 $n_0$를 찾을 수 있다면, 이 $n_0$는 $R \to S$에
대해서도 조건을 만족한다. 실제로 Morphisms, Lemma
\ref{morphisms-lemma-dimension-fibre-after-base-change}에 의해
$f_{a^n + b} : \Spec(S) \to \mathbf{A}^1_R$의 각 올은
$f_{a_\lambda^n + b_\lambda} : \Spec(S_\lambda) \to \mathbf{A}^1_{R_\lambda}$의
어떤 올과 같은 차원을 가진다.
% P05-EN-CH37-0145
이로써 다음 문단에서 다루는 경우로 환원된다.

\medskip\noindent
$R$과 $S$가 $\mathbf{Z}$ 위에서 유한형이라고 가정하자.
특히 $R$의 차원은 유한하므로 $\dim(R)$에 관한 귀납법을 사용할 수 있다.
따라서 보조정리에서와 같은 $R' \to S'$, $a$, $b$를 갖추고
$R'$과 $S'$가 $\mathbf{Z}$ 위에서 유한형이며
$\dim(R') < \dim(R)$인 모든 상황에 대하여 결과가 성립한다고 가정할 수 있다.

\medskip\noindent
이 명제는 $S$의 스펙트럼의 위상에 관한 것이므로 $S$가 축약이라고
가정할 수 있다. $S^\nu$를 $S$의 정규화라 하자.
$S$가 우수환이므로 $S \subset S^\nu$는 유한 확대이다. Algebra,
Proposition \ref{algebra-proposition-ubiquity-nagata}와
Morphisms, Lemma \ref{morphisms-lemma-nagata-normalization}을 참조하라.
따라서 $\Spec(S^\nu) \to \Spec(S)$는 전사이고 유한이다
(Algebra, Lemma \ref{algebra-lemma-integral-overring-surjective}).
그러므로 $R \to S^\nu$와 $S^\nu$에서의 $a$, $b$의 상에 대하여
결과가 성립하면 $R \to S$, $a$, $b$에 대해서도 결과가 성립한다.
% P05-EN-CH37-0146
(간단한 세부 사항은 생략한다.) 이로써 다음 문단에서 다루는 경우로 환원된다.

\medskip\noindent
$R$과 $S$가 $\mathbf{Z}$ 위에서 유한형이고 $S$가 정규라고 가정하자.
그러면 어떤 정규 정역들 $S_i$에 대하여 $S = S_1 \times \ldots \times S_r$이다.
각 $R \to S_i$와 $S_i$에서의 $a$, $b$의 상에 대하여 결과가 성립하면,
$R \to S$, $a$, $b$에 대해서도 결과가 성립한다.
% P05-EN-CH37-0146
(간단한 세부 사항은 생략한다.)
이로써 다음 문단에서 다루는 경우로 환원된다.

\medskip\noindent
% P05-EN-CH37-0147
$R$과 $S$가 $\mathbf{Z}$ 위에서 유한형이고 $S$가 정규 정역이라고 가정하자.
$R$을 $S$에서의 $R$의 상으로 바꾸어도 된다
(이렇게 해도 $R$의 차원은 증가하지 않는다).
이로써 다음 문단에서 다루는 경우로 환원된다.

\medskip\noindent
% P05-EN-CH37-0147
$R \subset S$가 $\mathbf{Z}$ 위에서 유한형이고 $S$가 정규 정역이라고
가정하자. 사상
$$
f_a : \Spec(S) \to \mathbf{A}^1_R
$$
을 생각하자. 가정에 의해 $f_a$의 올들의 차원은 $\leq d$이다.
따라서 $f : \Spec(S) \to \Spec(R)$의 올들의 차원은 $\leq d + 1$이다
(Morphisms, Lemma \ref{morphisms-lemma-dimension-fibre-at-a-point-additive}).
$x$를 $a$로, $y$를 $b$로 보내는 $R$-대수 준동형 $R[x, y] \to S$에
대응하는 정수적 스킴의 사상
$$
\phi : \Spec(S) \to \mathbf{A}^2_R = \Spec(R[x, y])
$$
을 생각하자. 다음 두 경우를 고려한다.
\begin{enumerate}
\item $\phi$가 지배적이다.
\item $\phi$가 지배적이지 않다.
\end{enumerate}
두 경우 모두 정수 $n_0$와 비어 있지 않은 열린집합 $V \subset \Spec(R)$가
존재하여, $n \geq n_0$일 때 점 $q \in \mathbf{A}^1_V$에서의
$f_{a^n + b}$의 올들의 차원이 $\leq d$라고 주장한다.

\medskip\noindent
경우 (1)에서 주장을 증명하자. $f_{a^n + b} = \pi_n \circ \phi$이다. 여기서
$$
\pi_n :  \mathbf{A}^2_R \to \mathbf{A}^1_R
$$
은 $x$를 $x^n + y$로 보내는 $R$-대수 준동형 $R[x] \to R[x, y]$에
대응하는 평탄 사상이다. $\phi$가 지배적이므로
$\phi|_U : U \to \mathbf{A}^2_R$가 평탄이 되도록 하는
조밀한 열린집합 $U \subset \Spec(S)$가 존재한다
(이는 예를 들어 일반 평탄성에서 따른다. Morphisms, Proposition
\ref{morphisms-proposition-generic-flatness} 참조).
그러면 합성
$$
f_{a^n + b}|_U :
U \xrightarrow{\phi|_U}
\mathbf{A}^2_R
\xrightarrow{\pi_n}
\mathbf{A}^1_R
$$
도 평탄이다. 따라서 Morphisms, Lemma
\ref{morphisms-lemma-dimension-fibre-at-a-point-additive}에 의해,
이 사상의 올들은 $f|_U : U \to \Spec(R)$의 올들 안에서
여차원이 적어도 $1$이다. 다시 말해 $f_{a^n + b}|_U$의 올들의 차원은
$\leq d$이다. 한편 $U$가 $\Spec(S)$에서 조밀하므로, 모든 $p \in V$에
대하여 $U \cap f^{-1}(p) \subset f^{-1}(p)$가 조밀하도록 하는
비어 있지 않은 열린집합 $V \subset \Spec(R)$를 찾을 수 있다
(예를 들어 Lemma \ref{more-morphisms-lemma-nowhere-dense-generic-fibre} 참조).
따라서 $\dim(f^{-1}(p) \setminus U \cap f^{-1}(p)) \leq d$이고,
% P05-EN-CH37-0148
주장이 성립함을 얻는다($f_{a^n + b} : \Spec(S) \to \mathbf{A}^1_R$의
각 올은 $f$의 어떤 올에 포함되기 때문이다).

\medskip\noindent
경우 (2). 이 경우 $\Im(\phi) \subset V(g)$를 만족하는
$R[x, y]$의 영이 아닌 원소 $g = \sum c_{ij} x^i y^j$를 찾을 수 있다.
실제로 $g$가 $\text{Frac}(R)$ 위에서 기약이라고 가정할 수 있다.
$g \in R[x]$이고 그 최고차항 계수가 $c$라면,
$V = D(c) \subset \Spec(R)$ 위에서 $f$의 올들의 차원은 이미 $\leq d$이다
($f_a$의 상은 $V(g) \subset \mathbf{A}^1_R$에 포함되며, 이 부분스킴은
$V$ 위에서 유한한 올들을 갖기 때문이다).
따라서 $g$가 $R[x]$에 속하지 않는다고 가정할 수 있다.
$s \geq 1$을 $y$의 다항식으로서 $g$의 차수라 하고,
$t$를 $x$의 다항식으로서 $\sum c_{is} x^i$의 차수라 하자.
그러면 $c_{ts}$는 영이 아니고,
$$
g(x, -x^n) = (-1)^s c_{ts} x^{t + sn} + l.o.t.
$$
는 $n$이 $x$의 다항식으로서 $g$의 차수보다 크면 성립한다
% P05-EN-CH37-0146
(간단한 세부 사항은 생략한다). 그러한 $n$에 대하여 다항식 $g(x, -x^n)$은
$x$의 영이 아닌 다항식이며,
% P05-EN-CH37-0149
$c_{st} \not \in \mathfrak p$인 모든 $\mathfrak p \subset R$에 대하여
$\kappa(\mathfrak p)[x]$의 영이 아닌 다항식으로 간다.
% P05-EN-CH37-0150
따라서 $V$를 $\Spec(R)$의 주 열린집합 $D(c_{ts})$로 택하면
주장이 성립한다.

\medskip\noindent
이제 $\dim(R)$에 관한 귀납법을 사용할 수 있다.
$I \subset R$을 $V(I) = \Spec(R) \setminus V$를 만족하는 아이디얼이라 하자.
$R$이 정역이므로 $\dim(R/I) < \dim(R)$임에 유의하자.
$n'_0$를 $R/I \to S/IS$와 $S/IS$에서의 $a$와 $b$의 상에 대하여
$\dim(R)$에 관한 귀납법으로 얻은 정수라 하자.
그러면 $\max(n_0, n'_0)$가 조건을 만족한다.
\end{proof}

\begin{lemma}
\label{more-morphisms-lemma-weak-relative-noether-normalization}
$R \to S$를 유한형 환 준동형이라 하자.
% P05-EN-CH37-0151
$d$를 $\Spec(S) \to \Spec(R)$의 올들의 차원의 최댓값이라 하자.
그러면 준유한 환 준동형 $R[t_1, \ldots, t_d] \to S$가 존재한다.
\end{lemma}

\begin{proof}
이 문단에서는 $R \to S$가 유한 표시인 경우로 환원한다.
즉, 어떤 아이디얼 $J \subset R[x_1, \ldots, x_n]$에 대하여
$S = R[x_1, \ldots, x_n]/J$라 쓰자.
그러면 $J$는 그 유한 생성 아이디얼들 $J_\lambda \subset J$의 합집합이다.
$S_\lambda = R[x_1, \ldots, x_n]/J_\lambda$로 놓자.
그러면 어떤 $\lambda$에 대하여 $\Spec(S_\lambda) \to \Spec(R)$의
올들의 차원이 $\leq d$이다. Limits, Lemma
\ref{limits-lemma-limit-dimension}을 참조하라.
그러한 $\lambda$를 고정하자. 준유한 준동형
$R[t_1, \ldots, t_d] \to S_\lambda$를 찾을 수 있다면,
물론 합성 $R[t_1, \ldots, t_d] \to S$도 준유한이다.
이로써 다음 문단에서 다루는 경우로 환원된다.

\medskip\noindent
$R \to S$가 유한 표시라고 가정하자. 이 문단에서는 $R$이
$\mathbf{Z}$ 위에서 유한형인 경우로 환원한다.
Algebra, Lemma \ref{algebra-lemma-limit-module-finite-presentation}에 의해,
유향집합 $\Lambda$와 $\Lambda$ 위의 환 준동형들의 계 $R_\lambda \to S_\lambda$를
찾을 수 있다. 그 여극한은 $R \to S$이고, $\mu \geq \lambda$에 대하여
$S_\mu = S_\lambda \otimes_{R_\lambda} R_\mu$이며, 각 $R_\lambda$와
$S_\lambda$는 $\mathbf{Z}$ 위에서 유한형이다.
그러면 충분히 큰 $\lambda$에 대하여
$\Spec(S_\lambda) \to \Spec(R_\lambda)$의 올들의 차원이 $\leq d$이다.
Limits, Lemma \ref{limits-lemma-descend-dimension-d}를 참조하라.
그러한 $\lambda$를 고정하자. 준유한 환 준동형
$R_\lambda[t_1, \ldots, t_d] \to S_\lambda$를 찾을 수 있다면,
밑변환 $R[t_1, \ldots, t_d] \to S$도 준유한이다
(Algebra, Lemma \ref{algebra-lemma-quasi-finite-base-change}).
% P05-EN-CH37-0145
이로써 다음 문단에서 다루는 경우로 환원된다.

\medskip\noindent
$R$과 $S$가 $\mathbf{Z}$ 위에서 유한형이라고 가정하자.
$d = 0$이면 환 준동형이 준유한이므로 증명이 끝난다.
$d > 0$이라고 가정하자. $R$-대수 준동형 $R[x] \to S$, $x \mapsto a$의
올들의 차원이 $< d$가 되도록 하는 원소 $a \in S$를 찾겠다.
그러면 귀납법으로 증명이 끝난다.

\medskip\noindent
$\dim(R)$에 관한 귀납법으로 $a$의 존재를 증명하겠다.

\medskip\noindent
$\mathfrak q_1, \ldots, \mathfrak q_r \subset S$를 $S$의 극소 소아이디얼
중 $\dim_{\mathfrak q_i}(S/R) = d$를 만족하는 것들이라 하자.
표기에 관해서는 Algebra, Definition
\ref{algebra-definition-relative-dimension}을 참조하라.
$\mathfrak q_i$가 소아이디얼 $\mathfrak p_i \subset R$ 위에 놓인다고 하자.
$\mathfrak q_i$는 그 올의 일반점이므로
$\text{trdeg}_{\kappa(\mathfrak p_i)}(\kappa(\mathfrak q_i)) = d$이다.
예를 들어 $S \otimes_R \kappa(\mathfrak p_i)$에
Algebra, Lemma \ref{algebra-lemma-dimension-at-a-point-finite-type-field}를
적용하면 된다. 따라서 Lemma \ref{more-morphisms-lemma-silly}에 의해,
$i = 1, \ldots, r$에 대하여 $\kappa(\mathfrak q_i)$에서 $a$의 상이
$\kappa(\mathfrak p_i)$ 위에서 초월적이 되도록 하는 원소 $a \in S$를
찾을 수 있다. 사상
$$
f_a : \Spec(S) \longrightarrow \mathbf{A}^1_R
$$
% P05-EN-CH37-0152
을 생각하자. 이는 $x$를 $a$로 보내는 $R$-대수 준동형 $R[x] \to S$에
대응한다. $U \subset \Spec(S)$를 올들의 차원이 $\leq d - 1$인
열린 부분집합이라 하자. Morphisms, Lemma
\ref{morphisms-lemma-openness-bounded-dimension-fibres}를 참조하라.
구성에 의해 $U$는 $\Spec(S)$의 모든 일반점을 포함한다.
특히 $\Spec(S) \to \Spec(R)$의 모든 일반올의 모든 일반점은
반드시 $\Spec(S)$의 일반점이므로, $U$가 이들을 모두 포함함을 알 수 있다.
$T = \Spec(S) \setminus U$로 놓고 이를 $\Spec(S)$의 축약 닫힌
부분스킴으로 생각하자. 방금 한 관찰과 가정 $\dim(S/R) \leq d$에 의해,
$T \to \Spec(R)$의 일반올들의 차원은 $\leq d - 1$이다.
따라서 Lemma \ref{more-morphisms-lemma-dimension-in-neighbourhood}을 여러 번 적용하여
$\Spec(R)$의 일반점들의 열린 근방을 만들면,
$T_V \to V$의 올들의 차원이 $\leq d - 1$인 조밀한 열린집합
$V \subset \Spec(R)$를 찾을 수 있다.
그러므로 $q \in \mathbf{A}^1_V$에 대하여 $q$ 위의 $f_a$의 올의 차원은
$\leq d - 1$이다($f_a|_U$의 올과 $f_a|_T$의 올의 차원을 모두
제한했기 때문이다).

\medskip\noindent
소아이디얼 회피에 의해 어떤 $f \in R$에 대하여 $V = D(f)$라고
가정할 수 있다. 그러면 환 준동형 $R_f[x] \to S_f$, $x \mapsto a$의
올들의 차원은 $\leq d - 1$이다.
$a$를 $fa$로 바꾸어 $a \in (f)$라고 가정할 수 있다.
$\dim(R)$에 관한 귀납법에 의해, 원소 $\overline{b} \in S/fS$를 찾아
% P05-EN-CH37-0153
$\Spec(S/fS) \to \Spec(R/fR[x])$, $x \mapsto \overline{b}$의
올들의 차원이 $\leq d - 1$이 되게 할 수 있다.
$b \in S$를 $\overline{b}$의 올림이라 하자.
Lemma \ref{more-morphisms-lemma-sum-is-ok}에 의해 $a^n + b$가 여전히 $R_f \to S_f$에
대하여 조건을 만족하도록 하는 $n > 0$이 존재한다.
한편 $S/fS$에서 $a^n + b$의 상은 $\overline{b}$이므로 증명이 끝난다.
\end{proof}













\section{베르티니 정리}
\label{more-morphisms-section-bertini}

\noindent
Varieties, Section \ref{varieties-section-bertini}에서 시작한 논의를
계속한다. 이 절에서는 \cite{Jou}에 제시된 놀라운 논증을 따라,
기하학적으로 기약인 대수다양체의 일반 초평면 단면이
기하학적으로 기약임을 증명한다.

\begin{lemma}
\label{more-morphisms-lemma-amazing-bertini-lemma}
\begin{reference}
\cite{Jou}의 71쪽과 72쪽을 참조하라.
\end{reference}
$K/k$를 기하학적으로 기약인 유한 생성 체 확대라 하자.
$n \geq 1$이라 하자.
$g_1, \ldots, g_n \in K$를, 다음 원소들이 $k$ 위에서 대수적으로
독립이 되도록 하는 $c_1, \ldots, c_n \in k$가 존재하게 하는
원소들이라 하자.
$$
x_1, \ldots, x_n, \sum g_ix_i, \sum c_ig_i \in K(x_1, \ldots, x_n)
$$
그러면 $K(x_1, \ldots, x_n)$은
$k(x_1, \ldots, x_n, \sum g_ix_i)$ 위에서 기하학적으로 기약이다.
\end{lemma}

\begin{proof}
$c_1, \ldots, c_n \in k$를 보조정리의 명제에서와 같이 택하자.
$\xi = \sum g_ix_i$, $\delta = \sum c_ig_i$라 쓰자.
$a \in k$에 대하여 $K$ 위에서는 항등이고 다음 규칙으로 주어지는
$K(x_1, \ldots, x_n)$의 자기동형 $\sigma_a$를 생각하자.
$$
\sigma_a(x_i) = x_i + a c_i
$$
$\sigma_a(\xi) = \xi + a \delta$이고 $\sigma_a(\delta) = \delta$임에
유의하자. 체들의 탑
$$
K_0 = k(x_1, \ldots, x_n) \subset
K_1 = K_0(\xi) \subset
K_2 = K_0(\xi, \delta) \subset K(x_1, \ldots, x_n) = \Omega
$$
을 생각하자. $\sigma_a(K_0) = K_0$이고 $\sigma_a(K_2) = K_2$임에
유의하자. $\theta \in \Omega$가 $K_1$ 위에서 분리 대수적이라고 하자.
Algebra, Lemma
\ref{algebra-lemma-geometrically-irreducible-separable-elements}에 의해
$\theta \in K_1$임을 보이면 된다.

\medskip\noindent
% P05-EN-CH37-0154
$K'_2$를 $\Omega$ 안에서 $K_2$의 분리 대수적 폐포라 하자.
$K'_2/K_2$는 유한이고(Algebra, Lemma
\ref{algebra-lemma-make-geometrically-irreducible}) 분리적이므로,
$K'_2$와 $K_2$ 사이에는 유한 개의 체만 존재한다
(Fields, Lemma \ref{fields-lemma-primitive-element}).
$k$가 무한체이면\footnote{유한체의 경우는 증명의 마지막 문단에서 다룬다.},
서로 다른 원소 $a_1, a_2$를 $k$에서 찾아 다음을 만족시킬 수 있다.
$$
K_2(\sigma_{a_1}(\theta)) = K_2(\sigma_{a_2}(\theta))
$$
이는 $\Omega$의 부분체로서의 등식이다.
$\theta_i = \sigma_{a_i}(\theta)$,
$\xi_i = \sigma_{a_i}(\xi) = \xi + a_i \delta$라 쓰자. 다음에 유의하자.
$$
K_2 = K_0(\xi_1, \xi_2)
$$
실제로 $\xi_i = \xi + a_i \delta$,
$\xi = (a_2 \xi_1 - a_1 \xi_2)/(a_2 - a_1)$,
$\delta = (\xi_1 - \xi_2)/(a_1 - a_2)$이다.
$K_2/K_0$가 차수 $2$의 순수 초월 확대이므로,
$\xi_1$과 $\xi_2$는 $K_0$ 위에서 대수적으로 독립이다.
$\theta_1$이 $K_0(\xi_1)$ 위에서 대수적이므로,
$\xi_2$는 $K_0(\xi_1, \theta_1)$ 위에서 초월적이다.

\medskip\noindent
가정에 의해 $K/k$는 기하학적으로 기약이다. 따라서
$K(x_1, \ldots, x_n)/K_0$는 기하학적으로 기약이다
(Algebra, Lemma
\ref{algebra-lemma-geometrically-irreducible-base-change-transcendental}).
그러면 부분확대로서 $K_0(\xi_1, \theta_1)/K_0$도
기하학적으로 기약이다
(Algebra, Lemma \ref{algebra-lemma-subalgebra-geometrically-irreducible}).
$\xi_2$가 $K_0(\xi_1, \theta_1)$ 위에서 초월적이므로,
$K_0(\xi_1, \xi_2, \theta_1)/K_0(\xi_2)$는 기하학적으로 기약이다
(Algebra, Lemma
\ref{algebra-lemma-geometrically-irreducible-add-transcendental}).
위에서 $a_1, a_2$를 택한 방식에 의해
$$
K_0(\xi_1, \xi_2, \theta_1) =
K_2(\sigma_{a_1}(\theta)) =
K_2(\sigma_{a_2}(\theta)) =
K_0(\xi_1, \xi_2, \theta_2)
$$
이다. $\theta_2$가 $K_0(\xi_2)$ 위에서 분리 대수적이므로,
Algebra, Lemma
\ref{algebra-lemma-geometrically-irreducible-separable-elements}를 다시
적용하면 $\theta_2 \in K_0(\xi_2)$이다.
% P05-EN-CH37-0155
이 관계에 $\sigma_{a_2}^{-1}$을 취하면 원하는 대로
$\theta \in K_0(\xi) = K_1$을 얻는다.

\medskip\noindent
이로써 $k$가 무한체인 경우의 증명이 끝난다. $k$가 유한체이면
변수 $t$를 택하고 확대 $K(t)/k(t)$를 생각할 수 있다.
이는 Algebra, Lemma
\ref{algebra-lemma-geometrically-irreducible-base-change-transcendental}에
의해 기하학적으로 기약이다.
% P05-EN-CH37-0156
$K(t, x_1, \ldots, x_n)$의 원소들
$x_1, \ldots, x_n, \sum g_ix_i, \sum c_ig_i$는 여전히 $k(t)$ 위에서
대수적으로 독립이므로, 이미 제시한 논증에 의해
$K(t, x_1, \ldots, x_n)$은
$k(t, x_1, \ldots, x_n, \sum g_ix_i)$ 위에서 기하학적으로 기약이다.
이제 Algebra, Lemma
\ref{algebra-lemma-geometrically-irreducible-base-change-transcendental}을
한 번 더 적용하면 증명이 끝난다.
\end{proof}

\begin{lemma}
\label{more-morphisms-lemma-algebraically-independent}
$A$를 체 $k$ 위의 유한형 정역이라 하자. $n \geq 2$라 하자.
$g_1, \ldots, g_n \in A$를 $V(g_1, g_2)$가 차원
$\dim(A) - 2$인 기약 성분을 갖도록 하는 원소들이라 하자.
그러면 다음 원소들이 $k$ 위에서 대수적으로 독립이 되도록 하는
$c_1, \ldots, c_n \in k$가 존재한다.
$$
x_1, \ldots, x_n, \sum g_ix_i, \sum c_ig_i \in
\text{Frac}(A)(x_1, \ldots, x_n)
$$
\end{lemma}

\begin{proof}
$k$ 위에서의 대수적 독립성은 $y = \sum g_ix_i$와
$z = \sum c_ig_i$로 주어지는 사상
$$
T = \Spec(A[x_1, \ldots, x_n])
\longrightarrow
\Spec(k[x_1, \ldots, x_n, y, z]) = S
$$
이 지배적이라는 뜻이다. $d = \dim(A)$로 놓자.
$T \to S$가 지배적이지 않으면 그 상의 차원은 $< n + 2$이고,
따라서 모든 올의 모든 기약 성분의 차원은
$> d + n - (n + 2) = d - 2$이다. Varieties, Lemma
\ref{varieties-lemma-dimension-fibres-locally-algebraic}를 참조하라.
$V(g_1, g_2)$의 차원 $d - 2$인 기약 성분에 포함되고 다른 어떤
성분에도 포함되지 않는 닫힌점 $u \in V(g_1, g_2)$를 택하자.
% P05-EN-CH37-0157
$u$ 위에 놓이는 $T$의 닫힌점 $t = (u, 1, 0, \ldots 0)$를 생각하자.
$(c_1, \ldots, c_n) = (0, 1, 0, \ldots, 0)$으로 놓자.
그러면 $t$는 $S$의 점 $s = (1, 0, \ldots, 0)$으로 간다.
$s$ 위의 $T \to S$의 올은
$$
x_1 - 1, x_2, \ldots, x_n, \sum x_ig_i, g_2
$$
에 의해 잘리며, 따라서 동치로
$$
x_1 - 1, x_2, \ldots, x_n, g_1, g_2
$$
에 의해 잘린다. $g_1, g_2$에 관한 조건에 의해 이 부분스킴은
차원 $d - 2$인 기약 성분을 갖는다.
\end{proof}

\begin{lemma}
\label{more-morphisms-lemma-bertini-irreducible}
\begin{reference}
\cite[Theorem 6.3 part 4)]{Jou}
\end{reference}
Varieties, Situation \ref{varieties-situation-family-divisors}에서 다음을
가정하자.
\begin{enumerate}
\item $X$는 $k$ 위에서 유한형이다.
\item $X$는 $k$ 위에서 기하학적으로 기약이다.
\item $v_1, v_2, v_3 \in V$와 $H_{v_2} \cap H_{v_3}$의 기약 성분
$Z$가 존재하여 $Z \not \subset H_{v_1}$이고
$\text{codim}(Z, X) = 2$이다.
\item $\bigcap_{v \in V} H_v$의 모든 기약 성분 $Y$에 대하여
$\text{codim}(Y, X) \geq 2$이다.
\end{enumerate}
그러면 일반적인 $v \in V \otimes_k k'$에 대하여 스킴 $H_v$는
$k'$ 위에서 기하학적으로 기약이다.
\end{lemma}

\begin{proof}
가정 (3)이 성립하려면 원소 $v_1, v_2, v_3$는 $V$에서 $k$-선형
독립이어야 한다.
% P05-EN-CH37-0162
(간단한 세부 사항은 생략한다.)
따라서 이 원소들을 처음 $3$개로 포함하는 $V$의 기저
$v_1, \ldots, v_r$를 택할 수 있다.
$H_{univ} \subset \mathbf{A}^r_k \times_k X$가 ``보편 인자''임을
상기하자. 스킴론적 올들이 인자 $H_v$인 사영
$q : H_{univ} \to \mathbf{A}^r_k$를 생각하자.
Lemma \ref{more-morphisms-lemma-geom-irreducible-generic-fibre}에 의해
$q$의 일반올이 기하학적으로 기약임을 보이면 충분하다.
이를 증명하기 위해 $X$를 그 축약으로 바꾸어도 되므로,
$X$가 $k$ 위의 유한형 정수적 스킴이라고 가정할 수 있다.

\medskip\noindent
$\mathcal{L}|_U \cong \mathcal{O}_U$가 되도록 하는 비어 있지 않은
아핀 열린집합 $U \subset X$를 택하자. $U = \Spec(A)$라 쓰자.
% P05-EN-CH37-0158
동형 $\mathcal{L}|_U \cong \mathcal{O}_U$를 통해 절단
$\psi(v_i)|_U$에 대응하는 원소를 $f_i \in A$라 하자.
그러면 $H_{univ} \cap (\mathbf{A}^r_k \times_k U)$는 다음으로 주어진다.
$$
H_U = \Spec(A[x_1, \ldots, x_r]/(x_1f_1 + \ldots + x_rf_r))
$$
% P05-EN-CH37-0159
기저를 택한 방식에 의해 $f_1$은 영일 수 없다. 실제로 영이라면
$v_1 = 0$이고 따라서 $H_{v_1} = X$가 되어 가정 (3)에 모순이다.
그러므로 $\sum x_if_i$는 $A[x_1, \ldots, x_r]$에서 영인자가 아니다.
따라서 $d = \dim(X) = \dim(A)$라고 할 때 $H_U$의 모든 기약 성분의
차원은 $d + r - 1$이다. $U' = U \cap D(f_1)$이면
$$
H_{U'} =
\Spec(A_{f_1}[x_1, \ldots, x_r]/(x_1f_1 + \ldots + x_rf_r)) \cong
% P05-EN-CH37-0160
\Spec(A_{f_1}[x_2, \ldots x_r]) =
\mathbf{A}^{r - 1}_k \times_k U'
$$
은 기약이다. 한편
$$
H_U \setminus H_{U'} =
\Spec(A/(f_1)[x_1, \ldots, x_r]/(x_2f_2 + \ldots + x_rf_r))
$$
의 차원은 많아야 $d + r - 2$이다. 실제로 $i \not = 1$에 대하여
스킴 $(H_U \setminus H_{U'}) \times_U D(f_i)$는 비어 있거나
($f_i = 0$인 경우), 위와 같은 논증에 의해 $\Spec(A/(f_1))$의 열린집합
위의 $r - 1$차원 아핀 공간과 동형이므로 그 차원은 많아야
$d + r - 2$이다. 한편
$(H_U \setminus H_{U'}) \times_U V(f_2, \ldots, f_r)$는
$\Spec(A/(f_1, \ldots, f_r))$ 위의 $r$차원 아핀 공간이므로,
가정 (4)에 의해 그 차원은 많아야 $d + r - 2$이다.
따라서 위와 같은 모든 $U$에 대하여 $H_U$는 기약이다.
그러므로 $H_{univ}$는 기약이다.

\medskip\noindent
따라서 $H_{univ}$의 일반점이 $\mathbf{A}^r_k$의 일반점 위에서
기하학적으로 기약임을 보이면 충분하다. Varieties, Lemma
\ref{varieties-lemma-geometrically-irreducible-function-field}를 참조하라.
가정 (3)에서 존재가 주어진 $H_{v_2} \cap H_{v_3}$의 기약 성분 $Z$와
만나고 $X \setminus H_{v_1}$에 포함되는 $X$의 비어 있지 않은 아핀
열린집합 $U = \Spec(A)$를 택하자. 위의 표기를 사용하면 체 확대
$$
\text{Frac}(A[x_1, \ldots, x_r]/(x_1f_1 + \ldots + x_rf_r))/
k(x_1, \ldots , x_r)
$$
이 기하학적으로 기약임을 증명해야 한다. $U$를 택한 방식에 의해
$f_1$은 $A$에서 가역임에 유의하자.
% P05-EN-CH37-0161
$K = \text{Frac}(A)$를 $A$의 분수체라 하자.
위와 같이 변수 $x_1$을 소거하면 체 확대
$$
K(x_2, \ldots, x_r)/
k(x_2, \ldots, x_r, -\sum\nolimits_{i = 2, \ldots, r} f_1^{-1}f_i x_i)
$$
가 기하학적으로 기약임을 보여야 한다.
Lemma \ref{more-morphisms-lemma-amazing-bertini-lemma}에 의해, 어떤
$c_2, \ldots, c_r \in k$에 대하여 다음 원소들이
$$
x_2, \ldots, x_r, \sum\nolimits_{i = 2, \ldots, r} f_1^{-1}f_i x_i,
\sum\nolimits_{i = 2, \ldots, r} c_if_1^{-1}f_i
$$
$A[x_2, \ldots, x_r]$의 분수체에서 $k$ 위에 대수적으로 독립임을
보이면 충분하다. 이는 Lemma \ref{more-morphisms-lemma-algebraically-independent}와
$Z \cap U$가 $V(f_1^{-1}f_2, f_1^{-1}f_3) \subset U$의 기약 성분이라는
사실에서 따른다.
\end{proof}

\begin{remark}
\label{more-morphisms-remark-geometric-proof-bertini-irreducible}
Lemma \ref{more-morphisms-lemma-bertini-irreducible}의 특수한 경우에 대한 ``기하학적''
증명을 개략적으로 설명하자. 즉, $k$가 대수적으로 닫힌 체이고
$X \subset \mathbf{P}^n_k$가 차원 $\geq 2$인 매끄러운 기약
스킴이라고 하자. 그러면 $X \cap H$가 매끄럽고 기약이 되도록 하는
초평면 $H \subset \mathbf{P}^n_k$가 존재한다고 주장한다.
실제로 Varieties, Lemma \ref{varieties-lemma-bertini}에 의해
일반적인 $v \in V = kT_0 \oplus \ldots \oplus kT_n$에 대하여
대응하는 초평면 단면 $X \cap H_v$는 매끄럽다.
한편 Enriques--Severi--Zariski에 의해 스킴 $X \cap H_v$는 연결이다.
Varieties, Lemma \ref{varieties-lemma-connectedness-ample-divisor}를
참조하라. 따라서 $X \cap H_v$는 매끄럽고 기약이다.
\end{remark}













\section{정육면체 정리}
\label{more-morphisms-section-theorem-cube}

\noindent
다음 보조정리는 그 안의 가정들 아래에서 Picard 함자의 대각 사상이
국소 닫힌 몰입들로 표현가능함을 말해 준다.

\begin{lemma}
\label{more-morphisms-lemma-diagonal-picard-flat-proper}
$f : X \to S$를 유한 표시인 평탄 고유 사상이라 하자.
$\mathcal{E}$를 유한 국소 자유 $\mathcal{O}_X$-가군이라 하자.
사상 $g : T \to S$에 대하여 밑변환 도식
$$
\xymatrix{
X_T \ar[d]_p \ar[r]_q & X \ar[d]^f \\
T \ar[r]^g & S
}
$$
을 생각하자. 모든 $g : T \to S$에 대하여
$\mathcal{O}_T \to p_*\mathcal{O}_{X_T}$가 동형이라고 가정하자.
그러면 다음 성질을 갖는 유한 표시 몰입 $j : Z \to S$가 존재한다.
사상 $g : T \to S$가 $Z$를 통해 인수분해될 필요충분조건은
$p^*\mathcal{N} \cong q^*\mathcal{E}$를 만족하는 유한 국소 자유
$\mathcal{O}_T$-가군 $\mathcal{N}$이 존재하는 것이다.
\end{lemma}

\begin{proof}
$H^0(X_s, \mathcal{O}_{X_s}) = \kappa(s)$라는 가정에 의해 $f$의
올 $X_s$는 연결임에 유의하자. 따라서 $\mathcal{E}$의 랭크는
올들 위에서 상수이다. $f$가 열린 사상이고
(Morphisms, Lemma \ref{morphisms-lemma-fppf-open}) 닫힌 사상이므로,
$S$를 열린 동시에 닫힌 부분스킴들로 분해한
$S = \coprod S_r$가 존재하여 $S_r$의 역상 위에서
$\mathcal{E}$의 랭크가 상수 $r$임을 얻는다.
따라서 $\mathcal{E}$의 랭크가 상수 $r$이라고 가정할 수 있다.
% P05-EN-CH37-0163
$\mathcal{E}^\vee = \SheafHom(\mathcal{E}, \mathcal{O}_X)$를
랭크 $r$인 쌍대 가군이라 하겠다.

\medskip\noindent
코호몰로지와 밑변환에 의해(더 정확히는 Derived Categories of Schemes,
Lemma \ref{perfect-lemma-flat-proper-perfect-direct-image-general}에 의해),
$E = Rf_*\mathcal{E}$는 $S$의 유도 범주의 완전 대상이고 그 구성은
임의의 밑변환과 가환한다. $E' = Rf_*\mathcal{E}^\vee$도 마찬가지이다.
차수 $< 0$에서는 코호몰로지가 전혀 없으므로, 어떤 $b$에 대하여
$E$와 $E'$의 Tor-진폭은 (국소적으로) $[0, b]$에 들어간다.
임의의 $g : T \to S$에 대하여
$p_*(q^*\mathcal{E}) = H^0(Lg^*E)$이고
$p_*(q^*\mathcal{E}^\vee) = H^0(Lg^*E')$임에 유의하자.
$j : Z \to S$와 $j' : Z' \to S$를 Derived Categories of Schemes,
Lemma \ref{perfect-lemma-locally-closed-where-H0-locally-free}에서
$E$, $E'$에 대해 $a = 0$, $r = r$로 구성한 유한 표시 몰입이라 하자.
대략 말하면 이들은 $H^0(Lj^*E)$와 $H^0((j')^*E')$가 당김과 가환하는
유한 국소 자유 가군이라는 성질로 특징지어진다.

\medskip\noindent
$g : T \to S$를 사상이라 하자. 보조정리에서와 같은
$\mathcal{N}$이 존재하면, 사영 공식 Cohomology, Lemma
\ref{cohomology-lemma-projection-formula}를 사용하여 다음 가군들을 얻는다.
$$
p_*(q^*\mathcal{E}) \cong
p_*(p^*\mathcal{N}) \cong
\mathcal{N} \otimes_{\mathcal{O}_T} p_*\mathcal{O}_{X_T} \cong
\mathcal{N}\quad\text{이고 마찬가지로 }\quad
p_*(q^*\mathcal{E}^\vee) \cong \mathcal{N}^\vee
$$
이들은 랭크 $r$인 유한 국소 자유 가군이고, 이후 임의로 밑변환
$T' \to T$를 해도 그러하다. 따라서 이 경우 $T \to S$는 $j$를 통해서도,
$j'$을 통해서도 인수분해된다. 그러므로 $S$를 $Z \times_S Z'$으로
바꾸어, $f_*\mathcal{E}$와 $f_*\mathcal{E}^\vee$가 랭크 $r$인
유한 국소 자유 $\mathcal{O}_S$-가군이고 그 구성이 임의의 밑변환과
가환한다고 가정할 수 있다(간단한 세부 사항은 생략한다).

\medskip\noindent
% P05-EN-CH37-0164
이 상황에서 $g : T \to S$가 사상이고 보조정리에서와 같은
$\mathcal{N}$이 존재하면, 다음 사상(차수 $0$의 컵곱)은
$$
p_*(q^*\mathcal{E})
\otimes_{\mathcal{O}_T}
p_*(q^*\mathcal{E}^\vee)
\longrightarrow \mathcal{O}_T
$$
완전 쌍이다. 역으로 이 컵곱 사상이 완전 쌍이면, $T$ 위에서 국소적으로
$p_*(q^*\mathcal{E})$ 안의 절단들의 기저
$\sigma_1, \ldots, \sigma_r$와 $p_*(q^*\mathcal{E}^\vee)$ 안의
$\tau_1, \ldots, \tau_r$를 택하여 그 곱들이
$\sigma_i \tau_j = \delta_{ij}$를 만족하게 할 수 있다.
$\sigma_i$를 $X_T$ 위의 $q^*\mathcal{E}$의 절단으로,
$\tau_j$를 $X_T$ 위의 $q^*\mathcal{E}^\vee$의 절단으로 생각하면
$$
\sigma_1, \ldots, \sigma_r :
\mathcal{O}_{X_T}^{\oplus r}
\longrightarrow
q^*\mathcal{E}
$$
가 동형이고 그 역은
$$
\tau_1, \ldots, \tau_r :
q^*\mathcal{E}
\longrightarrow
\mathcal{O}_{X_T}^{\oplus r}
$$
로 주어짐을 얻는다. 다시 말해
$p^*p_*q^*\mathcal{E} \cong q^*\mathcal{E}$이다.
한편 컵곱이 비퇴화라는 조건은 역콤팩트 열린 부분스킴
(즉 적절한 행렬식이 영이 아닌 자취)을 골라내므로 증명이 끝난다.
\end{proof}

\noindent
특히 위 보조정리는 고유 공간들의 고유 평탄 족 위의 벡터 다발이
올들 위에서 자명하면 어떤 벡터 다발의 당김임을 말해 준다.
이를 조금 더 자세히 적어 보자.

\begin{lemma}
\label{more-morphisms-lemma-trivial-on-fibres}
$f : X \to S$를 유한 표시인 평탄 고유 사상이라 하자.
$f_*\mathcal{O}_X = \mathcal{O}_S$이고 이는 임의의 밑변환 후에도
성립한다고 하자. $\mathcal{E}$를 유한 국소 자유
$\mathcal{O}_X$-가군이라 하자. 다음을 가정하자.
\begin{enumerate}
\item 모든 $s \in S$에 대하여
$\mathcal{E}|_{X_s}$는 $\mathcal{O}_{X_s}^{\oplus r_s}$와 동형이다.
\item $S$는 축약이다.
\end{enumerate}
그러면 어떤 유한 국소 자유 $\mathcal{O}_S$-가군 $\mathcal{N}$에
대하여 $\mathcal{E} = f^*\mathcal{N}$이다.
\end{lemma}

\begin{proof}
실제로 이 경우 Lemma \ref{more-morphisms-lemma-diagonal-picard-flat-proper}의
국소 닫힌 몰입 $j : Z \to S$는 전단사이므로 닫힌 몰입이다.
그러나 $S$가 축약이므로 $j$는 동형이다.
\end{proof}

\begin{lemma}
\label{more-morphisms-lemma-triviality-generic-fibre-valuation-ring}
$f : X \to S$를 유한 표시인 고유 평탄 사상이라 하자.
$\mathcal{L}$을 가역 $\mathcal{O}_X$-가군이라 하자. 다음을 가정하자.
\begin{enumerate}
\item $S$는 값매김환의 스펙트럼이다.
\item $\mathcal{L}$은 $f$의 일반올 $X_\eta$ 위에서 자명하다.
\item $f$의 닫힌올 $X_0$는 정수적이다.
\item $H^0(X_\eta, \mathcal{O}_{X_\eta})$는 $S$의 함수체와 같다.
\end{enumerate}
그러면 $\mathcal{L}$은 자명하다.
\end{lemma}

\begin{proof}
$S = \Spec(A)$라 쓰자. 먼저 $A$가 이산 값매김환인 경우에 보조정리를
증명하겠다(실제로 가장 자주 사용되는 경우이다).
$\pi \in A$를 균등화원이라 하자.
자명화 절단 $s \in \Gamma(X_\eta, \mathcal{L}_\eta)$를 택하자.
필요하면 $s$를 $\pi^n s$로 바꾸어
$s \in \Gamma(X, \mathcal{L})$이라고 가정할 수 있다.
$s|_{X_0} = 0$이면 $s$는 $\pi$로 나누어진다
($X_0$가 스킴론적 올이고 $X$가 $A$ 위에서 평탄하기 때문이다).
따라서 $s|_{X_0}$가 영이 아니라고 가정할 수 있다.
그러면 $s$의 영점 자취 $Z(s)$는 $X_0$에 포함되지만,
$X_0$의 일반점을 포함하지는 않는다($X_0$가 정수적이기 때문이다).
이는 $Z(s)$의 $X$에서의 여차원이 $\geq 2$라는 뜻이고,
$Z(s) = \emptyset$가 아니면 Divisors, Lemma
\ref{divisors-lemma-effective-Cartier-codimension-1}에 모순이다.
따라서 원하는 결론을 얻는다.

\medskip\noindent
일반적인 경우의 증명. 값매김환 $A$는 정합환이므로
(Algebra, Example \ref{algebra-example-valuation-ring-coherent}),
$H^0(X, \mathcal{L})$은 정합 $A$-가군이다. Derived Categories of
Schemes, Lemma \ref{perfect-lemma-cohomology-over-coherent-ring}을 참조하라.
동치로 $H^0(X, \mathcal{L})$은 유한 표시 $A$-가군이다
(Algebra, Lemma \ref{algebra-lemma-coherent-ring}).
$H^0(X, \mathcal{L})$은 꼬임이 없으므로($X$가 $A$ 위에서 평탄하기
때문이다), More on Algebra, Lemma
\ref{more-algebra-lemma-generalized-valuation-ring-modules}에 의해
어떤 $n$에 대하여 $H^0(X, \mathcal{L}) = A^{\oplus n}$이다.
평탄 밑변환(Cohomology of Schemes, Lemma
\ref{coherent-lemma-flat-base-change-cohomology})에 의해
$$
K = H^0(X_\eta, \mathcal{O}_{X_\eta}) \cong
H^0(X_\eta, \mathcal{L}_\eta) =
H^0(X, \mathcal{L}) \otimes_A K
$$
이다. 여기서 $K$는 $A$의 분수체이다. 따라서 $n = 1$이다.
생성원 $s \in H^0(X, \mathcal{L})$을 택하자.
$\mathfrak m \subset A$를 극대 아이디얼이라 하자.
% P05-EN-CH37-0165
그러면 $\kappa = A/\mathfrak m = \colim A/\pi$이다. 여기서 여극한은
영이 아닌 $\pi \in \mathfrak m$에 관한 여과 여극한이다
($A$가 값매김환이라는 사실을 사용했다).
따라서 $X_0 = \lim X \times_S \Spec(A/\pi)$이다.
$s|_{X_0}$가 영이면 어떤 $\pi$에 대하여 $s$는
$X \times_S \Spec(A/\pi)$ 위에서 영으로 제한된다.
Limits, Lemma \ref{limits-lemma-descend-section}을 참조하라.
그러나 이 경우 $\pi^{-1} s$는 $\mathcal{L}$의 대역 절단이 되어,
$s$가 $H^0(X, \mathcal{L})$의 생성원이라는 사실에 모순이다.
따라서 $s|_{X_0}$는 영이 아니다. $Z(s) \subset X$를 $s$의 영점
스킴이라 하자. $s|_{X_0}$가 영이 아니고 $X_0$가 정수적이므로
$Z(s)_0 \subset X_0$는 유효 Cartier 인자이다.
$f$가 고유이고 $S$가 국소 스킴이므로, $Z(s)$의 모든 점은
$Z(s)_0$의 어떤 점으로 특수화된다. 따라서 Divisors, Lemma
\ref{divisors-lemma-fibre-Cartier}의 (3)에 의해 $Z(s)$는 상대 유효
Cartier 인자이며, 특히 $Z(s) \to S$는 평탄이다.
% P05-EN-CH37-0166
그러므로 $Z(s)$가 비어 있지 않으면 $Z(s)_\eta$도 비어 있지 않은데,
이는 $s|_{X_\eta}$가 $\mathcal{L}_\eta$의 자명화라는 사실에 모순이다.
따라서 원하는 대로 $Z(s) = \emptyset$이다.
\end{proof}

\begin{lemma}
\label{more-morphisms-lemma-get-a-closed}
$f : X \to S$와 $\mathcal{E}$가 Lemma
\ref{more-morphisms-lemma-diagonal-picard-flat-proper}에서와 같다고 하고,
추가로 $\mathcal{E}$가 가역 $\mathcal{O}_X$-가군이라고 가정하자.
더욱이 $f$의 기하학적 올들이 정수적이면 $Z$는 $S$에서 닫혀 있다.
\end{lemma}

\begin{proof}
$j : Z \to S$가 유한 표시이므로 다음을 보이면 충분하다.
$A$가 분수체 $K$를 갖는 값매김환이고 $g(\Spec(K)) \in j(Z)$인
임의의 사상 $g : \Spec(A) \to S$에 대하여
$g(\Spec(A)) \subset j(Z)$이다. Morphisms, Lemma
\ref{morphisms-lemma-reach-points-scheme-theoretic-image}를 참조하라.
이는 Lemma \ref{more-morphisms-lemma-triviality-generic-fibre-valuation-ring}과
Lemma \ref{more-morphisms-lemma-diagonal-picard-flat-proper}의
$j : Z \to S$에 대한 특징화에서 따른다.
\end{proof}

\begin{lemma}
\label{more-morphisms-lemma-H1-O-picard-flat-proper}
스킴의 가환 도식
$$
\xymatrix{
X' \ar[rr] \ar[dr]_{f'} & & X \ar[dl]^f \\
& S
}
$$
을 생각하자. 여기서 $f' : X' \to S$와 $f : X \to S$는
Lemma \ref{more-morphisms-lemma-diagonal-picard-flat-proper}의 가정들을 만족한다.
$\mathcal{L}$을 가역 $\mathcal{O}_X$-가군이라 하고,
$\mathcal{L}'$을 $X'$으로의 당김이라 하자.
$Z \subset S$, 각각 $Z' \subset S$를
Lemma \ref{more-morphisms-lemma-diagonal-picard-flat-proper}에서
$(f, \mathcal{L})$, 각각 $(f', \mathcal{L}')$에 대해 구성한
국소 닫힌 부분스킴이라 하자. 그러면 $Z \subset Z'$이다.
$s \in Z$이고
$$
H^1(X_s, \mathcal{O}) \longrightarrow H^1(X'_s, \mathcal{O})
$$
가 단사이면, $s$의 어떤 열린 근방 $U$에 대하여
$Z \cap U = Z' \cap U$이다.
\end{lemma}

\begin{proof}
$S$를 $Z'$으로 바꾸어도 된다. $S$를 $s$의 아핀 열린 근방으로
축소하여 $\mathcal{L}' = \mathcal{O}_{X'}$이라고 가정할 수 있다.
$E = Rf_*\mathcal{L}$,
$E' = Rf'_*\mathcal{L}' = Rf'_*\mathcal{O}_{X'}$로 놓자.
이들은 그 구성이 임의의 밑변환과 가환하는 완전 복합체이다
(Derived Categories of Schemes, Lemma
\ref{perfect-lemma-flat-proper-perfect-direct-image-general}).
특히
$$
E \otimes_{\mathcal{O}_S}^\mathbf{L} \kappa(s) =
R\Gamma(X_s, \mathcal{L}_s) = R\Gamma(X_s, \mathcal{O}_{X_s})
$$
이다. 두 번째 등식은 $s \in Z$이기 때문이다.
$h_i = \dim_{\kappa(s)} H^i(X_s, \mathcal{O}_{X_s})$로 놓자.
$S$를 축소하면 $E$를 복합체
$$
\mathcal{O}_S \to \mathcal{O}_S^{\oplus h_1} \to
\mathcal{O}_S^{\oplus h_2} \to \ldots
$$
로 나타낼 수 있다. More on Algebra, Lemma
\ref{more-algebra-lemma-lift-perfect-from-residue-field}를 참조하라
(엄밀히 말하면 Derived Categories of Schemes, Lemmas
\ref{perfect-lemma-affine-compare-bounded}와
\ref{perfect-lemma-perfect-affine}도 사용한다).
마찬가지로 $E'$이 복합체
$$
\mathcal{O}_S \to \mathcal{O}_S^{\oplus h'_1} \to
\mathcal{O}_S^{\oplus h'_2} \to \ldots
$$
로 나타나며, 여기서
$h'_i = \dim_{\kappa(s)} H^i(X'_s, \mathcal{O}_{X'_s})$라고 가정할 수 있다.
코호몰로지의 함자성에 의해 $D(\mathcal{O}_S)$에서 사상
$$
E \longrightarrow E'
$$
이 존재하고 그 구성은 밑변환과 가환한다.
$E$를 나타내는 복합체가 유한 자유 가군들의 유한 복합체이고
$S$가 아핀이므로, 주어진 사상 $E \to E'$을 나타내는 복합체 사상
$$
\xymatrix{
\mathcal{O}_S \ar[r]_d \ar[d]_a &
\mathcal{O}_S^{\oplus h_1} \ar[r] \ar[d]_b &
\mathcal{O}_S^{\oplus h_2} \ar[r] \ar[d]_c & \ldots \\
\mathcal{O}_S \ar[r]^{d'} &
\mathcal{O}_S^{\oplus h'_1} \ar[r] &
\mathcal{O}_S^{\oplus h'_2} \ar[r] & \ldots
}
$$
을 택할 수 있다. $s \in Z$이므로 $\mathcal{L}_s$의 자명화 절단은
$\mathcal{L}'_s = \mathcal{O}_{X'_s}$의 자명화 절단으로 당겨진다.
따라서 $a \otimes \kappa(s)$는 동형이고, $S$를 축소하면
$a$가 동형이라고 가정할 수 있다.
끝으로 $H^1(X_s, \mathcal{O}) \to H^1(X'_s, \mathcal{O})$가
단사라는 가정에 의해, $b$를 정의하는 행렬에는
$\kappa(s)$의 영이 아닌 원소로 가는 $h_1 \times h_1$ 소행렬식이
존재한다. 따라서 $S$를 축소하여 $b$가 단사라고 가정할 수 있다.
그러나 $\mathcal{L}' = \mathcal{O}_{X'}$이므로 $d' = 0$이다.
따라서 $d = 0$이다. 이와 같이 $\mathcal{L}_s$의 자명화 절단이
$X$ 위의 $\mathcal{L}$의 절단으로 올라감을 알 수 있다.
% P05-EN-CH37-0167
직접적인 위상적 논증(생략한다)에 의해, 필요하면 $S$를 조금 더
축소했을 때 $\mathcal{L}$이 자명하다는 결론을 얻는다.
\end{proof}

\begin{lemma}
\label{more-morphisms-lemma-H1-O-multiple-picard-flat-proper}
스킴의 가환 도식 $n$개
$$
\xymatrix{
X_i \ar[rr] \ar[dr]_{f_i} & & X \ar[dl]^f \\
& S
}
$$
을 생각하자. 여기서 $f_i : X_i \to S$와 $f : X \to S$는
Lemma \ref{more-morphisms-lemma-diagonal-picard-flat-proper}의 가정들을 만족한다.
$\mathcal{L}$을 가역 $\mathcal{O}_X$-가군이라 하고,
$\mathcal{L}_i$를 $X_i$로의 당김이라 하자.
$Z \subset S$, 각각 $Z_i \subset S$를
Lemma \ref{more-morphisms-lemma-diagonal-picard-flat-proper}에서
$(f, \mathcal{L})$, 각각 $(f_i, \mathcal{L}_i)$에 대하여 구성한
국소 닫힌 부분스킴이라 하자. 그러면
$Z \subset \bigcap_{i = 1, \ldots, n} Z_i$이다. $s \in Z$이고
$$
H^1(X_s, \mathcal{O}) \longrightarrow
\bigoplus\nolimits_{i = 1, \ldots, n} H^1(X_{i, s}, \mathcal{O})
$$
가 단사이면, $s$의 어떤 열린 근방 $U$에 대하여
$Z \cap U = (\bigcap_{i = 1, \ldots, n} Z_i) \cap U$이다
(스킴론적 교집합).
\end{lemma}

\begin{proof}
이 보조정리는 Lemma \ref{more-morphisms-lemma-H1-O-picard-flat-proper}의 변형이다.
독자에게 먼저 그 증명을 읽을 것을 강하게 권한다. 이 증명은
약간 수정한 그 증명의 사본과 거의 같다.
$Z$와 $Z_i$의 (스킴값) 점들에 대한 기술로부터, 교집합을
스킴론적으로 취했을 때 $Z \subset \bigcap_{i = 1, \ldots, n} Z_i$임을
얻는다. 따라서 $S$를 스킴론적 교집합
$\bigcap_{i = 1, \ldots, n} Z_i$로 바꾸어도 된다.
$S$를 $s$의 아핀 열린 근방으로 축소하여
$i = 1, \ldots, n$에 대하여
$\mathcal{L}_i = \mathcal{O}_{X_i}$라고 가정할 수 있다.
$E = Rf_*\mathcal{L}$,
$E_i = Rf_{i, *}\mathcal{L}_i = Rf_{i, *}\mathcal{O}_{X_i}$로 놓자.
이들은 그 구성이 임의의 밑변환과 가환하는 완전 복합체이다
(Derived Categories of Schemes, Lemma
\ref{perfect-lemma-flat-proper-perfect-direct-image-general}).
특히
$$
E \otimes_{\mathcal{O}_S}^\mathbf{L} \kappa(s) =
R\Gamma(X_s, \mathcal{L}_s) = R\Gamma(X_s, \mathcal{O}_{X_s})
$$
이다. 두 번째 등식은 $s \in Z$이기 때문이다.
$h_j = \dim_{\kappa(s)} H^j(X_s, \mathcal{O}_{X_s})$로 놓자.
$S$를 축소하면 $E$를 복합체
$$
\mathcal{O}_S \to \mathcal{O}_S^{\oplus h_1} \to
\mathcal{O}_S^{\oplus h_2} \to \ldots
$$
로 나타낼 수 있다. More on Algebra, Lemma
\ref{more-algebra-lemma-lift-perfect-from-residue-field}를 참조하라
(엄밀히 말하면 Derived Categories of Schemes, Lemmas
\ref{perfect-lemma-affine-compare-bounded}와
\ref{perfect-lemma-perfect-affine}도 사용한다).
마찬가지로 $E_i$가 복합체
$$
\mathcal{O}_S \to \mathcal{O}_S^{\oplus h_{i, 1}} \to
\mathcal{O}_S^{\oplus h_{i, 2}} \to \ldots
$$
로 나타나며, 여기서
$h_{i, j} = \dim_{\kappa(s)} H^j(X_{i, s}, \mathcal{O}_{X_{i, s}})$라고
가정할 수 있다. 코호몰로지의 함자성에 의해
$D(\mathcal{O}_S)$에서 사상
$$
E \longrightarrow E_i
$$
이 존재하고 그 구성은 밑변환과 가환한다.
$E$를 나타내는 복합체가 유한 자유 가군들의 유한 복합체이고
$S$가 아핀이므로, 주어진 사상 $E \to E_i$를 나타내는 복합체 사상
$$
\xymatrix{
\mathcal{O}_S \ar[r]_d \ar[d]_{a_i} &
\mathcal{O}_S^{\oplus h_1} \ar[r] \ar[d]_{b_i} &
\mathcal{O}_S^{\oplus h_2} \ar[r] \ar[d]_{c_i} & \ldots \\
\mathcal{O}_S \ar[r]^{d_i} &
\mathcal{O}_S^{\oplus h_{i, 1}} \ar[r] &
\mathcal{O}_S^{\oplus h_{i, 2}} \ar[r] & \ldots
}
$$
을 택할 수 있다. $s \in Z$이므로 $\mathcal{L}_s$의 자명화 절단은
$\mathcal{L}_{i, s} = \mathcal{O}_{X_{i, s}}$의 자명화 절단으로
당겨진다. 따라서 $a_i \otimes \kappa(s)$는 동형이고,
$S$를 축소하면 $a_i$가 동형이라고 가정할 수 있다. 끝으로
  $H^1(X_s, \mathcal{O}) \to
\bigoplus_{i = 1, \ldots, n} H^1(X_{i, s}, \mathcal{O})$가
단사라는 가정을 사용하면, $\oplus b_i$를 정의하는 행렬에는
$\kappa(s)$의 영이 아닌 원소로 가는 $h_1 \times h_1$ 소행렬식이
존재한다. 따라서 $S$를 축소하여
% P05-EN-CH37-0168
$(b_1, \ldots, b_n) : \mathcal{O}_S^{h_1}
\to \bigoplus_{i = 1, \ldots, n} \mathcal{O}_S^{h_{i, 1}}$
가 단사라고 가정할 수 있다. 그러나 $\mathcal{L}_i = \mathcal{O}_{X_i}$이므로
$i = 1, \ldots n$에 대하여 $d_i = 0$이다.
$(b_1, \ldots, b_n) \circ d = (\oplus d_i) \circ (a_1, \ldots, a_n)$이므로
$d = 0$이다. 이와 같이 $\mathcal{L}_s$의 자명화 절단이
$X$ 위의 $\mathcal{L}$의 절단으로 올라감을 알 수 있다.
% P05-EN-CH37-0167
직접적인 위상적 논증(생략한다)에 의해, 필요하면 $S$를 조금 더
축소했을 때 $\mathcal{L}$이 자명하다는 결론을 얻는다.
\end{proof}

\begin{lemma}
\label{more-morphisms-lemma-pic-of-product}
$f : X \to S$와 $g : Y \to S$를
Lemma \ref{more-morphisms-lemma-diagonal-picard-flat-proper}의 가정들을 만족하는
스킴 사상이라 하자. $\sigma : S \to X$와 $\tau : S \to Y$를
$f$와 $g$의 단면이라 하고, $s \in S$라 하자.
$\mathcal{L}$을 $X \times_S Y$ 위의 가역층이라 하자.
$X$ 위의 $(1 \times \tau)^*\mathcal{L}$, $Y$ 위의
$(\sigma \times 1)^*\mathcal{L}$, 그리고
$\mathcal{L}|_{(X \times_S Y)_s}$가 자명하면,
어떤 $s$의 열린 근방 $U$에 대하여
$(X \times_S Y)_U$ 위에서 $\mathcal{L}$이 자명하다.
\end{lemma}

\begin{proof}
K\"unneth 정리(Varieties, Lemma \ref{varieties-lemma-kunneth})에 의해 사상
$$
H^1(X_s \times_{\Spec(\kappa(s))} Y_s, \mathcal{O}) \to
H^1(X_s, \mathcal{O}) \oplus H^1(Y_s, \mathcal{O})
$$
은 단사이다. 따라서 두 사상
$$
1 \times \tau : X \to X \times_S Y
\quad\text{이고}\quad
\sigma \times 1 : Y \to X \times_S Y
$$
에 Lemma \ref{more-morphisms-lemma-H1-O-multiple-picard-flat-proper}를 적용하면
결론을 얻는다.
\end{proof}

\begin{theorem}[정육면체 정리]
\label{more-morphisms-theorem-of-the-cube}
$S$를 스킴이라 하고, $X$, $Y$, $Z$를 $S$ 위의 스킴이라 하자.
$x : S \to X$와 $y : S \to Y$를 구조 사상들의 단면이라 하자.
$\mathcal{L}$을 $X \times_S Y \times_S Z$ 위의 가역 가군이라 하자. 다음을 가정하자.
\begin{enumerate}
\item $X \to S$와 $Y \to S$는 유한 표시인 평탄하고 고유한 사상이고,
그 올들은 기하학적으로 정역이다.
\item $x \times \text{id}_Y \times \text{id}_Z$와
$\text{id}_X \times y \times \text{id}_Z$에 의한 $\mathcal{L}$의 당김은
각각 $Y \times_S Z$와 $X \times_S Z$ 위에서 자명하다.
\item 어떤 점 $z \in Z$에 대하여 $\mathcal{L}$의
$X \times_S Y \times_S z$로의 제한은 자명하다.
\item $Z$는 연결이다.
\end{enumerate}
그러면 $\mathcal{L}$은 자명하다.
\end{theorem}

\noindent
자주 사용되는 특수한 경우는 다음과 같다.
$k$를 체라 하고, $X, Y, Z$를 $k$-유리점 $x, y, z$를 갖는
다양체라 하자. $\mathcal{L}$을 $X \times Y \times Z$ 위의
가역 가군이라 하자. 다음을 가정하자.
\begin{enumerate}
\item $\mathcal{L}$은 $x \times Y \times Z$, $X \times y \times Z$,
$X \times Y \times z$ 위에서 자명하다.
\item $X$와 $Y$는 $k$ 위에서 기하학적으로 정역이고 고유하다.
\end{enumerate}
그러면 $\mathcal{L}$은 자명하다.

\begin{proof}
% P05-EN-CH37-0169
$X \times_S Y \to S$는 유한 표시인 평탄하고 고유한 사상이며,
그 올들은 기하학적으로 정역임에 유의하자
(올에 관한 명제는 Varieties, Lemmas \ref{varieties-lemma-geometrically-integral},
\ref{varieties-lemma-bijection-irreducible-components},
and \ref{varieties-lemma-geometrically-reduced-any-base-change}를 참조하라).
Derived Categories of Schemes, Lemma
\ref{perfect-lemma-proper-flat-geom-red-connected}
에 의해 $X \to S$, $Y \to S$, 또는 $X \times_S Y \to S$에 의한
구조층의 직상은 $S$의 구조층이고, 이는 임의의 밑변환 뒤에도
참임을 알 수 있다. 따라서 사상
$$
p : X \times_S Y \times_S Z \longrightarrow Z
$$
과 가역 가군 $\mathcal{L}$에
Lemma \ref{more-morphisms-lemma-diagonal-picard-flat-proper}를 적용하면,
$\mathcal{L}|_{X \times_S Y \times_S Z'}$가 $Z'$ 위의 어떤 가역 가군
$\mathcal{N}$의 당김이 되게 하는 ``보편적'' 국소 닫힌 부분스킴
$Z' \subset Z$를 얻는다. $z$의 존재에 의해 $Z'$는 비어 있지 않다.
Lemma \ref{more-morphisms-lemma-get-a-closed}에 의해 $Z' \subset Z$는 닫힌 부분스킴이다.
$z' \in Z'$를 한 점이라 하자. $p$를 곱사상
$$
(X \times_S Z) \times_Z (Y \times_S Z) \longrightarrow Z
$$
으로 쓸 수 있음에 유의하자. 따라서 사상 $p$, 점 $z'$,
그리고 $x$와 $y$가 주는 단면
$\sigma : Z \to X \times_S Z$와 $\tau : Z \to Y \times_S Z$에
Lemma \ref{more-morphisms-lemma-pic-of-product}를 적용할 수 있다.
그러므로 $Z'$는 열린집합이다. 따라서 $Z' = Z$이고,
$Z$ 위의 어떤 가역 가군 $\mathcal{N}$에 대하여
$\mathcal{L} = p^*\mathcal{N}$이다. 사상
$x \times y \times \text{id}_Z : Z  \to X \times_S Y \times_S Z$
에 따라 당기면 한편으로 $\mathcal{N}$을 얻고, 다른 한편으로는
가정 (2)에 의해 자명한 가역 가군을 얻는다.
따라서 $\mathcal{N} = \mathcal{O}_Z$이고 증명이 끝난다.
\end{proof}









\section{극한 논법}
\label{more-morphisms-section-limits}

\noindent
스킴의 극한과 뇌터 근사에 관한 몇 가지 보조정리를 제시한다.
대부분 아핀 경우만 다룬다. 이들 중 일부는 극한에 관한 장의
보조정리들의 특수한 경우이다.

\begin{lemma}
\label{more-morphisms-lemma-Noetherian-approximation}
$f : X \to S$를 유한 표시인 아핀 스킴의 사상이라 하자.
그러면 데카르트 도식
$$
\xymatrix{
X_0 \ar[d]_{f_0} & X \ar[l]^g \ar[d]^f \\
S_0 & S \ar[l]
}
$$
이 존재하여 다음을 만족한다.
\begin{enumerate}
\item $X_0$, $S_0$는 아핀 스킴이다.
\item $S_0$는 $\mathbf{Z}$ 위에서 유한형이다.
\item $f_0$는 유한형이다.
\end{enumerate}
\end{lemma}

\begin{proof}
$S = \Spec(A)$와 $X = \Spec(B)$로 쓰자.
$f$가 유한 표시이므로 $B$는 $A$-대수로서 유한 표시이다.
Morphisms,
Lemma \ref{morphisms-lemma-locally-finite-presentation-characterize}를 참조하라.
따라서 이 보조정리는 Algebra, Lemma
\ref{algebra-lemma-limit-module-finite-presentation}에서 따른다.
\end{proof}

\begin{lemma}
\label{more-morphisms-lemma-Noetherian-approximation-module}
$f : X \to S$를 유한 표시인 아핀 스킴의 사상이라 하자.
$\mathcal{F}$를 유한 표시인 준연접 $\mathcal{O}_X$-가군이라 하자.
그러면 Lemma \ref{more-morphisms-lemma-Noetherian-approximation}에서와 같은 도식을
$g^*\mathcal{F}_0 = \mathcal{F}$를 만족하는 정합
$\mathcal{O}_{X_0}$-가군 $\mathcal{F}_0$가 존재하도록 택할 수 있다.
\end{lemma}

\begin{proof}
$S = \Spec(A)$, $X = \Spec(B)$,
$\mathcal{F} = \widetilde{M}$으로 쓰자.
$f$가 유한 표시이므로 $B$는 $A$-대수로서 유한 표시이다.
Morphisms,
Lemma \ref{morphisms-lemma-locally-finite-presentation-characterize}를 참조하라.
$\mathcal{F}$가 $\mathcal{O}_X$ 위에서 유한 표시이므로
$M$은 $B$-가군으로서 유한 표시이다.
Properties, Lemma \ref{properties-lemma-finite-presentation-module}를 참조하라.
따라서 이 보조정리는 Algebra, Lemma
\ref{algebra-lemma-limit-module-finite-presentation}에서 따른다.
\end{proof}

\begin{lemma}
\label{more-morphisms-lemma-Noetherian-approximation-flat-module}
$f : X \to S$를 유한 표시인 아핀 스킴의 사상이라 하자.
$\mathcal{F}$를 유한 표시인 준연접 $\mathcal{O}_X$-가군이며
$S$ 위에서 평탄하다고 하자. 그러면
Lemma \ref{more-morphisms-lemma-Noetherian-approximation-module}에서와 같은 도식과
층 $\mathcal{F}_0$를, 추가로 $\mathcal{F}_0$가
$S_0$ 위에서 평탄하도록 택할 수 있다.
\end{lemma}

\begin{proof}
$S = \Spec(A)$, $X = \Spec(B)$,
$\mathcal{F} = \widetilde{M}$으로 쓰자.
$f$가 유한 표시이므로 $B$는 $A$-대수로서 유한 표시이다.
Morphisms,
Lemma \ref{morphisms-lemma-locally-finite-presentation-characterize}를 참조하라.
$\mathcal{F}$가 $\mathcal{O}_X$ 위에서 유한 표시이므로
$M$은 $B$-가군으로서 유한 표시이다.
Properties, Lemma \ref{properties-lemma-finite-presentation-module}를 참조하라.
$\mathcal{F}$가 $S$ 위에서 평탄하므로 $M$은 $A$ 위에서 평탄하다.
Morphisms, Lemma \ref{morphisms-lemma-flat-module-characterize}를 참조하라.
따라서 이 보조정리는 Algebra, Lemma
\ref{algebra-lemma-flat-finite-presentation-limit-flat}에서 따른다.
\end{proof}

\begin{lemma}
\label{more-morphisms-lemma-Noetherian-approximation-flat}
$f : X \to S$를 유한 표시이고 평탄한 아핀 스킴의 사상이라 하자.
그러면 Lemma \ref{more-morphisms-lemma-Noetherian-approximation}에서와 같은 도식을
추가로 $f_0$가 평탄하도록 택할 수 있다.
\end{lemma}

\begin{proof}
이는 Lemma \ref{more-morphisms-lemma-Noetherian-approximation-flat-module}의
특수한 경우이다.
\end{proof}

\begin{lemma}
\label{more-morphisms-lemma-Noetherian-approximation-smooth}
$f : X \to S$를 매끄러운 아핀 스킴의 사상이라 하자.
그러면 Lemma \ref{more-morphisms-lemma-Noetherian-approximation}에서와 같은 도식을
추가로 $f_0$가 매끄럽도록 택할 수 있다.
\end{lemma}

\begin{proof}
$S = \Spec(A)$, $X = \Spec(B)$로 쓰자.
$f$가 매끄러우므로 $B$는 $A$-대수로서 매끄럽다.
Morphisms,
Lemma \ref{morphisms-lemma-smooth-characterize}를 참조하라.
따라서 이 보조정리는 Algebra, Lemma
\ref{algebra-lemma-finite-presentation-fs-Noetherian}에서 따른다.
\end{proof}

\begin{lemma}
\label{more-morphisms-lemma-Noetherian-approximation-geometrically-reduced}
$f : X \to S$를 유한 표시인 아핀 스킴의 사상이라 하고,
그 올들이 기하학적으로 축약이라고 하자.
그러면 Lemma \ref{more-morphisms-lemma-Noetherian-approximation}에서와 같은 도식을
추가로 $f_0$의 올들이 기하학적으로 축약이도록 택할 수 있다.
\end{lemma}

\begin{proof}
Lemma \ref{more-morphisms-lemma-Noetherian-approximation}를 적용하여 아핀 스킴의
데카르트 도식
$$
\xymatrix{
X_0 \ar[d]_{f_0} & X \ar[l]^g \ar[d]^f \\
S_0 & S \ar[l]_h
}
$$
을 얻자. 여기서 $X_0 \to S_0$는 $\mathbf{Z}$ 위에서 유한형인
스킴 사이의 유한형 사상이다.
Lemma \ref{more-morphisms-lemma-geometrically-reduced-constructible}
에 의해 $f_0$의 올이 기하학적으로 축약인 점들의 집합
$E \subset S_0$는 구성가능 부분집합이다.
Lemma \ref{more-morphisms-lemma-base-change-fibres-geometrically-reduced}
에 의해 $h(S) \subset E$이다. $S_0 = \Spec(A_0)$와
$S = \Spec(A)$로 쓰자. 유한형 $A_0$-대수들의 유향 극한
$A = \colim_i A_i$로 쓰자.
Limits, Lemma \ref{limits-lemma-limit-contained-in-constructible}
에 의해 어떤 $i$에 대하여 $\Spec(A_i) \to S_0$의 상이
$E$에 포함된다. $S_0$를 $\Spec(A_i)$로, $X_0$를
$X_0 \times_{S_0} \Spec(A_i)$로 바꾸면
$f_0$의 모든 올이 기하학적으로 축약임을 알 수 있다.
\end{proof}

\begin{lemma}
\label{more-morphisms-lemma-Noetherian-approximation-geometrically-irreducible}
$f : X \to S$를 유한 표시인 아핀 스킴의 사상이라 하고,
그 올들이 기하학적으로 기약이라고 하자.
그러면 Lemma \ref{more-morphisms-lemma-Noetherian-approximation}에서와 같은 도식을
추가로 $f_0$의 올들이 기하학적으로 기약이도록 택할 수 있다.
\end{lemma}

\begin{proof}
Lemma \ref{more-morphisms-lemma-Noetherian-approximation}를 적용하여 아핀 스킴의
데카르트 도식
$$
\xymatrix{
X_0 \ar[d]_{f_0} & X \ar[l]^g \ar[d]^f \\
S_0 & S \ar[l]_h
}
$$
을 얻자. 여기서 $X_0 \to S_0$는 $\mathbf{Z}$ 위에서 유한형인
스킴 사이의 유한형 사상이다.
Lemma \ref{more-morphisms-lemma-nr-geom-irreducible-components-constructible}
에 의해 $f_0$의 올이 기하학적으로 기약인 점들의 집합
$E \subset S_0$는 구성가능 부분집합이다.
Lemma \ref{more-morphisms-lemma-base-change-fibres-geometrically-irreducible}
에 의해 $h(S) \subset E$이다. $S_0 = \Spec(A_0)$와
$S = \Spec(A)$로 쓰자. 유한형 $A_0$-대수들의 유향 극한
$A = \colim_i A_i$로 쓰자.
Limits, Lemma \ref{limits-lemma-limit-contained-in-constructible}
에 의해 어떤 $i$에 대하여 $\Spec(A_i) \to S_0$의 상이
$E$에 포함된다. $S_0$를 $\Spec(A_i)$로, $X_0$를
$X_0 \times_{S_0} \Spec(A_i)$로 바꾸면
$f_0$의 모든 올이 기하학적으로 기약임을 알 수 있다.
\end{proof}

\begin{lemma}
\label{more-morphisms-lemma-Noetherian-approximation-geometrically-connected}
$f : X \to S$를 유한 표시인 아핀 스킴의 사상이라 하고,
그 올들이 기하학적으로 연결이라고 하자.
그러면 Lemma \ref{more-morphisms-lemma-Noetherian-approximation}에서와 같은 도식을
추가로 $f_0$의 올들이 기하학적으로 연결이도록 택할 수 있다.
\end{lemma}

\begin{proof}
Lemma \ref{more-morphisms-lemma-Noetherian-approximation}를 적용하여 아핀 스킴의
데카르트 도식
$$
\xymatrix{
X_0 \ar[d]_{f_0} & X \ar[l]^g \ar[d]^f \\
S_0 & S \ar[l]_h
}
$$
을 얻자. 여기서 $X_0 \to S_0$는 $\mathbf{Z}$ 위에서 유한형인
스킴 사이의 유한형 사상이다.
Lemma \ref{more-morphisms-lemma-nr-geom-connected-components-constructible}
에 의해 $f_0$의 올이 기하학적으로 연결인 점들의 집합
$E \subset S_0$는 구성가능 부분집합이다.
Lemma \ref{more-morphisms-lemma-base-change-fibres-geometrically-connected}
에 의해 $h(S) \subset E$이다. $S_0 = \Spec(A_0)$와
$S = \Spec(A)$로 쓰자. 유한형 $A_0$-대수들의 유향 극한
$A = \colim_i A_i$로 쓰자.
Limits, Lemma \ref{limits-lemma-limit-contained-in-constructible}
에 의해 어떤 $i$에 대하여 $\Spec(A_i) \to S_0$의 상이
$E$에 포함된다. $S_0$를 $\Spec(A_i)$로, $X_0$를
$X_0 \times_{S_0} \Spec(A_i)$로 바꾸면
$f_0$의 모든 올이 기하학적으로 연결임을 알 수 있다.
\end{proof}

\begin{lemma}
\label{more-morphisms-lemma-Noetherian-approximation-dimension-d}
$d \geq 0$을 정수라 하자. $f : X \to S$를 유한 표시인
아핀 스킴의 사상이라 하고, 그 모든 올의 차원이 $d$라고 하자.
그러면 Lemma \ref{more-morphisms-lemma-Noetherian-approximation}에서와 같은 도식을
추가로 $f_0$의 모든 올의 차원이 $d$가 되도록 택할 수 있다.
\end{lemma}

\begin{proof}
Lemma \ref{more-morphisms-lemma-Noetherian-approximation}를 적용하여 아핀 스킴의
데카르트 도식
$$
\xymatrix{
X_0 \ar[d]_{f_0} & X \ar[l]^g \ar[d]^f \\
S_0 & S \ar[l]_h
}
$$
을 얻자. 여기서 $X_0 \to S_0$는 $\mathbf{Z}$ 위에서 유한형인
스킴 사이의 유한형 사상이다.
Lemma \ref{more-morphisms-lemma-dimension-fibres-constructible}
에 의해 $f_0$의 올의 차원이 $d$인 점들의 집합
$E \subset S_0$는 구성가능 부분집합이다.
Lemma \ref{more-morphisms-lemma-base-change-dimension-fibres}
에 의해 $h(S) \subset E$이다. $S_0 = \Spec(A_0)$와
$S = \Spec(A)$로 쓰자. 유한형 $A_0$-대수들의 유향 극한
$A = \colim_i A_i$로 쓰자.
Limits, Lemma \ref{limits-lemma-limit-contained-in-constructible}
에 의해 어떤 $i$에 대하여 $\Spec(A_i) \to S_0$의 상이
$E$에 포함된다. $S_0$를 $\Spec(A_i)$로, $X_0$를
$X_0 \times_{S_0} \Spec(A_i)$로 바꾸면
$f_0$의 모든 올의 차원이 $d$임을 알 수 있다.
\end{proof}

\begin{lemma}
\label{more-morphisms-lemma-Noetherian-approximation-standard-syntomic}
$f : X \to S$를 표준 신토믹인 아핀 스킴의 사상이라 하자
(Morphisms, Definition \ref{morphisms-definition-syntomic}을 참조하라).
그러면 Lemma \ref{more-morphisms-lemma-Noetherian-approximation}에서와 같은 도식을
추가로 $f_0$가 표준 신토믹이도록 택할 수 있다.
\end{lemma}

\begin{proof}
이 보조정리는 Algebra, Lemma
\ref{algebra-lemma-relative-global-complete-intersection-Noetherian}의
사본이다.
\end{proof}

\begin{lemma}
\label{more-morphisms-lemma-Noetherian-approximation-combine}
(뇌터 근사와 성질들의 결합.)
$P$, $Q$를 밑변환에 대해 안정한 스킴 사상의 성질이라 하자.
$f : X \to S$를 유한 표시인 아핀 스킴의 사상이라 하자.
데카르트 도식
$$
\vcenter{
\xymatrix{
X_1 \ar[d]_{f_1} & X \ar[l] \ar[d]^f \\
S_1 & S \ar[l]
}
}
\quad\text{and}\quad
\vcenter{
\xymatrix{
X_2 \ar[d]_{f_2} & X \ar[l] \ar[d]^f \\
S_2 & S \ar[l]
}
}
$$
을 찾을 수 있다고 가정하자. 여기서 $S_1$, $S_2$는
$\mathbf{Z}$ 위에서 유한형이고 $f_1$, $f_2$는 유한형이며,
$f_1$은 성질 $P$를, $f_2$는 성질 $Q$를 갖는다.
그러면 아핀 스킴의 데카르트 도식
$$
\xymatrix{
X_0 \ar[d]_{f_0} & X \ar[l] \ar[d]^f \\
S_0 & S \ar[l]
}
$$
을 찾을 수 있다. 여기서 $S_0$는 $\mathbf{Z}$ 위에서 유한형이고
$f_0$는 유한형이며, $f_0$는 성질 $P$와 성질 $Q$를 모두 갖는다.
\end{lemma}

\begin{proof}
주어진 두 도식은 환들의 쌍대 데카르트 도식
$$
\vcenter{
\xymatrix{
B_1 \ar[r] & B \\
A_1 \ar[u] \ar[r] & A \ar[u]
}
}
\quad\text{and}\quad
\vcenter{
\xymatrix{
B_2 \ar[r] & B \\
A_2 \ar[u] \ar[r] & A \ar[u]
}
}
$$
에 대응한다. $A_0 \subset A$를 $A_1 \to A$와 $A_2 \to A$의 상을
모두 포함하는 $A$의 유한형 $\mathbf{Z}$-부분대수라 하자.
가정에 의해 $A_1$과 $A_2$가 모두 $\mathbf{Z}$ 위에서 유한형이므로
이러한 부분대수가 존재한다. 환
$B_{0, 1} = B_1 \otimes_{A_1} A_0$와
$B_{0, 2} = B_2 \otimes_{A_2} A_0$는 다음을 만족하는
유한형 $A_0$-대수임에 유의하자.
$B_{0, 1} \otimes_{A_0} A \cong B \cong B_{0, 2} \otimes_{A_0} A$
이 동형들은 $A$-대수의 동형이다. $A$는 그 유한형
$A_0$-부분대수들의 유향 극한이므로
Limits, Lemma \ref{limits-lemma-descend-finite-presentation}
에 의해 $A_0$를 확대하여 $A_0$-대수의 동형
$B_{0, 1} \cong B_{0, 2}$가 존재한다고 가정할 수 있다.
성질 $P$와 $Q$가 밑변환에 대해 안정하다고 가정했으므로,
$S_0 = \Spec(A_0)$와
$$
X_0 = X_1 \times_{S_1} S_0 =
\Spec(B_{0, 1}) \cong \Spec(B_{0, 2}) = X_2 \times_{S_2} S_0
$$
로 놓으면 된다.
\end{proof}












\section{\'Etale 근방}
\label{more-morphisms-section-etale-neighbourhoods}

\noindent
사상의 몇몇 성질은 \'etale 밑변환을 한 뒤에 더 쉽게 연구할 수 있음이
밝혀진다. 다음 용어를 도입하는 것이 편리하다.

\begin{definition}
\label{more-morphisms-definition-etale-neighbourhood}
$S$를 스킴이라 하고 $s \in S$를 한 점이라 하자.
\begin{enumerate}
\item {\it $(S, s)$의 \'etale 근방}은 쌍 $(U, u)$와
$\varphi(u) = s$를 만족하는 스킴의 \'etale 사상
$\varphi : U \to S$로 이루어진다.
\item $(S, s)$의 {\it \'etale 근방의 사상}
$f : (V, v) \to (U, u)$은 단지 $f(v) = u$를 만족하는
$S$-스킴의 사상 $f : V \to U$이다.
\item {\it 기본 \'etale 근방}은 $\kappa(s) = \kappa(u)$를 만족하는
\'etale 근방 $\varphi : (U, u) \to (S, s)$이다.
\end{enumerate}
\end{definition}

\noindent
% P05-EN-CH37-0171
기본 \'etale 근방이라는 개념은 문헌에서 여러 다른 이름으로 불린다.
예를 들어 이를 ``\'etale 근방''(\cite[Page 36]{Milne}) 또는
``강한 \'etale''(\cite[Page 108]{KPR})이라고도 한다.
여기서는 논문 \cite{GruRay}의 관례를 따라
기본 \'etale 근방이라 부른다.

\medskip\noindent
$f : (V, v) \to (U, u)$가 \'etale 근방의 사상이면
$f$는 자동으로 \'etale이다. Morphisms, Lemma
\ref{morphisms-lemma-etale-permanence}를 참조하라.
따라서 이는 $(V, v)$를 $(U, u)$의 \'etale 근방으로 만든다.
역으로, \'etale 사상의 합성은 \'etale이므로(Morphisms, Lemma
\ref{morphisms-lemma-composition-etale}), $(U, u)$의 임의의
\'etale 근방 $(V, v)$는 $(S, s)$의 \'etale 근방이기도 하다.
또한 $U \subset S$가 $s$의 열린 근방이면
$(U, s) \to (S, s)$는 \'etale 근방이다.
이는 열린 몰입이 \'etale이라는 사실(Morphisms, Lemma
\ref{morphisms-lemma-open-immersion-etale})에서 따른다.
이 절 전체에서 이러한 언급들을 별도 설명 없이 사용한다.

\medskip\noindent
$(U, u) \to (S, s)$가 \'etale 근방이면
$\kappa(u)/\kappa(s)$는 유한 분리 확대임에 유의하자.
Morphisms, Lemma \ref{morphisms-lemma-etale-at-point}를 참조하라.

\begin{lemma}
\label{more-morphisms-lemma-realize-prescribed-residue-field-extension-etale}
$S$를 스킴이라 하고 $s \in S$라 하자.
$k/\kappa(s)$를 유한 분리 확대체라 하자.
그러면 확대체 $\kappa(u)/\kappa(s)$가 $k/\kappa(s)$와 동형인
\'etale 근방 $(U, u) \to (S, s)$가 존재한다.
\end{lemma}

\begin{proof}
$S$가 아핀이라고 가정해도 된다. 이 경우 보조정리는
Algebra, Lemma
\ref{algebra-lemma-make-etale-map-prescribed-residue-field}에서 따른다.
\end{proof}

\begin{lemma}
\label{more-morphisms-lemma-etale-neighbourhoods-not-quite-filtered}
$S$를 스킴이라 하고 $s$를 $S$의 한 점이라 하자.
\'etale 근방의 범주는 다음 성질들을 갖는다.
\begin{enumerate}
\item $(U_i, u_i)_{i=1, 2}$를 $S$에서 $s$의 두 \'etale 근방이라 하자.
그러면 세 번째 \'etale 근방 $(U, u)$와 사상
$(U, u) \to (U_i, u_i)$, $i = 1, 2$가 존재한다.
\item $h_1, h_2: (U, u) \to (U', u')$를 $s$의 \'etale 근방 사이의
두 사상이라 하자. $h_1$, $h_2$가 잉여체의 같은 사상
$\kappa(u') \to \kappa(u)$를 유도한다고 가정하자.
그러면 \'etale 근방 $(U'', u'')$와 $h_1$, $h_2$를 같게 만드는 사상
$h : (U'', u'') \to (U, u)$가 존재한다. 즉,
$h_1 \circ h = h_2 \circ h$이다.
\end{enumerate}
\end{lemma}

\begin{proof}
(1)에 대하여 올곱 $U = U_1 \times_S U_2$를 생각하자.
\'etale 사상은 밑변환에 의해 보존되므로 이는 $U_1$과 $U_2$
양쪽 위에서 \'etale이다. Morphisms, Lemma
\ref{morphisms-lemma-base-change-etale}를 참조하라.
예를 들어 Schemes, Lemma \ref{schemes-lemma-points-fibre-product}의
올곱의 점에 대한 기술에 의해, $u_1$과 $u_2$ 양쪽으로 가는
$U$의 점이 존재한다. (2)에 대하여 $U''$를 올곱
$$
\xymatrix{
U'' \ar[r] \ar[d] & U \ar[d]^{(h_1, h_2)} \\
U' \ar[r]^-\Delta & U' \times_S U'.
}
$$
으로 정의하자.
% P05-EN-CH37-0172
$h_1$과 $h_2$가 잉여체의 같은 사상
$\kappa(u') \to \kappa(u)$를 유도하므로, $u'$ 위에 놓이고
$\kappa(u'') = \kappa(u')$를 만족하는 점 $u'' \in U''$가 존재한다.
특히 $U'' \not = \emptyset$이다. 또한 $U'$가 $S$ 위에서
\'etale이므로 올곱 $U'\times_S U'$도 그러하다(Morphisms, Lemmas
\ref{morphisms-lemma-base-change-etale}와
\ref{morphisms-lemma-composition-etale}를 참조하라).
따라서 Morphisms, Lemma \ref{morphisms-lemma-etale-permanence}에 의해
수직 화살표 $(h_1, h_2)$는 \'etale이다.
그러므로 밑변환에 의해 $U''$는 $U'$ 위에서 \'etale이고, 따라서
$S$ 위에서도 \'etale이다(\'etale 사상의 합성은 \'etale이기 때문이다).
따라서 $(U'', u'')$가 문제의 해이다.
\end{proof}

\begin{lemma}
\label{more-morphisms-lemma-elementary-etale-neighbourhoods}
$S$를 스킴이라 하고 $s$를 $S$의 한 점이라 하자.
$(S, s)$의 기본 \'etale 근방의 범주는 공여과이다
(Categories, Definition \ref{categories-definition-codirected}을 참조하라).
\end{lemma}

\begin{proof}
이는 정의와 Lemma
\ref{more-morphisms-lemma-etale-neighbourhoods-not-quite-filtered}에서 바로 따른다.
\end{proof}

\begin{lemma}
\label{more-morphisms-lemma-describe-henselization}
$S$를 스킴이라 하고 $s \in S$라 하자. 그러면
$$
\mathcal{O}_{S, s}^h =
\colim_{(U, u)} \mathcal{O}(U)
$$
이다. 여기서 여극한은 $(S, s)$의 기본 \'etale 근방 $(U, u)$의
범주의 반대범주인 여과 범주에 걸쳐 취한다.
\end{lemma}

\begin{proof}
$\Spec(A) \subset S$를 $s$의 아핀 근방이라 하고,
$\mathfrak p \subset A$를 $s$에 대응하는 소아이디얼이라 하자.
이 선택들에 대하여 표준 동형
$\mathcal{O}_{S, s} = A_{\mathfrak p}$와
$\kappa(s) = \kappa(\mathfrak p)$가 있다.
기본 \'etale 근방의 공종계는 $U$가 아핀이고 $U \to S$가
$\Spec(A)$를 통해 인수분해되는 기본 \'etale 근방 $(U, u)$들로 주어진다.
달리 말해 우변은 $\colim_{(B, \mathfrak q)} B$와 같다.
여기서 여극한은 $\mathfrak p$ 위에 놓이고
$\kappa(\mathfrak p) = \kappa(\mathfrak q)$를 만족하는
소아이디얼 $\mathfrak q$가 주어진 \'etale $A$-대수 $B$들에 걸쳐 취한다.
따라서 이 보조정리는 Algebra, Lemma
\ref{algebra-lemma-henselization-different}에서 따른다.
\end{proof}

\noindent
올 위 점의 \'etale 근방을 전체 공간으로 올릴 수 있다.

\begin{lemma}
\label{more-morphisms-lemma-lift-etale-neighbourhood-fibre}
\begin{slogan}
올 위 점의 \'etale 근방을 전체 공간으로 올린다.
\end{slogan}
$X \to S$를 스킴의 사상이라 하고, $x \in X$의 상을 $s \in S$라 하자.
$(V, v) \to (X_s, x)$를 \'etale 근방이라 하자.
그러면 \'etale 근방 $(U, u) \to (X, x)$가 존재하여,
$(X_s, x)$의 \'etale 근방의 사상 $(U_s, u) \to (V, v)$ 중
열린 몰입인 것이 존재한다.
\end{lemma}

\begin{proof}
$X$, $V$, $S$가 아핀이라고 가정해도 된다.
% P05-EN-CH37-0173
사상 $X \to S$는 $A \to B$로, 점 $x$는 소아이디얼
$\mathfrak q \subset B$로, 점 $s$는 $\mathfrak p = A \cap \mathfrak q$로,
사상 $V \to X_s$는 $B \otimes_A \kappa(\mathfrak p) \to C$로
주어진다고 하자. $\kappa(\mathfrak p)$는 $A/\mathfrak p$의
국소화이므로, 다음을 만족하는 $f \in A$, $f \not \in \mathfrak p$와
\'etale 환 사상 $B \otimes_A (A/\mathfrak p)_f \to D$가 존재한다.
$$
C = (B \otimes_A \kappa(\mathfrak p))
\otimes_{B \otimes_A (A/\mathfrak p)_f} D
$$
Algebra, Lemma \ref{algebra-lemma-etale}의 (9)를 참조하라.
$A$를 $A_f$로, $B$를 $B_f$로 바꾸어
$D$가 $B \otimes_A A/\mathfrak p = B/\mathfrak p B$ 위에서
\'etale이라고 가정할 수 있다. 이제 Algebra, Lemma
\ref{algebra-lemma-lift-etale}를 적용할 수 있다.
이로써 보조정리가 증명된다.
\end{proof}














\section{\'Etale 근방과 분지}
\label{more-morphisms-section-etale-nbhds-branches}

\noindent
한 점에서 스킴의 (기하학적) 분지 수는
Properties, Section \ref{properties-section-unibranch}에서 정의했다.
Varieties, Section \ref{varieties-section-number-of-branches}에서는
이를 정규화 사상의 올과 관련지었다.
이 절에서는 이 수를 \'etale 근방으로 특징짓는다.

\begin{lemma}
\label{more-morphisms-lemma-nr-minimal-primes}
$R = \colim R_i$를 전이 사상들이 충실히 평탄한 환의 유향계의
여극한이라 하자. 그러면 $R$의 극소 소아이디얼 수를
$\{0, 1, 2, \ldots, \infty\}$의 원소로 보았을 때,
이는 $R_i$들의 극소 소아이디얼 수의 상한이다.
\end{lemma}

\begin{proof}
$A \to B$가 평탄한 환 사상이면, 하강 정리에 의해
$\Spec(B) \to \Spec(A)$는 극소 소아이디얼을 극소 소아이디얼로 보낸다
(Algebra, Lemma \ref{algebra-lemma-flat-going-down}).
$A \to B$가 충실히 평탄하면 모든 극소 소아이디얼은
어떤 극소 소아이디얼의 상이다(Algebra, Lemmas
\ref{algebra-lemma-ff-rings}와
\ref{algebra-lemma-minimal-prime-image-minimal-prime}에 의해).
% P05-EN-CH37-0174
따라서 $i \leq i'$이면 $R_i$의 극소 소아이디얼 수는
$R_{i'}$의 극소 소아이디얼 수 이상이다($\geq$).
Algebra, Lemma \ref{algebra-lemma-colimit-faithfully-flat}에 의해
각 사상 $R_i \to R$은 충실히 평탄하고, 또한 $R$의 극소 소아이디얼 수는
$R_i$의 극소 소아이디얼 수보다 $\geq$임을 알 수 있다.
끝으로 $\mathfrak q_1, \ldots, \mathfrak q_n$을 $R$의 서로 다른
극소 소아이디얼들이라 하자. 그러면
$R_i \cap \mathfrak q_1, \ldots, R_i \cap \mathfrak q_n$이
서로 다른(집합으로서, 따라서 소아이디얼로서) 어떤 $i$를 찾을 수 있다.
이로부터 보조정리가 따른다.
\end{proof}

\begin{lemma}
\label{more-morphisms-lemma-nr-branches}
$X$를 스킴이라 하고 $x \in X$를 한 점이라 하자. 그러면 다음이 성립한다.
\begin{enumerate}
\item $x$에서 $X$의 분지 수는 기본 \'etale 근방
$(U, u) \to (X, x)$들에 걸쳐 취한, $u$를 지나는 $U$의
기약 성분 수의 상한과 같다.
\item $x$에서 $X$의 기하학적 분지 수는 \'etale 근방
$(U, u) \to (X, x)$들에 걸쳐 취한, $u$를 지나는 $U$의
기약 성분 수의 상한과 같다.
\item $X$가 $x$에서 단일분지일 필요충분조건은, 모든 기본 \'etale 근방
$(U, u) \to (X, x)$에 대하여 $u$를 지나는 $U$의 기약 성분이
정확히 하나인 것이다.
\item $X$가 $x$에서 기하학적으로 단일분지일 필요충분조건은,
모든 \'etale 근방 $(U, u) \to (X, x)$에 대하여
$u$를 지나는 $U$의 기약 성분이 정확히 하나인 것이다.
\end{enumerate}
\end{lemma}

\begin{proof}
(3)과 (4)는 Properties, Lemma
\ref{properties-lemma-number-of-branches-1}를 통해 (1)과 (2)에서 따른다.

\medskip\noindent
(1)의 증명. $\Spec(A)$를 $x$의 아핀 열린 근방이라 하고,
$\mathfrak p \subset A$를 $x$에 대응하는 소아이디얼이라 하자.
$X$를 $\Spec(A)$로 바꾸어도 된다. 아핀 기본 \'etale 근방
$(U, u)$들은 공종 부분계를 이루므로 상한에서는 이들만 고려하면 충분하다.
헨젤화 $A_\mathfrak p^h$는, $B$가 \'etale $A$-대수이고
$\mathfrak q$가 $\mathfrak p$ 위에 놓이며
$\kappa(\mathfrak q) = \kappa(\mathfrak p)$를 만족하는 소아이디얼인
쌍 $(B, \mathfrak q)$의 범주에 걸쳐 취한 환 $B_\mathfrak q$들의
여극한임을 상기하자. Algebra, Lemma
\ref{algebra-lemma-henselization-different}를 참조하라.
환과 아핀 스킴 사이의 대응에 의해 이 쌍 $(B, \mathfrak q)$들은
아핀 기본 \'etale 근방 $(U, u)$들과 정확히 대응한다.
$\mathfrak q$를 지나는 $\Spec(B)$의 기약 성분들은 정확히
$B_\mathfrak q$의 극소 소아이디얼들이다.
Properties, Definition \ref{properties-definition-number-of-branches}에 의해
$A_\mathfrak p^h$의 극소 소아이디얼 수는 $x$에서 $X$의 분지 수이다.
이 계의 전이 사상 $B_\mathfrak q \to B'_{\mathfrak q'}$는 모두 평탄하다.
평탄한 국소 환 사상은 충실히 평탄하므로(Algebra, Lemma
\ref{algebra-lemma-local-flat-ff}), 보조정리는
Lemma \ref{more-morphisms-lemma-nr-minimal-primes}에서 따른다.

\medskip\noindent
(2)의 증명. Algebra, Lemma
\ref{algebra-lemma-strict-henselization-different}를 사용한다는 점을
제외하면 (1)의 증명과 같다. 작은 차이가 하나 있다.
$\kappa(x)$의 분리 대수적 폐포 $\kappa^{sep}$가 주어졌을 때,
$\kappa(u)/\kappa(x)$가 유한 분리이므로(Morphisms, Lemma
\ref{morphisms-lemma-etale-at-point}), 모든 \'etale 근방
$(U, u)$에 대하여 $\kappa(x)$-매장
$\phi : \kappa(u) \to \kappa^{sep}$를 택할 수 있다.
따라서 $(U, u) \to (X, x)$가 아핀 \'etale 근방이고
$\phi : \kappa(u) \to \kappa^{sep}$가 $\kappa(x)$-매장인
모든 삼중쌍 $(U, u, \phi)$에 걸쳐 상한을 취할 수 있다.
이 삼중쌍들은 Algebra, Lemma
\ref{algebra-lemma-strict-henselization-different}의 삼중쌍들과
정확히 대응하며, 나머지 증명은 완전히 같다.
\end{proof}

\noindent
위 보조정리의 상대적 변형이 필요하다.

\begin{lemma}
\label{more-morphisms-lemma-nr-branches-fibre}
$X \to S$를 스킴의 사상이라 하고 $x \in X$를 상이 $s$인 점이라 하자.
그러면 다음이 성립한다.
\begin{enumerate}
\item $x$에서 올 $X_s$의 분지 수는 기본 \'etale 근방
$(U, u) \to (X, x)$들에 걸쳐 취한, $u$를 지나는 올 $U_s$의
기약 성분 수의 상한과 같다.
\item $x$에서 올 $X_s$의 기하학적 분지 수는 \'etale 근방
$(U, u) \to (X, x)$들에 걸쳐 취한, $u$를 지나는 올 $U_s$의
기약 성분 수의 상한과 같다.
\item 올 $X_s$가 $x$에서 단일분지일 필요충분조건은, 모든 기본
\'etale 근방 $(U, u) \to (X, x)$에 대하여 $u$를 지나는 올 $U_s$의
기약 성분이 정확히 하나인 것이다.
\item % P05-EN-CH37-0175
$X$가 $x$에서 기하학적으로 단일분지일 필요충분조건은, 모든
\'etale 근방 $(U, u) \to (X, x)$에 대하여 $u$를 지나는 $U_s$의
기약 성분이 정확히 하나인 것이다.
\end{enumerate}
\end{lemma}

\begin{proof}
Lemmas \ref{more-morphisms-lemma-nr-branches}와
\ref{more-morphisms-lemma-lift-etale-neighbourhood-fibre}를 결합하라.
\end{proof}

\begin{lemma}
\label{more-morphisms-lemma-number-of-branches-and-smooth}
$X \to S$를 매끄러운 스킴 사상이라 하고,
$x \in X$의 상을 $s \in S$라 하자. 그러면 다음이 성립한다.
\begin{enumerate}
\item $x$에서 $X$의 기하학적 분지 수는
$s$에서 $S$의 기하학적 분지 수와 같다.
\item $\kappa(x)/\kappa(s)$가 순수 비분리\footnote{실제로는
$\kappa(x)$가 $\kappa(s)$ 위에서 기하학적으로 기약이면 충분하다.
이 사실이 필요해지면 상세한 증명을 추가할 것이다.}
확대체이면, $x$에서 $X$의 분지 수는 $s$에서 $S$의 분지 수와 같다.
% P05-EN-CH37-0383
\end{enumerate}
\end{lemma}

\begin{proof}
이는 More on Algebra, Lemma
\ref{more-algebra-lemma-invariance-number-branches-smooth}
와 정의에서 바로 따른다.
\end{proof}




\section{비분기 사상과 \'etale 사상}
\label{more-morphisms-section-unramified-is-etale}

\noindent
비분기 사상이 자동으로 \'etale인 경우가 있다.

\begin{lemma}
\label{more-morphisms-lemma-unramified-dominant-unibranch-is-etale}
$f : X \to Y$를 스킴의 사상이라 하자. $x \in X$의 상을
$y \in Y$라 하자. 다음을 가정하자.
\begin{enumerate}
\item $Y$는 정수적이고 $y$에서 기하학적으로 단일분지이다.
\item $f$는 국소 유한형이다.
\item $f(x')$가 $Y$의 일반점인 특수화 $x' \leadsto x$가 존재한다.
\item $f$는 $x$에서 비분기이다.
\end{enumerate}
그러면 $f$는 $x$에서 \'etale이다.
\end{lemma}

\begin{proof}
$X$와 $Y$를 각각 $x$와 $y$의 적절한 아핀 열린 근방으로 바꾸어도 된다.
그러면 $Y$는 정역 $A$의 스펙트럼이고 $X$는 유한형 $A$-대수 $B$의
스펙트럼이다. $\mathfrak q \subset B$를 $x$에 대응하는 소 아이디얼이라 하고,
$\mathfrak p \subset A$를 $y$에 대응하는 소 아이디얼이라 하자.
국소환 $A_\mathfrak p = \mathcal{O}_{Y, y}$는 기하학적으로 단일분지이다.
환 사상 $A \to B$는 $\mathfrak q$에서 비분기이다. 또한 (3)의 점 $x'$은
$A \cap \mathfrak q' = (0)$을 만족하는 소 아이디얼
$\mathfrak q' \subset \mathfrak q$에 대응한다. 따라서
$A_\mathfrak p \to B_\mathfrak q$는 단사이다. 결론은 More on Algebra, Lemma
\ref{more-algebra-lemma-local-unramified-extension-unibranch-domain-is-etale}에서 따른다.
\end{proof}

\begin{lemma}
\label{more-morphisms-lemma-global-unramified-dominant-unibranch-is-etale}
\begin{reference}
\cite[Expose I, Corollary 9.11]{SGA1}
\end{reference}
$f : X \to Y$를 스킴의 사상이라 하자. 다음을 가정하자.
\begin{enumerate}
\item $Y$는 정수적이고 기하학적으로 단일분지이다.
\item $X$의 기약 성분 가운데 적어도 하나는 $Y$ 위에서 우세하다.
\item $f$는 비분기이다.
\item $X$는 연결이다.
\end{enumerate}
그러면 $f$는 \'etale이고 $X$는 기약이다.
\end{lemma}

\begin{proof}
$X' \subset X$를 $Y$ 위에서 우세한 기약 성분이라 하자. 이는 $X'$의
일반점이 $Y$의 일반점으로 간다는 뜻이다(예를 들어 Morphisms, Lemma
\ref{morphisms-lemma-dominant-finite-number-irreducible-components}을 보라).
Lemma \ref{more-morphisms-lemma-unramified-dominant-unibranch-is-etale}에 의해 $f$는
$X'$의 모든 점에서 \'etale이다. 특히 $f$가 \'etale인 열린 부분스킴
$U \subset X$는 $X'$을 포함한다. $U$의 모든 준콤팩트 열린집합은 유한 개의
기약 성분을 가짐에 유의하라. Descent, Lemma
\ref{descent-lemma-locally-finite-nr-irred-local-fppf}를 보라.
한편 $Y$가 기하학적으로 단일분지이고 $U$가 $Y$ 위에서 \'etale이므로,
스킴 $U$는 기하학적으로 단일분지이다. 특히 $U$의 각 점을 지나는
기약 성분은 많아야 하나이다. 이제 간단한 위상적 논증으로
$X' \subset U$가 열리고 닫혀 있음을 알 수 있다. 그러면 물론 $X'$은
$X$에서도 열리고 닫혀 있으며, 연결성에 의해 원하는 대로
$X = U = X'$을 얻는다.
\end{proof}

\begin{lemma}
\label{more-morphisms-lemma-weird-permanence-etale}
$f : X \to Y$와 $g : Y \to Z$를 스킴의 사상이라 하자.
$x \in X$의 상을 $y \in Y$라 하자. 다음을 가정하자.
\begin{enumerate}
\item $Y$는 정수적이고 $y$에서 기하학적으로 단일분지이다.
\item $g \circ f$는 $x$에서 \'etale이다.
\item $f(x')$가 $Y$의 일반점인 특수화 $x' \leadsto x$가 존재한다.
\end{enumerate}
그러면 $f$는 $x$에서 \'etale이고 $g$는 $y$에서 \'etale이다.
\end{lemma}

\begin{proof}
$X$를 $x$의 열린 근방으로 바꾸어 $g \circ f$가 \'etale이라고
가정할 수 있다. 그러면 Morphisms, Lemmas
\ref{morphisms-lemma-unramified-permanence}
와 \ref{morphisms-lemma-etale-smooth-unramified}에 의해 $f$는 비분기이다.
따라서 Lemma \ref{more-morphisms-lemma-unramified-dominant-unibranch-is-etale}에 의해
$f$는 $x$에서 \'etale이다. 이제 \'etale 내림에 의해 $g$가 $y$에서
\'etale임을 알 수 있다. 예를 들어 Descent, Lemma
\ref{descent-lemma-syntomic-smooth-etale-permanence}를 보라.
\end{proof}

\begin{lemma}
\label{more-morphisms-lemma-global-weird-permanence-etale}
$f : X \to Y$와 $g : Y \to Z$를 스킴의 사상이라 하자.
다음을 가정하자.
\begin{enumerate}
\item $Y$는 정수적이고 기하학적으로 단일분지이다.
\item $g \circ f$는 \'etale이다.
\item $X$의 모든 기약 성분은 $Y$ 위에서 우세하다.
\end{enumerate}
그러면 $f$는 \'etale이고 $g$는 $f$의 상에 속하는 모든 점에서
\'etale이다.
\end{lemma}

\begin{proof}
이는 점별 판인 Lemma \ref{more-morphisms-lemma-weird-permanence-etale}에서 바로 따른다.
\end{proof}










\section{매끄러운 사상의 절단}
\label{more-morphisms-section-etale-over-smooth}

\noindent
이 절에서는 대략 스킴 $S$의 매끄러운 덮개를 \'etale 덮개로
세분할 수 있다는 결과를 설명한다. 그 증명 기법은 절단 논증에 의존한다.

\begin{lemma}
\label{more-morphisms-lemma-slice-smooth-given-element}
$f : X \to S$를 스킴의 사상이라 하자.
$x \in X$를 상이 $s \in S$인 점이라 하자.
$h \in \mathfrak m_x \subset \mathcal{O}_{X, x}$라 하자.
다음을 가정하자.
\begin{enumerate}
\item $f$는 $x$에서 매끄럽다.
\item $\text{d}h$의 다음 공간에서의 상 $\text{d}\overline{h}$는
$$
\Omega_{X_s/s, x} \otimes_{\mathcal{O}_{X_s, x}} \kappa(x) =
\Omega_{X/S, x} \otimes_{\mathcal{O}_{X, x}} \kappa(x)
$$
영이 아니다.
\end{enumerate}
그러면 $x$의 아핀 열린 근방 $U \subset X$가 존재하여, $h$는
$h \in \Gamma(U, \mathcal{O}_U)$에서 오고, $D = V(h)$는 $U$의
유효 Cartier 인자이며 $x \in D$이고 $D \to S$는 매끄럽다.
\end{lemma}

\begin{proof}
$f$가 $x$에서 매끄러우므로 $X$를 $x$의 열린 근방으로 바꾸어
$f$가 매끄럽다고 가정할 수 있다. 특히 $f$는 평탄하고 국소 유한 표시이다.
Lemma \ref{more-morphisms-lemma-slice-given-element}에 의해, $x$의 열린 근방
$U \subset X$가 존재하여 $h$가 $h \in \Gamma(U, \mathcal{O}_U)$에서
오고, $D = V(h)$가 $U$의 유효 Cartier 인자이며 $x \in D$이고,
$D \to S$가 평탄하고 유한 표시임을 이미 안다. Morphisms, Lemma
\ref{morphisms-lemma-differentials-relative-immersion}에 의해 짧은 완전열
$$
\mathcal{C}_{D/U} \to i^*\Omega_{U/S} \to \Omega_{D/S} \to 0
$$
을 얻는다. 여기서 $i : D \to U$는 닫힌 몰입이고
$\mathcal{C}_{D/U}$는 $U$ 안의 $D$의 여법선 층이다. $D$는
$h \in \Gamma(U, \mathcal{O}_U)$로 잘라낸 유효 Cartier 인자이므로
% P05-EN-CH37-0176
$\mathcal{C}_{D/U} = h \cdot \mathcal{O}_S$이다. $U \to S$가 매끄러우므로
층 $\Omega_{U/S}$는 유한 국소 자유이고, 따라서 그 당김
$i^*\Omega_{U/S}$도 유한 국소 자유이다. 이 완전열의 첫 화살표는 자유 생성원
$h$를 $i^*\Omega_{U/S}$의 단면 $\text{d}h|_D$로 보내며, 이 단면의 올
$\Omega_{U/S, x} \otimes \kappa(x)$에서의 값은 가정에 의해 영이 아니다.
$\otimes \kappa(x)$의 오른쪽 완전성에 의해 다음을 얻는다.
$$
\dim_{\kappa(x)} \left( \Omega_{D/S, x} \otimes \kappa(x) \right)
=
\dim_{\kappa(x)} \left( \Omega_{U/S, x} \otimes \kappa(x) \right) - 1.
$$
Morphisms, Lemma \ref{morphisms-lemma-smooth-at-point}에 의해
$\Omega_{U/S, x} \otimes \kappa(x)$는 많아야 $\dim_x(U_s)$개의 원소로
생성된다. 위의 공식에 의해 $\Omega_{D/S, x} \otimes \kappa(x)$는 많아야
$\dim_x(U_s) - 1$개의 원소로 생성된다. 예를 들어 Algebra, Lemma
\ref{algebra-lemma-one-equation}에 의해
$\dim(\mathcal{O}_{D_s, x}) = \dim(\mathcal{O}_{U_s, x}) - 1$이고
($D_s \subset U_s$도 유효 Cartier이다. Divisors, Lemma
\ref{divisors-lemma-relative-Cartier}를 보라), 이어서 Morphisms, Lemma
\ref{morphisms-lemma-dimension-fibre-at-a-point}을 사용하면
$\dim_x(D_s) = \dim_x(U_s) - 1$임에 유의하라. 따라서
$\Omega_{D/S, x} \otimes \kappa(x)$는 많아야 $\dim_x(D_s)$개의
원소로 생성되며, Morphisms, Lemma \ref{morphisms-lemma-smooth-at-point}을
다시 적용하면 $D \to S$가 $x$에서 매끄러움을 얻는다. $U$를 줄이면
$D \to S$가 매끄럽게 되어 결론을 얻는다.
\end{proof}

\begin{lemma}
\label{more-morphisms-lemma-slice-smooth-once}
$f : X \to S$를 스킴의 사상이라 하자.
$x \in X$를 상이 $s \in S$인 점이라 하자.
다음을 가정하자.
\begin{enumerate}
\item $f$는 $x$에서 매끄럽다.
\item 사상
$$
\Omega_{X_s/s, x} \otimes_{\mathcal{O}_{X_s, x}} \kappa(x)
\longrightarrow
\Omega_{\kappa(x)/\kappa(s)}
$$
의 핵은 영이 아니다.
\end{enumerate}
그러면 $x$의 아핀 열린 근방 $U \subset X$와 $x$를 포함하는
유효 Cartier 인자 $D \subset U$가 존재하여 $D \to S$는 매끄럽다.
\end{lemma}

\begin{proof}
$k = \kappa(s)$ 및 $R = \mathcal{O}_{X_s, x}$라 쓰자.
$\mathfrak m$을 $R$의 극대 아이디얼이라 하고
$\kappa = R/\mathfrak m$라 놓으면 $\kappa = \kappa(x)$이다.
미분 가군을 만드는 것은 국소화와 가환하므로(Algebra, Lemma
\ref{algebra-lemma-differentials-localize} 참조)
$\Omega_{X_s/s, x} = \Omega_{R/k}$이다. Algebra, Lemma
\ref{algebra-lemma-differential-seq}에 의해 완전열
$$
\mathfrak m/\mathfrak m^2 \xrightarrow{\text{d}}
\Omega_{R/k} \otimes_R \kappa \to
\Omega_{\kappa/k} \to 0.
$$
이 존재한다. 따라서 (2)가 성립하면
$\text{d}\overline{h}$가 영이 아닌 원소
$\overline{h} \in \mathfrak m$이 존재한다. $\overline{h}$의 올림
$h \in \mathcal{O}_{X, x}$를 택하고
Lemma \ref{more-morphisms-lemma-slice-smooth-given-element}을 적용한다.
\end{proof}

\begin{remark}
\label{more-morphisms-remark-necessary-condition-slice-smooth}
Lemma \ref{more-morphisms-lemma-slice-smooth-once}의 둘째 조건은 $x$가 양의 차원인
올의 닫힌점일 때에도 필요하다. 다음이 한 예이다. 표수가 $p > 0$인
비완전체 $k$를 택하고, $p$제곱이 아닌 원소 $a \in k$를 택하자.
$\mathfrak m = (x, y^p - a) \subset k[x, y]$라 하자. 이는
$X = \mathbf{A}^2_k$의 닫힌점 $w$에 대응한다.
$S = \mathbf{A}^1_k$라 놓고, $f : X \to S$를
$k[x] \to k[x, y]$에 대응하는 사상이라 하자. 그러면 $g$가 \'etale이고
$w$가 $h$의 상에 속하는 다음 꼴의 가환 도식은 존재하지 않는다.
$$
\xymatrix{
S' \ar[rr]_h \ar[rd]_g & & X \ar[ld]^f \\
& S
}
$$
실제로 잉여체 확대 $\kappa(w)/\kappa(f(w))$는 순수 비분리이지만,
$g(s') = f(w)$인 임의의 $s' \in S'$에 대해 확대
$\kappa(s')/\kappa(f(w))$는 분리여야 하므로 이는 명백하다.
\end{remark}

\noindent
잉여체 확대가 분리라고 가정하면
Remark \ref{more-morphisms-remark-necessary-condition-slice-smooth}의 현상은
일어나지 않는다. 정확한 결과는 다음과 같다.

\begin{lemma}
\label{more-morphisms-lemma-slice-smooth-once-separable-residue-field-extension}
$f : X \to S$를 스킴의 사상이라 하자.
$x \in X$를 상이 $s \in S$인 점이라 하자.
다음을 가정하자.
\begin{enumerate}
\item $f$는 $x$에서 매끄럽다.
\item 잉여체 확대 $\kappa(x)/\kappa(s)$는 분리이다.
\item $x$는 $X_s$의 일반점이 아니다.
\end{enumerate}
그러면 $x$의 아핀 열린 근방 $U \subset X$와 $x$를 포함하는
유효 Cartier 인자 $D \subset U$가 존재하여 $D \to S$는 매끄럽다.
\end{lemma}

\begin{proof}
$k = \kappa(s)$ 및 $R = \mathcal{O}_{X_s, x}$라 쓰자.
$\mathfrak m$을 $R$의 극대 아이디얼이라 하고
$\kappa = R/\mathfrak m$라 놓으면 $\kappa = \kappa(x)$이다.
미분 가군을 만드는 것은 국소화와 가환하므로(Algebra, Lemma
\ref{algebra-lemma-differentials-localize} 참조)
$\Omega_{X_s/s, x} = \Omega_{R/k}$이다. 가정 (2)와 Algebra, Lemma
\ref{algebra-lemma-computation-differential}에 의해 사상
$$
\text{d} :
\mathfrak m/\mathfrak m^2
\longrightarrow
\Omega_{R/k} \otimes_R \kappa(\mathfrak m)
$$
은 단사이다. 가정 (3)은
$\mathfrak m/\mathfrak m^2 \not = 0$을 함의한다.
따라서 $\text{d}\overline{h}$가 영이 아닌 원소
$\overline{h} \in \mathfrak m$이 존재한다. $\overline{h}$의 올림
$h \in \mathcal{O}_{X, x}$를 택하고
Lemma \ref{more-morphisms-lemma-slice-smooth-given-element}을 적용한다.
\end{proof}

\noindent
다음 보조정리에서 구성하는 부분스킴 $Z$는 실제로 $x$의 한 아핀 열린
근방에서 완전 교차이다. 이 사실이 필요해지면 이를 명시하는 별도의
보조정리를 세울 것이다.

\begin{lemma}
\label{more-morphisms-lemma-slice-smooth}
$f : X \to S$를 스킴의 사상이라 하자.
$x \in X$를 상이 $s \in S$인 점이라 하자.
다음을 가정하자.
\begin{enumerate}
\item $f$는 $x$에서 매끄럽다.
\item $x$는 $X_s$의 닫힌점이고 $\kappa(s) \subset \kappa(x)$는 분리이다.
\end{enumerate}
그러면 $x$를 포함하는 몰입 $Z \to X$가 존재하여 다음이 성립한다.
\begin{enumerate}
\item $Z \to S$는 \'etale이다.
\item 집합론적으로 $Z_s = \{x\}$이다.
\end{enumerate}
\end{lemma}

\begin{proof}
$S$를 $s$의 한 아핀 열린 근방으로 바꾼다. 또한 $X$를 $x$의 한 아핀
열린 근방으로 바꾸어 $X \to S$가 매끄럽다고 가정한다.
아핀 스킴 사이의 매끄러운 사상에 대해 $d = \dim_x(X_s)$에 대한
귀납법으로 보조정리를 증명한다.

\medskip\noindent
$d = 0$인 경우. 이 경우 $Z$를 $x$의 한 열린 근방으로 택할 수 있음을
보인다. 실제로 $d = 0$이면 $X \to S$는 $x$에서 준유한이다
(Morphisms, Lemma
\ref{morphisms-lemma-locally-quasi-finite-rel-dimension-0} 참조).
따라서 $U \to S$가 준유한인 아핀 열린 근방 $U \subset X$가 존재한다
(Morphisms, Lemma \ref{morphisms-lemma-quasi-finite-points-open} 참조).
그러므로 $X$를 $U$로 바꾸면 $X$는 $S$ 위에서 준유한이고 매끄러우며,
따라서 $S$ 위에서 상대 차원 $0$으로 매끄럽고, 따라서 $S$ 위에서
\'etale이다. 더욱이 올 $X_s$는 유한 이산 집합이다.
% P05-EN-CH37-0177
따라서 $X$를 $X$의 더 작은 아핀 열린 근방으로 바꾸면
$f^{-1}(\{s\}) = \{x\}$임을 알 수 있다($X_s$의 위상은 $X$의 위상에서
유도되기 때문이다. Schemes, Lemma
\ref{schemes-lemma-fibre-topological} 참조). 이는 이 경우의
보조정리를 증명한다.

\medskip\noindent
이제 $d > 0$이라 가정하자. $x$는 그 올의 닫힌점이므로 확대
$\kappa(x)/\kappa(s)$는 유한이다(Hilbert 영점정리에 의하며,
Morphisms, Lemma
\ref{morphisms-lemma-closed-point-fibre-locally-finite-type} 참조).
대수적 분리 체 확대에 대해서는 성립하므로
$\Omega_{\kappa(x)/\kappa(s)} = 0$이다. 따라서
Lemma \ref{more-morphisms-lemma-slice-smooth-once}을 적용하여 다음 도식을 얻는다.
$$
\xymatrix{
D \ar[r] \ar[rrd] & U \ar[r] \ar[rd] & X \ar[d] \\
& & S
}
$$
여기서 $x \in D$이다. 예를 들어 Algebra, Lemma
\ref{algebra-lemma-one-equation}에 의해
$\dim(\mathcal{O}_{D_s, x}) = \dim(\mathcal{O}_{X_s, x}) - 1$이고
($D_s \subset X_s$도 유효 Cartier이다. Divisors, Lemma
\ref{divisors-lemma-relative-Cartier}를 보라), 이어서 Morphisms, Lemma
\ref{morphisms-lemma-dimension-fibre-at-a-point}을 사용하면
$\dim_x(D_s) = \dim_x(X_s) - 1$임에 유의하라. 따라서 사상
$D \to S$는 매끄럽고
$\dim_x(D_s) = \dim_x(X_s) - 1 = d - 1$이다. 귀납 가정에 따라
원하는 몰입 $Z \to D$를 찾을 수 있으며, 이로써 증명이 끝난다.
\end{proof}

\begin{lemma}
\label{more-morphisms-lemma-etale-nbhd-dominates-smooth}
\begin{slogan}
매끄러운 사상은 \'etale 국소 단면을 갖는다.
\end{slogan}
\begin{reference}
See \cite[Corollaire 17.16.3 (ii)]{EGA4}.
\end{reference}
$f : X \to S$를 스킴의 매끄러운 사상이라 하자.
$s \in S$를 $f$의 상에 속하는 점이라 하자.
그러면 \'etale 근방 $(S', s') \to (S, s)$와
$S$-사상 $S' \to X$가 존재한다.
\end{lemma}

\begin{proof}[Lemma \ref{more-morphisms-lemma-etale-nbhd-dominates-smooth}의 첫째 증명]
가정에 따라 $X_s \not = \emptyset$이다. Varieties, Lemma
\ref{varieties-lemma-smooth-separable-closed-points-dense}에 의해
$\kappa(x)$가 $\kappa(s)$의 유한 분리 체 확대인 닫힌점
$x \in X_s$가 존재한다. 따라서 Lemma \ref{more-morphisms-lemma-slice-smooth}에 의해
$Z \to S$가 \'etale이고 $x \in Z$인 몰입 $Z \to X$가 존재한다.
$(S' , s') = (Z, x)$로 택한다.
\end{proof}

\begin{proof}[Lemma \ref{more-morphisms-lemma-etale-nbhd-dominates-smooth}의 둘째 증명]
$f(x) = s$인 점 $x \in X$를 택한다. 다음 도식을 택한다.
$$
\xymatrix{
X \ar[d] & U \ar[l] \ar[d] \ar[r]_-\pi & \mathbf{A}^d_V \ar[ld] \\
S & V \ar[l]
}
$$
여기서 $\pi$는 \'etale이고, $x \in U$이며,
$V = \Spec(R)$는 아핀이다. Morphisms, Lemma
\ref{morphisms-lemma-smooth-etale-over-affine-space}을 보라.
특히 $s \in V$이다. 사상 $\pi : U \to \mathbf{A}^d_V$는 열린 사상이다.
Morphisms, Lemma \ref{morphisms-lemma-etale-open}을 보라.
따라서 $W = \pi(U) \cap \mathbf{A}^d_s$는
$\mathbf{A}^d_s$의 공집합이 아닌 열린 부분집합이다.
$\kappa(s) \subset \kappa(w)$가 유한 분리인 점 $w \in W$를 택한다.
Varieties, Lemma \ref{varieties-lemma-affine-space-over-field}을 보라.
Algebra, Lemma \ref{algebra-lemma-dim-affine-space}에 의해
$\kappa(s)[x_1, \ldots, x_d]$에서 $w$에 대응하는 극대 아이디얼을
생성하는 $d$개의 원소
$\overline{f}_1, \ldots, \overline{f}_d \in \kappa(s)[x_1, \ldots, x_d]$
가 존재한다. $R$을 한 주 국소화로 바꾸어, 그 상이
$\overline{f}_1, \ldots, \overline{f}_d \in \kappa(s)[x_1, \ldots, x_d]$
인 $f_1, \ldots, f_d \in R[x_1, \ldots, x_d]$가 존재한다고
가정할 수 있다. $R$-대수
$$
R' = R[x_1, \ldots, x_d]/(f_1, \ldots, f_d)
$$
를 생각하고 $S' = \Spec(R')$라 놓자. 구성상 $V$ 위의 닫힌 몰입
$j : S' \to \mathbf{A}^d_V$가 있다. 구성상 $S' \to V$의 $s$ 위의 올은
잉여체가 $\kappa(s)$ 위에서 유한 분리인 한 점 $s'$이다. 이에 대응하는
소 아이디얼을 $\mathfrak q' \subset R'$라 하자. Algebra, Lemma
\ref{algebra-lemma-localize-relative-complete-intersection}에 의해
어떤 $g \in R'$, $g \not \in \mathfrak q$에 대하여 $(R')_g$는
$R$ 위의 상대 전역 완전 교차이다.
% P05-EN-CH37-0178
따라서 $S' \to V$는 $s'$의 한 근방에서 평탄하고 유한 표시이다.
Algebra, Lemma
\ref{algebra-lemma-relative-global-complete-intersection}을 보라.
구성상 $S' \to V$의 $s$ 위의 스킴론적 올은
$\Spec(\kappa(s'))$이다. 따라서 Morphisms, Lemma
\ref{morphisms-lemma-etale-at-point}에 의해 $S' \to S$는
$s'$에서 \'etale이다. 다음과 같이 놓자.
$$
S'' = U \times_{\pi, \mathbf{A}^d_V, j} S'.
$$
구성상 사영 $p : S'' \to S'$에 의해 $s'$로 가는 점
$s'' \in S''$가 존재한다. $p$는 \'etale 사상 $\pi$의 밑변환이므로
\'etale임에 유의하라. Morphisms, Lemma
\ref{morphisms-lemma-base-change-etale}를 보라. $s''$의 작은 아핀 근방
$S''' \subset S''$을 택하여, $S' \to S$가 \'etale인
$s' \in S'$의 공집합이 아닌 열린 근방 안으로 가게 하자.
그러면 \'etale 근방 $(S''', s'') \to (S, s)$가 보조정리의 문제를 푼다.
\end{proof}

\noindent
다음 보조정리는 매끄러운 위상의 층이 \'etale 위상의 층과
같은 것임을 보인다.

\begin{lemma}
\label{more-morphisms-lemma-etale-dominates-smooth}
$S$를 스킴이라 하자. $\mathcal{U} = \{S_i \to S\}_{i \in I}$를
$S$의 매끄러운 덮개라 하자(Topologies, Definition
\ref{topologies-definition-smooth-covering} 참조).
그러면 $\mathcal{U}$를 세분하는(Sites, Definition
\ref{sites-definition-morphism-coverings} 참조) \'etale 덮개
$\mathcal{V} = \{T_j \to S\}_{j \in J}$가 존재한다(Topologies,
Definition \ref{topologies-definition-etale-covering} 참조).
\end{lemma}

\begin{proof}
모든 $s \in S$에 대하여 $s$가 $S_i \to S$의 상에 속하는
$i \in I$가 존재한다. Lemma \ref{more-morphisms-lemma-etale-nbhd-dominates-smooth}에
따라 $s \in g_s(T_s)$이고 $g_s$가 $S_i \to S$를 통해 분해되는
\'etale 사상 $g_s : T_s \to S$를 찾을 수 있다. 따라서
$\{T_s \to S\}$는 $\mathcal{U}$를 세분하는 $S$의 \'etale 덮개이다.
\end{proof}

\begin{lemma}
\label{more-morphisms-lemma-cover-smooth-by-special}
$f : X \to S$를 스킴의 매끄러운 사상이라 하자. 그러면 \'etale 덮개
$\{U_i \to X\}_{i \in I}$가 존재하여, $U_i \to S$는
$U_i \to V_i \to S$로 분해된다. 여기서 $V_i \to S$는 \'etale이고,
$U_i \to V_i$는 아핀 스킴 사이의 매끄러운 사상으로서 단면을 가지며
기하학적으로 연결인 올들을 가진다.
\end{lemma}

\begin{proof}
$s \in S$라 하자. Varieties, Lemma
\ref{varieties-lemma-smooth-separable-closed-points-dense}에 의해
$\kappa(x)/\kappa(s)$가 분리인 닫힌점 $x \in X_s$들의 집합은
$X_s$에서 조밀하다. 따라서 $x$가 상에 속하고 $U \to S$가
보조정리에 서술된 방식으로 분해되는 \'etale 사상 $U \to X$를
구성하면 충분하다. 이를 위해 $x$를 지나며 $Z \to S$가 \'etale인
몰입 $Z \to X$를 택한다(Lemma \ref{more-morphisms-lemma-slice-smooth}).
$S$를 $Z$로, $X$를 $Z \times_S X$로 바꾸면, $X \to S$가
$\sigma(s) = x$를 만족하는 단면 $\sigma : S \to X$를 가진다고
가정할 수 있다. 이제 먼저 $S$를 $s$의 한 아핀 열린 근방으로,
다음으로 $X$를 $x$의 한 아핀 열린 근방으로 바꿀 수 있다.
끝으로 Section \ref{more-morphisms-section-connected-components}의 부분집합
$X^0 \subset X$를 생각한다. Lemmas
\ref{more-morphisms-lemma-connected-along-section-open}와
\ref{more-morphisms-lemma-connected-along-section-locally-constructible}에 의해
이는 $\sigma$를 포함하는 역콤팩트 열린 부분스킴이고,
$X^0 \to S$의 올들은 기하학적으로 연결이다. $X^0$가 아핀이 아니면
$x$를 포함하는 아핀 열린집합 $U \subset X^0$를 택한다.
$X^0 \to S$가 매끄러우므로 $U$의 상은 열린집합이다.
$\sigma^{-1}(U)$와 $U \to S$의 상에 포함되는 $s$의 아핀 열린 근방
$V \subset S$를 택한다. 끝으로 $U \cap f^{-1}(V) \to V$가 원하는
모든 성질을 가짐을 독자가 확인할 수 있다. 예를 들어
$U \cap f^{-1}(V)$는 $U \times_S V$와 같고, 아핀 스킴들의
올곱이므로 아핀이다. 또한 $U \cap f^{-1}(V) \to V$의 기하학적
올들은 $X^0 \to S$의 기약인 올들의 공집합이 아닌 열린 부분스킴이므로
연결이다. 몇몇 세부사항은 생략한다.
\end{proof}











\section{\'Etale 근방과 Artin 근사}
\label{more-morphisms-section-etale-nbhds-artin}

\noindent
이 절에서는 두 점을 가진 스킴의 완비 국소환이 동형이면 이들이 서로
동형인 \'etale 근방을 가진다는 꼴의 결과들을 증명한다. Popescu 정리에
의존할 것이다. Smoothing Ring Maps, Theorem
\ref{smoothing-theorem-popescu}을 보라.

\begin{lemma}
\label{more-morphisms-lemma-map-approximation}
$S$를 국소 뇌터 스킴이라 하자. $X$, $Y$를 $S$ 위에서 국소 유한형인
스킴이라 하자. $x \in X$와 $y \in Y$를 같은 점 $s \in S$ 위에 놓인
점들이라 하자. $\mathcal{O}_{S, s}$가 G-환이라고 가정하자. 더 나아가
국소 $\mathcal{O}_{S, s}$-대수 사상
$$
\varphi : \mathcal{O}_{Y, y} \longrightarrow \mathcal{O}_{X, x}^\wedge
$$
이 주어졌다고 가정하자. 모든 $N \geq 1$에 대하여 기본 \'etale 근방
$(U, u) \to (X, x)$와 $u$를 $y$로 보내는 $S$-사상
$f : U \to Y$가 존재하여, 다음 도식은 $\mathfrak m_u^N$을 법으로
가환한다.
$$
\xymatrix{
\mathcal{O}_{X, x}^\wedge \ar[r] &
\mathcal{O}_{U, u}^\wedge \\
\mathcal{O}_{Y, y} \ar[r]^{f^\sharp_u} \ar[u]^\varphi &
\mathcal{O}_{U, u} \ar[u]
}
$$
\end{lemma}

\begin{proof}
문제는 $X$ 위에서 국소적이므로 $X$, $Y$, $S$가 아핀이라고 가정할 수 있다.
$S = \Spec(R)$, $X = \Spec(A)$, $Y = \Spec(B)$라 하고,
$B = R[x_1, \ldots, x_n]/(f_1, \ldots, f_m)$라 쓰자.
$\mathfrak p \subset A$를 $x$에 대응하는 소 아이디얼이라 하자.
More on Algebra, Proposition
\ref{more-algebra-proposition-finite-type-over-G-ring}에 의해 국소환
$\mathcal{O}_{X, x} = A_\mathfrak p$는 G-환이다.
사상 $\varphi$는 다음 사상이다.
$$
B_\mathfrak q^\wedge \longrightarrow A_\mathfrak p^\wedge
$$
여기서 $\mathfrak q \subset B$는 $y$에 대응하는 소 아이디얼이다.
$a_1, \ldots, a_n \in A_\mathfrak p^\wedge$를 사상
$R[x_1, \ldots, x_n] \to B \to B_\mathfrak q^\wedge \to A_\mathfrak p^\wedge$
에 의한 $x_1, \ldots, x_n$의 상들이라 하자. 이제 Smoothing Ring Maps,
Lemma \ref{smoothing-lemma-approximation-property-variant}을 적용하면
\'etale 환 사상 $A \to A'$, 소 아이디얼 $\mathfrak p' \subset A'$,
그리고 $b_1, \ldots, b_n \in A'$를 얻으며, 이들은
$\kappa(\mathfrak p) = \kappa(\mathfrak p')$,
$a_i - b_i \in (\mathfrak p')^N(A'_{\mathfrak p'})^\wedge$, 그리고
$j = 1, \ldots, n$에 대해
$f_j(b_1, \ldots, b_n) = 0$을 만족한다.
% P05-EN-CH37-0179
이는 $x_i$의 류를 $b_i \in A'$로 보내는 $R$-대수 사상
$B \to A'$를 정한다. $U = \Spec(A') \to \Spec(B)$를 사상 $f$로,
$u = \mathfrak p'$로 택하면 증명이 끝난다.
\end{proof}

\begin{lemma}
\label{more-morphisms-lemma-isomorphism-approximation}
$S$를 국소 뇌터 스킴이라 하자. $X$, $Y$를 $S$ 위에서 국소 유한형인
스킴이라 하자. $x \in X$와 $y \in Y$를 같은 점 $s \in S$ 위에 놓인
점들이라 하자. $\mathcal{O}_{S, s}$가 G-환이라고 가정하자.
완비 국소환 사이의 $\mathcal{O}_{S, s}$-대수 동형
$$
\varphi : \mathcal{O}_{Y, y}^\wedge \longrightarrow \mathcal{O}_{X, x}^\wedge
$$
이 주어졌다고 가정하자. 그러면 모든 $N \geq 1$에 대하여
$S$ 위의 점을 가진 스킴들의 사상
$$
(X, x) \leftarrow (U, u) \rightarrow (Y, y)
$$
들이 존재하여, 두 화살표 모두 기본 \'etale 근방을 정하고 다음 도식은
$\mathfrak m_u^N$을 법으로 가환한다.
$$
\xymatrix{
& \mathcal{O}_{U, u}^\wedge \\
\mathcal{O}_{Y, y}^\wedge \ar[rr]^\varphi \ar[ru] & &
\mathcal{O}_{X, x}^\wedge \ar[lu]
}
$$
\end{lemma}

\begin{proof}
$N \geq 2$라고 가정할 수 있다. Lemma \ref{more-morphisms-lemma-map-approximation}을
적용하여 $(U, u) \to (X, x)$와
$f : (U, u) \to (Y, y)$를 얻는다. $f$가 $u$에서 \'etale임을
보이면 증명이 끝난다. 실제로 유도된 사상
$\mathcal{O}_{Y, y}^\wedge \to \mathcal{O}_{U, u}^\wedge$
이 동형임을 보일 것이다. 이를 증명하면 Lemma
\ref{more-morphisms-lemma-lifting-along-artinian-at-point}에 의해 $f$는 $u$에서
매끄럽고, 물론 $f$는 $u$에서 비분기이기도 하므로 Morphisms, Lemma
\ref{morphisms-lemma-etale-smooth-unramified}에 의해 $f$는
$u$에서 \'etale이다. 국소환 $(R, \mathfrak m)$에 대하여 다음과 같이 놓자.
$\text{Gr}_\mathfrak m(R) =
\bigoplus_{n \geq 0} \mathfrak m^n/\mathfrak m^{n + 1}$.
주장을 증명하기 위해 유도된 등급환 도식
$$
\xymatrix{
& \text{Gr}_{\mathfrak m_u}(\mathcal{O}_{U, u}) \\
\text{Gr}_{\mathfrak m_y}(\mathcal{O}_{Y, y}) \ar[rr]^\varphi \ar[ru] & &
\text{Gr}_{\mathfrak m_x}(\mathcal{O}_{X, x}) \ar[lu]
}
$$
을 살펴보자. 표시한 등급 대수들은 차수 $1$에서 생성되므로
$N \geq 2$이면 이 도식은 실제로 가환한다. 가정에 따라 아래 화살표는
동형이다. 예를 들어 More on Algebra, Lemma
\ref{more-algebra-lemma-flat-unramified}에 의해 사상
$\mathcal{O}_{X, x}^\wedge \to \mathcal{O}_{U, u}^\wedge$
은 동형이고, 따라서 도식의 북서쪽 화살표도 동형이다. 그러므로
$f$는 다음 동형을 유도한다.
% P05-EN-CH37-0180
$\text{Gr}_{\mathfrak m_x}(\mathcal{O}_{X, x}) \to
\text{Gr}_{\mathfrak m_y}(\mathcal{O}_{U, u})$.
두 국소환에 대해 귀납법과 짧은 완전열
$$
0 \to \text{Gr}^n_{\mathfrak m}(R) \to R/\mathfrak m^{n + 1} \to
R/\mathfrak m^n \to 0
$$
을 사용하면, 뱀 보조정리로부터 $f$가 모든 $n$에 대해 동형
$\mathcal{O}_{Y, y}/\mathfrak m_y^n \to \mathcal{O}_{U, u}/\mathfrak m_u^n$
을 유도함을 얻는다. 이것이 보이려던 바이다.
\end{proof}

\begin{lemma}
\label{more-morphisms-lemma-relative-map-approximation-pre}
$X \to S$, $Y \to T$, $x$, $s$, $y$, $t$, $\sigma$, $y_\sigma$ 및
$\varphi$가 다음과 같이 주어졌다고 하자. 스킴의 사상
$$
\vcenter{
\xymatrix{
X \ar[d] & Y \ar[d] \\
S & T
}
}
\quad\text{점들을 함께 표시하면}\quad
\vcenter{
\xymatrix{
x \ar[d] & y \ar[d] \\
s & t
}
}
$$
들이 있고 표시된 점들이 있다고 하자. 여기서 $S$는 국소 뇌터이고
$T$는 $\mathbf{Z}$ 위에서 유한형이다. 사상 $X \to S$와
$Y \to T$는 국소 유한형이다. 국소환 $\mathcal{O}_{S, s}$는 G-환이다.
사상
$$
\sigma : \mathcal{O}_{T, t} \longrightarrow \mathcal{O}_{S, s}^\wedge
$$
은 국소 준동형이다.
$Y_\sigma = Y \times_{T, \sigma} \Spec(\mathcal{O}_{S, s}^\wedge)$라 놓자.
다음으로 $y_\sigma$는 $y$와
$\Spec(\mathcal{O}_{S, s}^\wedge)$의 닫힌점으로 가는
$Y_\sigma$의 점이다. 끝으로
$$
\varphi :
\mathcal{O}_{X, x}^\wedge
\longrightarrow
\mathcal{O}_{Y_\sigma, y_\sigma}^\wedge
$$
는 $\mathcal{O}_{S, s}^\wedge$-대수의 동형이다. 이 상황에서
가환 도식
$$
\xymatrix{
X \ar[d] &
W \ar[l] \ar[rd] \ar[rr] & &
Y \times_{T, \tau} V \ar[r] \ar[ld] & Y \ar[d] \\
S & &
V \ar[ll] \ar[rr]^\tau & &
T
}
$$
과 점 $w \in W$, $v \in V$가 존재하여 다음이 성립한다.
\begin{enumerate}
\item $(V, v) \to (S, s)$는 기본 \'etale 근방이다.
\item $(W, w) \to (X, x)$는 기본 \'etale 근방이다.
\item $\tau(v) = t$이다.
\end{enumerate}
$y_\tau \in Y \times_T V$를 식별
$(Y_\sigma)_s = (Y \times_T V)_v$ 아래에서 $y_\sigma$에
대응하는 점이라 하자. 그러면
\begin{enumerate}
\item[(4)] $(W, w) \to (Y \times_{T, \tau} V, y_\tau)$는
기본 \'etale 근방이다.
\end{enumerate}
\end{lemma}

\begin{proof}
$X_\sigma = X \times_S \Spec(\mathcal{O}_{S, s}^\wedge)$라 하고,
$x_\sigma \in X_\sigma$를 $x$ 위에 놓인 유일한 점이라 하자.
More on Algebra, Proposition
\ref{more-algebra-proposition-Noetherian-complete-G-ring}에 의해
$\mathcal{O}_{S, s}^\wedge$는 G-환임에 유의하라. Lemma
\ref{more-morphisms-lemma-isomorphism-approximation}에 의해
$$
(X_\sigma, x_\sigma) \leftarrow (U, u) \rightarrow (Y_\sigma, y_\sigma)
$$
를 택할 수 있으며, 여기서 두 화살표 모두 기본 \'etale 근방이다.

\medskip\noindent
$S$를 $s$의 한 열린 근방으로 바꾸어 $S = \Spec(R)$가 아핀이라고
가정할 수 있다. $\mathcal{O}_{S, s}$가 G-환이므로 Smoothing Ring Maps,
Theorem \ref{smoothing-theorem-popescu}에 의해 환
$\mathcal{O}_{S, s}^\wedge$는 매끄러운 $R$-대수들의 여과 쌍극한이다.
따라서
$$
\Spec(\mathcal{O}_{S, s}^\wedge) = \lim S_i
$$
로 쓸 수 있다. 여기서 이는 $S$ 위에서 매끄러운 아핀 스킴
$S_i$들의 유향 극한이다. $s_i \in S_i$를
$\Spec(\mathcal{O}_{S, s}^\wedge)$의 닫힌점의 상이라 하자.
$\kappa(s) = \kappa(s_i)$임에 유의하라.
$X_i = X \times_S S_i$라 놓고, $x_i \in X_i$를 $x$로 가는 유일한
점이라 하자. $\kappa(x) = \kappa(x_i)$임에 유의하라.
$T$가 $\mathbf{Z}$ 위에서 유한형이므로 Limits, Proposition
\ref{limits-proposition-characterize-locally-finite-presentation}에 의해
어떤 $i$와 점을 가진 스킴의 사상
$\sigma_i : (S_i, s_i) \to (T, t)$를 택하여,
$\Spec(\mathcal{O}_{S, s}^\wedge) \to S_i$와의 합성이 $\sigma$와
같게 할 수 있다. $Y_i = Y \times_T S_i$라 놓고, $y_i$를
$y_\sigma$의 상이라 하자. $\kappa(y_i) = \kappa(y_\sigma)$임에
유의하라. Limits, Lemma
\ref{limits-lemma-descend-finite-presentation}에 의해 어떤 $i$와 도식
$$
\xymatrix{
X_i \ar[rd] &
U_i \ar[l] \ar[d] \ar[r] &
Y_i \ar[ld] \\
& S_i
}
$$
을 택하여, 이를 $\Spec(\mathcal{O}_{S, s}^\wedge)$로 밑변환하면
$X_\sigma \leftarrow U \rightarrow Y_\sigma$를 복원하게 할 수 있다.
Limits, Lemma \ref{limits-lemma-descend-etale}에 의해 $i$를 늘려
사상들 $X_i \leftarrow U_i \rightarrow Y_i$가 \'etale이라고
가정할 수 있다. $u_i \in U_i$를 $u$의 상이라 하자.
그러면 $u_i \mapsto x_i$이므로
$\kappa(x) = \kappa(x_\sigma) = \kappa(u) \supset \kappa(u_i) \supset
\kappa(x_i) = \kappa(x)$이고, $\kappa(u_i) = \kappa(x_i)$임을
알 수 있다. 따라서
$(X_i, x_i) \leftarrow (U_i, u_i)$는 기본 \'etale 근방이다.
또한 $\kappa(y_i) = \kappa(y_\sigma) = \kappa(u)$이므로
$(U_i, u_i) \to (Y_i, y_i)$도 기본 \'etale 근방임을 알 수 있다.

\medskip\noindent
이제 보조정리의 진술과 같은 도식
$$
\xymatrix{
X \ar[d] &
X \times_S S_i \ar[l] \ar[rd] &
U_i \ar[l] \ar[r] \ar[d] &
Y \times_T S_i \ar[r] \ar[ld] &
Y \ar[d] \\
S & &
S_i \ar[ll] \ar[rr] & &
T
}
$$
을 구성했다. 다만 $S_i \to S$가 매끄럽다는 차이가 있다.
Lemma \ref{more-morphisms-lemma-slice-smooth}에 의해, 그리고 $S_i$를 줄여,
$s_i$를 지나며 $V \to S$가 \'etale인 닫힌 부분스킴
$V \subset S_i$가 존재한다고 가정할 수 있다.
$W$를 $U_i$에서 $V$의 스킴론적 역상으로 놓으면 결론을 얻는다.
\end{proof}

\noindent
독자에게 이 절의 나머지를 건너뛰기를 강력히 권한다.

\begin{lemma}
\label{more-morphisms-lemma-relative-map-approximation}
다음 도식을 생각하자.
$$
\vcenter{
\xymatrix{
X \ar[d] & Y \ar[d] \\
S & T \ar[l]
}
}
\quad\text{점들을 함께 표시하면}\quad
\vcenter{
\xymatrix{
x \ar[d] & y \ar[d] \\
s & t \ar[l]
}
}
$$
여기서 $S$는 국소 뇌터 스킴이고 사상들은 국소 유한형이다.
$\mathcal{O}_{S, s}$가 G-환이라고 가정하자. 더 나아가 국소
$\mathcal{O}_{S, s}$-대수 사상
$$
\sigma : \mathcal{O}_{T, t} \longrightarrow \mathcal{O}_{S, s}^\wedge
$$
과 국소 $\mathcal{O}_{S, s}$-대수 사상
$$
\varphi :
\mathcal{O}_{X, x}
\longrightarrow
\mathcal{O}_{Y_\sigma, y_\sigma}^\wedge
$$
이 주어졌다고 가정하자. 여기서
$Y_\sigma = Y \times_{T, \sigma} \Spec(\mathcal{O}_{S, s}^\wedge)$이고
$y_\sigma$는 $y$ 위에 놓인 $Y_\sigma$의 유일한 점이다.
모든 $N \geq 1$에 대하여 $S$ 위의 스킴들의 가환 도식
$$
\xymatrix{
X \ar[d] & X \times_S V \ar[l] \ar[rd] &
W \ar[l]^-f \ar[r] \ar[d] &
Y \times_{T, \tau} V \ar[r] \ar[ld] & Y \ar[d] \\
S & & V \ar[ll] \ar[rr]^\tau & & T
}
$$
과 점 $w \in W$, $v \in V$가 존재하여 다음이 성립한다.
\begin{enumerate}
\item $v \mapsto s$, $\tau(v) = t$, $f(w) = (x, v)$이고
$w \mapsto (y, v)$이다.
\item $(V, v) \to (S, s)$는 기본 \'etale 근방이다.
\item 도식
$$
\xymatrix{
\mathcal{O}_{S, s}^\wedge \ar[r] & \mathcal{O}_{V, v}^\wedge \\
\mathcal{O}_{T, t} \ar[r]^{\tau^\sharp_v} \ar[u]_\sigma &
\mathcal{O}_{V, v} \ar[u]
}
$$
은 $\mathfrak m_v^N$을 법으로 가환한다.
% P05-EN-CH37-0181
\item $(W, w) \to (Y \times_{T, \tau} V, (y, v))$는
기본 \'etale 근방이다.
\item 도식
$$
\xymatrix{
\mathcal{O}_{X, x} \ar[r]_\varphi &
\mathcal{O}_{Y_\sigma, y_\sigma}^\wedge \ar[r] &
\mathcal{O}_{Y_\sigma, y_\sigma}/\mathfrak m_{y_\sigma}^N \ar@{=}[r] &
\mathcal{O}_{Y \times_{T, \tau} V, (y, v)}/\mathfrak m_{(y, v)}^N
\ar[d]_{\cong} \\
\mathcal{O}_{X, x} \ar[r] \ar@{=}[u] &
\mathcal{O}_{X \times_S V, (x, v)} \ar[r]^{f^\sharp_w} &
\mathcal{O}_{W, w} \ar[r] &
\mathcal{O}_{W, w}/\mathfrak m_w^N
}
$$
은 가환한다. 등식은 (2)와 (3)에 의해 $Y_\sigma$와
$Y \times_{T, \tau} V$가 다음 환 위에서 표준적으로 동형이라는
사실에서 나온다.
$\mathcal{O}_{V, v}/\mathfrak m_v^N = \mathcal{O}_{S, s}/\mathfrak m_s^N$
\end{enumerate}
\end{lemma}

\begin{proof}
$X$, $S$, $T$, $Y$를 아핀 열린 부분스킴으로 바꾸어, 보조정리의
진술에 있는 도식이 뇌터 환과 유한형 환 사상들의 도식
$$
\vcenter{
\xymatrix{
A & B \\
R \ar[u] \ar[r] & C \ar[u]
}
}
\quad\text{소 아이디얼들을 함께 표시하면}\quad
\vcenter{
\xymatrix{
\mathfrak p_A & \mathfrak p_B \\
\mathfrak p_R \ar@{-}[u] \ar@{-}[r] & \mathfrak p_C \ar@{-}[u]
}
}
$$
에 $\Spec$을 적용하여 얻어진다고 가정할 수 있다. 이 증명에서 모든 환
$E$는 $\mathfrak p_R$ 위에 놓인 소 아이디얼 $\mathfrak p_E$가 주어진
뇌터 $R$-대수이고, 모든 환 사상은 주어진 소 아이디얼들과 양립하는
$R$-대수 사상이다. 더욱이 $E^\wedge$라 쓰면 $E$를
$\mathfrak p_E$에서 국소화한 환의 완비화를 뜻한다. 또한
$\kappa(\mathfrak p_{E_1}) = \kappa(\mathfrak p_{E_2})$를 만족하는
\'etale 환 사상 $E_1 \to E_2$가 More on Algebra, Lemma
\ref{more-algebra-lemma-flat-unramified}에 의해 동형
$E_1^\wedge = E_2^\wedge$를 유도한다는 사실을 앞으로 언급 없이 사용한다.

\medskip\noindent
이 표기에서 $\sigma$와 $\varphi$는 환 사상
$$
\sigma : C \to R^\wedge
\quad\text{그리고}\quad
\varphi : A \longrightarrow (B \otimes_{C, \sigma} R^\wedge)^\wedge
$$
에 대응한다. 그림으로 나타내면 다음과 같다.
$$
\xymatrix{
A \ar@/^1em/[rrr]^\varphi &
B \ar[r] &
B \otimes_{C, \sigma} R^\wedge \ar[r] &
(B \otimes_{C, \sigma} R^\wedge)^\wedge \\
R \ar[r] \ar[u] &
C \ar[r]^\sigma \ar[u] & R^\wedge \ar[u] \ar[ru]
}
$$
More on Algebra, Proposition
\ref{more-algebra-proposition-Noetherian-complete-G-ring}에 의해
$R^\wedge$는 G-환임에 유의하라. 따라서 More on Algebra, Proposition
\ref{more-algebra-proposition-finite-type-over-G-ring}에 의해
$B \otimes_{C, \sigma} R^\wedge$도 G-환이다.
Lemma \ref{more-morphisms-lemma-map-approximation}을 대수로 옮겨 적용하면, 동형
$\kappa(\mathfrak p_{B \otimes_{C, \sigma} R^\wedge})
\to \kappa(\mathfrak p_{B'})$
을 유도하는 \'etale 환 사상
$B \otimes_{C, \sigma} R^\wedge \to B'$와, 합성
$$
A \to B' \to (B')^\wedge = (B \otimes_{C, \sigma} R^\wedge)^\wedge
$$
이 $(\mathfrak p_{(B \otimes_{C, \sigma} R^\wedge)^\wedge})^N$을
법으로 $\varphi$와 같게 되는 $R$-대수 사상 $A \to B'$가 존재한다.
보조정리의 결론에서 $\varphi$가 쓰이는 곳은 진술의 (5)에 있는 가환성
조건뿐이므로, $\varphi$를 이 합성으로 바꾸어도 된다. 그림은 다음과 같다.
$$
\xymatrix{
& & B' \ar[r] & (B')^\wedge \ar@{=}[d] \\
A \ar[rru] &
B \ar[r] &
B \otimes_{C, \sigma} R^\wedge \ar[r] \ar[u] &
(B \otimes_{C, \sigma} R^\wedge)^\wedge \\
R \ar[r] \ar[u] &
C \ar[r]^\sigma \ar[u] & R^\wedge \ar[u] \ar[ru]
}
$$
다음으로 가정에 의해 $R_{\mathfrak p_R}$가 G-환이므로
$R^\wedge$가 매끄러운 $R$-대수들의 여과 쌍극한이라는 사실을 사용한다
(Smoothing Ring Maps, Theorem \ref{smoothing-theorem-popescu}).
$C$가 $R$ 위에서 유한 표시이므로 어떤 매끄러운 $R \to R'$에 대하여
인수분해
$$
C \to R' \to R^\wedge
$$
를 얻는다. Algebra, Lemma
\ref{algebra-lemma-characterize-finite-presentation}을 보라.
$R'$을 늘려, $B \otimes_{C, \sigma} R^\wedge$로 밑변환하면 $B'$이
되는 \'etale $B \otimes_C R'$-대수 $B''$가 존재한다고 가정할 수 있다.
Algebra, Lemma \ref{algebra-lemma-etale}를 보라. 그러면 $B'$은
이러한 $B''$들의 여과 쌍극한이므로, $R'$을 더 늘려
$A \to B'' \to B'$이 앞에서 구성한 사상이 되는 $R$-대수 사상
$A \to B''$가 존재한다고 가정할 수 있다(위와 같은 참고문헌).
그림은 다음과 같다.
$$
\xymatrix{
& & B'' \ar[r] & B' \ar[r] & (B')^\wedge \ar@{=}[d] \\
A \ar[rru] &
B \ar[r] &
B \otimes_C R' \ar[r] \ar[u] &
B \otimes_{C, \sigma} R^\wedge \ar[r] \ar[u] &
(B \otimes_{C, \sigma} R^\wedge)^\wedge \\
R \ar[r] \ar[u] &
C \ar[r] \ar[u] &
R' \ar[r] \ar[u] &
R^\wedge \ar[u] \ar[ru]
}
$$
그리고
$$
B' = B'' \otimes_{(B \otimes_C R')} (B \otimes_{C, \sigma} R^\wedge)
$$
이다. 이는 $C$를 $R'$으로, $\sigma : C \to R^\wedge$를
$R' \to R^\wedge$로, $B$를 $B''$로 바꾸어 다음 도식으로
단순화할 수 있음을 뜻한다.
$$
\xymatrix{
A \ar[r] &
B \ar[r] &
B \otimes_{C, \sigma} R^\wedge \\
R \ar[r] \ar[u] &
C \ar[r]^\sigma \ar[u] & R^\wedge \ar[u]
}
$$
여기서 $\varphi$는 수평 화살표들의 합성에 이어
$B \otimes_{C, \sigma} R^\wedge$에서 그 완비화로 가는 표준 사상을
합성한 것과 같다. 증명의 마지막 단계에서는
$\Spec(R)$ 위의 $\Spec(C)$와 $\Spec(R)$ 및 사상 $C \to R^\wedge$에
Lemma \ref{more-morphisms-lemma-map-approximation} 또는 그 증명을 한 번 더 적용한다.
이 보조정리는 $R \to D$가 \'etale이고
$\kappa(\mathfrak p_R) = \kappa(\mathfrak p_D)$이며 합성
$$
C \to D \to D^\wedge = R^\wedge
$$
이 $(\mathfrak p_{R^\wedge})^N$을 법으로
$\sigma : C \to R^\wedge$와 같은 환 사상 $C \to D$를 준다.
이제 문제의 해로
$$
V = \Spec(D)
\quad\text{그리고}\quad
W = \Spec(B \otimes_C D)
$$
를 택할 수 있다. 실제로 도식
$$
\xymatrix{
A \ar[r] &
B \otimes_{C, \sigma} R^\wedge \ar[r] &
B \otimes_{C, \sigma} R^\wedge/(\mathfrak p_{R^\wedge})^N \ar@{=}[r] &
B \otimes_C D/(\mathfrak p_D)^N \\
A \ar@{=}[u] \ar[r] &
A \otimes_R D \ar[r] &
B \otimes_R D \ar[r] &
B \otimes_C D/(\mathfrak p_D)^N \ar@{=}[u]
}
$$
은 $C \to D \to D^\wedge = R^\wedge$가
$(\mathfrak p_{R^\wedge})^N$을 법으로 $\sigma$와 같으므로 가환한다.
이는 (5)를 증명하며, 다른 성질들은 구성에서 바로 따른다.
\end{proof}

\begin{lemma}
\label{more-morphisms-lemma-control-agreement}
$T \to S$를 뇌터 스킴의 유한형 사상이라 하자.
$t \in T$가 $s \in S$로 간다고 하고,
$\sigma : \mathcal{O}_{T, t} \to \mathcal{O}_{S, s}^\wedge$
를 국소 $\mathcal{O}_{S, s}$-대수 사상이라 하자. 모든 $N \geq 1$에
대하여 유한형 사상 $(T', t') \to (T, t)$가 존재하여
$\sigma$는
$\mathcal{O}_{T, t} \to \mathcal{O}_{T', t'}$
를 통해 분해되고, 모든 국소
$\mathcal{O}_{S, s}$-대수 사상
$\sigma' : \mathcal{O}_{T, t} \to \mathcal{O}_{S, s}^\wedge$
이 $\mathcal{O}_{T, t} \to \mathcal{O}_{T', t'}$를 통해
분해될 때 $\sigma$와 $\sigma'$은 $\mathfrak m_s^N$을 법으로 일치한다.
\end{lemma}

\begin{proof}
$S$와 $T$가 아핀이라고 가정할 수 있다. $S = \Spec(R)$ 및
$T = \Spec(C)$라 하자. $c_1, \ldots, c_n \in C$를 $R$-대수 $C$의
생성원들이라 하자. $\mathfrak p \subset R$를 $s$에 대응하는
소 아이디얼이라 하고 $\mathfrak p = (f_1, \ldots, f_m)$라 하자.
$R_\mathfrak p$의 분모를 없애기 위해 $R$을 한 주 국소화로 바꾸어,
$\sum i_j = N$인 $I = (i_1, \ldots, i_m)$에 대해
$r_1, \ldots, r_n \in R$ 및
$a_{i, I} \in \mathcal{O}_{S, s}^\wedge$가 존재하여
$$
\sigma(c_i) = r_i + \sum\nolimits_I a_{i, I} f_1^{i_1} \ldots f_m^{i_m}
$$
가 $\mathcal{O}_{S, s}^\wedge$에서 성립한다고 가정할 수 있다.
이제
$$
C' = C[t_{i, I}]/
\left(c_i - r_i - \sum\nolimits_I t_{i, I} f_1^{i_1} \ldots f_m^{i_m}\right)
$$
를 생각하고 $\mathfrak p' = \mathfrak pC' + (t_{i, I})$라 놓자.
$t_{i, I}$를 $a_{i, I}$로 보내어
$\sigma : C \to \mathcal{O}_{S, s}^\wedge$를 $C'$을 통해 분해한다.
% P05-EN-CH37-0182
$T' = \Spec(C')$로 택하면 된다. 실제로 보조정리의 진술에 나오는
임의의 $\sigma'$은 $c_i$를 극대 아이디얼의 $N$제곱을 법으로
$r_i$로 보낸다.
\end{proof}

\begin{lemma}
\label{more-morphisms-lemma-control-graded}
$Y \to T \to S$를 뇌터 스킴의 유한형 사상들이라 하자.
$t \in T$가 $s \in S$로 간다고 하고,
$\sigma : \mathcal{O}_{T, t} \to \mathcal{O}_{S, s}^\wedge$
를 국소 $\mathcal{O}_{S, s}$-대수 사상이라 하자. 유한형 사상
$(T', t') \to (T, t)$가 존재하여 $\sigma$는
$\mathcal{O}_{T, t} \to \mathcal{O}_{T', t'}$
를 통해 분해되고, 모든 국소
$\mathcal{O}_{S, s}$-대수 사상
$\sigma' : \mathcal{O}_{T, t} \to \mathcal{O}_{S, s}^\wedge$
이 $\mathcal{O}_{T, t} \to \mathcal{O}_{T', t'}$를 통해 분해될 때
닫힌 몰입들
$$
Y \times_{T, \sigma} \Spec(\mathcal{O}_{S, s}^\wedge) = Y_\sigma
\longleftarrow Y_t \longrightarrow
Y_{\sigma'} =
Y \times_{T, \sigma'} \Spec(\mathcal{O}_{S, s}^\wedge)
$$
은 서로 동형인 여법선 대수를 가진다.
\end{lemma}

\begin{proof}
먼저 $\sigma$의 존재에 의해 $\kappa(s) = \kappa(t)$임에 유의하라.
$s \in \Spec(\mathcal{O}_{S, s}^\wedge)$ 위의 $Y_\sigma$와
$Y_{\sigma'}$의 올이 모두 $Y_t$와 표준적으로 동형이므로 진술은
뜻이 있다. “$\sigma'$이
$\mathcal{O}_{T, t} \to \mathcal{O}_{T', t'}$를 통해 분해된다”는
성질을 $\sigma'$에 대한 제약으로 생각하자. 이러한 제약이 여러 개,
이를테면 $i = 1, \ldots, n$에 대해
$(T'_i, t'_i) \to (T, t)$로 주어지면,
$(T'_1 \times_T \ldots \times_T T'_n, (t'_1, \ldots, t'_n)) \to
(T, t)$
를 생각하여 이들을 결합할 수 있다. 앞으로 이를 더 언급하지 않고 사용한다.

\medskip\noindent
Lemma \ref{more-morphisms-lemma-control-agreement}에 의해 보조정리의 진술에 나오는
모든 $\sigma'$이 $\mathfrak m_s^2$을 법으로 $\sigma$와 같다고
가정할 수 있다. $Y_\sigma$ 안의 $Y_t$의 여법선 대수는 바로
준연접 등급 $\mathcal{O}_{Y_t}$-대수
$$
\bigoplus\nolimits_{n \geq 0}
\mathfrak m_s^n\mathcal{O}_{Y_\sigma}/
\mathfrak m_s^{n + 1}\mathcal{O}_{Y_\sigma}
$$
이고, $Y_{\sigma'}$에 대해서도 마찬가지이다. $\sigma$와 $\sigma'$이
$\mathfrak m_s^2$을 법으로 일치하므로 이 두 대수는 차수 $0$과
$1$에서 같다. 한편 이 여법선 대수들은 차수 $0$ 위에서 차수 $1$로
생성된다. 따라서 방금 차수 $0$과 $1$에서 구성한 동형을 확장하는
동형이 존재한다면 그것은 유일하다.

\medskip\noindent
$S$와 $T$가 아핀이라고 가정할 수 있다.
$Y = Y_1 \cup \ldots \cup Y_n$을 아핀 열린 덮개라 하자.
각 $Y_i \to T \to S$와 $\sigma$에 대해 원하는 동형(앞 문단 참조)이
존재하도록 보조정리의 $(T_i', t'_i) \to (T, t)$를 구성할 수 있다면,
유일성에 의해 이 동형들을 붙여 $Y \to T$에 대한 결과를 얻는다.
따라서 $Y$가 아핀이라고 가정할 수 있다.

\medskip\noindent
$S = \Spec(R)$, $T = \Spec(C)$, $Y = \Spec(B)$라 쓰자.
표시 $B = C[x_1, \ldots, x_n]/(f_1, \ldots, f_m)$를 택하고,
$R^\wedge = \mathcal{O}_{S, s}^\wedge$라 하자.
$a_{kj} \in R^\wedge[x_1, \ldots, x_n]$를 다음 조건을 만족하는
다항식들이라 하자.
$$
\sum\nolimits_{j = 1, \ldots, m} a_{kj}\sigma(f_j) = 0,\quad
\text{여기서 }k = 1, \ldots, K
$$
이들은 $\sigma(f_j) \in R^\wedge[x_1, \ldots, x_n]$ 사이의
관계 가군의 생성원 집합이다. 따라서 완전열
\begin{equation}
\label{more-morphisms-equation-resolution}
R^\wedge[x_1, \ldots, x_n]^{\oplus K} \to
R^\wedge[x_1, \ldots, x_n]^{\oplus m} \to
R^\wedge[x_1, \ldots, x_n] \to B \otimes_{C, \sigma} R^\wedge \to 0
\end{equation}
을 얻는다. $c$를 이 완전열의 첫째와 둘째 사상 및 아이디얼
$\mathfrak m_{R^\wedge}R^\wedge[x_1, \ldots, x_n]$에 대해
Artin--Rees 보조정리에서 작동하는 정수라 하자. 이는 More on Algebra,
Section \ref{more-algebra-section-artin-rees}에서 정의되어 있다.
다중지표 표기로
$$
a_{kj} = \sum\nolimits_{I \in \Omega} a_{kj, I} x^I
\quad\text{그리고}\quad
f_j = \sum\nolimits_{I \in \Omega} f_{j, I} x^I
$$
라 쓰자. 여기서 $a_{kj, I} \in R^\wedge$, $f_{j, I} \in C$이고,
$\Omega$는 다중지표들의 유한 집합이다. 그러면 $R^\wedge$에서
$$
\sum\nolimits_{j = 1, \ldots, m,\ I, I' \in \Omega,\ I + I' = I''}
a_{kj, I} \sigma(f_{j, I'}) = 0,\quad
I''\text{는 다중지표}
$$
임을 알 수 있다. 따라서
% P05-EN-CH37-0183
$$
C' = C[t_{jk, I}]/
\left(
\sum\nolimits_{j = 1, \ldots, m,\ I, I' \in \Omega,\ I + I' = I''}
t_{kj, I} f_{j, I'},\ I''\text{는 다중지표}\right)
$$
로 택한다. 그러면 $\sigma$는 $t_{kj, I}$를 $a_{jk, I}$로 보내는 사상
$\tilde\sigma : C' \to R^\wedge$를 통해 분해된다.
따라서 $T' = \Spec(C')$에는 점 $t' \in T'$가 있어서
$\sigma$는 $\mathcal{O}_{T, t} \to \mathcal{O}_{T', t'}$를 통해
분해된다. $C'[x_1, \ldots, x_n]$에서
$t_{kj} = \sum t_{kj, I} x^I$라 하자. 그러면 복합체
\begin{equation}
\label{more-morphisms-equation-resolution-new}
C'[x_1, \ldots, x_n]^{\oplus K} \to
C'[x_1, \ldots, x_n]^{\oplus m} \to
C'[x_1, \ldots, x_n] \to B \otimes_C C' \to 0
\end{equation}
를 얻는다. 이는 $C'[x_1, \ldots, x_n]$에서 완전하고
$\tilde\sigma$에 의한 밑변환은 (\ref{more-morphisms-equation-resolution})을 준다.

\medskip\noindent
Lemma \ref{more-morphisms-lemma-control-agreement}에 의해 사상
$(T'', t'') \to (T', t')$를 더 찾을 수 있어서, $\tilde\sigma$는
$\mathcal{O}_{T', t'} \to \mathcal{O}_{T'', t''}$를 통해 분해되고,
$\sigma' : C \to R^\wedge$가 $\mathcal{O}_{T'', t''}$를 통해
분해되면 유도된 사상 $\tilde \sigma' : C' \to R^\wedge$는
$\mathfrak m_s^{c + 1}$을 법으로 $\tilde \sigma$와 일치한다.
따라서 $\sigma'$이 이러한 사상이면 복합체
$$
R^\wedge[x_1, \ldots, x_n]^{\oplus K} \to
R^\wedge[x_1, \ldots, x_n]^{\oplus m} \to
R^\wedge[x_1, \ldots, x_n] \to B \otimes_{C, \sigma'} R^\wedge \to 0
$$
를 얻는데, 이는 다항식 $t_{kj}$와 $f_j$에 $\tilde\sigma'$을 적용하여
얻는 $R^\wedge[x_1, \ldots, x_n]$ 위의 복합체이다. 달리 말하면 이는
복합체 (\ref{more-morphisms-equation-resolution-new})을 $\tilde\sigma'$으로
밑변환한 것이다. $\tilde \sigma$와 $\tilde \sigma'$이
$\mathfrak m_s^{c + 1}$을 법으로 합동이므로, 이 복합체를 정의하는
행렬들은 (\ref{more-morphisms-equation-resolution})을 정의하는 행렬들과
$\mathfrak m_s^{c + 1}$을 법으로 합동이다.
(\ref{more-morphisms-equation-resolution})이 완전하므로 More on Algebra, Lemma
\ref{more-algebra-lemma-approximate-complex-graded}을 적용하여
$$
\text{Gr}_{\mathfrak m_s}(B \otimes_{C, \sigma'} R^\wedge)
\cong
\text{Gr}_{\mathfrak m_s}(B \otimes_{C, \sigma} R^\wedge)
$$
를 얻는다.
\end{proof}

\begin{lemma}
\label{more-morphisms-lemma-relative-isomorphism-approximation}
Lemma \ref{more-morphisms-lemma-relative-map-approximation}과 같은 표기와 가정 아래에서
$\varphi$가 완비화 사이의 동형을 유도한다고 가정하자.
그러면 $f$가 \'etale이 되도록 도식을 택할 수 있다.
\end{lemma}

\begin{proof}
$N \geq 2$라고 가정할 수 있고, Lemma \ref{more-morphisms-lemma-control-graded}에서와
같이 $(T, t)$를 $(T', t')$로 바꾸어도 된다.
$(V, v) \to (S, s)$가 기본 \'etale 근방이므로
$(X \times_S V, (x, v)) \to (X, x)$도 그러하다. 따라서 More on Algebra,
Lemma \ref{more-algebra-lemma-flat-unramified}에 의해
$\mathcal{O}_{X, x} \to \mathcal{O}_{X \times_S V, (x, v)}$는
완비화 사이의 동형을 유도한다.
$\mathcal{O}_{X, x} \to \mathcal{O}_{W, w}$가 완비화 사이의 동형을
유도한다고 주장한다. 이를 증명하면 Lemma
\ref{more-morphisms-lemma-lifting-along-artinian-at-point}에 의해 $f$는 $w$에서
매끄럽고, 물론 $f$는 $u$에서도 비분기이므로 Morphisms, Lemma
\ref{morphisms-lemma-etale-smooth-unramified}에 의해 $f$는
$w$에서 \'etale이다.
% P05-EN-CH37-0184

\medskip\noindent
먼저 Lemma \ref{more-morphisms-lemma-relative-map-approximation}의 (5)의 가환성을
사용하면, $i \leq N$에 대해 가환 도식
$$
\xymatrix{
\text{Gr}^i_{\mathfrak m_x}(\mathcal{O}_{X, x})
\ar[r]_-\varphi &
\text{Gr}^i_{\mathfrak m_{y_\sigma}}(\mathcal{O}_{Y_\sigma, y_\sigma}^\wedge)
\ar@{=}[r] &
\text{Gr}^i_{\mathfrak m_{(y, v)}}(\mathcal{O}_{Y \times_{T, \tau} V, (y, v)})
\ar[d]_{\cong} \\
\text{Gr}^i_{\mathfrak m_x}(\mathcal{O}_{X, x})
\ar[r]^-{\cong} \ar@{=}[u] &
\text{Gr}^i_{\mathfrak m_{(x, v)}}(\mathcal{O}_{X \times_S V, (x, v)})
\ar[r]^{f^\sharp_w} &
\text{Gr}^i_{\mathfrak m_w}(\mathcal{O}_{W, w})
}
$$
이 존재함을 알 수 있다. 이는 $f^\sharp_w$가 잉여체 사이의 동형
$\kappa(x) \to \kappa(w)$와 여접공간 사이의 동형
$\mathfrak m_x/\mathfrak m_x^2 \to \mathfrak m_w/\mathfrak m_w^2$
을 정함을 함의한다. 따라서 $f^\sharp_w$는 완비 국소환 사이의 전사
$\mathcal{O}_{X, x}^\wedge \to \mathcal{O}_{W, w}^\wedge$를 정한다.

\medskip\noindent
Lemma \ref{more-morphisms-lemma-control-graded}에 의해
$\text{Gr}_{\mathfrak m_s}(\mathcal{O}_{(Y \times_{T, \tau} V, (y, v)})$
와
$\text{Gr}_{\mathfrak m_s}(\mathcal{O}_{Y_\sigma, y_\sigma})$.
사이의 동형이 존재한다.
% P05-EN-CH37-0185
이는 여법선 층들의 동형에서 점 $y$에서의 줄기를 취하면 따른다.
우리의 국소환들은 뇌터이고, 한 아이디얼에 관한 완비화는 뇌터 환 위의
유한 가군에 대한 완전 함자이므로, $\mathfrak m_s$에 관한 연관 등급을
취하는 것은 완비화와 가환한다. 이로부터 다음 등급환 도식의 오른쪽
수직 동형을 얻는다.
$$
\xymatrix{
\text{Gr}_{\mathfrak m_s}(\mathcal{O}_{W, w}^\wedge) &
\text{Gr}_{\mathfrak m_s}
(\mathcal{O}_{(Y \times_{T, \tau} V, (y, v)}^\wedge) \ar[l] \\
\text{Gr}_{\mathfrak m_s}(\mathcal{O}_{X, x}^\wedge)
\ar[r]^\varphi \ar[u] &
\text{Gr}_{\mathfrak m_s}(\mathcal{O}_{Y_\sigma, y_\sigma}^\wedge)
\ar[u]_{\cong}
}
$$
이 도식이 가환한다고 주장하지는 않는다. 앞 문단의 결과에 의해
왼쪽 화살표는 전사이고, 나머지 세 화살표는 동형이다. 따라서 왼쪽
화살표는 서로 동형인 뇌터 환 사이의 전사이다. 그러므로 Algebra, Lemma
\ref{algebra-lemma-surjective-endo-noetherian-ring-is-iso}에 의해
동형이다(Hilbert 함수로 직접 논증할 수도 있다). 특히
$\mathcal{O}_{X, x}^\wedge \to \mathcal{O}_{W, w}^\wedge$는
전사일 뿐 아니라 단사여야 하며, 이로써 증명이 끝난다.
\end{proof}















\section{유한 국소 자유 덮개가 국소적으로 \'etale 덮개를 세분한다}
\label{more-morphisms-section-finite-free-over-etale}

\noindent
이 절에서는 대략 스킴 $S$의 \'etale 덮개를 $S$의 유한 국소 자유
덮개들의 Zariski 덮개로 세분할 수 있다는 결과를 설명한다.

\begin{lemma}
\label{more-morphisms-lemma-dominate-etale-neighbourhood-finite-flat}
$S$를 스킴이라 하고 $s \in S$라 하자.
$f : (U, u) \to (S, s)$를 \'etale 근방이라 하자.
아핀 열린 근방 $s \in V \subset S$와 전사인 유한 국소 자유 사상
$\pi : T \to V$가 존재하여, 모든 $t \in \pi^{-1}(s)$에 대해
열린 근방 $t \in W_t \subset T$와 가환 도식
$$
\xymatrix{
T \ar[d]^\pi & W_t \ar[l] \ar[rr]_{h_t} \ar[rd] & & U \ar[dl] \\
V \ar[rr] & & S
}
$$
이 존재하며 $h_t(t) = u$이다.
\end{lemma}

\begin{proof}
문제는 $S$ 위에서 국소적이므로 $S$를 $s$의 임의의 열린 근방으로
바꾸어도 된다. 또한 $U$를 $u$의 한 열린 근방으로 바꾸어도 된다.
따라서 Morphisms, Lemma
\ref{morphisms-lemma-etale-locally-standard-etale}에 의해
$U \to S$가 아핀 스킴 사이의 표준 \'etale 사상이라고 가정할 수 있다.
이 경우 보조정리는 $V = S$로 택하면 Algebra, Lemma
\ref{algebra-lemma-standard-etale-finite-flat-Zariski}에서 따른다.
\end{proof}

\begin{lemma}
\label{more-morphisms-lemma-dominate-etale-affine-finite-flat}
$f : U \to S$를 아핀 스킴 사이의 전사 \'etale 사상이라 하자.
전사인 유한 국소 자유 사상 $\pi : T \to S$와 유한 열린 덮개
$T = T_1 \cup \ldots \cup T_n$이 존재하여, 각 $T_i \to S$는
$U \to S$를 통해 분해된다. 도식은 다음과 같다.
$$
\xymatrix{
& \coprod T_i  \ar[rd] \ar[ld] & \\
T \ar[rd]^\pi & & U \ar[ld]_f \\
& S &
}
$$
여기서 남서쪽 화살표는 Zariski 덮개이다.
\end{lemma}

\begin{proof}
이는 Algebra, Lemma \ref{algebra-lemma-etale-finite-flat-zariski}를
다시 서술한 것이다.
\end{proof}

\begin{remark}
\label{more-morphisms-remark-topologies}
위상의 관점에서 Lemmas
\ref{more-morphisms-lemma-dominate-etale-neighbourhood-finite-flat}와
\ref{more-morphisms-lemma-dominate-etale-affine-finite-flat}의 뜻은 다음과 같다.
$S$를 임의의 스킴이라 하고 $\{f_i : U_i \to S\}$를 $S$의
\'etale 덮개라 하자. Zariski 열린 덮개 $S = \bigcup V_j$, 각 $j$에
대한 전사인 유한 국소 자유 사상 $W_j \to V_j$, 그리고 각 $j$에 대한
Zariski 열린 덮개 $\{W_{j, k} \to W_j\}$가 존재하여,
족 $\{W_{j, k} \to S\}$은 주어진 \'etale 덮개
$\{f_i : U_i \to S\}$를 세분한다. 이것은 실제로 무엇을 뜻하는가?
예를 들어 Zariski 덮개와, $\pi$가 전사인 유한 국소 자유 사상인 꼴의
fppf 덮개 $\{\pi : T \to S\}$에 대해서 풀 수 있는 내림 문제가 있다고
하자. 그러면 이 내림 문제는 \'etale 덮개에 대해서도 긍정적인 답을
가진다. Gabber는 아핀 스킴 $X$에 대해
$\text{Br}(X) = \text{Br}'(X)$임을 증명할 때 이 방법을 사용했다.
\cite{Hoobler}를 보라.
\end{remark}




\section{준유한 사상의 \'etale 국소화}
\label{more-morphisms-section-etale-localization}

\noindent
이제 ``준유한 사상은 \'etale 국소화 뒤에 유한 사상이 된다''라는
주제를 둘러싼 일련의 보조정리를 다룬다. 일반적인 착상은 다음과 같다.
스킴의 사상 $f : X \to S$와 점 $s \in S$가 주어졌다고 하자.
$\varphi : (U, u) \to (S, s)$를 $S$ 안의 $s$의 \'etale 근방이라 하자.
올곱 $X_U = U \times_S X$와 기본 도식
\begin{equation}
\label{more-morphisms-equation-basic-diagram}
\vcenter{
\xymatrix{
V \ar[r] \ar[dr] & X_U \ar[d] \ar[r] & X \ar[d]^f \\
& U \ar[r]^\varphi & S
}
}
\end{equation}
을 생각하자. 여기서 $V \subset X_U$는 열린 부분이다.
사상 $f_U : X_U \to U$, 또는 적당한 열린 부분 $V$에 대한 사상
$V \to U$에는 어떤 표준 모형이 있는가? 물론 일반적으로는 그렇지 않다.
그러나 준유한 사상에 대해서는 몇 가지를 말할 수 있다.

\begin{lemma}
\label{more-morphisms-lemma-etale-makes-quasi-finite-finite-at-point}
$f : X \to S$를 스킴의 사상이라 하자. $x \in X$라 하고
$s = f(x)$라 두자. 다음을 가정하자.
\begin{enumerate}
\item $f$는 국소 유한형이고,
\item $x \in X_s$는 고립점이다\footnote{(1)이 성립할 때 이것은
$f$가 $x$에서 준유한이라는 뜻이다. Morphisms, Lemma
\ref{morphisms-lemma-quasi-finite-at-point-characterize}를 보라.}.
\end{enumerate}
그러면 다음이 존재한다.
\begin{enumerate}
\item[(a)] 기본 \'etale 근방 $(U, u) \to (S, s)$,
\item[(b)] 열린 부분스킴 $V \subset X_U$
(\ref{more-morphisms-equation-basic-diagram}을 보라).
\end{enumerate}
이들은 다음을 만족한다.
\begin{enumerate}
\item[(\romannumeral1)] $V \to U$는 유한 사상이다.
\item[(\romannumeral2)] $U$의 $u$로 가는 $V$의 점 $v$가 유일하게 존재한다.
\item[(\romannumeral3)] 점 $v$는 사상 $X_U \to X$ 아래에서 $x$로 가고,
$\kappa(x) = \kappa(v)$를 유도한다.
\end{enumerate}
더욱이 임의의 기본 \'etale 근방 $(U', u') \to (U, u)$에 대해
$V' = U' \times_U V \subset X_{U'}$라 두면, 삼중항 $(U', u', V')$도
성질 (\romannumeral1), (\romannumeral2), (\romannumeral3)을 만족한다.
\end{lemma}

\begin{proof}
$f(Y) \subset W$이고 $x \in Y$인 아핀 열린 부분 $Y \subset X$와
$W \subset S$를 택하자. $x$는 사상 $f|_Y : Y \to W$의 올에서도
고립점이다. $f|_Y : Y \to W$에 대해 정리를 증명할 수 있으면 $f$에
대한 정리도 따른다. 따라서 $f$가 아핀 스킴 사이의 사상인 경우로
귀착된다. 이 경우는 Algebra, Lemma
\ref{algebra-lemma-etale-makes-quasi-finite-finite-one-prime}이다.
\end{proof}

\noindent
앞의 보조정리와 다음 보조정리에서는 사상 $f$가 분리라고 가정하지 않는다.
따라서 그 안에서 만든 열린 부분 $V$, $V_i$는 $X_U$ 안에서 닫혀 있을
필요가 없다. 또한 근방 $U$를 아핀으로 택하면 각 $V_i$는 아핀이지만,
비분리인 경우 교집합 $V_i \cap V_j$는 아핀이 아닐 수 있다.

\begin{lemma}
\label{more-morphisms-lemma-etale-makes-quasi-finite-finite-multiple-points}
$f : X \to S$를 스킴의 사상이라 하자.
$x_1, \ldots, x_n \in X$를 $S$에서 같은 상 $s$를 갖는 점들이라 하자.
다음을 가정하자.
\begin{enumerate}
\item $f$는 국소 유한형이고,
\item $i = 1, \ldots, n$에 대해 $x_i \in X_s$는 고립점이다.
\end{enumerate}
그러면 다음이 존재한다.
\begin{enumerate}
\item[(a)] 기본 \'etale 근방 $(U, u) \to (S, s)$,
\item[(b)] 각 $i$에 대한 열린 부분스킴 $V_i \subset X_U$.
\end{enumerate}
각 $i$에 대해 다음이 성립한다.
\begin{enumerate}
\item[(\romannumeral1)] $V_i \to U$는 유한 사상이다.
\item[(\romannumeral2)] $U$의 $u$로 가는 $V_i$의 점 $v_i$가 유일하게 존재한다.
\item[(\romannumeral3)] 점 $v_i$는 $X$의 $x_i$로 가고
$\kappa(x_i) = \kappa(v_i)$이다.
\end{enumerate}
\end{lemma}

\begin{proof}
$n$에 대한 귀납법을 사용한다. 즉 $(U, u) \to (S, s)$와
$V_i \subset X_U$, $i = 1, \ldots, n - 1$이
$x_1, \ldots, x_{n - 1}$에 대해 조건을 만족한다고 하자.
$\kappa(s) = \kappa(u)$이므로 올 $(X_U)_u = X_s$이다. 따라서
% P05-EN-CH37-0186
$x_n \in X_s$에 대응하는 유일한 점 $x'_n \in X_u \subset X_U$가
존재한다. 또한 $x'_n$은 $X_u$에서 고립되어 있다. 따라서 Lemma
\ref{more-morphisms-lemma-etale-makes-quasi-finite-finite-at-point}에 의해 기본 \'etale 근방
$(U', u') \to (U, u)$와 $x'_n$에 대해, 따라서 $x_n$에 대해 조건을
만족하는 열린 부분 $V_n \subset X_{U'}$가 존재한다. Lemma
\ref{more-morphisms-lemma-etale-makes-quasi-finite-finite-at-point}의 마지막 주장에 의해,
$i = 1, \ldots, n - 1$에 대한 열린 부분스킴
$V'_i = U'\times_U V_i$는 $x_1, \ldots, x_{n - 1}$에 대해 여전히
조건을 만족한다. 이로써 증명이 끝난다.
\end{proof}

\noindent
비자명한 체 확대 $\kappa(u)/\kappa(s)$, 즉 일반적인 \'etale 근방을
허용하면 다음과 같이 점들을 분리할 수 있다.

\begin{lemma}
\label{more-morphisms-lemma-etale-makes-quasi-finite-finite-multiple-points-var}
$f : X \to S$를 스킴의 사상이라 하자.
$x_1, \ldots, x_n \in X$를 $S$에서 같은 상 $s$를 갖는 점들이라 하자.
다음을 가정하자.
\begin{enumerate}
\item $f$는 국소 유한형이고,
\item $i = 1, \ldots, n$에 대해 $x_i \in X_s$는 고립점이다.
\end{enumerate}
그러면 다음이 존재한다.
\begin{enumerate}
\item[(a)] \'etale 근방 $(U, u) \to (S, s)$,
\item[(b)] 각 $i$에 대한 정수 $m_i$와 열린 부분스킴
$V_{i, j} \subset X_U$, $j = 1, \ldots, m_i$.
\end{enumerate}
이들은 다음을 만족한다.
\begin{enumerate}
\item[(\romannumeral1)] 각 $V_{i, j} \to U$는 유한 사상이다.
\item[(\romannumeral2)] $U$의 $u$로 가는 $V_{i, j}$의 점
$v_{i, j}$가 유일하게 존재하고, $\kappa(u) \subset \kappa(v_{i, j})$는
유한 순수 비분리 확대이다.
% P05-EN-CH37-0187
\item[(\romannumeral4)] $v_{i, j} = v_{i', j'}$이면 $i = i'$이고
$j = j'$이다.
\item[(\romannumeral3)] 점 $v_{i, j}$는 $X$의 $x_i$로 가며,
$(X_U)_u$의 다른 어떤 점도 $x_i$로 가지 않는다.
\end{enumerate}
\end{lemma}

\begin{proof}
이 증명은 Algebra, Lemma
\ref{algebra-lemma-etale-makes-quasi-finite-finite-variant}의 증명을
스킴의 언어로 옮긴 변형이다. Morphisms, Lemma
\ref{morphisms-lemma-quasi-finite-at-point-characterize}에 의해 사상 $f$는
각 점 $x_i$에서 준유한이다. 따라서 각 $i$에 대해
$\kappa(s) \subset \kappa(x_i)$는 유한이다(Morphisms, Lemma
\ref{morphisms-lemma-residue-field-quasi-finite}). 각 $i$에 대해
$L_i/\kappa(s)$가 분리이고 $\kappa(x_i)/L_i$가 순수 비분리가 되도록 하는
부분체 $\kappa(s) \subset L_i \subset \kappa(x_i)$를 택하자.
$i = 1, \ldots, n$에 대해 $\kappa(s)$-매장 $L_i \to L$이 존재하도록
유한 Galois 확대 $L/\kappa(s)$를 택하자. $\kappa(s)$-확대로서
$L \cong \kappa(u)$가 되도록 하는 \'etale 근방
$(U, u) \to (S, s)$를 택하자(Lemma
\ref{more-morphisms-lemma-realize-prescribed-residue-field-extension-etale}).

\medskip\noindent
$y_{i, j}$, $j = 1, \ldots, m_i$를 $x_i \in X$와 $u \in U$ 위에 놓이는
$X_U$의 점들이라 하자. Schemes, Lemma
\ref{schemes-lemma-points-fibre-product}에 의해 이 점들 $y_{i, j}$는 환
$\kappa(u) \otimes_{\kappa(s)} \kappa(x_i)$의 소아이디얼들과 정확히
대응한다. 이것은 점들이 유한 개임도 설명한다. 실제로
$m_i = [L_i : \kappa(s)]$이지만 이 사실은 필요하지 않다. $L$을 택한
방법과 초등적인 체론에 의해 각 쌍 $i, j$에 대해
$\kappa(u) \subset \kappa(y_{i, j})$가 유한 순수 비분리임을 알 수 있다.
또한 예를 들어 Morphisms, Lemma
\ref{morphisms-lemma-base-change-quasi-finite}에 의해 사상 $X_U \to U$는
모든 $i, j$에 대해 점 $y_{i, j}$에서 준유한이다.

\medskip\noindent
사상 $X_U \to U$, 점 $u \in U$, 점들 $y_{i, j} \in (X_U)_u$에
Lemma \ref{more-morphisms-lemma-etale-makes-quasi-finite-finite-multiple-points}를 적용하자.
그러면 $\kappa(u) = \kappa(u')$인 \'etale 근방
$(U', u') \to (U, u)$와, 그 보조정리의 성질 (\romannumeral1),
(\romannumeral2), (\romannumeral3)을 만족하는 열린 부분
$V_{i, j} \subset X_{U'}$를 얻는다. \'etale 근방
$(U', u') \to (S, s)$와 열린 부분 $V_{i, j} \subset X_{U'}$가 이
보조정리에서 제시한 문제의 해라고 주장한다. 확인은 생략한다.
\end{proof}

\begin{lemma}
\label{more-morphisms-lemma-etale-splits-off-quasi-finite-part-technical}
$f : X \to S$를 스킴의 사상이라 하자. $s \in S$라 하고
$x_1, \ldots, x_n \in X_s$라 하자. 다음을 가정하자.
\begin{enumerate}
\item $f$는 국소 유한형이다.
\item $f$는 분리이다.
\item $x_1, \ldots, x_n$은 $X_s$의 서로 다른 고립점들이다.
\end{enumerate}
그러면 기본 \'etale 근방 $(U, u) \to (S, s)$와 분해
$$
U \times_S X = W \amalg V_1 \amalg \ldots \amalg V_n
$$
가 존재한다. 각 항은 열리고 닫힌 부분스킴이며, 사상 $V_i \to U$는
유한이고, $u$ 위의 $V_i \to U$의 올은 한원소 집합 $\{v_i\}$이며,
각 $v_i$는 $x_i$로 가고 $\kappa(x_i) = \kappa(v_i)$이며,
$u$ 위의 $W \to U$의 올은 어느 $x_i$로 가는 점도 포함하지 않는다.
\end{lemma}

\begin{proof}
Lemma \ref{more-morphisms-lemma-etale-makes-quasi-finite-finite-multiple-points}에서처럼
$(U, u) \to (S, s)$와 $V_i \subset X_U$를 택하자. $X_U \to U$는
분리이고(Schemes, Lemma \ref{schemes-lemma-separated-permanence}),
$V_i \to U$는 유한이므로 고유이다(Morphisms, Lemma
\ref{morphisms-lemma-finite-proper}). 따라서 Morphisms, Lemma
\ref{morphisms-lemma-image-proper-scheme-closed}에 의해
$V_i \subset X_U$는 닫혀 있다. 그러므로 $V_i \cap V_j$는 $v_i$를
포함하지 않는 $V_i$의 닫힌 부분집합이다. 따라서 $U$에서
$V_i \cap V_j$의 상은 닫힌 집합이고($V_i \to U$가 고유이므로) $u$를
포함하지 않는다. 그러므로 $U$를 축소한 뒤에는
% P05-EN-CH37-0188
모든 $i, j$에 대해 $V_i \cap V_j = \emptyset$이라고 가정할 수 있다.
이것이 보조정리의 분해를 준다.
\end{proof}

\noindent
다음은 순수 비분리 체 확대로 귀착시키는 변형이다.

\begin{lemma}
\label{more-morphisms-lemma-etale-splits-off-quasi-finite-part-technical-variant}
$f : X \to S$를 스킴의 사상이라 하자. $s \in S$라 하고
$x_1, \ldots, x_n \in X_s$라 하자. 다음을 가정하자.
\begin{enumerate}
\item $f$는 국소 유한형이다.
\item $f$는 분리이다.
\item $x_1, \ldots, x_n$은 $X_s$의 서로 다른 고립점들이다.
\end{enumerate}
그러면 \'etale 근방 $(U, u) \to (S, s)$와 분해
$$
U \times_S X =
W \amalg
\ \coprod\nolimits_{i = 1, \ldots, n}
\ \coprod\nolimits_{j = 1, \ldots, m_i}
V_{i, j}
$$
가 존재한다. 각 항은 열리고 닫힌 부분스킴이며, 사상
$V_{i, j} \to U$는 유한이고, $u$ 위의 $V_{i, j} \to U$의 올은
한원소 집합 $\{v_{i, j}\}$이며, 각 $v_{i, j}$는 $x_i$로 가고,
$\kappa(u) \subset \kappa(v_{i, j})$는 순수 비분리이며, $u$ 위의
$W \to U$의 올은 어느 $x_i$로 가는 점도 포함하지 않는다.
\end{lemma}

\begin{proof}
증명은 Lemma \ref{more-morphisms-lemma-etale-splits-off-quasi-finite-part-technical}의
증명과 완전히 같다. 다만 Lemma
\ref{more-morphisms-lemma-etale-makes-quasi-finite-finite-multiple-points} 대신 Lemma
\ref{more-morphisms-lemma-etale-makes-quasi-finite-finite-multiple-points-var}를 사용한다.
\end{proof}

\noindent
다음 형태가 조금 더 이해하기 쉬울 수 있다.

\begin{lemma}
\label{more-morphisms-lemma-etale-splits-off-quasi-finite-part}
$f : X \to S$를 스킴의 사상이라 하자. $s \in S$라 하고 다음을 가정하자.
\begin{enumerate}
\item $f$는 국소 유한형이다.
\item $f$는 분리이다.
\item $X_s$의 고립점은 많아야 유한 개이다.
\end{enumerate}
그러면 기본 \'etale 근방 $(U, u) \to (S, s)$와 분해
$$
U \times_S X = W \amalg V
$$
가 존재한다. 두 항은 열리고 닫힌 부분스킴이며, 사상 $V \to U$는
유한이고, 사상 $W \to U$의 올 $W_u$는 고립점을 포함하지 않는다.
특히 $f^{-1}(s)$가 유한 집합이면 $W_u = \emptyset$이다.
\end{lemma}

\begin{proof}
$x_1, \ldots, x_n$을 $X_s$의 모든 고립점의 집합으로 택하고
$V = \bigcup V_i$라 두면, Lemma
\ref{more-morphisms-lemma-etale-splits-off-quasi-finite-part-technical}에서 바로 따른다.
\end{proof}






\section{정수적 사상의 \'etale 국소화}
\label{more-morphisms-section-etale-localization-integral}

\noindent
Section \ref{more-morphisms-section-etale-localization}의 결과 가운데 정수적 사상의
경우에 대한 몇 가지 변형을 제시한다.

\begin{lemma}
\label{more-morphisms-lemma-etale-makes-integral-split}
$R \to S$를 정수적 환 사상이라 하자. $\mathfrak p \subset R$를
소아이디얼이라 하자. 다음을 가정하자.
\begin{enumerate}
\item $\mathfrak p$ 위에 놓인 소아이디얼
$\mathfrak q_1, \ldots, \mathfrak q_n$은 유한 개이고,
\item 각 $i$에 대해 극대 분리 부분확대
% P05-EN-CH37-0189
$\kappa(\mathfrak q)/\kappa(\mathfrak q_i)_{sep}/\kappa(\mathfrak p)$
(Fields, Lemma \ref{fields-lemma-separable-first})는
$\kappa(\mathfrak p)$ 위에서 유한이다.
\end{enumerate}
그러면 \'etale 환 사상 $R \to R'$과 $\mathfrak p$ 위에 놓인 소아이디얼
$\mathfrak p'$이 존재하여
$$
S \otimes_R R' = A_1 \times \ldots \times A_m
$$
이고, 정수적 사상 $R' \to A_j$는 $\mathfrak p'$ 위에 놓인 유일한
소아이디얼 $\mathfrak r_j$를 가지며
$\kappa(\mathfrak r_j)/\kappa(\mathfrak p')$는 순수 비분리이다.
\end{lemma}

\begin{proof}[첫째 증명]
이 증명에서는 Algebra, Lemma
\ref{algebra-lemma-etale-makes-quasi-finite-finite-variant}를 사용한다.
즉 이 체의 $\kappa(\mathfrak p)$ 위의 생성원
$\theta_i \in \kappa(\mathfrak q_i)_{sep}$를 택하자
(Fields, Lemma \ref{fields-lemma-primitive-element}).
올환 $S \otimes_R \kappa(\mathfrak p)$의 스펙트럼은 유한 이산이고,
그 점들은 $\mathfrak q_1, \ldots, \mathfrak q_n$에 대응한다.
중국인의 나머지 정리(Algebra, Lemma
\ref{algebra-lemma-chinese-remainder})에 의해
$S \otimes_R \kappa(\mathfrak p) \to \prod \kappa(\mathfrak q_i)$는
전사이다. 따라서 어떤 $g \in R$, $g \not \in \mathfrak p$에 대해
$R$을 $R_g$로 바꾼 뒤에는
$(0, \ldots, 0, \theta_i, 0, \ldots, 0) \in \prod \kappa(\mathfrak q_i)$
가 어떤 $x_i \in S$의 상이라고 가정할 수 있다. $S' \subset S$를
택한 $x_i$들이 생성하는 $R$-부분대수라 하자.
$\Spec(S) \to \Spec(S')$가 전사이므로(Algebra, Lemma
\ref{algebra-lemma-integral-overring-surjective}),
$\mathfrak q_i' = S' \cap \mathfrak q_i$는 $\mathfrak p$ 위에 놓인
$S'$의 소아이디얼들이다. $x_i$를 택한 방법에 의해 이 소아이디얼들은
% P05-EN-CH37-0190
서로 다르고
$\kappa(\mathfrak q'_i)_{sep} = \kappa(\mathfrak q_i)_{sep}$.
특히 체 확대 $\kappa(\mathfrak q_i)/\kappa(\mathfrak q'_i)$는 순수
비분리이다. $R \to S'$이 유한이므로
Algebra, Lemma \ref{algebra-lemma-etale-makes-quasi-finite-finite-variant}.
% P05-EN-CH37-0191
를 적용하여 $R \to R'$, $\mathfrak p'$, 그리고 분해
$$
S' \otimes_R R' = A'_1 \times \ldots \times A'_m \times B'
$$
를 얻는다. 여기서 정수적 사상 $R' \to A'_j$는 $\mathfrak p'$ 위에 놓인
유일한 소아이디얼 $\mathfrak r'_j$를 가지고,
$\kappa(\mathfrak r'_j)/\kappa(\mathfrak p')$는 순수 비분리이며,
$B'$은 $\mathfrak p'$ 위에 놓인 소아이디얼을 갖지 않는다.
$R \to S'$이 유한이므로 $R' \to B'$도 유한이다. 따라서 어떤
$g' \in R'$, $g' \not \in \mathfrak p'$에서 $R'$을 국소화한 뒤에는
$B' = 0$이라고 가정할 수 있다. 사상
$S' \otimes_R R' \to S \otimes_R R'$을 통해 $S$에 대한 대응하는
분해를 얻는다.
\end{proof}

\begin{proof}[둘째 증명]
이 증명에서는 강한 헨젤화를 사용한다. 먼저
% P05-EN-CH37-0192
$R$이 극대 아이디얼 $\mathfrak p$를 갖는 강한 헨젤화라고 가정하자.
그러면 $S/\mathfrak p S$는 $\mathfrak q_1, \ldots, \mathfrak q_n$에
대응하는 유한 개의 소아이디얼을 가지며, 이들은 모두 극대이고 각 잉여체는
$\kappa(\mathfrak p)$ 위에서 순수 비분리이다. 따라서
% P05-EN-CH37-0193
$S/\mathfrak p S$는 $\prod (S/\mathfrak p S)_{\mathfrak p_i}$와 같다.
More on Algebra, Lemma
\ref{more-algebra-lemma-characterize-henselian-pair}에 의해 이 곱분해를
주장에 나온 것과 같은 $S$의
% P05-EN-CH37-0194
곱 합성으로 올릴 수 있다.

\medskip\noindent
일반적인 경우에는 $R^{sh}$를 $R_\mathfrak p$의 강한 헨젤화라 하자.
그러면 첫째 문단의 결과를 $R^{sh} \to S \otimes_R R^{sh}$에 적용할 수
있다. 곱분해에 대응하는 $S \otimes_R R^{sh}$의 서로 직교하는 $m$개의
멱등원을 생각하자. Algebra, Lemma
\ref{algebra-lemma-strict-henselization-different}에 의해 $R^{sh}$는
\'etale 환 사상 $(R, \mathfrak p) \to (R', \mathfrak p')$의 여과
쌍극한이므로, 이 멱등원들은 원하는 대로 어떤 $R'$로 내려감을 알 수 있다.
\end{proof}







\section{Zariski의 주정리}
\label{more-morphisms-section-application-etale-neighbourhoods}

\noindent
이 절에서는 Grothendieck이 다시 정식화한 Zariski의 주정리를 증명한다.
% P05-EN-CH37-0195
% P05-EN-CH37-0196
이 문맥에서 ``Zariski의 주정리''라고 할 때는 흔히 Lemma
\ref{more-morphisms-lemma-finite-type-separated}, Lemma
\ref{more-morphisms-lemma-quasi-finite-separated-quasi-affine}, 또는 Lemma
\ref{more-morphisms-lemma-quasi-finite-separated-pass-through-finite} 가운데 하나를 뜻한다.
대부분의 문헌에서는 이 가운데 마지막 것을 Zariski의 주정리라고 부른다.

\medskip\noindent
우리는 이미 Algebra, Theorem \ref{algebra-theorem-main-theorem}에서
대수적 형태를 증명했고, 이를 스킴의 언어로도 다시 진술했다. 그 진술은
Morphisms, Theorem \ref{morphisms-theorem-main-theorem}에서 볼 수 있다. 이 절의 형태는 더
미묘하다. 완전한 결과를 얻기 위해 Section
\ref{more-morphisms-section-etale-localization}의 \'etale 국소화 기법을 사용하여
대수적 경우로 귀착시킨다.

\begin{lemma}
\label{more-morphisms-lemma-finite-type-separated}
$f : X \to S$를 스킴의 사상이라 하자. $f$가 유한형이고 분리라고
가정하자. $S'$을 $X$에서 $S$의 정규화라 하자. Morphisms, Definition
\ref{morphisms-definition-normalization-X-in-Y}을 보라. 도식은 다음과 같다.
$$
\xymatrix{
X \ar[rd]_f \ar[rr]_{f'} & & S' \ar[ld]^\nu \\
& S &
}
$$
그러면 다음을 만족하는 열린 부분스킴 $U' \subset S'$이 존재한다.
\begin{enumerate}
\item $(f')^{-1}(U') \to U'$은 동형사상이다.
\item $(f')^{-1}(U') \subset X$는 $f$가 준유한인 점들의 집합이다.
\end{enumerate}
\end{lemma}

\begin{proof}
Morphisms, Lemma \ref{morphisms-lemma-quasi-finite-points-open}에 의해
$f$가 준유한인 점들의 부분집합 $U \subset X$는 열려 있다.
보조정리는 다음과 동치이다.
\begin{enumerate}
\item[(a)] $U' = f'(U) \subset S'$은 열려 있다.
\item[(b)] $U = (f')^{-1}(U')$이다.
\item[(c)] $U \to U'$은 동형사상이다.
\end{enumerate}
임의의 $x \in U$를 택하자. $(f')^{-1}V \to V$가 동형사상이 되는 열린
근방 $f'(x) \in V \subset S'$이 존재한다고 주장한다. 먼저 이 주장이
보조정리를 함의함을 보이자. 실제로 그러면 $(f')^{-1}V \cong V$는
$X$의 열린 부분스킴이므로 $S$ 위에서 국소 유한형이고, 각 $v \in V$에
대해 잉여체 확대 $\kappa(v)/\kappa(\nu(v))$는 대수적이다. 이는
$V \subset S'$이고 $S'$이 $S$ 위에서 정수적이기 때문이다. 따라서
$V \to S$의 올은 이산이고(Morphisms, Lemma
\ref{morphisms-lemma-algebraic-residue-field-extension-closed-point-fibre}),
$(f')^{-1}V \to S$는 국소 준유한이다(Morphisms, Lemma
\ref{morphisms-lemma-locally-quasi-finite-fibres}). 이는
$(f')^{-1}V \subset U$와 $V \subset U'$을 함의한다. $x$가 임의였으므로
(a), (b), (c)가 성립한다.

\medskip\noindent
$s = f(x)$라 하자. $(T, t) \to (S, s)$를 기본 \'etale 근방이라 하자.
아래첨자 ${}_T$로 $T$로의 밑변환을 나타내자. $y = (x, t) \in X_T$를
$x$ 위에 놓인 올 $X_t$의 유일한 점이라 하자. $U_T \subset X_T$는
$f_T$가 준유한인 점들의 집합임에 유의하자. Morphisms, Lemma
\ref{morphisms-lemma-base-change-quasi-finite}를 보라. 또한
$$
X_T \xrightarrow{f'_T} S'_T \xrightarrow{\nu_T} T
$$
는 $X_T$에서 $T$의 정규화이다. Lemma
\ref{more-morphisms-lemma-normalization-smooth-localization}을 보라.
$y \in U_T \subset X_T \to S'_T \to T$에 대해 주장이 성립한다고 하자.
즉 $(f'_T)^{-1}V' \to V'$가 동형사상이 되는 열린 근방
$f'_T(y) \in V' \subset S'_T$을 찾을 수 있다고 하자. 사상
$S'_T \to S'$은 \'etale이므로 $V'$의 상 $V \subset S'$은 열려 있다.
$f'_T(y) \in V'$이므로 $f'(x) \in V$임에 유의하자. 도식
$$
\xymatrix{
(f'_T)^{-1}V' \ar[r] \ar[d] & (f')^{-1}(V) \ar[d] \\
V' \ar[r] & V
}
$$
은 올곱 도식이다. 실제로 $S'_T \times_{S'} X = X_T$이다.
왼쪽 세로 화살표는 동형사상이고 $\{V' \to V\}$는 \'etale 덮개이므로,
Descent, Lemma \ref{descent-lemma-descending-property-isomorphism}에 의해
오른쪽 세로 화살표도 동형사상이다. 다시 말해
$x \in U \subset X \to S' \to S$에 대해 주장이 성립한다.

\medskip\noindent
앞 문단의 결과에 의해 주장을 증명할 때 $S$를 $s = f(x)$의 기본
\'etale 근방으로 바꾸어도 된다. 따라서 분해
$$
X = V \amalg W
$$
가 존재한다고 가정할 수 있다. 여기서 각 항은 열리고 닫힌 부분스킴이고,
$V \to S$는 유한이며 $x \in V$이다. Lemma
\ref{more-morphisms-lemma-etale-splits-off-quasi-finite-part-technical}을 보라.
$X$는 $S$ 위에서 $V$와 $W$의 서로소 합이고 $V \to S$는 유한이므로,
$X$에서 $S$의 정규화는 사상
$$
X = V \amalg W \longrightarrow V \amalg W' \longrightarrow S
$$
이다. 여기서 $W'$은 $W$에서 $S$의 정규화이다. Morphisms, Lemmas
\ref{morphisms-lemma-normalization-in-disjoint-union},
\ref{morphisms-lemma-finite-integral},
\ref{morphisms-lemma-normalization-in-integral}을 보라.
주장이 따르므로 증명이 끝난다.
\end{proof}

\begin{lemma}
\label{more-morphisms-lemma-quasi-finite-separated-quasi-affine}
\begin{slogan}
준유한 분리 사상은 준아핀이다
\end{slogan}
$f : X \to S$를 스킴의 사상이라 하자. $f$가 준유한이고 분리라고
가정하자. $S'$을 $X$에서 $S$의 정규화라 하자. Morphisms, Definition
\ref{morphisms-definition-normalization-X-in-Y}을 보라. 도식은 다음과 같다.
$$
\xymatrix{
X \ar[rd]_f \ar[rr]_{f'} & & S' \ar[ld]^\nu \\
& S &
}
$$
그러면 $f'$은 준콤팩트 열린 몰입이고 $\nu$는 정수적이다.
특히 $f$는 준아핀이다.
\end{lemma}

\begin{proof}
이는 Lemma \ref{more-morphisms-lemma-finite-type-separated}에서 따른다. 실제로 그
보조정리에 의해 $(f')^{-1}(U') = X$이고 $X \to U'$이 동형사상인 열린
부분스킴 $U' \subset S'$이 존재한다. 다시 말해 $f'$은 열린 몰입이다.
$f$가 준콤팩트이고 $\nu : S' \to S$가 분리이므로(Schemes, Lemma
\ref{schemes-lemma-quasi-compact-permanence}) $f'$도 준콤팩트임에 유의하자.
따라서 Morphisms, Lemma
\ref{morphisms-lemma-characterize-quasi-affine}에 의해 $f$는 준아핀이다.
\end{proof}

\begin{lemma}[Zariski의 주정리]
\label{more-morphisms-lemma-quasi-finite-separated-pass-through-finite}
\begin{reference}
\cite[IV Corollary 18.12.13]{EGA}
\end{reference}
$f : X \to S$를 스킴의 사상이라 하자. $f$가 준유한이고 분리이며,
$S$가 준콤팩트이고 준분리라고 가정하자. 그러면 인수분해
$$
\xymatrix{
X \ar[rd]_f \ar[rr]_j & & T \ar[ld]^\pi \\
& S &
}
$$
가 존재한다. 여기서 $j$는 준콤팩트 열린 몰입이고 $\pi$는 유한이다.
\end{lemma}

\begin{proof}
$X \to S' \to S$를 Lemma
\ref{more-morphisms-lemma-quasi-finite-separated-quasi-affine}의 결론에 나온 것으로 택하자.
Properties, Lemma
\ref{properties-lemma-integral-algebra-directed-colimit-finite}에 의해
$\nu_*\mathcal{O}_{S'} = \colim_{i \in I} \mathcal{A}_i$를 유한 준연접
$\mathcal{O}_S$-대수 $\mathcal{A}_i \subset \nu_*\mathcal{O}_{S'}$들의
유향 쌍극한으로 쓸 수 있다. 그러면 각 $i$에 대해
$\pi_i : T_i = \underline{\Spec}_S(\mathcal{A}_i) \to S$는 유한 사상이다.
전이 사상 $T_{i'} \to T_i$는 아핀이고 $S' = \lim T_i$임에 유의하자.

\medskip\noindent
Limits, Lemma \ref{limits-lemma-descend-opens}에 의해 어떤 $i$와
$S'$에서의 역상이 $f'(X)$인 준콤팩트 열린 부분 $U_i \subset T_i$가
존재한다. $i' \geq i$에 대해 $U_{i'}$를 $T_{i'}$에서 $U_i$의 역상이라
하자. 그러면 $X \cong f'(X) = \lim_{i' \geq i} U_{i'}$이다. Limits,
Lemma \ref{limits-lemma-directed-inverse-system-has-limit}을 보라. Limits,
Lemma \ref{limits-lemma-finite-type-eventually-closed}에 의해 어떤
$i' \geq i$에 대해서는 $X \to U_{i'}$가 닫힌 몰입이다. 실제로 충분히
큰 $i'$에 대해 $X \cong U_{i'}$이지만, 이 사실은 필요하지 않다.
따라서 $X \to T_{i'}$는 몰입이다. Morphisms, Lemma
\ref{morphisms-lemma-factor-quasi-compact-immersion}에 의해 이를
$X \to T \to T_{i'}$로 인수분해할 수 있으며, 첫째 화살표는 열린 몰입이고
둘째 화살표는 닫힌 몰입이다. 이로써 증명이 끝난다.
\end{proof}

\begin{lemma}
\label{more-morphisms-lemma-quasi-finite-separated-pass-through-finite-addendum}
% P05-EN-CH37-0197
Lemma \ref{more-morphisms-lemma-quasi-finite-separated-pass-through-finite}의 표기와 가정
아래에서, $f$가 국소 유한 표시라고 더 가정하자. 그러면 $T$가 $S$ 위에서
유한이고 유한 표시가 되도록 인수분해를 택할 수 있다.
\end{lemma}

\begin{proof}
Limits, Lemma
\ref{limits-lemma-finite-in-finite-and-finite-presentation}에 의해
$T = \lim T_i$로 쓸 수 있다. 여기서 모든 $T_i$는 $S$ 위에서 유한이고
유한 표시이며, 전이 사상 $T_{i'} \to T_i$는 닫힌 몰입이다. Limits,
Lemma \ref{limits-lemma-descend-opens}에 의해 어떤 $i$와, $T$에서의 역상이
$X$인 열린 부분스킴 $U_i \subset T_i$가 존재한다. Limits, Lemma
\ref{limits-lemma-finite-type-eventually-closed}에 의해 충분히 큰 $i$에 대해
$X \cong U_i$이다. $T$를 $T_i$로 바꾸면 증명이 끝난다.
\end{proof}








\section{Zariski의 주정리의 응용, I}
\label{more-morphisms-section-applications-zmt}

\noindent
첫째 응용은 유한 사상을 올이 유한한 고유 사상으로 특징짓는 것이다.

\begin{lemma}
\label{more-morphisms-lemma-characterize-finite}
$f : X \to S$를 스킴의 사상이라 하자. 다음은 서로 동치이다.
\begin{enumerate}
\item $f$는 유한이다.
\item $f$는 고유이고 올이 유한하다.
\item $f$는 고유이고 국소 준유한이다.
\item $f$는 보편 닫힌이고 분리이며 국소 유한형이고 올이 유한하다.
\end{enumerate}
\end{lemma}

\begin{proof}
Morphisms, Lemmas \ref{morphisms-lemma-finite-proper},
\ref{morphisms-lemma-quasi-finite},
\ref{morphisms-lemma-finite-quasi-finite}에 의해 (1)은 (2)를 함의한다.
Morphisms, Lemma \ref{morphisms-lemma-finite-fibre}에 의해 (2)는 (3)을
함의한다. 고유 사상의 정의와 Morphisms, Lemmas
\ref{morphisms-lemma-quasi-finite-locally-quasi-compact},
\ref{morphisms-lemma-quasi-finite}에 의해 (3)은 (4)를 함의한다.

\medskip\noindent
(4)를 가정하고 $s \in S$를 택하자. Morphisms, Lemma
\ref{morphisms-lemma-finite-fibre}에 의해 $X_s$의 유한 개의 모든 점은
$X_s$에서 고립되어 있다. Lemma
\ref{more-morphisms-lemma-etale-splits-off-quasi-finite-part}에서처럼 기본 \'etale 근방
$(U, u) \to (S, s)$와 분해 $X_U = V \amalg W$를 택하자. $X_s$의 모든
점이 고립되어 있으므로 $W_u = \emptyset$이다. $f$가 보편 닫힌이므로
$U$에서 $W$의 상은 $u$를 포함하지 않는 닫힌 집합이다. $U$를 축소한
뒤에는 $W = \emptyset$이라고 가정할 수 있다. 다시 말해
$X_U = V$는 $U$ 위에서 유한이다. $s \in S$가 임의였으므로, 상들이
$S$를 덮고 밑변환 $X_{U_i} \to U_i$가 유한인 \'etale 사상들의 족
$\{U_i \to S\}$이 존재한다. $\{U_i \to S\}$은 \'etale 덮개이다.
Topologies, Definition \ref{topologies-definition-etale-covering}을 보라.
따라서 이는 fpqc 덮개이다. Topologies, Lemma
\ref{topologies-lemma-zariski-etale-smooth-syntomic-fppf-fpqc}를 보라.
그러므로 Descent, Lemma \ref{descent-lemma-descending-property-finite}에
의해 $f$는 유한이다.
\end{proof}

\noindent
그 결과 다음의 유용한 결과들을 얻는다.

\begin{lemma}
\label{more-morphisms-lemma-proper-finite-fibre-finite-in-neighbourhood}
$f : X \to S$를 스킴의 사상이라 하고 $s \in S$라 하자.
$f$가 고유이고 $f^{-1}(\{s\})$가 유한 집합이라고 가정하자. 그러면
$f|_{f^{-1}(V)} : f^{-1}(V) \to V$가 유한인 $s$의 열린 근방
$V \subset S$가 존재한다.
\end{lemma}

\begin{proof}
Morphisms, Lemma \ref{morphisms-lemma-finite-fibre}에 의해 사상 $f$는
$f^{-1}(\{s\})$의 모든 점에서 준유한이다. Morphisms, Lemma
\ref{morphisms-lemma-quasi-finite-points-open}에 의해 $f$가 준유한인 점들의
집합은 열린 부분 $U \subset X$이다. $Z = X \setminus U$라 하자. 그러면
$s \not \in f(Z)$이다. $f$가 고유이므로 집합 $f(Z) \subset S$는 닫혀
있다. $f(Z) \cap V = \emptyset$인 $s$의 임의의 열린 근방
$V \subset S$를 택하자. 그러면 $f^{-1}(V) \to V$는 국소 준유한이고
고유이다. 따라서 이는 준유한이고(Morphisms, Lemma
\ref{morphisms-lemma-quasi-finite-locally-quasi-compact}), 올이 유한하며
(Morphisms, Lemma \ref{morphisms-lemma-quasi-finite}), Lemma
\ref{more-morphisms-lemma-characterize-finite}에 의해 유한이다.
\end{proof}

\begin{lemma}
\label{more-morphisms-lemma-flat-proper-family-cannot-collapse-fibre}
스킴들의 가환 도식
$$
\xymatrix{
X \ar[rr]_h \ar[rd]_f & & Y \ar[ld]^g \\
& S
}
$$
을 생각하자. $s \in S$라 하고 다음을 가정하자.
\begin{enumerate}
\item $X \to S$는 고유 사상이다.
\item $Y \to S$는 분리이고 국소 유한형이다.
\item $X_s \to Y_s$의 상은 유한하다.
\end{enumerate}
그러면 $s$를 포함하는 열린
% P05-EN-CH37-0198
부분스킴 $U \subset S$가 존재하여 $X_U \to Y_U$는 $U$ 위에서 유한인
닫힌 부분스킴 $Z \subset Y_U$를 통해 인수분해된다.
\end{lemma}

\begin{proof}
$Z \subset Y$를 $h$의 스킴론적 상이라 하자. Morphisms, Section
\ref{morphisms-section-scheme-theoretic-image}을 보라. Morphisms, Lemma
\ref{morphisms-lemma-scheme-theoretic-image-is-proper}에 의해 사상
$X \to Z$는 전사이고 $Z \to S$는 고유이다. 따라서 $X_s \to Z_s$는
전사이다. (3)은 $Z_s$가 유한임을 함의한다.
% P05-EN-CH37-0199
따라서 Lemma \ref{more-morphisms-lemma-proper-finite-fibre-finite-in-neighbourhood}에 의해
$Z \to S$는 $s$의 어떤 열린 근방에서 유한이다.
\end{proof}





\section{Zariski의 주정리의 응용, II}
\label{more-morphisms-section-applications-zmt-II}

\noindent
이 절에서는 준유한 사상의 구조에 관한 Zariski의 주정리의 결과를
몇 가지 더 제시한다.

\begin{lemma}
\label{more-morphisms-lemma-separated-locally-quasi-finite-over-affine}
$f : X \to Y$를 분리이고 국소 준유한인 사상이라 하고 $Y$가 아핀이라고
하자. 그러면 $X$의 점들의 모든 유한 집합은 $X$의 어떤 아핀 열린 부분에
포함된다.
\end{lemma}

\begin{proof}
$x_1, \ldots, x_n \in X$라 하자. $x_i \in U$인 준콤팩트 열린 부분
$U \subset X$를 택하자. Lemma
\ref{more-morphisms-lemma-quasi-finite-separated-quasi-affine}에 의해 $U \to Y$는
준아핀이다. 따라서 Properties, Lemma
\ref{properties-lemma-ample-finite-set-in-affine}에 의해
$x_1, \ldots, x_n$을 포함하는 아핀 열린 부분 $V \subset U$가 존재한다.
\end{proof}

\begin{lemma}
\label{more-morphisms-lemma-quasi-finite-finite-over-dense-open}
$f : Y \to X$를 준유한 사상이라 하자. 그러면
$f|_{f^{-1}(U)} : f^{-1}(U) \to U$가 유한이 되는 조밀 열린 부분
$U \subset X$가 존재한다.
\end{lemma}

\begin{proof}
$U_i \subset X$, $i \in I$가 제한
$f|_{f^{-1}(U_i)} : f^{-1}(U_i) \to U_i$가 유한인 열린 부분들의 모임이면,
$U = \bigcup U_i$에 대한 제한 $f|_{f^{-1}(U)} : f^{-1}(U) \to U$도
유한이다. Morphisms, Lemma \ref{morphisms-lemma-finite-local}을 보라.
따라서 문제는 $X$ 위에서 국소적이고, $X$가 아핀이라고 가정할 수 있다.

\medskip\noindent
$X$가 아핀이라고 가정하자. 각 $V_j$가 아핀이 되도록
$Y = \bigcup_{j = 1, \ldots, m} V_j$로 쓰자. $f$가 준유한이고, 특히
준콤팩트이므로 이렇게 할 수 있다. 각 $V_j \to X$는 준유한이고 분리이다.
$\eta \in X$를 $X$의 한 기약 성분의 일반점이라 하자. Morphisms, Lemmas
\ref{morphisms-lemma-quasi-finite},
\ref{morphisms-lemma-generically-finite}에 의해
$f^{-1}(U_\eta) \cap V_j \to U_\eta$가 유한인 열린 근방
$\eta \in U_\eta$가 존재한다. 각 $j = 1, \ldots, m$에 대해 성립하도록
$U_\eta$를 택할 수 있다. $X$의 일반점들의 모임은 $X$에서 조밀함에
유의하자. 따라서 각 $f^{-1}(W) \cap V_j \to W$가 유한인 조밀 열린 부분
$W = \bigcup_\eta U_\eta$가 존재한다. 제한
$f|_{f^{-1}(U)} : f^{-1}(U) \to U$가 유한인 조밀 열린 부분
$U \subset W$가 존재함을 보이면 충분하다. 그러므로 $X$를 $W$의 아핀
열린 부분스킴으로 바꾸고 각 $V_j \to X$가 유한이라고 가정할 수 있다.

\medskip\noindent
$X$가 아핀이고, 각 $V_j$가 아핀인
$Y = \bigcup_{j = 1, \ldots, m} V_j$가 주어지며, 제한
$f|_{V_j} : V_j \to X$가 유한이라고 가정하자. 다음과 같이 두자.
$$
\Delta_{ij} =
\Big(\overline{V_i \cap V_j} \setminus V_i \cap V_j\Big) \cap V_j.
$$
이는 $V_j$ 안의 열린 부분집합 $V_i \cap V_j$의 경계이므로 $V_j$의
어디에서도 조밀하지 않은 닫힌 부분집합이다. Morphisms, Lemma
\ref{morphisms-lemma-image-nowhere-dense-finite}에 의해 상
$f(\Delta_{ij})$는 $X$의 어디에서도 조밀하지 않은 닫힌 부분집합이다.
Topology, Lemma \ref{topology-lemma-nowhere-dense}에 의해 합집합
$T = \bigcup f(\Delta_{ij})$도 $X$의 어디에서도 조밀하지 않은 닫힌
부분집합이다. 따라서 $U = X \setminus T$는 $X$의 조밀 열린 부분집합이다.
$f|_{f^{-1}(U)} : f^{-1}(U) \to U$가 유한이라고 주장한다. 이를 보이기
위해 아핀 열린 부분 $U' \subset U$를 택하자. 다음과 같이 두자.
$Y' = f^{-1}(U') = U' \times_X Y$,
$V_j' = Y' \cap V_j = U' \times_X V_j$. $f$의 제한
$$
f' = f|_{Y'} : Y' \longrightarrow U'
$$
을 생각하자. 이 사상에 대해 $Y' = \bigcup_{j = 1, \ldots, m} V'_j$는
아핀 열린 덮개이고, 각 $V'_j \to U'$은 유한이며,
$V_i' \cap V_j'$는 $V'_i$와 $V'_j$ 양쪽에서 (열리고) 닫혀 있다.
따라서 $V_i' \cap V_j'$는 아핀이고, 사상
$$
\mathcal{O}(V'_i) \otimes_{\mathbf{Z}} \mathcal{O}(V'_j)
\longrightarrow
\mathcal{O}(V'_i \cap V'_j)
$$
은 전사이다. 이는 $Y'$이 분리임을 함의한다. Schemes, Lemma
\ref{schemes-lemma-characterize-separated}를 보라. 마지막으로 가환 도식
$$
\xymatrix{
\coprod_{j = 1, \ldots, m} V'_j \ar[rd] \ar[rr] & & Y' \ar[ld] \\
& U' &
}
$$
을 생각하자. 남동쪽 화살표는 유한이고 따라서 고유이며, 가로 화살표는
전사이고, 남서쪽 화살표는 분리이다. 그러므로 Morphisms, Lemma
\ref{morphisms-lemma-image-proper-is-proper}에 의해 $Y' \to U'$은 고유이다.
이 사상은 또한 준유한이므로 Lemma \ref{more-morphisms-lemma-characterize-finite}에 의해
유한이다. 이로써 증명이 끝난다.
\end{proof}

\begin{lemma}
\label{more-morphisms-lemma-stratify-flat-fp-lqf-universally-bounded}
$f : X \to S$가 평탄하고 국소 유한 표시이며 분리이고 국소 준유한이며,
그 올들이 보편적으로 유계라고 하자. 그러면 닫힌 부분집합들
$$
\emptyset = Z_{-1} \subset Z_0 \subset Z_1 \subset Z_2 \subset
\ldots \subset Z_n = S
$$
이 존재한다. $S_r = Z_r \setminus Z_{r - 1}$로 두었을 때 층화
$S = \coprod_{r = 0, \ldots, n} S_r$는 다음 보편 성질로 특징지어진다.
$g : T \to S$가 주어지면, 사영 $X \times_S T \to T$가 차수 $r$인
유한 국소 자유일 필요충분조건은 $g(T) \subset S_r$가 집합론적으로
성립하는 것이다.
\end{lemma}

\begin{proof}
$n$을 $X \to S$의 올들의 차수를 제한하는 정수라 하자. Morphisms, Lemma
\ref{morphisms-lemma-base-change-universally-bounded}에 의해 임의의
밑변환에서도 올들의 차수는 $n$으로 제한된다. 특히 보조정리의 명제에
나타나는 모든 정수 $r$은 $\leq n$이다. $n$에 대한 귀납법으로
보조정리를 증명한다. 시작 경우 $n = 0$은 자명하다.

\medskip\noindent
$\deg_{\kappa(s)}(X_s) = n$인 점 $s \in S$들의 집합은 열린 부분집합
$S_n \subset S$이고 $X \times_S S_n \to S_n$은 차수 $n$인 유한 국소
자유라고 주장한다. 실제로 $s \in S$가 그런 점이라고 하자. Lemma
\ref{more-morphisms-lemma-etale-splits-off-quasi-finite-part}에서처럼 기본 \'etale 사상
$(U, u) \to (S, s)$와 분해 $U \times_S X = W \amalg V$를 택하자.
$V \to U$가 유한이고 평탄하며 국소 유한 표시이므로, $V \to U$는
유한 국소 자유이다. Morphisms, Lemma \ref{morphisms-lemma-finite-flat}을 보라.
$U$를 $u$의 더 작은 근방으로 축소한 뒤 $V \to U$가 어떤 차수 $d$인
유한 국소 자유라고 가정할 수 있다. Morphisms, Lemma
\ref{morphisms-lemma-finite-locally-free}를 보라. $u \mapsto s$이고
$W_u = \emptyset$이므로 $d = n$이다. $n$이 올의 최대 차수이므로
$W = \emptyset$이다. 따라서 $U \times_S X \to U$는 차수 $n$인 유한
국소 자유이다. Descent, Lemma
\ref{descent-lemma-descending-property-finite-locally-free}에 의해 $X \to S$는
$s$의 열린 근방 $\Im(U \to S)$ 위에서 차수 $n$인 유한 국소 자유이다
(Morphisms, Lemma \ref{morphisms-lemma-etale-open}). 이로써 주장이 증명된다.

\medskip\noindent
$S' = S \setminus S_n$에 유도된 축약 스킴 구조를 주고
$X' = X \times_S S'$라 하자. $X' \to S'$의 올들의 차수는 보편적으로
$n - 1$로 제한된다. 귀납법에 의해 사상 $X' \to S'$에 맞는 층화
$S' = S_0 \amalg \ldots \amalg S_{n - 1}$을 얻는다. 층화
$S = \coprod_{r = 0, \ldots, n} S_r$가 사상 $X \to S$에 대해 원하는
성질을 갖는다고 주장한다. $g : T \to S$를 스킴의 사상이라 하고
$X \times_S T \to T$가 차수 $r$인 유한 국소 자유라고 가정하자. 위에서
말했듯이 이는 $r \leq n$을 함의한다. $r = n$이면 $T \to S$가 $S_n$을
통해 인수분해됨은 명백하다. $r < n$이면 집합론적으로
% P05-EN-CH37-0200
$g(T) \subset S' = S \setminus S_d$이므로 $T_{red} \to S$는 $S'$을 통해
인수분해된다. Schemes, Lemma \ref{schemes-lemma-map-into-reduction}을 보라.
$X \times_S T_{red} \to T_{red}$도 밑변환으로서 차수 $r$인 유한 국소
자유임에 유의하자. 층화 $S' = \coprod_{r = 0, \ldots, n - 1} S_r$의
보편 성질에 의해 $g(T) = g(T_{red})$는 $S_r$에 포함된다.
거꾸로 집합론적으로 $g(T) \subset S_r$인 $g : T \to S$가 주어졌다고
하자. $r = n$이면 $g$는 $S_n$을 통해 인수분해되고, 밑변환
$X \times_S T \to T$가 차수 $n$인 유한 국소 자유임은 명백하다.
$r < n$이면 $X \times_S T \to T$는 분리이고 평탄하며 국소 유한 표시인
사상이고, $T_{red}$로의 제한은 차수 $r$인 유한 국소 자유이다.
$T_{red} \to T$가 보편 위상동형사상이므로
$X \times_S T_{red} \to X \times_S T$도 보편 위상동형사상이다. 따라서
$X \times_S T \to T$는 보편 닫힌이다. 실제로 유한 사상
$X \times_S T_{red} \to T_{red}$에 대해 이 성질이 성립한다. 그러므로
예를 들어 Lemma \ref{more-morphisms-lemma-characterize-finite}에 의해
$X \times_S T \to T$는 유한이다. 이제 Morphisms, Lemma
\ref{morphisms-lemma-finite-flat}에 의해 $X \times_S T \to T$는 유한
국소 자유이다. 마지막으로 모든 올의 차수가 $r$이므로 그 차수는 $r$이다.
\end{proof}

\begin{lemma}
\label{more-morphisms-lemma-stratify-flat-fp-qf}
$f : X \to S$를 평탄하고 국소 유한 표시이며 분리이고 준유한인 스킴의
사상이라 하자. 그러면 닫힌 부분집합들
$$
\emptyset = Z_{-1} \subset Z_0 \subset Z_1 \subset Z_2 \subset
\ldots \subset S
$$
이 존재한다. $S_r = Z_r \setminus Z_{r - 1}$로 두었을 때 층화
$S = \coprod S_r$는 다음 보편 성질로 특징지어진다. 사상
$g : T \to S$가 주어지면, 사영 $X \times_S T \to T$가 차수 $r$인
유한 국소 자유일 필요충분조건은 $g(T) \subset S_r$가 집합론적으로
성립하는 것이다. 더욱이 포함 사상 $S_r \to S$는 준콤팩트이다.
\end{lemma}

\begin{proof}
문제는 $S$ 위에서 국소적이므로 $S$가 아핀이라고 가정할 수 있다.
이 경우 Morphisms, Lemma
\ref{morphisms-lemma-locally-quasi-finite-qc-source-universally-bounded}에
의해 $f$의 올들은 보편적으로 유계이다. 따라서 층화의 존재는 Lemma
\ref{more-morphisms-lemma-stratify-flat-fp-lqf-universally-bounded}에서 따른다.

\medskip\noindent
각 $r \geq 0$에 대해 $U_r = S \setminus Z_r \to S$가 준콤팩트임을
보이자. 그러면 초등적인 위상론으로 마지막 주장이 증명된다. 준콤팩트
사상들의 합성은 준콤팩트이므로 $U_r \to U_{r - 1}$이 준콤팩트임을
보이면 충분하다. 아핀 열린 부분 $W \subset U_{r - 1}$을 택하고
$W = \Spec(A)$라 쓰자. 그러면 어떤 아이디얼 $I \subset A$에 대해
$Z_r \cap W = V(I)$이고, $X \times_S \Spec(A/I) \to \Spec(A/I)$는
차수 $r$인 유한 국소 자유이다. $I_i \subset I$가 유한 생성
아이디얼들을 달릴 때 $A/I = \colim A/I_i$임에 유의하자. Limits, Lemma
\ref{limits-lemma-descend-finite-locally-free}에 의해 어떤 $i$에 대해
$X \times_S \Spec(A/I_i) \to \Spec(A/I_i)$는 차수 $r$인 유한 국소
자유이다. 여기서는 $X \to S$가 유한 표시라는 사실을 사용한다. 실제로
이는 국소 유한 표시이고 분리이며 준콤팩트이다. 따라서
$\Spec(A/I_i) \to \Spec(A) = W$는 집합론적으로 $Z_r \cap W$를 통해
인수분해된다. 그러므로 $Z_r \cap W = V(I_i)$는 $A$의 원소들의 어떤
유한 부분집합의 영점집합이다. 이는 $W \setminus Z_r$가 표준 열린집합들의
유한 합집합이고, 따라서 원하는 대로 준콤팩트임을 뜻한다.
\end{proof}

\begin{lemma}
\label{more-morphisms-lemma-stratify-flat-fp-lqf}
$f : X \to S$를 평탄하고 국소 유한 표시이며 분리이고 국소 준유한인
스킴의 사상이라 하자. 그러면 열린 부분스킴들
$$
S = U_0 \supset U_1 \supset U_2 \supset \ldots
$$
이 존재하여, 체 $k$에 대한 사상 $\Spec(k) \to S$가 $U_d$를 통해
인수분해될 필요충분조건은 $X \times_S \Spec(k)$의 $k$ 위 차수가
$\geq d$인 것이다.
\end{lemma}

\begin{proof}
이 명제는 올의 차수가 $\geq d$인 점들의 모임이 열려 있다는 뜻일
뿐이다. 따라서 $S$ 위에서 국소적으로 작업하여 $S$가 아핀이라고 가정할
수 있다. 이 경우 각 준콤팩트 열린 부분 $W \subset X$에 대해,
$W \to S$의 올의 차수가 $\geq d$인 점들의 집합 $U_d(W)$는 Lemma
\ref{more-morphisms-lemma-stratify-flat-fp-qf}에 의해 열려 있다.
$U_d = \bigcup_W U_d(W)$이므로 결과가 따른다.
\end{proof}

\begin{lemma}
\label{more-morphisms-lemma-go-down-with-annihilators}
$f : X \to S$를 평탄하고 국소 유한 표시이며 국소 준유한인 스킴의
사상이라 하자.
% P05-EN-CH37-0201
$g \in \Gamma(X, \mathcal{O}_X)$가 영이 아니라고 하자. 그러면
$g|_V \not = 0$인 열린 부분 $V \subset X$, 다음 가환 도식에 들어맞는
열린 부분 $U \subset S$,
$$
\xymatrix{
V \ar[r] \ar[d]_\pi & X \ar[d]^f \\
U \ar[r] & S,
}
$$
준연접 부분층 $\mathcal{F} \subset \mathcal{O}_U$, 정수 $r > 0$,
그리고 상이 $g|_V$를 포함하는 단사 $\mathcal{O}_U$-가군 사상
$\mathcal{F}^{\oplus r} \to \pi_*\mathcal{O}_V$
가 존재한다.
\end{lemma}

\begin{proof}
우리는 $X$와 $S$가 아핀이라고 가정해도 된다. Lemmas
\ref{more-morphisms-lemma-stratify-flat-fp-lqf-universally-bounded} 및
\ref{more-morphisms-lemma-stratify-flat-fp-qf}에서와 같은 여과
$\emptyset = Z_{-1} \subset Z_0 \subset Z_1 \subset Z_2 \subset \ldots
\subset Z_n = S$를 얻는다.
$T \subset X$를 유한 $\mathcal{O}_X$-가군
$\Im(g : \mathcal{O}_X \to \mathcal{O}_X)$의 스킴론적 받침이라 하자.
$T$는 $\mathcal{O}_X$의 단면으로서의 $g$의 받침이고
(Modules, Definition \ref{modules-definition-support}), 임의의 열린 부분
$V \subset X$에 대해 $g|_V \not = 0$일 필요충분조건은
$V \cap T \not = \emptyset$임에 유의하자.
$f(T) \subset Z_r$가 집합론적으로 성립하도록 하는 가장 작은 정수를
$r$이라 하자. $f(\xi) \not \in Z_{r - 1}$인(따라서
$f(\xi) \in Z_r$인) $T$의 한 기약 성분의 일반점을
$\xi \in T$라 하자. $S$를 $S \setminus Z_{r - 1}$에 포함된
$f(\xi)$의 아핀 근방으로 대체할 수 있다. $S = \Spec(A)$라 쓰고,
$V(I) = Z_r$가 집합론적으로 성립하는 유한 생성 아이디얼
$I = (a_1, \ldots, a_m) \subset A$를 택하자(Algebra, Lemma
\ref{algebra-lemma-qc-open} 참조).
$r$의 선택에 의해 $g$의 받침은 $f^{-1}V(I)$에 포함되므로,
$j = 1, \ldots, m$에 대해 $a_j^N g = 0$인 정수 $N$이 존재한다.
% P05-EN-CH37-0202
$a_j$를 $a_j^r$로 대체하여 $Ig = 0$이라고 가정해도 된다.
임의의 $A$-가군 $M$에 대해 $M[I]$로 $M$의 $I$-꼬임 부분,
즉 $M[I] = \{m \in M \mid Im = 0\}$을 나타내자.
$X = \Spec(B)$라 쓰면 $g \in B[I]$이다. $A \to B$가 평탄하므로
$$
B[I] = A[I] \otimes_A B \cong A[I] \otimes_{A/I} B/IB
$$
이다. $Z_r$의 선택에 의해 $A/I$-가군 $B/IB$는 계수 $r$인 유한 국소
자유 가군이다. 따라서 $S$를 $f(\xi)$의 더 작은 아핀 열린 근방으로
대체하여, $A/I$-가군으로서
$B/IB \cong (A/IA)^{\oplus r}$이라고 가정해도 된다.
$I$로 나눈 뒤 앞 문장의 동형사상이 되는 사상
$\psi : A^{\oplus r} \to B$를 택하자. 그러면 유도된 사상
$$
A[I]^{\oplus r} \longrightarrow B[I]
$$
은 동형사상이다. $\mathcal{F}$를 $A$-가군 $A[I]$에 대응하는 준연접
층으로 택하고, $\mathcal{F}^{\oplus r} \to \pi_*\mathcal{O}_V$를
$A[I]^{\oplus r} \subset A^{\oplus r} \to B$에 대응하는 사상으로
택하면 보조정리가 따른다.
\end{proof}

\begin{lemma}
\label{more-morphisms-lemma-there-is-a-scheme-integral-over}
$U \to X$를 스킴의 전사 \'etale 사상이라 하자. $X$가 준콤팩트이고
준분리라고 가정하자. 그러면 임의의 $y \in Y$에 대해 $V \to X$가
$U$를 통해 분해되는 열린 근방 $V \subset Y$가 존재하도록 하는
전사 정수적 사상 $Y \to X$가 존재한다. 실제로 $Y \to X$가 유한이고
유한 표시라고 가정할 수 있다.
\end{lemma}

\begin{proof}
$X$가 준콤팩트이므로 $U' = \coprod U_i \to X$가 전사이도록 하는
유한 개의 아핀 열린 부분 $U_i \subset U$가 존재한다. $U$를 $U'$로
대체하여 $U$가 아핀이라고 가정해도 된다. 특히 $U \to X$는
분리 사상이다(Schemes, Lemma \ref{schemes-lemma-affine-separated}).
그러면 $U \to X$의 기하학적 올들의 차수에 대한 상계인 정수 $d$가
존재한다(Morphisms, Lemma
\ref{morphisms-lemma-locally-quasi-finite-qc-source-universally-bounded}).
$X$ 위로 전사이고 \'etale인 사상을 갖는 모든 준콤팩트 분리 스킴
$U$에 대해 $d$에 관한 귀납법으로 보조정리를 증명하겠다.
$d = 1$이면 $U = X$이고 $Y = U$로 두어 결과가 성립한다.
$d > 1$이라고 가정하자.

\medskip\noindent
Lemma \ref{more-morphisms-lemma-quasi-finite-separated-quasi-affine}를 적용하여 분해
$$
\xymatrix{
U \ar[rr]_j \ar[rd] & & Y \ar[ld]^\pi \\
& X
}
$$
를 얻는데, 여기서 $\pi$는 정수적이고 $j$는 준콤팩트 열린 몰입이다.
$j(U)$가 $Y$에서 스킴론적으로 조밀하다고 가정할 수 있고 그렇게 하자.
다음에 유의하자.
$$
U \times_X Y = U \amalg W
$$
여기서 첫 번째 성분은 $U \to U \times_X Y$의 상이고(Schemes, Lemma
\ref{schemes-lemma-semi-diagonal}에 의해 닫혀 있으며, $Y$ 위에서
\'etale인 스킴들 사이의 사상으로서 \'etale이므로 열려 있다), 두 번째
성분은 그 열리고 닫힌 여집합이다. $W$의 상 $V \subset Y$는
$Y \setminus U$를 포함하는 열린 부분스킴이다.

\medskip\noindent
\'Etale 사상 $W \to Y$의 기하학적 올들의 기수는 $< d$이다.
실제로 $U \subset Y$의 기하학적 점들에 대해서는 직접 보면 명백하다.
$U \subset Y$가 조밀하므로 이는 예를 들어 Lemma
\ref{more-morphisms-lemma-stratify-flat-fp-lqf-universally-bounded}
에 의해 $Y$의 모든 기하학적 점에 대해서도 성립한다(준콤팩트 분리
\'etale 사상의 올들의 차수는 특수화에서 증가하지 않는다). 따라서
$W \to V$에 귀납 가정을 적용하여, $Z$가 스킴이고 Zariski 국소적으로
$W$를 통해 분해되는 전사 정수적 사상 $Z \to V$를 찾을 수 있다.
$Z' \to Y$가 정수적이고 $Z \to Z'$가 열린 몰입인 분해
$Z \to Z' \to Y$를 택하자(Lemma
\ref{more-morphisms-lemma-quasi-finite-separated-quasi-affine}). $Z'$를 $Z'$ 안에서
$Z$의 스킴론적 폐포로 대체하여 $Z$가 $Z'$에서 스킴론적으로 조밀하다고
가정해도 된다. 이렇게 하면 $Z' \times_Y V = Z$이다. 마지막으로
$T \subset Y$를 $Y \setminus V$ 위에 유도된 축약 닫힌 부분스킴
구조라 하자. 사상
$$
Z' \amalg T \longrightarrow X
$$
을 생각하자. 이는 구성에 의해 전사 정수적 사상이다.
$T \subset U$이므로 사상 $T \to X$가 $U$를 통해 분해됨은 명백하다.
한편 점 $z \in Z'$를 택하자. $z \not \in Z$이면 $z$는
$Y \setminus V \subset U$의 한 점으로 가며, 사상이 $U$를 통해
분해되는 $z$의 근방을 찾을 수 있다. $z \in Z$이면
$W \subset U \times_X Y$를 통해, 따라서 $U$를 통해 분해되는 근방
$\Omega \subset Z$가 존재한다. 이로써 존재성이 증명되었다.

\medskip\noindent
Zariski 국소적으로 $U$를 통해 분해되는 정수적 전사 사상
$Y \to X$를 찾았다고 가정하자. $V_j \to X$가 $U$를 통해 분해되는
유한 아핀 열린 덮개 $Y = \bigcup V_j$를 택하자. $Y_i \to X$가 유한이고
유한 표시이도록 $Y = \lim Y_i$라고 쓸 수 있다(Limits, Lemma
\ref{limits-lemma-integral-limit-finite-and-finite-presentation}).
충분히 큰 $i$에 대해 $Y$로의 역상이 $V_j$가 되는 아핀 열린 부분
$V_{i, j} \subset Y_i$를 찾을 수 있다(Limits, Lemma
\ref{limits-lemma-descend-opens}). $i$를 더 크게 하면 $X$ 위의 사상
$V_j \to U$는 $X$ 위의 사상 $V_{i, j} \to U$에서 온다(Limits,
Proposition
\ref{limits-proposition-characterize-locally-finite-presentation}).
이로써 증명이 끝난다.
\end{proof}







\section{연결 올을 갖는 사상에의 응용}
\label{more-morphisms-section-application-geometrically-connected}

\noindent
이 절에서는 모든 올이 기하학적으로 연결이거나 기하학적으로 정역인
사상을 만드는 몇 가지 보조정리를 증명한다. 이는 뒤에서 유한형 사상의
국소 구조를 연구할 때 유용할 것이다.

\begin{lemma}
\label{more-morphisms-lemma-descent-connected-fibres}
스킴 사상들의 도식
$$
\xymatrix{
Z \ar[r]_{\sigma} \ar[rd] & X \ar[d] \\
& Y
}
$$
% P05-EN-CH37-0203
과 점 $y \in Y$를 생각하자. 다음을 가정하자.
\begin{enumerate}
\item $X \to Y$는 유한 표시이고 평탄하다.
\item $Z \to Y$는 유한 국소 자유이다.
\item $Z_y \not = \emptyset$이다.
\item $X \to Y$의 모든 올은 기하학적으로 축약이다.
\item $X_y$는 $\kappa(y)$ 위에서 기하학적으로 연결이다.
\end{enumerate}
그러면 $X^0_y = X_y$이고 $X^0 \to Y$의 모든 비어 있지 않은 올이
기하학적으로 연결이 되는 준콤팩트 열린 부분 $X^0 \subset X$가 존재한다.
\end{lemma}

\begin{proof}
이 증명에서는 평탄, 유한 표시, 유한 국소 자유라는 성질이 밑변환과
합성에 의해 보존된다는 사실을 사용한다. 또한 유한 국소 자유 사상이
열린 사상이면서 닫힌 사상이라는 사실도 사용한다. 이 사실들은
Morphisms, Lemmas
\ref{morphisms-lemma-base-change-flat},
\ref{morphisms-lemma-base-change-finite-presentation},
\ref{morphisms-lemma-base-change-finite-locally-free},
\ref{morphisms-lemma-composition-flat},
\ref{morphisms-lemma-composition-finite-presentation},
\ref{morphisms-lemma-composition-finite-locally-free},
\ref{morphisms-lemma-fppf-open}, 그리고
\ref{morphisms-lemma-finite-proper}에 있다.

\medskip\noindent
% P05-EN-CH37-0204
$X_Z \to Z$는 $\sigma$에서 오는 단면 $s$를 갖는 유한 표시 평탄
사상임에 유의하자. $X_Z^0$로 Situation
\ref{more-morphisms-situation-connected-along-section}에서 정의한 $X_Z$의 부분집합을
나타내자. Lemma \ref{more-morphisms-lemma-connected-along-section-open}에 의해 이는
$X_Z$의 열린 부분집합이다.

\medskip\noindent
$X$를 $Z \times_Y Z$로 당겨 되얻은 $X_{Z \times_Y Z}$에는 두 단면
$s_0, s_1$이 주어지는데, 이들은 $s$를
$\text{pr}_0, \text{pr}_1 : Z \times_Y Z \to Z$로 밑변환한 것이다.
Situation \ref{more-morphisms-situation-connected-along-section}의 구성은 두 부분집합
$(X_{Z \times_Y Z})_{s_0}^0$와 $(X_{Z \times_Y Z})_{s_1}^0$를 준다.
Lemma \ref{more-morphisms-lemma-base-change-connected-along-section}에 의해 이들은 각각
사상 $1_X \times \text{pr}_0, 1_X \times \text{pr}_1 :
X_{Z \times_Y Z} \to X_Z$에 의한 $X_Z^0$의 역상이다. 특히 이
부분집합들은 열려 있다.

\medskip\noindent
$(Z \times_Y Z)_y = \{z_1, \ldots, z_n\}$이라 하자.
$X_y$가 기하학적으로 연결이므로 $(X_{Z \times_Y Z})_{s_0}^0$와
$(X_{Z \times_Y Z})_{s_1}^0$의 각 $z_i$ 위 올들은 일치한다(전체 올과
같다). 이를 달리 말하면
$$
s_0(z_i) \in (X_{Z \times_Y Z})_{s_1}^0
\quad\text{and}\quad
s_1(z_i) \in (X_{Z \times_Y Z})_{s_0}^0.
$$
이다. 집합 $(X_{Z \times_Y Z})_{s_0}^0$와
$(X_{Z \times_Y Z})_{s_1}^0$가 $X_{Z \times_Y Z}$에서 열려 있으므로,
다음을 만족하는 $(Z \times_Y Z)_y$의 열린 근방
$W \subset Z \times_Y Z$가 존재한다.
$$
s_0(W) \subset (X_{Z \times_Y Z})_{s_1}^0
\quad\text{and}\quad
s_1(W) \subset (X_{Z \times_Y Z})_{s_0}^0.
$$
그러면 Situation \ref{more-morphisms-situation-connected-along-section}의 구성에서
곧바로
$$
p^{-1}(W) \cap (X_{Z \times_Y Z})_{s_0}^0
=
p^{-1}(W) \cap (X_{Z \times_Y Z})_{s_1}^0
$$
가 따른다. 여기서
% P05-EN-CH37-0205
$p : X_{Z \times_Y Z} \to Z \times_W Z$는 사영이다.
$Z \times_Y Z \to Y$가 유한 국소 자유이고, 따라서 열리고 닫힌
사상이므로, $q^{-1}(V) \subset W$를 만족하는 $y$의 아핀 열린 근방
$V \subset Y$가 존재한다. 여기서 $q : Z \times_Y Z \to Y$는 구조
사상이다. 보조정리를 증명하기 위해 $Y$를 $V$로 대체해도 된다.
이렇게 한 뒤,
% P05-EN-CH37-0206
$X_Z^0 \subset Y_Z$는 다음을 만족하는 열린 부분임을 알 수 있다.
$$
(1_X \times \text{pr}_0)^{-1}(X_Z^0) =
(1_X \times \text{pr}_1)^{-1}(X_Z^0).
$$
이는 $X_Z^0$의 상 $X^0 \subset X$가
$(X_Z \to X)^{-1}(X^0) = X_Z^0$을 만족하는 열린 부분이라는 뜻이다.
Descent, Lemma \ref{descent-lemma-open-fpqc-covering}을 보라.
마지막으로 Lemma \ref{more-morphisms-lemma-connected-along-section-locally-constructible}에
의해 $X_Z^0$가 준콤팩트이므로 $X^0$도 준콤팩트이다(이 단계에서 $Y$가
아핀이므로 $X$가 준콤팩트이고 준분리이며, 따라서 국소 구성 가능은
구성 가능과 같고 특히 준콤팩트라는 사실을 사용한다. 자세한 내용은
생략한다). 이로써 $X^0$가 원하는 모든 성질을 가짐을 알 수 있다.
\end{proof}

\begin{lemma}
\label{more-morphisms-lemma-fibre-geometrically-connected-reduced}
$h : Y \to S$를 스킴의 사상이라 하고, $s \in S$를 점이라 하며,
$T \subset Y_s$를 열린 부분스킴이라 하자. 다음을 가정하자.
\begin{enumerate}
\item $h$는 평탄하고 유한 표시이다.
\item $h$의 모든 올은 기하학적으로 축약이다.
\item $T$는 $\kappa(s)$ 위에서 기하학적으로 연결이다.
\end{enumerate}
그러면 아핀 기본 \'etale 근방
$(S', s') \to (S, s)$
와 다음을 만족하는 준콤팩트 열린 부분 $V \subset Y_{S'}$를 찾을 수 있다.
\begin{enumerate}
\item[(a)] $V \to S'$의 모든 올은 기하학적으로 연결이다.
\item[(b)] $V_{s'} = T \times_s s'$이다.
\end{enumerate}
\end{lemma}

\begin{proof}
문제는 명백히 $S$ 위에서 국소적이므로 $S$를 $s$의 아핀 열린 근방으로
대체해도 된다. $Y_s$의 위상은 $Y$의 위상에서 유도된다. Schemes,
Lemma \ref{schemes-lemma-fibre-topological}을 보라. 따라서
$V_s = T$를 만족하는 준콤팩트 열린 부분 $V \subset Y$를 찾을 수 있다.
$h$의 $V$로의 제한은 준콤팩트이고($S$가 아핀이고 $V$가 준콤팩트이므로),
준분리이며, 국소 유한 표시이고, 평탄하므로 유한 표시 평탄이다. 따라서
$Y$를 $V$로 대체하여 (1), (2)에 더해 $Y_s = T$이고 $S$가 아핀이라고
가정해도 된다.

\medskip\noindent
$h$가 $y$에서 Cohen--Macaulay가 되는 닫힌점 $y \in Y_s$를 택하자.
Lemma \ref{more-morphisms-lemma-flat-finite-presentation-CM-open}을 보라. Lemma
\ref{more-morphisms-lemma-slice-CM}에 의해 도식
$$
\xymatrix{
Z \ar[r] \ar[rd] & Y \ar[d] \\
& S
}
$$
이 존재하는데, $Z \to S$는 평탄하고 국소 유한 표시이며 국소 준유한이고
$Z_s = \{y\}$이다. Lemma
\ref{more-morphisms-lemma-etale-makes-quasi-finite-finite-at-point}를 적용하여 기본 근방
$(S', s') \to (S, s)$와 열린 부분
$Z' \subset Z_{S'} = S' \times_S Z$를 찾을 수 있는데, $Z' \to S'$는
유한이고 $s$ 위에 놓이는 점 $z' \in Z'$가 유일하다. $Z' \to S'$는
또한 국소 유한 표시이고 평탄하므로($Z \to S$의 밑변환의 열린 부분이므로)
$Z' \to S'$는 유한 국소 자유이다. Morphisms, Lemma
\ref{morphisms-lemma-finite-flat}을 보라. $Y_{S'} \to S'$는 $h$의
밑변환으로서 유한 표시 평탄이고 올들이 기하학적으로 축약이다.
또한 $Y_{s'} = Y_s$는 기하학적으로 연결이다. Lemma
\ref{more-morphisms-lemma-descent-connected-fibres}를 $S'$ 위의
$Z' \to Y_{S'}$에 적용하면,
% P05-EN-CH37-0207
조건 (2)를 만족하고 $S'$ 위의 올들이 비어 있거나 기하학적으로 연결인
준콤팩트 열린 부분 $V \subset Y_{S'}$를 얻는다. $V \to S'$가 열린
사상이므로(Morphisms, Lemma \ref{morphisms-lemma-fppf-open}),
$S'$를 $s'$의 아핀 열린 근방으로 대체하여 $V \to S'$가 전사라고
가정해도 된다. 따라서
% P05-EN-CH37-0384
(1)이 성립한다.
\end{proof}

\begin{lemma}
\label{more-morphisms-lemma-cover-by-geometrically-connected}
$f : X \to S$를 국소 유한 표시이고 평탄하며 올들이 기하학적으로
축약인 스킴의 사상이라 하자. 그러면 $X_i \to S$가
$X_i \to S_i \to S$로 분해되고, $S_i \to S$가 \'etale이며,
$X_i \to S_i$가 유한 표시 평탄이고 올들이 기하학적으로 연결이며
기하학적으로 축약인 \'etale 덮개 $\{X_i \to X\}_{i \in I}$가
존재한다.
\end{lemma}

\begin{proof}
상이 $s \in S$인 점 $x \in X$를 택하자. 도식
$$
\xymatrix{
X' \ar[r] \ar[rd] & S' \times_S X \ar[r] \ar[d] & X \ar[d] \\
& S' \ar[r] & S
}
$$
과 점 $s' \in S'$, $x' \in X'$, $y \in S' \times_S X$를 만들겠다.
여기서 $x'$은 $x$로 가고, $(S', s') \to (S, s)$는 \'etale 근방이며,
$(X', x') \to (S' \times_S X, y)$는 \'etale
근방이고\footnote{실제로 증명은 열린 부분
$X' \subset S' \times_S X$를 준다.}, $X' \to S'$의 올들은
기하학적으로 연결이다. 모든 $x \in X$에 대해 이렇게 할 수 있으면
보조정리가 따른다(이렇게 만든 \'etale 사상 $X' \to X$들의 모임을
덮개의 원소들로 삼는다).
% P05-EN-CH37-0208
첫 단계로 $X$와 $S$를 각각 $x$와 $s$의 아핀 열린 근방으로 대체하면,
$X$와 $S$가 아핀인 경우로 환원된다(따라서 $f$는 유한 표시이다).

\medskip\noindent
% P05-EN-CH37-0209
$\kappa(s)$의 분리 대수적 확대 $\overline{k}$를 택하자.
$X_{\overline{k}}$로 $X_s$의 밑변환을 나타내자.
$x \in X_s$로 가는 $X_{\overline{k}}$의 점 $\overline{x}$를 택하자.
$\overline{x}$의 연결 준콤팩트 열린 근방
$\overline{V} \subset X_{\overline{k}}$를 택하자. (이는 체 위에서
국소 유한형인 임의의 스킴은 국소 Noether 위상공간으로서 국소 연결이기
때문에 가능하다.) Varieties, Lemma
\ref{varieties-lemma-Galois-action-quasi-compact-open}에 의해 유한 분리
확대 $k'/\kappa(s)$와, 밑변환이 $\overline{V}$인 준콤팩트 열린 부분
$V' \subset X_{k'}$를 찾을 수 있다. 특히 $V'$은 $k'$ 위에서
기하학적으로 연결이다. Varieties, Lemma
\ref{varieties-lemma-characterize-geometrically-connected}를 보라. Lemma
\ref{more-morphisms-lemma-realize-prescribed-residue-field-extension-etale}에 의해
$\kappa(s)$의 확대인 $k'$와 $\kappa(s')$가 동형이 되는 \'etale 근방
$(S', s') \to (S, s)$를 찾을 수 있다.
$x' \in (S' \times_S X)_{s'} = X_{k'}$로 $\overline{x}$의 상을
나타내자. 따라서 $(S, s)$를 $(S', s')$로, $(X, x)$를
$(S' \times_S X, x')$로 대체하면 다음 문단에서 다루는 경우로
환원된다.

\medskip\noindent
$x$를 포함하며 기하학적으로 연결인 준콤팩트 열린 부분
$V \subset X_s$가 존재한다고 가정하자. 그러면 Lemma
\ref{more-morphisms-lemma-fibre-geometrically-connected-reduced}를 적용하여,
$X' \to S'$의 올들이 기하학적으로 연결이고 $X'$이 $x$로 가는 점을
포함하도록 하는 아핀 \'etale 근방 $(S', s') \to (S, s)$와
준콤팩트 열린 부분 $X' \subset S' \times_S X$를 찾을 수 있다.
이로써 증명이 끝난다.
\end{proof}

\begin{lemma}
\label{more-morphisms-lemma-normal-morphism-irreducible}
$h : Y \to S$를 스킴의 사상이라 하고, $s \in S$를 점이라 하며,
$T \subset Y_s$를 열린 부분스킴이라 하자. 다음을 가정하자.
\begin{enumerate}
\item $h$는 유한 표시이다.
\item $h$는 정규 사상이다.
\item $T$는 $\kappa(s)$ 위에서 기하학적으로 기약이다.
\end{enumerate}
그러면 아핀 기본 \'etale 근방 $(S', s') \to (S, s)$와 다음을 만족하는
준콤팩트 열린 부분 $V \subset Y_{S'}$를 찾을 수 있다.
\begin{enumerate}
\item[(a)] $V \to S'$의 모든 올은 기하학적으로 정역이다.
\item[(b)] $V_{s'} = T \times_s s'$이다.
\end{enumerate}
\end{lemma}

\begin{proof}
Lemma \ref{more-morphisms-lemma-fibre-geometrically-connected-reduced}를 적용하여,
$V \to S'$의 모든 올이 기하학적으로 연결이고
$V_{s'} = T \times_s s'$가 되도록 하는 아핀 기본 \'etale 근방
$(S', s') \to (S, s)$와 준콤팩트 열린 부분
$V \subset Y_{S'}$를 찾자. $V$는 $h$의 밑변환의 열린 부분이므로
$V \to S'$의 모든 올은 기하학적으로 정규이다. Lemma
\ref{more-morphisms-lemma-normal}을 보라. 특히 이 올들은 기하학적으로 축약이다.
증명을 끝내려면 기하학적으로 기약임을 보이면 된다. 그런데
$t \in S'$이면 $V_t$는 $\kappa(t)$ 위에서 유한형이고, 따라서
$V_t \times_{\kappa(t)} \overline{\kappa(t)}$는
$\overline{\kappa(t)}$ 위에서 유한형이므로 Noether이다.
$S' \to S$의 선택에 의해 스킴
$V_t \times_{\kappa(t)} \overline{\kappa(t)}$는 연결이다. 따라서
Properties, Lemma \ref{properties-lemma-normal-Noetherian}에 의해
$V_t \times_{\kappa(t)} \overline{\kappa(t)}$는 기약이고, 증명이 끝난다.
\end{proof}








\section{유한형 사상의 구조에의 응용}
\label{more-morphisms-section-application-finite-type}

\noindent
이 절의 결과는 \cite{GruRay}에서 찾을 수 있다. 대략 말하면, 유한형
사상은 원천과 목표에서 \'etale 국소적으로, 유한 사상과 기하학적으로
연결인 올들을 갖는 매끄러운 사상의 합성이다. 이 매끄러운 사상의 상대
차원은 원래 사상의 올 차원과 같다.

\begin{lemma}
\label{more-morphisms-lemma-local-structure-finite-type}
$f : X \to S$를 사상이라 하자. $x \in X$를 택하고 $s = f(x)$라 두자.
$f$가 국소 유한형이고 $n = \dim_x(X_s)$라고 가정하자. 그러면 가환 도식
$$
\xymatrix{
X \ar[dd] & X' \ar[l]^g \ar[d]^\pi & x \ar@{|->}[dd] &
x' \ar@{|->}[l]  \ar@{|->}[d] \\
& Y \ar[d]^h & & y \ar@{|->}[d] \\
S \ar@{=}[r] & S & s & s \ar@{=}[l]
}
$$
과 $g(x') = x$인 점 $x' \in X'$가 존재하여, $y = \pi(x')$라 두면
다음이 성립한다.
\begin{enumerate}
\item $h : Y \to S$는 상대 차원 $n$인 매끄러운 사상이다.
\item $g : (X', x') \to (X, x)$는 기본 \'etale 근방이다.
\item $\pi$는 유한이고, $\pi^{-1}(\{y\}) = \{x'\}$이다.
\item $\kappa(y)$는 $\kappa(s)$의 순수 초월 확대이다.
\end{enumerate}
더욱이 $f$가 국소 유한 표시이면 $\pi$는 유한 표시이다.
\end{lemma}

\begin{proof}
문제는 $X$와 $S$ 위에서 국소적이므로 $X$와 $S$가 아핀이라고 가정해도
된다. Algebra, Lemma
\ref{algebra-lemma-refined-quasi-finite-over-polynomial-algebra}에 의해
$X$를 $X$ 안의 $x$의 표준 열린 근방으로 대체하여 분해
$$
\xymatrix{
X \ar[r]^\pi & \mathbf{A}^n_S \ar[r] & S
}
$$
가 존재하고, $\pi$가 준유한이며 $\kappa(\pi(x))$가 $\kappa(s)$
위에서 순수 초월이라고 가정해도 된다. Lemma
\ref{more-morphisms-lemma-etale-makes-quasi-finite-finite-at-point}에 의해 기본
\'etale 근방
$$
(Y, y) \to (\mathbf{A}^n_S, \pi(x))
$$
과 $y$ 위에 놓이는 유일한 점 $x'$을 포함하고 $X' \to Y$가 유한인
열린 부분 $X' \subset X \times_{\mathbf{A}^n_S} Y$가 존재한다.
% P05-EN-CH37-0210
이로써 (1)--(4)가 성립함이 증명된다. 마지막 명제에는 Morphisms,
Lemma \ref{morphisms-lemma-finite-presentation-permanence}을 적용한다.
\end{proof}

\begin{lemma}
\label{more-morphisms-lemma-local-local-structure-finite-type}
\begin{slogan}
유한형 사상은 \'etale 근방에서 매끄러운 사상 위의 유한 사상이다.
\end{slogan}
$f : X \to S$를 사상이라 하자. $x \in X$를 택하고 $s = f(x)$라 두자.
$f$가 국소 유한형이고 $n = \dim_x(X_s)$라고 가정하자. 그러면 가환 도식
$$
\xymatrix{
X \ar[dd] & X' \ar[l]^g \ar[d]^\pi & x \ar@{|->}[dd] &
x' \ar@{|->}[l] \ar@{|->}[d] \\
& Y' \ar[d]^h & & y' \ar@{|->}[d] \\
S & S' \ar[l]_e & s & s' \ar@{|->}[l]
}
$$
과 $g(x') = x$인 점 $x' \in X'$가 존재하여, $y' = \pi(x')$,
$s' = h(y')$라 두면 다음이 성립한다.
\begin{enumerate}
\item $h : Y' \to S'$는 상대 차원 $n$인 매끄러운 사상이다.
\item $Y' \to S'$의 모든 올은 기하학적으로 정역이다.
\item $g : (X', x') \to (X, x)$는 기본 \'etale 근방이다.
\item $\pi$는 유한이고, $\pi^{-1}(\{y'\}) = \{x'\}$이다.
\item $\kappa(y')$는 $\kappa(s')$의 순수 초월 확대이다.
\item $e : (S', s') \to (S, s)$는 기본 \'etale 근방이다.
\end{enumerate}
더욱이 $f$가 국소 유한 표시이면 $\pi$는 유한 표시이다.
\end{lemma}

\begin{proof}
문제는 $S$ 위에서 국소적이므로 $S$를 $s$의 아핀 열린 근방으로 대체해도
된다. 다음으로 Lemma \ref{more-morphisms-lemma-local-structure-finite-type}를 적용하여
가환 도식
$$
\xymatrix{
X \ar[dd] & X' \ar[l]^g \ar[d]^\pi & x \ar@{|->}[dd] &
x' \ar@{|->}[l] \ar@{|->}[d] \\
& Y \ar[d]^h & & y \ar@{|->}[d] \\
S \ar@{=}[r] & S & s & s \ar@{=}[l]
}
$$
을 얻는다. 여기서 $h$는 상대 차원 $n$인 매끄러운 사상이고,
$\kappa(y)$는 $\kappa(s)$의 순수 초월 확대이다. 문제는 $X$ 위에서도
국소적이므로 $Y$를 $y$의 아핀 근방으로 대체하고($X'$은 이를 $\pi$로
당긴 역상으로 대체한다), $S$가 아핀이므로 $Y \to S$가 준콤팩트이고
분리이며 매끄럽고, 특히 유한 표시임이 보장된다.
$T$를 $y$를 포함하는 $Y_s$의 연결 성분이라 하자. $Y_s$가
Noether이므로 $T$는 열려 있다. 또한 Varieties, Lemma
\ref{varieties-lemma-geometrically-connected-if-connected-and-point}에 의해
$T$는 $\kappa(s)$ 위에서 기하학적으로 연결이다. $T$는
$\kappa(s)$ 위에서 매끄럽기도 하므로 기하학적으로 정규이다. Varieties,
Lemma \ref{varieties-lemma-smooth-geometrically-normal}을 보라. 연결인
Noether 정규 스킴은 기약이므로(Properties, Lemma
\ref{properties-lemma-normal-Noetherian}), $T$는 $\kappa(s)$ 위에서
기하학적으로 기약이라는 결론을 얻는다. 마지막으로 Lemma
\ref{more-morphisms-lemma-smooth-normal}에 의해 매끄러운 사상 $h$는 정규임에 유의하자.
이제
% P05-EN-CH37-0211
사상 $h : Y \to S$와 열린 부분 $T \subset Y_s$에 대해 Lemma
\ref{more-morphisms-lemma-normal-morphism-irreducible}의 모든 가정이 성립함을
확인하였다. Lemma \ref{more-morphisms-lemma-normal-morphism-irreducible}를 적용하면
$e : S' \to S$, $s' \in S'$, $Y'$와 가환 도식
$$
\xymatrix{
X \ar[dd] & X' \ar[l]^g \ar[d]^\pi & X' \times_Y Y' \ar[l] \ar[d] &
x \ar@{|->}[dd] & x' \ar@{|->}[l] \ar@{|->}[d] &
(x', s') \ar@{|->}[l] \ar@{|->}[d] \\
& Y \ar[d]^h & Y' \ar[d] \ar[l] & & y \ar@{|->}[d] &
(y, s') \ar@{|->}[l] \ar@{|->}[d] \\
S \ar@{=}[r] & S & S' \ar[l]_e & s & s \ar@{=}[l] & s' \ar@{|->}[l]
}
$$
을 얻는다. 여기서 $e : (S', s') \to (S, s)$는 기본 \'etale 근방이고,
$Y' \subset Y_{S'}$는 $S'$ 위의 모든 올이 기하학적으로 기약이며
$Y_s = Y_{S', s'}$라는 식별 아래 $Y'_{s'} = T$를 만족하는 열린
근방이다. $(y, s') \in Y'$를 $y \in T$에 대응하는 점이라 하자.
이는 $Y \times_S S'$에서 $y$ 위에 놓이고 잉여체가 $\kappa(y)$와 같으며
$S'$의 $s'$로 가는 유일한 점이기도 하다. 마찬가지로
$(x', s') \in X' \times_Y Y' \subset X' \times_S S'$를 $x'$ 위에 놓이고
잉여체가 $\kappa(x')$와 같으며 $s'$ 위에 놓이는 유일한 점이라 하자.
그러면 이 도식의 바깥 부분은 보조정리에서 제기한 문제의 해이다.
몇 가지 사소한 세부사항은 생략한다.
\end{proof}

\begin{lemma}
\label{more-morphisms-lemma-local-local-structure-finite-type-affine}
% P05-EN-CH37-0212
가정과 표기는 Lemma \ref{more-morphisms-lemma-local-local-structure-finite-type}에서와
같다. 성질 (1)--(6)에 더해 다음도 성립하게 할 수 있다.
\begin{enumerate}
\item[(7)] $S'$, $Y'$, $X'$은 아핀이다.
\end{enumerate}
\end{lemma}

\begin{proof}
$Y'$이 아핀이면 $\pi$가 유한이므로 $X'$도 아핀임에 유의하자.
$s'$의 아핀 열린 근방 $U' \subset S'$를 택하고, $y'$의 아핀 열린 근방
$V' \subset h^{-1}(U')$를 택하자. $W' = h(V')$라 두자. 이는 $U'$에
포함된 $S'$ 안의 $s'$의 열린 근방이다. Morphisms, Lemma
\ref{morphisms-lemma-smooth-open}을 보라. $s'$의 아핀 열린 근방
$U'' \subset W'$를 택하자. 그러면 $h^{-1}(U'') \cap V'$는
$U'' \times_{U'} V'$와 같으므로 아핀이다. 구성에 의해
$h^{-1}(U'') \cap V' \to U''$는 전사 매끄러운 사상이고, 그 올들은
$Y' \to S'$의 기하학적으로 정역인 올들의 비어 있지 않은 열린
부분스킴이므로 기하학적으로 정역이다. 따라서 $S'$를 $U''$로,
$Y'$을 $h^{-1}(U'') \cap V'$로 대체해도 된다.
\end{proof}

\noindent
성질 $\pi^{-1}(\{y'\}) = \{x'\}$의 의미는 다음 보조정리로 일부
설명된다.

\begin{lemma}
\label{more-morphisms-lemma-finite-morphism-single-point-in-fibre}
$\pi : X \to Y$를 유한 사상이라 하자. $y = \pi(x)$이고
$\pi^{-1}(\{y\}) = \{x\}$인 $x \in X$를 택하자. 그러면
\begin{enumerate}
\item $X$ 안의 $x$의 임의의 근방 $U \subset X$에 대해,
$\pi^{-1}(V) \subset U$인 $y$의 근방 $V \subset Y$가 존재한다.
\item 환 사상 $\mathcal{O}_{Y, y} \to \mathcal{O}_{X, x}$는 유한이다.
\item $\pi$가 유한 표시이면
$\mathcal{O}_{Y, y} \to \mathcal{O}_{X, x}$는 유한 표시이다.
\item 임의의 준연접 $\mathcal{O}_X$-가군 $\mathcal{F}$에 대해
% P05-EN-CH37-0213
$\mathcal{F}_x = \pi_*\mathcal{F}_y$가 $\mathcal{O}_{Y, y}$-가군의
등식으로 성립한다.
\end{enumerate}
\end{lemma}

\begin{proof}
첫 번째 명제는 순전히 위상적이다. $\pi^{-1}(\{y\}) = \{x\}$를
만족하는 $\pi$가 연속 닫힌 사상임을 사용한다. 두 번째와 세 번째
명제를 증명하기 위해 $X = \Spec(B)$이고 $Y = \Spec(A)$라고 가정해도
된다. 그러면 $A \to B$는 유한 환 사상이고, $y$는 $A$의 소아이디얼
$\mathfrak p$에 대응하며, $\mathfrak p$ 위에 놓이는 $B$의 소아이디얼
$\mathfrak q$는 유일하다. 그러면
$B_{\mathfrak q} = B_{\mathfrak p}$이다.
Algebra, Lemma \ref{algebra-lemma-unique-prime-over-localize-below}를
보라. 다시 말해, 사상 $A_{\mathfrak p} \to B_{\mathfrak q}$는
$A \to B$를 $\mathfrak p$에서 국소화하여 얻는 사상
$A_{\mathfrak p} \to B_{\mathfrak p}$와 같다. 따라서 (2), (3)은
국소화의 간단한 성질에서 따른다(몇 가지 세부사항은 생략한다).
마지막 명제를 위해 어떤 $B$-가군 $M$에 대해
$\mathcal{F} = \widetilde M$이라고 하자. 그러면
% P05-EN-CH37-0214
$\mathcal{F} = M_{\mathfrak q}$이고
$\pi_*\mathcal{F}_y = M_{\mathfrak p}$이다. 위의 내용에 의해 이 두
국소화는 일치한다. 또는 (1)과 줄기의 정의를 사용하여
$\mathcal{F}_x = \pi_*\mathcal{F}_y$를 직접 확인할 수도 있다.
\end{proof}






\section{fppf 위상에의 응용}
\label{more-morphisms-section-application-fppf}

\noindent
위의 \'etale 국소화 기법을 사용하여 다음 결과를 증명할 수 있다.
이 결과는 fppf 위상이 Zariski 덮개들과, $f$가 전사 유한 국소 자유인
$\{f : T \to S\}$ 꼴의 덮개들에 의해 ``생성되는'' 위상과 같음을
서술한다.

\begin{lemma}
\label{more-morphisms-lemma-dominate-fppf-etale-locally}
$S$를 스킴이라 하고, $\{S_i \to S\}_{i \in I}$를 fppf 덮개라 하자.
그러면 다음이 존재한다.
\begin{enumerate}
\item \'etale 덮개 $\{S'_a \to S\}$,
\item 전사 유한 국소 자유 사상 $V_a \to S'_a$.
\end{enumerate}
이때 fppf 덮개 $\{V_a \to S\}$는 주어진 덮개
$\{S_i \to S\}$를 세분한다.
\end{lemma}

\begin{proof}
각 $S_i \to S$가 국소 준유한이라고 가정해도 된다. Lemma
\ref{more-morphisms-lemma-qf-fp-flat-dominates-fppf}를 보라.

\medskip\noindent
점 $s \in S$를 고정하자. $i \in I$와 $s$로 가는 점
$s_i \in S_i$를 택하자.
% P05-EN-CH37-0215
다음 열린 부분이 존재하는 기본 \'etale 근방
$(S', s) \to (S, s)$를 택하자.
$$
S_i \times_S S'  \supset V
$$
이 열린 부분은 $s \in S'$로 가는 유일한 점 $v \in V$를 포함하고
$V \to S'$는 유한이다. Lemma
\ref{more-morphisms-lemma-etale-makes-quasi-finite-finite-at-point}를 보라.
이 사상이 유한이고, $S_i \times_S S' \to S'$가
$S_i \to S$의 밑변환으로서 평탄하고 국소 유한 표시이므로,
$V \to S'$는 유한 국소 자유이다. Morphisms, Lemmas
\ref{morphisms-lemma-base-change-finite-presentation},
\ref{morphisms-lemma-base-change-flat}, 그리고
\ref{morphisms-lemma-finite-flat}을 보라. 따라서 $V \to S'$는 열린
사상이고, $S'$를 축소하여 $V \to S'$가 전사 유한 국소 자유라고
가정해도 된다. $S$의 모든 점에 대해 이렇게 할 수 있으므로,
$\{S_i \to S\}$는 각 $V_a \to S$가
$V_a \to S'_a \to S$로 분해되고, $S'_a \to S$가 \'etale이며
$V_a \to S'_a$가 전사 유한 국소 자유인 덮개
$\{V_a \to S\}_{a \in A}$로 세분될 수 있다.
\end{proof}

\begin{lemma}
\label{more-morphisms-lemma-dominate-fppf}
$S$를 스킴이라 하고, $\{S_i \to S\}_{i \in I}$를 fppf 덮개라 하자.
그러면 다음이 존재한다.
\begin{enumerate}
\item Zariski 열린 덮개 $S = \bigcup U_j$,
\item 전사 유한 국소 자유 사상 $W_j \to U_j$,
\item Zariski 열린 덮개 $W_j = \bigcup_k W_{j, k}$,
\item 전사 유한 국소 자유 사상 $T_{j, k} \to W_{j, k}$.
\end{enumerate}
이때 fppf 덮개 $\{T_{j, k} \to S\}$는 주어진 덮개
$\{S_i \to S\}$를 세분한다.
\end{lemma}

\begin{proof}
Lemma \ref{more-morphisms-lemma-dominate-fppf-etale-locally}에서 찾은 fppf 덮개를
$\{V_a \to S\}_{a \in A}$라 하자. 다시 말해, 이 덮개는
$\{S_i \to S\}$를 세분하고, 각 $V_a \to S$는
$V_a \to S'_a \to S$로 분해되는데 $S'_a \to S$는 \'etale이고
$V_a \to S'_a$는 전사 유한 국소 자유이다.

\medskip\noindent
Remark \ref{more-morphisms-remark-topologies}에 의해 Zariski 열린 덮개
$S = \bigcup U_j$, 각 $j$에 대한 유한 국소 자유 전사 사상
$W_j \to U_j$, 그리고 각 $j$에 대한 Zariski 열린 덮개
$\{W_{j, k} \to W_j\}$가 존재하여, 족
$\{W_{j, k} \to S\}$는 \'etale 덮개 $\{S'_a \to S\}$를 세분한다.
즉 각 쌍 $j, k$에 대해 $a(j, k)$와 사상 $W_{j, k} \to S$의 분해
% P05-EN-CH37-0216
$W_{j, k} \to S'_a \to S$가 존재한다.
$T_{j, k} = W_{j, k} \times_{S'_a} V_a$라 두면 모든 것이 명백하다.
\end{proof}

\begin{lemma}
\label{more-morphisms-lemma-extend-integral-surjective-morphisms}
$S$를 스킴이라 하자. $U \subset S$가 열려 있고 $V \to U$가 전사
정수적 사상이면, $\overline{V} \times_S U$가 $U$ 위의 스킴으로서
$V$와 동형이 되는 전사 정수적 사상 $\overline{V} \to S$가 존재한다.
\end{lemma}

\begin{proof}
% P05-EN-CH37-0217
$V' \to S$를 $U$에서 $S$의 정규화라 하자. Morphisms, Section
\ref{morphisms-section-normalization-X-in-Y}를 보라. 구성에 의해
$V' \to S$는 정수적이다. Morphisms, Lemmas
\ref{morphisms-lemma-normalization-localization}와
\ref{morphisms-lemma-normalization-in-integral}에 의해 $V'$에서 $U$의
역상은 $V$이다. $Z$를 $S \setminus U$ 위에 유도된 축약 스킴 구조라
하자. 그러면 $\overline{V} = V' \amalg Z$가 원하는 스킴이다.
\end{proof}

\begin{lemma}
\label{more-morphisms-lemma-extend-finite-surjective-morphisms}
$S$를 준콤팩트 준분리 스킴이라 하자. $U \subset S$가 준콤팩트 열린
부분이고 $V \to U$가 전사 유한 사상이면,
$\overline{V} \times_S U$가 $U$ 위의 스킴으로서 $V$와 동형이 되는
전사 유한 사상 $\overline{V} \to S$가 존재한다.
\end{lemma}

\begin{proof}
Zariski의 주정리(Lemma
\ref{more-morphisms-lemma-quasi-finite-separated-pass-through-finite})에 의해 $V$가
$S$ 위에서 유한인 스킴 $V'$의 준콤팩트 열린 부분이라고 가정해도 된다.
$V'$를 $V$의 스킴론적 상으로 대체하여 $V$가 $V'$에서 조밀하다고
가정해도 된다. $V$가 $U$ 위에서 유한이므로
$V \to V' \times_S U$는 닫힌 사상이고, 따라서
$V' \times_S U = V$이다. $Z$를 $S \setminus U$ 위에 유도된 축약
스킴 구조라 하자. 그러면 $\overline{V} = V' \amalg Z$가 원하는
스킴이다.
\end{proof}

\begin{lemma}
\label{more-morphisms-lemma-dominate-fppf-integral}
$S$를 스킴이라 하고, $\{S_i \to S\}_{i \in I}$를 fppf 덮개라 하자.
그러면 전사 정수적 사상 $S' \to S$와 열린 덮개
$S' = \bigcup U'_\alpha$가 존재하여, 각 $\alpha$에 대해 사상
$U'_\alpha \to S$는 어떤 $i$에 대한 $S_i \to S$를 통해 분해된다.
\end{lemma}

\begin{proof}
Lemma \ref{more-morphisms-lemma-dominate-fppf}에서와 같은 $S = \bigcup U_j$,
$W_j \to U_j$, $W_j = \bigcup W_{j, k}$, 그리고
$T_{j, k} \to W_{j, k}$를 택하자. Lemma
\ref{more-morphisms-lemma-extend-integral-surjective-morphisms}에 의해 $W_j \to U_j$를
전사 정수적 사상 $\overline{W}_j \to S$로 연장할 수 있다. 그런 다음
$T_{j, k} \to W_{j, k}$를 전사 정수적 사상
$\overline{T}_{j, k} \to \overline{W}_j$로 연장할 수 있다.
$\overline{T}_j$를 $\overline{W}_j$ 위의 모든 스킴
$\overline{T}_{j, k}$의 곱과 같게 두자(Limits, Lemma
\ref{limits-lemma-infinite-product}). 이어서 $S'$을 $S$ 위의 모든
스킴 $\overline{T}_j$의 곱과 같게 두자. $x \in S'$이면 $S$에서
$x$의 상이 $U_j$에 놓이는 $j$가 존재한다. 따라서 사영
$S' \to \overline{W}_j$에 의한 $x$의 상이 $W_{j, k}$에 놓이는
$k$가 존재한다. 그러므로 사영
$S' \to \overline{T}_j \to \overline{T}_{j, k}$에 의해 점 $x$는
$T_{j, k}$로 가고, $T_{j, k} \to S$는 어떤 $i$에 대한 $S_i$를
통해 분해된다. 마지막으로 Limits, Lemmas
\ref{limits-lemma-infinite-product-integral}와
\ref{limits-lemma-infinite-product-surjective}에 의해 사상
$S' \to S$는 정수적이고 전사이다.
\end{proof}

\begin{lemma}
\label{more-morphisms-lemma-dominate-fppf-finite}
$S$를 준콤팩트 준분리 스킴이라 하고,
$\{S_i \to S\}_{i \in I}$를 fppf 덮개라 하자. 그러면 유한 표시인
전사 유한 사상 $S' \to S$와 열린 덮개 $S' = \bigcup U'_\alpha$가
존재하여, 각 $\alpha$에 대해 사상 $U'_\alpha \to S$는 어떤 $i$에
대한 $S_i \to S$를 통해 분해된다.
\end{lemma}

\begin{proof}
% P05-EN-CH37-0218
$Y \to X$를 Lemma \ref{more-morphisms-lemma-dominate-fppf-integral}에서 찾은 정수적
전사 사상이라 하자. $V_j \to X$가 $S_{i(j)}$를 통해 분해되는 유한
아핀 열린 덮개 $Y = \bigcup V_j$를 택하자. $Y_\lambda \to X$가
유한이고 유한 표시이도록 $Y = \lim Y_\lambda$라고 쓸 수 있다.
Limits, Lemma
\ref{limits-lemma-integral-limit-finite-and-finite-presentation}를 보라.
충분히 큰 $\lambda$에 대해 $Y$로의 역상이 $V_j$가 되는 아핀 열린
부분 $V_{\lambda, j} \subset Y_\lambda$를 찾을 수 있다. Limits,
Lemma \ref{limits-lemma-descend-opens}를 보라. $\lambda$를 더 크게 하면
$X$ 위의 사상 $V_j \to S_{i(j)}$는 $X$ 위의 사상
$V_{\lambda, j} \to S_{i(j)}$에서 온다. Limits, Proposition
\ref{limits-proposition-characterize-locally-finite-presentation}을 보라.
이 $\lambda$에 대해 $S' = Y_\lambda$로 두면 증명이 끝난다.
\end{proof}

\begin{lemma}
\label{more-morphisms-lemma-fppf-ph}
스킴의 fppf 덮개는 ph 덮개이다.
\end{lemma}

\begin{proof}
$\{T_i \to T\}$를 스킴의 fppf 덮개라 하자. Topologies, Definition
\ref{topologies-definition-fppf-covering}을 보라. $T_i \to T$는 국소
유한형임에 유의하자. $U \subset T$를 아핀 열린 부분이라 하자.
$\{T_i \times_T U \to U\}$가 표준 ph 덮개로 세분될 수 있음을 보이면
충분하다. Topologies, Definition
\ref{topologies-definition-ph-covering}을 보라. 이는 Lemma
\ref{more-morphisms-lemma-dominate-fppf-finite}와 유한 사상이 고유하다는 사실에서
곧바로 따른다(Morphisms, Lemma
\ref{morphisms-lemma-finite-proper}).
\end{proof}

\begin{remark}
\label{more-morphisms-remark-change-topologies}
Lemma \ref{more-morphisms-lemma-fppf-ph}의 결과로 비교 사상
$$
\epsilon : (\Sch/S)_{ph} \longrightarrow (\Sch/S)_{fppf}
$$
을 얻는다. 이는 밑 범주들 사이의 항등 함자가 주는 위치의 사상이다
(Topologies, Remark \ref{topologies-remark-choice-sites}에서와 같이
위치들을 적절히 택한다). 함자 $\epsilon_*$는 밑 전층들 위에서
항등이고,
% P05-EN-CH37-0219
함자 $\epsilon^{-1}$는 fppf 층에 그 ph 층화를 대응시킨다. 합성에
의해 ph 위상을 syntomic, 매끄러운, \'etale, 그리고 Zariski 위상과도
비교할 수 있다.
\end{remark}




\section{준사영 스킴}
\label{more-morphisms-section-quasi-projective}

\noindent
``준사영 스킴''이라는 용어는 아직 정의하지 않았다. 한 가지 가능한
정의는 풍부한 가역층을 갖는 스킴일 것이다. 그러나 $X$가 밑스킴 $S$
위의 스킴일 때, 사상 $X \to S$가 준사영적이면 {\it $X$는 $S$ 위에서
준사영적이다}라고 한다(Morphisms, Definition
\ref{morphisms-definition-quasi-projective}). 임의의 스킴의 항등사상은
준사영적이므로, $S$ 위에서 준사영적인 스킴이 반드시 풍부한 가역층을
갖는 것은 아니다. 이런 이유로 ``준사영 스킴''이라는 용어는 정의하지
않는 편이 나아 보인다.

\begin{lemma}
\label{more-morphisms-lemma-quasi-projective}
$S$를 풍부한 가역층을 갖는 스킴이라 하고, $f : X \to S$를 스킴의
사상이라 하자. 다음은 동치이다.
\begin{enumerate}
\item $X \to S$는 준사영적이다.
\item $X \to S$는 H-준사영적이다.
\item $X' \to S$가 사영적인, $S$ 위 스킴의 준콤팩트 열린 몰입
$X \to X'$가 존재한다.
\item $X \to S$는 유한형이고 $X$는 풍부한 가역층을 갖는다.
\item $X \to S$는 유한형이고 $f$-매우 풍부한 가역층이 존재한다.
\end{enumerate}
\end{lemma}

\begin{proof}
(2) $\Rightarrow$ (1)은 Morphisms, Lemma
\ref{morphisms-lemma-H-quasi-projective-quasi-projective}이다.
(1) $\Rightarrow$ (2)는 Morphisms, Lemma
\ref{morphisms-lemma-projective-over-quasi-projective-is-H-projective}이다.
(2) $\Rightarrow$ (3)은 Morphisms, Lemma
\ref{morphisms-lemma-H-quasi-projective-open-H-projective}이다.

\medskip\noindent
(3)에서와 같은 $X \subset X'$을 가정하자. 특히 $X \to S$는 유한형이다.
% P05-EN-CH37-0220
Morphisms, Lemma
\ref{morphisms-lemma-H-quasi-projective-open-H-projective}에 의해 사상
$X \to S$는 H-사영적이다. 따라서 준콤팩트 몰입
$i : X \to \mathbf{P}^n_S$가 존재한다. 그러므로
$\mathcal{L} = i^*\mathcal{O}_{\mathbf{P}^n_S}(1)$은
$f$-매우 풍부하다. $X \to S$가 준콤팩트이므로 Morphisms, Lemma
\ref{morphisms-lemma-ample-very-ample}에 의해 $\mathcal{L}$은
$f$-풍부하다. 따라서 정의에 의해 $X \to S$는 준사영적이다.

\medskip\noindent
(4) $\Rightarrow$ (2)는 Morphisms, Lemma
\ref{morphisms-lemma-quasi-projective-finite-type-over-S}이다.

\medskip\noindent
(5) $\Rightarrow$ (4)는 Morphisms, Lemma
\ref{morphisms-lemma-pullback-ample-tensor-relatively-ample}에서 따른다.

\medskip\noindent
(1)과 (5)의 동치는 Morphisms, Lemmas
\ref{morphisms-lemma-ample-very-ample}와
\ref{morphisms-lemma-finite-type-ample-very-ample}에서 따른다.
\end{proof}

\begin{lemma}
\label{more-morphisms-lemma-category-quasi-projective}
$S$를 풍부한 가역층을 갖는 스킴이라 하자. $\text{QP}_S$를
Lemma \ref{more-morphisms-lemma-quasi-projective}의 동치 조건들을 만족하는 $S$ 위
스킴들의 범주의 충만 부분범주라 하자.
\begin{enumerate}
\item $S' \to S$가 스킴의 사상이고 $S'$이 풍부한 가역층을 가지면,
밑변환은 함자 $\text{QP}_S \to \text{QP}_{S'}$를 정한다.
\item $X \in \text{QP}_S$이고 $Y \in \text{QP}_X$이면
$Y \in \text{QP}_S$이다.
\item 범주 $\text{QP}_S$는 올곱에 대해 닫혀 있다.
\item 범주 $\text{QP}_S$는 유한 분리합집합에 대해 닫혀 있다.
\item $X \to S$가 사영적이면 $X \in \text{QP}_S$이다.
\item $X \to S$가 유한형 준아핀이면 $X$는 $\text{QP}_S$에 속한다.
\item $X \to S$가 준유한이고 분리이면 $X \in \text{QP}_S$이다.
\item $X \to S$가 준콤팩트 몰입이면 $X \in \text{QP}_S$이다.
% P05-EN-CH37-0221
\item 여기에 더 추가한다.
\end{enumerate}
\end{lemma}

\begin{proof}
(1)은 Morphisms, Lemma
\ref{morphisms-lemma-base-change-quasi-projective}에서 따른다.

\medskip\noindent
(2)는 Lemma \ref{more-morphisms-lemma-quasi-projective}의 네 번째 특징짓기에서 따른다.

\medskip\noindent
$X \to S$와 $Y \to S$가 준사영적이면 Morphisms, Lemma
\ref{morphisms-lemma-base-change-quasi-projective}에 의해
$X \times_S Y \to Y$는 준사영적이다. 따라서 (3)은 (2)에서 따른다.

\medskip\noindent
$X = Y \amalg Z$가 스킴들의 분리합집합이고 $\mathcal{L}$이
$\mathcal{L}|_Y$와 $\mathcal{L}|_Z$가 풍부한 가역
$\mathcal{O}_X$-가군이면, $\mathcal{L}$은 풍부하다(세부사항은
생략한다). 따라서 (4)는 Lemma \ref{more-morphisms-lemma-quasi-projective}의 네 번째
특징짓기에서 따른다.

\medskip\noindent
(5)는 Morphisms, Lemma
\ref{morphisms-lemma-projective-quasi-projective}에서 따른다.

\medskip\noindent
(6)은 Morphisms, Lemma
\ref{morphisms-lemma-quasi-affine-finite-type-quasi-projective}에서 따른다.

\medskip\noindent
(7)은 (6)과 Lemma
\ref{more-morphisms-lemma-quasi-finite-separated-quasi-affine}에서 따른다.

\medskip\noindent
(8)은 (7)과 Morphisms, Lemma
\ref{morphisms-lemma-immersion-locally-quasi-finite}에서 따른다.
\end{proof}

\noindent
다음 보조정리는 사실 이 절에 속하지 않지만, 다른 어디에도 적절히 둘
곳이 없어 보인다.

\begin{lemma}
\label{more-morphisms-lemma-integral-over-quasi-affine}
$X$를 준아핀 스킴이라 하고, $f : U \to X$를 정수적 사상이라 하자.
그러면 $U$는 준아핀이고 다음 도식은
$$
\xymatrix{
U \ar[r] \ar[d] & \Spec(\Gamma(U, \mathcal{O}_U)) \ar[d] \\
X \ar[r] & \Spec(\Gamma(X, \mathcal{O}_X))
}
$$
Cartesian이다.
\end{lemma}

\begin{proof}
$U$가 준아핀인 것은 정수적 사상은 아핀이고, 아핀 사상은 준아핀이며,
스킴이 준아핀일 필요충분조건은 $\Spec(\mathbf{Z})$로 가는 구조사상이
준아핀인 것이고, 준아핀 사상들의 합성은 준아핀이기 때문이다. 앞의 두
명제는 정의에서 바로 따르고 세 번째 명제는 Morphisms, Lemma
\ref{morphisms-lemma-composition-quasi-affine}이다. 이제
$U' =
X \times_{\Spec(\Gamma(X, \mathcal{O}_X))} \Spec(\Gamma(U, \mathcal{O}_U))$
로 놓고 확장된 도식
$$
\xymatrix{
U \ar[r]_j \ar[rd] & U' \ar[d] \ar[r] &
\Spec(\Gamma(U, \mathcal{O}_U)) \ar[d] \\
& X \ar[r] & \Spec(\Gamma(X, \mathcal{O}_X))
}
$$
을 생각하자. 정수적 사상이 보편적으로 닫혔다는 사실
(Morphisms, Lemma \ref{morphisms-lemma-integral-universally-closed})과
% P05-EN-CH37-0222
도식의 수직 화살표들이 분리 사상이라는 사실을 Morphisms, Lemma
\ref{morphisms-lemma-image-proper-scheme-closed}와 함께 적용하면 사상
$j$는 닫힌 사상이다. 한편 Properties, Lemma
\ref{properties-lemma-quasi-affine}에 의해 보조정리의 도식에 있는 수평
화살표들이 열린 사상이므로 $j$는 열린 사상이다. 따라서 $j$는 $U$를
$U'$의 열리고 닫힌 부분스킴과 동일시한다. $U \not = U'$이면 $U$는
$U'$에서 조밀하지 않고, 따라서 더욱이 $\Gamma(U, \mathcal{O}_U)$의
스펙트럼에서도 조밀하지 않다. 그러나
$\Spec(\Gamma(U, \mathcal{O}_U))$에서 $U$의 스킴론적 상은
$\Spec(\Gamma(U, \mathcal{O}_U))$이다. 실제로 $U$가 그것을 통해
분해되는 닫힌 부분스킴을 잘라 내는 $\Gamma(U, \mathcal{O}_U)$의 임의의
아이디얼은 영이어야 한다. 그러므로 예를 들어 Morphisms, Lemma
\ref{morphisms-lemma-quasi-compact-scheme-theoretic-image}에 의해 $U$는
$\Spec(\Gamma(U, \mathcal{O}_U))$에서 조밀하다. 따라서 $U = U'$이고
증명이 끝난다.
\end{proof}









\section{사영 스킴}
\label{more-morphisms-section-projective}

\noindent
이 절은 사영 사상에 관한 Section \ref{more-morphisms-section-quasi-projective}의
대응물이다.

\begin{lemma}
\label{more-morphisms-lemma-projective}
$S$를 풍부한 가역층을 갖는 스킴이라 하고, $f : X \to S$를 스킴의
사상이라 하자. 다음은 동치이다.
\begin{enumerate}
\item $X \to S$는 사영적이다.
\item $X \to S$는 H-사영적이다.
\item $X \to S$는 준사영적이고 고유하다.
\item $X \to S$는 H-준사영적이고 고유하다.
\item $X \to S$는 고유하고 $X$는 풍부한 가역층을 갖는다.
\item $X \to S$는 고유하고 $f$-풍부한 가역층이 존재한다.
\item $X \to S$는 고유하고 $f$-매우 풍부한 가역층이 존재한다.
\item $\mathcal{A}_1$이 유한형 $\mathcal{O}_S$-가군이고
$\mathcal{A}_0$ 위에서 $\mathcal{A}_1$에 의해 생성되며 다음을 만족하는
준연접 등급 $\mathcal{O}_S$-대수 $\mathcal{A}$가 존재한다.
$X = \underline{\text{Proj}}_S(\mathcal{A})$. 
\end{enumerate}
\end{lemma}

\begin{proof}
먼저 각 경우에 사상 $f$가 고유함에 유의하자. Morphisms, Lemmas
\ref{morphisms-lemma-H-projective}와
\ref{morphisms-lemma-locally-projective-proper}을 보라. 따라서 $f$가
고유 사상인 경우에 이 개념들의 동치만 증명하면 충분하다. 이하에서는
이를 따로 언급하지 않고 사용한다.

\medskip\noindent
(1) $\Leftrightarrow$ (3)과 (2) $\Leftrightarrow$ (4)의 동치는
Morphisms, Lemma
\ref{morphisms-lemma-projective-is-quasi-projective-proper}이다.

\medskip\noindent
(2) $\Rightarrow$ (1)은 Morphisms, Lemma
\ref{morphisms-lemma-H-projective}이다.

\medskip\noindent
(1) $\Rightarrow$ (2)와 (3) $\Rightarrow$ (4)는 Morphisms, Lemma
\ref{morphisms-lemma-projective-over-quasi-projective-is-H-projective}이다.

\medskip\noindent
(1) $\Rightarrow$ (7)은 Morphisms, Definitions
\ref{morphisms-definition-projective}와
\ref{morphisms-definition-very-ample}에서 바로 따른다.

\medskip\noindent
(3)과 (6)은 Morphisms, Definition
\ref{morphisms-definition-quasi-projective}에 의해 동치이다.

\medskip\noindent
따라서 (1)--(4)와 (6)은 동치이고 (7)을 함의한다. Lemma
\ref{more-morphisms-lemma-quasi-projective}에 의해 (3), (5), (7)은 동치이다.
그러므로 (1)--(7)이 동치이다.

\medskip\noindent
Divisors, Lemma \ref{divisors-lemma-relative-proj-projective}에 의해 (8)은
(1)을 함의한다. 반대로 (2)가 성립하면 닫힌 몰입
$$
i :
X
\longrightarrow
\mathbf{P}^n_S = \underline{\text{Proj}}_S(\mathcal{O}_S[T_0, \ldots, T_n]).
$$
을 택할 수 있다. 등식은 Constructions, Lemma
\ref{constructions-lemma-projective-space-bundle}을 보라. Divisors,
Lemma \ref{divisors-lemma-closed-subscheme-proj}에 의해 $X$는
$\mathcal{O}_S[T_0, \ldots, T_n]$의 준연접 등급 몫대수 $\mathcal{A}$의
상대 Proj이다. 이때 $\mathcal{A}$는 (8)의 조건들을 만족한다.
\end{proof}

\begin{lemma}
\label{more-morphisms-lemma-category-projective}
$S$를 풍부한 가역층을 갖는 스킴이라 하자. $\text{P}_S$를 Lemma
\ref{more-morphisms-lemma-projective}의 동치 조건들을 만족하는 $S$ 위 스킴들의 범주의
충만 부분범주라 하자.
\begin{enumerate}
\item $S' \to S$가 스킴의 사상이고 $S'$이 풍부한 가역층을 가지면,
밑변환은 함자 $\text{P}_S \to \text{P}_{S'}$를 정한다.
\item $X \in \text{P}_S$이고 $Y \in \text{P}_X$이면
$Y \in \text{P}_S$이다.
\item 범주 $\text{P}_S$는 올곱에 대해 닫혀 있다.
\item 범주 $\text{P}_S$는 유한 분리합집합에 대해 닫혀 있다.
\item $X \to S$가 유한이면 $X$는 $\text{P}_S$에 속한다.
% P05-EN-CH37-0223
\item 여기에 더 추가한다.
\end{enumerate}
\end{lemma}

\begin{proof}
(1)은 Morphisms, Lemma
\ref{morphisms-lemma-base-change-projective}에서 따른다.

\medskip\noindent
(2)는 Lemma \ref{more-morphisms-lemma-projective}의 다섯 번째 특징짓기와 고유 사상들의
합성이 고유하다는 사실(Morphisms, Lemma
\ref{morphisms-lemma-composition-proper})에서 따른다.

\medskip\noindent
$X \to S$와 $Y \to S$가 사영적이면 Morphisms, Lemma
\ref{morphisms-lemma-base-change-projective}에 의해
$X \times_S Y \to Y$는 사영적이다. 따라서 (3)은 (2)에서 따른다.

\medskip\noindent
$X = Y \amalg Z$가 스킴들의 분리합집합이고 $\mathcal{L}$이
$\mathcal{L}|_Y$와 $\mathcal{L}|_Z$가 풍부한 가역
$\mathcal{O}_X$-가군이면, $\mathcal{L}$은 풍부하다(세부사항은 생략한다).
따라서 (4)는 Lemma \ref{more-morphisms-lemma-projective}의 다섯 번째 특징짓기에서
따른다.

\medskip\noindent
(5)는 Morphisms, Lemma
\ref{morphisms-lemma-finite-projective}에서 따른다.
\end{proof}

\noindent
다음은 조금 다른 유형의 결과이다.

\begin{lemma}
\label{more-morphisms-lemma-ample-in-neighbourhood}
\begin{reference}
\cite[IV Corollary 9.6.4]{EGA}
\end{reference}
$f : X \to Y$를 스킴의 고유 사상이라 하고, $\mathcal{L}$을 가역
$\mathcal{O}_X$-가군이라 하자. $y \in Y$를 $\mathcal{L}_y$가 $X_y$
위에서 풍부한 점이라 하자. 그러면 $\mathcal{L}|_{f^{-1}(V)}$가
$f^{-1}(V)/V$ 위에서 풍부한 $y$의 열린 근방 $V \subset Y$가 존재한다.
\end{lemma}

\begin{proof}
$Y$가 아핀이라고 가정해도 된다. 그러면 $Y_i$는 $\mathbf{Z}$ 위에서
유한형이고, 전이 사상 $X_i \to X_{i'}$와 $Y_i \to Y_{i'}$는
아핀이며, $X_i \to Y_i$는 고유하고,
$X \to Y = \lim (X_i \to Y_i)$가 되는 유향 집합 $I$와 스킴 사상
$X_i \to Y_i$의 역계를 얻는다. Limits, Lemma
\ref{limits-lemma-proper-limit-of-proper-finite-presentation-noetherian}.
$I$를 축소한 뒤, $\mathcal{L}$로 당겨지는 가역
$\mathcal{O}_{X_i}$-가군 $\mathcal{L}_i$의 호환되는 계가 있다고 가정할
수 있다. Limits, Lemma \ref{limits-lemma-descend-invertible-modules}를
보라. $y_i \in Y_i$를 $y$의 상이라 하자. 그러면
$\kappa(y) = \colim \kappa(y_i)$이다. 따라서 어떤 $i$에 대하여 Limits,
Lemma \ref{limits-lemma-limit-ample}에 의해 $\mathcal{L}_{i, y_i}$는
$X_{i, y_i}$ 위에서 풍부하다. Cohomology of Schemes, Lemma
\ref{coherent-lemma-ample-in-neighbourhood}에 의해 $\mathcal{L}_i$의
$f_i^{-1}(V_i)$로의 제한이 $V_i$에 대해 상대적으로 풍부한 $y_i$의
열린 근방 $V_i \subset Y_i$를 얻는다. $V \subset Y$를 $V_i$의
역상으로 놓으면 증명이 끝난다(힌트: Morphisms, Lemma
\ref{morphisms-lemma-ample-base-change}, $X \to Y \times_{Y_i} X_i$가
아핀이라는 사실, 그리고 Morphisms, Lemma
\ref{morphisms-lemma-pullback-ample-tensor-relatively-ample}에 의해
풍부한 가역층을 아핀 사상으로 당기면 풍부하다는 사실을 사용하라).
\end{proof}







\section{Proj와 Spec}
\label{more-morphisms-section-proj-spec}

\noindent
이 절에서는 등급환의 Proj와 스펙트럼 사이의 관계를 명확히 한다.

\medskip\noindent
$R$을 환이라 하고 $A$를 등급 $R$-대수라 하자. Algebra, Section
\ref{algebra-section-graded}를 보라. $m \geq 0$에 대하여
$A_{\geq m} = \bigoplus_{d \geq m} A_d$로 쓴다. 등급환
$$
B = \bigoplus\nolimits_{d \geq 0} A_{\geq d}
$$
을 생각하자. $d' \geq d$이고 $a \in A_{d'}$일 때, $a$에 대응하는
$B_d$의 원소를 $a^{(d)} \in B$로 쓰자. 자명한 두 환 준동형을
$\sigma : A \to B$와 $\psi : A \to B$로 쓰자. 즉 $a \in A_d$이면
$\sigma(a) = a^{(0)}$이고 $\psi(a) = a^{(d)}$이다. 그러면 $\psi$는
등급환 준동형이고 $\sigma$는 $B$를 $A$ 위의 등급대수로 만든다. 또한
$d' \geq d$이고 $a \in A_{d'}$일 때, $d' > d$이면 $a^{(d)}$를 $0$으로
보내고 $d' = d$이면 $a$로 보내는 전사 등급환 준동형
$\tau : B \to A$가 있다.

\medskip\noindent
아핀 스킴과 스펙트럼. $X = \Spec(A)$로 놓자. 무관 아이디얼 $A_+$는
닫힌 부분스킴 $Z = V(A_+) = \Spec(A/A_+) = \Spec(A_0)$을 잘라 낸다.
$U = X \setminus Z$로 놓자.
$$
U \longrightarrow X \longrightarrow Z
$$
사영 스킴과 Proj. $P = \text{Proj}(A)$로 놓자. $P$를
$\Spec(A_0) = Z$ 위의 스킴으로 보아도 되고, 그렇게 하기로 한다.
$L = \text{Proj}(B)$로 놓자. $L$을
$\Spec(B_0) = \Spec(A) = X$ 위의 스킴으로 보아도 되고, 그렇게 하기로
한다. $B_0$와 $A$의 동일시는 $\sigma$로 주어짐에 유의하자. 전사
$\tau$는 닫힌 몰입 $0 : P \to L$을 정한다.
$A \xrightarrow{\sigma} B \to A$가 사상 $A \to A_0 \to A$와 같으므로
다음 도식은
$$
\xymatrix{
P \ar[d] \ar[r]_0 & L \ar[d] \\
Z \ar[r] & X
}
$$
가환이다.

\medskip\noindent
$\psi$가 사상 $L \to P$를 정의한다고 주장한다. 이를 보이기 위해서는
Constructions, Lemma \ref{constructions-lemma-morphism-proj}에 의해
$B_+ \not \subset \mathfrak p$인 모든 동차 소아이디얼
$\mathfrak p \subset B$에 대하여
$\psi(A_+) \not \subset \mathfrak p$임을 확인하면 충분하다. 먼저
$g \in B_+$를 $g \not \in \mathfrak p$인 동차원소로 택하자. 그러면
어떤 $d_i \geq d$와 $a_i \in A_{d_i}$에 대하여 $g$를 유한합
$g = \sum a_i^{(d)}$로 쓸 수 있다. 따라서 $a^{(d)} \not \in \mathfrak p$인
$d' \geq d$와 $a \in A_{d'}$가 존재한다. 이제
$$
(a^{(d)})^{d'} =
(a^{d'})^{(d'd)} =
a^{(d)} (a^{d' - 1})^{(d(d' - 1))} =
\psi(a) (a^{d' - 1})^{(d(d' - 1))}
$$
는(표기가 썩 만족스럽지는 않지만) $\mathfrak p$에 속하지 않는다.
따라서 $\psi(a) \not \in \mathfrak p$이고 주장이 증명된다. 그러므로
위의 도식을 가환도식
$$
\xymatrix{
P \ar[d] \ar[r]_0 & L \ar[d] \ar[r]_\pi & P \ar[d] \\
Z \ar[r] & X \ar[r] & Z
}
$$
으로 확장할 수 있다. 여기서 $X \to Z$는 $A_0 \to A$로 주어진다.
$\tau \circ \psi = \text{id}_A$이므로
$\pi \circ 0 = \text{id}_P$이다.

\medskip\noindent
$\pi$가 아핀 사상임에 유의하자. 이는 Constructions, Lemma
\ref{constructions-lemma-morphism-proj}의 구성에서 명백하다. 실제로 어떤
$d > 0$에 대하여 $f \in A_d$이면, $g = \psi(f)$로 놓을 때
$\pi^{-1}(D_+(f)) = D_+(g)$이다. 이 경우 동차 성분들의 등식
$$
(B[1/g])_{m'} = \bigoplus\nolimits_{m \geq m'} (A[1/f])_m
$$
을 얻는다. 이 동형은 추가 국소화와 호환된다. $m' = 0$으로 놓으면
$\pi_*\mathcal{O}_L$은 $m \geq 0$인 $\mathcal{O}_P(m)$들의
직합이다\footnote{마찬가지로
$\pi_*\mathcal{O}_L(i) = \bigoplus_{m \geq -i} \mathcal{O}_P(m)$임을
얻는다.}. 따라서 $L$은 상대 스펙트럼
$$
L = \underline{\Spec}_P
\left(
\bigoplus\nolimits_{m \geq 0} \mathcal{O}_P(m)
\right)
$$
과 동일시된다. 특히 $L \to P$는 원뿔이다\footnote{$L$은 흔히 $P$
위의 선다발이다. 아래를 보라.}. Constructions, Section
\ref{constructions-section-cone}을 보라. 또한 $0 : P \to L$이 원뿔의
꼭짓점임은 명백하다.

\medskip\noindent
앞 단락에서처럼 어떤 $d > 0$에 대하여 $f \in A_d$이고
$g = \psi(f) \in B_d$라 하자. 다음 환 준동형들의 구조를 살펴보자.
$$
\xymatrix{
A_0 \ar[r] \ar[d] & A \ar[d]^\sigma \ar[r] & A_0 \ar[d] \\
(A[1/f])_0 \ar[r]^-\psi &
(B[1/g])_0 = \bigoplus\nolimits_{m \geq 0} (A[1/f])_m
\ar[r]^-\tau &
(A[1/f])_0
}
$$
% P05-EN-CH37-0224
등급환에서 계산하면\footnote{(1)과 (2)는 명백하다. (3)을 보이려면
$a \in A_d$일 때
$\sigma(a) = \sigma(f) \psi(a/f)$임에 유의하라. (4)에 대해서는
$b/g^m$이 $\tau$의 핵에 속할 필요충분조건이
$b \in A_{\geq md}$의 $A_{md}$에서의 상이 영인 것임에 유의하라. 따라서
$m' > md$이고 $a \in A_{m'}$이면 $a^{(md)}/g^m$의 어떤 거듭제곱이
$\sigma(f)$가 생성하는 아이디얼에 속함을 보이면 충분하다.
$em' - emd \geq d$인 $e$를 택하자. 그러면
$$
(a^{(md)}/g^m)^e = (a^e)^{(emd)}/g^{em} =
(fa^e)^{(emd + d)}/g^{em + 1} = \sigma(f) \cdot (a^e)^{(emd + d)}/g^{em + 1}
$$
이고 이것이 원하는 바이다(끔찍한 표기를 양해하기 바란다). (5)는
앞에서처럼 논증하고, $d = 1$이면
$a^{(md)}/g^m = \sigma(f) \cdot a^{(md + 1)}/g^{m + 1}$임에 유의하라.}
다음을 얻는다.
\begin{enumerate}
\item $\sigma(A_+)(B[1/g])_0 \subset \Ker(\tau : (B[1/g])_0 \to (A[1/f])_0)$,
\item $\sigma(f) \in (B[1/g])_0$는 영인자가 아닌 원소이다.
\item 아이디얼로서
$\sigma(f) (B[1/g])_0 = \sigma(A_d) (B[1/g])_0$이다.
\item $\sigma(f) (B[1/g])_0$ and
$\Ker(\tau : (B[1/g])_0 \to (A[1/f])_0)$은 같은 근기를 갖는다.
\item $d = 1$이면
$\sigma(f) (B[1/g])_0 = \Ker(\tau : (B[1/g])_0 \to (A[1/f])_0)$.
\end{enumerate}
특히 집합론적으로
$$
0(D_+(f)) = V(\sigma(f)) \subset D_+(g) = \Spec((B[1/g])_0)
$$
임을 알 수 있다. 다시 말해 $\sigma(A_d)$가 생성하는 아이디얼은
$D_+(g)$ 위에서 유효 Cartier 인자를 잘라 내며, 이는 집합론적으로 닫힌
몰입 $0 : P \to L$의 상과 같다.

\medskip\noindent
$L \to X$가 $U$ 위에서 동형이라고 주장한다. 실제로 어떤 $d > 0$에
대하여 $f \in A_d$이면
$$
\Spec(A_f) \times_X L = \text{Proj}(A_f \otimes_A B) =
\text{Proj}(B_{\sigma(f)})
$$
이다. 각 $e$에 대하여
% P05-EN-CH37-0225
$(B_{\sigma(f)})_e = A_f \otimes_B B_e = A_f \otimes_A A_{\geq e} = A_f$,
이며 마지막 등식은 단사 $A_{\geq e} \subset A$가 유도한다. 따라서
$B_{\sigma(f)} \cong A_f[T]$이고 $T$의 차수는 $1$이다.
$\text{Proj}(A_f[T]) \to \Spec(A_f)$가 동형이므로 주장이 증명된다.
이제부터 $U$를 $L$의 대응하는 열린집합과 동일시한다.

\medskip\noindent
앞 단락에서 한 동일시에 의해 제한 $\pi|_U : U \to P$를 생각할 수 있다.
위에서 여러 번 했듯이 어떤 $d > 0$에 대하여 $f \in A_d$를 택하고
$g = \psi(f) \in B_d$라 하자. 그러면
$$
U \cap \pi^{-1}(D_+(f)) = U \cap D_+(g)
$$
은 $D_+(g)$를 $(B[1/g])_0$의 스펙트럼과 동일시했을 때
$\sigma(f) \in (B[1/g])_0$의 영점자취의 여집합이다. 이것은 위의 주장
(4)이다. 따라서 $U \cap D_+(g)$는 아핀이고
$$
\mathcal{O}_L(U \cap D_+(g)) = (B[1/g])_0[1/\sigma(f)] =
\bigoplus\nolimits_{m \in \mathbf{Z}} (A[1/f])_m
$$
이다. 여기서 마지막 등호는 위에서 한 동일시
$(B[1/g])_0 = \bigoplus_{m \geq 0} (A[1/f])_m$의 자연스러운
확장이다. 앞에서 $\pi : L \to P$에 대해 했던 것과 똑같이,
$\pi|_U : U \to P$가 아핀이고
$$
U = \underline{\Spec}_P
\left(
\bigoplus\nolimits_{m \in \mathbf{Z}} \mathcal{O}_P(m)
\right)
$$
임을 얻는다. 여기서 등식은 $P$ 위의 스킴으로서의 등식이다.

\medskip\noindent
위 내용을 요약하면, 우리의 구성은 스킴의 가환도식
\begin{equation}
\label{more-morphisms-equation-proj-and-spec}
\vcenter{
\xymatrix{
\underline{\Spec}_P \left(
\bigoplus\nolimits_{m \in \mathbf{Z}} \mathcal{O}_P(m)
\right)
\ar[r] \ar@{=}[d] &
L =
\underline{\Spec}_P \left(
\bigoplus\nolimits_{m \geq 0} \mathcal{O}_P(m)
\right) \ar[d]^\sigma \ar[r]_-\pi & P \ar[d] \\
U \ar[r] & X \ar[r] & Z
}
}
\end{equation}
을 준다. 여기서 $\pi$는 원뿔이고, 그 영단면 $0 : P \to L$의 상은
집합론적으로 $L$에서 $Z$의 역상이다.

\medskip\noindent
$W \subset P$를 $\mathcal{O}_P(1)|_W$가 가역이고 자연스러운 사상들이
모든 $m \in \mathbf{Z}$에 대하여 동형
$\mathcal{O}_P(m)|_W \cong \mathcal{O}_P(1)^{\otimes m}|_W$
을 유도하는 가장 큰 열린집합이라 하자. 즉 $d = 1$에 대한
Constructions, Lemma \ref{constructions-lemma-where-invertible}의
열린집합이다. 그러면 $L|_W = \pi^{-1}(W) \to W$는 계수 $1$인 벡터
다발이다(Constructions, Section
\ref{constructions-section-vector-bundle}). 구체적으로 Grothendieck식
표기로
$$
L|_W = \mathbf{V}(\mathcal{O}_P(1)|_W)
$$
이다. 이는 위에서 $L|_W$가 $\mathcal{O}_P(1)|_W$ 위의
$\mathcal{O}_W$-대칭대수의 상대 스펙트럼과 같음을 보였으므로 바로
따른다. 이때 사상 $0|_W : W \to L|_W$는 이 벡터 다발의 영단면임이
명백하다. 특히 $0(W)$는 $L|_W$ 위의 유효 Cartier 인자이다. 또한
열린집합 $U|_W = (\pi|_U)^{-1}(W)$는 영단면의 여집합이다.

\medskip\noindent
$A$가 $A_0$ 위에서 $f_1, \ldots, f_r \in A_1$에 의해 생성되면
모든 $m \geq 0$에 대하여
$(f_1, \ldots, f_r)^m = A_{\geq m}$이고, 따라서 위의 $B$는
$A_+ = (f_1, \ldots, f_r)$에 관한 리스 대수이다. 그러므로 이 경우
$L \to X$는 $Z$의 블로업이고, 앞 단락의 $W$에 대하여 $W = P$이다.

\medskip\noindent
$P$가 준콤팩트이면 충분히 나누어지는 $d$에 대하여
$\sigma(A_d)\mathcal{O}_L$이 잘라 내는 닫힌 부분스킴 $D \subset L$은
유효 Cartier 인자이고, $0 : P \to L$은 $D$를 통해 분해되며, 집합론적으로
$0(P) = D$이다. 이는 Constructions, Lemma
\ref{constructions-lemma-proj-quasi-compact}와 위에서 증명한 (1), (2),
(3), (4)에서 따른다. (보조정리에서 얻은 원소들의 차수의 최소공배수로
나누어지는 임의의 $d$를 택하라.)

\medskip\noindent
$P$가 준콤팩트이라는 가정을 계속 유지하자. $\mathcal{F}$를 준연접
$\mathcal{O}_P$-가군이라 하고
$\mathcal{F}_U = \pi^*\mathcal{F}|_U$로 놓자. 그러면
\begin{equation}
\label{more-morphisms-equation-cohomology-torsor}
R\Gamma(U, \mathcal{F}_U) =
\bigoplus\nolimits_{m \in \mathbf{Z}}
R\Gamma(P, \mathcal{F} \otimes_{\mathcal{O}_P} \mathcal{O}_P(m))
\end{equation}
이다. 또한 이 직합 분해는 $\mathcal{F}$에 대해 함자적이고, 오른쪽에
유도된 $A$-가군 구조는 $U \subset X$에서 오는 왼쪽의 $A$-가군 구조와
같다. 이 공식을 증명하자. $\pi|_U$가 아핀이고
$(\pi|_U)_*\mathcal{O}_U = \bigoplus_{m \in \mathbf{Z}} \mathcal{O}_P(m)$
이므로
\begin{align*}
R(\pi|_U)_*\mathcal{F}_U
& =
(\pi|_U)_*\mathcal{F}_U \\
& =
(\pi|_U)_*(\pi|_U)^*\mathcal{F} \\
& =
\mathcal{F} \otimes_{\mathcal{O}_P} 
\bigoplus\nolimits_{m \in \mathbf{Z}} \mathcal{O}_P(m) \\
& =
\bigoplus\nolimits_{m \in \mathbf{Z}}
\mathcal{F} \otimes_{\mathcal{O}_P} \mathcal{O}_P(m)
\end{align*}
을 얻는다. Leray에 의해
$R\Gamma(U, \mathcal{F}_U) =
R\Gamma(P, R(\pi|_U)_*\mathcal{F}_U)$, see
임을 얻는다. Cohomology, Lemma \ref{cohomology-lemma-apply-Leray}를 보라.
이 경우 코호몰로지를 취하는 것은 직합과 가환하므로 증명이 끝난다.
Derived Categories of Schemes, Lemma
\ref{perfect-lemma-quasi-coherence-pushforward-direct-sums}를 보라.
여기서 $P$가 준콤팩트이라는 것을 사용했다. Constructions, Lemma
\ref{constructions-lemma-proj-separated}에 의해 $P$는 분리이다.

\begin{lemma}
\label{more-morphisms-lemma-apply-proj-spec}
$R$을 환이라 하자. $P$를 $R$ 위의 고유 스킴이라 하고 $\mathcal{L}$을
풍부한 가역 $\mathcal{O}_P$-가군이라 하자.
$A = \bigoplus_{m \geq 0} \Gamma(P, \mathcal{L}^{\otimes m})$로 놓자.
그러면 $P = \text{Proj}(A)$이고 도식 (\ref{more-morphisms-equation-proj-and-spec})은
다음 도식이 된다.
$$
\xymatrix{
\underline{\Spec}_P \left(
\bigoplus\nolimits_{m \in \mathbf{Z}} \mathcal{L}^{\otimes m}
\right)
\ar[r] \ar@{=}[d] &
L =
\underline{\Spec}_P \left(
\bigoplus\nolimits_{m \geq 0} \mathcal{L}^{\otimes m}
\right) \ar[d]^\sigma \ar[r]_-\pi & P \ar[d] \\
U \ar[r] & X \ar[r] & Z
}
$$
이 도식은 위에서 설명한 성질들을 갖는다.
\end{lemma}

\begin{proof}
Morphisms, Lemma \ref{morphisms-lemma-proper-ample-is-proj}에 의해
$P = \text{Proj}(A)$이다. 또한 Properties, Lemma
\ref{properties-lemma-ample-gcd-is-one}에 의해 이 동일시 아래에서 모든
$m \in \mathbf{Z}$에 대하여
$\mathcal{O}_P(m) = \mathcal{L}^{\otimes m}$이다.
\end{proof}









\section{올의 닫힌점}
\label{more-morphisms-section-closed-points-fibres}

\noindent
이 절의 내용 일부는 프리프린트 \cite{Osserman-Payne}에서 가져왔다.

\begin{lemma}
\label{more-morphisms-lemma-locally-principal-vertical}
$f : X \to S$를 스킴의 사상이라 하고, $Z \subset X$를 닫힌
부분스킴이라 하며, $s \in S$라 하자. 다음을 가정하자.
\begin{enumerate}
\item $S$는 기약이고 일반점은 $\eta$이다.
\item $X$는 기약이다.
\item $f$는 우세하다.
\item $f$는 국소 유한형이다.
\item $\dim(X_s) \leq \dim(X_\eta)$,
\item $Z$는 $X$에서 국소 주이다.
\item $Z_\eta = \emptyset$.
\end{enumerate}
그러면 올 $Z_s$는 집합론적으로 $X_s$의 기약성분들의 합집합이다.
\end{lemma}

\begin{proof}
$X_{red}$를 $X$의 축약이라 하자. 그러면 $Z \cap X_{red}$는
$X_{red}$의 국소 주 닫힌 부분스킴이다. Divisors, Lemma
\ref{divisors-lemma-pullback-locally-principal}를 보라. 따라서 $X$가
축약이라고 가정해도 된다. 다시 말해 $X$는 정역 스킴이다. Properties,
Lemma \ref{properties-lemma-characterize-integral}을 보라. 이 경우 사상
$X \to S$는 $S_{red}$를 통해 분해된다. Schemes, Lemma
\ref{schemes-lemma-map-into-reduction}을 보라. 그러므로 $S$를
$S_{red}$로 바꾸고 $S$도 정역이라고 가정해도 된다.

\medskip\noindent
$f$가 우세하다는 주장은 $X$의 일반점이 $\eta$로 간다는 뜻이다.
Morphisms, Lemma
\ref{morphisms-lemma-dominant-finite-number-irreducible-components}를 보라.
또한 스킴 $X_\eta$는 체 $\kappa(\eta)$ 위에서 국소 유한형인 정역
스킴이다. 따라서 $X_\eta$의 모든 점 $\xi$에 대하여
$d = \dim(X_\eta) \geq 0$은 $\dim_\xi(X_\eta)$와 같다. Algebra, Lemmas
\ref{algebra-lemma-dimension-spell-it-out}와
\ref{algebra-lemma-dimension-at-a-point-finite-type-over-field}를 보라.
Morphisms, Lemma \ref{morphisms-lemma-openness-bounded-dimension-fibres}와
조건 (5)에 의해 모든 $x \in X_s$에 대하여
$\dim_x(X_s) = d$임을 얻는다.

\medskip\noindent
뇌터인 경우에는 다음과 같이 증명할 수 있다.
% P05-EN-CH37-0226
보조정리가 거짓이면 $Z_s$의 한 기약성분의 일반점이지만 $X_s$의 어느
기약성분의 일반점도 아닌 $x \in Z_s$가 존재한다. 이때
$\dim_x(X_s) = d$이고 $X_s$에서 $x$의 한 근방에서는 닫힌 부분스킴
$Z_s$가 $X_s$의 어느 기약성분도 포함하지 않으므로
$\dim_x(Z_s) \leq d - 1$이다. 따라서 $X$를 $x$의 열린 근방으로 바꾼
뒤 모든 $z \in Z$에 대하여
$\dim_z(Z_{f(z)}) \leq d - 1$이라고 가정해도 된다. Morphisms, Lemma
\ref{morphisms-lemma-openness-bounded-dimension-fibres}를 보라.
$\xi' \in Z$를 $Z$의 한 기약성분의 일반점이라 하고,
% P05-EN-CH37-0227
$s' = f(\xi)$로 놓자. $Z \not = X$가 국소 주이므로
% P05-EN-CH37-0228
$\dim(\mathcal{O}_{X, \xi}) = 1$이다. Algebra, Lemma
\ref{algebra-lemma-minimal-over-1}을 보라(여기서 $X$가 뇌터임을
사용한다). $\xi \in X$를 $X$의 일반점이라 하고, $\xi_1$을
$\xi'$을 포함하는 $X_{s'}$의 임의의 기약성분의 일반점이라 하자.
그러면 특수화
$$
\xi \leadsto \xi_1 \leadsto \xi'.
$$
를 얻는다. $\dim(\mathcal{O}_{X, \xi}) = 1$이므로 두 특수화 중 하나는
등식이어야 한다. 가정에 의해 $s' \not = \eta$이므로 첫 번째 특수화는
등식이 아니다. 따라서 $\xi' = \xi_1$은 $X_{s'}$의 한 기약성분의
일반점이다. Morphisms, Lemma
\ref{morphisms-lemma-openness-bounded-dimension-fibres}를 한 번 더 적용하면
$\dim_{\xi'}(Z_{s'}) = \dim_{\xi'}(X_{s'}) \geq \dim(X_\eta) = d$
를 얻으며, 이는 원하는 모순이다.

\medskip\noindent
일반적인 경우에는 다음과 같이 뇌터인 경우로 환원한다. 보조정리가
거짓이면 점 $x \in X$가 존재하여, $x$는 $s$ 위에 놓이고 $Z_s$의 한
기약성분의 일반점이지만 $X_s$의 어느 기약성분의 일반점도 아니다.
$U \subset S$를
$s$의 아핀 근방이라 하고, $V \subset X$를 $f(V) \subset U$인 $x$의
아핀 근방이라 하자. $U = \Spec(A)$와 $V = \Spec(B)$로 써서 $f|_V$가
환 준동형 $A \to B$로 주어지게 하자. $\mathfrak q \subset B$와
$\mathfrak p \subset A$를 각각 $x$와 $s$에 대응하는 소아이디얼이라
하자. 필요하면 $V$를 축소하여 $Z \cap V$가 어떤 원소 $g \in B$로
잘려 나온다고 가정해도 된다. $K$를 $A$의 분수체라 하자. 이 시점에서
알고 있는 것은 다음과 같다.
\begin{enumerate}
\item $A \subset B$는 정역들의 유한 생성 확대이다.
\item 원소 $g \otimes 1$은 $B \otimes_A K$에서 가역이다.
\item $d = \dim(B \otimes_A K) = \dim(B \otimes_A \kappa(\mathfrak p))$,
\item $g \otimes 1$은 $B \otimes_A \kappa(\mathfrak p)$의 단위원이 아니다.
\item $g \otimes 1$은 $B \otimes_A \kappa(\mathfrak p)$의 어느
극소 소아이디얼에도 속하지 않는다.
\end{enumerate}
모순을 이끌어 내고자 한다.

\medskip\noindent
$A$ 위에서 $B$를 생성하는 원소 $x_1, \ldots, x_n \in B$를 택하자.
유한 생성 $\mathbf{Z}$-대수 $A_0 \subset A$에 대하여 $B_0 \subset B$를
$x_1, \ldots, x_n$이 생성하는 $A_0$-부분대수라 하고, $K_0$를 $A_0$의
분수체라 하며, $\mathfrak p_0 = A_0 \cap \mathfrak p$로 놓자. $A_0$가
충분히 크면 계 $(A_0 \subset B_0, g, \mathfrak p_0)$에 대해서도
(1)--(5)가 성립한다고 주장한다.

\medskip\noindent
각 조건을 차례로 증명하자. (1)은 구성에 의해 성립한다. (2)를 위해
어떤 $h \otimes 1/a \in B \otimes_A K$에 대하여
% P05-EN-CH37-0229
$(g \otimes 1) h = 1$로 쓰자. 어떤 계수 $a_I, a'_I \in A$에 대하여
$g = \sum a_I x^I$, $h = \sum a'_I x^I$로 쓰자(다중지표 표기).
$A_0$가 $a$와 모든 $a_I, a'_I$를 포함하면 (2)가 성립한다. 실제로 단사
준동형 $B_0 \to B$의 국소화로서
$B_0 \otimes_{A_0} K_0 \subset B \otimes_A K$이다. (3)을 얻기 위해
완전열
$$
0 \to I \to A[X_1, \ldots, X_n] \to B \to 0
$$
을 생각하자. 여기서 두 번째 사상은 $X_i$를 $x_i$로 보내며 이 완전열은
$I$를 정의한다. $\otimes$가 오른쪽 완전하므로 $I \otimes_A K$와
$I \otimes_A \kappa(\mathfrak p)$는 각각 전사
$K[X_1, \ldots, X_n] \to B \otimes_A K$와
$\kappa(\mathfrak p)[X_1, \ldots, X_n] \to
B \otimes_A \kappa(\mathfrak p)$의 핵이다. 체 위의 다항식환은
뇌터이므로 $I \otimes_A K$와 $I \otimes_A \kappa(\mathfrak p)$를
생성하는 유한개의 원소 $h_j \in I$, $j = 1, \ldots, m$이 존재한다.
$h_j = \sum a_{j, I}X^I$로 쓰자. $A_0$가 모든 $a_{j, I}$를 포함하면
$$
B_0 \otimes_{A_0} K_0 \otimes_{K_0} K = B \otimes_A K
\quad\text{and}\quad
B_0 \otimes_{A_0} \kappa(\mathfrak p_0)
\otimes_{\kappa(\mathfrak p_0)} \kappa(\mathfrak p)
=
B \otimes_A \kappa(\mathfrak p).
$$
Morphisms, Lemma \ref{morphisms-lemma-dimension-fibre-after-base-change} 또는
Algebra, Lemma \ref{algebra-lemma-dimension-preserved-field-extension}에
의해 (3)의 차원 등식들이 성립한다. (4)는 자명하다.
$B_0 \otimes_{A_0} \kappa(\mathfrak p_0) \subset
B \otimes_A \kappa(\mathfrak p)$이므로 Algebra, Lemma
\ref{algebra-lemma-image-dense-generic-points}에 의해
$B_0 \otimes_{A_0} \kappa(\mathfrak p_0)$의 각 극소 소아이디얼은
$B \otimes_A \kappa(\mathfrak p)$의 어떤 극소 소아이디얼 아래에 놓인다.
따라서 (5)가 성립한다. 이와 같이 문제를 위에서 다룬 뇌터인 경우로
환원한다.
\end{proof}

\noindent
다음은 위 보조정리의 대수적 응용이다. $A \to B$가 평탄하면 보조정리의
네 번째 가정이 성립한다. Lemma \ref{more-morphisms-lemma-equality-dimensions}를 보라.

\begin{lemma}
\label{more-morphisms-lemma-horizontal}
$A \to B$를 국소환의 국소 준동형이라 하고
$g \in \mathfrak m_B$라 하자. 다음을 가정하자.
\begin{enumerate}
\item $A$와 $B$는 정역이고 $A \subset B$이다.
\item $B$는 $A$ 위에서 본질적으로 유한형이다.
\item $g$는 $\mathfrak m_AB$ 위의 어느 극소 소아이디얼에도 속하지 않는다.
\item $\dim(B/\mathfrak m_AB) +
\text{trdeg}_{\kappa(\mathfrak m_A)}(\kappa(\mathfrak m_B)) =
\text{trdeg}_A(B)$.
\end{enumerate}
그러면 $A \subset B/gB$이다. 즉 $\Spec(A)$의 일반점은 사상
$\Spec(B/gB) \to \Spec(A)$의 상에 속한다.
\end{lemma}

\begin{proof}
두 주장은 Algebra, Lemma
\ref{algebra-lemma-image-dense-generic-points}에 의해 동치임에 유의하자.
증명을 시작하기 위해 $C$를 유한형 $A$-대수라 하고 $\mathfrak q$를
$B = C_{\mathfrak q}$인 $C$의 소아이디얼이라 하자. 물론 $C$가 정역이고
$g \in C$라고 가정해도 된다. $C$를 한 국소화로 바꾸면
$\dim(C/\mathfrak m_AC) = \dim(B/\mathfrak m_AB) +
\text{trdeg}_{\kappa(\mathfrak m_A)}(\kappa(\mathfrak m_B))$임을 알 수 있다.
Morphisms, Lemma \ref{morphisms-lemma-dimension-fibre-at-a-point}를 보라.
$K$를 $A$의 분수체로 놓으면 같은 보조정리에 의해
$\dim(C \otimes_A K) = \text{trdeg}_A(B)$이다. 따라서 가정 (4)는 사상
$\Spec(C) \to \Spec(A)$의 일반올과 닫힌올의 차원이 같다는 뜻이다.

\medskip\noindent
보조정리가 거짓이라고 가정하자. 그러면
$(B/gB) \otimes_A K = 0$이다. 이는 $g \otimes 1$이
$B \otimes_A K = C_{\mathfrak q} \otimes_A K$에서 가역이라는 뜻이다.
% P05-EN-CH37-0230
$C_{\mathfrak q}$는 주 국소화들의 극한이므로 어떤
$h \in C$, $h \not \in \mathfrak q$에 대하여 $g \otimes 1$이
$C_h \otimes_A K$에서 가역임을 얻는다. 따라서 $C$를 $C_h$로 바꾸고
$(C/gC) \otimes_A K = 0$이라고 가정해도 된다. $C$를 한 번 더 바꾸어
$C/\mathfrak m_AC$의 극소 소아이디얼들과 $B/\mathfrak m_AB$의 극소
소아이디얼들이 일대일로 대응하게 하자. 이제
$X = \Spec(C) \to \Spec(A) = S$와
% P05-EN-CH37-0231
국소 닫힌 부분스킴 $Z = \Spec(C/gC)$에 Lemma
\ref{more-morphisms-lemma-locally-principal-vertical}를 적용한다. $Z_K = \emptyset$이므로
$Z \otimes \kappa(\mathfrak m_A)$는
$X \otimes \kappa(\mathfrak m_A) = \Spec(C/\mathfrak m_AC)$의 한
기약성분을 포함해야 한다. 그러나 이는 $g$가 $\mathfrak m_AB$ 위의
어느 극소 소아이디얼에도 속하지 않는다는 가정에 모순이다. 따라서
보조정리가 성립한다.
\end{proof}

\begin{lemma}
\label{more-morphisms-lemma-equality-dimensions}
$A \to B$를 국소환의 국소 준동형이라 하자. 다음을 가정하자.
\begin{enumerate}
\item $A$와 $B$는 정역이고 $A \subset B$이다.
\item $B$는 $A$ 위에서 본질적으로 유한형이다.
\item $B$는 $A$ 위에서 평탄하다.
\end{enumerate}
그러면
$$
\dim(B/\mathfrak m_AB) +
\text{trdeg}_{\kappa(\mathfrak m_A)}(\kappa(\mathfrak m_B)) =
\text{trdeg}_A(B).
$$
이다.
\end{lemma}

\begin{proof}
$C$를 유한형 $A$-대수라 하고 $\mathfrak q$를
$B = C_{\mathfrak q}$인 $C$의 소아이디얼이라 하자. $C$가 정역이라고
가정해도 된다. 다음 등식이 성립한다.
$\dim_{\mathfrak q}(C/\mathfrak m_AC) = \dim(B/\mathfrak m_AB) +
\text{trdeg}_{\kappa(\mathfrak m_A)}(\kappa(\mathfrak m_B))$.
Morphisms, Lemma \ref{morphisms-lemma-dimension-fibre-at-a-point}를 보라.
$K$를 $A$의 분수체로 놓으면 같은 보조정리에 의해
$\dim(C \otimes_A K) = \text{trdeg}_A(B)$.
따라서 실제로 증명할 것은
$\dim_{\mathfrak q}(C/\mathfrak m_AC) = \dim(C \otimes_A K)$.
이다. $A$를 우세하는 $K$의 값매김환 $A'$을 택하자. Algebra, Lemma
\ref{algebra-lemma-dominate}를 보라. $C' = C \otimes_A A'$로 놓자.
$\mathfrak q$ 위에 놓이는 $C'$의 소아이디얼 $\mathfrak q'$을 택하자.
그러한 소아이디얼이 존재하는 것은
$$
C'/\mathfrak m_{A'}C' =
C/\mathfrak m_AC \otimes_{\kappa(\mathfrak m_A)} \kappa(\mathfrak m_{A'})
$$
이고, 이 등식은 $C/\mathfrak m_AC \to C'/\mathfrak m_{A'}C'$가 충실히
평탄함을 보이기 때문이다. 또한 이는
$\dim_{\mathfrak q}(C/\mathfrak m_AC) =
\dim_{\mathfrak q'}(C'/\mathfrak m_{A'}C')$를 보인다. Algebra, Lemma
\ref{algebra-lemma-dimension-at-a-point-preserved-field-extension}를 보라.
$B' = C'_{\mathfrak q'}$은 $B \otimes_A A'$의 국소화임에 유의하자.
따라서 $B'$은 $A'$ 위에서 평탄하다. 일반올
$B' \otimes_{A'} K$는 $B \otimes_A K$의 국소화이다. 따라서 $B'$은
정역이다. $A' \subset B'$에 대해 보조정리를 증명하면 등식
$\dim_{\mathfrak q'}(C'/\mathfrak m_{A'}C') = \dim(C' \otimes_{A'} K)$
를 얻고, 위에서 말한 바에 의해 이는 원하는 등식
$\dim_{\mathfrak q}(C/\mathfrak m_AC) = \dim(C \otimes_A K)$
를 함의한다. 이로써 보조정리는 $A$가 값매김환인 경우로 환원된다.

\medskip\noindent
$A \subset B$가 보조정리에서와 같고 $A$가 값매김환이라고 하자.
앞에서처럼 $A$ 위의 어떤 유한형 정역 $C$에 대하여
$B = C_{\mathfrak q}$로 쓰자. Algebra, Lemma
\ref{algebra-lemma-finite-type-domain-over-valuation-ring-dim-fibres}에 의해
$\dim(C/\mathfrak m_AC) = \dim(C \otimes_A K)$를 얻고 증명이 끝난다.
\end{proof}

\begin{lemma}
\label{more-morphisms-lemma-closed-point-nearby-fibre}
$f : X \to S$를 스킴의 사상이라 하고
$x \leadsto x'$을 $X$의 점들의 특수화라 하자.
$s = f(x)$와 $s' = f(x')$로 놓자. 다음을 가정하자.
\begin{enumerate}
\item $x'$은 $X_{s'}$의 닫힌점이다.
\item $f$는 국소 유한형이다.
\end{enumerate}
그러면 집합
$$
\{x_1 \in X
\text{ 중에서 }
f(x_1) = s
\text{이고 }
x_1\text{은 }X_s\text{에서 닫혀 있으며 }
\text{또한 }
x \leadsto x_1 \leadsto x'
\}
$$
은 $X_s$에서 $x$의 폐포 안에서 조밀하다.
\end{lemma}

\begin{proof}
특수화 $x \leadsto x'$에 Schemes, Lemma
\ref{schemes-lemma-points-specialize}를 적용한다. 그러면 값매김환 $B$와
사상 $\varphi : \Spec(B) \to X$를 얻는데, $\varphi$는 일반점을 $x$로
보내고 닫힌점을 $x'$으로 보낸다. 또한 $\kappa(x)$가 $B$의 분수체라고
가정해도 된다. $A = B \cap \kappa(s)$로 놓자. 이는 값매김환이고
(Algebra, Lemma \ref{algebra-lemma-valuation-ring-cap-field}를 보라)
$\mathcal{O}_{S, s'} \to \kappa(s)$의 상을 지배함에 유의하자. 가환도식
$$
\xymatrix{
\Spec(B) \ar[rd] \ar[r] &
X_A \ar[d] \ar[r] & X \ar[d] \\
& \Spec(A) \ar[r] & S
}
$$
을 생각하자. $B$의 일반점(각각 닫힌점)은 $\Spec(A)$의 일반점(각각
닫힌점) 위에 놓이는 $X_A$의 점 $x_A$(각각 $x'_A$)로 간다. Morphisms,
Lemma
\ref{morphisms-lemma-base-change-closed-point-fibre-locally-finite-type}에
의해 $x'_A$는 $X_A$의 특수올의 닫힌점임에 유의하자. 또한
$X_A \to \Spec(A)$의 일반올은 $X_s$와 동형이다. 따라서 보조정리는
$S$가 값매김환의 스펙트럼이고 $s = \eta \in S$가 일반점이며
$s' \in S$가 닫힌점인 경우로 환원된다.

\medskip\noindent
$\dim_x(X_\eta)$에 관한 귀납법으로 보조정리를 증명한다.
$\dim_x(X_\eta) = 0$이면 $x$로 특수화되는 $X_\eta$의 다른 점은 없고
$x$는 그 올에서 닫혀 있다. Morphisms, Lemma
\ref{morphisms-lemma-quasi-finite-at-point-characterize}를 보라. 따라서
결과가 성립한다. 이제 $\dim_x(X_\eta) > 0$이라고 가정하자.

\medskip\noindent
$X' \subset X$를 $X$의 기약 닫힌 부분스킴 $\overline{\{x\}}$ 위의 유도된
축약 스킴 구조라 하자. Schemes, Definition
\ref{schemes-definition-reduced-induced-scheme}를 보라. $X$를 $X'$으로
바꾸면 $\dim_x(X_\eta)$는 감소하기만 하므로 보조정리를 증명할 때 그렇게
해도 된다. 따라서 $X$가 정역 스킴이고 $x$가 그 일반점이라고 가정해도
된다. 또한 $X$를 $x'$의 아핀 근방으로 바꾸어도 된다. 그러므로
$X = \Spec(B)$이고 $A \subset B$는 정역들의 유한형 확대이다. 이 경우
$\dim_x(X_\eta) = \dim(X_\eta) = \dim(X_{s'})$이고, 실제로 $X_{s'}$는
등차원적이다. Algebra, Lemma
\ref{algebra-lemma-finite-type-domain-over-valuation-ring-dim-fibres}를 보라.

\medskip\noindent
$W \subset X_\eta$를 진닫힌 부분집합이라 하자(이것이 우리가
``피하고자'' 하는 부분집합이다). $X_s$는 체 위에서 유한형이므로 $W$는
유한개의 기약성분을 가지며
$W = W_1 \cup \ldots \cup W_n$으로 쓸 수 있다.
% P05-EN-CH37-0232
$\mathfrak q_j \subset B$, $j = 1, \ldots, r$
를 대응하는 소아이디얼들이라 하자. $\mathfrak q \subset B$를 점 $x'$에
대응하는 극대아이디얼이라 하자.
$\mathfrak p_1, \ldots, \mathfrak p_s \subset B$를
$\mathfrak m_AB$ 위에 놓이는 극소 소아이디얼들이라 하자. 이들은 뇌터
스킴 $X_{s'}$의 기약성분들에 대응하므로 유한개이다. 또한 이
기약성분들은 각각 차원이 $> 0$이므로(위를 보라) 모든 $i$에 대하여
$\mathfrak p_i \not = \mathfrak q$이다. 이제 모든 $j$에 대하여
$g \not \in \mathfrak q_j$이고 모든 $i$에 대하여
$g \not \in \mathfrak p_i$인 원소 $g \in \mathfrak q$를 택하자. Algebra,
Lemma \ref{algebra-lemma-silly}를 보라. $Z \subset X$를 $g$가 정의하는
국소 주 닫힌 부분스킴이라 하자.
$Z_\eta = Z_{1, \eta} \cup \ldots \cup Z_{n, \eta}$, $n \geq 0$을
$Z$의 일반올의 기약성분 분해라 하자(일반올이 뇌터이므로 유한개이다).
$Z_i \subset X$를 $Z_{i, \eta}$의 폐포라 하자. $X$를 더 작은 아핀
근방으로 바꾸고 각 $i = 1, \ldots, n$에 대하여
$x' \in Z_i$라고 가정해도 된다. 구성에 의해 $Z \cap X_{s'}$는
$X_{s'}$의 어느 기약성분도 포함하지 않는다. 따라서 Lemma
\ref{more-morphisms-lemma-locally-principal-vertical}에 의해
$Z_\eta \not = \emptyset$이다! 다시 말해 $n \geq 1$이다.
$x_1 \in Z_1$을 일반점으로 놓으면 $x_1 \leadsto x'$이고
$f(x_1) = \eta$이다. 또한 구성에 의해
$Z_{1, \eta} \cap W_j \subset W_j$는 진닫힌 부분집합이다. 따라서
$Z_{1, \eta} \cap W_j$의 각 기약성분은 $X_\eta$에서 여차원이
$\geq 2$인 반면, Algebra, Lemma \ref{algebra-lemma-minimal-over-1}에 의해
$\text{codim}(Z_{1, \eta}, X_\eta) = 1$이다. 그러므로
$W \cap Z_{1, \eta}$는 진닫힌 부분집합이다. 이제 귀납 가설을
$Z_1 \to S$와 특수화 $x_1 \leadsto x'$에 적용할 수 있다. 그러면
$W$에 속하지 않고 $x'$으로 특수화되는 $Z_{1, \eta}$의 닫힌점 $x_2$를
얻는다. 따라서 원하는 대로 $x \leadsto x_2 \leadsto x'$이고, 점 $x_2$는
$X_\eta$에서 닫혀 있으며, $x_2 \not \in W$이다.
\end{proof}

\begin{remark}
\label{more-morphisms-remark-full-specialization-sequence}
Lemma \ref{more-morphisms-lemma-closed-point-nearby-fibre}의 증명은 실제로 특수화의 열
$$
x \leadsto x_1 \leadsto x_2 \leadsto \ldots \leadsto x_d \leadsto x'
$$
이 존재함을 보인다. 여기서 모든 $x_i$는 올 $X_s$에 속하고 각 특수화는
직접 특수화이며, $x_d$는 $X_s$의 닫힌점이다. 정수
$d = \text{trdeg}_{\kappa(s)}(\kappa(x)) = \dim(\overline{\{x\}})$
이며 여기서 폐포는 $X_s$ 안에서 취한다. 또한 점 $x$를 포함하지 않는
$X_s$의 임의의 닫힌 부분집합을 피하도록 점 $x_i$들을 택할 수 있다.
\end{remark}

\noindent
Examples, Section \ref{examples-section-topology-finite-type}는 $A$가
뇌터라는 가정을 하지 않으면 다음 보조정리가 거짓임을 보인다.

\begin{lemma}
\label{more-morphisms-lemma-quasi-finite-quasi-section-meeting-nearby-open}
$\varphi : A \to B$를 국소환의 국소환 준동형이라 하자.
$V \subset \Spec(B)$를 $\mathfrak m_A$ 위에 놓이지 않는 소아이디얼을
적어도 하나 포함하는 열린 부분스킴이라 하자. $A$는 뇌터이고,
$\varphi$는 본질적으로 유한형이며,
$A/\mathfrak m_A \subset B/\mathfrak m_B$는 유한이라고 가정하자.
그러면 다음을 만족하는 $\mathfrak q \in V$가 존재한다.
$\mathfrak m_A \not = \mathfrak q \cap A$이고
$A \to B/\mathfrak q$는 준유한 환 준동형의 국소화이다.
\end{lemma}

\begin{proof}
$A$가 뇌터이고 $A \to B$가 본질적으로 유한형이므로 $B$도 뇌터이다.
Properties, Lemma \ref{properties-lemma-complement-closed-point-Jacobson}에
의해 위상공간 $\Spec(B) \setminus \{\mathfrak m_B\}$는 Jacobson이다.
따라서 공집합이 아닌 열린집합
$$
V \setminus \{\mathfrak q \subset B \mid \mathfrak m_A = \mathfrak q \cap A\}.
$$
에 속하는 닫힌점 $\mathfrak q$를 택할 수 있다. (가정에 의해
공집합이 아니고, $\{\mathfrak m_A\}$가 $\Spec(A)$의 닫힌 부분집합이므로
열려 있다.) 그러면 $\Spec(B/\mathfrak q)$는 두 점
$\mathfrak m_B$와 $\mathfrak q$를 가지며, $\mathfrak q$는
$\mathfrak m_A$ 위에 놓이지 않는다. 어떤 유한형 $A$-대수 $C$와
소아이디얼 $\mathfrak m$에 대하여
$B/\mathfrak q = C_{\mathfrak m}$으로 쓰자. Algebra, Lemma
\ref{algebra-lemma-isolated-point-fibre} (2)에 의해 $A \to C$는
$\mathfrak m$에서 준유한이다. 따라서 Algebra, Lemma
\ref{algebra-lemma-quasi-finite-open}에 의해 $C$를 주 국소화로 바꾸면
환 준동형 $A \to C$가 준유한이다.
\end{proof}

\begin{lemma}
\label{more-morphisms-lemma-quasi-finite-quasi-section-meeting-nearby-open-X}
$f : X \to S$를 스킴의 사상이라 하고, $x \in X$의 상을 $s \in S$라
하자. $U \subset X$를 열린 부분스킴이라 하자. $f$가 국소 유한형이고,
$S$가 국소 뇌터이며, $x$가 $X_s$의 닫힌점이고,
$x' \leadsto x$와 $f(x') \not = s$를 만족하는 점 $x' \in U$가 존재한다고
가정하자. 그러면 다음을 만족하는 닫힌 부분스킴 $Z \subset X$가 존재한다.
(a) $x \in Z$, (b) $f|_Z : Z \to S$는 $x$에서 준유한이고,
(c) $z \in U$, $z \leadsto x$, $f(z) \not = s$를 만족하는
$z \in Z$가 존재한다.
\end{lemma}

\begin{proof}
이는 Lemma \ref{more-morphisms-lemma-quasi-finite-quasi-section-meeting-nearby-open}를
다시 진술한 것이다. 실제로 $A = \mathcal{O}_{S, s}$와
$B = \mathcal{O}_{X, x}$로 놓자. $V \subset \Spec(B)$를 $U$의 역상이라
하자. 환 준동형 $f^\sharp : A \to B$는 본질적으로 유한형이다. 가정에
의해 $\Spec(A)$의 닫힌점으로 가지 않는 $V$의 점이 적어도 하나
존재한다. 따라서 Lemma
\ref{more-morphisms-lemma-quasi-finite-quasi-section-meeting-nearby-open}의 모든 가정이
성립하며, $\mathfrak m_A$ 위에 놓이지 않고
$A \to B/\mathfrak q$가 준유한 환 준동형의 국소화가 되는 소아이디얼
$\mathfrak q \subset B$를 얻는다. $z \in X$를 표준 사상
$\Spec(B) \to X$ 아래에서 점 $\mathfrak q$의 상이라 하자.
$Z = \overline{\{z\}}$에 유도된 축약 스킴 구조를 주자.
$z \leadsto x$이므로 $x \in Z$이고
$\mathcal{O}_{Z, x} = B/\mathfrak q$이다. 구성에 의해 $Z \to S$는
$x$에서 준유한이다.
\end{proof}

\begin{remark}
\label{more-morphisms-remark-alternative-closed-point-nearby-fibre}
Lemma \ref{more-morphisms-lemma-quasi-finite-quasi-section-meeting-nearby-open} 또는 그
변형인
Lemma \ref{more-morphisms-lemma-quasi-finite-quasi-section-meeting-nearby-open-X}
를 사용하면 $S$가 국소 뇌터인 경우 Lemma
\ref{more-morphisms-lemma-closed-point-nearby-fibre}의 다른 증명을 줄 수 있다. 대략의
개요는 다음과 같다. 먼저 $S$를 $s'$에서의 국소환의 스펙트럼으로
바꾼다. 그러면 $\dim(S)$에 관한 귀납법을 사용할 수 있다.
$\dim(S) = 0$인 경우에는 $s' = s$이므로 자명하다. $X$를
$\overline{\{x\}}$ 위의 유도된 축약 스킴 구조로 바꾸자.
$X \to S$, $x' \mapsto s'$, 그리고 $x$를 포함하는 임의의 공집합이 아닌
열린집합 $U \subset X$에 Lemma
\ref{more-morphisms-lemma-quasi-finite-quasi-section-meeting-nearby-open-X}를 적용한다.
그러면
% P05-EN-CH37-0233
닫힌 부분스킴 $x' \in Z \subset X$와 점 $z \in Z$를 얻는데,
$Z \to S$는 $x'$에서 준유한이고 $f(z) \not = s'$이다. 이때 $z$는
$X_{f(z)}$의 닫힌점이고 $z \leadsto x'$이다. $f(z) \not = s'$이므로
$\dim(\mathcal{O}_{S, f(z)}) < \dim(S)$이다. $x$는 $X$의 일반점이므로
$x \leadsto z$이고, 따라서 $s = f(x) \leadsto f(z)$이다.
% P05-EN-CH37-0234
귀납 가설을 $s \leadsto f(z)$와 $z \mapsto f(z)$에 적용하면 증명이
끝난다.
\end{remark}

\begin{lemma}
\label{more-morphisms-lemma-change-hypotheses}
$f : X \to S$가 국소 유한형이고, $S$가 국소 뇌터이며, $x \in X$가
그 올 $X_s$의 닫힌점이고, $U \subset X$가
$U \cap X_s = \emptyset$와 $x \in \overline{U}$를 만족하는 열린
부분스킴이라고 가정하자. 그러면
Lemma \ref{more-morphisms-lemma-quasi-finite-quasi-section-meeting-nearby-open-X}
의 결론들이 성립한다.
\end{lemma}

\begin{proof}
다음과 같이 인용한 보조정리로 환원할 수 있다. 먼저 $X$와 $S$를 각각
$x$와 $s$의 아핀 근방으로 바꾼다. 그러면 $X$는 뇌터이고, 특히 $U$는
준콤팩트이다(Morphisms, Lemma
\ref{morphisms-lemma-finite-type-noetherian}와 Topology, Lemmas
\ref{topology-lemma-Noetherian},
\ref{topology-lemma-Noetherian-quasi-compact}를 보라). 따라서
$x' \in U$인 특수화 $x' \leadsto x$가 존재한다(Morphisms, Lemma
\ref{morphisms-lemma-reach-points-scheme-theoretic-image}를 보라).
$f(x') \not = s$임에 유의하자. 그러므로 보조정리의 모든 가정이
성립하고 증명이 끝난다.
\end{proof}



\section{Stein 분해}
\label{more-morphisms-section-stein-factorization}

\noindent
Stein 분해는 $f_*\mathcal{O}_X = \mathcal{O}_S$인 고유 사상
$f : X \to S$가 연결된 올들을 갖는다는 명제이다.

\begin{lemma}
\label{more-morphisms-lemma-stein-universally-closed}
$S$를 스킴이라 하고, $f : X \to S$를 보편적으로 닫힌 준분리 사상이라
하자. 다음 분해가 존재한다.
$$
\xymatrix{
X \ar[rr]_{f'} \ar[rd]_f & & S' \ar[dl]^\pi \\
& S &
}
$$
이 분해는 다음 성질들을 갖는다.
\begin{enumerate}
\item 사상 $f'$은 보편적으로 닫혔고 준콤팩트이며 준분리이고 전사이다.
\item 사상 $\pi : S' \to S$는 정수적이다.
\item $f'_*\mathcal{O}_X = \mathcal{O}_{S'}$이다.
\item $S' = \underline{\Spec}_S(f_*\mathcal{O}_X)$이다.
\item $S'$은 $X$에서 $S$의 정규화이다. Morphisms, Definition
\ref{morphisms-definition-normalization-X-in-Y}를 보라.
\end{enumerate}
분해 $f = \pi \circ f'$의 형성은 평탄 밑변환과 가환한다.
\end{lemma}

\begin{proof}
Morphisms, Lemma
\ref{morphisms-lemma-universally-closed-quasi-compact}에 의해 사상 $f$는
준콤팩트이다. 따라서 $X$에서 $S$의 정규화 $S'$이 정의되고(Morphisms,
Definition \ref{morphisms-definition-normalization-X-in-Y}), 분해
$X \to S' \to S$를 얻는다. Morphisms, Lemma
\ref{morphisms-lemma-normalization-in-universally-closed}에 의해 (2),
(4), (5)가 성립한다. 사상 $f'$은 Morphisms, Lemma
\ref{morphisms-lemma-image-proper-scheme-closed}에 의해 보편적으로
닫혔고, Schemes, Lemma
\ref{schemes-lemma-quasi-compact-permanence}에 의해 준콤팩트이며,
Schemes, Lemma \ref{schemes-lemma-compose-after-separated}에 의해
준분리이다.

\medskip\noindent
남은 명제들을 보이기 위해 밑스킴 $S$가 아핀이라고 가정해도 된다.
$S = \Spec(R)$이라 하자. 그러면 $A = \Gamma(X, \mathcal{O}_X)$가 정수적
$R$-대수인 $S' = \Spec(A)$이다. 따라서
$f'_*\mathcal{O}_X$가 $\mathcal{O}_{S'}$임은 명백하다. 실제로 Schemes,
Lemma \ref{schemes-lemma-push-forward-quasi-coherent}에 의해
$f'_*\mathcal{O}_X$는 준연접이고, 따라서 $\widetilde{A}$와 같다. 이로써
(3)이 증명된다.

\medskip\noindent
$f'$이 전사임을 보이자. $f'$은 보편적으로 닫혔으므로(위를 보라)
$f'$의 상은 닫힌 부분집합 $V(I) \subset S' = \Spec(A)$이다.
$h \in I$를 택하자. 그러면
% P05-EN-CH37-0235
$h|_X = f^\sharp(h)$는 모든 점에서 소멸하는 $X$의 구조층의 대역단면이다.
$X$가 준콤팩트이므로 이는 $h|_X$가 멱영 단면이라는 뜻이다. 즉 어떤
$n > 0$에 대하여
% P05-EN-CH37-0236
$h^n|X = 0$이다. 그런데 $A = \Gamma(X, \mathcal{O}_X)$이므로
$h^n = 0$이다. $I$의 모든 원소가 멱영이므로 원하는 대로
$V(I) = S'$임을 얻는다.

\medskip\noindent
Cohomology of Schemes, Lemma
\ref{coherent-lemma-flat-base-change-cohomology}에 의해
$f_*\mathcal{O}_X$의 형성은 평탄 밑변환과 가환한다. Constructions,
Lemma \ref{constructions-lemma-spec-properties}에 의해 상대 스펙트럼의
형성은 임의의 밑변환과 가환한다. 따라서 이 분해의 형성은 평탄
밑변환과 가환한다.
\end{proof}

\begin{lemma}
\label{more-morphisms-lemma-stein-universally-closed-residue-fields}
Lemma \ref{more-morphisms-lemma-stein-universally-closed}에서 $f$가 국소 유한형이라고
추가로 가정하자. 그러면 $s \in S$에 대하여 올
$\pi^{-1}(\{s\}) = \{s_1, \ldots, s_n\}$은 유한하고, 체 확대
$\kappa(s_i)/\kappa(s)$도 유한하다.
\end{lemma}

\begin{proof}
$\pi^{-1}(\{s\})$의 점들 사이에는 특수화가 없음을 상기하자. Algebra,
Lemma \ref{algebra-lemma-integral-no-inclusion}를 보라.
$f'$이 전사이므로 $|X_s| \to \pi^{-1}(\{s\})$도 전사이다.
$X_s$는 체 위의 준분리 유한형 스킴이다(준콤팩트성은 인용한 보조정리의
증명에서 보였다). 따라서 $X_s$는 뇌터 스킴이다(Morphisms, Lemma
\ref{morphisms-lemma-finite-type-noetherian}). 이제 위상적 논증(생략함)에
의해 $\pi^{-1}(\{s\})$가 유한함을 알 수 있다. 각 $i$에 대하여 $s_i$로
가는 유한형 점 $x_i \in X_s$를 택할 수 있다(Morphisms, Lemma
\ref{morphisms-lemma-enough-finite-type-points}). 이로부터
$\kappa(s_i)/\kappa(s)$가 유한함을 얻는다. 실제로 $x_i$는 유한형 사상
$\Spec(k_i) \to X_s$로 나타낼 수 있으므로(유한형 점의 정의에 의해),
$\Spec(k_i) \to s = \Spec(\kappa(s))$도 유한형 사상들의 합성으로서
유한형이다. 따라서 $k_i/\kappa(s)$는 유한하다(Morphisms, Lemma
\ref{morphisms-lemma-point-finite-type}).
\end{proof}

\begin{lemma}
\label{more-morphisms-lemma-characterize-geometrically-connected-fibres}
$f : X \to S$를 스킴의 사상이라 하고 $s \in S$라 하자. 그러면
$X_s$가 기하적으로 연결일 필요충분조건은 모든 에탈 근방
$(U, u) \to (S, s)$에 대하여 밑변환 $X_U \to U$의 올 $X_u$가
연결인 것이다.
\end{lemma}

\begin{proof}
$X_s$가 기하적으로 연결이면 그 임의의 밑변환도 연결이다. 역으로
$X_s$가 기하적으로 연결이 아니라고 하자. 그러면 Varieties, Lemma
\ref{varieties-lemma-characterize-geometrically-disconnected}에 의해
어떤 유한 분리 체 확대 $k/\kappa(s)$에 대하여
$X_s \times_{\Spec(\kappa(s))} \Spec(k)$가 연결이 아니다. Lemma
\ref{more-morphisms-lemma-realize-prescribed-residue-field-extension-etale}에 의해
$\kappa(u)/\kappa(s)$가 $k/\kappa(s)$와 동일시되는 아핀 에탈 근방
$(U, u) \to (S, s)$가 존재한다. 이 경우 $X_u$는 연결이 아니다.
\end{proof}

\begin{theorem}[Stein 분해; 뇌터 경우]
\label{more-morphisms-theorem-stein-factorization-Noetherian}
$S$를 국소 뇌터 스킴이라 하고, $f : X \to S$를 고유 사상이라 하자.
다음 분해가 존재한다.
$$
\xymatrix{
X \ar[rr]_{f'} \ar[rd]_f & & S' \ar[dl]^\pi \\
& S &
}
$$
이 분해는 다음 성질들을 갖는다.
\begin{enumerate}
\item 사상 $f'$은 고유이고 그 올들은 기하적으로 연결이다.
\item 사상 $\pi : S' \to S$는 유한이다.
\item $f'_*\mathcal{O}_X = \mathcal{O}_{S'}$이다.
\item $S' = \underline{\Spec}_S(f_*\mathcal{O}_X)$이다.
\item $S'$은 $X$에서 $S$의 정규화이다. Morphisms, Definition
\ref{morphisms-definition-normalization-X-in-Y}를 보라.
\end{enumerate}
\end{theorem}

\begin{proof}
Lemma \ref{more-morphisms-lemma-stein-universally-closed}의 분해를
$f = \pi \circ f'$이라 하자. Lemma
\ref{more-morphisms-lemma-stein-universally-closed}의 결론들에 더하여 $f'$은
분리 사상이고(Schemes, Lemma
\ref{schemes-lemma-compose-after-separated}) 유한형 사상이다(Morphisms,
Lemma \ref{morphisms-lemma-permanence-finite-type}). 따라서 $f'$은
고유이다. Cohomology of Schemes, Proposition
\ref{coherent-proposition-proper-pushforward-coherent}에 의해
$f_*\mathcal{O}_X$는 연접 $\mathcal{O}_S$-가군이다. 따라서 $\pi$가
유한임을, 즉 (2)를 얻는다.

\medskip\noindent
이로써 가장 흥미로운 명제, 즉 $f'$의 모든 올이 기하적으로 연결이라는
명제를 제외한 나머지는 모두 증명되었다. 위 논의에서 $S$를 $S'$으로
대체해도 됨은 명백하다. 따라서 $S$가 뇌터 아핀 스킴이고,
$f : X \to S$가 고유이며, $f_*\mathcal{O}_X = \mathcal{O}_S$라고
가정할 수 있다. $s \in S$를 $S$의 한 점이라 하자. $X_s$가 기하적으로
연결임을 보여야 한다. Lemma
\ref{more-morphisms-lemma-characterize-geometrically-connected-fibres}에 의해 모든
에탈 근방 $(U, u) \to (S, s)$에 대하여 $X_u$가 연결임을 보이면
충분하다. $U$는 아핀이라고 가정해도 된다. 그러면 $U$는 뇌터이고
(Morphisms, Lemma \ref{morphisms-lemma-finite-type-noetherian}), 밑변환
$f_U : X_U \to U$는 고유이며(Morphisms, Lemma
\ref{morphisms-lemma-base-change-proper}),
$(f_U)_*\mathcal{O}_{X_U} = \mathcal{O}_U$도 성립한다(Cohomology of
Schemes, Lemma \ref{coherent-lemma-flat-base-change-cohomology}).
따라서 $(f : X \to S, s)$를 밑변환 $(f_U : X_U \to U, u)$로 대체하면,
$f_*\mathcal{O}_X = \mathcal{O}_S$일 때 올 $X_s$가 연결임을 증명하는
것으로 충분하다. 이는 $X_s$의 구조층의 멱등원들을 살펴서 Derived
Categories of Schemes, Lemma
\ref{perfect-lemma-proper-idempotent-on-fibre}로부터 따르지만, 아래에서
직접적인 논증도 제시하겠다.

\medskip\noindent
구체적으로 형식함수 정리, 더 정확히는 Cohomology of Schemes, Lemma
\ref{coherent-lemma-formal-functions-stalk}를 적용한다. 이 보조정리는
$$
\mathcal{O}^\wedge_{S, s} = (f_*\mathcal{O}_X)_s^\wedge =
\lim_n H^0(X_n, \mathcal{O}_{X_n})
$$
임을 말한다. 여기서 $X_n$은 $X_s$의 $n$차 무한소 근방이다.
$X_n$의 바탕 위상공간은 $X_s$의 바탕 위상공간과 같으므로,
$X_s = T_1 \amalg T_2$가 공집합이 아닌 열린닫힌 부분스킴들의
서로소 합이면 모든 $n$에 대하여 마찬가지로
$X_n = T_{1, n} \amalg T_{2, n}$이다. 따라서
$H^0(X_n, \mathcal{O}_{X_n})$에는 자명하지 않은 멱등원 $e_{1, n}$,
즉 $T_{1, n}$에서 항등적으로 $1$이고 $T_{2, n}$에서 항등적으로
$0$인 함수가 들어 있다. $e_{1, n + 1}$을 $X_n$에 제한하면
$e_{1, n}$이 됨은 명백하다. 따라서 $e_1 = \lim e_{1, n}$은 이
극한의 자명하지 않은 멱등원이다. 이는 $\mathcal{O}^\wedge_{S, s}$가
국소환이라는 사실과 모순이다. 그러므로 가정이 잘못되었고, 즉
$X_s$는 연결이므로 증명이 끝난다.
\end{proof}

\begin{theorem}[Stein 분해; 일반적인 경우]
\label{more-morphisms-theorem-stein-factorization-general}
$S$를 스킴이라 하고 $f : X \to S$를 고유 사상이라 하자. 다음 분해가
존재한다.
$$
\xymatrix{
X \ar[rr]_{f'} \ar[rd]_f & & S' \ar[dl]^\pi \\
& S &
}
$$
이 분해는 다음 성질들을 갖는다.
\begin{enumerate}
\item 사상 $f'$은 고유이고 그 올들은 기하적으로 연결이다.
\item 사상 $\pi : S' \to S$는 정수적이다.
\item $f'_*\mathcal{O}_X = \mathcal{O}_{S'}$이다.
\item $S' = \underline{\Spec}_S(f_*\mathcal{O}_X)$이다.
\item $S'$은 $X$에서 $S$의 정규화이다. Morphisms, Definition
\ref{morphisms-definition-normalization-X-in-Y}를 보라.
\end{enumerate}
\end{theorem}

\begin{proof}
Lemma \ref{more-morphisms-lemma-stein-universally-closed}를 적용하여 사상
$f' : X \to S'$을 얻을 수 있다. Lemma
\ref{more-morphisms-lemma-stein-universally-closed}의 결론들에 더하여 $f'$은 분리
사상이고(Schemes, Lemma
\ref{schemes-lemma-compose-after-separated}) 유한형 사상이다(Morphisms,
Lemma \ref{morphisms-lemma-permanence-finite-type}). 따라서 $f'$은
고유이다. 이 시점에서 $f'$의 올들이 기하적으로 연결이라는 명제를
제외한 모든 명제가 증명되었다.

\medskip\noindent
$S = \Spec(R)$가 아핀이라고 가정해도 된다.
$R' = \Gamma(X, \mathcal{O}_X)$로 놓자. 그러면 $S' = \Spec(R')$이다.
따라서 $S$를 $S'$으로 대체하여
% P05-EN-CH37-0237
$S = \Spec(R)$가 아핀이고 $R = \Gamma(X, \mathcal{O}_X)$라고 가정할
수 있다. 이제 $s \in S$를 한 점이라 하자. $U \to S$를 아핀 스킴들의
에탈 사상이라 하고, $u \in U$를 $s$로 가는 점이라 하자.
$X_U \to U$를 $X$의 밑변환이라 하자. Lemma
\ref{more-morphisms-lemma-characterize-geometrically-connected-fibres}에 의해
$X_U \to U$의 $u$ 위의 올이 연결임을 보이면 충분하다. Cohomology of
Schemes, Lemma \ref{coherent-lemma-flat-base-change-cohomology}에 의해
$\Gamma(X_U, \mathcal{O}_{X_U}) = \Gamma(U, \mathcal{O}_U)$이다.
따라서 다음을 보여야 한다.
% P05-EN-CH37-0238
$S = \Spec(R)$가 아핀이고, $X \to S$가
$\Gamma(X, \mathcal{O}_X) = R$을 만족하는 고유 사상이며, $s \in S$가
한 점일 때, 올 $X_s$는 연결이다.

\medskip\noindent
이를 위해서는 $e \in H^0(X_s, \mathcal{O}_{X_s})$인 멱등원이 $0$과
$1$뿐임을 보이면 충분하다(Lemma
\ref{more-morphisms-lemma-stein-universally-closed}에 의해 $X_s$가 공집합이 아님은
이미 알고 있다). Derived Categories of Schemes, Lemma
\ref{perfect-lemma-proper-idempotent-on-fibre}에 의해 $R$을 한 주원소
국소화로 대체한 뒤 $e$가 $R$의 한 원소의 상이라고 가정할 수 있다.
$R \to H^0(X_s, \mathcal{O}_{X_s})$는 $\kappa(s)$를 통해 인수분해되므로
결론이 따른다.
\end{proof}

\noindent
다음은 한 응용이다.

\begin{lemma}
\label{more-morphisms-lemma-geometrically-connected-fibres-towards-normal}
$f : X \to S$를 스킴의 사상이라 하자. 다음을 가정하자.
\begin{enumerate}
\item $f$는 고유이다.
\item $S$는 일반점이 $\xi$인 정역 스킴이다.
\item $S$는 정규이다.
\item $X$는 축약이다.
\item $X$의 각 기약성분의 모든 일반점은 $\xi$로 간다.
\item $H^0(X_\xi, \mathcal{O}) = \kappa(\xi)$이다.
\end{enumerate}
그러면 $f_*\mathcal{O}_X = \mathcal{O}_S$이고, $f$의 올들은
기하적으로 연결이다.
\end{lemma}

\begin{proof}
Theorem \ref{more-morphisms-theorem-stein-factorization-general}을 적용하여 분해
$X \to S' \to S$를 얻는다. $S' = S$임을 보이면 충분하며, 이는
Morphisms, Lemma \ref{morphisms-lemma-finite-birational-over-normal}에서
따를 것이다. 실제로 $X$가 축약이므로 $S'$도 축약이다(Morphisms,
Lemma \ref{morphisms-lemma-normalization-in-reduced}). 위에서 인용한
정리에 의해 사상 $S' \to S$는 정수적이다. Morphisms, Lemma
\ref{morphisms-lemma-normalization-generic}와 가정 (5)에 의해
$S'$의 모든 일반점은 $\xi$ 위에 놓인다. 한편 $S'$은
$f_*\mathcal{O}_X$의 상대 스펙트럼이므로, 스킴론적 올 $S'_\xi$는
$H^0(X_\xi, \mathcal{O})$의 스펙트럼이고, 가정에 의해 이는
$\kappa(\xi)$와 같다. 따라서 $S'$은 함수체가 $S$의 함수체와 같은
정역 스킴이다. 이로써 증명이 끝난다.
\end{proof}

\noindent
또 다른 응용은 다음과 같다.

\begin{lemma}
\label{more-morphisms-lemma-proper-flat-nr-geom-conn-comps-lower-semicontinuous}
$X \to S$를 유한 표시인 평탄 고유 사상이라 하자.
% P05-EN-CH37-0239
$n_{X/S}$를 Lemma
\ref{more-morphisms-lemma-base-change-fibres-nr-geometrically-connected-components}에서
도입한, $f$의 올들의 기하적 연결 성분의 개수를 세는 $S$ 위의 함수라
하자. 그러면 $n_{X/S}$는 하반연속이다.
\end{lemma}

\begin{proof}
$s \in S$라 하고 $n = n_{X/S}(s)$로 놓자. $s$에서 $X \to S$의 기하적
올은 체 위의 고유 스킴이므로 뇌터이고, 따라서 연결 성분이 유한
개이므로 $n < \infty$임에 유의하자. $n_{X/S}|_V \geq n$을 만족하는
$s$의 열린 근방 $V$를 찾아야 한다. Theorem
\ref{more-morphisms-theorem-stein-factorization-general}에서와 같이
$X \to S' \to S$를 Stein 분해라 하자. Lemma
\ref{more-morphisms-lemma-stein-universally-closed-residue-fields}에 의해 $s$ 위에 놓인
점 $s'_1, \ldots, s'_m \in S'$은 유한 개이고, 확대
$\kappa(s'_i)/\kappa(s)$는 유한하다. 이제 Lemma
\ref{more-morphisms-lemma-etale-makes-integral-split}에 의해 $S$를 $s$의 한 에탈
근방으로 대체한 뒤, 스킴으로서 $S' = V_1 \amalg \ldots \amalg V_m$이고
$s'_i \in V_i$이며 $\kappa(s'_i)/\kappa(s)$가 순비분리라고 가정할 수
있다. 그러면 스킴 $X_{s_i'}$는 $\kappa(s)$ 위에서 기하적으로
연결이므로 $m = n$이다. 스킴 $X_i = (f')^{-1}(V_i)$,
$i = 1, \ldots, n$은 $S$ 위에서 평탄이고 유한 표시이다. 따라서
$X_i \to S$의 상은 열려 있다(Morphisms, Lemma
\ref{morphisms-lemma-fppf-open}). 그러므로 $s$의 한 근방에서
$n_{X/S}$는 적어도 $n$이다.
\end{proof}

\begin{lemma}
\label{more-morphisms-lemma-proper-flat-geom-red}
$f : X \to S$를 스킴의 사상이라 하자. 다음을 가정하자.
\begin{enumerate}
\item $f$는 고유이고 평탄이며 유한 표시이다.
\item $f$의 기하적 올들은 축약이다.
\end{enumerate}
% P05-EN-CH37-0240
그러면 $f$의 올들의 기하적 연결 성분의 개수를 세는 함수
$n_{X/S} : S \to \mathbf{Z}$는 국소 상수이다.
\end{lemma}

\begin{proof}
Lemma \ref{more-morphisms-lemma-proper-flat-nr-geom-conn-comps-lower-semicontinuous}에
의해 함수
% P05-EN-CH37-0241
$n_{X/S}$는 하반연속이다. $s \in S$에 대하여 다음
$\kappa(s)$-대수를 생각하자.
$$
A_s = H^0(X_s, \mathcal{O}_{X_s})
$$
% P05-EN-CH37-0242
$s$를 $\dim_{\kappa(s)} A_s$로 보내는 함수
$\beta_0 : X \to \mathbf{Z}$를 정의하자. Varieties, Lemma
\ref{varieties-lemma-proper-geometrically-reduced-global-sections}와
$X_s$가 기하적으로 축약이라는 사실에 의해 $A_s$는
$\kappa(s)$의 유한 분리 확대들의 유한 곱이다. 따라서
$A_s \otimes_{\kappa(s)} \kappa(\overline{s})$는
$\kappa(\overline{s})$의 $\beta_0(s)$개 복사본의 곱이다. 그러므로
$X_{\overline{s}}$에는 $\beta_0(s)$개의 연결 성분이 있다. 다시 말해
$S$ 위의 함수로서 $n_{X/S} = \beta_0$이다. 따라서 Derived Categories
of Schemes, Lemma \ref{perfect-lemma-jump-loci-geometric}에 의해
$n_{X/S}$는 상반연속이다. 이로써 증명이 끝난다.
\end{proof}

\noindent
마지막 응용이다.

\begin{lemma}
\label{more-morphisms-lemma-split-off-proper-part-henselian}
\begin{reference}
아딕 뇌터 밑인 경우의 참고문헌은
\cite[III, Proposition 5.5.1]{EGA}이다.
\end{reference}
$(A, I)$를 헨젤 쌍이라 하자. $X \to \Spec(A)$가 분리이고 유한형이라고
하자. $X_0 = X \times_{\Spec(A)} \Spec(A/I)$로 놓자.
$Y \to \Spec(A/I)$가 고유인 열린닫힌 부분스킴 $Y \subset X_0$를
생각하자. 그러면 $A$ 위에서 고유이고
$W \times_{\Spec(A)} \Spec(A/I) = Y$를 만족하는 열린닫힌 부분스킴
$W \subset X$가 존재한다.
\end{lemma}

\begin{proof}
$\Spec(A/I) \to \Spec(A)$에 의한 밑변환을 $T \mapsto T_0$로
표기하겠다. Chow 보조정리(Limits, Lemma
\ref{limits-lemma-chow-finite-type}의 형태)에 의해, $X'$이
$\mathbf{P}^n_A$로의 몰입을 갖게 하는 전사 고유 사상
$\varphi : X' \to X$가 존재한다. $Y' = \varphi^{-1}(Y)$로 놓자.
이는 $X'_0$의 열린닫힌 부분스킴이다. 이 보조정리가 $(X', Y')$에
대하여 성립한다고 하자. $Y' = W'_0$을 만족하며 $A$ 위에서 고유인
열린닫힌 부분스킴을 $W' \subset X'$이라 하자. Morphisms, Lemma
\ref{morphisms-lemma-image-proper-scheme-closed}에 의해
$W = \varphi(W') \subset X$와
$Q = \varphi(X' \setminus W') \subset X$는 닫힌 부분집합이고,
Morphisms, Lemma \ref{morphisms-lemma-image-proper-is-proper}에 의해
$W$는 $A$ 위에서 고유이다. $W \cap Q$의 $\Spec(A)$에서의 상은 닫혀
있다. $(A, I)$가 헨젤이므로 $W \cap Q$가 공집합이 아니면
$W \cap Q$에는 $\Spec(A/I)$ 위에 놓인 점이 있다. 그러나
$W'_0 = Y' = \varphi^{-1}(Y)$이므로 이는 불가능하다. 따라서 $W$는
$W_0 = Y$를 만족하며 $A$ 위에서 고유인 $X$의 열린닫힌 부분스킴이다.
이로써 다음 문단에서 설명하는 경우로 환원된다.

\medskip\noindent
$A$ 위의 몰입 $j : X \to \mathbf{P}^n_A$가 존재한다고 하자.
$\overline{X}$를 $j$의 스킴론적 상이라 하자. $j$는 준콤팩트
사상이므로(Schemes, Lemma
\ref{schemes-lemma-quasi-compact-permanence}),
$j : X \to \overline{X}$는 열린 몰입이다(Morphisms, Lemma
\ref{morphisms-lemma-quasi-compact-immersion}). 따라서 밑변환
$j_0 : X_0 \to \overline{X}_0$도 열린 몰입이다. 그러므로
$j_0(Y) \subset \overline{X}_0$는 열려 있고, Morphisms, Lemma
\ref{morphisms-lemma-image-proper-scheme-closed}에 의해 닫혀 있기도
하다. 이 보조정리가 $(\overline{X}, j_0(Y))$에 대하여 성립한다고
하자. $j_0(Y) = \overline{W}_0$을 만족하며 $A$ 위에서 고유인 대응하는
열린닫힌 부분스킴을 $\overline{W} \subset \overline{X}$라 하자.
그러면 $T = \overline{W} \setminus j(X)$는 $\overline{W}$에서 닫혀
있으므로, $\overline{W}$가 $A$ 위에서 고유임에 의해 그
$\Spec(A)$에서의 상은 닫혀 있다. $(A, I)$가 헨젤이므로 $T$가
공집합이 아니면 $T$에는 $\Spec(A/I)$로 가는 점이 있다. 그러나
$j_0(Y) = \overline{W}_0$이 $j(X)$에 포함되므로 이는 불가능하다.
따라서 $\overline{W}$는 $j(X)$에 포함되고, $W \subset X$를 $j$에 의해
$\overline{W}$와 동형으로 가는 유일한 열린닫힌 부분스킴으로 둘 수
있다. 이로써 다음 문단에서 설명하는 경우로 환원된다.

\medskip\noindent
$X \subset \mathbf{P}^n_A$가 닫힌 부분스킴이라고 하자. 그러면
$X \to \Spec(A)$는 고유 사상이다. $Z = X_0 \setminus Y$로 놓자.
이는 $X_0$의 열린닫힌 부분스킴이고 $X_0 = Y \amalg Z$이다. Theorem
\ref{more-morphisms-theorem-stein-factorization-general}에서와 같이
$X \to X' \to \Spec(A)$를 Stein 분해라 하자. $Y$와 $Z$의 상을 각각
$Y' \subset X'_0$와 $Z' \subset X'_0$라 하자.
% P05-EN-CH37-0243
$X \to Z$의 올들이 기하적으로 연결이므로
$Y' \cap Z' = \emptyset$이다. $X \to X'$이 전사이므로
$X'_0 = Y' \amalg Z'$이다. $X' \to \Spec(A)$가 정수적이므로 $X'$은
$A$ 위에서 정수적인 $A$-대수의 스펙트럼이다. 스펙트럼의 열린닫힌
부분집합은 대응하는 환의 멱등원과 $1$-대-$1$로 대응함을 상기하자.
Algebra, Lemma \ref{algebra-lemma-disjoint-decomposition}를 보라.
따라서 More on Algebra, Lemma
\ref{more-algebra-lemma-characterize-henselian-pair}에 의해
$X' = W' \amalg V'$로 쓸 수 있다. 여기서 $W'$과 $V'$은 열리고 닫혀
있으며, $Y' = W'_0$이고 $Z' = V'_0$이다.
% P05-EN-CH37-0244
$W$를 $X$ 안의 역상으로 두면 증명이 끝난다.
\end{proof}









\section{일반 평탄성 층화}
\label{more-morphisms-section-generic-flatness-stratification}

\noindent
일반 평탄성을 사용하여 주어진 가군이 각 층 위에서 평탄해지는 밑의
층화를 구성할 수 있다.

\begin{lemma}[일반 평탄성 층화]
\label{more-morphisms-lemma-generic-flatness-stratification}
$f : X \to S$를 준콤팩트 준분리 스킴들 사이의 유한 표시 사상이라
하자. $\mathcal{F}$를 유한 표시 $\mathcal{O}_X$-가군이라 하자. 그러면
어떤 $t \geq 0$과 닫힌 부분스킴들
$$
S \supset S_0 \supset S_1 \supset \ldots \supset S_t = \emptyset
$$
이 존재하여, $S_i \to S$는 유한형 아이디얼층으로 정의되고,
$S_0 \subset S$는 비후화이며, $\mathcal{F}$를
$X \times_S (S_i \setminus S_{i + 1})$로 당긴 가군은
$S_i \setminus S_{i + 1}$ 위에서 평탄하다.
\end{lemma}

\begin{proof}
다음 Cartesian 도식을 찾을 수 있다.
$$
\xymatrix{
X \ar[d] \ar[r] & X_0 \ar[d] \\
S \ar[r] & S_0
}
$$
또한 $X_0$와 $S_0$가 $\mathbf{Z}$ 위에서 유한형이고, 당겨서
$\mathcal{F}$가 되는 유한 표시 $\mathcal{O}_{X_0}$-가군
$\mathcal{F}_0$를 찾을 수 있다. Limits, Proposition
\ref{limits-proposition-approximate}와 Lemmas
\ref{limits-lemma-descend-finite-presentation},
\ref{limits-lemma-descend-modules-finite-presentation}를 보라. 따라서
$X$와 $S$가 $\mathbf{Z}$ 위에서 유한형이고, $\mathcal{F}$가 연접
$\mathcal{O}_X$-가군이라고 가정해도 된다.

\medskip\noindent
$X$와 $S$가 $\mathbf{Z}$ 위에서 유한형이고 $\mathcal{F}$가 연접
$\mathcal{O}_X$-가군이라고 가정하자. 이 경우 모든 준연접 아이디얼은
유한형이므로 $S_i$가 유한형 아이디얼로 잘려 나온다는 조건은 확인할
필요가 없다. $S_0 = S_{red}$를 $S$의 축약으로 놓자. Morphisms,
Proposition \ref{morphisms-proposition-generic-flatness-reduced}의 일반
평탄성에 의해 조밀한 열린집합 $U_0 \subset S_0$가 존재하여,
$\mathcal{F}$를 $X \times_S U_0$로 당긴 가군은 $U_0$ 위에서 평탄하다.
$S_1 \subset S_0$를 바탕 닫힌 부분집합이 $S \setminus U_0$인 축약
닫힌 부분스킴이라 하자. $S_1 \not = \emptyset$이면 이 과정을 계속하여
$S_0 \supset S_1 \supset \ldots$를 얻는다. $S$가 뇌터이므로 닫힌
부분집합들의 모든 내림 사슬은 안정화한다. 따라서 어떤 $t \geq 0$에
대하여 $S_t = \emptyset$이다.
\end{proof}

\begin{lemma}
\label{more-morphisms-lemma-generic-flatness-stratification-scheme}
$f : X \to S$를 준콤팩트 준분리 스킴들 사이의 유한 표시 사상이라
하자. 그러면 어떤 $t \geq 0$과 닫힌 부분스킴들
$$
S \supset S_0 \supset S_1 \supset \ldots \supset S_t = \emptyset
$$
이 존재하여, $S_i \to S$는 유한형 아이디얼층으로 정의되고,
$S_0 \subset S$는 비후화이며,
$X \times_S (S_i \setminus S_{i + 1})$는
$S_i \setminus S_{i + 1}$ 위에서 평탄하다.
\end{lemma}

\begin{proof}
$\mathcal{F} = \mathcal{O}_X$로 두고 Lemma
\ref{more-morphisms-lemma-generic-flatness-stratification}을 적용한다.
\end{proof}

\begin{lemma}[Longke Tang]
\label{more-morphisms-lemma-flatness-criterion}
$f : X \to S$를 전사 유한 표시 사상이라 하자.
$M$을 $i > 0$일 때 $H^i(M) = 0$을 만족하는
$D_\QCoh(\mathcal{O}_S)$의 대상이라 하자. 다음은 서로 동치이다.
\begin{enumerate}
\item $M$은 차수 $0$에 놓인 평탄 $\mathcal{O}_S$-가군과 동형이다.
\item $Lf^*M$은 차수 $0$에 놓인 평탄 $\mathcal{O}_X$-가군과 동형이다.
\end{enumerate}
\end{lemma}

\begin{proof}
(1)이 (2)를 함의한다는 증명은 생략한다. (2)를 가정하자. (1)의 증명은
국소적인 문제이므로 $S$가 아핀이라고 가정해도 된다. Lemma
\ref{more-morphisms-lemma-generic-flatness-stratification-scheme}에서와 같은 층화
$S \supset S_0 \supset S_1 \supset \ldots \supset S_t = \emptyset$
를 택하자.
% P05-EN-CH37-0245
$T_i = S_i \setminus S_{i - 1}$와 $Y_i = X \times_S T_i$로 놓자.
포함 사상을 $g_i : T_i \to S$와 $h_i : Y_i \to X$로 나타내자.
$f$가 전사이므로 사상
% P05-EN-CH37-0246
$f_i : X_i \to T_i$는 전사이고 평탄하다.
$Lh_i^* \circ Lf^* = Lf_i^* \circ Lg_i^*$이므로 (2)에 의해
$Lf_i^* Lg_i^*M$은 차수 $0$에 놓인 평탄
$\mathcal{O}_{Y_i}$-가군과 동형이다. $f_i$가 평탄이고 전사이므로
$Lg_i^*M$은 차수 $0$에 놓인 평탄 $\mathcal{O}_{T_i}$-가군과
동형이다.

\medskip\noindent
이로부터 결론이 따름을 $t$에 대한 귀납법으로 증명하겠다. 다음 표기를
사용한다. $S = \Spec(A)$이고, $M$은 $D(A)$의 $N$에 대응하는
$D_\QCoh(\mathcal{O}_S)$의 대상이다. Derived Categories of Schemes,
Lemma \ref{perfect-lemma-affine-compare-bounded}를 보라.

\medskip\noindent
기초 단계는 $t = 1$인 경우이다. $S_0 = \Spec(A/I)$로 쓰자. 여기서
% P05-EN-CH37-0247
$I$는 유한 생성 아이디얼이다. 집합론적으로 $S_0 = S$이므로 $I$는
멱영이다. 가정에 의해 $N \otimes_A^\mathbf{L} A/I$의 Tor 진폭은
$[0, 0]$ 안에 있다. More on Algebra, Lemma
\ref{more-algebra-lemma-check-tor-dimension-modulo-nilpotent}에 의해
$N$에 대해서도 마찬가지이다.

\medskip\noindent
귀납 단계에서는 $t > 1$이라고 하자.
$S_{t - 1} = \Spec(A/I)$로 쓰자. 여기서
$I = (f_1, \ldots, f_r)$는 유한 생성 아이디얼이다. $r$에 대한
귀납법으로 논증하겠다. 귀납 가정에 의해 $D(A_{f_r})$에서
$N_{f_r} = N \otimes_A^\mathbf{L} A_{f_r}$의 Tor 진폭은
$[0, 0]$ 안에 있다. 실제로 주 열린집합 $D(f_r)$ 위의 층화 길이는
기껏해야 $t - 1$이다. 한편 $D(A/f_rA)$에서
$N' = N \otimes_A^\mathbf{L} A/f_rA$를 생각하자.
$S' = \Spec(A/f_rA)$로 놓으면 층화
$$
S' \supset S' \cap S_0 \supset S' \cap S_1 \supset \ldots \supset
S' \cap S_{t - 1} \supset
S' \cap S_t = \emptyset
$$
를 얻는다. 여기서 $S' \cap S_{t - 1}$은 $S'$ 안에서 $r - 1$개의
방정식으로 잘려 나온다. 앞에서와 마찬가지로 $M$을 각 부분
$S' \cap S_i \setminus S' \cap S_{i + 1}$로 유도 당긴 대상의 Tor
진폭은 $[0, 0]$ 안에 있다. 따라서 $r$에 대한 귀납법으로 $N'$의
Tor 진폭이 $[0, 0]$ 안에 있음을 얻는다. 마지막으로 More on Algebra,
Lemma \ref{more-algebra-lemma-f-flat}에 의해 $N$의 Tor 진폭도
$[0, 0]$ 안에 있다.
\end{proof}

\begin{lemma}
\label{more-morphisms-lemma-cokernel-flat}
$R$을 뇌터 정역이라 하고, $R \to A \to B$를 유한형 환 준동형이라
하자. $M$을 유한 $A$-가군이라 하고,
% P05-EN-CH37-0248
$N$을 유한 $B$-가군이라 하자. $M \to N$을 $A$-선형 사상이라 하자.
% P05-EN-CH37-0249
$M_f \to N_f$의 여핵이 평탄 $R_f$-가군이 되게 하는 0 아닌
$f \in R$이 존재한다.
\end{lemma}

\begin{proof}
$M$을 $M \to N$의 상으로 대체하여 $M \subset N$이라고 가정해도 된다.
어떤 소 아이디얼 $\mathfrak q_i \subset B$에 대하여
$N_i/N_{i - 1} = B/\mathfrak q_i$가 되게 하는 여과
$0 = N_0 \subset N_1 \subset \ldots \subset N_t = N$을 택하자. Algebra,
Lemma \ref{algebra-lemma-filter-Noetherian-module}를 보라.
$M_i = M \cap N_i$로 놓자. 그러면 $Q = N/M$은 부분가군
$Q_i = N_i/M_i$들에 의한 여과를 갖는다.
% P05-EN-CH37-0250
0 아닌 원소 $f$에서 국소화한 뒤 $Q_i/Q_{i - 1}$이 평탄해짐을 보이면
충분하다. 실제로 Algebra, Lemma \ref{algebra-lemma-flat-ses}에 의해
평탄 가군들의 확대는 평탄하다.
% P05-EN-CH37-0251
$Q_i/Q_{i - 1}$은 사상
$M_i/M_{i - 1} \to N_i/N_{i - 1}$의 여핵과 동형이므로, 다음 문단에서
논의하는 경우로 환원된다.

\medskip\noindent
$B$가 정역이고 $M \subset N = B$라고 하자. $A$를 $B$에서의 $A$의
상으로 대체하여 $A \subset B$라고 가정해도 된다. 일반 평탄성에 의해
$A$와 $B$가 $R$ 위에서 평탄이라고 가정해도 된다(Algebra, Lemma
\ref{algebra-lemma-generic-flatness-Noetherian}). 이제 $R$을 한 주원소
국소화로 대체한 뒤 $M \to B$가 $R$-보편 단사가 됨을 보이면 충분하다
(Algebra, Lemma \ref{algebra-lemma-ui-flat-domain}). 일반 자유성에 의해
$B_g$가 자유 $A_g$-가군이 되게 하는 0 아닌 $g \in A$를 찾을 수 있다
(Algebra, Lemma \ref{algebra-lemma-generic-flatness-Noetherian}).
따라서 $A_g$-가군으로서 직합인자 $M' \subset B_g$를 택할 수 있다.
이는 유한 자유 $A_g$-가군이고, $M \to B \to B_g$는 $M'$을 통해
인수분해된다. $R$을 한 주원소 국소화로 대체한 뒤 $M \to M'$이
$R$-보편 단사가 됨을 보이면 충분함은 명백하다.

\medskip\noindent
$M' = A_g^{\oplus n}$이라 하자. $M \subset M'$은 유한 $A$-가군이므로
어떤 $m \geq 0$에 대하여 $M$은 $(1/g^m)A^{\oplus n}$에 포함된다.
$M'$의 기저를 바꾸어 $M \subset A^{\oplus n}$이라고 가정해도 된다.
그러면 $R$을 한 주원소 국소화로 대체한 뒤 $A^{\oplus n}/M$과
$A_g/A$가 $R$-평탄이 됨을 보이면 충분하다. 실제로 이때 Algebra,
Lemma \ref{algebra-lemma-flat-tor-zero}에 의해
$M \to A^{\oplus n}$과 $A^{\oplus n} \to A_g^{\oplus n}$은 보편
단사이고, 따라서 그 합성 $M \to M' = A_g^{\oplus n}$도 보편 단사이다.

\medskip\noindent
일반 평탄성(위의 인용을 보라)에 의해 가군 $A^{\oplus n}/M$이
$R$-평탄이라고 가정해도 된다. 몫 $A_g/A$에 대해서는 다음 사실을
사용한다.
$$
A_g/A = \colim (1/g^m)A/A \cong \colim A/g^mA
$$
가군 $A/g^mA$는 길이 $m$인 여과를 가지며, 그 연속 몫들은
$A/gA$와 동형이다. 다시 일반 평탄성에 의해 $A/gA$가 $R$-평탄이라고
가정해도 된다. 따라서 각 $A/g^mA$는 $R$-평탄이고, 그러므로
$A_g/A$도 평탄이다.
\end{proof}

\noindent
$f : X \to Y$를 밑스킴 $S$ 위의 스킴들의 사상이라 하자.
$Z \subset Y$를 $f$의 스킴론적 상이라 하자. Morphisms, Section
\ref{morphisms-section-scheme-theoretic-image}를 보라.
$g : S' \to S$를 스킴의 사상이라 하고,
$f' : X \times_S S' \to Y \times_S S'$을 $g$에 의한 $f$의 밑변환이라
하자. $Z \times_S S' \subset Y \times_S S'$이 항상 $f'$의 스킴론적
상인 것은 아니다. 위와 같은 모든 $g$에 대하여 $f'$의 스킴론적 상이
$Z \times_S S'$과 같으면, {\it $f/S$의 스킴론적 상의 형성은 임의의
밑변환과 가환한다}고 하자.

\begin{lemma}
\label{more-morphisms-lemma-base-change-scheme-theoretic-image}
$S$를 준콤팩트 준분리 스킴이라 하자. $f : X \to Y$를 $S$ 위의
스킴들의 사상이라 하고, $X$와 $Y$가 모두 $S$ 위에서 유한 표시라고
하자. 그러면 어떤 $t \geq 0$과 닫힌 부분스킴들
$$
S \supset S_0 \supset S_1 \supset \ldots \supset S_t = \emptyset
$$
이 존재하여 다음 성질들을 갖는다.
\begin{enumerate}
\item $S_i \to S$는 유한형 아이디얼층으로 정의된다.
\item $S_0 \subset S$는 비후화이다.
\item $T_i = S_i \setminus S_{i + 1}$로 놓고 $f_i$를 $f$의
$T_i$로의 밑변환이라 하면, $f_i/T_i$의 스킴론적 상의 형성은 임의의
밑변환과 가환한다(보조정리 앞의 논의를 보라).
\end{enumerate}
\end{lemma}

\begin{proof}
각 정사각형이 Cartesian이고 $U$, $V$, $W$가 $\mathbf{Z}$ 위에서
유한형인 다음 가환 도식을 찾을 수 있다.
$$
\xymatrix{
X \ar[d] \ar[r] & Y \ar[d] \ar[r] & S \ar[d] \\
U \ar[r] & V \ar[r] & W
}
$$
실제로 먼저 Limits, Proposition \ref{limits-proposition-approximate}을
사용하여 $S$를 아핀 전이 사상들을 갖는 $\mathbf{Z}$ 위의 유한형
스킴들의 여과쌍대 극한으로 쓰고, 이어서 Limits, Lemma
\ref{limits-lemma-descend-finite-presentation}을 사용하여 사상
$X \to Y$를 강하시킨다. 이로써 다음 문단에서 논의하는 경우로
환원된다.

\medskip\noindent
$S$가 뇌터라고 하자. 이 경우 모든 준연접 아이디얼은 유한형이므로
$S_i$가 유한형 아이디얼로 잘려 나온다는 조건은 확인할 필요가 없다.
$S_0 = S_{red}$를 $S$의 축약으로 놓자. $\eta \in S_0$를 $S_0$의 한
기약성분의 일반점이라 하자. $S_0$의 바탕 위상공간에 대한 뇌터
귀납법으로, $\eta$를 포함하지 않는 $S_0$의 모든 닫힌 부분스킴에
대하여 결론이 성립한다고 가정해도 된다. 따라서 $f$의 $U_0$로의
밑변환 $f_0$가 성질 (3)을 갖게 하는 열린 근방 $U_0 \subset S_0$가
존재함을 보이면 충분하다.

\medskip\noindent
$R$을 뇌터 정역이라 하고, $f : X \to Y$를 $R$ 위의 유한형 스킴들의
사상이라 하자. 앞 문단의 논의에 의해 어떤 0 아닌 $g \in R$에 대하여
$R$을 $R_g$로, $X$, $Y$를 각각 $R_g$로의 밑변환으로 대체한 뒤
$f/R$의 스킴론적 상의 형성이 임의의 밑변환과 가환함을 보이면
충분하다.

\medskip\noindent
% P05-EN-CH37-0252
$Y = V_1 \cup \ldots V_n$을 아핀 열린 덮개라 하자.
$U_i = f^{-1}(V_i)$로 놓자. 각 $R$ 위의 사상 $U_i \to V_i$에 대하여
명제가 성립하면 $f$에 대해서도 성립한다. 실제로 Morphisms, Lemma
\ref{morphisms-lemma-quasi-compact-scheme-theoretic-image}에 의해
$U_i \to V_i$의 스킴론적 상은 $V_i$와 $f : X \to Y$의 스킴론적
상의 교집합이다. 따라서 $Y$가 아핀이라고 가정해도 된다.

\medskip\noindent
% P05-EN-CH37-0253
$X = U_1 \cup \ldots U_n$을 아핀 열린 덮개라 하자. 그러면
$X \to Y$의 스킴론적 상은
% P05-EN-CH37-0254
$\coprod U_i \to Y$의 스킴론적 상과 같다. 따라서 $X$가 아핀이라고
가정해도 된다.

\medskip\noindent
% P05-EN-CH37-0255
$X = \Spec(A)$, $Y = \Spec(B)$라 하고, $f$가 $R$-대수 준동형
$\varphi : A \to B$에 대응한다고 하자. 그러면 $f$의 스킴론적 상은
$\Spec(A/\Ker(\varphi))$이고, 밑변환 뒤에도 마찬가지이다(아핀 사상에
의한 밑변환에 대한 말이지만, 이들만 확인하면 충분하다). 따라서 모든
환 준동형 $R \to R'$에 대하여
$\Ker(\varphi \otimes_R R') = \Ker(\varphi) \otimes_R R'$이면
스킴론적 상의 형성은 밑변환과 가환한다.

\medskip\noindent
$R$의 적당한 0 아닌 $g$에 대하여 $R$, $A$, $B$를 각각 $R_g$, $A_g$,
$B_g$로 대체하여 $A$와 $B$가 $R$ 위에서 평탄이라고 가정해도 된다.
Lemma \ref{more-morphisms-lemma-cokernel-flat}에 의해 $B/A$도 평탄 $R$-가군이라고
가정해도 된다. 그러면
$0 \to \Ker(\varphi) \to A \to B \to B/A \to 0$은 평탄 $R$-가군들의
완전열이고, 이는 원하는 밑변환 명제를 함의한다.
\end{proof}






\section{사상의 층화}
\label{more-morphisms-section-stratifying-morphisms}

\noindent
$f : X \to S$를 준콤팩트 준분리 스킴들의 유한 표시 사상이라 하자.
Section \ref{more-morphisms-section-generic-flatness-stratification}에서 $X$가 각 층
위에서 평탄해지도록 $S$를 층화할 수 있음을 보았다.
% P05-EN-CH37-0256
이 절에서는 매끄러운 층들을 얻도록 $S$와 $X$를 모두 층화하고자 한다.
그대로는 완전히 성립하지 않아서 유한 국소 자유 사상에 의한 밑변환도
필요하다.

\begin{lemma}
\label{more-morphisms-lemma-smoothness-stratification-at-generic-point}
$f : X \to S$를 스킴들의 유한 표시 사상이라 하자.
$\eta \in S$를 $S$의 한 기약성분의 일반점이라 하자. $S$가 축약이라고
가정하자. 그러면 다음이 존재한다.
\begin{enumerate}
\item $\eta$를 포함하는 열린 부분스킴 $U \subset S$.
\item 전사이고 보편 단사이며 유한 국소 자유인 사상 $V \to U$.
\item 어떤 $t \geq 0$과 닫힌 부분스킴들
$$
X \times_S V \supset Z_0 \supset Z_1 \supset \ldots \supset Z_t = \emptyset
$$
로서, $Z_i \to X \times_S V$는 유한형 아이디얼층으로 정의되고,
$Z_0 \subset X \times_S V$는 비후화이며, 사상
$Z_i \setminus Z_{i + 1} \to V$는 매끄러운 것.
\end{enumerate}
\end{lemma}

\begin{proof}
$S$를 $\eta$의 한 열린 근방으로, $X$를 이 열린집합으로의 제한으로
대체해도 됨은 명백하다. 따라서 $S = \Spec(A)$이고 $A$가 축약환이며,
$\eta$가 극소 소 아이디얼 $\mathfrak p$에 대응한다고 가정해도 된다.
이 경우 국소환 $\mathcal{O}_{S, \eta} = A_\mathfrak p$는
$\kappa(\mathfrak p)$와 같음을 상기하자. Algebra, Lemma
\ref{algebra-lemma-minimal-prime-reduced-ring}를 보라.

\medskip\noindent
$k = \kappa(\eta)$ 위의 스킴 $X_\eta$에 Varieties, Lemma
\ref{varieties-lemma-smooth-stratification}을 적용한다. 이로부터 얻는
순비분리 체 확대를 $k'/k$로 나타내자. 다음 문단에서 $k' = k$인
경우로 환원한다. 이 단계는 보조정리의 명제에서 사상 $V \to U$를 찾는
것에 해당한다. 특히 $\kappa(\mathfrak p)$의 표수가 0이면 $V = U$로
둘 수 있다.

\medskip\noindent
$k = \kappa(\mathfrak p)$의 표수가 0이면 $k' = k$이다.
$k = \kappa(\mathfrak p)$의 표수가 $p > 0$이면, $p$는
$A_\mathfrak p = \kappa(\mathfrak p)$에서 0으로 간다. 따라서 $A$를
한 주원소 국소화로 대체하여(즉 $S$를 축소하여) $A$에서 $p = 0$이라고
가정해도 된다. $k' \not = k$이면,
% P05-EN-CH37-0257
$\beta \in k'$, $\beta \not \in k$, $\beta^p \in k$를 만족하는 원소가
존재한다. $A$를 한 주원소 국소화로 대체하여 $\beta^p = a$를 만족하는
$a \in A$가 존재한다고 가정해도 된다. $A' = A[x]/(x^p - a)$로 놓자.
그러면 $S' = \Spec(A') \to \Spec(A) = S$는 유한 국소 자유이고
전사이며 보편 단사이다. 또한 $\mathfrak p' \subset A'$를
$\mathfrak p$ 위에 놓인 유일한 소 아이디얼이라 하면,
$A'_{\mathfrak p'} = k(\beta)$이고 $k'/k(\beta)$의 차수는 더 작다.
따라서 $S$를 $S'$으로, $\eta$를 $\mathfrak p'$에 대응하는 점
$\eta'$으로 대체하면, $\eta$의 잉여체 위에서 $k'$의 차수가 감소한다.
이 과정을 계속하여 귀납법으로
$k' = \kappa(\mathfrak p) = \kappa(\eta)$인 경우로 환원된다.

\medskip\noindent
따라서 $S$가 아핀이고 축약이며, 어떤 $t \geq 0$과 닫힌 부분스킴들
$$
X_\eta \supset Z_{\eta, 0} \supset Z_{\eta, 1}
\supset \ldots \supset Z_{\eta, t} = \emptyset
$$
이 존재한다고 가정해도 된다. 여기서
$Z_{\eta, 0} = (X_\eta)_{red}$이고 모든 $i$에 대하여
$Z_{\eta, i} \setminus Z_{\eta, i + 1}$은 $\eta$ 위에서 매끄럽다.
$\kappa(\eta) = \kappa(\mathfrak p) = A_\mathfrak p$는
$a \in A$, $a \not \in \mathfrak p$에 대한 $A_a$들의 여과 귀납극한임을
상기하자. Algebra, Lemma \ref{algebra-lemma-localization-colimit}를
보라. 따라서 어떤 $a \in A$, $a \not \in \mathfrak p$에 대하여 위
도식을 $\Spec(A_a)$ 위의 대응하는 도식으로 강하시킬 수 있다.
더 정확히 말해 $S$를 $\Spec(A_a)$로 대체한 뒤, 어떤 $t \geq 0$과
닫힌 부분스킴들
$$
X \supset Z_0 \supset Z_1 \supset \ldots \supset Z_t = \emptyset
$$
이 존재한다고 가정해도 된다. 여기서 $Z_i \to X$는 유한 표시 닫힌
몰입이고, $Z_0 \to X$는 비후화이며,
$Z_i \setminus Z_{i + 1}$은 $S$ 위에서 매끄럽다. 다시 말해 보조정리가
성립한다. 더 정확히는 먼저 Limits, Lemma
\ref{limits-lemma-descend-finite-presentation}을 사용하여 $S$ 위의
유한 표시 사상들
$$
Z_t \to Z_{t - 1} \to \ldots \to Z_0 \to X
$$
을 얻는다. 이를 $\eta$로 밑변환하면 위에 주어진 닫힌 부분스킴들 사이의
포함들이 된다. $S$를 더 축소하여 각 사상이 닫힌 몰입이라고 가정해도
된다. Limits, Lemma
\ref{limits-lemma-descend-closed-immersion-finite-presentation}을 보라.
$S$를 축소하여 $Z_0 \to X$가 전사이고 따라서 비후화라고 가정해도
된다. Limits, Lemma \ref{limits-lemma-descend-surjective}를 보라.
$S$를 한 번 더 축소하여 $Z_i \setminus Z_{i + 1} \to S$가
매끄럽다고 가정해도 된다. Limits, Lemma
\ref{limits-lemma-descend-smooth}를 보라. 이로써 증명이 끝난다.
\end{proof}

\begin{lemma}
\label{more-morphisms-lemma-smoothness-stratification}
$f : X \to S$를 준콤팩트 준분리 스킴들 사이의 유한 표시 사상이라
하자. 그러면 어떤 $t \geq 0$과 닫힌 부분스킴들
$$
S \supset S_0 \supset S_1 \supset \ldots \supset S_t = \emptyset
$$
이 존재하여 다음을 만족한다.
\begin{enumerate}
\item $S_i \to S$는 유한형 아이디얼층으로 정의된다.
\item $S_0 \subset S$는 비후화이다.
\item 각 $i$에 대하여 전사 유한 국소 자유 사상
$T_i \to S_i \setminus S_{i + 1}$가 존재한다.
\item 각 $i$에 대하여 어떤 $t_i \geq 0$과 닫힌 부분스킴들
$$
X_i = X \times_S T_i \supset Z_{i, 0}
\supset Z_{i, 1} \supset \ldots \supset Z_{i, t_i} = \emptyset
$$
로서, $Z_{i, j} \to X_i$는 유한형 아이디얼층으로 정의되고,
$Z_{i, 0} \subset X_i$는 비후화이며, 사상
$Z_{i, j} \setminus Z_{i, j + 1} \to T_i$는 매끄러운 것.
\end{enumerate}
\end{lemma}

\begin{proof}
다음 Cartesian 도식을 찾을 수 있다.
$$
\xymatrix{
X \ar[d] \ar[r] & X_0 \ar[d] \\
S \ar[r] & S_0
}
$$
여기서 $X_0$와 $S_0$는 $\mathbf{Z}$ 위에서 유한형이다. Limits,
Proposition \ref{limits-proposition-approximate}과 Lemma
\ref{limits-lemma-descend-finite-presentation}을 보라. 따라서 $X$와
$S$가 $\mathbf{Z}$ 위에서 유한형이라고 가정해도 된다. 실제로 이
보조정리가 제기한 $X_0 \to S_0$에 대한 문제의 해를 밑변환하면
$S$ 위의 해가 된다. 세부사항은 생략한다.

\medskip\noindent
$X$와 $S$가 $\mathbf{Z}$ 위에서 유한형이라고 하자. 이 경우 모든
준연접 아이디얼은 유한형이므로 $S_i$가 유한형 아이디얼로 잘려
나온다는 조건은 확인할 필요가 없다. $S_0 = S_{red}$를 $S$의
축약으로 놓자. $\eta \in S_0$를 한 기약성분의 일반점이라 하자.
Lemma \ref{more-morphisms-lemma-smoothness-stratification-at-generic-point}에 의해
열린 부분스킴 $U \subset S_0$, 전사이고 보편 단사이며 유한 국소
자유인 사상 $V \to U$, 어떤 $t_0 \geq 0$, 그리고 닫힌 부분스킴들
$$
X \times_S V \supset Z_{0, 0} \supset Z_{0, 1} \supset \ldots \supset
Z_{0, t_0} = \emptyset
$$
을 찾을 수 있다. 여기서 $Z_{0, i} \to X \times_S V$는 유한형
아이디얼층으로 정의되고, $Z_{0, 0} \subset X \times_S V$는
비후화이며, 사상 $Z_{0, i} \setminus Z_{0, i + 1} \to V$는
매끄럽다. 이제 $S_1 \subset S_0$을 $S_0 \setminus U$ 위의 유도된
축약 부분스킴 구조라 하자. $S$의 바탕 위상공간에 대한 뇌터 귀납법으로
$X \times_S S_1 \to S_1$에 대하여 보조정리가 성립한다고 가정해도
된다. 이로부터 어떤 $t \geq 1$과
$$
S_1 = S_1 \supset S_2 \supset \ldots \supset S_t = \emptyset
$$
및 보조정리의 명제에서와 같은 $t_i$와 $Z_{i, j}$를 얻는다. 이로써
보조정리가 증명된다.
\end{proof}


















\section{상대차원 1인 사상의 개선}
\label{more-morphisms-section-make-good-curves}

\noindent
바탕체를 확대하고 콤팩트화하고 정규화하면 임의의 곡선을 매끄럽고
사영적으로 만들 수 있다. 이 사실은 일반올의 차원이 $1$인 유한형
사상들에 관한 결과도 함의한다.

\begin{lemma}
\label{more-morphisms-lemma-make-good-curves}
$f : X \to S$를 스킴의 사상이라 하고, $\eta \in S$를 $S$의 한
기약성분의 일반점이라 하자. $f$가 분리이고 유한 표시이며
$\dim(X_\eta) \leq 1$이라고 가정하자. 그러면 다음 가환 도식이
존재한다.
$$
\xymatrix{
\overline{Y}_1 \amalg \ldots \amalg \overline{Y}_n \ar[rd] &
Y_1 \amalg \ldots \amalg Y_n \ar[r]_-\nu \ar[d] \ar[l]^j &
X_V \ar[r] \ar[d] &
X_U \ar[r] \ar[d] &
X \ar[d]^f \\
& T_1 \amalg \ldots \amalg T_n \ar[r] &
V \ar[r] &
U \ar[r] &
S
}
$$
이 스킴들은 다음 성질들을 갖는다.
\begin{enumerate}
\item $U \subset S$는 $\eta$의 열린 근방이다.
\item $V \to U$는 유한이고 전사이며 보편 단사인 사상이다.
\item $X_U = U \times_S X$와 $X_V = V \times_S X$는 밑변환이다.
\item $\nu$는 유한이고 전사이며, 다음을 만족하는 열린집합
$W \subset X_V$가 존재한다.
\begin{enumerate}
\item $W$는 $X_V \to V$의 모든 올에서 조밀하다.
\item $\nu^{-1}(W) \cap Y_i$는 $Y_i \to T_i$의 모든 올에서 조밀하다.
\item $\nu^{-1}(W) \to W$는 비후화이다.
\end{enumerate}
\item $j$는 열린 몰입이다.
\item $T_i \to V$는 유한 에탈이다.
\item $Y_i \to T_i$는 전사이고 매끄럽다.
\item $\overline{Y}_i \to T_i$는 매끄럽고 고유이며, 그 올들은
기하적으로 연결이고 차원이 $\leq 1$이다.
\end{enumerate}
\end{lemma}

\begin{proof}
$S$를 $\eta$의 한 열린 근방으로, $X$를 이 열린집합으로의 제한으로
대체해도 됨은 명백하다. 또한 $S$를 그 축약으로, $X$를 이 축약으로의
밑변환으로 대체해도 된다. 따라서 $S = \Spec(A)$이고 $A$가 축약환이며
$\eta$가 극소 소 아이디얼 $\mathfrak p$에 대응한다고 가정해도 된다.
이 경우 국소환 $\mathcal{O}_{S, \eta} = A_\mathfrak p$는
$\kappa(\mathfrak p)$와 같음을 상기하자. Algebra, Lemma
\ref{algebra-lemma-minimal-prime-reduced-ring}를 보라.

\medskip\noindent
$k = \kappa(\eta)$ 위의 스킴 $X_\eta$에 Varieties, Lemma
\ref{varieties-lemma-dim-1-projective-completion-after-insep}를 적용한다.
이로부터 얻는 순비분리 체 확대를 $k'/k$로 나타내자. 다음 문단에서
$k' = k$인 경우로 환원한다. 이 단계는 보조정리의 명제에서 사상
$V \to U$를 찾는 것에 해당한다. 특히 $\kappa(\mathfrak p)$의 표수가
0이면 $V = U$로 둘 수 있다.

\medskip\noindent
$k = \kappa(\mathfrak p)$의 표수가 0이면 $k' = k$이다.
$k = \kappa(\mathfrak p)$의 표수가 $p > 0$이면, $p$는
$A_\mathfrak p = \kappa(\mathfrak p)$에서 0으로 간다. 따라서 $A$를
한 주원소 국소화로 대체하여(즉 $S$를 축소하여) $A$에서 $p = 0$이라고
가정해도 된다. $k' \not = k$이면,
% P05-EN-CH37-0258
$\beta \in k'$, $\beta \not \in k$, $\beta^p \in k$를 만족하는 원소가
존재한다. $A$를 한 주원소 국소화로 대체하여 $\beta^p = a$를 만족하는
$a \in A$가 존재한다고 가정해도 된다. $A' = A[x]/(x^p - a)$로 놓자.
그러면 $S' = \Spec(A') \to \Spec(A) = S$는 유한이고 전사이며 보편
단사이다. 또한 $\mathfrak p' \subset A'$를 $\mathfrak p$ 위에 놓인
유일한 소 아이디얼이라 하면, $A'_{\mathfrak p'} = k(\beta)$이고
$k'/k(\beta)$의 차수는 더 작다. 따라서 $S$를 $S'$으로, $\eta$를
$\mathfrak p'$에 대응하는 점 $\eta'$으로 대체하면 $\eta$의 잉여체
위에서 $k'$의 차수가 감소한다. 이 과정을 계속하여 귀납법으로
$k' = \kappa(\mathfrak p) = \kappa(\eta)$인 경우로 환원된다.

\medskip\noindent
따라서 $S$가 아핀이고 축약이며 다음 도식이 주어졌다고 가정해도 된다.
$$
\xymatrix{
\overline{Y}_{1, \eta} \amalg \ldots \amalg \overline{Y}_{n, \eta} \ar[rd] &
Y_{1, \eta} \amalg \ldots \amalg Y_{n, \eta} \ar[r]_-\nu \ar[d] \ar[l]^j &
X_\eta \ar[d] \\
& \Spec(k_1) \amalg \ldots \amalg \Spec(k_n) \ar[r] &
\eta
}
$$
이 스킴들은 다음 성질들을 갖는다.
\begin{enumerate}
\item $\nu$는 $X_\eta$의 정규화이다.
\item $j$는 상이 조밀한 열린 몰입이다.
\item $i = 1, \ldots, n$에 대하여 $k_i/\kappa(\eta)$는 유한 분리
확대이다.
\item $\overline{Y}_{i, \eta}$는 $k_i$ 위에서 매끄럽고 사영적이며
기하적으로 기약이고 차원이 $\leq 1$이다.
\end{enumerate}
$\kappa(\eta) = \kappa(\mathfrak p) = A_\mathfrak p$는
$a \in A$, $a \not \in \mathfrak p$에 대한 $A_a$들의 여과 귀납극한임을
상기하자. Algebra, Lemma \ref{algebra-lemma-localization-colimit}를
보라. 따라서 어떤 $a \in A$, $a \not \in \mathfrak p$에 대하여 위
도식을 $\Spec(A_a)$ 위의 대응하는 도식으로 강하시킬 수 있다.
더 정확히는 $S$를 $\Spec(A_a)$로 대체한 뒤 다음 가환 도식이
주어졌다고 가정해도 된다.
$$
\xymatrix{
\overline{Y}_1 \amalg \ldots \amalg \overline{Y}_n \ar[rd] &
Y_1 \amalg \ldots \amalg Y_n \ar[r]_-\nu \ar[d] \ar[l]^j &
X \ar[d] \\
& T_1 \amalg \ldots \amalg T_n \ar[r] & S
}
$$
이 도식을 $\eta$로 밑변환하면 위 도식이 되며, 다음 성질들이
성립한다.
\begin{enumerate}
\item $\nu$는 유한 전사 사상이다.
\item $j$는 열린 몰입이다.
\item $i = 1, \ldots, n$에 대하여 $T_i \to S$는 유한 에탈이다.
\item $Y_i \to T_i$는 매끄럽고 전사이다.
\item $\overline{Y}_i \to T_i$는 매끄럽고 고유이며, 그 올들은
기하적으로 연결이고 차원이 $\leq 1$이다.
\end{enumerate}
이를 위해 먼저 Limits, Lemma
\ref{limits-lemma-descend-finite-presentation}을 사용하여 밑변환하면
앞 도식이 되는 도식을 얻는다. 이어서 Limits, Lemmas
\ref{limits-lemma-descend-etale},
\ref{limits-lemma-descend-smooth},
\ref{limits-lemma-descend-finite-finite-presentation},
\ref{limits-lemma-limit-affine},
\ref{limits-lemma-descend-open-immersion},
\ref{limits-lemma-eventually-proper}, and
\ref{limits-lemma-descend-surjective}
을 사용한다. 이로써 $\nu$가 유한이고 전사이며, $j$가 열린 몰입이고,
$T_i \to S$가 유한 에탈이며,
% P05-EN-CH37-0259
$Y_i \to T$가 매끄럽고, $\overline{Y}_i \to T_i$가 고유이고
매끄러움을 얻는다. $Y_i$는 공집합일 수 없고, 매끄러운 사상은 열려
있으며, $T_i \to S$는 유한 에탈이므로 $S$를 축소한 뒤
$Y_i \to T_i$가 전사라고 가정해도 된다. 마지막으로 $T_i$에서
$\eta$ 위에 놓인 유일한 점 $\eta_i = \Spec(k_i)$ 위의
$\overline{Y}_i \to T_i$의 올은 기하적으로 연결이다. 따라서 한 번 더
축소하여 모든 올에 대해서도 마찬가지라고 가정해도 된다. Lemma
\ref{more-morphisms-lemma-proper-flat-geom-red}를 보라.

\medskip\noindent
이제 (a), (b), (c)를 만족하는 열린집합 $W \subset X$의 존재만
증명하면 된다. $\nu_\eta : \coprod Y_{i, \eta} \to X_\eta$는 정규화
사상이므로, Varieties, Lemma
\ref{varieties-lemma-normalization-locally-algebraic}에 의해 조밀한
열린집합 $W_\eta \subset X_\eta$가 존재하여,
$\nu^{-1}(W_\eta) \to W_\eta$는 $W_\eta$의 축약을 $W_\eta$에 넣는
포함과 같다. $W \subset X$를 $\eta$ 위의 올이 방금 찾은
$W_\eta$인 준콤팩트 열린집합이라 하자.
$A = \Gamma(S, \mathcal{O}_S)$를 또 다른 국소화로 대체하여
$\nu^{-1}(W) \to W$가 닫힌 몰입이라고 가정해도 된다. Limits, Lemma
\ref{limits-lemma-descend-closed-immersion-finite-presentation}을 보라.
$\nu$는 전사이기도 하므로 $\nu^{-1}(W) \to W$는 비후화이다.
$W_i = \nu^{-1}(W) \cap Y_i$로 놓자. $S$를 한 번 더 축소하여 모든
$i$에 대하여 $W_i \to T_i$가 전사라고 가정해도 된다(위와 같은
논증이다). $Y_i \to T_i$의 올들이 기하적으로 기약이므로
$W_i \subset Y_i$는 $Y_i \to T_i$의 모든 올에서 조밀하다.
$\nu$가 유한이고 전사이므로
$W = \nu(\nu^{-1}(W))$도 $X \to S$의 모든 올에서 조밀하다.
\end{proof}











\section{분리 국소 준유한 사상의 강하}
\label{more-morphisms-section-separated-locally-quasi-finite}

\noindent
이 절에서는 ``분리 국소 준유한 사상은 fppf-덮개에 대하여 강하를
만족한다''는 것을 보인다. 용어는 Descent, Definition
\ref{descent-definition-descending-types-morphisms}을 보라. 이 결과는
Raynaud와 Gruson의 여러 의미에서 훌륭한 논문
\cite[Lemma 5.7.1]{GruRay}의 증명 속에 숨어 있다. 사상이 국소 유한
표시라는 추가 가정 아래에서는 \cite{Murre-representation}과
\cite[Expos\'e X, Lemma 5.4]{SGA3}에서도 찾을 수 있다. 정식 명제는
다음과 같다.

\begin{lemma}
\label{more-morphisms-lemma-separated-locally-quasi-finite-morphisms-fppf-descend}
$S$를 스킴이라 하고, $\{X_i \to S\}_{i\in I}$를 fppf-덮개라 하자.
Topologies, Definition \ref{topologies-definition-fppf-covering}을 보라.
$(V_i/X_i, \varphi_{ij})$를 $\{X_i \to S\}$에 대한 강하 자료라 하자.
각 사상 $V_i \to X_i$가 분리이고 국소 준유한이면, 이 강하 자료는
유효이다.
\end{lemma}

\begin{proof}
분리성과 국소 준유한성은 임의의 밑변환으로 보존되는 스킴 사상의
성질이다. Schemes, Lemma \ref{schemes-lemma-separated-permanence}와
Morphisms, Lemma \ref{morphisms-lemma-base-change-quasi-finite}를 보라.
따라서 Descent, Lemma \ref{descent-lemma-descending-types-morphisms}을
적용할 수 있고, fppf-덮개가 아핀 스킴 사이의 하나의
% P05-EN-CH37-0260
평탄 전사 유한 표시 사상 $\{X \to S\}$로 주어지는 경우에 보조정리의
명제를 증명하면 충분하다. $X = \Spec(A)$와 $S = \Spec(R)$라 하자.
그러면 $R \to A$는 충실히 평탄한 환 준동형이다. $(V, \varphi)$를
$S$ 위의 $X$에 대한 강하 자료라 하고, $\pi : V \to X$가 분리이고
국소 준유한이라고 가정하자.

\medskip\noindent
$W^1 \subset V$를 임의의 아핀 열린집합이라 하자.
$W = \text{pr}_1(\varphi(W^1 \times_S X)) \subset V$를 생각하자.
다음 도식을 보자.
$$
\xymatrix{
W^1 \times_S X \ar[rrrrr] \ar[ddd] \ar[rd]
& & & & &
\varphi(W^1 \times_S X) \ar[ddd] \ar[ld] \\
& V \times_S X \ar[rrr]^\varphi \ar[rd] \ar[dd]
& & &
X \times_S V \ar[ld] \ar[dd] & \\
& &
X \times_S X \ar[r]^1 \ar[d]_{\text{pr}_0}
&
X \times_S X \ar[d]^{\text{pr}_1}
& & \\
W^1 \ar[r] &
V \ar[r] &
X &
X &
V \ar[l] &
W \ar[l]
}
$$
이제 $X \to S$는 평탄이고 유한 표시이므로 보편적으로 열린 사상이다
(Morphisms, Lemma \ref{morphisms-lemma-fppf-open}). 따라서 $W$는
열려 있다. 또한 $W$가 준콤팩트이고 $W^1 \subset W$임도 명백하다.
$\varphi$의 코사이클 조건에 의해
$\varphi(W \times_S X) = X \times_S W$임에 유의하자. 따라서
$\varphi$를 $W \times_S X$로 제한하여 새 강하 자료
$(W, \varphi')$를 얻는다. 사상 $W \to X$는 준콤팩트이고 분리이며
국소 준유한이다. 따라서 정의에 의해 분리 준유한 사상이다. 그러므로
Lemma \ref{more-morphisms-lemma-quasi-finite-separated-quasi-affine}에 의해
준아핀이다. 따라서 Descent, Lemma \ref{descent-lemma-quasi-affine}에
의해 강하 자료 $(W, \varphi')$는 유효이다.

\medskip\noindent
다시 말해 강하 사상 $\varphi$에 대하여 안정적인 준콤팩트 열린집합
$W_i$들로 이루어진 열린 덮개 $V = \bigcup W_i$가 존재한다.
또한 이와 같은 각 준콤팩트 열린집합 $W \subset V$에 대응하는 강하
% P05-EN-CH37-0261
자료 $(W, \varphi')$는 유효이다. 따라서 열린집합 $W_i$들을 강하시켜
얻은 스킴들을 붙이면 원래 강하 자료가 유효임을 알 수 있다. Descent,
Lemma \ref{descent-lemma-effective-for-fpqc-is-local-upstairs}를 보라.
\end{proof}






\section{상대 유한 표시}
\label{more-morphisms-section-finite-type-finite-presentation}

\noindent
$R \to A$를 유한형 환 준동형이라 하고 $M$을 $A$-가군이라 하자. More
on Algebra, Section
\ref{more-algebra-section-relative-finite-presentation}에서 $M$이
$R$에 상대적으로 유한 표시라는 것의 의미를 정의했다. 또한 이 개념이
좋은 국소화 성질을 가지며 붙여짐을 증명했다. 따라서 대응하는 대역
개념을 다음과 같이 정의할 수 있다.

\begin{definition}
\label{more-morphisms-definition-relatively-finitely-presented-sheaf}
$f : X \to S$를 국소 유한형인 스킴의 사상이라 하고, $\mathcal{F}$를
준연접 $\mathcal{O}_X$-가군이라 하자. 아핀 열린 덮개
$S = \bigcup V_i$와 각 $i$에 대한 아핀 열린 덮개
$f^{-1}(V_i) = \bigcup_j U_{ij}$가 존재하여, $\mathcal{F}(U_{ij})$가
$\mathcal{O}_S(V_i)$에 상대적으로 유한 표시인
$\mathcal{O}_X(U_{ij})$-가군이면, $\mathcal{F}$가
{\it $S$에 상대적으로 유한 표시} 또는 {\it $S$에 상대적으로 유한
표시형}이라고 한다.
\end{definition}

\noindent
이 조건은 $\mathcal{F}$가 유한형 $\mathcal{O}_X$-가군임을 함의한다.
$X \to S$가 단지 국소 유한형이면, $X \to S$가 국소 유한 표시가
아니면서도 $\mathcal{F}$는 $S$에 상대적으로 유한 표시일 수 있다.
$X \to S$가 국소 유한 표시일 필요충분조건은 $\mathcal{O}_X$가
$S$에 상대적으로 유한 표시인 것임을 보게 될 것이다.

\begin{lemma}
\label{more-morphisms-lemma-relative-finite-presentation-characterize}
$f : X \to S$를 국소 유한형인 스킴의 사상이라 하고, $\mathcal{F}$를
준연접 $\mathcal{O}_X$-가군이라 하자. 다음은 서로 동치이다.
\begin{enumerate}
\item $\mathcal{F}$는 $S$에 상대적으로 유한 표시이다.
\item % P05-EN-CH37-0262
$f(U) \subset V$인 모든 아핀 열린집합 $U \subset X$, $V \subset S$에
대하여, $\mathcal{O}_X(U)$-가군 $\mathcal{F}(U)$는
$\mathcal{O}_S(V)$에 상대적으로 유한 표시이다.
\end{enumerate}
또한 이 조건들이 성립하면,
% P05-EN-CH37-0263
$f(U) \subset V$인 모든 열린 부분스킴 $U \subset X$와 $V \subset S$에
대하여 제한 $\mathcal{F}|_U$는 $V$에 상대적으로 유한 표시이다.
\end{lemma}

\begin{proof}
마지막 명제는 (1)과 (2)의 동치에서 명백하다. (2)가 (1)을 함의함도
명백하다. (1)을 가정하자. Definition
\ref{more-morphisms-definition-relatively-finitely-presented-sheaf}에서와 같은 아핀
열린 덮개들을 $S = \bigcup V_i$와
$f^{-1}(V_i) = \bigcup U_{ij}$라 하자. $U \subset X$와 $V \subset S$를
(2)에서와 같이 택하자. More on Algebra, Lemma
\ref{more-algebra-lemma-glue-relative-finite-presentation}에 의해
$\mathcal{F}(U_k)$가 $\mathcal{O}_S(V)$에 상대적으로 유한 표시가 되게
하는 $U$의 표준 열린 덮개 $U = \bigcup U_k$를 찾으면 충분하다.
다시 말해 모든 $u \in U$에 대하여, $\mathcal{F}(U')$가
$\mathcal{O}_S(V)$에 상대적으로 유한 표시가 되게 하는 표준 아핀
열린집합 $u \in U' \subset U$를 찾으면 충분하다. $f(u) \in V_i$인
$i$를 택하고, 이어서 $u \in U_{ij}$인 $j$를 택하자. Schemes, Lemma
\ref{schemes-lemma-standard-open-two-affines}에 의해
% P05-EN-CH37-0264
$V'$와 $V_i$에서 표준 아핀 열린집합인
$v \in V' \subset V \cap V_i$를 찾을 수 있다. 그러면 각각
$f^{-1}V'  \cap U$와 $f^{-1}V' \cap U_{ij}$는 각각 $U$와 $U_{ij}$의
표준 아핀 열린집합이다. 이 보조정리를 다시 적용하여
$f^{-1}V'  \cap U$와 $f^{-1}V' \cap U_{ij}$ 모두에서 표준 아핀
열린집합인 $u \in U' \subset f^{-1}V' \cap U \cap U_{ij}$를 찾을 수
있다. 따라서 $U'$은 $U$와 $U_{ij}$ 모두의 표준 아핀 열린집합이다.
More on Algebra, Lemma
\ref{more-algebra-lemma-localize-relative-finite-presentation}에 의해,
$\mathcal{F}(U_{ij})$가 $\mathcal{O}_S(V_i)$에 상대적으로 유한
표시라는 가정은 $\mathcal{F}(U')$가 $\mathcal{O}_S(V_i)$에
상대적으로 유한 표시임을 함의한다.
$\mathcal{O}_X(U') =
\mathcal{O}_X(U') \otimes_{\mathcal{O}_S(V_i)} \mathcal{O}_S(V')
$
이므로 More on Algebra, Lemma
\ref{more-algebra-lemma-base-change-relative-finite-presentation}에 의해
$\mathcal{F}(U')$는 $\mathcal{O}_S(V')$에 상대적으로 유한 표시이다.
More on Algebra, Lemma
\ref{more-algebra-lemma-localize-relative-finite-presentation}을 다시
적용하면 $\mathcal{F}(U')$는 $\mathcal{O}_S(V)$에 상대적으로 유한
표시이다. 이로써 증명이 끝난다.
\end{proof}

\begin{lemma}
\label{more-morphisms-lemma-relative-finite-presentation}
$f : X \to S$를 국소 유한형인 스킴의 사상이라 하고, $\mathcal{F}$를
준연접 $\mathcal{O}_X$-가군이라 하자.
\begin{enumerate}
\item $f$가 국소 유한 표시이면, $\mathcal{F}$가 $S$에 상대적으로
유한 표시일 필요충분조건은 $\mathcal{F}$가 유한 표시인 것이다.
\item 사상 $f$가 국소 유한 표시일 필요충분조건은 $\mathcal{O}_X$가
$S$에 상대적으로 유한 표시인 것이다.
\end{enumerate}
\end{lemma}

\begin{proof}
정의에서 곧바로 따른다. More on Algebra, Definition
\ref{more-algebra-definition-relatively-finitely-presented} 뒤의 논의를
보라.
\end{proof}

\begin{lemma}
\label{more-morphisms-lemma-finite-morphism-relative-finite-presentation}
$\pi : X \to Y$를 밑스킴 $S$ 위에서 국소 유한형인 스킴들의 유한
사상이라 하고, $\mathcal{F}$를 준연접 $\mathcal{O}_X$-가군이라 하자.
그러면 $\mathcal{F}$가 $S$에 상대적으로 유한 표시일 필요충분조건은
$\pi_*\mathcal{F}$가 $S$에 상대적으로 유한 표시인 것이다.
\end{lemma}

\begin{proof}
More on Algebra, Lemma \ref{more-algebra-lemma-finite-extension}의 결과를
스킴의 언어로 옮긴 것이다.
\end{proof}

\begin{lemma}
\label{more-morphisms-lemma-base-change-relative-finite-presentation}
$f : X \to S$를 국소 유한형인 스킴의 사상이라 하고, $\mathcal{F}$를
준연접 $\mathcal{O}_X$-가군이라 하자. $S' \to S$를 스킴의 사상이라
하고, $X' = X \times_S S'$로 놓으며,
% P05-EN-CH37-0265
$\mathcal{F}'$로 $\mathcal{F}$를 $X'$으로 당긴 가군을 나타내자.
$\mathcal{F}$가 $S$에 상대적으로 유한 표시이면, $\mathcal{F}'$는
$S'$에 상대적으로 유한 표시이다.
\end{lemma}

\begin{proof}
More on Algebra, Lemma
\ref{more-algebra-lemma-base-change-relative-finite-presentation}의 결과를
스킴의 언어로 옮긴 것이다.
\end{proof}

\begin{lemma}
\label{more-morphisms-lemma-pull-relative-finite-presentation}
$X \to Y \to S$를 국소 유한형인 스킴의 사상들이라 하고,
$\mathcal{G}$를 준연접 $\mathcal{O}_Y$-가군이라 하자.
$f : X \to Y$가 국소 유한 표시이고,
% P05-EN-CH37-0266
$\mathcal{G}$가 $S$에 상대적으로 유한 표시이면,
$f^*\mathcal{G}$는 $S$에 상대적으로 유한 표시이다.
\end{lemma}

\begin{proof}
More on Algebra, Lemma
\ref{more-algebra-lemma-pull-relative-finite-presentation}의 결과를
스킴의 언어로 옮긴 것이다.
\end{proof}

\begin{lemma}
\label{more-morphisms-lemma-composition-relative-finite-presentation}
$X \to Y \to S$를 국소 유한형인 스킴의 사상들이라 하고,
$\mathcal{F}$를 준연접 $\mathcal{O}_X$-가군이라 하자.
$Y \to S$가 국소 유한 표시이고 $\mathcal{F}$가 $Y$에 상대적으로
유한 표시이면, $\mathcal{F}$는 $S$에 상대적으로 유한 표시이다.
\end{lemma}

\begin{proof}
More on Algebra, Lemma
\ref{more-algebra-lemma-composition-relative-finite-presentation}의 결과를
스킴의 언어로 옮긴 것이다.
\end{proof}

\begin{lemma}
\label{more-morphisms-lemma-ses-relatively-finite-presentation}
$X \to S$를 국소 유한형인 스킴의 사상이라 하자.
$0 \to \mathcal{F}' \to \mathcal{F} \to \mathcal{F}'' \to 0$을 준연접
$\mathcal{O}_X$-가군들의 짧은 완전열이라 하자.
\begin{enumerate}
\item $\mathcal{F}', \mathcal{F}''$가 $S$에 상대적으로 유한 표시이면
$\mathcal{F}$도 그러하다.
\item $\mathcal{F}'$가 유한형 $\mathcal{O}_X$-가군이고
$\mathcal{F}$가 $S$에 상대적으로 유한 표시이면, $\mathcal{F}''$도
$S$에 상대적으로 유한 표시이다.
\end{enumerate}
\end{lemma}

\begin{proof}
More on Algebra, Lemma
\ref{more-algebra-lemma-ses-relatively-finite-presentation}의 결과를
스킴의 언어로 옮긴 것이다.
\end{proof}

\begin{lemma}
\label{more-morphisms-lemma-sum-relatively-finite-presentation}
\begin{slogan}
가군의 직합인자는 밑에 상대적으로 유한 표시라는 성질을 물려받는다.
\end{slogan}
$X \to S$를 국소 유한형인 스킴의 사상이라 하고,
$\mathcal{F}, \mathcal{F}'$를 준연접 $\mathcal{O}_X$-가군들이라 하자.
$\mathcal{F} \oplus \mathcal{F}'$가 $S$에 상대적으로 유한 표시이면,
$\mathcal{F}$와 $\mathcal{F}'$도 그러하다.
\end{lemma}

\begin{proof}
More on Algebra, Lemma
\ref{more-algebra-lemma-sum-relatively-finite-presentation}의 결과를
스킴의 언어로 옮긴 것이다.
\end{proof}










\section{상대 유사연접성}
\label{more-morphisms-section-relative-pseudo-coherence}

\noindent
이 절은 스킴에 대한
More on Algebra, Section \ref{more-algebra-section-relative-pseudo-coherent}의
대응물이다. 독자에게 먼저 그 절을 살펴볼 것을 강력히 권한다.
이 절의 내용과 유사연접 복합체에 관한 내용을
Cohomology, Sections \ref{cohomology-section-strictly-perfect},
\ref{cohomology-section-pseudo-coherent},
\ref{cohomology-section-tor}, and
\ref{cohomology-section-perfect}
에서 임의의 $\mathcal{O}_X$-가군 복합체에 대해 전개했지만,
$X$가 스킴이면 $D_\QCoh(\mathcal{O}_X)$의 대상만을 다루는 것이
% P05-EN-CH37-0267
많은 보조정리와 논증을 크게 단순화하며, 흔히 당면한 문제를 곧바로
그 대수적 대응물로 환원한다.
% P05-EN-CH37-0268
더욱이 우리가 가장 먼저 보이는 사실 가운데 하나는 상대 유사연접성이
코호몰로지 층들의 준연접성을 함의한다는 것이다. Lemma
\ref{more-morphisms-lemma-relative-pseudo-coherence}를 보라.
따라서 처음 읽을 때에는 $D_\QCoh(\mathcal{O}_X)$의 대상만을
다룰 것을 권한다.

\begin{lemma}
\label{more-morphisms-lemma-relatively-pseudo-coherent}
$X \to S$를 아핀 스킴의 유한형 사상이라 하자.
$E$를 $D(\mathcal{O}_X)$의 대상으로 하고,
$m \in \mathbf{Z}$라 하자.
다음 조건들은 서로 동치이다.
\begin{enumerate}
\item 어떤 닫힌 몰입 $i : X \to \mathbf{A}^n_S$에 대해
$D(\mathcal{O}_{\mathbf{A}^n_S})$의 대상 $Ri_*E$가
$m$-유사연접이고,
\item 모든 닫힌 몰입 $i : X \to \mathbf{A}^n_S$에 대해
$D(\mathcal{O}_{\mathbf{A}^n_S})$의 대상 $Ri_*E$가
$m$-유사연접이다.
\end{enumerate}
\end{lemma}

\begin{proof}
$S = \Spec(R)$ 및 $X = \Spec(A)$라 하자. $i$가 전사
$\alpha : R[x_1, \ldots, x_n] \to A$에 대응하고,
% P05-EN-CH37-0269
$X \to \mathbf{A}^m_S$가 $\beta : R[y_1, \ldots, y_m] \to A$에
대응한다고 하자. $\alpha(f_j) = \beta(y_j)$인
$f_j \in R[x_1, \ldots, x_n]$와 $\beta(g_i) = \alpha(x_i)$인
$g_i \in R[y_1, \ldots, y_m]$를 택한다. 그러면 가환 그림
$$
\xymatrix{
R[x_1, \ldots, x_n, y_1, \ldots, y_m]
\ar[d]^{x_i \mapsto g_i} \ar[rr]_-{y_j \mapsto f_j} & &
R[x_1, \ldots, x_n] \ar[d] \\
R[y_1, \ldots, y_m] \ar[rr] & & A
}
$$
을 얻으며, 이는 닫힌 몰입의 가환 그림
$$
\xymatrix{
\mathbf{A}^{n + m}_S & \mathbf{A}^n_S \ar[l] \\
\mathbf{A}^m_S \ar[u] & X \ar[u] \ar[l]
}
$$
에 대응한다. 따라서 닫힌 몰입
$$
f : \mathbf{A}^m_S \to \mathbf{A}^{n + m}_S
$$
에 대해 $D(\mathcal{O}_{\mathbf{A}^m_S})$의 대상 $E$가
$m$-유사연접일 필요충분조건이 $Rf_*E$가 $m$-유사연접인 것임을
보이면 충분하다. 이는 Derived Categories of Schemes, Lemma
\ref{perfect-lemma-closed-push-pseudo-coherent}
와 $f_*\mathcal{O}_{\mathbf{A}^m_S}$가 유사연접
$\mathcal{O}_{\mathbf{A}^{n + m}_S}$-가군이라는 사실에서 따른다.
$f_*\mathcal{O}_{\mathbf{A}^m_S}$의 유사연접성을 직접 보이는 것은
어렵지 않지만, Derived Categories of Schemes, Lemma
\ref{perfect-lemma-pseudo-coherent-affine} 및 More on Algebra, Lemma
\ref{more-algebra-lemma-relatively-pseudo-coherent}에서도 따른다.
\end{proof}

\noindent
% P05-EN-CH37-0270
$f : X \to S$가 국소 유한형인 스킴의 사상이면,
$f(U) \subset V$를 만족하는 모든 아핀 열린집합 쌍 $U \subset X$,
$V \subset S$에 대해 환 준동형
$\mathcal{O}_S(V) \to \mathcal{O}_X(U)$가 유한형임을 상기하자
(Morphisms, Lemma \ref{morphisms-lemma-locally-finite-type-characterize}).
따라서 닫힌 몰입 $U \to \mathbf{A}^n_V$가 항상 존재하므로
다음 정의는 의미가 있다.

\begin{definition}
\label{more-morphisms-definition-relative-pseudo-coherence}
$f : X \to S$를 국소 유한형인 스킴의 사상이라 하자.
$E$를 $D(\mathcal{O}_X)$의 대상으로 하고, $\mathcal{F}$를
$\mathcal{O}_X$-가군이라 하자. $m \in \mathbf{Z}$를 고정한다.
\begin{enumerate}
\item 아핀 열린 덮개 $S = \bigcup V_i$와, 각 $i$마다 아핀 열린 덮개
$f^{-1}(V_i) = \bigcup U_{ij}$가 존재하여 모든 쌍
$(U_{ij} \to V_i, E|_{U_{ij}})$에 대해 Lemma
\ref{more-morphisms-lemma-relatively-pseudo-coherent}의 동치 조건들이 성립하면,
$E$가 {\it $S$에 상대적으로 $m$-유사연접}이라고 한다.
\item 모든 $m \in \mathbf{Z}$에 대해 $E$가 $S$에 상대적으로
$m$-유사연접이면, $E$가 {\it $S$에 상대적으로 유사연접}이라고 한다.
\item $D(\mathcal{O}_X)$의 대상으로 본 $\mathcal{F}$가 $S$에
상대적으로 $m$-유사연접이면, $\mathcal{F}$가 {\it $S$에 상대적으로
$m$-유사연접}이라고 한다.
\item $D(\mathcal{O}_X)$의 대상으로 본 $\mathcal{F}$가 $S$에
상대적으로 유사연접이면, $\mathcal{F}$가 {\it $S$에 상대적으로
유사연접}이라고 한다.
\end{enumerate}
\end{definition}

\noindent
$X$가 준콤팩트이고 어떤 $m$에 대해 $E$가 $S$에 상대적으로
$m$-유사연접이면 $E$는 위로 유계이다. $E$가 $S$에 상대적으로
유사연접이면 $E$의 코호몰로지 층들은 준연접이다.

\begin{lemma}
\label{more-morphisms-lemma-relative-pseudo-coherence}
$f : X \to S$를 국소 유한형인 스킴의 사상이라 하자.
$D(\mathcal{O}_X)$의 $E$가 $S$에 상대적으로 $m$-유사연접이면,
$i > m$에 대해 $H^i(E)$는 준연접 $\mathcal{O}_X$-가군이다.
$E$가 $S$에 상대적으로 유사연접이면 $E$는
$D_\QCoh(\mathcal{O}_X)$의 대상이다.
\end{lemma}

\begin{proof}
아핀 열린 덮개 $S = \bigcup V_i$와, 각 $i$마다 아핀 열린 덮개
$f^{-1}(V_i) = \bigcup U_{ij}$를 택하여 모든 쌍
$(U_{ij} \to V_i, E|_{U_{ij}})$에 대해 Lemma
\ref{more-morphisms-lemma-relatively-pseudo-coherent}의 동치 조건들이 성립하게 한다.
준연접성은 $X$ 위에서 국소적이므로, 닫힌 몰입
% P05-EN-CH37-0271
$i : X \to \mathbf{A}^n_S$가 존재하여 $Ri_*E$가
$\mathbf{A}^n_S$ 위에서 $m$-유사연접이라고 가정할 수 있다.
Derived Categories of Schemes, Lemma \ref{perfect-lemma-pseudo-coherent}에
의하면 이는 $q > m$에 대해 $H^q(Ri_*E)$가 준연접이라는 뜻이다.
$i_*$는 완전함자이므로 $i_*H^q(E) = H^q(Ri_*E)$이고, 이는
$\mathbf{A}^n_S$ 위에서 준연접이다. Morphisms, Lemma
\ref{morphisms-lemma-i-star-equivalence}에 의해 원하는 대로
$H^q(E)$가 준연접이다. 엄밀히 말하면 이 보조정리는 준연접
$\mathcal{O}_X$-가군 $\mathcal{F}$가 존재하여
$i_*\mathcal{F} = i_*H^q(E)$가 됨을 뜻하고, 이어 Modules, Lemma
\ref{modules-lemma-i-star-equivalence}가
$\mathcal{F} \cong H^q(E)$임을 주므로 결론이 따른다.
\end{proof}

\noindent
다음으로 상대 유사연접성 조건이 국소화와 잘 양립함을 보인다.

\begin{lemma}
\label{more-morphisms-lemma-localize-relative-pseudo-coherent}
$S$를 아핀 스킴이라 하고 $V \subset S$를 표준 열린집합이라 하자.
$X \to V$를 아핀 스킴의 유한형 사상이라 하자.
$U \subset X$를 아핀 열린집합이라 하고 $E$를
$D(\mathcal{O}_X)$의 대상으로 하자. 쌍 $(X \to V, E)$에 대해
Lemma \ref{more-morphisms-lemma-relatively-pseudo-coherent}의 동치 조건들이 성립하면,
쌍 $(U \to S, E|_U)$에 대해서도 Lemma
\ref{more-morphisms-lemma-relatively-pseudo-coherent}의 동치 조건들이 성립한다.
\end{lemma}

\begin{proof}
$S = \Spec(R)$, $V = D(f)$, $X = \Spec(A)$, $U = D(g)$라 쓰자.
쌍 $(X \to V, E)$에 대해 Lemma
\ref{more-morphisms-lemma-relatively-pseudo-coherent}의 동치 조건들이 성립한다고 하자.

\medskip\noindent
전사 $R_f[x_1, \ldots, x_n] \to A$를 택한다.
$R_f = R[x_0]/(fx_0 - 1)$라 쓰자. 그러면
$R[x_0, x_1, \ldots, x_n] \to A$는 전사이고,
$R_f[x_1, \ldots, x_n]$는 $R[x_0, \ldots, x_n]$-가군으로서
유사연접이다. 따라서
$$
X \to \mathbf{A}^n_V \to \mathbf{A}^{n + 1}_S
$$
를 얻는다. Derived Categories of Schemes, Lemma
\ref{perfect-lemma-closed-push-pseudo-coherent}를 적용하면 $E$를
$\mathbf{A}^{n + 1}_S$로 전진시킨 $E'$가 $m$-유사연접임을 알 수 있다.

\medskip\noindent
$g \in A$로 가는 원소 $g' \in R[x_0, x_1, \ldots, x_n]$를 택한다.
전사
$R[x_0, \ldots, x_{n + 1}] \to R[x_0, \ldots, x_n, 1/g']$를
생각하면 다음 그림을 얻는다.
$$
\xymatrix{
X \ar[d] & U \ar[d] \ar[l] \ar[dr] \\
\mathbf{A}^{n + 1}_S & D(g')\ar[l] \ar[r] & \mathbf{A}^{n + 2}_S
}
$$
여기서 아래 왼쪽 화살표는 열린 몰입이고 아래 오른쪽 화살표는 닫힌 몰입이다.
앞과 같이 $E'|_{D(g')}$를 $\mathbf{A}^{n + 2}_S$로 전진시킨 대상이
$m$-유사연접임을 알 수 있다. 이것은 $E|_U$를
$\mathbf{A}^{n + 2}_S$로 전진시킨 대상이기도 하므로 보조정리가 따른다.
\end{proof}

\begin{lemma}
\label{more-morphisms-lemma-glue-relative-pseudo-coherent}
$X \to S$를 아핀 스킴의 유한형 사상이라 하고, $E$를
$D(\mathcal{O}_X)$의 대상으로 하자. $m \in \mathbf{Z}$라 하자.
$X = \bigcup U_i$를 표준 아핀 열린 덮개라 하자.
다음 조건들은 서로 동치이다.
\begin{enumerate}
\item 쌍들 $(U_i \to S, E|_{U_i})$에 대해 Lemma
\ref{more-morphisms-lemma-relatively-pseudo-coherent}의 동치 조건들이 성립한다.
\item 쌍 $(X \to S, E)$에 대해 Lemma
\ref{more-morphisms-lemma-relatively-pseudo-coherent}의 동치 조건들이 성립한다.
\end{enumerate}
\end{lemma}

\begin{proof}
(2) $\Rightarrow$ (1)은 Lemma
\ref{more-morphisms-lemma-localize-relative-pseudo-coherent}이다.
(1)을 가정하자. $S = \Spec(R)$, $X = \Spec(A)$,
$U_i = D(f_i)$라 하고, $A$에서 $1 = \sum f_ig_i$라 쓰자.
전사들
$$
R[x_i, y_i, z_i] \to R[x_i, y_i, z_i]/(\sum y_iz_i - 1) \to A.
$$
을 생각하되, $y_i$는 $f_i$로, $z_i$는 $g_i$로 보낸다.
$R[x_i, y_i, z_i]/(\sum y_iz_i - 1)$는
$R[x_i, y_i, z_i]$-가군으로서 유사연접이다. 따라서 Derived Categories
of Schemes, Lemma \ref{perfect-lemma-closed-push-pseudo-coherent}에 의해
$E$를 $T = \Spec(R[x_i, y_i, z_i]/(\sum y_iz_i - 1))$로 전진시킨
대상이 $m$-유사연접임을 보이면 충분하다. 각 $i_0$에 대해 이 전진의
$W_{i_0} = \Spec(R[x_i, y_i, z_i, 1/y_{i_0}]/(\sum y_iz_i - 1))$
로의 제한이 $m$-유사연접임을 보이면 충분하다. 가환 그림
$$
\xymatrix{
X \ar[d] & U_{i_0} \ar[l] \ar[d] \\
T & W_{i_0} \ar[l]
}
$$
이 있으며, 이 그림은 $E$를 $T$로 전진시킨 뒤 $W_{i_0}$로 제한한
대상이 $E|_{U_{i_0}}$를 $W_{i_0}$로 전진시킨 대상임을 뜻한다.
$R[x_i, y_i, z_i, 1/y_{i_0}]/(\sum y_iz_i - 1)$는 $R$ 위의
다항식환과 동형이므로 원하는 결론을 얻는다.
\end{proof}

\begin{lemma}
\label{more-morphisms-lemma-relative-pseudo-coherence-characterize}
$f : X \to S$를 국소 유한형인 스킴의 사상이라 하자.
$E$를 $D(\mathcal{O}_X)$의 대상으로 하고 $m \in \mathbf{Z}$를
고정한다. 다음 조건들은 서로 동치이다.
\begin{enumerate}
\item $E$가 $S$에 상대적으로 $m$-유사연접이다.
% P05-EN-CH37-0272
\item $f(U) \subset V$인 모든 아핀 열린집합 쌍 $U \subset X$,
$V \subset S$에 대해 쌍 $(U \to V, E|_U)$가 Lemma
\ref{more-morphisms-lemma-relatively-pseudo-coherent}의 동치 조건들을 만족한다.
\end{enumerate}
더욱이 이 조건들이 성립하면,
% P05-EN-CH37-0273
$f(U) \subset V$인 모든 열린 부분스킴 쌍 $U \subset X$,
$V \subset S$에 대해 제한 $E|_U$는 $V$에 상대적으로
$m$-유사연접이다.
\end{lemma}

\begin{proof}
마지막 명제는 (1)과 (2)의 동치에서 명백하다. (2)가 (1)을 함의하는
것도 명백하다. (1)을 가정하자. Definition
\ref{more-morphisms-definition-relative-pseudo-coherence}에서처럼
$S = \bigcup V_i$와 $f^{-1}(V_i) = \bigcup U_{ij}$를 아핀 열린
덮개로 두자. (2)와 같은 $U \subset X$, $V \subset S$를 택한다.
Lemma \ref{more-morphisms-lemma-glue-relative-pseudo-coherent}에 의해, 쌍들
$(U_k \to V, E|_{U_k})$에 대해 Lemma
\ref{more-morphisms-lemma-relatively-pseudo-coherent}의 동치 조건들이 성립하는
$U$의 표준 열린 덮개 $U = \bigcup U_k$를 찾으면 충분하다.
다시 말해 모든 $u \in U$에 대해, 쌍 $(U' \to V, E|_{U'})$가 Lemma
\ref{more-morphisms-lemma-relatively-pseudo-coherent}의 동치 조건들을 만족하는 표준
아핀 열린집합 $u \in U' \subset U$를 찾으면 충분하다.
$f(u) \in V_i$인 $i$를 택하고, 이어 $u \in U_{ij}$인 $j$를 택한다.
Schemes, Lemma \ref{schemes-lemma-standard-open-two-affines}에 의해
% P05-EN-CH37-0274
$V'$와 $V_i$ 안에서 표준 아핀 열린집합인
$v \in V' \subset V \cap V_i$를 찾을 수 있다. 그러면 각각
$f^{-1}V'  \cap U$와 $f^{-1}V' \cap U_{ij}$는 각각 $U$와
$U_{ij}$의 표준 아핀 열린집합이다. 이 보조정리를 다시 적용하면,
$f^{-1}V'  \cap U$와 $f^{-1}V' \cap U_{ij}$ 둘 다에서 표준 아핀
열린집합인 $u \in U' \subset f^{-1}V' \cap U \cap U_{ij}$를 찾을 수
있다. 따라서 $U'$는 $U$와 $U_{ij}$의 표준 아핀 열린집합이기도 하다.
Lemma \ref{more-morphisms-lemma-localize-relative-pseudo-coherent}에 의해, 쌍
$(U_{ij} \to V_i, E|_{U_{ij}})$가 Lemma
\ref{more-morphisms-lemma-relatively-pseudo-coherent}의 동치 조건들을 만족한다는
가정으로부터 쌍 $(U' \to V, E|_{U'})$도 Lemma
\ref{more-morphisms-lemma-relatively-pseudo-coherent}의 동치 조건들을 만족함이 따른다.
\end{proof}

\noindent
코호몰로지 층들이 준연접인 유도범주의 대상에 대해서는 상대
$m$-유사연접성을 More on Algebra, Definition
\ref{more-algebra-definition-relatively-pseudo-coherent}에서 정의한
개념과 연관 지을 수 있다. 아핀 스킴 $U = \Spec(A)$에 대해 함자
$R\Gamma(U, -)$가 $D_\QCoh(\mathcal{O}_U)$와 $D(A)$ 사이의 동치를
유도한다는 사실을 사용할 것이다. Derived Categories of Schemes, Lemma
\ref{perfect-lemma-affine-compare-bounded}를 보라. 이 함자는 당김과
양립한다. 즉 $E$가 $D_\QCoh(\mathcal{O}_U)$의 대상이고 $A \to B$가
아핀 스킴의 사상 $g : V = \Spec(B) \to \Spec(A) = U$에 대응하는
환 준동형이면
$R\Gamma(V, Lg^*E) = R\Gamma(U, E) \otimes_A^\mathbf{L} B$이다.
Derived Categories of Schemes, Lemma
\ref{perfect-lemma-quasi-coherence-pullback}를 보라.

\begin{lemma}
\label{more-morphisms-lemma-qcoh-relative-pseudo-coherence-characterize}
$f : X \to S$를 국소 유한형인 스킴의 사상이라 하자.
$E$를 $D_\QCoh(\mathcal{O}_X)$의 대상으로 하고
$m \in \mathbf{Z}$를 고정한다. 다음 조건들은 서로 동치이다.
\begin{enumerate}
\item $E$가 $S$에 상대적으로 $m$-유사연접이다.
\item 아핀 열린 덮개 $S = \bigcup V_i$와, 각 $i$마다 아핀 열린 덮개
$f^{-1}(V_i) = \bigcup U_{ij}$가 존재하여
$\mathcal{O}_X(U_{ij})$-가군 복합체 $R\Gamma(U_{ij}, E)$가
$\mathcal{O}_S(V_i)$에 상대적으로 $m$-유사연접이다.
% P05-EN-CH37-0275
\item $f(U) \subset V$인 모든 아핀 열린집합 쌍 $U \subset X$,
$V \subset S$에 대해 $\mathcal{O}_X(U)$-가군 복합체
$R\Gamma(U, E)$가 $\mathcal{O}_S(V)$에 상대적으로
$m$-유사연접이다.
\end{enumerate}
\end{lemma}

\begin{proof}
(2)와 같은 $U$와 $V$를 택하고 닫힌 몰입
$i : U \to \mathbf{A}^n_V$를 택한다. Lemma
\ref{more-morphisms-lemma-relative-pseudo-coherence-characterize}를 사용하는 형식적
논증에 의해, $Ri_*(E|_U)$가 $m$-유사연접일 필요충분조건이
$R\Gamma(U, E)$가 $\mathcal{O}_S(V)$에 상대적으로
$m$-유사연접인 것임을 보이면 충분하다. $U = \Spec(A)$,
$V = \Spec(R)$라 하고,
% P05-EN-CH37-0276
$\mathbf{A}^n_V = \Spec(R[x_1, \ldots, x_n]$라 하자. 보조정리 앞의
설명에 의해 $E|_U$는 $D(A)$의 대상 $R\Gamma(U, E)$에 대응하는
$U$ 위의 준연접 층 복합체와 준동형이다. $i$는 닫힌 몰입이고 따라서
$i_*$는 완전함자이므로
$R\Gamma(U, E) = R\Gamma(\mathbf{A}^n_V, Ri_*(E|_U))$임에 유의하자.
% P05-EN-CH37-0277
따라서 $Ri_*E$는 $D(R[x_1, \ldots, x_n])$의 대상으로 본
$R\Gamma(U, E)$에 대응한다. Derived Categories of Schemes, Lemma
\ref{perfect-lemma-pseudo-coherent-affine}에 의해 $Ri_*(E|_U)$의
$m$-유사연접성은 $D(R[x_1, \ldots, x_n])$에서
$R\Gamma(U, E)$의 $m$-유사연접성과 동치이고, 이는 정의에 의해
$R\Gamma(U, E)$가 $R = \mathcal{O}_S(V)$에 상대적으로
$m$-유사연접인 것과 동치이다.
\end{proof}

\begin{lemma}
\label{more-morphisms-lemma-closed-morphism-relative-pseudo-coherence}
% P05-EN-CH37-0278
$i : X \to Y$를 밑스킴 $S$ 위에서 국소 유한형인 스킴의 사상이라 하자.
$i$가 $X$와 $Y$의 어떤 닫힌 부분집합 사이의 위상동형을 유도한다고
가정하자. $E$를 $D(\mathcal{O}_X)$의 대상으로 하자. 그러면 $E$가
$S$에 상대적으로 $m$-유사연접일 필요충분조건은 $Ri_*E$가
$S$에 상대적으로 $m$-유사연접인 것이다.
\end{lemma}

\begin{proof}
Morphisms, Lemma \ref{morphisms-lemma-homeomorphism-affine}에 의해 사상
$i$는 아핀이다. 따라서 $S$, $Y$, $X$가 아핀이라고 가정할 수 있다.
$S = \Spec(R)$, $Y = \Spec(A)$, $X = \Spec(B)$라 하자.
% P05-EN-CH37-0279
이 조건은 $A/\text{rad}(A) \to B/\text{rad}(B)$가 전사임을 뜻한다.
여기서 $\text{rad}(A)$와 $\text{rad}(B)$는 각각 $A$와 $B$의
Jacobson 근기이다. $B$가 $A$ 위에서 유한형이므로, $B$를
$A$-대수로서 생성하는 $b_1, \ldots, b_m \in \text{rad}(B)$를 찾을
수 있다. 모든 $j$에 대해 $b_j^N = 0$이라 하자. 환의 그림
$$
\xymatrix{
B & R[x_i, y_j]/(y_j^N) \ar[l] & R[x_i, y_j] \ar[l] \\
A \ar[u] & R[x_i] \ar[l] \ar[u] \ar[ru]
}
$$
을 생각하면, 이는 아핀 스킴의 그림
$$
\xymatrix{
X \ar[d] \ar[r] & T \ar[d] \ar[r] & \mathbf{A}^{n + m}_S \ar[ld] \\
Y \ar[r] & \mathbf{A}^n_S
}
$$
으로 옮겨진다. Lemma \ref{more-morphisms-lemma-relative-pseudo-coherence-characterize}에
의해 $E$가 $S$에 상대적으로 $m$-유사연접일 필요충분조건은 이를
$\mathbf{A}^{n + m}_S$로 전진시킨 대상이 $m$-유사연접인 것이다.
Derived Categories of Schemes, Lemma
\ref{perfect-lemma-closed-push-pseudo-coherent}에 의해 이는 다시 이를
$T$로 전진시킨 대상이 $m$-유사연접인 것과 동치이다. 같은 보조정리에
의해 이는 $\mathbf{A}^n_S$로의 전진이 $m$-유사연접인 것과 동치이다.
마지막으로 Lemma \ref{more-morphisms-lemma-relative-pseudo-coherence-characterize}에
의해 이는 $Ri_*E$가 $S$에 상대적으로 $m$-유사연접인 것과 동치이다.
\end{proof}

\begin{lemma}
\label{more-morphisms-lemma-finite-morphism-relative-pseudo-coherence}
$\pi : X \to Y$를 밑스킴 $S$ 위에서 국소 유한형인 스킴의 유한
사상이라 하자. $E$를 $D_\QCoh(\mathcal{O}_X)$의 대상으로 하자.
그러면 $E$가 $S$에 상대적으로 $m$-유사연접일 필요충분조건은
$R\pi_*E$가 $S$에 상대적으로 $m$-유사연접인 것이다.
\end{lemma}

\begin{proof}
More on Algebra, Lemma
\ref{more-algebra-lemma-finite-extension-pseudo-coherent}의 결과를
스킴의 언어로 옮긴 것이다. Derived Categories of Schemes, Lemma
\ref{perfect-lemma-quasi-coherence-direct-image}에 의해 $R\pi_*$가 실제로
$D_\QCoh(\mathcal{O}_X)$를 $D_\QCoh(\mathcal{O}_Y)$로 보냄에 유의하자.
이 번역에는 Lemma \ref{more-morphisms-lemma-relative-pseudo-coherence-characterize}를
사용한다.
\end{proof}

\begin{lemma}
\label{more-morphisms-lemma-cone-relatively-pseudo-coherent}
$f : X \to S$를 국소 유한형인 스킴의 사상이라 하자.
$(E, E', E'')$를 $D(\mathcal{O}_X)$의 완전 삼각형이라 하고,
$m \in \mathbf{Z}$라 하자.
\begin{enumerate}
\item $E$가 $S$에 상대적으로 $(m + 1)$-유사연접이고 $E'$가
$S$에 상대적으로 $m$-유사연접이면, $E''$는 $S$에 상대적으로
$m$-유사연접이다.
\item $E, E''$가 $S$에 상대적으로 $m$-유사연접이면, $E'$도
$S$에 상대적으로 $m$-유사연접이다.
\item $E'$가 $S$에 상대적으로 $(m + 1)$-유사연접이고 $E''$가
$S$에 상대적으로 $m$-유사연접이면, $E$는 $S$에 상대적으로
$(m + 1)$-유사연접이다.
\end{enumerate}
% P05-EN-CH37-0280
더욱이 $E, E', E''$ 가운데 둘이 $S$에 상대적으로 유사연접이면
나머지 하나도 그러하다.
\end{lemma}

\begin{proof}
Lemma \ref{more-morphisms-lemma-relative-pseudo-coherence-characterize}와 Cohomology,
Lemma \ref{cohomology-lemma-cone-pseudo-coherent}에서 즉시 따른다.
\end{proof}

\begin{lemma}
\label{more-morphisms-lemma-rel-n-pseudo-module}
$X \to S$를 국소 유한형인 스킴의 사상이라 하고,
$\mathcal{F}$를 $\mathcal{O}_X$-가군이라 하자. 그러면
\begin{enumerate}
\item 모든 $m > 0$에 대해 $\mathcal{F}$는 $S$에 상대적으로
$m$-유사연접이고,
\item $\mathcal{F}$가 $S$에 상대적으로 $0$-유사연접일 필요충분조건은
$\mathcal{F}$가 유한형 $\mathcal{O}_X$-가군인 것이며,
\item $\mathcal{F}$가 $S$에 상대적으로 $(-1)$-유사연접일
필요충분조건은 $\mathcal{F}$가 준연접이고 $S$에 상대적으로 유한
표시인 것이다.
\end{enumerate}
\end{lemma}

\begin{proof}
(1)은 정의에서 즉시 따른다. (3)을 보이기 위해 $X$ 위에서 국소적으로
논증할 수 있다. 두 성질 모두 국소적이기 때문이다. 따라서 $X$와 $S$가
아핀이라고 가정할 수 있다. 닫힌 몰입 $i : X \to \mathbf{A}^n_S$를
택한다. 그러면 $\mathcal{F}$가 $S$에 상대적으로 $(-1)$-유사연접일
필요충분조건은 $i_*\mathcal{F}$가 $(-1)$-유사연접인 것이고, 이는 다시
$i_*\mathcal{F}$가 유한 표시 $\mathcal{O}_{\mathbf{A}^n_S}$-가군인 것과
동치이다. Cohomology, Lemma \ref{cohomology-lemma-finite-cohomology}를
보라. 유한 표시 가군은 준연접이다. Modules, Lemma
\ref{modules-lemma-finite-presentation-quasi-coherent}를 보라. Morphisms,
Lemma \ref{morphisms-lemma-i-star-equivalence}에 의해 $\mathcal{F}$가
준연접일 필요충분조건은 $i_*\mathcal{F}$가 준연접인 것이다. 따라서
(3)이 따른다. (2)의 증명은 이와 비슷하지만 더 간단하다.
\end{proof}

\begin{lemma}
\label{more-morphisms-lemma-summands-relative-pseudo-coherent}
$X \to S$를 국소 유한형인 스킴의 사상이라 하자.
$m \in \mathbf{Z}$라 하고 $E, K$를 $D(\mathcal{O}_X)$의 대상으로 하자.
$E \oplus K$가 $S$에 상대적으로 $m$-유사연접이면 $E$와 $K$도
그러하다.
\end{lemma}

\begin{proof}
Cohomology, Lemma \ref{cohomology-lemma-summands-pseudo-coherent}와
정의에서 따른다.
\end{proof}

\begin{lemma}
\label{more-morphisms-lemma-complex-relative-pseudo-coherent-modules}
$X \to S$를 국소 유한형인 스킴의 사상이라 하고,
$m \in \mathbf{Z}$라 하자. $\mathcal{F}^\bullet$를 (국소적으로) 위로
유계인 $\mathcal{O}_X$-가군 복합체라 하되, 모든 $i$에 대해
$\mathcal{F}^i$가 $S$에 상대적으로 $(m - i)$-유사연접이라고 하자.
그러면 $\mathcal{F}^\bullet$는 $S$에 상대적으로 $m$-유사연접이다.
\end{lemma}

\begin{proof}
Cohomology, Lemma \ref{cohomology-lemma-complex-pseudo-coherent-modules}와
정의에서 따른다.
\end{proof}

\begin{lemma}
\label{more-morphisms-lemma-cohomology-relative-pseudo-coherent}
$X \to S$를 국소 유한형인 스킴의 사상이라 하고,
$m \in \mathbf{Z}$라 하자. $E$를 $D(\mathcal{O}_X)$의 대상으로 하자.
$E$가 (국소적으로) 위로 유계이고 모든 $i$에 대해 $H^i(E)$가
$S$에 상대적으로 $(m - i)$-유사연접이면, $E$는 $S$에 상대적으로
$m$-유사연접이다.
\end{lemma}

\begin{proof}
Cohomology, Lemma \ref{cohomology-lemma-cohomology-pseudo-coherent}와
정의에서 따른다.
\end{proof}

\begin{lemma}
\label{more-morphisms-lemma-base-change-relative-pseudo-coherent}
$X \to S$를 국소 유한형인 스킴의 사상이라 하자.
$m \in \mathbf{Z}$라 하고, $E$를 $S$에 상대적으로 $m$-유사연접인
$D(\mathcal{O}_X)$의 대상으로 하자. $S' \to S$를 스킴의 사상이라
하고 $X' = X \times_S S'$라 놓자.
% P05-EN-CH37-0281
$E'$로 $E$의 $X'$로의 유도 당김을 나타내자. $S'$와 $X$가
$S$ 위에서 Tor-독립이면, $E'$는 $S'$에 상대적으로
$m$-유사연접이다.
\end{lemma}

\begin{proof}
문제는 $X$와 $X'$ 위에서 국소적이므로 $X$, $S$, $S'$, $X'$가
아핀이라고 가정할 수 있다. 닫힌 몰입 $i : X \to \mathbf{A}^n_S$를
택하고,
% P05-EN-CH37-0282
$i' : X' \to \mathbf{A}^n_{S'}$로 이를 $S'$로 밑변환한 사상을
나타내자.
% P05-EN-CH37-0283
$g : X' \to X$와 $g' : \mathbf{A}^n_{S'} \to \mathbf{A}^n_S$로
사영들을 나타내면 $E' = Lg^*E$이다. $X$와 $S'$가 $S$ 위에서
Tor-독립이므로 밑변환 사상(Cohomology, Remark
\ref{cohomology-remark-base-change})은 동형
$$
Ri'_*(Lg^*E) = L(g')^*Ri_*E
$$
을 유도한다. 실제로 $x \in X$ 위에 놓인 점 $x' \in X'$에 대해,
$x'$에서 줄기 사이의 밑변환 사상은 닫힌 몰입 $i$와 $i'$에서 오는 사상
$$
E_x \otimes_{\mathcal{O}_{\mathbf{A}^n_S, x}}^\mathbf{L}
\mathcal{O}_{\mathbf{A}^n_{S'}, x'}
\longrightarrow
E_x \otimes_{\mathcal{O}_{X, x}}^\mathbf{L}
\mathcal{O}_{X', x'}
$$
이다. 그 원천은
$E_x \otimes_{\mathcal{O}_{S, s}}^\mathbf{L} \mathcal{O}_{S', s'}$의
어떤 국소화와 준동형임에 유의하자. 가정에 의해
$\mathcal{O}_{X', x'}$가
$\mathcal{O}_{X, x} \otimes_{\mathcal{O}_{S, s}}^\mathbf{L}
\mathcal{O}_{S', s'}$의 (같은) 국소화와 동형이므로, 이는 표적과
동형이다. Cohomology, Lemma
\ref{cohomology-lemma-pseudo-coherent-pullback}을 적용하면 보조정리가
따른다.
\end{proof}

\begin{lemma}
\label{more-morphisms-lemma-pull-relative-pseudo-coherent}
$f : X \to Y$를 밑스킴 $S$ 위에서 국소 유한형인 스킴의 사상이라
하자. $m \in \mathbf{Z}$라 하고 $E$를 $D(\mathcal{O}_Y)$의 대상으로
하자. 다음을 가정한다.
\begin{enumerate}
\item $\mathcal{O}_X$가 $Y$에 상대적으로 유사연접이다\footnote{이는
$f$가 유사연접이라는 뜻이다. Definition
\ref{more-morphisms-definition-pseudo-coherent}를 보라.}.
\item $E$가 $S$에 상대적으로 $m$-유사연접이다.
\end{enumerate}
그러면 $Lf^*E$는 $S$에 상대적으로 $m$-유사연접이다.
\end{lemma}

\begin{proof}
문제는 $X$ 위에서 국소적이다. 따라서 $X$, $Y$, $S$가 아핀이라고
가정할 수 있다. More on Algebra, Lemma
\ref{more-algebra-lemma-pull-relative-pseudo-coherent}의 증명과 같이
논증하면 가환 그림
% P05-EN-CH37-0284
$$
\xymatrix{
X \ar[r]_i \ar[d]_f &
\mathbf{A}^d_Y \ar[r]_j \ar[ld]^p &
\mathbf{A}^{n + d}_S \ar[ld] \\
Y \ar[r] &
\mathbf{A}^n_S
}
$$
을 찾을 수 있다. 다음에 유의하자.
$$
Ri_* Lf^*E = Ri_* Li^* Lp^*E =
Lp^*E \otimes_{\mathcal{O}_{\mathbf{A}_Y^n}}^\mathbf{L} Ri_*\mathcal{O}_X
$$
이는 Cohomology, Lemma
\ref{cohomology-lemma-projection-formula-closed-immersion}에서 따른다.
가정과 $Y$가 아핀이라는 사실에 의해
% P05-EN-CH37-0285
$Ri_*\mathcal{O}_X = i_*\mathcal{O}_X$를 유한 자유
$\mathcal{O}_{\mathbf{A}_Y^n}$-가군의 복합체 $\mathcal{F}^\bullet$로
나타낼 수 있으며, $q > 0$에 대해 $\mathcal{F}^q = 0$이다. 세부사항은
생략한다. Derived Categories of Schemes, Lemma
\ref{perfect-lemma-pseudo-coherent-affine}
와 More on Algebra, Lemma
\ref{more-algebra-lemma-rel-n-pseudo-module}를 사용하라. 가정에 의해
$E$는 위로 유계이다. $q > a$에 대해 $H^q(E) = 0$이라 하자.
$q > a$에 대해 $\mathcal{E}^q = 0$인 $\mathcal{O}_Y$-가군 복합체
$\mathcal{E}^\bullet$로 $E$를 나타내자. 그러면 위의 유도 텐서곱은
$\text{Tot}(p^*\mathcal{E}^\bullet
\otimes_{\mathcal{O}_{\mathbf{A}_Y^n}} \mathcal{F}^\bullet)$로 나타난다.

\medskip\noindent
$j$는 닫힌 몰입이므로 함자 $j_*$는 완전하고, $Rj_*$는 대상을 나타내는
임의의 층 복합체에 $j_*$를 적용하여 계산된다. 따라서
$j_*\text{Tot}(p^*\mathcal{E}^\bullet \otimes_{\mathcal{O}_{\mathbf{A}_Y^n}}
\mathcal{F}^\bullet)$가 $\mathcal{O}_{\mathbf{A}^{n + m}_S}$-가군의
복합체로서 $m$-유사연접임을 보여야 한다. 복합체
$\text{Tot}(p^*\mathcal{E}^\bullet \otimes_{\mathcal{O}_{\mathbf{A}_Y^n}}
\mathcal{F}^\bullet)$는 부분복합체들에 의한 여과를 가지며, 그 연속
몫들은 복합체
$p^*\mathcal{E}^\bullet
\otimes_{\mathcal{O}_{\mathbf{A}_Y^n}} \mathcal{F}^q[-q]$이다.
$q \ll 0$에 대해 복합체
$p^*\mathcal{E}^\bullet \otimes_{\mathcal{O}_{\mathbf{A}_Y^n}}
\mathcal{F}^q[-q]$
는 차수 $\leq m$에서 코호몰로지가 영이므로 $m$-유사연접임에 유의하자.
따라서 Lemma \ref{more-morphisms-lemma-cone-relatively-pseudo-coherent}와 귀납법을
적용하면 모든 $q$에 대해
$p^*\mathcal{E}^\bullet \otimes_{\mathcal{O}_{\mathbf{A}_Y^n}}
\mathcal{F}^q[-q]$
% P05-EN-CH37-0286
이 $S$에 상대적으로 유사연접임을 보이면 충분하다.
$q > 0$에 대해 $\mathcal{F}^q = 0$이고 $\mathcal{F}^q$가 유한
자유이기도 하므로, 이는 $p^*\mathcal{E}^\bullet$가 $S$에 상대적으로
$m$-유사연접임을 보이는 것으로 환원된다. 예를 들어 Lemma
\ref{more-morphisms-lemma-base-change-relative-pseudo-coherent}에서 이 사실이 따른다.
\end{proof}

\begin{lemma}
\label{more-morphisms-lemma-composition-relative-pseudo-coherent}
$f : X \to Y$를 밑스킴 $S$ 위에서 국소 유한형인 스킴의 사상이라
하자. $m \in \mathbf{Z}$라 하고 $E$를 $D(\mathcal{O}_X)$의 대상으로
하자. $\mathcal{O}_Y$가 $S$에 상대적으로 유사연접이라고 가정하자
\footnote{이는 $Y \to S$가 유사연접이라는 뜻이다. Definition
\ref{more-morphisms-definition-pseudo-coherent}를 보라.}. 그러면 다음은 서로 동치이다.
\begin{enumerate}
\item $E$가 $Y$에 상대적으로 $m$-유사연접이다.
\item $E$가 $S$에 상대적으로 $m$-유사연접이다.
\end{enumerate}
\end{lemma}

\begin{proof}
문제는 $X$ 위에서 국소적이므로 $X$, $Y$, $S$가 아핀이라고 가정할
수 있다. More on Algebra, Lemma
\ref{more-algebra-lemma-pull-relative-pseudo-coherent}의 증명과 같이
논증하면 가환 그림
% P05-EN-CH37-0287
$$
\xymatrix{
X \ar[r]_i \ar[d]_f &
\mathbf{A}^m_Y \ar[r]_j \ar[ld]^p &
\mathbf{A}^{n + m}_S \ar[ld] \\
Y \ar[r] &
\mathbf{A}^n_S
}
$$
을 찾을 수 있다. $\mathcal{O}_Y$가 $S$에 상대적으로 유사연접이라는
가정은 $\mathcal{O}_{\mathbf{A}^m_Y}$가 $\mathbf{A}^m_S$에 상대적으로
유사연접임을 함의한다. 이는 평탄 밑변환에 의한 것이며, 예를 들어
Lemma \ref{more-morphisms-lemma-base-change-relative-pseudo-coherent}를 사용해 볼 수
있다. 이어
% P05-EN-CH37-0288
$j_*\mathcal{O}_{\mathbf{A}^n_Y}$가 유사연접
$\mathcal{O}_{\mathbf{A}^{n + m}_S}$-가군임이 따른다. 이제 보조정리의
동치는 Derived Categories of Schemes, Lemma
\ref{perfect-lemma-closed-push-pseudo-coherent}에서 따른다.
\end{proof}

\begin{lemma}
\label{more-morphisms-lemma-check-relative-pseudo-coherence-on-charts}
다음
$$
\xymatrix{
X \ar[rd] \ar[rr]_i & & P \ar[ld] \\
& S
}
$$
를 스킴의 가환 그림이라 하자. $i$가 닫힌 몰입이고,
% P05-EN-CH37-0289
$P \to S$가 평탄하고 국소 유한 표시라고 가정하자. $E$를
$D(\mathcal{O}_X)$의 대상으로 하자. 그러면 다음은 서로 동치이다.
\begin{enumerate}
\item $E$가 $S$에 상대적으로 $m$-유사연접이다.
\item $Ri_*E$가 $S$에 상대적으로 $m$-유사연접이다.
\item $Ri_*E$가 $P$ 위에서 $m$-유사연접이다.
\end{enumerate}
\end{lemma}

\begin{proof}
(1)과 (2)의 동치는 Lemma
\ref{more-morphisms-lemma-finite-morphism-relative-pseudo-coherence}이다.
$\mathcal{O}_P$가 $S$에 상대적으로 유사연접임을 보일 수 있으면,
$\text{id} : P \to P$에 Lemma
\ref{more-morphisms-lemma-composition-relative-pseudo-coherent}를 적용하여 (2)와 (3)의
동치를 얻는다. 필요한 사실은 More on Algebra, Lemma
\ref{more-algebra-lemma-flat-finite-presentation-perfect}와 정의에서 따른다.
\end{proof}









\section{유사연접 사상}
\label{more-morphisms-section-pseudo-coherent}

\noindent
% P05-EN-CH37-0290
어떤 대가를 치르더라도 이 절은 읽지 말라.
이 내용 가운데 일부가 필요하다면 먼저 대응하는 대수학 절들을 살펴보라.
More on Algebra, Sections \ref{more-algebra-section-pseudo-coherent},
\ref{more-algebra-section-relative-pseudo-coherent}, and
\ref{more-algebra-section-pseudo-coherent-perfect-ring-map}를 보라.
지금 알아야 할 것은 환 준동형 $A \to B$가 유사연접일 필요충분조건이
$B = A[x_1, \ldots, x_n]/I$이고, $A[x_1, \ldots, x_n]$-가군으로서
$B$가 유한 자유 $A[x_1, \ldots, x_n]$-가군들에 의한 분해를 갖는
것이라는 사실뿐이다.

\begin{lemma}
\label{more-morphisms-lemma-pseudo-coherent}
$f : X \to S$를 스킴의 사상이라 하자. 다음 조건들은 서로 동치이다.
\begin{enumerate}
% P05-EN-CH37-0291
\item 아핀 열린 덮개 $S = \bigcup V_j$와, 각 $j$마다 아핀 열린 덮개
$f^{-1}(V_j) = \bigcup U_{ji}$가 존재하여 환 준동형
$\mathcal{O}_S(V_j) \to \mathcal{O}_X(U_{ij})$가 유사연접이다.
\item $f(U) \subset V$인 모든 아핀 열린집합 쌍 $U \subset X$,
$V \subset S$에 대해 환 준동형
$\mathcal{O}_S(V) \to \mathcal{O}_X(U)$가 유사연접이다.
\item $f$가 국소 유한형이고 $\mathcal{O}_X$가 $S$에 상대적으로
유사연접이다.
\end{enumerate}
\end{lemma}

\begin{proof}
(1)과 (2)의 동치를 보이려면, 유사연접 환 준동형이라는 성질에 대해
Morphisms, Definition \ref{morphisms-definition-property-local}의
조건 (1)(a), (b), (c)를 확인하면 충분하다. 이 성질들은 국소화가
평탄하다는 사실을 사용하여 More on Algebra, Lemmas
\ref{more-algebra-lemma-base-change-relative-pseudo-coherent},
\ref{more-algebra-lemma-localize-relative-pseudo-coherent}, and
\ref{more-algebra-lemma-glue-relative-pseudo-coherent}에서 따른다.

\medskip\noindent
(1)이 성립하면 유사연접 환 준동형은 정의상 유한형이므로 $f$는
국소 유한형이다. 또한 (1)은 Lemma
\ref{more-morphisms-lemma-qcoh-relative-pseudo-coherence-characterize}와 정의를 통해
$\mathcal{O}_X$가 $S$에 상대적으로 유사연접임을 함의한다. 반대로
(3)이 성립한다고 하자. (2)와 같은 모든 $U$, $V$에 대해 환
$\mathcal{O}_X(U)$는 $\mathcal{O}_S(V)$ 위에서 유한형이고, 또한
$\mathcal{O}_X(U)$는 가군으로서 $\mathcal{O}_S(V)$에 상대적으로
유사연접이다. Lemmas
\ref{more-morphisms-lemma-relative-pseudo-coherence-characterize} and
\ref{more-morphisms-lemma-qcoh-relative-pseudo-coherence-characterize}를 보라. 이것이
바로 유사연접 환 준동형의 정의이므로 (2)와 (1)이 성립한다.
\end{proof}

\begin{definition}
\label{more-morphisms-definition-pseudo-coherent}
스킴의 사상 $f : X \to S$가 Lemma \ref{more-morphisms-lemma-pseudo-coherent}의 동치
조건들을 만족하면 이를 {\it 유사연접}이라고 한다. 이 경우 $X$가
$S$ 위에서 유사연접이라고도 한다.
\end{definition}

\noindent
유사연접 사상의 밑변환이 일반적으로는 유사연접이 아님에 주의하라.

\begin{lemma}
\label{more-morphisms-lemma-flat-base-change-pseudo-coherent}
유사연접 사상의 평탄 밑변환은 유사연접이다.
\end{lemma}

\begin{proof}
이는 다음 대수학 결과로 옮겨진다. $A \to B$를 유사연접 환 준동형이라
하고 $A \to A'$가 평탄하다고 하자. 그러면
$A' \to B \otimes_A A'$는 유사연접이다. 이는 더 일반적인 More on
Algebra, Lemma
\ref{more-algebra-lemma-base-change-relative-pseudo-coherent}에서 따른다.
\end{proof}

\begin{lemma}
\label{more-morphisms-lemma-composition-pseudo-coherent}
스킴의 유사연접 사상들의 합성은 유사연접이다.
\end{lemma}

\begin{proof}
이는 다음 대수학 결과로 옮겨진다. $A \to B \to C$가 합성 가능한
유사연접 환 준동형들이면 $A \to C$는 유사연접이다. 이는 More on
Algebra, Lemma \ref{more-algebra-lemma-pull-relative-pseudo-coherent}
또는 More on Algebra, Lemma
\ref{more-algebra-lemma-composition-relative-pseudo-coherent}에서 따른다.
\end{proof}

\begin{lemma}
\label{more-morphisms-lemma-pseudo-coherent-finite-presentation}
유사연접 사상은 국소 유한 표시이다.
\end{lemma}

\begin{proof}
정의에서 즉시 따른다.
\end{proof}

\begin{lemma}
\label{more-morphisms-lemma-flat-finite-presentation-pseudo-coherent}
평탄하고 국소 유한 표시인 사상은 유사연접이다.
\end{lemma}

\begin{proof}
유한 표시인 평탄 환 준동형이 유사연접이고 (더 나아가 완전하다는)
사실에서 따른다. More on Algebra, Lemma
\ref{more-algebra-lemma-flat-finite-presentation-perfect}를 보라.
\end{proof}

\begin{lemma}
\label{more-morphisms-lemma-permanence-pseudo-coherent}
$f : X \to Y$를 밑스킴 $S$ 위에서 유사연접인 스킴들의 사상이라
하자. 그러면 $f$는 유사연접이다.
\end{lemma}

\begin{proof}
% P05-EN-CH37-0292
이는 다음 대수학 결과로 옮겨진다. $R \to A \to B$가 합성 가능한 환
준동형들이고 $R \to A$, $R \to B$가 유사연접이면, $R \to B$는
유사연접이다. 이는 More on Algebra, Lemma
\ref{more-algebra-lemma-composition-relative-pseudo-coherent}에서 따른다.
\end{proof}

\begin{lemma}
\label{more-morphisms-lemma-finite-pseudo-coherent}
$f : X \to S$를 스킴의 유한 사상이라 하자. 그러면 $f$가
유사연접일 필요충분조건은 $f_*\mathcal{O}_X$가
$\mathcal{O}_S$-가군으로서 유사연접인 것이다.
\end{lemma}

\begin{proof}
대수학으로 옮기면 이 보조정리는 다음을 말한다. $R \to A$가 유한 환
준동형이면, $R \to A$가 환 준동형으로서 유사연접일 필요충분조건은
$A$가 $R$-가군으로서 유사연접인 것이다. 여기서 전자는 정의에 의해
$A$가 $A$-가군으로서 $R$에 상대적으로 유사연접이라는 뜻이다.
이는 더 일반적인 More on Algebra, Lemma
\ref{more-algebra-lemma-finite-extension-pseudo-coherent}에서 따른다.
\end{proof}

\begin{lemma}
\label{more-morphisms-lemma-Noetherian-pseudo-coherent}
$f : X \to S$를 스킴의 사상이라 하자. $S$가 국소 Noether이면,
$f$가 유사연접일 필요충분조건은 $f$가 국소 유한형인 것이다.
\end{lemma}

\begin{proof}
% P05-EN-CH37-0293
이는 다음 대수학 결과로 옮겨진다. $R$가 Noether이고 $R \to A$가
유한형 환 준동형이면, $R \to A$가 유사연접일 필요충분조건은
$R \to A$가 유한형인 것이다. 이를 보이기 위해 유사연접 환 준동형은
정의상 유한형임에 유의하자. 반대로 $R \to A$가 유한형이면
$A = R[x_1, \ldots, x_n]/I$라 쓸 수 있고, More on Algebra, Lemma
\ref{more-algebra-lemma-Noetherian-pseudo-coherent}에 의해 $A$는
$R[x_1, \ldots, x_n]$-가군으로서 유사연접이다. 즉 $R \to A$는
유사연접 환 준동형이다.
\end{proof}

\begin{lemma}
\label{more-morphisms-lemma-descending-property-pseudo-coherent}
성질 $\mathcal{P}(f) =$``$f$는 유사연접이다''는 밑에서 fpqc
국소적이다.
\end{lemma}

\begin{proof}
이를 증명하기 위해 Descent, Lemma
\ref{descent-lemma-descending-properties-morphisms}의 판정법을 사용할
것이다. Definition \ref{more-morphisms-definition-pseudo-coherent}에 의해 유사연접성은
밑에서 Zariski 국소적이다. Lemma
\ref{more-morphisms-lemma-flat-base-change-pseudo-coherent}에 의해 유사연접성은 평탄
밑변환으로 보존된다. Descent, Lemma
\ref{descent-lemma-descending-properties-morphisms}의 마지막 가정 (3)은
다음 대수학 명제로 옮겨진다. $A \to B$를 충실히 평탄한 환 준동형이라
하고 $C = A[x_1, \ldots, x_n]/I$를 $A$-대수라 하자.
% P05-EN-CH37-0294
$C \otimes_A B$가 $B[x_1, \ldots, x_n]$-가군으로서 유사연접이면,
% P05-EN-CH37-0295
$C$는 $A[x_1, \ldots, x_n]$-가군으로서 유사연접이다. 이것이 More on
Algebra, Lemma \ref{more-algebra-lemma-flat-descent-pseudo-coherent}이다.
\end{proof}

\begin{lemma}
\label{more-morphisms-lemma-quotient-of-flat-finitely-presented}
$A \to B$를 유한 표시인 평탄 환 준동형이라 하고 $I \subset B$를
아이디얼이라 하자. 그러면 $A \to B/I$가 유사연접일 필요충분조건은
$I$가 $B$-가군으로서 유사연접인 것이다.
\end{lemma}

\begin{proof}
표시 $B = A[x_1, \ldots, x_n]/J$를 택한다. Lemma
\ref{more-morphisms-lemma-flat-finite-presentation-pseudo-coherent}에 의해 $A \to B$가
유사연접 환 준동형이므로, $B$는 $A[x_1, \ldots, x_n]$-가군으로서
유사연접임에 유의하자. $A \to B/I$가 유사연접일 필요충분조건은
$B/I$가 $A[x_1, \ldots, x_n]$-가군으로서 유사연접인 것이다.
More on Algebra, Lemma
\ref{more-algebra-lemma-finite-push-pseudo-coherent}에 의해 이는
% P05-EN-CH37-0296
$B/I$가 $B$-가군으로서 유사연접이라는 조건과 동치이다. 짧은 완전열
$0 \to I \to B \to B/I \to 0$은 $I$가 유사연접일 필요충분조건이
$B/I$가 유사연접인 것임을 보인다. More on Algebra, Lemma
\ref{more-algebra-lemma-two-out-of-three-pseudo-coherent}를 보라. 따라서
보조정리가 증명된다.
\end{proof}

\noindent
다음 보조정리는 더 강한 Lemma
\ref{more-morphisms-lemma-pseudo-coherent-fppf-local-source}에 의해 불필요해질 것이다.

\begin{lemma}
\label{more-morphisms-lemma-pseudo-coherent-syntomic-local-source}
성질 $\mathcal{P}(f) =$``$f$는 유사연접이다''는 원천에서 syntomic
국소적이다.
\end{lemma}

\begin{proof}
이를 증명하기 위해 Descent, Lemma
\ref{descent-lemma-properties-morphisms-local-source}의 판정법을 사용할
것이다. Lemmas \ref{more-morphisms-lemma-flat-finite-presentation-pseudo-coherent} and
\ref{more-morphisms-lemma-composition-pseudo-coherent}에 의해, 유사연접성은 국소 유한
표시인 평탄 사상과 전합성할 때 보존되고, 특히 syntomic 사상과
전합성할 때 보존된다. Morphisms, Lemmas
\ref{morphisms-lemma-syntomic-flat} and
\ref{morphisms-lemma-syntomic-locally-finite-presentation}를 보라.
Definition \ref{more-morphisms-definition-pseudo-coherent}에서 유사연접성이 원천과
표적에서 Zariski 국소적임은 명백하다. 따라서 앞서 언급한 Descent,
Lemma \ref{descent-lemma-properties-morphisms-local-source}에 의해 다음을
증명하면 충분하다. $X' \to X \to Y$가 아핀 스킴의 사상들이고,
% P05-EN-CH37-0297
$X' \to X$가 syntomic이며 $X' \to Y$가 유사연접이라고 하자. 그러면
$X \to Y$가 유사연접이다. 실제로 Descent, Lemma
\ref{descent-lemma-flat-finitely-presented-permanence-algebra}에 의해
$X \to Y$는 어쨌든 유한 표시이다. 닫힌 몰입
$X \to \mathbf{A}^n_Y$를 택한다. Algebra, Lemma
\ref{algebra-lemma-lift-syntomic}에 의해
% P05-EN-CH37-0298
아핀 열린 덮개 $X' = \bigcup_{i = 1, \ldots, n} X'_i$와 사상들
$X'_i \to X$를 들어 올리는 syntomic 사상들
$W_i \to \mathbf{A}^n_Y$를 찾을 수 있다. 즉, 올곱 그림들
$$
\xymatrix{
X'_i \ar[d] \ar[r] & W_i \ar[d] \\
X \ar[r] & \mathbf{A}^n_Y
}
$$
이 존재한다. $X'$를 $\coprod X'_i$로 대체하고 $W = \coprod W_i$라
놓으면 올곱 그림
$$
\xymatrix{
X' \ar[d] \ar[r] & W \ar[d]^h \\
X \ar[r] & \mathbf{A}^n_Y
}
$$
을 얻는다. 여기서 $W \to \mathbf{A}^n_Y$는 평탄하고 유한 표시이며
$X' \to Y$는 여전히 유사연접이다. $W \to \mathbf{A}^n_Y$는 열린
사상이고(Morphisms, Lemma \ref{morphisms-lemma-fppf-open}를 보라),
$X' \to X$는 전사이므로
$X \subset D(f) \subset \Im(h)$를 만족하는
$f \in \Gamma(\mathbf{A}^n_Y, \mathcal{O})$를 찾을 수 있다.
$Y = \Spec(R)$, $X = \Spec(A)$, $X' = \Spec(A')$, $W = \Spec(B)$,
$A = R[x_1, \ldots, x_n]/I$, $A' = B/IB$라 쓰자. 그러면
$R \to A'$는 유사연접이다. 그림은 다음과 같다.
$$
\xymatrix{
A' = B/IB & B \ar[l] \\
A = R[x_1, \ldots, x_n]/I \ar[u] & R[x_1, \ldots, x_n] \ar[l] \ar[u]
}
$$
Lemma \ref{more-morphisms-lemma-quotient-of-flat-finitely-presented}에 의해 $IB$는
$B$-가군으로서 유사연접이다. 위에서 $f$를 택한 방식에 의해 환 준동형
$R[x_1, \ldots, x_n]_f \to B_f$는 충실히 평탄하다. 따라서
$I_f \subset R[x_1, \ldots, x_n]_f$는 유사연접이다. More on Algebra,
Lemma \ref{more-algebra-lemma-flat-descent-pseudo-coherent}를 보라.
Lemma \ref{more-morphisms-lemma-quotient-of-flat-finitely-presented}를 한 번 더 적용하면
$R \to A$가 유사연접임을 알 수 있다.
\end{proof}

\begin{lemma}
\label{more-morphisms-lemma-pseudo-coherent-fppf-local-source}
성질 $\mathcal{P}(f) =$``$f$는 유사연접이다''는 원천에서 fppf
국소적이다.
\end{lemma}

\begin{proof}
$f : X \to S$를 스킴의 사상이라 하자. 각 합성 $f \circ g_i$가
유사연접인 fppf 덮개 $\{g_i : X_i \to X\}$를 택한다. Lemma
\ref{more-morphisms-lemma-dominate-fppf}에 의하면 다음이 존재한다.
\begin{enumerate}
\item Zariski 열린 덮개 $X = \bigcup U_j$,
\item 전사인 유한 국소 자유 사상 $W_j \to U_j$,
\item Zariski 열린 덮개 $W_j = \bigcup_k W_{j, k}$,
\item 전사인 유한 국소 자유 사상 $T_{j, k} \to W_{j, k}$.
\end{enumerate}
이들은 fppf 덮개 $\{h_{j, k} : T_{j, k} \to X\}$가 주어진 덮개
$\{X_i \to X\}$를 세분하도록 택한다.
% P05-EN-CH37-0299
$\psi_{j, k} : T_{j, k} \to X_{\alpha(j, k)}$로
$\{T_{j, k} \to X\}$가 주어진 덮개 $\{X_i \to X\}$를 세분한다는
사실을 증명하는 사상들을 나타내자. $T_{j, k} \to X$는 평탄하고 국소
유한 표시인 사상이므로, Lemma
\ref{more-morphisms-lemma-flat-finite-presentation-pseudo-coherent}에 의해 $X_i$와
$T_{j, k}$ 둘 다 $X$ 위에서 유사연접이다. 따라서
$\psi_{j, k} : T_{j, k} \to X_i$는 유사연접이다. Lemma
\ref{more-morphisms-lemma-permanence-pseudo-coherent}를 보라.
% P05-EN-CH37-0300
따라서 $T_{j, k} \to S$는 $\psi_{j, k}$와
$f \circ g_{\alpha(j, k)}$의 합성이므로 유사연접이다. Lemma
\ref{more-morphisms-lemma-composition-pseudo-coherent}를 보라. 이로써 보조정리는
Zariski 열린 덮개의 경우(이는 문제없다)와 하나의 전사인 유한 국소 자유
사상이 주는 덮개의 경우로 환원되었으며, 후자는 다음 문단에서 다룬다.

\medskip\noindent
$X' \to X \to S$가 스킴의 사상열이고, $X' \to X$가 전사인 유한
국소 자유 사상이며,
% P05-EN-CH37-0301
$X' \to Y$가 유사연접이라고 가정하자. 목표는 $X \to S$가
유사연접임을 보이는 것이다. Descent, Lemma
\ref{descent-lemma-flat-finitely-presented-permanence}에 의해 사상
$X \to S$가 국소 유한 표시임에 유의하자. 문제는 $X'$, $X$, $S$가
아핀이고 $X' \to X$가 어떤 계수 $r > 0$의 자유 사상인 경우로
환원됨이 명백하다. 대응하는 대수학 문제는 다음과 같다.
$R \to A \to A'$가 환 준동형들이고, $R \to A'$가 유사연접이며,
$R \to A$가 유한 표시이고, $A$-가군으로서
$A' \cong A^{\oplus r}$라고 하자. 목표는 $R \to A$가 유사연접임을
보이는 것이다. $R \to A'$가 유사연접이라는 가정은 $A'$가
$A'$-가군으로서 $R$에 상대적으로 유사연접이라는 뜻이다. More on
Algebra, Lemma
\ref{more-algebra-lemma-finite-extension-pseudo-coherent}에 의해 $A'$는
$A$-가군으로서도 $R$에 상대적으로 유사연접이다.
$A' \cong A^{\oplus r}$가 $A$-가군의 동형이므로 $A$는 $A$-가군으로서
$R$에 상대적으로 유사연접이다. More on Algebra, Lemma
\ref{more-algebra-lemma-summands-relative-pseudo-coherent}를 보라. 이는
정의에 의해 $R \to A$가 유사연접이라는 뜻이므로 증명이 끝난다.
\end{proof}








\section{완전 사상}
\label{more-morphisms-section-perfect}

\noindent
이 절의 내용을 이해하려면 유사연접 사상에 관한 앞 절의 내용을 조금은
이해해야 한다. 지금 알아야 할 것은 환 준동형 $A \to B$가 완전일
필요충분조건이 그것이 유사연접이고 $B$가 $A$-가군으로서 유한
Tor 차원을 갖는 것이라는 사실뿐이다.

\begin{lemma}
\label{more-morphisms-lemma-perfect}
$f : X \to S$를 국소 유한형인 스킴의 사상이라 하자.
다음 조건들은 서로 동치이다.
\begin{enumerate}
% P05-EN-CH37-0302
\item 아핀 열린 덮개 $S = \bigcup V_j$와, 각 $j$마다 아핀 열린 덮개
$f^{-1}(V_j) = \bigcup U_{ji}$가 존재하여 환 준동형
$\mathcal{O}_S(V_j) \to \mathcal{O}_X(U_{ij})$가 완전이다.
\item $f(U) \subset V$인 모든 아핀 열린집합 쌍 $U \subset X$,
$V \subset S$에 대해 환 준동형
$\mathcal{O}_S(V) \to \mathcal{O}_X(U)$가 완전이다.
\end{enumerate}
\end{lemma}

\begin{proof}
(1)을 가정하고 (2)와 같은 $U, V$를 택한다. Lemma
\ref{more-morphisms-lemma-pseudo-coherent}에 의해
$\mathcal{O}_S(V) \to \mathcal{O}_X(U)$는 유사연접이다. 따라서 환
준동형이 ``유한 Tor 차원을 갖는다''는 성질이 Morphisms, Definition
\ref{morphisms-definition-property-local}의 조건 (1)(a), (b), (c)를
만족함을 보이면 충분하다. 이 성질들은 More on Algebra, Lemmas
\ref{more-algebra-lemma-flat-push-tor-amplitude},
\ref{more-algebra-lemma-flat-base-change-finite-tor-dimension}, and
\ref{more-algebra-lemma-glue-tor-amplitude}에서 따른다. 일부 세부사항은
생략한다.
\end{proof}

\begin{definition}
\label{more-morphisms-definition-perfect}
스킴의 사상 $f : X \to S$가 Lemma \ref{more-morphisms-lemma-perfect}의 동치 조건들을
만족하면 이를 {\it 완전}이라고 한다. 이 경우 $X$가 $S$ 위에서
완전이라고도 한다.
\end{definition}

\noindent
완전 사상은 특히 유사연접이고, 따라서 국소 유한 표시임에 유의하자.
완전 사상의 밑변환은 일반적으로 완전이 아님에 주의하라.

\begin{lemma}
\label{more-morphisms-lemma-flat-base-change-perfect}
완전 사상의 평탄 밑변환은 완전이다.
\end{lemma}

\begin{proof}
이는 다음 대수학 결과로 옮겨진다. $A \to B$를 완전 환 준동형이라 하고
$A \to A'$가 평탄하다고 하자. 그러면 $A' \to B \otimes_A A'$는
완전이다. 유사연접 환 준동형에 대한 이 결과는 Lemma
\ref{more-morphisms-lemma-flat-base-change-pseudo-coherent}에서 보았다. 유한 Tor 차원에
대한 대응하는 사실은 More on Algebra, Lemma
\ref{more-algebra-lemma-flat-base-change-finite-tor-dimension}에서 따른다.
\end{proof}

\begin{lemma}
\label{more-morphisms-lemma-composition-perfect}
스킴의 완전 사상들의 합성은 완전이다.
\end{lemma}

\begin{proof}
이는 다음 대수학 결과로 옮겨진다. $A \to B \to C$가 합성 가능한
완전 환 준동형들이면 $A \to C$는 완전이다. 유사연접성에 대해서는
Lemma \ref{more-morphisms-lemma-composition-pseudo-coherent}와 그 증명에서 이 사실을
보았다. 가정에 의해 정수 $n$, $m$이 존재하여 $B$는 $A$ 위에서
Tor 차원 $\leq n$을 갖고 $C$는 $B$ 위에서 Tor 차원 $\leq m$을 갖는다.
그러면 임의의 $A$-가군 $M$에 대해
$$
M \otimes_A^{\mathbf{L}} C =
(M \otimes_A^{\mathbf{L}} B) \otimes_B^{\mathbf{L}} C
$$
이고, More on Algebra, Example \ref{more-algebra-example-tor}의 스펙트럼
열은 원하는 대로 $p > n + m$에 대해
$\text{Tor}^A_p(M, C) = 0$임을 보인다.
\end{proof}

\begin{lemma}
\label{more-morphisms-lemma-flat-finite-presentation-perfect}
$f : X \to S$를 스킴의 사상이라 하자. 다음 조건들은 서로 동치이다.
\begin{enumerate}
\item $f$가 평탄하고 완전이다.
\item $f$가 평탄하고 국소 유한 표시이다.
\end{enumerate}
\end{lemma}

\begin{proof}
(2) $\Rightarrow$ (1)은 More on Algebra, Lemma
\ref{more-algebra-lemma-flat-finite-presentation-perfect}이다. 역은
유사연접 사상이 국소 유한 표시라는 사실에서 따른다. Lemma
\ref{more-morphisms-lemma-pseudo-coherent-finite-presentation}를 보라.
\end{proof}

\begin{lemma}
\label{more-morphisms-lemma-regular-target-perfect}
$f : X \to S$를 스킴의 사상이라 하자. $S$가 정칙이고 $f$가
국소 유한형이라고 가정하자. 그러면 $f$는 완전이다.
\end{lemma}

\begin{proof}
More on Algebra, Lemma
\ref{more-algebra-lemma-regular-perfect-ring-map}을 보라.
\end{proof}

\begin{lemma}
\label{more-morphisms-lemma-regular-immersion-perfect}
스킴의 정칙 몰입은 완전이다.
스킴의 Koszul-정칙 몰입은 완전이다.
\end{lemma}

\begin{proof}
정칙 몰입은 Koszul-정칙 몰입이므로(Divisors, Lemma
\ref{divisors-lemma-regular-quasi-regular-immersion}를 보라), 두 번째
명제를 증명하면 충분하다. 이는 다음 대수학 명제로 옮겨진다.
$I \subset A$가 $A$의 Koszul-정칙 열 $f_1, \ldots, f_r$로 생성되는
아이디얼이라고 하자. 그러면 $A \to A/I$는 완전 환 준동형이다.
$A \to A/I$가 전사이므로 이것은 $A$ 위의 다항식대수에 의한
$A/I$의 표시이다. 따라서 $A/I$가 $A$-가군으로서 유사연접이고 유한
Tor 차원을 가짐을 보이면 충분하다. Koszul 열의 정의에 의해 Koszul
복합체 $K(A, f_1, \ldots, f_r)$는 $A/I$의 유한 자유 분해이다.
따라서 $A/I$는 $A$-가군의 완전 복합체이고 증명이 끝난다.
\end{proof}

\begin{lemma}
\label{more-morphisms-lemma-perfect-permanence}
다음
$$
\xymatrix{
X \ar[rr]_f \ar[rd] & & Y \ar[ld] \\
& S
}
$$
을 스킴 사상들의 가환 그림이라 하자. $Y \to S$가 매끄럽고
$X \to S$가 완전이라고 가정하자. 그러면 $f : X \to Y$는 완전이다.
\end{lemma}

\begin{proof}
$f$를 합성
$$
X \longrightarrow X \times_S Y \longrightarrow Y
$$
으로 분해할 수 있다. 여기서 첫째 사상은 $i = (1, f)$이고 둘째 사상은
사영이다. $Y \to S$는 평탄하므로(Morphisms, Lemma
\ref{morphisms-lemma-smooth-flat}를 보라), Lemma
\ref{more-morphisms-lemma-flat-base-change-perfect}에 의해 $X \times_S Y \to Y$는
완전이다. $Y \to S$가 매끄러우므로 $X \times_S Y \to X$도
매끄럽다. Morphisms, Lemma \ref{morphisms-lemma-base-change-smooth}를
보라. 따라서 $i$는 매끄러운 사상의 단면이고, 그러므로 $i$는 정칙 몰입이다.
Divisors, Lemma \ref{divisors-lemma-section-smooth-regular-immersion}를
보라. 이에 따라 $i$는 완전이다. Lemma
\ref{more-morphisms-lemma-regular-immersion-perfect}를 보라. 완전 사상들의 합성은
완전이므로 $f$가 완전이라는 결론을 얻는다. Lemma
\ref{more-morphisms-lemma-composition-perfect}를 보라.
\end{proof}

\begin{remark}
\label{more-morphisms-remark-perfect-permanence}
% P05-EN-CH37-0303
밑스킴 $S$ 위에서 완전인 스킴 $X, Y$ 사이의 사상이 반드시 완전인
것은 아니다. 예를 들어 $S = \Spec(k)$, $X = \Spec(k)$,
% P05-EN-CH37-0304
$Y = \Spec(k[x]/(x^2)$로 두고 $X \to Y$를 유일한 $S$-사상으로 잡는다.
\end{remark}

\begin{lemma}
\label{more-morphisms-lemma-descending-property-perfect}
성질 $\mathcal{P}(f) =$``$f$는 완전이다''는 밑에서 fpqc 국소적이다.
\end{lemma}

\begin{proof}
이를 증명하기 위해 Descent, Lemma
\ref{descent-lemma-descending-properties-morphisms}의 판정법을 사용할
것이다. Definition \ref{more-morphisms-definition-perfect}에 의해 완전성은 밑에서
Zariski 국소적이다. Lemma \ref{more-morphisms-lemma-flat-base-change-perfect}에 의해
완전성은 평탄 밑변환으로 보존된다. Descent, Lemma
\ref{descent-lemma-descending-properties-morphisms}의 마지막 가정 (3)은
다음 대수학 명제로 옮겨진다. $A \to B$를 충실히 평탄한 환 준동형이라
하고 $C = A[x_1, \ldots, x_n]/I$를 $A$-대수라 하자.
% P05-EN-CH37-0305
$C \otimes_A B$가 $B[x_1, \ldots, x_n]$-가군으로서 완전이면,
% P05-EN-CH37-0306
$C$는 $A[x_1, \ldots, x_n]$-가군으로서 완전이다. 이것이 More on
Algebra, Lemma \ref{more-algebra-lemma-flat-descent-perfect}이다.
\end{proof}

\begin{lemma}
\label{more-morphisms-lemma-check-perfect-stalks}
$f : X \to S$를 스킴의 유사연접 사상이라 하자.
다음 조건들은 서로 동치이다.
\begin{enumerate}
\item $f$가 완전이다.
\item $\mathcal{O}_X$가 $f^{-1}\mathcal{O}_S$-가군의 층으로서
국소적으로 유한 Tor 차원을 갖는다.
\item 모든 $x \in X$에 대해 환 $\mathcal{O}_{X, x}$가
$\mathcal{O}_{S, f(x)}$-가군으로서 유한 Tor 차원을 갖는다.
\end{enumerate}
\end{lemma}

\begin{proof}
문제는 $X$와 $S$ 위에서 국소적이다. 따라서 $X = \Spec(B)$,
$S = \Spec(A)$이고 $f$가 유사연접 환 준동형 $A \to B$에 대응한다고
가정할 수 있다.

\medskip\noindent
(1)이 성립하면 $B$는 $A$-가군으로서 유한 Tor 차원 $d$를 갖는다.
그러면 $\mathfrak p = A \cap \mathfrak q$를 만족하는 모든 소 아이디얼
$\mathfrak q \subset B$에 대해 $B_\mathfrak q$는
$A_\mathfrak p$-가군으로서 Tor 차원 $d$를 갖는다. More on Algebra,
Lemma \ref{more-algebra-lemma-tor-amplitude-localization}을 보라. 이어
Cohomology, Lemma \ref{cohomology-lemma-tor-amplitude-stalk}에 의해
$\mathcal{O}_X$는 $f^{-1}\mathcal{O}_S$-가군의 층으로서 Tor 차원
$d$를 갖는다. 따라서 (1)은 (2)를 함의한다.

\medskip\noindent
Cohomology, Lemma \ref{cohomology-lemma-tor-amplitude-stalk}에 의해 (2)는
(3)을 함의한다.

\medskip\noindent
(3)을 가정하자. $B_\mathfrak q$의 $A_\mathfrak p$ 위 Tor 차원이
$\mathfrak q$와 무관하게 유계라는 조건이 주어지지 않았으므로 More on
Algebra, Lemma \ref{more-algebra-lemma-tor-amplitude-localization}을 직접
사용해 결론을 얻을 수는 없다. 표시 $A[x_1, \ldots, x_n] \to B$를
택한다. 그러면
% P05-EN-CH37-0307
$B$는 $A[x_1, \ldots, x_n]$-가군으로서 유사연접이다.
$\mathfrak q \subset A[x_1, \ldots, x_n]$를
$\mathfrak p \subset A$ 위에 놓인 소 아이디얼이라 하자. 그러면
$B_\mathfrak q$가 영이거나, 가정에 의해 $A_\mathfrak p$-가군으로서
유한 Tor 차원을 갖는다. $A \to A[x_1, \ldots, x_n]$의 올들은 유한
대역차원을 가지므로, More on Algebra, Lemma
\ref{more-algebra-lemma-perfect-over-polynomial-ring}을
$A_\mathfrak p \to A[x_1, \ldots, x_n]_\mathfrak q$에 적용하면
$B_\mathfrak q$가 완전 $A[x_1, \ldots, x_n]_\mathfrak q$-가군임을 알
수 있다. 따라서 More on Algebra, Lemma
\ref{more-algebra-lemma-check-perfect-stalks}에 의해 $B$는 완전
$A[x_1, \ldots, x_n]$-가군이다. 그러므로 정의상 $A \to B$는 완전
환 준동형이다.
\end{proof}

\begin{lemma}
\label{more-morphisms-lemma-perfect-closed-immersion-perfect-direct-image}
$i : Z \to X$를 스킴의 완전 닫힌 몰입이라 하자. 그러면
$i_*\mathcal{O}_Z$는 완전 $\mathcal{O}_X$-가군이다. 즉
$D(\mathcal{O}_X)$의 완전 대상이다.
\end{lemma}

\begin{proof}
이는 거의 정의에서 즉시 따른다. 실제로 $U = \Spec(A)$를 $X$의 아핀
열린집합이라 하자. 그러면 어떤 아이디얼 $I \subset A$에 대해
$i^{-1}(U) = \Spec(A/I)$이고, More on Algebra, Lemma
\ref{more-algebra-lemma-perfect-ring-map}에 의해 $A/I$는 유한 생성 사영
$A$-가군들에 의한 유한 분해를 갖는다. 따라서
$i_*\mathcal{O}_Z|_U$는 유한 국소 자유 $\mathcal{O}_U$-가군들의 유한
길이 복합체로 나타낼 수 있다. 이것이 보일 바였다. Cohomology, Section
\ref{cohomology-section-perfect}를 보라.
\end{proof}

\begin{lemma}
\label{more-morphisms-lemma-perfect-proper-perfect-direct-image}
$S$를 Noether 스킴이라 하고 $f : X \to S$를 스킴의 완전 고유
사상이라 하자. $E \in D(\mathcal{O}_X)$가 완전이라고 하자. 그러면
$Rf_*E$는 $D(\mathcal{O}_S)$의 완전 대상이다.
\end{lemma}

\begin{proof}
Derived Categories of Schemes, Lemma
\ref{perfect-lemma-perfect-direct-image}를 적용할 수 있다고 주장한다.
조건 (1), (2)는 즉시 성립한다. 조건 (3)은 $X$ 위에서 국소적이다.
따라서 $X$와 $S$가 아핀이고 $E$가 엄밀히 완전한
$\mathcal{O}_X$-가군 복합체로 나타난다고 가정할 수 있다. 그러므로
$\mathcal{O}_X$가 $f^{-1}\mathcal{O}_S$-가군의 층으로서 유한 Tor
차원을 가짐을 보이면 충분하다. Lemma \ref{more-morphisms-lemma-check-perfect-stalks}에
의해 이는 완전성과 동치이다.
\end{proof}

\begin{lemma}
\label{more-morphisms-lemma-perfect-fppf-local-source}
성질 $\mathcal{P}(f) =$``$f$는 완전이다''는 원천에서 fppf 국소적이다.
\end{lemma}

\begin{proof}
$\{g_i : X_i \to X\}_{i \in I}$를 스킴의 fppf 덮개라 하고,
$f : X \to S$를 모든 $f \circ g_i$가 완전인 사상이라 하자. Lemma
\ref{more-morphisms-lemma-pseudo-coherent-fppf-local-source}에 의해 $f$는 유사연접이다.
따라서 Lemma \ref{more-morphisms-lemma-check-perfect-stalks}에 의해 모든 $x \in X$에
대해 $\mathcal{O}_{X, x}$가 유한 Tor 차원의
$\mathcal{O}_{S, f(x)}$-가군임을 확인하면 충분하다. $i \in I$와
$x$로 가는 $x_i \in X_i$를 택한다. 그러면
$\mathcal{O}_{X_i, x_i}$는 $\mathcal{O}_{S, f(x)}$ 위에서 유한 Tor
차원을 갖고, $\mathcal{O}_{X, x} \to \mathcal{O}_{X_i, x_i}$는
충실히 평탄하다. 원하는 결론은 More on Algebra, Lemma
\ref{more-algebra-lemma-flat-descent-tor-amplitude}에서 따른다.
\end{proof}

\begin{lemma}
\label{more-morphisms-lemma-factor-regular-immersion}
$i : Z \to Y$와 $j : Y \to X$를 스킴의 몰입이라 하자. 다음을 가정한다.
\begin{enumerate}
\item $X$가 국소 Noether이다.
\item $j \circ i$가 정칙 몰입이다.
\item $i$가 완전이다.
\end{enumerate}
그러면 $i$와 $j$는 정칙 몰입이다.
\end{lemma}

\begin{proof}
$X$와 따라서 $Y$가 국소 Noether이므로 네 종류의 정칙 몰입이 모두
일치하고, 더 나아가 사상이 정칙 몰입인지 여부를 국소환의 수준에서
확인할 수 있다. Divisors, Lemma
\ref{divisors-lemma-Noetherian-scheme-regular-ideal}을 보라. 따라서 결과는
Divided Power Algebra, Lemma \ref{dpa-lemma-regular-sequence}에서 따른다.
\end{proof}









\section{국소 완전교차 사상}
\label{more-morphisms-section-lci}

\noindent
Divisors, Section \ref{divisors-section-regular-immersions}에서 네 종류의
정칙 몰입, 즉 정칙, Koszul-정칙, $H_1$-정칙, 준정칙 몰입을 정의했다.
이 절에서는 $X$ 위에서 국소적으로 다음과 같이 분해되는 사상
$f : X \to S$를 생각한다.
$$
\xymatrix{
X \ar[rr]_i \ar[rd] & & \mathbf{A}^n_S \ar[ld] \\
& S
}
$$
여기서 $* \in \{\emptyset, Koszul, H_1, quasi\}$에 대해 $i$는
$*$-정칙 몰입이다. 그러나 $* = \emptyset$인 경우, 즉 $i$가 정칙
몰입이어야 하는 경우에는 이 조건이 분해와 무관함을 증명하는 법을 알지
못한다. 한편 국소 완전교차 사상은 완전이기를 원하는데, 이는
$* = Koszul$ 또는 $* = \emptyset$일 때에만 참이다. 따라서 $X$ 위에서
국소적으로 위와 같이 분해되고 $i$가 Koszul-정칙 몰입인 스킴의 사상
$f : X \to S$를 {\it 국소 완전교차 사상} 또는 {\it Koszul 사상}이라
정의할 것이다. 이 정의가 제대로 작동함을 보이기 위해 먼저 이 조건이
택한 분해와 무관함을 증명한다.

\begin{lemma}
\label{more-morphisms-lemma-koszul-independence-factorization}
$S$를 스킴이라 하고 $U$, $P$, $P'$를 $S$ 위의 스킴이라 하자.
$u \in U$라 하고 $i : U \to P$, $i' : U \to P'$를 $S$ 위의 몰입이라
하자. $P$와 $P'$가 $S$ 위에서 매끄럽다고 가정하자. 다음은 서로
동치이다.
\begin{enumerate}
\item $i$가 $u$의 근방에서 Koszul-정칙 몰입이다.
\item $i'$가 $u$의 근방에서 Koszul-정칙 몰입이다.
\end{enumerate}
\end{lemma}

\begin{proof}
$i$가 $u$의 근방에서 Koszul-정칙 몰입이라고 가정하자. 사상
$j = (i, i') : U \to P \times_S P' = P''$를 생각한다.
$P'' = P \times_S P' \to P$가 매끄러우므로 Divisors, Lemma
\ref{divisors-lemma-lift-regular-immersion-to-smooth}에 의해 $j$는
Koszul-정칙 몰입이다. 그러면 Divisors, Lemma
\ref{divisors-lemma-push-regular-immersion-thru-smooth}에 의해 $i'$도
Koszul-정칙 몰입이다.
\end{proof}

\noindent
정의를 진술하기 전에 다음 간단한 사실을 언급하자. $f : X \to S$를
국소 유한형인 스킴의 사상이라 하고 $x \in X$라 하자. 그러면 열린
근방 $U \subset X$가 존재하여 $f|_U$는 몰입
$i : U \to \mathbf{A}^n_S$와 그 뒤에 오는 매끄러운 사영
$\mathbf{A}^n_S \to S$의 합성으로 분해된다. 그림으로는
$$
\xymatrix{
X \ar[rd] & U \ar[l] \ar[d] \ar[r]_-i & \mathbf{A}^n_S = P \ar[ld]^\pi \\
& S
}
$$
이다. 실제로 $X$ 안에서 $x$의 임의의 아핀 열린 근방 $U$에 대해 이렇게
할 수 있다. Morphisms, Lemma
\ref{morphisms-lemma-quasi-affine-finite-type-over-S}를 보라.

\begin{definition}
\label{more-morphisms-definition-lci}
$f : X \to S$를 스킴의 사상이라 하자.
\begin{enumerate}
\item $x \in X$라 하자. $f$가 {\it $x$에서 Koszul}이라고 하는 것은
$f$가 $x$에서 유한형이고,
% P05-EN-CH37-0308
열린 근방과 $f|_U$의 $\pi \circ i$로의 분해가 존재하여
$i : U \to P$가 Koszul-정칙 몰입이고 $\pi : P \to S$가 매끄러우면,
그렇게 된다는 뜻이다.
\item $f$가 {\it Koszul 사상}이라고, 또는 $f$가 {\it 국소 완전교차
사상}이라고 하는 것은 $f$가 모든 점에서 Koszul이라는 뜻이다.
\end{enumerate}
\end{definition}

\noindent
위에서 분해 $f|_U = \pi \circ i$의 선택은 무관함을 보았다. 즉
$f|_U$가 몰입 $i$와 그 뒤에 오는 매끄러운 사상 $\pi$의 합성으로
분해되었을 때, $i$가 $x$의 근방에서 Koszul-정칙인지 여부는 $x$에서
$f$의 내재적 성질이다. 나중에 인용할 수 있도록 이를 보조정리로
명시해 둔다.

\begin{lemma}
\label{more-morphisms-lemma-lci}
$f : X \to S$를 국소 완전교차 사상이라 하자. $P$를 $S$ 위에서
매끄러운 스킴이라 하고 $U \subset X$를 열린 부분스킴이라 하자.
$i : U \to P$가 $S$ 위의 스킴의 몰입이면 $i$는 Koszul-정칙 몰입이다.
\end{lemma}

\begin{proof}
% P05-EN-CH37-0309
이는 국소 완전교차 사상의 정의 성질이다. 위의 논의를 보라.
\end{proof}

\noindent
모든 Koszul 사상이 공통으로 갖는 몇 가지 성질을 여기에 모아 두는 것이
좋겠다.

\begin{lemma}
\label{more-morphisms-lemma-lci-properties}
$f : X \to S$를 국소 완전교차 사상이라 하자. 그러면
\begin{enumerate}
\item $f$는 국소 유한 표시이고,
\item $f$는 유사연접이며,
\item $f$는 완전이다.
\end{enumerate}
\end{lemma}

\begin{proof}
완전 사상은 유사연접이고(완전 환 준동형이 유사연접이기 때문이다),
유사연접 사상은 국소 유한 표시이므로(유사연접 환 준동형이 유한 표시이기
때문이다), 마지막 명제를 증명하면 충분하다. 완전성은 국소적 성질이므로
$f$가 $\pi \circ i$로 분해되고, 여기서 $\pi$가 매끄럽고 $i$가
Koszul-정칙 몰입이라고 가정할 수 있다. Koszul-정칙 몰입은 완전이다.
Lemma \ref{more-morphisms-lemma-regular-immersion-perfect}를 보라. 매끄러운 사상은
평탄하고 국소 유한 표시이므로 완전이다. Lemma
\ref{more-morphisms-lemma-flat-finite-presentation-perfect}를 보라. 마지막으로 완전
사상들의 합성은 완전이다. Lemma \ref{more-morphisms-lemma-composition-perfect}를 보라.
\end{proof}

\begin{lemma}
\label{more-morphisms-lemma-affine-lci}
$f : X = \Spec(B) \to S = \Spec(A)$를 아핀 스킴의 사상이라 하자.
그러면 $f$가 국소 완전교차 사상일 필요충분조건은 $A \to B$가 국소
완전교차 준동형인 것이다. More on Algebra, Definition
\ref{more-algebra-definition-local-complete-intersection}를 보라.
\end{lemma}

\begin{proof}
정의에서 즉시 따른다.
\end{proof}

\noindent
Koszul 사상의 밑변환은 일반적으로 Koszul이 아님에 주의하라.

\begin{lemma}
\label{more-morphisms-lemma-flat-base-change-lci}
국소 완전교차 사상의 평탄 밑변환은 국소 완전교차 사상이다.
\end{lemma}

\begin{proof}
생략한다. 힌트: 매끄러운 사상의 밑변환은 매끄럽고 Koszul-정칙 몰입의
평탄 밑변환은 Koszul-정칙 몰입이기 때문이다. Divisors, Lemma
\ref{divisors-lemma-regular-immersion-noetherian}을 보라.
\end{proof}

\begin{lemma}
\label{more-morphisms-lemma-composition-lci}
국소 완전교차 사상들의 합성은 국소 완전교차 사상이다.
\end{lemma}

\begin{proof}
$g : Y \to S$와 $f : X \to Y$를 국소 완전교차 사상이라 하자.
$x \in X$라 하고 $y = f(x)$라 놓는다. $y$의 열린 근방
$V \subset Y$와 분해 $g|_V = \pi \circ i$를 택하되,
$i : V \to P$는 Koszul-정칙 몰입이고 $\pi : P \to S$는 매끄럽게
한다. 다음으로 $x \in X$의 열린 근방 $U$와 분해
$f|_U = \pi' \circ i'$를 택하되, $i' : U \to P'$는 Koszul-정칙
몰입이고 $\pi' : P' \to Y$는 매끄럽게 한다. 실제로
$P' = \mathbf{A}^n_V$라고 가정할 수 있다. Definition
\ref{more-morphisms-definition-lci} 앞뒤의 논의를 보라. 그림은 다음과 같다.
$$
\xymatrix{
X \ar[d] & U \ar[l] \ar[r]_-{i'} & P' = \mathbf{A}^n_V \ar[d] \\
Y \ar[d] &  & V \ar[ll] \ar[r]_i & P \ar[d] \\
S & & & S \ar[lll]
}
$$
$P'' = \mathbf{A}^n_P$라 놓자. 그러면 $U \to P' \to P''$는 두
Koszul-정칙 몰입, 즉 $i'$와 $P'' \to P$를 통한 $i$의 평탄 밑변환의
합성이므로 Koszul-정칙 몰입이다. Divisors, Lemma
\ref{divisors-lemma-regular-immersion-noetherian} and Divisors, Lemma
\ref{divisors-lemma-composition-regular-immersion}을 보라. 또한
$P'' \to P \to S$는 매끄러운 사상들의 합성이므로 매끄럽다. Morphisms,
Lemma \ref{morphisms-lemma-composition-smooth}를 보라. 따라서 원하는 대로
$X \to S$가 $x$에서 Koszul이라는 결론을 얻는다.
\end{proof}

\begin{lemma}
\label{more-morphisms-lemma-flat-lci}
\begin{slogan}
사상이 평탄이고 lci일 필요충분조건은 그것이 syntomic인 것이다.
\end{slogan}
$f : X \to S$를 스킴의 사상이라 하자. 다음은 서로 동치이다.
\begin{enumerate}
\item $f$는 평탄이며 국소 완전교차 사상이다.
\item $f$는 syntomic이다.
\end{enumerate}
\end{lemma}

\begin{proof}
아핀 국소적으로 작업하면 이는 More on Algebra, Lemma
\ref{more-algebra-lemma-syntomic-lci}이다. 더 기하적인 증명도 제시한다.

\medskip\noindent
(2)를 가정하자. Morphisms, Lemma
\ref{morphisms-lemma-syntomic-locally-standard-syntomic}에 의해 $X$의
모든 점 $x$에 대하여 $x$의 아핀 열린 근방 $U$와 $f(x)$의 아핀 열린 근방
$V$가 존재하여 $f|_U : U \to V$는 표준 syntomic이다. 이는
$U = \Spec(R[x_1, \ldots, x_n]/(f_1, \ldots, f_c))
\to V = \Spec(R)$이고
$R[x_1, \ldots, x_n]/(f_1, \ldots, f_c)$가 $R$ 위의 상대 전역
완전교차임을 뜻한다. Algebra, Lemma
\ref{algebra-lemma-relative-global-complete-intersection-conormal}에 의해,
모든 소 아이디얼 $\mathfrak q \supset (f_1, \ldots, f_c)$에 대하여 열
$f_1, \ldots, f_c$는 각 국소환
$R[x_1, \ldots, x_n]_{\mathfrak q}$에서 정칙이다. 호몰로지 군이
$H_i = H_i(K_\bullet)$인 Koszul 복합체
$K_\bullet = K_\bullet(R[x_1, \ldots, x_n], f_1, \ldots, f_c)$
를 생각하자. More on Algebra, Lemma
\ref{more-algebra-lemma-regular-koszul-regular}에 의해 위와 같은 모든
$\mathfrak q$에 대하여 $(H_i)_{\mathfrak q} = 0$, $i > 0$임을 알 수 있다.
한편 More on Algebra, Lemma
\ref{more-algebra-lemma-homotopy-koszul}에 의해 $H_i$는
$(f_1, \ldots, f_c)$에 의해 소멸된다. 따라서 $H_i = 0$, $i > 0$이고
$f_1, \ldots, f_c$는 Koszul-정칙 열이다. 이는 원하는 대로 $U \to V$가
Koszul-정칙 몰입 $U \to \mathbf{A}^n_V$와 그 뒤의 매끄러운 사상의
합성으로 분해됨을 증명한다.

\medskip\noindent
(1)을 가정하자.
% P05-EN-CH37-0310
그러면 $f$는 평탄이고 국소 유한 표시이다(Lemma
\ref{more-morphisms-lemma-lci-properties}). 따라서 Morphisms, Lemma
\ref{morphisms-lemma-syntomic-locally-standard-syntomic}에 따라 국소환
$\mathcal{O}_{X_s, x}$들이 국소 완전교차 환임을 보이면 충분하다.
$X$ 위에서 국소적으로 어떤 Koszul-정칙 몰입 $i : X \to P$와 매끄러운
사상 $\pi : P \to S$에 대한 분해 $f = \pi \circ i$를 택하자. $X \to P$는
$S$ 위의 상대 준정칙 몰입이다. Divisors, Definition
\ref{divisors-definition-relative-H1-regular-immersion}을 보라. 따라서
Divisors, Lemma
\ref{divisors-lemma-relative-regular-immersion-flat-in-neighbourhood}에 의해
$X \to P$는 정칙 몰입이고, 임의의 밑변환 뒤에도 이 성질이 유지된다.
그러므로 각 올은 정칙 몰입이며, 따라서 $X$의 모든 올의 모든 국소환은
국소 완전교차이다.
\end{proof}

\begin{lemma}
\label{more-morphisms-lemma-regular-immersion-lci}
스킴의 정칙 몰입은 국소 완전교차 사상이다. 스킴의 Koszul-정칙 몰입도
국소 완전교차 사상이다.
\end{lemma}

\begin{proof}
정칙 몰입은 Koszul-정칙 몰입이므로(Divisors, Lemma
\ref{divisors-lemma-regular-quasi-regular-immersion} 참조), 둘째 명제만
증명하면 충분하다. 둘째 명제는 정의에서 즉시 따른다.
\end{proof}

\begin{lemma}
\label{more-morphisms-lemma-lci-permanence}
스킴 사상들의 가환 그림
$$
\xymatrix{
X \ar[rr]_f \ar[rd] & & Y \ar[ld] \\
& S
}
$$
가 주어졌다고 하자.
% P05-EN-CH37-0311
$Y \to S$가 매끄럽고 $X \to S$가 국소 완전교차 사상이라고 가정하자.
그러면 $f : X \to Y$는 국소 완전교차 사상이다.
\end{lemma}

\begin{proof}
정의에서 즉시 따른다.
\end{proof}

\begin{lemma}
\label{more-morphisms-lemma-morphism-regular-schemes-is-lci}
$f : X \to Y$를 스킴의 사상이라 하자. $f$가 국소 유한형이고 $X$와
$Y$가 정칙이면, $f$는 국소 완전교차 사상이다.
\end{lemma}

\begin{proof}
첫째 화살표가 몰입인 분해 $X \to \mathbf{A}^n_Y \to Y$가 존재한다고
가정해도 좋다. $Y$가 정칙이므로 Algebra, Lemma
\ref{algebra-lemma-regular-goes-up}에 의해 $\mathbf{A}^n_Y$도 정칙이다.
따라서 Divisors, Lemma
\ref{divisors-lemma-immersion-regular-regular-immersion}에 의해
$X \to \mathbf{A}^n_Y$는 정칙 몰입이다.
\end{proof}

\noindent
다음 보조정리는 성격이 다르다.

\begin{lemma}
\label{more-morphisms-lemma-lci-avramov}
스킴 사상들의 가환 그림
$$
\xymatrix{
X \ar[rr]_f \ar[rd] & & Y \ar[ld] \\
& S
}
$$
가 주어졌다고 하자. 다음을 가정하자.
\begin{enumerate}
\item $S$는 국소 Noether이다.
\item $Y \to S$는 국소 유한형이다.
\item $f : X \to Y$는 완전이다.
\item $X \to S$는 국소 완전교차 사상이다.
\end{enumerate}
그러면 $X \to Y$는 국소 완전교차 사상이고, 모든 $x \in X$에 대하여
$Y \to S$는 $f(x)$에서 Koszul이다.
\end{lemma}

\begin{proof}
% P05-EN-CH37-0312
이 증명에서 모든 스킴은 국소 Noether이고 모든 환은 Noether라고 한다.
이 상황에서는 정칙 열과 Koszul-정칙 열이 일치한다는 사실을 더 언급하지
않고 사용한다. More on Algebra, Lemma
\ref{more-algebra-lemma-noetherian-finite-all-equivalent}를 보라. 또한
아이디얼(각각 아이디얼 층)이 정칙인지 여부는
국소환(각각 줄기)에서 확인할 수 있다. Algebra, Lemma
\ref{algebra-lemma-regular-sequence-in-neighbourhood}(각각 Divisors, Lemma
\ref{divisors-lemma-Noetherian-scheme-regular-ideal})를 보라.

\medskip\noindent
문제는 국소적이다. 따라서 $S$, $X$, $Y$가 아핀이라고 가정해도 좋다.
이 상황에서 가로 화살표들이 닫힌 몰입인 가환 그림
$$
\xymatrix{
\mathbf{A}^{n + m}_S \ar[d] & X \ar[l] \ar[d] \\
\mathbf{A}^n_S \ar[d] & Y \ar[l] \ar[ld] \\
S
}
$$
을 택할 수 있다. 점 $x \in X$를 잡고 이에 대응하는 국소환들의 가환 그림
$$
\xymatrix{
J \ar[r] &
\mathcal{O}_{\mathbf{A}^{n + m}_S, x} \ar[r] &
\mathcal{O}_{X, x} \\
I \ar[r] \ar[u] &
\mathcal{O}_{\mathbf{A}^n_S, f(x)} \ar[r] \ar[u] &
\mathcal{O}_{Y, f(x)} \ar[u]
}
$$
을 생각하자. 여기서 $J$와 $I$는 가로 화살표들의 핵이다. $X \to S$가
국소 완전교차 사상이므로 아이디얼 $J$는 정칙 열에 의해 생성된다.
$X \to Y$가 완전이므로 환 $\mathcal{O}_{X, x}$는
$\mathcal{O}_{Y, f(x)}$ 위에서 유한 Tor 차원을 가진다. 따라서 Divided
Power Algebra, Lemma \ref{dpa-lemma-perfect-map-ci}를 적용하여 $I$와
$J/I$가 정칙 열들에 의해 생성된다는 결론을 얻는다. 처음의 언급에 의해
증명이 끝난다.
\end{proof}

\begin{lemma}
\label{more-morphisms-lemma-lci-to-regular}
스킴 사상들의 가환 그림
$$
\xymatrix{
X \ar[rr]_f \ar[rd] & & Y \ar[ld] \\
& S
}
$$
가 주어졌다고 하자. $S$는 국소 Noether이고, $Y \to S$는 국소 유한형이며,
$Y$는 정칙이고, $X \to S$는 국소 완전교차 사상이라고 가정하자. 그러면
$f : X \to Y$는 국소 완전교차 사상이고, 모든 $x \in X$에 대하여
$Y \to S$는 $f(x)$에서 Koszul이다.
\end{lemma}

\begin{proof}
이는 Lemma \ref{more-morphisms-lemma-regular-target-perfect}(및 Morphisms, Lemma
\ref{morphisms-lemma-permanence-finite-type})를 고려하면 Lemma
\ref{more-morphisms-lemma-lci-avramov}의 특수한 경우이다.
\end{proof}

\begin{lemma}
\label{more-morphisms-lemma-perfect-conormal-free-lci}
$i : X \to Y$를 몰입이라 하자. 다음을 가정하자.
\begin{enumerate}
\item $i$는 완전이다.
\item $Y$는 국소 Noether이다.
\item 여법선 층 $\mathcal{C}_{X/Y}$는 유한 국소 자유이다.
\end{enumerate}
그러면 $i$는 정칙 몰입이다.
\end{lemma}

\begin{proof}
대수로 옮기면 이는 Divided Power Algebra, Proposition
\ref{dpa-proposition-regular-ideal}이다.
\end{proof}

\begin{lemma}
\label{more-morphisms-lemma-lci-NL}
% P05-EN-CH37-0313
$f : X \to Y$를 국소 완전교차 준동형이라 하자. 그러면 소박한 여접
복합체 $\NL_{X/Y}$는 $[-1, 0]$ 안의 Tor 진폭을 가지는
$D(\mathcal{O}_X)$의 완전 대상이다.
\end{lemma}

\begin{proof}
대수로 옮기면 이는 More on Algebra, Lemma
\ref{more-algebra-lemma-lci-NL}이다. 이 번역에는 Lemmas
\ref{more-morphisms-lemma-affine-lci} 및 \ref{more-morphisms-lemma-NL-affine}와 Derived Categories of
Schemes, Lemmas
\ref{perfect-lemma-affine-compare-bounded},
\ref{perfect-lemma-tor-dimension-affine} and
\ref{perfect-lemma-perfect-affine}를 사용한다.
\end{proof}

\begin{lemma}
\label{more-morphisms-lemma-perfect-NL-lci}
$f : X \to Y$를 국소 Noether 스킴들의 완전 사상이라 하자. 다음은 서로
동치이다.
\begin{enumerate}
\item $f$는 국소 완전교차 사상이다.
\item $\NL_{X/Y}$는 $[-1, 0]$ 안의 Tor 진폭을 가진다.
\item $\NL_{X/Y}$는 완전이고 $[-1, 0]$ 안의 Tor 진폭을 가진다.
\end{enumerate}
\end{lemma}

\begin{proof}
대수로 옮기면 이는 Divided Power Algebra, Lemma
\ref{dpa-lemma-perfect-NL-lci}이다. 이 번역에는 Lemmas
\ref{more-morphisms-lemma-affine-lci} 및 \ref{more-morphisms-lemma-NL-affine}와 Derived Categories of
Schemes, Lemmas
\ref{perfect-lemma-affine-compare-bounded},
\ref{perfect-lemma-tor-dimension-affine} and
\ref{perfect-lemma-perfect-affine}를 사용한다.
\end{proof}

\begin{lemma}
\label{more-morphisms-lemma-flat-fp-NL-lci}
$f : X \to Y$를 유한 표시인 평탄 사상이라 하자. 다음은 서로 동치이다.
\begin{enumerate}
\item $f$는 국소 완전교차 사상이다.
\item $f$는 syntomic이다.
\item $\NL_{X/Y}$는 $[-1, 0]$ 안의 Tor 진폭을 가진다.
\item $\NL_{X/Y}$는 완전이고 $[-1, 0]$ 안의 Tor 진폭을 가진다.
\end{enumerate}
\end{lemma}

\begin{proof}
대수로 옮기면 이는 Divided Power Algebra, Lemma
\ref{dpa-lemma-flat-fp-NL-lci}이다. 이 번역에는 Lemmas
\ref{more-morphisms-lemma-affine-lci} 및 \ref{more-morphisms-lemma-NL-affine}와 Derived Categories of
Schemes, Lemmas
\ref{perfect-lemma-affine-compare-bounded},
\ref{perfect-lemma-tor-dimension-affine} and
\ref{perfect-lemma-perfect-affine}를 사용한다.
\end{proof}

\noindent
다음 보조정리는 매끄러운 사상을 대각 사상이 완전인 평탄 사상으로
특징짓는다.

\begin{lemma}
\label{more-morphisms-lemma-smooth-diagonal-perfect}
$f : X \to Y$를 국소 Noether 스킴들의 유한형 사상이라 하자.
% P05-EN-CH37-0314
대각 사상을 $\Delta : X \to X \times_Y X$로 나타내자. 다음은 서로
동치이다.
\begin{enumerate}
\item $f$는 매끄럽다.
\item $f$는 평탄이고 $\Delta : X \to X \times_Y X$는 정칙 몰입이다.
\item $f$는 평탄이고 $\Delta : X \to X \times_Y X$는 국소 완전교차
사상이다.
\item $f$는 평탄이고 $\Delta : X \to X \times_Y X$는 완전이다.
\end{enumerate}
\end{lemma}

\begin{proof}
(1)을 가정하자. Morphisms, Lemma \ref{morphisms-lemma-smooth-flat}에 의해
$f$는 평탄이다. Morphisms, Lemma
\ref{morphisms-lemma-base-change-smooth}에 의해 사영들
$X \times_Y X \to X$는 매끄럽다. 따라서 대각 사상은 매끄러운 사상의
절단이고, 따라서 정칙 몰입이다. Divisors, Lemma
\ref{divisors-lemma-section-smooth-regular-immersion}을 보라. 그러므로
(1) $\Rightarrow$ (2)이다. 함의 (2) $\Rightarrow$ (3)은 Lemma
\ref{more-morphisms-lemma-regular-immersion-lci}이고, 함의 (3) $\Rightarrow$ (4)는 Lemma
\ref{more-morphisms-lemma-lci-properties}이다. 흥미로운 함의 (4) $\Rightarrow$ (1)은
Divided Power Algebra, Lemma \ref{dpa-lemma-perfect-diagonal}에서 즉시 따른다.
\end{proof}

\begin{lemma}
\label{more-morphisms-lemma-descending-property-lci}
성질 $\mathcal{P}(f) =$``$f$는 국소 완전교차 사상이다''는 밑에서 fpqc
국소적이다.
\end{lemma}

\begin{proof}
$f : X \to S$를 스킴의 사상이라 하고, $\{S_i \to S\}$를 $S$의 fpqc
덮개라 하자. $f$의 각 밑변환 $f_i : X_i \to S_i$가 국소 완전교차
사상이라고 가정하자. 특히 이는 $f$가 국소 유한형임을 함의한다. Lemma
\ref{more-morphisms-lemma-lci-properties}와 Descent, Lemma
\ref{descent-lemma-descending-property-locally-finite-type}을 보라.
$x \in X$라 하자. $x$의 열린 근방 $U$와 $S$ 위의 몰입
$j : U \to \mathbf{A}^n_S$를 택하자(Definition \ref{more-morphisms-definition-lci} 앞의
논의 참조). $j$가 Koszul-정칙 몰입임을 보여야 한다. $f_i$가 국소
완전교차 사상이므로 밑변환
$j_i : U \times_S S_i \to \mathbf{A}^n_{S_i}$는 Koszul-정칙 몰입이다.
Lemma \ref{more-morphisms-lemma-lci}를 보라. $\{\mathbf{A}^n_{S_i} \to \mathbf{A}^n_S\}$가
fpqc 덮개이므로 Descent, Lemma
\ref{descent-lemma-descending-property-regular-immersion}에서 원하는 대로
$j$가 Koszul-정칙 몰입임을 알 수 있다.
\end{proof}

\begin{lemma}
\label{more-morphisms-lemma-lci-syntomic-local-source}
성질 $\mathcal{P}(f) =$``$f$는 국소 완전교차 사상이다''는 원천에서
syntomic 국소적이다.
\end{lemma}

\begin{proof}
이를 증명하기 위해 Descent, Lemma
\ref{descent-lemma-properties-morphisms-local-source}의 판정법을 사용한다.
Lemmas \ref{more-morphisms-lemma-flat-lci} 및 \ref{more-morphisms-lemma-composition-lci}에 의해 국소
완전교차 사상이라는 성질은 syntomic 사상과 전합성할 때 보존된다.
Definition \ref{more-morphisms-definition-lci}에서 국소 완전교차 사상이라는 성질이
원천과 표적에서 Zariski 국소적임은 명백하다. 따라서 앞서 말한 Descent,
Lemma \ref{descent-lemma-properties-morphisms-local-source}에 따라 다음을
증명하면 충분하다. $X' \to X \to Y$를 아핀 스킴들의 사상이라 하고,
$X' \to X$는 syntomic이며 $X' \to Y$는 국소 완전교차 사상이라고 하자.
그러면 $X \to Y$는 국소 완전교차 사상이다. 이를 보이기 위해, 어느
경우든 Descent, Lemma
\ref{descent-lemma-flat-finitely-presented-permanence-algebra}에 의해
$X \to Y$가 유한 표시임에 유의하자. 닫힌 몰입
$X \to \mathbf{A}^n_Y$를 택하자. Algebra, Lemma
\ref{algebra-lemma-lift-syntomic}에 의해 아핀 열린 덮개
$X' = \bigcup_{i = 1, \ldots, n} X'_i$와 사상들 $X'_i \to X$를 들어 올리는
syntomic 사상들 $W_i \to \mathbf{A}^n_Y$를 찾을 수 있다. 즉, 올곱 그림
$$
\xymatrix{
X'_i \ar[d] \ar[r] & W_i \ar[d] \\
X \ar[r] & \mathbf{A}^n_Y
}
$$
들이 존재한다. $X'$를 $\coprod X'_i$로 바꾸고 $W = \coprod W_i$라 놓으면
아핀 스킴들의 올곱 그림
$$
\xymatrix{
X' \ar[d] \ar[r] & W \ar[d]^h \\
X \ar[r] & \mathbf{A}^n_Y
}
$$
을 얻으며, 여기서 $h : W \to \mathbf{A}^n_Y$는 syntomic이고
$X' \to Y$는 여전히 국소 완전교차 사상이다. Morphisms, Lemma
\ref{morphisms-lemma-fppf-open}에 의해 $W \to \mathbf{A}^n_Y$는 열린
사상이고 $X' \to X$는 전사이므로, $X$는 $W \to \mathbf{A}^n_Y$의 상에
포함된다. $\mathbf{A}^n_Y$ 위의 닫힌 몰입
$W \to \mathbf{A}^{n + m}_Y$를 택하자. 이제 그림은 다음과 같다.
$$
\xymatrix{
X' \ar[d] \ar[r] & W \ar[d]^h \ar[r] & \mathbf{A}^{n + m}_Y \ar[ld] \\
X \ar[r] & \mathbf{A}^n_Y
}
$$
$h$가 syntomic이고 따라서 국소 완전교차 사상이므로(위 참조), 사상
$W \to \mathbf{A}^{n + m}_Y$는 Koszul-정칙 몰입이다. $X' \to Y$가 국소
완전교차 사상이므로 사상 $X' \to \mathbf{A}^{n + m}_Y$도 Koszul-정칙
몰입이다. Divisors, Lemma
\ref{divisors-lemma-permanence-regular-immersion}에 의해 $X' \to W$가
Koszul-정칙 몰입이라는 결론을 얻는다. Koszul-정칙 몰입이라는 성질은
표적에서 fpqc 국소적이므로(Descent, Lemma
\ref{descent-lemma-descending-property-regular-immersion} 참조),
$X \to \mathbf{A}^n_Y$가 Koszul-정칙 몰입이라는 결론을 얻는다. 이것이
보여야 할 바였다.
\end{proof}

\begin{lemma}
\label{more-morphisms-lemma-base-change-lci-fibres}
$S$를 스킴이라 하고 $f : X \to Y$를 $S$ 위의 스킴들의 사상이라 하자.
$X$와 $Y$가 모두 $S$ 위에서 평탄이고 국소 유한 표시라고 가정하자.
그러면 집합
$$
\{x \in X \mid f\text{ Koszul at }x\}.
$$
은 $X$에서 열려 있고, 그 구성은 임의의 밑변환 $S' \to S$와 가환한다.
\end{lemma}

\begin{proof}
이 집합은 정의에 의해 열려 있다(Definition \ref{more-morphisms-definition-lci} 참조).
$S' \to S$를 스킴의 사상이라 하자. $X' = S' \times_S X$,
$Y' = S' \times_S Y$라 놓고, $f$의 밑변환을 $f' : X' \to Y'$로 나타내자.
$f'$가 $x'$에서 Koszul인 점 $x' \in X'$를 잡자. $x'$의 상들을 각각
$s' \in S'$, $x \in X$, $y' \in Y'$, $y \in Y$, $s \in S$로 나타내자.
$f$가 국소 유한 표시임에 유의하라. Morphisms, Lemma
\ref{morphisms-lemma-finite-presentation-permanence}를 보라. 따라서 $x$의
아핀 근방 $U \subset X$와 몰입 $i : U \to \mathbf{A}^n_Y$를 택할 수 있다.
$U' = S' \times_S U$라 놓고, $i$의 밑변환을
$i' : U' \to \mathbf{A}^n_{Y'}$로 나타내자. $f'$가 $x'$에서 Koszul이라는
가정은 $i'$가 $x'$의 어떤 근방에서 Koszul-정칙 몰입임을 함의한다. Lemma
\ref{more-morphisms-lemma-lci}를 보라. 스킴 $X'$는 $X$의 밑변환이므로 $S'$ 위에서
평탄이고 국소 유한 표시이다(Morphisms, Lemmas
\ref{morphisms-lemma-base-change-flat} and
\ref{morphisms-lemma-base-change-finite-presentation} 참조). 따라서 $i'$는
$x'$의 어떤 근방에서 $S'$ 위의 상대 $H_1$-정칙 몰입이다(Divisors,
Definition \ref{divisors-definition-relative-H1-regular-immersion} 참조).
그러므로 밑변환 $i'_{s'} : U'_{s'} \to \mathbf{A}^n_{Y'_{s'}}$는 $x'$의
어떤 열린 근방에서 $H_1$-정칙 몰입이다. Divisors, Lemma
\ref{divisors-lemma-relative-regular-immersion} 및 Divisors, Definition
\ref{divisors-definition-relative-H1-regular-immersion} 뒤의 논의를 보라.
$s' = \Spec(\kappa(s')) \to \Spec(\kappa(s)) = s$는 전사이고 평탄이며
보편적으로 열린 사상이므로(Morphisms, Lemma
\ref{morphisms-lemma-scheme-over-field-universally-open} 참조), 밑변환
$i_s : U_s \to \mathbf{A}^n_{Y_s}$가 $x$의 어떤 근방에서 $H_1$-정칙
몰입이라는 결론을 얻는다. Descent, Lemma
\ref{descent-lemma-descending-property-regular-immersion}을 보라. 마지막으로
$\mathbf{A}^n_Y$가 $S$ 위에서 평탄이고 국소 유한 표시임에 유의하자.
따라서 Divisors, Lemma \ref{divisors-lemma-fibre-quasi-regular}에 의해
원하는 대로 $i$는 $x$의 어떤 근방에서 (Koszul-)정칙 몰입이다.
\end{proof}

\begin{lemma}
\label{more-morphisms-lemma-unramified-lci}
$f : X \to Y$를 스킴들의 국소 완전교차 사상이라 하자. 그러면 $f$가
비분기일 필요충분조건은 $f$가 형식적으로 비분기인 것이고, 이 경우
여법선 층 $\mathcal{C}_{X/Y}$는 $X$ 위에서 유한 국소 자유이다.
\end{lemma}

\begin{proof}
첫째 주장은 Lemma \ref{more-morphisms-lemma-unramified-formally-unramified}와 국소
완전교차 사상이 국소 유한형이라는 사실에서 즉시 따른다. $f$의 여법선
층을 계산하기 위해 $X$ 위에서 국소적으로 $f$를 $f = p \circ i$로
분해하자. 여기서 $i : X \to V$는 Koszul-정칙 몰입이고 $V \to Y$는
매끄럽다. Lemma \ref{more-morphisms-lemma-two-unramified-morphisms-formally-smooth}에 의해
$\mathcal{C}_{X/Y}$는 $\mathcal{C}_{X/V}$의 국소 직합인자임을 알 수 있다.
$i$가 Koszul-정칙(따라서 준정칙) 몰입이므로 이 층은 유한 국소 자유이다.
Divisors, Lemma
\ref{divisors-lemma-quasi-regular-immersion}을 보라.
\end{proof}

\begin{lemma}
\label{more-morphisms-lemma-transitivity-conormal-lci}
$Z \to Y \to X$를 형식적으로 비분기인 스킴 사상들이라 하자.
$Z \to Y$가 국소 완전교차 사상이라고 가정하자. Lemma
\ref{more-morphisms-lemma-transitivity-conormal}의 완전열
$$
0 \to i^*\mathcal{C}_{Y/X} \to
\mathcal{C}_{Z/X} \to
\mathcal{C}_{Z/Y} \to 0
$$
은 짧은 완전열이다.
\end{lemma}

\begin{proof}
문제는 $Z$ 위에서 국소적이므로, 사상 $Z \to Y$의 분해
$Z \to \mathbf{A}^n_Y \to Y$가 존재한다고 가정해도 좋다. 그러면 가환
그림
$$
\xymatrix{
Z \ar[r]_{i'} \ar@{=}[d] &
\mathbf{A}^n_Y \ar[r] \ar[d] &
\mathbf{A}^n_X \ar[d] \\
Z \ar[r]^i & Y \ar[r] & X
}
$$
을 얻는다. $Z \to Y$가 국소 완전교차 사상이므로
$Z \to \mathbf{A}^n_Y$는 Koszul-정칙 몰입이다. 따라서 Divisors, Lemma
\ref{divisors-lemma-transitivity-conormal-quasi-regular}에 의해 열
$$
0 \to (i')^*\mathcal{C}_{\mathbf{A}^n_Y/\mathbf{A}^n_X} \to
\mathcal{C}_{Z/\mathbf{A}^n_X} \to
\mathcal{C}_{Z/\mathbf{A}^n_Y} \to 0
$$
은 완전이고 국소적으로 분할된다. Lemma
\ref{more-morphisms-lemma-universal-thickening-fibre-product-flat}에 의해
$i^*\mathcal{C}_{Y/X} = (i')^*\mathcal{C}_{\mathbf{A}^n_Y/\mathbf{A}^n_X}$
임에 유의하고, 그림
$$
\xymatrix{
(i')^*\mathcal{C}_{\mathbf{A}^n_Y/\mathbf{A}^n_X} \ar[r] &
\mathcal{C}_{Z/\mathbf{A}^n_X} \\
i^*\mathcal{C}_{Y/X} \ar[u]^{\cong} \ar[r] & \mathcal{C}_{Z/X} \ar[u]
}
$$
이 가환임에 유의하자. 따라서 아래 가로 화살표는 국소적으로 분할되는
단사이다. 이로써 보조정리가 증명된다.
\end{proof}













\section{미분층과 여법선 층의 완전열}
\label{more-morphisms-section-exact}

\noindent
이 절에서는 여법선 층과 미분층의 완전열에 관한 몇 가지 결과를 모은다.
어떤 의미에서는 이들 모두가 합성 가능한 두 스킴 사상에 결부된 여접
복합체의 삼각형을 실현한 것이다.

\medskip\noindent
$g : Z \to Y$와 $f : Y \to X$를 스킴의 사상들이라 하자.
\begin{enumerate}
\item 표준 완전열
$$
g^*\Omega_{Y/X} \to \Omega_{Z/X} \to \Omega_{Z/Y} \to 0,
$$
이 존재한다. Morphisms, Lemma
\ref{morphisms-lemma-triangle-differentials}를 보라. $g : Z \to Y$가
매끄럽거나 더 일반적으로 형식적으로 매끄러우면, 이 열은 짧은 완전열이다.
Morphisms, Lemma \ref{morphisms-lemma-triangle-differentials-smooth} 또는
Lemma \ref{more-morphisms-lemma-triangle-differentials-formally-smooth}를 보라.
\item $g$가 몰입이거나 더 일반적으로 형식적으로 비분기이면, 표준 완전열
$$
\mathcal{C}_{Z/Y} \to g^*\Omega_{Y/X} \to \Omega_{Z/X} \to 0,
$$
이 존재한다. Morphisms, Lemma
\ref{morphisms-lemma-differentials-relative-immersion} 또는 Lemma
\ref{more-morphisms-lemma-universally-unramified-differentials-sequence}를 보라.
$f \circ g : Z \to X$가 매끄럽거나 더 일반적으로 형식적으로 매끄러우면,
이 열은 짧은 완전열이다. Morphisms, Lemma
\ref{morphisms-lemma-differentials-relative-immersion-smooth} 또는 Lemma
\ref{more-morphisms-lemma-differentials-formally-unramified-formally-smooth}를 보라.
\item $g$와 $f \circ g$가 몰입이거나 더 일반적으로 형식적으로 비분기이면,
표준 완전열
$$
\mathcal{C}_{Z/X} \to \mathcal{C}_{Z/Y} \to g^*\Omega_{Y/X} \to 0,
$$
이 존재한다. Morphisms, Lemma \ref{morphisms-lemma-two-immersions} 또는
Lemma \ref{more-morphisms-lemma-two-unramified-morphisms}를 보라. $f : Y \to X$가
매끄럽거나 더 일반적으로 형식적으로 매끄러우면, 이 열은 짧은 완전열이다.
Morphisms, Lemma \ref{morphisms-lemma-two-immersions-smooth} 또는 Lemma
\ref{more-morphisms-lemma-two-unramified-morphisms-formally-smooth}를 보라.
\item $g$와 $f$가 몰입이거나 더 일반적으로 형식적으로 비분기이면, 표준
완전열
$$
g^*\mathcal{C}_{Y/X} \to \mathcal{C}_{Z/X} \to \mathcal{C}_{Z/Y} \to 0.
$$
이 존재한다. Morphisms, Lemma
\ref{morphisms-lemma-transitivity-conormal} 또는 Lemma
\ref{more-morphisms-lemma-transitivity-conormal}를 보라. $g : Z \to Y$가 정칙
몰입이거나\footnote{$g$가 $H_1$-정칙 몰입이면 충분하다. 국소 완전교차
사상인 몰입은 Koszul-정칙임에 유의하라.} 더 일반적으로 국소 완전교차
사상이면, 이 열은 짧은 완전열이다. Divisors, Lemma
\ref{divisors-lemma-transitivity-conormal-quasi-regular} 또는 Lemma
\ref{more-morphisms-lemma-transitivity-conormal-lci}를 보라.
\end{enumerate}










\section{약한 \'etale 사상}
\label{more-morphisms-section-weakly-etale}

\noindent
환 준동형 $A \to B$가 {it 약한 \'etale}라는 것은 $A \to B$와
$B \otimes_A B \to B$가 모두 평탄이라는 뜻이다. More on Algebra,
Definition \ref{more-algebra-definition-weakly-etale}를 보라. 스킴의
사상에 대한 유사한 개념은 다음과 같다.

\begin{definition}
\label{more-morphisms-definition-weakly-etale}
스킴의 사상 $X \to Y$에서 $X \to Y$와 대각 사상
$X \to X \times_Y X$가 모두 평탄이면 이 사상을 {\it 약한 \'etale} 또는
{\it 절대 평탄}이라고 한다.
\end{definition}

\noindent
\'Etale 사상은 약한 \'etale이고, 역으로 약한 \'etale 사상은 실제로
어느 정도 \'etale 사상과 비슷함이 드러난다. 예를 들어
$X \to Y$가 약한 \'etale이면 Cotangent, Lemma
\ref{cotangent-lemma-when-zero}에 의해 $L_{X/Y} = 0$이다. 약한 \'etale
사상과 \'etale 사상을 연결하는 매우 정밀한 결과를 뒤에서 증명할 것이다
(Pro-\'etale Cohomology, Section \ref{proetale-section-weakly-etale} 참조).
이 절에서는 기본적인 내용만 다룬다.

\begin{lemma}
\label{more-morphisms-lemma-check-weakly-etale-stalks}
$f : X \to Y$를 스킴의 사상이라 하자. 다음은 서로 동치이다.
\begin{enumerate}
\item $X \to Y$는 약한 \'etale이다.
\item 모든 $x \in X$에 대하여 환 준동형
$\mathcal{O}_{Y, f(x)} \to \mathcal{O}_{X, x}$는 약한 \'etale이다.
\end{enumerate}
\end{lemma}

\begin{proof}
(1)과 (2)의 어느 가정 아래서도 사상 $f$가 평탄임에 유의하라. 따라서
$f$가 평탄이라고 가정해도 좋다. $x \in X$라 하고 $Y$에서의 상을
$y = f(x)$라 하자. 환의 표준 사상들
$$
\mathcal{O}_{X, x} \otimes_{\mathcal{O}_{Y, y}} \mathcal{O}_{X, x}
\longrightarrow
\mathcal{O}_{X \times_Y X, \Delta_{X/Y}(x)}
\longrightarrow
\mathcal{O}_{X, x}
$$
이 존재한다. 첫째 사상은 국소화(따라서 평탄)이고 둘째 사상은 전사이다.
조건 (1)은 모든 $x$에 대하여 둘째 화살표가 평탄이라는 뜻이다. 조건 (2)는
모든 $x$에 대하여 그 합성이 평탄이라는 뜻이다.
% P05-EN-CH37-0317
따라서 동치는 Algebra, Lemma \ref{algebra-lemma-flat-localization}의 (2)에서
따른다.
\end{proof}

\begin{lemma}
\label{more-morphisms-lemma-key}
$X \to Y$를 $X \to X \times_Y X$가 평탄인 스킴의 사상이라 하자.
$\mathcal{F}$를 $\mathcal{O}_X$-가군이라 하자. $\mathcal{F}$가 $Y$ 위에서
평탄이면 $\mathcal{F}$는 $X$ 위에서 평탄이다.
\end{lemma}

\begin{proof}
% P05-EN-CH37-0386
$x \in X$라 하고 $Y$에서의 상을 $y = f(x)$라 하자.
$X \to X \times_Y X$가 평탄이므로
$\mathcal{O}_{X, x} \otimes_{\mathcal{O}_{Y, y}} \mathcal{O}_{X, x} \to
\mathcal{O}_{X, x}$가 평탄임을 알 수 있다. 따라서 결과는 More on Algebra,
Lemma \ref{more-algebra-lemma-key}와 정의들에서 따른다.
\end{proof}

\begin{lemma}
\label{more-morphisms-lemma-weakly-etale-characterize}
$f : X \to S$를 스킴의 사상이라 하자. 다음은 서로 동치이다.
\begin{enumerate}
\item 사상 $f$는 약한 \'etale이다.
% P05-EN-CH37-0315
\item $f(U) \subset V$인 모든 아핀 열린부분스킴 쌍 $U \subset X$,
$V \subset S$에 대하여 환 준동형
$\mathcal{O}_S(V) \to \mathcal{O}_X(U)$는 약한 \'etale이다.
\item 열린 덮개 $S = \bigcup_{j \in J} V_j$와 열린 덮개들
$f^{-1}(V_j) = \bigcup_{i \in I_j} U_i$가 존재하여 각 사상
$U_i \to V_j$, $j\in J, i\in I_j$가 약한 \'etale이다.
\item 아핀 열린 덮개 $S = \bigcup_{j \in J} V_j$와 아핀 열린 덮개들
$f^{-1}(V_j) = \bigcup_{i \in I_j} U_i$가 존재하여 모든
$j\in J, i\in I_j$에 대하여 환 준동형
$\mathcal{O}_S(V_j) \to \mathcal{O}_X(U_i)$가 약한 \'etale이다.
\end{enumerate}
더 나아가 $f$가 약한 \'etale이면, $f(U) \subset V$인 임의의 열린부분스킴
$U \subset X$, $V \subset S$에 대하여 제한 $f|_U : U \to V$는 약한
\'etale이다.
\end{lemma}

\begin{proof}
% P05-EN-CH37-0385
$f(U) \subset V$인 열린부분스킴 $U \subset X$, $V \subset S$가 주어졌다고
하자. 그러면 $U \times_V U \subset X \times_Y X$는 열려 있고(Schemes,
Lemma \ref{schemes-lemma-open-fibre-product} 참조), $f|_U : U \to V$의 대각
$\Delta_{U/V}$는 제한 $\Delta_{X/Y}|_U : U \to U \times_V U$이다. 평탄성은
스킴 사상의 국소 성질이므로(Morphisms, Lemma
\ref{morphisms-lemma-flat-characterize} 참조),
% P05-EN-CH37-0316
보조정리의 마지막 명제와 (1)과 (3)의 동치가 따른다. $X$와 $Y$가
아핀이면, $X \to Y$가 약한 \'etale일 필요충분조건은
$\mathcal{O}_Y(Y) \to \mathcal{O}_X(X)$가 약한 \'etale인 것이다(다시
Morphisms, Lemma \ref{morphisms-lemma-flat-characterize}를 사용하라).
따라서 (1)과 (3)은 각각 (2)와 (4)에도 동치이다.
\end{proof}

\begin{lemma}
\label{more-morphisms-lemma-composition-weakly-etale}
$X \to Y \to Z$를 스킴의 사상들이라 하자.
\begin{enumerate}
\item $X \to X \times_Y X$와 $Y \to Y \times_Z Y$가 평탄이면
$X \to X \times_Z X$도 평탄이다.
\item $X \to Y$와 $Y \to Z$가 약한 \'etale이면 $X \to Z$도 약한
\'etale이다.
\end{enumerate}
\end{lemma}

\begin{proof}
(1)은 $Z$ 위의 $X$의 대각 사상의 분해
$$
X \to X \times_Y X \to X \times_Z X
$$
와 등식
$$
X \times_Y X = (X \times_Z X) \times_{(Y \times_Z Y)} Y,
$$
그리고 평탄 사상의 밑변환이 평탄이며 평탄 사상들의 합성이 평탄이라는
사실에서 따른다(Morphisms, Lemmas
\ref{morphisms-lemma-base-change-flat} and
\ref{morphisms-lemma-composition-flat}). (2)는 (1)과 방금 사용한 평탄
사상들의 합성이 평탄이라는 사실에서 따른다.
\end{proof}

\begin{lemma}
\label{more-morphisms-lemma-base-change-weakly-etale}
$X \to Y$와 $Y' \to Y$를 스킴의 사상들이라 하고
$X' = Y' \times_Y X$를 $X$의 밑변환이라 하자.
\begin{enumerate}
\item $X \to X \times_Y X$가 평탄이면 $X' \to X' \times_{Y'} X'$도
평탄이다.
\item $X \to Y$가 약한 \'etale이면 $X' \to Y'$도 약한 \'etale이다.
\end{enumerate}
\end{lemma}

\begin{proof}
$X \to X \times_Y X$가 평탄이라고 가정하자. 사상
$X' \to X' \times_{Y'} X'$는 $Y' \to Y$에 의한
$X \to X \times_Y X$의 밑변환이다. 따라서 Morphisms, Lemmas
\ref{morphisms-lemma-base-change-flat}에 의해 평탄이다. 이로써 (1)이
증명된다. (2)는 (1)과 방금 사용한 평탄 사상의 밑변환이 평탄이라는
사실에서 따른다.
\end{proof}

\begin{lemma}
\label{more-morphisms-lemma-go-down}
$X \to Y \to Z$를 스킴의 사상들이라 하자. $X \to Y$가 평탄이고
전사이며 $X \to X \times_Z X$가 평탄이라고 가정하자. 그러면
$Y \to Y \times_Z Y$는 평탄이다.
\end{lemma}

\begin{proof}
가환 그림
$$
\xymatrix{
X \ar[r] \ar[d] & X \times_Z X \ar[d] \\
Y \ar[r] & Y \times_Z Y
}
$$
을 생각하자. 위 가로 화살표와 세로 화살표들은 평탄이다. 따라서 $X$는
$Y \times_Z Y$ 위에서 평탄이다. Morphisms, Lemma
\ref{morphisms-lemma-flat-permanence}에 의해 $Y$가 $Y \times_Z Y$ 위에서
평탄임을 알 수 있다.
\end{proof}

\begin{lemma}
\label{more-morphisms-lemma-weakly-etale-formally-unramified}
$f : X \to Y$를 스킴의 약한 \'etale 사상이라 하자. 그러면 $f$는
형식적으로 비분기이다. 즉, $\Omega_{X/Y} = 0$이다.
\end{lemma}

\begin{proof}
Lemma \ref{more-morphisms-lemma-formally-unramified-differentials}에 의해 $f$가 형식적으로
비분기일 필요충분조건은 $\Omega_{X/Y} = 0$임을 상기하라. Lemma
\ref{more-morphisms-lemma-weakly-etale-characterize}와 Morphisms, Lemma
\ref{morphisms-lemma-differentials-affine}를 통하여 이는 환의 경우인 More on
Algebra, Lemma \ref{more-algebra-lemma-formally-unramified}에서 따른다.
\end{proof}

\begin{lemma}
\label{more-morphisms-lemma-when-weakly-etale}
$f : X \to Y$를 스킴의 사상이라 하자. 다음 각 경우에 $X \to Y$는 약한
\'etale이다.
\begin{enumerate}
\item $X \to Y$는 평탄 모노사상이다.
\item $X \to Y$는 열린 몰입이다.
\item $X \to Y$는 평탄이고 비분기이다.
\item $X \to Y$는 \'etale이다.
\end{enumerate}
\end{lemma}

\begin{proof}
(1)이 성립하면 $\Delta_{X/Y}$는 동형이므로(Schemes, Lemma
\ref{schemes-lemma-monomorphism}), 당연히 $f$는 약한 \'etale이다. (2)는
(1)의 특수한 경우이다. 비분기 사상의 대각 사상은 열린 몰입이므로
(Morphisms, Lemma \ref{morphisms-lemma-diagonal-unramified-morphism}),
평탄이다. 따라서 평탄 비분기 사상은 약한 \'etale이다. \'Etale 사상은
평탄이고 비분기이므로(Morphisms, Lemma
\ref{morphisms-lemma-etale-smooth-unramified}), (4)는 (3)에서 따른다.
\end{proof}

\begin{lemma}
\label{more-morphisms-lemma-reduced-goes-up}
$f : X \to Y$를 스킴의 사상이라 하자. $Y$가 축약이고 $f$가 약한
\'etale이면 $X$도 축약이다.
\end{lemma}

\begin{proof}
Lemma \ref{more-morphisms-lemma-weakly-etale-characterize}를 통하여 이는 환의 경우인
More on Algebra, Lemma
\ref{more-algebra-lemma-absolutely-flat-over-absolutely-flat}에서 따른다.
\end{proof}

\noindent
다음 보조정리는 약한 \'etale 환 준동형에 관한 자명하지 않은 결과를
사용한다.

\begin{lemma}
\label{more-morphisms-lemma-weakly-etale-strictly-henselian-local-rings}
$f : X \to Y$를 스킴의 사상이라 하자. 다음은 서로 동치이다.
\begin{enumerate}
\item $f$는 약한 \'etale이다.
\item $x \in X$에 대하여 국소환 준동형
$\mathcal{O}_{Y, f(x)} \to \mathcal{O}_{X, x}$는 강한 헨젤화들 사이의
동형을 유도한다.
\end{enumerate}
\end{lemma}

\begin{proof}
$x \in X$를 점이라 하고 $Y$에서의 상을 $y = f(x)$라 하자.
$\kappa(x)$의 분리 대수적 폐포 $\kappa^{sep}$를 택하자.
$\mathcal{O}_{X, x}^{sh}$를 $\kappa^{sep}$에 대응하는 강한 헨젤화라 하고,
$\mathcal{O}_{Y, y}^{sh}$를 $\kappa^{sep}$ 안의 $\kappa(y)$의 분리 대수적
폐포에 대한 강한 헨젤화라 하자. 가환 그림
$$
\xymatrix{
\mathcal{O}_{X, x} \ar[r] & \mathcal{O}_{X, x}^{sh} \\
\mathcal{O}_{Y, y} \ar[u] \ar[r] & \mathcal{O}_{Y, y}^{sh} \ar[u]
}
$$
은 국소환들의 국소 준동형들로 이루어진다. Algebra, Lemma
\ref{algebra-lemma-strictly-henselian-functorial}을 보라. 강한 헨젤화는
\'etale 환 준동형들의 여과 쌍극한이므로, More on Algebra, Lemma
\ref{more-algebra-lemma-when-weakly-etale}에 의해 가로 사상들은 약한
\'etale이다. 더 나아가 More on Algebra, Lemma
\ref{more-algebra-lemma-dumb-properties-henselization}에 의해 가로 사상들은
충실히 평탄이다.

\medskip\noindent
$f$가 약한 \'etale이라고 가정하자. Lemma
\ref{more-morphisms-lemma-check-weakly-etale-stalks}에 의해 왼쪽 세로 화살표는 약한
\'etale이다. More on Algebra, Lemmas
\ref{more-algebra-lemma-composition-weakly-etale} and
\ref{more-algebra-lemma-weakly-etale-permanence}에 의해 오른쪽 세로 화살표도
약한 \'etale이다. More on Algebra, Theorem
\ref{more-algebra-theorem-olivier}에 의해 오른쪽 세로 사상이 동형이라는
결론을 얻는다.

\medskip\noindent
$\mathcal{O}_{Y, y}^{sh} \to \mathcal{O}_{X, x}^{sh}$가 동형이라고
가정하자. 그러면 $\mathcal{O}_{Y, y} \to \mathcal{O}_{X, x}^{sh}$는 약한
\'etale이다. $\mathcal{O}_{X, x} \to \mathcal{O}_{X, x}^{sh}$가 충실히
평탄이므로 More on Algebra, Lemma \ref{more-algebra-lemma-go-down}에 의해
$\mathcal{O}_{Y, y} \to \mathcal{O}_{X, x}$가 약한 \'etale이라는 결론을
얻는다. 따라서 Lemma \ref{more-morphisms-lemma-check-weakly-etale-stalks}에 의해 (2)는
(1)을 함의한다.
\end{proof}

\begin{lemma}
\label{more-morphisms-lemma-normal-goes-up}
$f : X \to Y$를 스킴의 사상이라 하자. $Y$가 정규 스킴이고 $f$가 약한
\'etale이면 $X$도 정규 스킴이다.
\end{lemma}

\begin{proof}
More on Algebra, Lemma \ref{more-algebra-lemma-henselization-normal}에 의해
스킴 $S$가 정규일 필요충분조건은 모든 $s \in S$에 대하여
$\mathcal{O}_{S, s}$의 강한 헨젤화가 정규 정역인 것이다. 따라서
보조정리는 Lemma
\ref{more-morphisms-lemma-weakly-etale-strictly-henselian-local-rings}에서 따른다.
\end{proof}

\begin{lemma}
\label{more-morphisms-lemma-weakly-etale-permanence}
$S$를 스킴이라 하고 $f : X \to Y$를 $S$ 위의 스킴들의 사상이라 하자.
$X$, $Y$가 $S$ 위에서 약한 \'etale이면 $f$도 약한 \'etale이다.
\end{lemma}

\begin{proof}
Morphisms, Lemmas \ref{morphisms-lemma-base-change-flat} and
\ref{morphisms-lemma-composition-flat}을 더 언급하지 않고 사용한다.
$X \to Y$를 합성 $X \to X \times_S Y \to Y$로 쓰자. 둘째 사상은 평탄
사상 $X \to S$의 밑변환이므로 평탄이다. 첫째 사상은 평탄 사상
$Y \to Y \times_S Y$를 사상 $X \times_S Y \to Y \times_S Y$에 의해
밑변환한 것이므로 평탄이다. 따라서 $X \to Y$는 평탄이다. 사상
$X \times_Y X \to X \times_S X$는 몰입이다.
% P05-EN-CH37-0318
그러므로 Lemma \ref{more-morphisms-lemma-key}에 의해, $X$가 $X \times_S X$ 위에서
평탄이므로 $X$는 $X \times_Y X$ 위에서도 평탄이다.
\end{proof}

\noindent
다음은 동시에 분리이고 순수 비분리인 체 확대는 동형이어야 한다는 관찰을
스킴 이론적으로 일반화한 것이다.

\begin{lemma}
\label{more-morphisms-lemma-weakly-etale-universal-homeomorphism}
$f : X \to Y$를 스킴의 사상이라 하자. $f$가 약한 \'etale이고 보편
위상동형이면 동형이다.
\end{lemma}

\begin{proof}
$f$가 보편 위상동형이므로 Morphisms, Lemmas
\ref{morphisms-lemma-homeomorphism-affine} and
\ref{morphisms-lemma-universally-injective}에 의해 대각 사상
$\Delta : X \to X \times_Y X$는 전사인 닫힌 몰입이다. $\Delta$는 또한
평탄이므로 Morphisms, Lemma
\ref{morphisms-lemma-characterize-flat-closed-immersions}에 의해 $\Delta$는
동형이어야 한다. 달리 말해 $f$는 모노사상이다(Schemes, Lemma
\ref{schemes-lemma-monomorphism}). $f$는 보편 위상동형이므로 당연히
준콤팩트이다. 따라서 Descent, Lemma
\ref{descent-lemma-flat-surjective-quasi-compact-monomorphism-isomorphism}에
의해 $f$가 동형임을 알 수 있다.
\end{proof}

\noindent
다음은 \'Etale Morphisms, Lemma
\ref{etale-lemma-relative-frobenius-etale}의 약한 \'etale 일반화이다.

\begin{lemma}
\label{more-morphisms-lemma-relative-frobenius-weakly-etale}
$U \to X$를 스킴의 약한 \'etale 사상이라 하고 $X$는 표수 $p$인 스킴이라
하자. 그러면 상대 Frobenius
$F_{U/X} : U \to U \times_{X, F_X} X$는 동형이다.
\end{lemma}

\begin{proof}
Varieties, Lemma \ref{varieties-lemma-relative-frobenius}에 의해 사상
$F_{U/X}$는 보편 위상동형이다. Lemma
\ref{more-morphisms-lemma-weakly-etale-permanence}에 의해 $F_{U/X}$는 $X$ 위에서 약한
\'etale인 스킴들 사이의 사상으로서 약한 \'etale이다. 따라서 Lemma
\ref{more-morphisms-lemma-weakly-etale-universal-homeomorphism}에 의해 $F_{U/X}$는
동형이다.
\end{proof}






\section{축약 올 정리}
\label{more-morphisms-section-reduced-fibre-theorem}

\noindent
이 절에서는 제목이 예고하는 유형 중 가장 단순한 정리를 논한다. 결과의
증명은 다소 번거롭지만, 본질적으로는 분기 제거에 관한 Epp의 결과에서
직접적으로 따른다. More on Algebra, Theorem
\ref{more-algebra-theorem-epp}를 보라.

\medskip\noindent
$A$를 분수체가 $K$인 Dedekind 정역이라 하고, $X$를 $A$ 위에서 평탄이고
유한형인 스킴이라 하자. $L$을 $K$의 유한 확대체라 하고, $B$를 $L$에서
$A$의 정수적 폐포라 하자. 그러면 $B$는 Dedekind 정역이다(Algebra,
Lemma \ref{algebra-lemma-integral-closure-Dedekind}).
$X_B = X \times_{\Spec(A)} \Spec(B)$를 밑변환이라 하자. 그러면
$X_B \to \Spec(B)$는 유한형이고(Morphisms, Lemma
\ref{morphisms-lemma-base-change-finite-type}), 따라서 $X_B$는 Noether이다
(Morphisms, Lemma \ref{morphisms-lemma-finite-type-noetherian}). 그러므로
정규화 $\nu : Y \to X_B$가 존재한다(Morphisms, Definition
\ref{morphisms-definition-normalization} 및 그 뒤의 논의 참조). 그림은
\begin{equation}
\label{more-morphisms-equation-normalized-base-change}
\xymatrix{
Y \ar[rd] \ar[r]_\nu & X_B \ar[r] \ar[d] & X \ar[d] \\
 & \Spec(B) \ar[r] & \Spec(A)
}
\end{equation}
와 같다. 때때로 $Y$를 $X$의 {\it 정규화된 밑변환}이라고 한다. 일반적으로
사상 $\nu$는 유한이 아닐 수 있다. 그러나 $A$가 Nagata 환이면(실제로는
거의 언제나 만족되는 조건이다), $\nu$는 유한이고 $Y$는 $B$ 위에서
유한형이다. Morphisms, Lemmas
\ref{morphisms-lemma-nagata-normalization} and
\ref{morphisms-lemma-finite-type-nagata}를 보라.

\medskip\noindent
정규화된 밑변환을 취하는 것은 합성과 가환한다. 더 정확히, $M/L/K$가
정수적 폐포 $A \subset B \subset C$를 갖는 체들의 유한 확대열이면,
$M/L$에 대한 $Y \to \Spec(B)$의 정규화된 밑변환 $Z$는 $M/K$에 대한
$X \to \Spec(A)$의 정규화된 밑변환과 같다.

\begin{theorem}
\label{more-morphisms-theorem-normalized-base-change-with-reduced-fibre}
$A$를 분수체가 $K$인 Dedekind 환이라 하자. $X$를 $A$ 위에서 평탄이고
유한형인 스킴이라 하자. $A$가 Nagata 환이라고 가정하자. 모든 올의 모든
일반점에서 정규화된 밑변환 $Y$가 $\Spec(B)$ 위에서 매끄럽게 되는 유한
확대 $L/K$가 존재한다.
\end{theorem}

\begin{proof}
증명 중에 $X \to \Spec(A)$와 같은 (평탄이고 유한 표시인) 사상이 매끄러운
점들의 집합을 구성하는 것이 밑변환과 가환한다는 사실을 반복해서
사용한다. Morphisms, Lemma
\ref{morphisms-lemma-set-points-where-fibres-smooth}를 보라.

\medskip\noindent
먼저 $(X_L)_{red}$가 $L$ 위에서 기하학적으로 축약이 되도록 유한 확대
$L/K$를 택한다. Varieties, Lemma
\ref{varieties-lemma-finite-extension-geometrically-reduced}를 보라.
$Y \to (X_B)_{red}$가 쌍유리이므로 Varieties, Lemma
\ref{varieties-lemma-generic-points-geometrically-reduced}를 적용하면
$Y_L$도 $L$ 위에서 기하학적으로 축약임을 알 수 있다. 따라서 Varieties,
Lemma \ref{varieties-lemma-geometrically-reduced-dense-smooth-open}에 의해
$Y_L \to \Spec(L)$는 조밀한 열린집합 $V \subset Y_L$ 위에서 매끄럽다.
그러므로 사상 $Y \to \Spec(B)$의 매끄러운 자취 $U \subset Y$는
열려 있고(Morphisms, Definition \ref{morphisms-definition-smooth}), 일반올에서
조밀하다. $A$를 $B$로, $X$를 $Y$로 바꾸면 다음 문단에서 다루는 경우로
환원된다.

\medskip\noindent
$X$가 정규이고 $X \to \Spec(A)$의 매끄러운 자취 $U \subset X$가
일반올에서 조밀하다고 가정하자. 그러면 $U$는 유한 개를 제외한 모든
올에서 조밀하다. Lemma \ref{more-morphisms-lemma-nowhere-dense-generic-fibre}를 보라.
$x_1, \ldots, x_r \in X \setminus U$를 $X \setminus U$의 기약 성분들의
일반점이면서 $X \to \Spec(A)$의 올들의 기약 성분의 일반점인 유한 개의
점들이라 하자. $\mathcal{O}_i = \mathcal{O}_{X, x_i}$라 놓고, $A_i$를
$\Spec(A)$에서 $x_i$의 상에 대응하는 극대 아이디얼에서 $A$를 국소화한
환이라 하자. More on Algebra, Proposition
\ref{more-algebra-proposition-epp-essentially-finite-type}에 의해 이산
값매김환의 확대 $A_i \to \mathcal{O}_i$에 대한 해인 유한 확대들 $K_i/K$가
존재한다. 모든 확대 $K_i/K$를 지배하는 유한 확대 $L/K$를 택하자. 그러면
More on Algebra, Lemma \ref{more-algebra-lemma-solution-goes-up}에 의해
$L/K$는 여전히 $A_i \to \mathcal{O}_i$에 대한 해이다.

\medskip\noindent
방금 만든 확대 $L/K$에 대한 그림
(\ref{more-morphisms-equation-normalized-base-change})을 생각하자. $U_B \subset X_B$는
$B$ 위에서 매끄럽고, 따라서 정규이다(예를 들어 Algebra, Lemma
\ref{algebra-lemma-normal-goes-up}를 사용하라). 그러므로 $Y \to X_B$는
$U_B$ 위에서 동형이다. $y \in Y$를 극대 아이디얼
$\mathfrak m \subset B$ 위에 놓인 $Y \to \Spec(B)$의 어떤 올의 기약
성분의 일반점이라 하자. $y \not \in U_B$라고 가정하자. 그러면 $y$는 점
$x_i$ 중 하나로 간다. 따라서 $\mathcal{O}_{Y, y}$는
$R(X) \otimes_K L$에서 $\mathcal{O}_i$의 정수적 폐포의 국소환이다
(세부사항은 생략한다). $L/K$가 $A_i \to \mathcal{O}_i$에 대한 해이므로
$B_\mathfrak m \to \mathcal{O}_{Y, y}$가 $\mathfrak m_y$-진 위상에서
형식적으로 매끄러움을 알 수 있다(이것이 ``해''라는 것의 정의이다).
달리 말해 $\mathfrak m\mathcal{O}_{Y, y} = \mathfrak m_y$이고 잉여체
확대는 분리이다. More on Algebra, Lemma
\ref{more-algebra-lemma-extension-dvrs-formally-smooth}를 보라. 따라서
$y$에서 올의 국소환은 $\kappa(y)$이다. 예를 들어 Algebra, Lemma
\ref{algebra-lemma-separable-smooth}에 의해 이는 올이 $y$에서
$\kappa(\mathfrak m)$ 위에서 매끄러움을 함의한다. 이로써 증명이 끝난다.
\end{proof}

\begin{lemma}[곡선 위의 변형]
\label{more-morphisms-lemma-normalized-base-change-with-reduced-fibre-over-curve}
$f : X \to S$를 평탄이고 유한형인 스킴 사상이라 하자. $S$는 Nagata이고
함수체가 $K$인 정역 스킴이며 차원이 $1$인 정칙 스킴이라고 가정하자.
그러면 그림
$$
\xymatrix{
Y \ar[rd]_g \ar[r]_-\nu & X \times_S T \ar[d] \ar[r] & X \ar[d]_f \\
& T \ar[r] & S
}
$$
에서 사상 $g$가 올들의 모든 일반점에서 매끄럽게 되는 유한 확대 $L/K$가
존재한다. 여기서 $T$는 $\Spec(L)$ 안에서 $S$의 정규화이고,
$\nu : Y \to X \times_S T$는 정규화이다.
\end{lemma}

\begin{proof}
유한 아핀 열린 덮개 $S = \bigcup \Spec(A_i)$를 택하자. 그러면 모든 $i$에
대하여 $K$는 $A_i$의 분수체와 같다. $X_i = X \times_S \Spec(A_i)$라
놓자. 사상 $X_i \to \Spec(A_i)$에 대하여 Theorem
\ref{more-morphisms-theorem-normalized-base-change-with-reduced-fibre}에서와 같은 $L_i/K$를
택하자. $B_i \subset L_i$를 $A_i$의 정수적 폐포라 하고, $Y_i$를 $B_i$로의
$X$의 정규화된 밑변환이라 하자. 각 $L_i$를 지배하는 유한 확대 $L/K$를
택하자. $T_i \subset T$를 $\Spec(A_i)$의 역상이라 하자. 각 $i$에 대하여
가환 그림
$$
\xymatrix{
g^{-1}(T_i) \ar[r] \ar[d] & Y_i \ar[r] \ar[d] & X \times_S \Spec(A_i) \ar[d] \\
T_i \ar[r] & \Spec(B_i) \ar[r] & \Spec(A_i)
}
$$
을 얻고, 실제로 왼쪽 정사각형은 절의 시작에서 논한 정규화된 밑변환이다.
Theorem \ref{more-morphisms-theorem-normalized-base-change-with-reduced-fibre}의 증명에서
$Y \to T$의 매끄러운 자취가 $Y_i$가 $B_i$ 위에서 매끄러운 점들의 집합의
$g^{-1}(T_i)$ 안에서의 역상을 포함함을 보았다. 이로써 보조정리가
증명된다.
\end{proof}

\begin{lemma}[분리 확대를 갖는 변형]
\label{more-morphisms-lemma-normalized-base-change-with-reduced-fibre-separable}
$A$를 분수체가 $K$인 Dedekind 환이라 하고, $X$를 $A$ 위에서 평탄이고
유한형인 스킴이라 하자. $A$가 Nagata 환이고, $X$의 기약 성분의 모든
일반점 $\eta$에 대하여 체 확대 $\kappa(\eta)/K$가 분리라고 가정하자.
그러면 모든 올의 모든 일반점에서 정규화된 밑변환 $Y$가 $\Spec(B)$ 위에서
매끄럽게 되는 유한 분리 확대 $L/K$가 존재한다.
\end{lemma}

\begin{proof}
몇 가지 사소한 변경을 제외하면 Theorem
\ref{more-morphisms-theorem-normalized-base-change-with-reduced-fibre}과 정확히 같은 방식으로
증명한다. 가장 중요한 변경은 More on Algebra, Proposition
\ref{more-algebra-proposition-epp-essentially-finite-type} 대신 More on
Algebra, Lemma
\ref{more-algebra-lemma-epp-essentially-finite-type-separable}를 사용하는
것이다. 증명 중에 $X \to \Spec(A)$와 같은 (평탄이고 유한 표시인) 사상이
매끄러운 점들의 집합을 구성하는 것이 밑변환과 가환한다는 사실을 반복해서
사용한다. Morphisms, Lemma
\ref{morphisms-lemma-set-points-where-fibres-smooth}를 보라.

\medskip\noindent
$X$가 $A$ 위에서 평탄이므로 $X$의 모든 일반점 $\eta$는 $\Spec(A)$의
일반점으로 간다. $X$를 그 축약으로 바꾸어 $X$가 축약이라고 가정해도
좋다. 이 경우 Varieties, Lemma
\ref{varieties-lemma-generic-points-geometrically-reduced}에 의해 $X_K$는
$K$ 위에서 기하학적으로 축약이다. 따라서 Varieties, Lemma
\ref{varieties-lemma-geometrically-reduced-dense-smooth-open}에 의해
$X_K \to \Spec(K)$는 어떤 조밀한 열린집합 위에서 매끄럽다. 그러므로 사상
$X \to \Spec(A)$의 매끄러운 자취 $U \subset X$는 열려 있고(Morphisms,
Definition \ref{morphisms-definition-smooth}), 일반올에서 조밀하다. 이로써
다음 문단의 상황으로 환원된다.

\medskip\noindent
$X$가 정규이고 $X \to \Spec(A)$의 매끄러운 자취 $U \subset X$가
일반올에서 조밀하다고 가정하자. 그러면 $U$는 유한 개를 제외한 모든
올에서 조밀하다. Lemma \ref{more-morphisms-lemma-nowhere-dense-generic-fibre}를 보라.
$x_1, \ldots, x_r \in X \setminus U$를 $X \setminus U$의 기약 성분들의
일반점이면서 $X \to \Spec(A)$의 올들의 기약 성분의 일반점인 유한 개의
점들이라 하자. $\mathcal{O}_i = \mathcal{O}_{X, x_i}$라 놓자.
$\mathcal{O}_i$의 분수체는 $X$의 어떤 일반점의 잉여체임에 유의하라.
$A_i$를 $\Spec(A)$에서 $x_i$의 상에 대응하는 극대 아이디얼에서 $A$를
국소화한 환이라 하자. More on Algebra, Lemma
\ref{more-algebra-lemma-epp-essentially-finite-type-separable}를 적용하면
$A_i \to \mathcal{O}_i$에 대한 해인 유한 분리 확대들 $K_i/K$를 얻는다.
모든 확대 $K_i/K$를 지배하는 유한 분리 확대 $L/K$를 택하자. 그러면
More on Algebra, Lemma \ref{more-algebra-lemma-solution-goes-up}에 의해 $L/K$는
여전히 $A_i \to \mathcal{O}_i$에 대한 해이다.

\medskip\noindent
방금 만든 확대 $L/K$에 대한 그림
(\ref{more-morphisms-equation-normalized-base-change})을 생각하자. $U_B \subset X_B$는
$B$ 위에서 매끄럽고, 따라서 정규이다(예를 들어 Algebra, Lemma
\ref{algebra-lemma-normal-goes-up}를 사용하라). 그러므로 $Y \to X_B$는
$U_B$ 위에서 동형이다. $y \in Y$를 극대 아이디얼
$\mathfrak m \subset B$ 위에 놓인 $Y \to \Spec(B)$의 어떤 올의 기약
성분의 일반점이라 하자. $y \not \in U_B$라고 가정하자. 그러면 $y$는 점
$x_i$ 중 하나로 간다. 따라서 $\mathcal{O}_{Y, y}$는
$R(X) \otimes_K L$에서 $\mathcal{O}_i$의 정수적 폐포의 국소환이다
(세부사항은 생략한다). $L/K$가 $A_i \to \mathcal{O}_i$에 대한 해이므로
$B_\mathfrak m \to \mathcal{O}_{Y, y}$가 형식적으로 매끄러움을 알 수
있다(이것이 ``해''라는 것의 정의이다). 달리 말해
$\mathfrak m\mathcal{O}_{Y, y} = \mathfrak m_y$이고 잉여체 확대는 분리이다.
따라서 $y$에서 올의 국소환은 $\kappa(y)$이다. 예를 들어 Algebra, Lemma
\ref{algebra-lemma-separable-smooth}에 의해 이는 올이 $y$에서
$\kappa(\mathfrak m)$ 위에서 매끄러움을 함의한다. 이로써 증명이 끝난다.
\end{proof}

\begin{lemma}[곡선 위의 분리 확대를 갖는 변형]
\label{more-morphisms-lemma-normalized-base-change-with-reduced-fibre-over-curve-separable}
$f : X \to S$를 평탄이고 유한형인 스킴 사상이라 하자. $S$는 Nagata이고
함수체가 $K$인 정역 스킴이며 차원이 $1$인 정칙 스킴이라고 가정하자.
$X$의 기약 성분의 모든 일반점 $\eta$에 대하여 체 확대
$\kappa(\eta)/K$가 분리라고 가정하자. 그러면 그림
$$
\xymatrix{
Y \ar[rd]_g \ar[r]_-\nu & X \times_S T \ar[d] \ar[r] & X \ar[d]_f \\
& T \ar[r] & S
}
$$
에서 사상 $g$가 올들의 모든 일반점에서 매끄럽게 되는 유한 분리 확대
$L/K$가 존재한다. 여기서 $T$는 $\Spec(L)$ 안에서 $S$의 정규화이고,
$\nu : Y \to X \times_S T$는 정규화이다.
\end{lemma}

\begin{proof}
Lemma \ref{more-morphisms-lemma-normalized-base-change-with-reduced-fibre-over-curve}가 Theorem
\ref{more-morphisms-theorem-normalized-base-change-with-reduced-fibre}에서 따르는 것과 정확히
같은 방식으로, 이는 Lemma
\ref{more-morphisms-lemma-normalized-base-change-with-reduced-fibre-separable}에서 따른다.
\end{proof}








\section{Ind-준아핀 사상}
\label{more-morphisms-section-ind-quasi-affine}

\noindent
뒤에서 사용할 약간의 이론이다.

\begin{definition}
\label{more-morphisms-definition-ind-quasi-affine}
스킴 $X$의 모든 준콤팩트 열린집합이 준아핀이면 $X$를
{\it ind-준아핀}이라고 한다. 마찬가지로 스킴의 사상 $X \to Y$에서 $Y$의
각 아핀 열린집합 $V$에 대하여
% P05-EN-CH37-0387
$f^{-1}(V)$가 ind-준아핀이면 이 사상을 {\it ind-준아핀}이라고 한다.
\end{definition}

\noindent
아핀 스킴의 열린부분스킴은 ind-준아핀 스킴의 한 예이다. 준아핀
열린집합들의 합집합 $X = \bigcup_{i \in I} U_i$에서 임의의 두 $U_i$가
셋째 것에 포함되면 $X$는 ind-준아핀이다. Ind-준아핀 스킴 $X$는
분리이다. 실제로 임의의 두 아핀 열린집합 $U, V$는 $X$의 분리인
열린부분스킴 $U \cup V$에 포함된다. 마찬가지로 ind-준아핀 사상은
분리이다.

\begin{lemma}
\label{more-morphisms-lemma-ind-quasi-affine-alternative-definition}
스킴의 사상 $f : X \to Y$에 대하여 다음은 서로 동치이다.
\begin{enumerate}
\item $f$는 ind-준아핀이다.
\item 모든 아핀 열린부분스킴 $V \subset Y$와 모든 준콤팩트 열린부분스킴
$U \subset f^{-1}(V)$에 대하여 유도된 사상 $U \to V$는 준아핀이다.
\item
$Y$의 준콤팩트 준분리 열린부분스킴 $V_j \subset Y$들로 된 어떤 덮개
$\{ V_j \}_{j \in J}$가 존재하여, 모든 $j \in J$와 모든 준콤팩트
열린부분스킴 $U \subset f^{-1}(V_j)$에 대하여 유도된 사상
$U \to V_j$는 준아핀이다.
\item 모든 준콤팩트 준분리 열린부분스킴 $V \subset Y$와 모든 준콤팩트
열린부분스킴 $U \subset f^{-1}(V)$에 대하여 유도된 사상 $U \to V$는
준아핀이다.
\end{enumerate}
특히 ind-준아핀 사상이라는 성질은 밑에서 Zariski 국소적이다.
\end{lemma}

\begin{proof}
(1) $\Leftrightarrow$ (2)는 정의들과 Morphisms, Lemma
\ref{morphisms-lemma-characterize-quasi-affine}에서 따른다. (2)
$\Rightarrow$ (4)에 대하여, $U$와 $V$를 (4)에서와 같이 택하자. Schemes,
Lemma \ref{schemes-lemma-quasi-compact-permanence}에 의해 유도된 사상
$U \to V$는 준콤팩트이다. 따라서 모든 아핀 열린집합 $V' \subset V$에
대하여 올곱 $V' \times_V U$는 준콤팩트이고, (2)에 의해 유도된 사상
$V' \times_V U \to V'$는 준아핀이다. 그러므로 Morphisms, Lemma
\ref{morphisms-lemma-characterize-quasi-affine}에 의해 $U \to V$도
준아핀이다. 이 논증은 (3) $\Rightarrow$ (4)도 준다. 실제로 같은 기호를
유지하면, $V_j$ 중 하나에 놓이는 아핀 열린집합들 $V' \subset V$가 $V$를
덮으므로 준콤팩트 사상 $V' \times_V U \to V'$가 준아핀임을 보이면 된다.
그런데 (3)에 의해 합성 $V' \times_V U \to V' \to V_j$는 준아핀이고,
Schemes, Lemma \ref{schemes-lemma-compose-after-separated}에 의해 사상
$V' \to V_j$는 준분리이다. 따라서 Morphisms, Lemma
\ref{morphisms-lemma-quasi-affine-permanence}에 의해
$V' \times_V U \to V'$는 준아핀이다. 마지막 함의들 (4) $\Rightarrow$ (2)와
(4) $\Rightarrow$ (3)은 명백하다.
\end{proof}

\begin{lemma}
\label{more-morphisms-lemma-ind-quasi-affine-composition}
Ind-준아핀 사상이라는 성질은 합성에 대해 안정적이다.
\end{lemma}

\begin{proof}
$f : X \to Y$와 $g : Y \to Z$를 ind-준아핀 사상들이라 하자.
$V \subset Z$와 $U \subset f^{-1}(g^{-1}(V))$를 준콤팩트 열린집합들이라
하고, $V$는 또한 준분리라고 하자. 상 $f(U)$는 $g^{-1}(V)$의 준콤팩트
부분집합이므로 어떤 준콤팩트 열린집합 $W \subset g^{-1}(V)$에 포함된다
(유한 개의 아핀들의 합집합을 취하면 된다). 분해 $U \to W \to V$를
얻는다. Lemma \ref{more-morphisms-lemma-ind-quasi-affine-alternative-definition}에 의해 사상
$W \to V$는 준아핀이므로, 특히 $W$는 준분리이다. 다시 Lemma
\ref{more-morphisms-lemma-ind-quasi-affine-alternative-definition}에 의해 $U \to W$도
준아핀이다. 따라서 Morphisms, Lemma
\ref{morphisms-lemma-composition-quasi-affine}에 의해 합성 $U \to V$도
준아핀이고, Lemma \ref{more-morphisms-lemma-ind-quasi-affine-alternative-definition}을 한
번 더 적용하면 된다.
\end{proof}

\begin{lemma}
\label{more-morphisms-lemma-ind-quasi-affine-examples}
모든 준아핀 사상은 ind-준아핀이다. 모든 몰입은 ind-준아핀이다.
\end{lemma}

\begin{proof}
첫째 주장은 정의들에서 즉시 따른다. 특히 닫힌 몰입과 같은 아핀 사상은
ind-준아핀이다. 따라서 Lemma \ref{more-morphisms-lemma-ind-quasi-affine-composition}에 의해
열린 몰입이 ind-준아핀임을 보이면 된다. 이는 정의들에서 즉시 따른다.
\end{proof}

\begin{lemma}
\label{more-morphisms-lemma-ind-quasi-affine-permanence}
$f : X \to Y$와 $g : Y \to Z$가 스킴의 사상들이고 $g \circ f$가
ind-준아핀이면 $f$도 ind-준아핀이다.
\end{lemma}

\begin{proof}
Lemma \ref{more-morphisms-lemma-ind-quasi-affine-alternative-definition}에 의해 먼저 $Z$
위에서, 그다음 $Y$ 위에서 Zariski 국소적으로 작업해도 된다. 따라서
$Z$가 아핀이고 이어서 $Y$도 아핀이라고 가정해도 일반성을 잃지 않는다.
그러면 $X$의 임의의 준콤팩트 열린집합은 준아핀이므로, 주장은 Lemma
\ref{more-morphisms-lemma-ind-quasi-affine-alternative-definition}에서 따른다.
\end{proof}

\begin{lemma}
\label{more-morphisms-lemma-base-change-ind-quasi-affine}
Ind-준아핀이라는 성질은 밑변환에 대해 안정적이다.
\end{lemma}

\begin{proof}
$f : X \to Y$를 ind-준아핀 사상이라 하자. $f$의 모든 밑변환이
ind-준아핀인지 확인하기 위해 Lemma
\ref{more-morphisms-lemma-ind-quasi-affine-alternative-definition}에 의해 $Y$ 위에서
Zariski 국소적으로 작업해도 된다. 따라서 $Y$가 아핀이라고 가정하자.
더 나아가 밑변환 사상 $Z \to Y$에서 스킴 $Z$도 아핀이라고 가정해도
된다. 밑변환 $X \times_Y Z \to X$는 아핀 사상이므로 Lemmas
\ref{more-morphisms-lemma-ind-quasi-affine-composition}와
\ref{more-morphisms-lemma-ind-quasi-affine-examples}에 의해 사상
$X \times_Y Z \to Y$는 ind-준아핀이다. 그러면 Lemma
\ref{more-morphisms-lemma-ind-quasi-affine-permanence}에 의해 밑변환
$X \times_Y Z \to Z$는 원하는 대로 ind-준아핀이다.
\end{proof}

\begin{lemma}
\label{more-morphisms-lemma-descending-property-ind-quasi-affine}
Ind-준아핀이라는 성질은 밑에서 fpqc 국소적이다.
\end{lemma}

\begin{proof}
Lemma \ref{more-morphisms-lemma-base-change-ind-quasi-affine}가 주는 밑변환에 대한
ind-준아핀성의 안정성이 한 방향을 준다. 역방향에 대하여, $f : X \to Y$를
스킴의 사상이라 하고 $\{g_i : Y_i \to Y\}$를 모든 $i$에 대해 밑변환
$f_i : X_i \to Y_i$가 ind-준아핀인 fpqc 덮개라 하자. $f$가 ind-준아핀임을
보여야 한다.

\medskip\noindent
Lemma \ref{more-morphisms-lemma-ind-quasi-affine-alternative-definition}에 의해 $Y$ 위에서
Zariski 국소적으로 작업해도 되므로 $Y$가 아핀이라고 가정하자. 이어서
Lemma \ref{more-morphisms-lemma-base-change-ind-quasi-affine}가 보장하는 밑변환에 대한
안정성을 사용하여 덮개를 세분하고, 하나의 아핀 충실히 평탄한 사상
$g : Y' \to Y$로 주어진다고 가정하자. 임의의 준콤팩트 열린집합
$U \subset X$에 대하여 그 $Y'$-밑변환
$U \times_Y Y' \subset X \times_Y Y'$도 준콤팩트이다. 이제 Descent,
Lemma \ref{descent-lemma-descending-property-quasi-affine}에 의해 사상
$U \to Y$가 준아핀일 필요충분조건은 $U \times_Y Y' \to Y'$가
준아핀이라는 것임을 관찰하면 된다.
\end{proof}

\begin{lemma}
\label{more-morphisms-lemma-etale-separated-ind-quasi-affine}
분리이고 국소 준유한인 스킴의 사상은 ind-준아핀이다.
\end{lemma}

\begin{proof}
$f : X \to Y$를 분리이고 국소 준유한인 스킴의 사상이라 하자.
$V \subset Y$를 아핀이라 하고 $U \subset f^{-1}(V)$를 준콤팩트
열린집합이라 하자. $U$가 준아핀임을 보여야 한다. $U \to V$가 분리이고
준유한인 스킴의 사상이므로 이는 Zariski 주정리에서 따른다. Lemma
\ref{more-morphisms-lemma-quasi-finite-separated-quasi-affine}를 보라.
\end{proof}




\section{스킴의 범주에서의 밀어내기, II}
\label{more-morphisms-section-pushouts-II}

\noindent
이 절은 Section \ref{more-morphisms-section-pushouts}의 연속이다. 여기서는 $Z \to X$가
닫힌 몰입이고 $Z \to Y$가 정수적이며 하나의 추가 조건이 만족되는 경우에
$Y \leftarrow Z \rightarrow X$의 밀어내기를 구성한다. 자세한 논의는
\cite{Ferrand-Conducteur}를 보라.

\begin{situation}
\label{more-morphisms-situation-pushout-along-closed-immersion-and-integral}
여기서 $S$는 스킴이고 $i : Z \to X$와 $j : Z \to Y$는 $S$ 위 스킴의
사상들이다. 다음을 가정한다.
\begin{enumerate}
\item $i$는 닫힌 몰입이다.
\item $j$는 스킴의 정수적 사상이다.
\item $y \in Y$에 대하여 $j^{-1}(\{y\}) \subset i^{-1}(U)$인 아핀
열린집합 $U \subset X$가 존재한다.
\end{enumerate}
\end{situation}

\begin{lemma}
\label{more-morphisms-lemma-prepare-pushout-along-closed-immersion-and-integral}
Situation \ref{more-morphisms-situation-pushout-along-closed-immersion-and-integral}에서
$y \in Y$에 대하여 $i^{-1}(U) = j^{-1}(V)$이고 $y \in V$인 아핀
열린집합 $U \subset X$와 $V \subset Y$가 존재한다.
\end{lemma}

\begin{proof}
$y \in Y$라 하자. $j^{-1}(\{y\}) \subset i^{-1}(U)$인 아핀 열린집합
$U \subset X$를 택한다(가정에 의해 가능하다). $y$의 아핀 열린 근방
$V \subset Y$ 가운데 $j^{-1}(V) \subset i^{-1}(U)$인 것을 택한다.
$j : Z \to Y$는 닫힌 사상이고(Morphisms, Lemma
\ref{morphisms-lemma-integral-universally-closed}) $i^{-1}(U)$가 $y$ 위의
올을 포함하므로 이것이 가능하다. $j$가 정수적이므로 스킴론적 올 $Z_y$는
체 위의 정수적 대수의 스펙트럼이다. Limits, Lemma
\ref{limits-lemma-ample-profinite-set-in-principal-affine}에 의해
$Z_y \subset D(\overline{f}) \subset j^{-1}(V)$인
$\overline{f} \in \Gamma(i^{-1}(U), \mathcal{O}_{i^{-1}(U)})$를 찾을 수 있다.
$i|_{i^{-1}(U)} : i^{-1}(U) \to U$는 아핀들 사이의 닫힌 몰입이므로
$i^{-1}(U)$로의 제한이 $\overline{f}$인 $f \in \Gamma(U, \mathcal{O}_U)$를
택할 수 있다. $U$를 주 열린집합 $D(f) \subset U$로 바꾸면
$y \in V \subset Y$와 $U \subset X$가 아핀 열린집합이고
$$
j^{-1}(\{y\}) \subset i^{-1}(U) \subset j^{-1}(V)
$$
인 것을 얻는다. 이제 (어떤 의미에서는) 이 논증을 반복한다. 즉,
$y \in D(g)$이고 $j^{-1}(D(g)) \subset i^{-1}(U)$인
$g \in \Gamma(V, \mathcal{O}_V)$를 택한다(위와 같은 논증으로 가능하다).
그런 다음 $i^{-1}(U) \to V$에 의한 $g$의 당김이 $i^{-1}(U)$로의 제한인
$f \in \Gamma(U, \mathcal{O}_U)$를 택할 수 있다(역시 위와 같은 이유로
가능하다). 따라서 마침내 아핀 열린집합
$y \in V' = D(g) \subset V \subset Y$와 $U' = D(f) \subset U \subset X$를
다음과 같이 얻는다.
% P05-EN-CH37-0319
$j^{-1}(V') = i^{-1}(V')$.
\end{proof}

\begin{proposition}
\label{more-morphisms-proposition-pushout-along-closed-immersion-and-integral}
\begin{reference}
\cite[Theorem 7.1 part iii]{Ferrand-Conducteur}
\end{reference}
Situation \ref{more-morphisms-situation-pushout-along-closed-immersion-and-integral}에서
밀어내기 $Y \amalg_Z X$는 스킴의 범주에서 존재한다. 그 도식은
$$
\xymatrix{
Z \ar[r]_i \ar[d]_j & X \ar[d]^a \\
Y \ar[r]^-b & Y \amalg_Z X
}
$$
이다. 이 도식은 올곱 도식이고, 사상 $a$는 정수적이며, 사상 $b$는 닫힌
몰입이고,
$$
\mathcal{O}_{Y \amalg_Z X} =
b_*\mathcal{O}_Y \times_{c_*\mathcal{O}_Z} a_*\mathcal{O}_X
$$
가 환의 층으로서 성립한다. 여기서 $c = a \circ i = b \circ j$이다.
\end{proposition}

\begin{proof}
위상공간으로서 $Y \amalg_Z X$를 위상공간의 범주에서 이 도식의 밀어내기로
정한다(Topology, Section \ref{topology-section-colimits}). 이는 바탕 집합들의
밀어내기에 몫위상을 준 것일 뿐이다(Topology, Lemma
\ref{topology-lemma-colimits}). $Y \amalg_Z X$ 위에는 환의 층의 사상
$$
b_*\mathcal{O}_Y \longrightarrow c_*\mathcal{O}_Z
\longleftarrow a_*\mathcal{O}_X
$$
들이 있고,
$$
\mathcal{O}_{Y \amalg_Z X} =
b_*\mathcal{O}_Y \times_{c_*\mathcal{O}_Z} a_*\mathcal{O}_X
$$
를 환의 층의 범주에서의 섬유곱으로 정의할 수 있다. 이로써 얻은 것이
스킴임을 보이려면 각 점이 아핀 열린 근방을 가짐을 보이면 된다. $c$의
상은 닫힌 부분집합이고 그 여집합은 환 달린 공간으로서
$(Y \setminus j(Z)) \amalg (X \setminus i(Z))$와 동형이므로, $c$의 상에
들지 않는 점들에 대해서는 명백하다.

\medskip\noindent
$c$의 상에 있는 점은 $j$의 상에 있는 유일한 $y \in Y$에 대응한다.
Lemma \ref{more-morphisms-lemma-prepare-pushout-along-closed-immersion-and-integral}에 의해
$y \in V$이고 $i^{-1}(U) = j^{-1}(V)$인 아핀 열린집합 $U \subset X$와
$V \subset Y$를 찾는다. 첫째 단락의 구성은 서로 양립하는 열린부분스킴으로의
제한과 명백히 양립하므로, 이 구성이 스킴을 준다는 것을 보일 때 $X$, $Y$,
$Z$가 아핀이라고 가정해도 된다.

\medskip\noindent
$X = \Spec(A)$, $Y = \Spec(B)$, $Z = \Spec(C)$가 아핀이면 More on
Algebra, Lemma \ref{more-algebra-lemma-points-of-fibre-product}에 의해
위상공간으로서 $Y \amalg_Z X = \Spec(B \times_C A)$이다.
% P05-EN-CH37-0320
$Y \times_Z X$가 스킴이라는 증명을 끝내려면 $\Spec(B \times_C A)$ 위의
구조층이 직상의 섬유곱임을 보이면 충분하다. 이는 $\Spec(B \times_C A)$의
주 아핀 열린집합들에 More on Algebra, Lemma
\ref{more-algebra-lemma-diagram-localize}를 적용하면 따른다.

\medskip\noindent
위의 논의에서 스킴 $Y \amalg_Z X$에는 아핀 열린 덮개
$Y \amalg_Z X = \bigcup W_i$가 존재하여 $U_i = a^{-1}(W_i)$,
$V_i = b^{-1}(W_i)$, $\Omega_i = c^{-1}(W_i)$가 각각 $X$, $Y$, $Z$의
아핀 열린집합임을 알 수 있다. 따라서 $a$와 $b$는 아핀이다. 더 나아가
$A_i$, $B_i$, $C_i$가 $U_i$, $V_i$, $\Omega_i$에 대응하는 환들이면
$A_i \to C_i$는 전사이고, $W_i$는 $B_i$로 전사하는
$A_i \times_{C_i} B_i$에 대응한다. 따라서 $b$는 닫힌 몰입이다. 환 사상
$A_i \times_{C_i} B_i \to A_i$는 More on Algebra, Lemma
\ref{more-algebra-lemma-fibre-product-integral}에 의해 정수적이므로 $a$는
정수적이다. 이 도식은 다음 이유로 Cartesian이다.
$$
C_i \cong B_i \otimes_{B_i \times_{C_i} A_i} A_i
$$
실제로 $B_i \times_{C_i} A_i \to B_i$와 $A_i \to C_i$는 핵이 같은
전사 사상들이다.

\medskip\noindent
마지막으로 Lemmas \ref{more-morphisms-lemma-basic-example-pushout}와
\ref{more-morphisms-lemma-pushout-fpqc-local}를 적용하면 우리의 구성이 스킴의 범주에서
밀어내기라고 결론지을 수 있다.
\end{proof}

\begin{lemma}
\label{more-morphisms-lemma-pushout-separated}
% P05-EN-CH37-0321
Situation \ref{more-morphisms-situation-pushout-along-closed-immersion-and-integral}에서
$X$와 $Y$가 분리이면 밀어내기 $Y \amalg_Z X$(Proposition
\ref{more-morphisms-proposition-pushout-along-closed-immersion-and-integral})도 분리이다.
``$S$ 위에서 분리'', ``준분리'', ``$S$ 위에서 준분리''에 대해서도
마찬가지이다.
\end{lemma}

\begin{proof}
사상 $Y \amalg X \to Y \amalg_Z X$는 전사이고 보편적으로 닫혀 있다. 따라서
Morphisms, Lemma \ref{morphisms-lemma-image-universally-closed-separated}를
적용할 수 있다.
\end{proof}

\begin{lemma}
\label{more-morphisms-lemma-pushout-finite-type}
Situation \ref{more-morphisms-situation-pushout-along-closed-immersion-and-integral}에서
$S$가 국소 Noether 스킴이고 $X$, $Y$, $Z$가 $S$ 위에서 국소 유한형이라고
가정하자. 그러면 밀어내기 $Y \amalg_Z X$(Proposition
\ref{more-morphisms-proposition-pushout-along-closed-immersion-and-integral})는 $S$ 위에서
국소 유한형이다.
\end{lemma}

\begin{proof}
아핀 열린집합들 위에서 살펴보면 More on Algebra, Lemma
\ref{more-algebra-lemma-fibre-product-finite-type}의 결과로 환원된다.
\end{proof}

\begin{lemma}
\label{more-morphisms-lemma-pushout-functor}
Situation \ref{more-morphisms-situation-pushout-along-closed-immersion-and-integral}에서
다음 가환 도식이 주어졌다고 하자.
$$
\xymatrix{
Y' \ar[d]^g & Z' \ar[l]^{j'} \ar[r]_{i'} \ar[d]^h & X' \ar[d]^f \\
Y & Z \ar[l] \ar[r] & X
}
$$
두 정사각형이 Cartesian이고 $f, g, h$가 분리이고 국소 준유한이라고 하자.
그러면
\begin{enumerate}
\item 밀어내기 $Y \amalg_Z X$와 $Y' \amalg_{Z'} X'$가 존재한다.
\item $Y' \amalg_{Z'} X' \to Y \amalg_Z X$는 분리이고 국소 준유한이다.
\item 다음 두 정사각형
$$
\xymatrix{
Y' \ar[r] \ar[d] & Y' \amalg_{Z'} X' \ar[d] & X' \ar[l] \ar[d] \\
Y \ar[r] & Y \amalg_Z X & X \ar[l]
}
$$
은 Cartesian이다.
\end{enumerate}
\end{lemma}

\begin{proof}
밀어내기 $Y \amalg_Z X$는 Proposition
\ref{more-morphisms-proposition-pushout-along-closed-immersion-and-integral}에 의해 존재한다.
밀어내기 $Y' \amalg_{Z'} X'$가 존재함을 보이기 위해 Situation
\ref{more-morphisms-situation-pushout-along-closed-immersion-and-integral}의 조건 (3)이
$(X', Y', Z', i', j')$에 대해 성립함을 확인하자. 즉 $y' \in Y'$라 하고
그 상을 $y \in Y$라 쓰자. $i(j^{-1}(y)) \subset U$인 아핀 열린집합
$U \subset X$를 택한다. $f^{-1}(U)$에 포함되고 준콤팩트 부분집합
$i'((j')^{-1}(\{y'\}))$를 포함하는 준콤팩트 열린집합 $U' \subset X'$를
택한다. Lemma \ref{more-morphisms-lemma-etale-separated-ind-quasi-affine}에 의해 $U'$는
준아핀이다. $Z'_{y'}$는 체 위의 정수적 대수의 스펙트럼이므로 Limits,
Lemma \ref{limits-lemma-ample-profinite-set-in-principal-affine}를 적용하면,
원하는 대로 $i'((j')^{-1}(\{y'\}))$를 포함하는 $U'$의 아핀 열린부분스킴이
존재한다.

\medskip\noindent
존재를 확인했으므로 나머지 주장들을 살펴보자. 아핀 국소적으로는
$B \to D$와 $A' \to C'$가 국소 준유한인 More on Algebra, Lemma
\ref{more-algebra-lemma-properties-algebras-over-fibre-product}의 상황과
정확히 같다\footnote{정확히는
$X, Y, Z, Y \amalg_Z X, X', Y', Z', Y' \amalg_{Z'} X'$가 각각
$A', B, A, B', C', D, C, D'$에 대응한다.}. 특히 사상
$Y' \amalg_{Z'} X' \to Y \amalg_Z X$는 국소 유한형이다.
% P05-EN-CH37-0322
도식의 정사각형들은 More on Algebra, Lemma
\ref{more-algebra-lemma-module-over-fibre-product}에 의해 Cartesian이다.
국소 준유한성은 올들 위에서 확인할 수 있으므로(Morphisms, Lemma
\ref{morphisms-lemma-quasi-finite-at-point-characterize}),
$Y' \amalg_{Z'} X' \to Y \amalg_Z X$가 국소 준유한이라고 결론짓는다.

\medskip\noindent
아직 $Y' \amalg_{Z'} X' \to Y \amalg_Z X$가 분리임을 확인해야 한다.
Proposition \ref{more-morphisms-proposition-pushout-along-closed-immersion-and-integral}에 의해
$Y' \amalg X' \to Y' \amalg_{Z'} X'$는 보편적으로 닫혀 있고 전사이다.
사상 $Y' \amalg X' \to Y \amalg_Z X$도
$Y' \amalg X' \to Y \amalg X \to Y \amalg_Z X$로 인수분해되므로 분리이다.
따라서 Morphisms, Lemma
\ref{morphisms-lemma-image-universally-closed-separated}에 의해 결론이 따른다.
\end{proof}

\begin{lemma}
\label{more-morphisms-lemma-pushout-functor-equivalence-flat}
Situation \ref{more-morphisms-situation-pushout-along-closed-immersion-and-integral}에서
밀어내기 $Y \amalg_Z X$ 위에서 평탄하고 분리이며 국소 준유한인 스킴들의
범주는 Lemma \ref{more-morphisms-lemma-pushout-functor}에서와 같은
$(X', Y', Z', i', j', f, g, h)$ 가운데 $f, g, h$가 평탄인 것들의 범주와
동치이다. ``평탄''을 ``\'etale''로 바꾸어도 마찬가지이다.
\end{lemma}

\begin{proof}
Lemma \ref{more-morphisms-lemma-pushout-functor}에서와 같은
$(X', Y', Z', i', j', f, g, h)$에서 시작하고 $f, g, h$가 평탄 또는
\'etale이면 More on Algebra, Lemma
\ref{more-algebra-lemma-properties-algebras-over-fibre-product}에 의해
$Y' \amalg_{Z'} X' \to Y \amalg_Z X$는 평탄 또는 \'etale이다.

\medskip\noindent
역으로 $W \to Y \amalg_Z X$를 분리이고 국소 준유한인 사상이라 하자.
$X' = W \times_{Y \amalg_Z X} X$, $Y' = W \times_{Y \amalg_Z X} Y$,
$Z' = W \times_{Y \amalg_Z X} Z$라 놓고 명백한 사상 $i', j', f, g, h$를
취한다. 밀어내기 $Y' \amalg_{Z'} X'$를 만들면
$$
Y' \amalg_{Z'} X' \longrightarrow W
$$
% P05-EN-CH37-0323
라는 $Y \amalg_X Z$ 위 스킴의 사상을 밀어내기의 보편 성질로부터 얻는다.
$W \to Y \amalg_Z X$가 평탄이라고 가정하지 않으면 일반적으로 이 사상은
동형이 아니다. (실제로 More on Algebra, Lemma
\ref{more-algebra-lemma-module-over-fibre-product-bis}에 따르면 표시된 화살표는
닫힌 몰입이지만 일반적으로 동형은 아니다.) 그러나
% P05-EN-CH37-0324
$W \to Y \times_Z X$가 평탄이면 More on Algebra, Lemma
\ref{more-algebra-lemma-properties-algebras-over-fibre-product}에 의해
이 사상은 동형이다.
\end{proof}

\noindent
다음으로 두 사상이 모두 닫힌 몰입인 경우의 존재성을 논의한다.

\begin{lemma}
\label{more-morphisms-lemma-pushout-along-closed-immersions}
$i : Z \to X$와 $j : Z \to Y$를 스킴의 닫힌 몰입들이라 하자. 그러면
밀어내기 $Y \amalg_Z X$는 스킴의 범주에서 존재한다. 그 도식은
$$
\xymatrix{
Z \ar[r]_i \ar[d]_j & X \ar[d]^a \\
Y \ar[r]^-b & Y \amalg_Z X
}
$$
이다. 이 도식은 올곱 도식이고, 사상 $a$와 $b$는 닫힌 몰입이며, 다음
짧은 완전열이 존재한다.
$$
0 \to \mathcal{O}_{Y \amalg_Z X} \to
a_*\mathcal{O}_X \oplus b_*\mathcal{O}_Y \to
c_*\mathcal{O}_Z \to 0
$$
여기서 $c = a \circ i = b \circ j$이다.
\end{lemma}

\begin{proof}
이는 Proposition \ref{more-morphisms-proposition-pushout-along-closed-immersion-and-integral}의
특별한 경우이다. $j$의 올들은 한원소집합이므로 Situation
\ref{more-morphisms-situation-pushout-along-closed-immersion-and-integral}의 가정 (3)은
즉시 성립한다. 마지막으로 화살표들의 역할을 뒤바꾸면 $a$와 $b$가 모두
닫힌 몰입이라고 결론지을 수 있다.
\end{proof}

\begin{lemma}
\label{more-morphisms-lemma-pushout-along-closed-immersions-properties-above}
$i : Z \to X$와 $j : Z \to Y$를 스킴의 닫힌 몰입들이라 하자.
$f : X' \to X$와 $g : Y' \to Y$를 스킴의 사상들이라 하고
$\varphi : X' \times_{X, i} Z \to Y' \times_{Y, j} Z$를 $Z$ 위 스킴의
동형이라 하자. 다음 사상을 생각하자.
$$
h :
X' \amalg_{X' \times_{X, i} Z, \varphi} Y'
\longrightarrow
X \amalg_Z Y
$$
그러면 다음이 성립한다.
\begin{enumerate}
\item $h$가 국소 유한형일 필요충분조건은 $f$와 $g$가 국소 유한형인 것이다.
\item $h$가 평탄일 필요충분조건은 $f$와 $g$가 평탄인 것이다.
\item $h$가 평탄이고 국소 유한 표시일 필요충분조건은 $f$와 $g$가 평탄이고
국소 유한 표시인 것이다.
\item $h$가 매끄러울 필요충분조건은 $f$와 $g$가 매끄러운 것이다.
\item $h$가 \'etale일 필요충분조건은 $f$와 $g$가 \'etale인 것이다.
% P05-EN-CH37-0325
\item 필요에 따라 여기에 더 추가한다.
\end{enumerate}
\end{lemma}

\begin{proof}
Lemma \ref{more-morphisms-lemma-pushout-along-closed-immersions}에 의해 밀어내기들이 존재함을
알고 있다. 특히 사상 $h$를 얻는다. 따라서 등장하는 모든 스킴을 아핀
스킴으로 바꾸어도 된다. 이 경우 보조정리의 주장들은 More on Algebra,
Lemma \ref{more-algebra-lemma-properties-algebras-over-fibre-product}의
대응하는 주장들과 동치이다.
\end{proof}





\section{상대 사상}
\label{more-morphisms-section-relative-morphisms}

\noindent
이 절에서는 표현가능성에 관한 결과를 증명한다. 이 결과는 Fundamental
Groups, Section \ref{pione-section-finite-etale}에서 스킴의 유한 \'etale
덮개들의 범주에 관한 결과를 증명하는 데 사용할 것이다. 이 절의 내용은
Criteria for Representability, Section
\ref{criteria-section-relative-morphisms}에서 적절한 일반성으로 논의된다.

\medskip\noindent
$S$를 스킴이라 하고 $Z$와 $X$를 $S$ 위의 스킴들이라 하자. $S$ 위의 스킴
$T$가 주어지면 $S$ 위의 사상
$b : T \times_S Z \to T \times_S X$를 생각할 수 있다. 그 도식은
\begin{equation}
\label{more-morphisms-equation-hom}
\vcenter{
\xymatrix{
T \times_S Z \ar[rd] \ar[rr]_b & &
T \times_S X \ar[ld] & Z \ar[rd] & & X \ar[ld] \\
& T \ar[rrr] & & & S
}
}
\end{equation}
이다. 물론 $b$를 다음 도식이 가환이 되게 하는 사상
$b : T \times_S Z \to X$로 생각할 수도 있다.
$$
\xymatrix{
T \times_S Z \ar[r] \ar[d] \ar@/^1pc/[rrr]_-b &
Z \ar[rd] & & X \ar[ld] \\
T \ar[rr] & & S
}
$$
이 상황에서 다음 함자를 정의할 수 있다.
\begin{equation}
\label{more-morphisms-equation-hom-functor}
\mathit{Mor}_S(Z, X) : (\Sch/S)^{opp} \longrightarrow \textit{Sets},
\quad
T \longmapsto \{b\text{ as above}\}
\end{equation}
다음은 기본적인 표현가능성 결과이다.

\begin{lemma}
\label{more-morphisms-lemma-hom-from-finite-free-into-affine}
$Z \to S$와 $X \to S$를 아핀 스킴들의 사상들이라 하자.
$\Gamma(Z, \mathcal{O}_Z)$가 유한 자유
$\Gamma(S, \mathcal{O}_S)$-가군이라고 가정하자. 그러면
$\mathit{Mor}_S(Z, X)$는 $S$ 위의 아핀 스킴으로 표현가능하다.
\end{lemma}

\begin{proof}
$S = \Spec(R)$라 쓰자. $\Gamma(Z, \mathcal{O}_Z)$의 $R$ 위 기저
$\{e_1, \ldots, e_m\}$를 택한다. 다음 표시도 택하자.
$$
\Gamma(X, \mathcal{O}_X) = R[\{x_i\}_{i \in I}]/(\{f_k\}_{k \in K}).
$$
이 몫에서 $x_i$의 상을 $\overline{x}_i$로 나타내자. 다음과 같이 쓰자.
$$
P = R[\{a_{ij}\}_{i \in I, 1 \leq j \leq m}].
$$
다음 $R$-대수 사상을 생각하자.
$$
\Psi :
R[\{x_i\}_{i \in I}]
\longrightarrow
P \otimes_R \Gamma(Z, \mathcal{O}_Z), \quad
x_i \longmapsto \sum\nolimits_j a_{ij} \otimes e_j.
$$
$c_{kj} \in P$를 사용하여 $\Psi(f_k) = \sum c_{kj} \otimes e_j$라 쓰자.
마지막으로 $J \subset P$를 원소 $c_{kj}$, $k \in K$,
$1 \leq j \leq m$이 생성하는 아이디얼이라 하자. $W = \Spec(P/J)$가 함자
$\mathit{Mor}_S(Z, X)$를 표현한다고 주장한다.

\medskip\noindent
먼저 구성에 의해 $P/J$는 $R$-대수이므로 사상 $W \to S$를 준다. 둘째로
구성에 의해 사상 $\Psi$는 $\Gamma(X, \mathcal{O}_X)$를 통해 분해되므로,
% P05-EN-CH37-0326
다음 $P/J$-대수 준동형을 얻는다.
$$
P/J \otimes_R \Gamma(X, \mathcal{O}_X)
\longrightarrow
P/J \otimes_R \Gamma(Z, \mathcal{O}_Z)
$$
이는 사상 $b_{univ} : W \times_S Z \to W \times_S X$를 정한다.
Yoneda 보조정리에 의해 $b_{univ}$는 함자들의 변환
$W \to \mathit{Mor}_S(Z, X)$를 정하며, 이것이 동형이라고 주장한다.
동형임을 보이려면 임의의 아핀 스킴 $T$에 대하여 집합들의 전단사
$W(T) \to \mathit{Mor}_S(Z, X)(T)$를 유도함을 보이면 충분하다.

\medskip\noindent
$T = \Spec(R')$가 $S$ 위의 아핀 스킴이고
$b \in \mathit{Mor}_S(Z, X)(T)$라고 하자. 구조 사상 $T \to S$는 $R'$ 위의
$R$-대수 구조를 정하고, $b$는 다음 $R'$-대수 사상을 정한다.
$$
b^\sharp :
R' \otimes_R \Gamma(X, \mathcal{O}_X)
\longrightarrow
R' \otimes_R \Gamma(Z, \mathcal{O}_Z).
$$
특히 어떤 $\alpha_{ij} \in R'$에 대하여
$b^\sharp(1 \otimes \overline{x}_i) = \sum \alpha_{ij} \otimes e_j$라
쓸 수 있다. 이는 규칙 $a_{ij} \mapsto \alpha_{ij}$로 정해지는 $R$-대수
사상 $P \to R'$에 대응한다. 아이디얼 $J$의 구성에 의해 이 사상은 몫
$P/J$를 통해 분해되어 사상 $P/J \to R'$를 준다. 이는 다시 사상
$T \to W$에 대응하며, 이때 $b$는 $b_{univ}$의 당김이다. 일부 세부사항은
생략한다.
\end{proof}

\begin{lemma}
\label{more-morphisms-lemma-hom-from-finite-locally-free-into-affine}
$Z \to S$와 $X \to S$를 스킴의 사상들이라 하자. $Z \to S$가 유한 국소
자유이고 $X \to S$가 아핀이면 $\mathit{Mor}_S(Z, X)$는 $S$ 위에서 아핀인
스킴으로 표현가능하다.
\end{lemma}

\begin{proof}
$\Gamma(Z \times_S U_i, \mathcal{O}_{Z \times_S U_i})$가
$\mathcal{O}_S(U_i)$ 위에서 유한 자유가 되는 아핀 열린 덮개
$S = \bigcup U_i$를 택한다. $F_i \subset \mathit{Mor}_S(Z, X)$를 다음
부분함자라 하자. 이는 $T \to S$가 $U_i$를 통해 분해되지 않으면 $T/S$에
공집합을 대응시키고, 그렇지 않으면 $\mathit{Mor}_S(Z, X)(T)$를 대응시킨다.
% P05-EN-CH37-0327
이 부분함자들의 모임은 Schemes, Lemma \ref{schemes-lemma-glue-functors}의
조건 (2)(a), (2)(b), (2)(c)를 만족하므로 보조정리가 증명된다. 조건 (2)(a)는
Lemma \ref{more-morphisms-lemma-hom-from-finite-free-into-affine}에서 따르고, 나머지 둘은
직접적인 논증에서 따른다.
\end{proof}

\noindent
아래 보조정리에서 사상 $f : X \to S$에 부과한 조건은 이와 같은 명제를
증명할 때 매우 유용하다. 다음 중 하나가 성립하면 이 조건이 만족된다.
$X$가 준아핀이거나, $f$가 준아핀이거나, $f$가 준사영적이거나, $f$가
국소 사영적이거나, $X$ 위에 풍부한 가역층이 존재하거나, $X$ 위에
$f$-풍부한 가역층이 존재하거나, $X$ 위에 $f$-매우 풍부한 가역층이
존재하는 경우이다.

\begin{lemma}
\label{more-morphisms-lemma-hom-from-finite-locally-free-representable}
$Z \to S$와 $X \to S$를 스킴의 사상들이라 하자. 다음을 가정하자.
\begin{enumerate}
\item $Z \to S$는 유한 국소 자유이다.
\item $s \in S$이고 $x_1, \ldots, x_d \in X_s$인 모든
$(s, x_1, \ldots, x_d)$에 대하여 $x_1, \ldots, x_d \in U$인 아핀
열린집합 $U \subset X$가 존재한다.
\end{enumerate}
그러면 $\mathit{Mor}_S(Z, X)$는 스킴으로 표현가능하다.
\end{lemma}

\begin{proof}
$U \subset X$와 $V \subset S$가 아핀 열린집합이고 $U \to S$가 $V$를 통해
분해되는 순서쌍 $(U, V)$들의 집합을 $I$라 하자. $i \in I$에 대응하는
순서쌍을 $(U_i, V_i)$라 쓰고 $F_i = \mathit{Mor}_{V_i}(Z_{V_i}, U_i)$라
놓자. $F_i$가 $\mathit{Mor}_S(Z, X)$의 부분함자인 것은 즉시 알 수 있다.
% P05-EN-CH37-0328
Schemes, Lemma \ref{schemes-lemma-glue-functors}의 조건 (2)(a), (2)(b),
(2)(c)가 만족된다고 주장한다. 이로써 보조정리가 증명된다.

\medskip\noindent
조건 (2)(a)는 Lemma
\ref{more-morphisms-lemma-hom-from-finite-locally-free-into-affine}에서 따른다.

\medskip\noindent
조건 (2)(b)를 확인하기 위해 $T/S$와
$b \in \mathit{Mor}_S(Z, X)(T)$를 생각하자. $b$를 사상
$T \times_S Z \to X$로 생각하면 열린집합
$b^{-1}(U_i) \subset T \times_S Z$를 얻는다. 명백히 $b \in F_i(T)$일
필요충분조건은 $b^{-1}(U_i) = T \times_S Z$인 것이다. 사영
$p : T \times_S Z \to T$는 유한이고 따라서 닫혀 있으므로,
$p^{-1}(\{t\}) \subset b^{-1}(U_i)$인 점 $t \in T$들의 집합
$U_{i, b} \subset T$는 열려 있다. 그러면
% P05-EN-CH37-0329
$f : T' \to T$가 $U_{i, b}$를 통해 분해될 필요충분조건은
$b \circ f \in F_i(T')$인 것이므로 조건 (2)(b)의 확인이 끝난다.

\medskip\noindent
마지막으로 조건 (2)(c)를 확인하자. 여기서 $X \to S$에 대한 조건을
사용한다. 즉 $T/S$와 $b \in \mathit{Mor}_S(Z, X)(T)$를 생각하자.
모든 $t \in T$가 앞 단락에서 정의한 열린집합 $U_{i, b}$ 중 하나에
포함됨을 보이면 충분하다. 이는 $p : T \times_S Z \to T$가 사영이고
$b : T \times_S Z \to X$가 주어진 사상일 때, 어떤 $i$에 대하여
$b(p^{-1}(\{t\})) \subset U_i$라는 조건과 동치이다. $p$가 유한이므로
집합 $b(p^{-1}(\{t\})) \subset X$는 유한하고, $t$의 $S$에서의 상을
$s$라 할 때 그 점 위의 $X \to S$의 올에 포함된다. 따라서 $X \to S$에
대한 우리의 조건은 적절한 순서쌍이 존재함을 정확히 보인다.
\end{proof}

\begin{lemma}
\label{more-morphisms-lemma-hom-from-finite-locally-free-separated-lqf}
$Z \to S$와 $X \to S$를 스킴의 사상들이라 하자. $Z \to S$가 유한 국소
자유이고 $X \to S$가 분리이고 국소 준유한이라고 가정하자. 그러면
$\mathit{Mor}_S(Z, X)$는 스킴으로 표현가능하다.
\end{lemma}

\begin{proof}
이는 Lemmas \ref{more-morphisms-lemma-hom-from-finite-locally-free-representable}와
\ref{more-morphisms-lemma-separated-locally-quasi-finite-over-affine}에서 따른다.
\end{proof}









\section{유사연접 복합체의 특징짓기, III}
\label{more-morphisms-section-characterize-pseudo-coherent}

\noindent
이 절에서는 코호몰로지를 이용한 유사연접 복합체의 특징짓기를 논한다.
이는 Derived Categories of Schemes, Section
\ref{perfect-section-pseudo-coherent}의 연속이다. 기본 도구는 Chow
보조정리의 유도 버전을 사용하여 사영공간의 경우로 환원하는 것이다.
Lemma \ref{more-morphisms-lemma-derived-chow}를 보라.

\begin{lemma}
\label{more-morphisms-lemma-case-of-tor-independence}
다음 스킴의 가환 도식을 생각하자.
$$
\xymatrix{
Z' \ar[d] \ar[r] & Y' \ar[d] \\
X' \ar[r] & S'
}
$$
$S \to S'$를 사상이라 하자. $X'$와 $Y'$의 $S$로의 밑변환을 각각 $X$와
$Y$로 나타내자. $Y' \to S'$와 $Z' \to X'$가 평탄이라고 가정하자.
그러면 $X \times_S Y$와 $Z'$는 $X' \times_{S'} Y'$ 위에서 Tor-독립이다.
\end{lemma}

\begin{proof}
문제는 국소적이므로 모든 스킴이 아핀이라고 가정해도 된다(일부 세부사항은
생략한다). 다음 도식은
$$
\xymatrix{
X \times_S Y \ar[r] \ar[d] & X' \times_{S'} Y' \ar[d] \\
X \ar[r] & X'
}
$$
수직 화살표들이 평탄인 Cartesian 도식이다. $X = \Spec(A)$,
$X' = \Spec(A')$, $X' \times_{S'} Y' = \Spec(B')$라 쓰자. 그러면
$X \times_S Y = \Spec(A \otimes_{A'} B')$이다. $Z' = \Spec(C')$라
쓰자. 다음을 보여야 한다.
$$
\text{Tor}_p^{B'}(A \otimes_{A'} B', C') = 0,
\quad\text{for } p > 0
$$
$A' \to B'$가 평탄이므로
$A \otimes_{A'} B' = A \otimes_{A'}^\mathbf{L} B'$이다. 따라서
$$
(A \otimes_{A'} B') \otimes_{B'}^\mathbf{L} C' =
(A \otimes_{A'}^\mathbf{L} B') \otimes_{B'}^\mathbf{L} C' =
A \otimes_{A'}^\mathbf{L} C' =
A \otimes_{A'} C'
$$
% P05-EN-CH37-0330
둘째 등식은 More on Algebra, Lemma
\ref{more-algebra-lemma-double-base-change}에서 따른다. 마지막 등식은
$A' \to C'$가 평탄이기 때문에 성립한다. 이로써 보조정리가 증명된다.
\end{proof}

\begin{lemma}[유도 Chow 보조정리]
\label{more-morphisms-lemma-derived-chow}
$A$를 환이라 하고 $X$를 $A$ 위에서 분리이고 유한 표시인 스킴이라 하자.
$x \in X$라 하자. 그러면 $x$의 열린 근방 $U \subset X$, $n \geq 0$,
열린집합 $V \subset \mathbf{P}^n_A$, 닫힌부분스킴
$Z \subset X \times_A \mathbf{P}^n_A$, 점 $z \in Z$, 그리고
$D(\mathcal{O}_{X \times_A \mathbf{P}^n_A})$의 대상 $E$가 존재하여
\begin{enumerate}
\item $Z \to X \times_A \mathbf{P}^n_A$는 유한 표시이고,
\item $b : Z \to X$는 $U$ 위에서 동형이고 $b(z) = x$이며,
\item $c : Z \to \mathbf{P}^n_A$는 $V$ 위에서 닫힌 몰입이고,
\item $b^{-1}(U) = c^{-1}(V)$이며, 특히 $c(z) \in V$이고,
\item $E|_{X \times_A V} \cong
(b, c)_*\mathcal{O}_Z|_{X \times_A V}$,
\item $E$는 유사연접이고 받침이 $Z$에 놓인다.
\end{enumerate}
\end{lemma}

\begin{proof}
유한형 $\mathbf{Z}$-부분대수 $A' \subset A$와 $A'$ 위에서 분리이고 유한
표시인 스킴 $X'$ 가운데 $A$로의 밑변환이 $X$인 것을 찾을 수 있다.
Limits, Lemmas \ref{limits-lemma-descend-finite-presentation}와
\ref{limits-lemma-descend-separated-finite-presentation}를 보라. $x$의 상을
$x' \in X'$라 하자. $x' \in X'/A'$에 대해 보조정리를 증명할 수 있으면
$x \in X/A$에 대한 보조정리도 따른다. 실제로
$U', n', V', Z', z', E'$가 $x' \in X'/A'$에 대한 해를 준다면,
$U \subset X$를 $U'$의 역상으로, $n = n'$으로,
$V \subset \mathbf{P}^n_A$를 $V'$의 역상으로,
$Z \subset X \times \mathbf{P}^n$를 $Z'$의 스킴론적 역상으로,
$z \in Z$를 $x$로 가는 유일한 점으로 놓고, $E$를 $E'$의 유도 당김으로
놓을 수 있다. Cohomology, Lemma
\ref{cohomology-lemma-pseudo-coherent-pullback}에 의해 $E$는 유사연접이다.
이제 (5)만 확인하면 된다. 이를 위해 $W = b^{-1}(U) = c^{-1}(V)$라 놓고,
% P05-EN-CH37-0331
$W' = (b')^{-1}(U) = (c')^{-1}(V')$라 놓은 뒤 다음 Cartesian 도식을
생각하자.
$$
\xymatrix{
W \ar[d]_{(b, c)} \ar[r] & W' \ar[d]^{(b', c')} \\
X \times_A V \ar[r] & X' \times_{A'} V'
}
$$
Lemma \ref{more-morphisms-lemma-case-of-tor-independence}에 의해 스킴
$X \times_A V$와 $W'$는 $X' \times_{A'} V'$ 위에서 Tor-독립이다. 따라서
Derived Categories of Schemes, Lemma
\ref{perfect-lemma-compare-base-change}에 의해
$(b', c')_*\mathcal{O}_{W'}$의 $X \times_A V$로의 유도 당김은
$(b, c)_*\mathcal{O}_W$이다. 여기서는 $(b', c')$가 닫힌 몰입이므로
$R(b', c')_*\mathcal{O}_{Z'} = (b', c')_*\mathcal{O}_{Z'}$이고
$(b, c)_*\mathcal{O}_Z$에 대해서도 마찬가지라는 사실도 사용한다.
$E'|_{U' \times_{A'} V'} = (b', c')_*\mathcal{O}_{W'}$이므로
$E|_{U \times_A V} = (b, c)_*\mathcal{O}_W$를 얻고 (5)가 성립한다.
이로써 다음 단락에서 설명하는 상황으로 환원된다.

\medskip\noindent
$A$가 $\mathbf{Z}$ 위에서 유한형이라고 가정하자. $x$의 아핀 열린 근방
$U \subset X$를 택한다. 그러면 $U$는 $A$ 위에서 유한형이다. 닫힌 몰입
$U \to \mathbf{A}^n_A$를 택하고, 이를 열린 몰입
$\mathbf{A}^n_A \to \mathbf{P}^n_A$와 합성하여 얻은 몰입을
$j : U \to \mathbf{P}^n_A$라 나타내자. 다음 사상의 스킴론적 폐포를 $Z$라
하자.
$$
(\text{id}_U, j) : U \longrightarrow X \times_A \mathbf{P}^n_A
$$
사영 $X \times \mathbf{P}^n \to X$가 분리이므로 Morphisms, Lemma
\ref{morphisms-lemma-scheme-theoretic-image-of-partial-section}에 의해
$b : Z \to X$는 $U$ 위에서 동형이라고 결론짓는다. $z \in Z$를 $x$ 위에
놓인 유일한 점이라 하자.

\medskip\noindent
$Y \subset \mathbf{P}^n_A$를 $j$의 스킴론적 폐포라 하자. 그러면
$Z \subset X \times_A Y$가
$(\text{id}_U, j) : U \to X \times_A Y$의 스킴론적 폐포인 것은 명백하다.
$X$가 분리이므로 사상 $X \times_A Y \to Y$도 분리이다. 따라서 위에서
사용한 같은 보조정리에 의해 $Z \to Y$는 열린부분스킴 $j(U) \subset Y$
위에서 동형이다. $V \cap Y = j(U)$인 열린집합
$V \subset \mathbf{P}^n_A$를 택하면 (3)과 (4)가 성립한다.

\medskip\noindent
$A$가 Noether이므로 $X$와 $X \times_A \mathbf{P}^n_A$는 Noether
스킴들이다. 따라서 이 경우 $E = (b, c)_*\mathcal{O}_Z$로 택할 수 있다.
Derived Categories of Schemes, Lemma
\ref{perfect-lemma-identify-pseudo-coherent-noetherian}를 보라. 이로써 증명이
끝난다.
\end{proof}

\begin{lemma}
\label{more-morphisms-lemma-compute-Fourier-Mukai-for-derived-chow}
$A$, $x \in X$, $U, n, V, Z, z, E$가 Lemma \ref{more-morphisms-lemma-derived-chow}에서와
같다고 하자. 임의의 $K \in D_\QCoh(\mathcal{O}_X)$에 대하여
$$
Rq_*(Lp^*K \otimes^\mathbf{L} E)|_V = R(U \to V)_*K|_U
$$
가 성립한다. 여기서 $p : X \times_A \mathbf{P}^n_A \to X$와
$q : X \times_A \mathbf{P}^n_A \to \mathbf{P}^n_A$는 사영들이고,
사상 $U \to V$는 유한 표시 닫힌 몰입 $c \circ (b|_U)^{-1}$이다.
\end{lemma}

\begin{proof}
$b^{-1}(U) = c^{-1}(V)$이고 $c$가 $V$ 위에서 닫힌 몰입이므로
$c \circ (b|_U)^{-1}$는 닫힌 몰입이다. $U$와 $V$가 $A$ 위에서 유한
표시이므로 이 사상도 유한 표시이다. Morphisms, Lemma
\ref{morphisms-lemma-finite-presentation-permanence}를 보라. 먼저
$$
Rq_*(Lp^*K \otimes^\mathbf{L} E)|_V =
Rq'_*\left((Lp^*K \otimes^\mathbf{L} E)|_{X \times_A V}\right)
$$
이다. 여기서 $q' : X \times_A V \to V$는 사영이며, 이는 전체 직상의
형성이 국소화와 가환하기 때문이다. $W = b^{-1}(U) = c^{-1}(V)$라 놓고
$i : W \to X \times_A V$로 닫힌 몰입 $i = (b, c)|_W$를 나타내자. 그러면
$$
Rq'_*\left((Lp^*K \otimes^\mathbf{L} E)|_{X \times_A V}\right) =
Rq'_*(Lp^*K|_{X \times_A V} \otimes^\mathbf{L} i_*\mathcal{O}_W)
$$
이며, 이는 성질 (5)에 의한 것이다. $i$가 닫힌 몰입이므로
$i_*\mathcal{O}_W = Ri_*\mathcal{O}_W$이다. Derived Categories of Schemes,
Lemma \ref{perfect-lemma-cohomology-base-change}를 사용하면 이를 다음과 같이
다시 쓸 수 있다.
$$
Rq'_* Ri_* Li^* Lp^*K|_{X \times_A V} =
R(q' \circ i)_* Lb^*K|_W =
R(U \to V)_* K|_U
$$
이는 원하는 바이다.
\end{proof}

\begin{lemma}
\label{more-morphisms-lemma-characterize-pseudo-coherent}
$A$를 환이라 하고 $X$를 $A$ 위에서 분리이고 유한 표시인 스킴이라 하자.
$K \in D_\QCoh(\mathcal{O}_X)$라 하자. $D(\mathcal{O}_X)$의 모든 유사연접
$E$에 대하여 $R\Gamma(X, E \otimes^\mathbf{L} K)$가 $D(A)$에서
유사연접이면, $K$는 $A$에 상대적으로 유사연접이다.
\end{lemma}

\begin{proof}
$K \in D_\QCoh(\mathcal{O}_X)$이고 $D(\mathcal{O}_X)$의 모든 유사연접
$E$에 대하여 $R\Gamma(X, E \otimes^\mathbf{L} K)$가 $D(A)$에서
유사연접이라고 가정하자. $x \in X$라 하자. $x$의 한 근방에서 $K$가
$A$에 상대적으로 유사연접임을 보일 것이며, 이로써 보조정리가 증명된다.

\medskip\noindent
$U, n, V, Z, z, E$를 Lemma \ref{more-morphisms-lemma-derived-chow}에서와 같이 택하자.
$p : X \times \mathbf{P}^n \to X$와
$q : X \times \mathbf{P}^n \to \mathbf{P}^n_A$로 사영들을 나타내자.
그러면 임의의 $i \in \mathbf{Z}$에 대하여
\begin{align*}
& R\Gamma(\mathbf{P}^n_A,
Rq_*(Lp^*K \otimes^\mathbf{L} E)
\otimes^\mathbf{L}
\mathcal{O}_{\mathbf{P}^n_A}(i)) \\
& =
R\Gamma(X \times \mathbf{P}^n,
Lp^*K \otimes^\mathbf{L} E \otimes^\mathbf{L}
Lq^*\mathcal{O}_{\mathbf{P}^n_A}(i)) \\
& =
R\Gamma(X,
K \otimes^\mathbf{L} Rp_*(E \otimes^\mathbf{L}
Lq^*\mathcal{O}_{\mathbf{P}^n_A}(i)))
\end{align*}
가 성립하며, 이는 Derived Categories of Schemes, Lemma
\ref{perfect-lemma-cohomology-base-change}에 의한 것이다. Derived Categories
of Schemes, Lemma
\ref{perfect-lemma-flat-proper-pseudo-coherent-direct-image-general}에 의해
복합체 $Rp_*(E \otimes^\mathbf{L} Lq^*\mathcal{O}_{\mathbf{P}^n_A}(i))$는
$X$ 위에서 유사연접이다. 따라서 가정에 의해 표시된 식은 $D(A)$의 유사연접
대상이다. Derived Categories of Schemes, Lemma
\ref{perfect-lemma-pseudo-coherent-on-projective-space}에 의해
$Rq_*(Lp^*K \otimes^\mathbf{L} E)$가 $\mathbf{P}^n_A$ 위에서
유사연접이라고 결론짓는다. Lemma
\ref{more-morphisms-lemma-compute-Fourier-Mukai-for-derived-chow}에 의해
$$
Rq_*(Lp^*K \otimes^\mathbf{L} E)|_{X \times_A V} =
R(U \to V)_*K|_U
$$
% P05-EN-CH37-0332
이다. $U \to V$가 $\mathbf{P}^n_A$의 열린부분스킴으로 가는 닫힌 몰입이므로
Lemma \ref{more-morphisms-lemma-check-relative-pseudo-coherence-on-charts}에 의해 이는
$K|_U$가 $A$에 상대적으로 유사연접임을 뜻한다.
\end{proof}

\begin{lemma}
\label{more-morphisms-lemma-characterize-pseudo-coh-improved}
$A$를 환이라 하고 $X$를 $A$ 위에서 분리이고 유한 표시인 스킴이라 하자.
$K \in D_\QCoh(\mathcal{O}_X).$ 모든 완전
$E \in D(\mathcal{O}_X)$에 대하여
$R \Gamma (X, E \otimes ^{\mathbf{L}} K)$가 $D(A)$에서 유사연접이면
$K$는 $A$에 상대적으로 유사연접이다.
\end{lemma}

\begin{proof}
Lemma \ref{more-morphisms-lemma-characterize-pseudo-coherent}에 비추어 보면, 모든 유사연접
$E \in D(\mathcal{O}_X)$에 대하여
$R \Gamma (X, E \otimes ^{\mathbf{L}} K)$가 $D(A)$에서 유사연접임을
보이면 충분하다. Derived Categories of Schemes, Proposition
\ref{perfect-proposition-detecting-bounded-above}에 의해
$K \in D^-_\QCoh (\mathcal{O}_X)$가 따른다. 이제 Derived Categories of
Schemes, Lemma \ref{perfect-lemma-perfect-enough}에서 결과가 따른다.
\end{proof}

\begin{lemma}
\label{more-morphisms-lemma-characterize-relatively-perfect}
$A$를 환이라 하고 $X$를 $A$ 위에서 분리이고 유한 표시이며 평탄인 스킴이라
하자. $K \in D_\QCoh(\mathcal{O}_X).$ 모든 완전
$E \in D(\mathcal{O}_X)$에 대하여
$R \Gamma (X, E \otimes^\mathbf{L} K)$가 $D(A)$에서 완전이면 $K$는
$\Spec(A)$-완전이다.
\end{lemma}

\begin{proof}
Lemma \ref{more-morphisms-lemma-characterize-pseudo-coh-improved}에 의해 $K$는 $A$에
상대적으로 유사연접이다. Lemma
\ref{more-morphisms-lemma-check-relative-pseudo-coherence-on-charts}에 의해 $K$는
$D( \mathcal{O}_X)$에서 유사연접이다. Derived Categories of Schemes,
Proposition \ref{perfect-proposition-detecting-bounded-below}에 의해 $K$는
$D^-(\mathcal{O}_X)$에 속한다. $\mathfrak{p}$를 $A$의 소아이디얼이라 하고,
$i : Y \to X$로 $\mathfrak{p}$ 위 스킴론적 올의 포함을 나타내자. 즉 $Y$는
$\kappa(\mathfrak p)$ 위의 스킴이다. Derived Categories of Schemes, Lemma
\ref{perfect-lemma-bounded-on-fibres}에 의해 $Li^*(K)$가 아래로 유계임을
보이면 증명이 끝난다. $G \in D_{perf} (\mathcal{O}_X)$를
$D_\QCoh (\mathcal{O}_X)$를 생성하는 완전 복합체라 하자. Derived Categories
of Schemes, Theorem \ref{perfect-theorem-bondal-van-den-Bergh}를 보라. 다음이
성립한다.
\begin{align*}
R\Hom _{\mathcal{O}_Y}(Li^*(G), Li^*(K))
& =
R\Gamma(Y, Li^*(G ^\vee \otimes ^\mathbf{L} K)) \\
& =
R\Gamma(X, G^\vee \otimes ^{\mathbf{L}} K)
\otimes^\mathbf{L}_A \kappa(\mathfrak{p})
\end{align*}
첫째 등식에서는 $Li^*$가 완전 대상과 쌍대를 보존한다는 사실 및 Cohomology,
Lemma \ref{cohomology-lemma-dual-perfect-complex}를 사용한다. 일부 세부사항은
생략한다. $X$가 $A$ 위에서 평탄이므로 둘째 등식은 Derived Categories of
Schemes, Lemma \ref{perfect-lemma-compare-base-change}에서 따른다. 가정에 의해
이는 $D(\kappa(\mathfrak{p}))$의 완전 대상이다. Derived Categories of Schemes,
Remark \ref{perfect-remark-pullback-generator}에 의해 대상
$Li^*(G) \in D_{perf}(\mathcal{O}_Y)$는 $D_\QCoh(\mathcal{O}_Y)$를 생성한다.
따라서 이제 Derived Categories of Schemes, Proposition
\ref{perfect-proposition-detecting-bounded-below}에 의해 $Li^*(K)$가 아래로
유계이고, 증명이 끝난다.
\end{proof}




\section{복합체의 유한성 성질의 강하}
\label{more-morphisms-section-descent-finiteness}

\noindent
본 절은 Derived Categories of Schemes, 절
\ref{perfect-section-descent-finiteness}의 계속이다.

\begin{lemma}
\label{more-morphisms-lemma-relative-pseudo-coherent-descends-fppf}
$X \to S$가 국소 유한형이라 하자.
$\{f_i : X_i \to X\}$가 스킴의 fppf 덮개라 하자.
$E \in D_\QCoh(\mathcal{O}_X)$ 및 $m \in \mathbf{Z}$라 하자.
그러면 $E$가 $S$에 상대적으로 $m$-유사연접일 필요충분조건은 각
$Lf_i^*E$가 $S$에 상대적으로 $m$-유사연접인 것이다.
\end{lemma}

\begin{proof}
$E$가 $S$에 상대적으로 $m$-유사연접이라고 가정하자.
보조정리 \ref{more-morphisms-lemma-flat-finite-presentation-pseudo-coherent}에 의해 사상
$f_i$들은 유사연접이다. 따라서 보조정리
\ref{more-morphisms-lemma-pull-relative-pseudo-coherent}에 의해 $Lf_i^*E$는
$S$에 상대적으로 $m$-유사연접이다.

\medskip\noindent
역으로, 각 $i$에 대하여 $Lf_i^*E$가 $S$에 상대적으로
$m$-유사연접이라고 가정하자. 보조정리 \ref{more-morphisms-lemma-dominate-fppf}에서와 같이
$S = \bigcup U_j$, $W_j \to U_j$, $W_j = \bigcup W_{j, k}$,
$T_{j, k} \to W_{j, k}$ 및 $S$ 위의 사상
$\alpha_{j, k} : T_{j, k} \to X_{i(j, k)}$를 택한다.
% P05-EN-CH37-0333: frozen source changes k to K and then uses the undefined T_{i,k}.
사상 $T_{j, K} \to S$가 평탄하고 유한 표시이므로, 보조정리
\ref{more-morphisms-lemma-permanence-pseudo-coherent}에 의해 $\alpha_{j, k}$가
유사연접임을 알 수 있다. 따라서
$$
L\alpha_{j, k}^*Lf_{i(j, k)}^*E = L(T_{i, k} \to S)^*E
$$
는 보조정리 \ref{more-morphisms-lemma-pull-relative-pseudo-coherent}에 의해 $S$에
상대적으로 $m$-유사연접이다. 이제 이 성질을 덮개
$\{T_{j, k} \to W_{j, k}\}$, $W_j = \bigcup W_{j, k}$,
$\{W_j \to U_j\}$ 및 $S = \bigcup U_j$를 따라 강하시키고자 한다.
Zariski 덮개에 대해서는 결과가 참이므로($S$에 상대적인
$m$-유사연접성의 정의에 의한다), $L\pi^*E$가 $S$에 상대적으로
유사연접인 하나의 전사 유한 국소 자유 사상 $\pi : Y \to X$가 있는
경우를 가정해도 된다는 뜻이다.
% P05-EN-CH37-0334: frozen source drops the m-qualifier in both descent assertions.
이 경우 보조정리 \ref{more-morphisms-lemma-finite-morphism-relative-pseudo-coherence}에 의해
$R\pi_*L\pi^*E$는 $S$에 상대적으로 유사연접이다(여기서 처음으로
$E$의 코호몰로지 층들이 준연접이라는 사실을 사용한다). 우리는
$R\pi_*L\pi^*E = E \otimes^\mathbf{L}_{\mathcal{O}_X} \pi_*\mathcal{O}_Y$
를 갖는다. 예를 들어 Derived Categories of Schemes, 보조정리
\ref{perfect-lemma-cohomology-base-change}에 의한다. 또한 $X$ 위에서
국소적으로 사상 $\mathcal{O}_X \to \pi_*\mathcal{O}_Y$는 직합인자의
포함이다. 따라서 보조정리
\ref{more-morphisms-lemma-summands-relative-pseudo-coherent}에 의해 결론을 얻는다.
\end{proof}

\begin{lemma}
\label{more-morphisms-lemma-relative-pseudo-coherent-post-compose}
$X \to T \to S$가 스킴의 사상들이라 하자. $T \to S$가 평탄하고
국소 유한 표시이며 $X \to T$가 국소 유한형이라고 가정하자.
$E \in D(\mathcal{O}_X)$ 및 $m \in \mathbf{Z}$라 하자. 그러면 $E$가
$S$에 상대적으로 $m$-유사연접일 필요충분조건은 $E$가 $T$에
상대적으로 $m$-유사연접인 것이다.
\end{lemma}

\begin{proof}
$X$ 위에서 국소적으로 닫힌 몰입 $i : X \to \mathbf{A}^n_T$를 택할 수 있다.
그러면 $\mathbf{A}^n_T \to S$는 평탄하고 국소 유한 표시이다. 따라서
보조정리 \ref{more-morphisms-lemma-composition-relative-pseudo-coherent}를 적용하면
동치가 성립함을 알 수 있다.
\end{proof}

\begin{lemma}
\label{more-morphisms-lemma-relative-pseudo-coherent-descends-fppf-base}
$f : X \to S$가 국소 유한형이라 하자. $\{S_i \to S\}$가 스킴의
fppf 덮개라 하자. $f$의 밑변환을 $f_i : X_i \to S_i$, 사영을
$g_i : X_i \to X$로 나타내자. $E \in D_\QCoh(\mathcal{O}_X)$ 및
$m \in \mathbf{Z}$라 하자. 그러면 $E$가 $S$에 상대적으로
$m$-유사연접일 필요충분조건은 각 $Lg_i^*E$가 $S_i$에 상대적으로
$m$-유사연접인 것이다.
\end{lemma}

\begin{proof}
이는 보조정리 \ref{more-morphisms-lemma-relative-pseudo-coherent-descends-fppf}와
\ref{more-morphisms-lemma-relative-pseudo-coherent-post-compose}에서 형식적으로 따른다.
실제로 $E$가 $S$에 상대적으로 $m$-유사연접이면, 첫 번째 보조정리에
의해 $Lg_i^*E$는 $S$에 상대적으로 $m$-유사연접이고, 따라서 두 번째
보조정리에 의해 $Lg_i^*E$는 $S_i$에 상대적으로 $m$-유사연접이다.
역으로 $Lg_i^*E$가 $S_i$에 상대적으로 $m$-유사연접이면, 두 번째
보조정리에 의해 $Lg_i^*E$는 $S$에 상대적으로 $m$-유사연접이고,
따라서 첫 번째 보조정리에 의해 $E$는 $S$에 상대적으로
$m$-유사연접이다.
\end{proof}






\section{상대 완전 대상}
\label{more-morphisms-section-relatively-perfect}

\noindent
본 절은 Derived Categories of Schemes, 절
\ref{perfect-section-relatively-perfect}의 논의의 계속이다.

\begin{lemma}
\label{more-morphisms-lemma-thickening-pseudo-coherent}
$i : X \to X'$가 스킴의 유한 차수 비후화라 하자.
$K' \in D(\mathcal{O}_{X'})$가 $K = Li^*K'$가 유사연접이 되는
대상이라 하자. 그러면 $K'$는 유사연접이다.
\end{lemma}

\begin{proof}
먼저 $K'$가 준연접 코호몰로지 층들을 가짐을 증명한다. 이를 위해
1차 비후화의 경우로 환원해도 된다. 절 \ref{more-morphisms-section-thickenings}를 보라.
$\mathcal{I} \subset \mathcal{O}_{X'}$를 $X$를 잘라 내는 준연접
아이디얼층이라 하자. 짧은 완전열
$$
0 \to \mathcal{I} \to \mathcal{O}_{X'} \to i_*\mathcal{O}_X \to 0
$$
과 $K'$의 텐서곱을 취하면 완전 삼각형
$$
K' \otimes_{\mathcal{O}_{X'}}^\mathbf{L} \mathcal{I}
\to K' \to
K' \otimes_{\mathcal{O}_{X'}}^\mathbf{L} i_*\mathcal{O}_X
\to
(K' \otimes_{\mathcal{O}_{X'}}^\mathbf{L} \mathcal{I})[1]
$$
을 얻는다. $i_* = Ri_*$이고, 1차 비후화를 다루고 있으므로
$\mathcal{I}$를 준연접 $\mathcal{O}_X$-가군으로 볼 수 있다. 따라서 이를
$$
i_*(K \otimes_{\mathcal{O}_X}^\mathbf{L} \mathcal{I})
\to K' \to
i_*K \to
i_*(K \otimes_{\mathcal{O}_X}^\mathbf{L} \mathcal{I})[1]
$$
로 다시 쓸 수 있다. 각 항을 식별하려면 Cohomology, 보조정리
\ref{cohomology-lemma-projection-formula-closed-immersion}를 사용한다.
$K$가 $D_\QCoh(\mathcal{O}_X)$에 속하므로 $K'$가
$D_\QCoh(\mathcal{O}_{X'})$에 속한다고 결론 내릴 수 있다. 여기에는
Derived Categories of Schemes, 보조정리
\ref{perfect-lemma-pseudo-coherent},
\ref{perfect-lemma-quasi-coherence-tensor-product} 및
\ref{perfect-lemma-quasi-coherence-direct-image}를 사용한다.

\medskip\noindent
$K'$가 $D_\QCoh(\mathcal{O}_{X'})$에 속한다고 가정하자. 문제는 $X'$ 위에서
국소적이므로 $X'$가 아핀이라고 가정해도 된다. $X' = \Spec(A')$,
$X = \Spec(A)$이고 $A = A'/I$이며 $I$가 멱영이라고 하자. 그러면 $K'$는
어떤 대상 $M' \in D(A')$에서 온다. Derived Categories of Schemes,
보조정리 \ref{perfect-lemma-affine-compare-bounded}를 보라. 따라서
Derived Categories of Schemes, 보조정리
\ref{perfect-lemma-pseudo-coherent-affine}와 $K$에 대한 가정에 의해
$M = M' \otimes_{A'}^\mathbf{L} A$는 $D(A)$의 유사연접 대상이다.
그러므로 More on Algebra, 보조정리
\ref{more-algebra-lemma-check-pseudo-coherent-modulo-nilpotent}에 의해
$M'$는 유사연접이다. 이로써 증명이 끝난다.
\end{proof}

\begin{lemma}
\label{more-morphisms-lemma-thickening-relatively-perfect}
스킴의 올곱 도식
$$
\xymatrix{
X \ar[r]_i \ar[d]_f & X' \ar[d]^{f'} \\
Y \ar[r]^j & Y'
}
$$
을 생각하자. $X' \to Y'$가 평탄하고 국소 유한 표시이며
$Y \to Y'$가 유한 차수 비후화라고 가정하자.
$E' \in D(\mathcal{O}_{X'})$라 하자. $E = Li^*(E')$가 $Y$-완전이면
$E'$는 $Y'$-완전이다.
\end{lemma}

\begin{proof}
$E$가 $Y$-완전이라는 것은 $E$가 유사연접이고
$f^{-1}\mathcal{O}_Y$-가군의 복합체로서 국소적으로 유한 Tor 차원을
가진다는 뜻임을 상기하자(Derived Categories of Schemes, 정의
\ref{perfect-definition-relatively-perfect}). 보조정리
\ref{more-morphisms-lemma-thickening-pseudo-coherent}에 의해 $E'$가 유사연접임을 얻는다.
특히 $E'$는 $D_\QCoh(\mathcal{O}_{X'})$에 속한다. Derived Categories of
Schemes, 보조정리 \ref{perfect-lemma-pseudo-coherent}를 보라. $E'$가
국소적으로 유한 Tor 차원을 가짐을 증명하기 위해 $X'$ 위에서 국소적으로
논해도 된다. 따라서
% P05-EN-CH37-0335: frozen source names S',S although the cartesian diagram uses Y',Y.
$X'$, $S'$, $X$, $S$가 아핀이며 각각 환 $A'$, $R'$, $A$, $R$로
주어진다고 가정해도 된다. 그러면 Derived Categories of Schemes,
보조정리 \ref{perfect-lemma-affine-locally-rel-perfect}에 의해 가환대수판으로
환원한다.
% P05-EN-CH37-0336: frozen source leaves the final commutative-algebra citation as a sentence fragment.
가환대수판은 More on Algebra, 보조정리
\ref{more-algebra-lemma-thickening-relatively-perfect}이다.
\end{proof}

\begin{lemma}
\label{more-morphisms-lemma-henselian-relatively-perfect}
$(R, I)$가 환과 Jacobson 근기에 포함되는 아이디얼 $I$로 이루어진
쌍이라 하자. $S = \Spec(R)$ 및 $S_0 = \Spec(R/I)$로 놓자.
$f : X \to S$가 고유하고 평탄하며 유한 표시라 하자.
$X_0 = S_0 \times_S X$로 나타내자. $E \in D(\mathcal{O}_X)$가
유사연접이라 하자. $E$의 $X_0$로의 유도 제한 $E_0$가
$S_0$-완전이면 $E$는 $S$-완전이다.
\end{lemma}

\begin{proof}
유한 아핀 열린 덮개 $X = U_1 \cup \ldots \cup U_n$을 택한다.
각 $i$에 대하여 닫힌 몰입 $U_i \to \mathbf{A}^{d_i}_S$를 택할 수 있다.
$U_{i, 0} = S_0 \times_S U_i$로 놓자. 각 $i$에 대하여 어떤
$a_i, b_i \in \mathbf{Z}$가 존재하여 복합체 $E_0|_{U_{i, 0}}$의
Tor 진폭은 $[a_i, b_i]$ 안에 있다. 점 $x \in X$를 택하자. 어떤 $i$에
대하여 $E_x$의 $R$ 위 Tor 진폭이 $[a_i - d_i, b_i]$ 안에 있음을
보이겠다. Cohomology, 보조정리
\ref{cohomology-lemma-tor-amplitude-stalk}에 의해 Tor 진폭을 줄기들에서
읽어 낼 수 있으므로, 이것으로 증명이 끝난다.

\medskip\noindent
$f$가 고유이므로 $f(\overline{\{x\}})$는 $S$의 닫힌 부분집합이다.
$I$가 Jacobson 근기에 포함되므로,
% P05-EN-CH37-0337: frozen source uses the participle "meeting" where the finite verb "meets" is required.
$f(\overline{\{x\}})$가 닫힌 부분집합 $S_0 \subset S$와 만난다는 것을 알 수 있다.
따라서 $x_0 \in X_0$인 특수화 $x \leadsto x_0$가 있다.
$x_0 \in U_i$인 $i$를 택하면 $x_0 \in U_{i, 0}$이다. 증명의 나머지에서는
이 $i$를 고정한다. $U_i = \Spec(A)$로 쓰자. 그러면 $A$는 평탄한 유한 표시
$R$-대수이며 $d_i$개 변수의 다항식 $R$-대수의 몫이다. 제한 $E|_{U_i}$는
Derived Categories of Schemes, 보조정리
\ref{perfect-lemma-affine-compare-bounded}와
\ref{perfect-lemma-pseudo-coherent-affine}에 의해 $D(A)$의 유사연접 대상
$K$에 대응한다. $E_0$는 $K \otimes_A^\mathbf{L} A/IA$에 대응함에 유의하자.
$x \leadsto x_0$에 대응하는 소 아이디얼들을
$\mathfrak q \subset \mathfrak q_0 \subset A$라 하자. 그러면
$E_x = K_{\mathfrak q}$이고 $K_{\mathfrak q}$는 $K_{\mathfrak q_0}$의
국소화이다. 따라서 $K_{\mathfrak q_0}$가 $R$-가군의 복합체로서
$[a_i - d_i, b_i]$ 안의 Tor 진폭을 가짐을 보이면 충분하다.
$I \subset \mathfrak p_0 \subset R$를 $f(x_0)$에 대응하는 소
아이디얼이라 하자. 그러면
\begin{align*}
K \otimes_R^\mathbf{L} \kappa(\mathfrak p_0)
& =
(K \otimes_R^\mathbf{L} R/I) \otimes_{R/I}^\mathbf{L}
\kappa(\mathfrak p_0) \\
& =
(K \otimes_A^\mathbf{L} A/IA) \otimes_{R/I}^\mathbf{L} \kappa(\mathfrak p_0)
\end{align*}
가 성립한다.
% P05-EN-CH37-0338: frozen source leaves "the second equality because" as a sentence fragment.
두 번째 등식은 $R \to A$가 평탄하기 때문에 성립한다. $a_i, b_i$의
선택에 의해 이 복합체의 코호몰로지는 구간 $[a_i, b_i]$에 속하는
차수에서만 0이 아닐 수 있다. 따라서 마지막으로 More on Algebra,
보조정리 \ref{more-algebra-lemma-lift-from-fibre-relatively-perfect}를
$R \to A$, $\mathfrak q_0$, $\mathfrak p_0$ 및 $K$에 적용하면
결론을 얻는다.
\end{proof}










\section{유리 곡선의 수축}
\label{more-morphisms-section-contracting}

\noindent
본 절에서는 올의 차원이 $\leq 1$이고 $R^1f_*\mathcal{O}_X = 0$인
고유 사상 $f : X \to Y$를 연구한다. 본 절의 제목을 이해하려면
Algebraic Curves, 절
\ref{curves-section-contracting-rational-tails},
\ref{curves-section-contracting-rational-bridges} 및
\ref{curves-section-contracting-to-stable}을 보라.

\begin{lemma}
\label{more-morphisms-lemma-check-h1-fibre-zero}
$f : X \to Y$가 스킴의 고유 사상이라 하자.
$y \in Y$가 $\dim(X_y) \leq 1$인 점이라 하자. 다음 중 하나가 성립한다고 하자.
\begin{enumerate}
\item $R^1f_*\mathcal{O}_X = 0$이다. 또는 더 일반적으로,
\item $y$가 $g$의 상에 속하고 $R^1f'_*\mathcal{O}_{X'} = 0$인 사상
$g : Y' \to Y$가 존재한다. 여기서 $f' : X' \to Y'$는 $g$에 의한
$f$의 밑변환이다.
\end{enumerate}
그러면 $H^1(X_y, \mathcal{O}_{X_y}) = 0$이다.
\end{lemma}

\begin{proof}
보조정리를 증명하기 위해 $Y$를 $y$의 열린 근방으로 바꾸어도 된다.
따라서 $Y$가 아핀이고 $f$의 모든 올의 차원이 $\leq 1$이라고 가정해도
된다. Morphisms, 보조정리
\ref{morphisms-lemma-openness-bounded-dimension-fibres}를 보라. 이 경우
$R^1f_*\mathcal{O}_X$는 유한형 준연접 $\mathcal{O}_Y$-가군이고 그 구성은
임의의 밑변환과 가환한다. Limits, 보조정리
\ref{limits-lemma-proper-top-cohomology-finite-type}와
\ref{limits-lemma-higher-direct-images-zero-above-dimension-fibre}를 보라.
보조정리는 곧바로 따른다.
\end{proof}

\begin{lemma}
\label{more-morphisms-lemma-h1-fibre-zero}
$f : X \to Y$가 스킴의 고유 사상이라 하자. $y \in Y$가
$\dim(X_y) \leq 1$이고 $H^1(X_y, \mathcal{O}_{X_y}) = 0$인 점이라 하자.
그러면 $R^1f_*\mathcal{O}_X|_V = 0$이고 임의의 $Y' \to V$에 의한
밑변환 뒤에도 같은 명제가 성립하는 $y$의 열린 근방 $V \subset Y$가 있다.
\end{lemma}

\begin{proof}
보조정리를 증명하기 위해 $Y$를 $y$의 열린 근방으로 바꾸어도 된다.
따라서 $Y$가 아핀이고 $f$의 모든 올의 차원이 $\leq 1$이라고 가정해도
된다. Morphisms, 보조정리
\ref{morphisms-lemma-openness-bounded-dimension-fibres}를 보라. 이 경우
$R^1f_*\mathcal{O}_X$는 유한형 준연접 $\mathcal{O}_Y$-가군이고 그 구성은
임의의 밑변환과 가환한다. Limits, 보조정리
\ref{limits-lemma-proper-top-cohomology-finite-type}와
\ref{limits-lemma-higher-direct-images-zero-above-dimension-fibre}를 보라.
$Y = \Spec(A)$이고, $y$가 소 아이디얼 $\mathfrak p \subset A$에 대응하며,
$R^1f_*\mathcal{O}_X$가 유한 $A$-가군 $M$에 대응한다고 하자. 그러면
밑변환에 관한 명제에 의해 $H^1(X_y, \mathcal{O}_{X_y}) = 0$이라는 것은
$\mathfrak pM_\mathfrak p = M_\mathfrak p$라는 뜻이다. Nakayama 보조정리에
의해 $M_\mathfrak p = 0$이라고 결론 내린다. $M$이 유한이므로
$M_f = 0$인 $f \in A$, $f \not \in \mathfrak p$를 찾을 수 있다.
따라서 $V$를 주 열린집합 $D(f)$로 택하면 원하는 결과를 얻는다.
\end{proof}

\begin{lemma}
\label{more-morphisms-lemma-globally-generated-vanishing}
$f : X \to Y$가 스킴의 고유 사상이고 모든 $y \in Y$에 대하여
$\dim(X_y) \leq 1$ 및 $H^1(X_y, \mathcal{O}_{X_y}) = 0$이라고 하자.
$\mathcal{F}$가 $X$ 위에서 준연접이라 하자. 그러면
\begin{enumerate}
\item $p > 1$에 대하여 $R^pf_*\mathcal{F} = 0$이고,
\item $Y$ 위의 준연접 가군 $\mathcal{G}$와 전사
$f^*\mathcal{G} \to \mathcal{F}$가 존재하면 $R^1f_*\mathcal{F} = 0$이다.
\end{enumerate}
$Y$가 아핀이면 다음도 성립한다.
\begin{enumerate}
\item[(3)] $p \not \in \{0, 1\}$에 대하여 $H^p(X, \mathcal{F}) = 0$이고,
\item[(4)] $\mathcal{F}$가 전역 생성이면 $H^1(X, \mathcal{F}) = 0$이다.
\end{enumerate}
\end{lemma}

\begin{proof}
(1)의 소멸은 Limits, 보조정리
\ref{limits-lemma-higher-direct-images-zero-above-dimension-fibre}이다.
(2)를 증명하기 위해 $Y$ 위에서 국소적으로 논하고 $Y$가 아핀이라고
가정해도 된다. 그러면 $R^1f_*\mathcal{F}$는 가군
$H^1(X, \mathcal{F})$에 대응하는 $Y$ 위의 준연접 가군이다. 여기서는
$Y$가 아핀이라는 사실, 고차 직상의 준연접성(Cohomology of Schemes,
보조정리 \ref{coherent-lemma-quasi-coherence-higher-direct-images}) 및
Cohomology of Schemes, 보조정리
\ref{coherent-lemma-quasi-coherence-higher-direct-images-application}을
사용한다. $Y$가 아핀이므로 준연접 가군 $\mathcal{G}$는 전역 생성이고,
따라서 $f^*\mathcal{G}$와 $\mathcal{F}$도 전역 생성이다. 이로써 (4)가
(2)를 함의함을 알 수 있다. (3)은 (1)과 방금 언급한 직상의 준연접성에서
따른다. 따라서
% P05-EN-CH37-0339: frozen source says "the prove (4)" instead of "to prove (4)".
남은 것은 (4)를 증명하는 일뿐이다. $\mathcal{F}$가 전역 생성이면 전사
$\bigoplus_{i \in I} \mathcal{O}_X \to \mathcal{F}$가 존재한다. (1)과
코호몰로지의 긴 완전열에 의해 이는 $H^1$ 위의 전사를 유도한다.
보조정리 \ref{more-morphisms-lemma-h1-fibre-zero}에 의해
$R^1f_*\mathcal{O}_X = 0$이므로 $H^1(X, \mathcal{O}_X) = 0$이고,
$H^1(X, -)$는 직합과 가환하므로(Cohomology, 보조정리
\ref{cohomology-lemma-quasi-separated-cohomology-colimit}) 결론을 얻는다.
\end{proof}

\begin{lemma}
\label{more-morphisms-lemma-h1-fibre-zero-check-h0-kappa}
$f : X \to Y$가 스킴의 고유 사상이라고 하자. 다음을 가정하자.
\begin{enumerate}
\item 모든 $y \in Y$에 대하여 $\dim(X_y) \leq 1$이고
$H^1(X_y, \mathcal{O}_{X_y}) = 0$이다.
\item $\mathcal{O}_Y \to f_*\mathcal{O}_X$는 전사이다.
\end{enumerate}
그러면 $f$의 임의의 밑변환 $f' : X' \to Y'$에 대하여
$\mathcal{O}_{Y'} \to f'_*\mathcal{O}_{X'}$는 전사이다.
\end{lemma}

\begin{proof}
$Y$와 $Y'$이 아핀이라고 가정해도 된다. 그러면 $Y'' \to Y$가 아핀
스킴들의 평탄 사상이 되는 닫힌 몰입 $Y' \to Y''$을 택할 수 있다.
평탄 밑변환(Cohomology of Schemes, 보조정리
\ref{coherent-lemma-flat-base-change-cohomology})에 의해 결과는
$X'' \to Y''$에 대하여 성립한다. 따라서 $Y'$이 $Y$의 닫힌 부분스킴이라고
가정해도 된다. $\mathcal{I} \subset \mathcal{O}_Y$를 $Y'$을 잘라 내는
아이디얼이라 하자. 그러면 짧은 완전열
$$
0 \to \mathcal{I}\mathcal{O}_X \to
\mathcal{O}_X \to \mathcal{O}_{X'} \to 0
$$
이 있다. 여기서 $\mathcal{O}_{X'}$를 $X$ 위의 준연접 가군으로 본다.
보조정리 \ref{more-morphisms-lemma-globally-generated-vanishing}에 의해
$H^1(X, \mathcal{I}\mathcal{O}_X) = 0$이다. 따라서
$$
H^0(Y, \mathcal{O}_Y) \to
H^0(Y, f_*\mathcal{O}_X) = H^0(X, \mathcal{O}_X) \to
H^0(X, \mathcal{O}_{X'})
$$
는 원한 대로 전사이다. 첫 번째 화살표는 $Y$가 아핀이고
$\mathcal{O}_Y \to f_*\mathcal{O}_X$가 전사라고 가정했으므로 전사이며,
두 번째 화살표는 위의 짧은 완전열에 딸린 코호몰로지 긴 완전열과 방금
증명한 소멸에 의해 전사이다.
\end{proof}

\begin{lemma}
\label{more-morphisms-lemma-h1-fibre-zero-h0-kappa}
스킴 사상들의 가환 도식
$$
\xymatrix{
X \ar[rr]_f \ar[rd] & & Y \ar[ld] \\
& S
}
$$
을 생각하자. $s \in S$가 점이라 하자. 다음을 가정하자.
\begin{enumerate}
\item $X \to S$는 국소 유한 표시이고 $X_s$의 점들에서 평탄하다.
\item $f$는 고유이다.
\item $f_s : X_s \to Y_s$의 올들의 차원은 $\leq 1$이고
$R^1f_{s, *}\mathcal{O}_{X_s} = 0$이다.
\item $\mathcal{O}_{Y_s} \to f_{s, *}\mathcal{O}_{X_s}$는 전사이다.
\end{enumerate}
그러면 다음을 만족하는 열린집합 $Y_s \subset V \subset Y$가 있다.
(a) $f^{-1}(V)$는 $S$ 위에서 평탄하고,
(b) $y \in V$에 대하여 $\dim(X_y) \leq 1$이고,
(c) $R^1f_*\mathcal{O}_X|_V = 0$이며,
(d) $\mathcal{O}_V \to f_*\mathcal{O}_X|_V$는 전사이다.
또한 (b), (c), (d)는 임의의 $Y' \to V$에 의한 밑변환 뒤에도 성립한다.
\end{lemma}

\begin{proof}
$y \in Y$가 $s$ 위의 점이라 하자. 원하는 성질들을 갖는 $y$의 열린
근방을 찾으면 충분하다. 첫 단계로 $Y$를 보조정리
\ref{more-morphisms-lemma-h1-fibre-zero}에서 얻은 열린집합 $V$로 바꾸어
$R^1f_*\mathcal{O}_X$가 보편적으로 0이 되게 한다(이 보조정리의 가정은
보조정리 \ref{more-morphisms-lemma-check-h1-fibre-zero}에 의해 성립한다). 또한 $Y$를
축소하여 $f$의 모든 올의 차원이 $\leq 1$이 되게 한다(Morphisms,
보조정리 \ref{morphisms-lemma-openness-bounded-dimension-fibres}와 $f$의
고유성을 사용한다). 따라서 $V = Y$일 때와 임의의 밑변환
$Y' \to Y$ 뒤에 (b), (c)가 성립한다고 가정해도 된다. 이제 보조정리
\ref{more-morphisms-lemma-h1-fibre-zero-check-h0-kappa}에 의해 $Y$ 위에서 (d)를
보이면 충분하다. $Y$를 더 축소할 수 있으므로, $Y$와 $S$가 아핀이라고
가정한다.

\medskip\noindent
정리 \ref{more-morphisms-theorem-openness-flatness}에 의해 $X \to S$가 평탄한 열린
부분집합 $U \subset X$가 있고, 가정에 의해 이는 $X_s$를 포함한다.
그러면 $f(X \setminus U)$는 $y$를 포함하지 않는 닫힌 부분집합이다.
따라서 $Y$를 축소한 뒤 $X$가 $S$ 위에서 평탄하다고 가정해도 된다.

\medskip\noindent
$S = \Spec(R)$라 하자. 어떤 집합 $E$에 대한 다항식환
$R[x_e; e \in E]$의 스펙트럼이 $Y'$이 되게 닫힌 몰입 $Y \to Y'$을
택한다. $f$와 $Y \to Y'$의 합성을 $f' : X \to Y'$으로 나타내자.
그러면 (1)--(4)와 (b), (c)는 $f'$과 $s$에 대해서도 성립한다.
% P05-EN-CH37-0340: frozen source duplicates “we” in the following conditional.
$y$의 열린 근방에서 $\mathcal{O}_{Y'} \to f'_*\mathcal{O}_X$가 전사임을 보이면
$\mathcal{O}_Y \to f_*\mathcal{O}_X$에 대해서도 마찬가지이다. 따라서
$Y$가 $R[x_e; e \in E]$의 스펙트럼이라고 가정해도 된다.

\medskip\noindent
이제 $X$와 $Y$는 $S$ 위에서 평탄하다. 그러면 $Y_s$와 $X$는 $Y$ 위에서
Tor-독립이다. 독자에게 직접 증명을 찾아 보기를 권하지만, 이는 꼭짓점이
$X, Y, S, S$인 정사각형과 $s \to S$에 의한 그 밑변환에 보조정리
\ref{more-morphisms-lemma-case-of-tor-independence}를 적용해서도 따른다. 따라서
$$
Rf_{s, *}\mathcal{O}_{X_s} = L(Y_s \to Y)^*Rf_*\mathcal{O}_X
$$
이다. Derived Categories of Schemes, 보조정리
\ref{perfect-lemma-compare-base-change}를 사용했다. 이미 확립한 소멸에 의해
이는 $f_{s, *}\mathcal{O}_{X_s} = (Y_s \to Y)^*f_*\mathcal{O}_X$를
함의한다. 따라서 $\mathcal{O}_Y \to f_*\mathcal{O}_X$는 준연접
$\mathcal{O}_Y$-가군들의 사상이고, $Y_s$로의 당김은 전사이다.
$f_*\mathcal{O}_X$가 유한형 $\mathcal{O}_Y$-가군이라고 주장한다.
이것이 참이면 $\mathcal{O}_Y \to f_*\mathcal{O}_X$의 여핵
$\mathcal{F}$는 $\mathcal{F}_y \otimes \kappa(y) = 0$을 만족하는 유한형
준연접 $\mathcal{O}_Y$-가군이다. Nakayama 보조정리(Algebra, 보조정리
\ref{algebra-lemma-NAK})에 의해 $\mathcal{F}_y = 0$이다. 따라서
$\mathcal{F}$는 $y$의 열린 근방에서 0이고(Modules, 보조정리
\ref{modules-lemma-finite-type-stalk-zero}), 증명이 완성된다.

\medskip\noindent
주장의 증명. 유한 부분집합 $E' \subset E$에 대하여
$Y' = \Spec(R[x_e; e \in E'])$로 놓자. 충분히 큰 $E'$에 대하여 사상
$f' : X \to Y \to Y'$은 고유이다. Limits, 보조정리
\ref{limits-lemma-finite-type-eventually-proper}를 보라. 이하에서는
$E'$과 $Y'$을 고정한다. 이를 그 유한형 $\mathbf{Z}$-부분대수들의
여과 여극한 $R = \colim R_i$로 쓰자. $S_i = \Spec(R_i)$ 및
$Y'_i = \Spec(R_i[x_e; e \in E'])$로 놓자. 충분히 큰 $i$에 대하여 도식
$$
\xymatrix{
X \ar[d] \ar[r]_{f'} & Y' \ar[d] \ar[r] & S \ar[d] \\
X_i \ar[r]^{f'_i} & Y'_i \ar[r] & S_i
}
$$
을 찾을 수 있다. 여기서 두 정사각형은 올곱이고, $X_i$는 $S_i$ 위에서
평탄하며 $X_i \to Y'_i$는 고유이다. Limits, 보조정리
\ref{limits-lemma-descend-finite-presentation},
\ref{limits-lemma-descend-flat-finite-presentation} 및
\ref{limits-lemma-eventually-proper}를 보라. 위와 같은 논증으로 $Y'$과
$X_i$는 $Y'_i$ 위에서 Tor-독립이고, 따라서
$$
R\Gamma(X, \mathcal{O}_X) =
R\Gamma(X_i, \mathcal{O}_{X_i})
\otimes^\mathbf{L}_{R_i[x_e; e \in E']}
R[x_e; e \in E']
$$
이다. 위와 같은 참고문헌을 사용했다. Cohomology of Schemes, 보조정리
\ref{coherent-lemma-proper-over-affine-cohomology-finite}에 의해 복합체
$R\Gamma(X_i, \mathcal{O}_{X_i})$는 Noether 환 $R_i[x_e; e \in E']$의
유도범주에서 유사연접이다(More on Algebra, 보조정리
\ref{more-algebra-lemma-Noetherian-pseudo-coherent}를 보라). 따라서
$R\Gamma(X, \mathcal{O}_X)$는 $R[x_e; e \in E']$의 유도범주에서
유사연접이다. More on Algebra, 보조정리
\ref{more-algebra-lemma-pull-pseudo-coherent}를 보라. 0이 아닐 수 있는
유일한 코호몰로지 가군은 $H^0(X, \mathcal{O}_X)$이므로 이는 유한
$R[x_e; e \in E']$-가군이라고 결론 내린다. More on Algebra, 보조정리
\ref{more-algebra-lemma-n-pseudo-module}를 보라. 이로써 증명이 끝난다.
\end{proof}

\begin{lemma}
\label{more-morphisms-lemma-h1-fibre-zero-h0-flat}
스킴 사상들의 가환 도식
$$
\xymatrix{
X \ar[rr]_f \ar[rd] & & Y \ar[ld] \\
& S
}
$$
을 생각하자. $X \to S$가 평탄하고, $f$가 고유이며, $y \in Y$에 대하여
$\dim(X_y) \leq 1$이고, $R^1f_*\mathcal{O}_X = 0$이라고 가정하자.
그러면 $f_*\mathcal{O}_X$는 $S$-평탄이고, $f_*\mathcal{O}_X$의 구성은
임의의 밑변환 $S' \to S$와 가환한다.
\end{lemma}

\begin{proof}
$Y$와 $S$가 아핀이고 $S = \Spec(A)$라고 가정해도 된다. 준연접
$\mathcal{O}_Y$-가군 $f_*\mathcal{O}_X$가 $S$에 상대적으로 평탄임을
보이려면 $H^0(X, \mathcal{O}_X)$가 $A$ 위에서 평탄임을 보이면 충분하다
(일부 세부사항은 생략한다). 보조정리
\ref{more-morphisms-lemma-globally-generated-vanishing}에 의해 모든 $A$-가군 $M$에 대하여
$H^1(X, \mathcal{O}_X \otimes_A M) = 0$이다. 또한 $\mathcal{O}_X$가
$A$ 위에서 평탄이므로 함자
$M \mapsto H^0(X, \mathcal{O}_X \otimes_A M)$가 완전임을 얻는다.
더욱이 Cohomology, 보조정리
\ref{cohomology-lemma-quasi-separated-cohomology-colimit}에 의해 이 함자는
직합과 가환한다. 그러면 $M$에 함자적으로
$H^0(X, \mathcal{O}_X \otimes_A M) = M \otimes_A H^0(X, \mathcal{O}_X)$
임을 확인하는 것은 연습문제이고, 이로부터 원하는 평탄성이 따른다.
마지막으로 아핀 스킴들의 사상 $S' \to S$가 환 사상 $A \to A'$로
주어지면, 방금 논한 아핀 경우에서
$$
H^0(X \times_S S', \mathcal{O}_{X \times_S S'}) =
H^0(X, \mathcal{O}_X \otimes_A A') = H^0(X, \mathcal{O}_X) \otimes_A A'
$$
를 얻는다. 이는 $f_*\mathcal{O}_X$의 구성이 임의의 밑변환
$S' \to S$와 가환함을 보인다. 일부 세부사항은 생략한다.
\end{proof}

\begin{lemma}
\label{more-morphisms-lemma-h1-fibre-zero-isom}
스킴 사상들의 가환 도식
$$
\xymatrix{
X \ar[rr]_f \ar[rd] & & Y \ar[ld] \\
& S
}
$$
을 생각하자. $s \in S$가 점이라 하자. 다음을 가정하자.
\begin{enumerate}
\item $X \to S$는 국소 유한 표시이고 $X_s$의 점들에서 평탄하다.
\item $Y \to S$는 국소 유한 표시이다.
\item $f$는 고유이다.
\item $f_s : X_s \to Y_s$의 올들의 차원은 $\leq 1$이고
$R^1f_{s, *}\mathcal{O}_{X_s} = 0$이다.
\item $\mathcal{O}_{Y_s} \to f_{s, *}\mathcal{O}_{X_s}$는 동형이다.
\end{enumerate}
그러면 다음을 만족하는 열린집합 $Y_s \subset V \subset Y$가 있다.
(a) $V$는 $S$ 위에서 평탄하고,
(b) $f^{-1}(V)$는 $S$ 위에서 평탄하고,
(c) $y \in V$에 대하여 $\dim(X_y) \leq 1$이고,
(d) $R^1f_*\mathcal{O}_X|_V = 0$이며,
(e) $\mathcal{O}_V \to f_*\mathcal{O}_X|_V$는 동형이다. 또한 (a)--(e)는
임의의 $S' \to S$에 의해 $f^{-1}(V) \to V$를 밑변환한 뒤에도 성립한다.
\end{lemma}

\begin{proof}
$y \in Y_s$라 하자. 언제든 $Y$를 $y$의 열린 근방으로 바꾸어도 된다.
따라서 $Y$와 $S$가 아핀이라고 가정해도 된다. 또한 $X$가 $S$ 위에서
평탄이고, $y \in Y$에 대하여 $\dim(X_y) \leq 1$이고,
$R^1f_*\mathcal{O}_X = 0$이 보편적으로 성립하며,
$\mathcal{O}_Y \to f_*\mathcal{O}_X$가 전사라고 가정해도 된다. 보조정리
\ref{more-morphisms-lemma-h1-fibre-zero-h0-kappa}를 보라. (이 모두를 사용하지는 않는다.)

\medskip\noindent
$S$와 $Y$가 아핀이라고 가정하자.
% P05-EN-CH37-0341: frozen source omits “limit” after “as a cofiltered”.
이를 아핀 Noether 스킴 $S_i$들의 여과 역극한 $S = \lim S_i$로 쓰자.
Limits, 보조정리 \ref{limits-lemma-descend-finite-presentation}에 의해
지표 $0 \in I$와 도식
$$
\xymatrix{
X_0 \ar[rr]_{f_0} \ar[rd] & & Y_0 \ar[ld] \\
& S_0
}
$$
이 존재한다. 이 도식은 스킴의 유한형 사상들로 이루어지고 $S$로
밑변환하면 보조정리의 도식이 된다. $0$을 증가시킨 뒤 $Y_0$가 아핀이고
% P05-EN-CH37-0342: frozen source says X_0 -> S_0 proper although only f_0 : X_0 -> Y_0 descends the proper map.
$X_0 \to S_0$가 고유하다고 가정해도 된다. Limits, 보조정리
\ref{limits-lemma-eventually-proper}와
\ref{limits-lemma-limit-affine}를 보라. $s_0 \in S_0$를 $s$의 상이라 하자.
$Y_s$가 아핀이므로 $R^1f_{s, *}\mathcal{O}_{X_s} = 0$일 필요충분조건은
$H^1(X_s, \mathcal{O}_{X_s}) = 0$인 것이다. $X_s$는 충실히 평탄한 사상
$\kappa(s_0) \to \kappa(s)$에 의한 $X_{0, s_0}$의 밑변환이므로
$H^1(X_{0, s_0}, \mathcal{O}_{X_{0, s_0}}) = 0$이고, 따라서
$R^1f_{0, *}\mathcal{O}_{X_{0, s_0}} = 0$이다. 마찬가지로
$\mathcal{O}_{Y_s} \to f_{s, *}\mathcal{O}_{X_s}$가 동형이므로
$\mathcal{O}_{Y_{0, s_0}} \to f_{0, *}\mathcal{O}_{X_{0, s_0}}$도
동형이다. $X_s \to Y_s$의 올들의 차원이 $1$ 이하이므로 사상
$X_{0, s_0} \to Y_{0, s_0}$에 대해서도 마찬가지이다. 마지막으로
$X \to S$가 평탄이므로 $0$을 증가시킨 뒤 $X_0$가 $S_0$ 위에서
평탄이라고 가정해도 된다. Limits, 보조정리
\ref{limits-lemma-descend-flat-finite-presentation}를 보라. 따라서
$X_0 \to Y_0 \to S_0$와 점 $s_0$에 대하여 보조정리를 증명하면 충분하다.

\medskip\noindent
위의 환원 논증들을 합치면 $S$와 $Y$가 아핀이고, $S$가 Noether이며,
$f$의 올들의 차원이 $\leq 1$이고, $R^1f_*\mathcal{O}_X = 0$이 보편적으로
성립하는 경우로 환원된다. $y \in Y_s$가 점이라 하자. 다음을 주장한다.
$$
\mathcal{O}_{Y, y} \longrightarrow (f_*\mathcal{O}_X)_y
$$
는 동형이다. 이 주장은 보조정리를 함의한다. 실제로
$f_*\mathcal{O}_X$는 연접이므로(Cohomology of Schemes, 명제
\ref{coherent-proposition-proper-pushforward-coherent}), 이 주장에 의해
$Y$를 $y$의 열린 근방으로 바꾸어 동형
$\mathcal{O}_Y \to f_*\mathcal{O}_X$를 얻을 수 있다. 그러면 보조정리
\ref{more-morphisms-lemma-h1-fibre-zero-h0-flat}에 의해 $Y$가 $S$ 위에서 평탄이라고
결론 내린다. 마지막으로 보조정리
\ref{more-morphisms-lemma-h1-fibre-zero-h0-flat}의 마지막 명제에 의해 동형
$\mathcal{O}_Y \to f_*\mathcal{O}_X$는 임의의 밑변환
$S' \to S$ 뒤에도 동형으로 남는다.

\medskip\noindent
주장의 증명. 이미
$\mathcal{O}_{Y, y} \longrightarrow (f_*\mathcal{O}_X)_y$가
전사이고(보조정리 \ref{more-morphisms-lemma-h1-fibre-zero-h0-kappa}),
$(f_*\mathcal{O}_X)_y$가 $\mathcal{O}_{S, s}$-평탄이며(보조정리
\ref{more-morphisms-lemma-h1-fibre-zero-h0-flat}), 유도된 사상
$$
\mathcal{O}_{Y_s, y} =  \mathcal{O}_{Y, y}/\mathfrak m_s\mathcal{O}_{Y, y}
\longrightarrow
(f_*\mathcal{O}_X)_y/\mathfrak m_s (f_*\mathcal{O}_X)_y
\to
(f_{s, *}\mathcal{O}_{X_s})_y
$$
이 보조정리의 가정에 의해 단사임을 안다. 그러면 Algebra, 보조정리
\ref{algebra-lemma-mod-injective}에서
$\mathcal{O}_{Y, y} \longrightarrow (f_*\mathcal{O}_X)_y$가 원한 대로
단사임이 따른다.
\end{proof}

\begin{lemma}
\label{more-morphisms-lemma-bijection-on-Pic}
$f : X \to Y$가 Noether 스킴들의 고유 사상이고,
$f_*\mathcal{O}_X = \mathcal{O}_Y$이며, $f$의 올들의 차원이 $\leq 1$이고,
$y \in Y$에 대하여 $H^1(X_y, \mathcal{O}_{X_y}) = 0$이라고 하자.
그러면 $f^* : \Pic(Y) \to \Pic(X)$는 모든 $y \in Y$에 대하여
$\mathcal{L}|_{X_y} \cong \mathcal{O}_{X_y}$를 만족하는
$\mathcal{L} \in \Pic(X)$들의 부분군 위로의 전단사이다.
\end{lemma}

\begin{proof}
사영 공식(Cohomology, 보조정리
\ref{cohomology-lemma-projection-formula})에 의해
$\mathcal{N} \in \Pic(Y)$에 대하여 $f_*f^*\mathcal{N} \cong \mathcal{N}$임을
알 수 있다. 모든 $y \in Y$에 대하여
$\mathcal{L}|_{X_y} \cong \mathcal{O}_{X_y}$를 만족하는
$\mathcal{L} \in \Pic(X)$에 대해 $\mathcal{N} = f_*\mathcal{L}$가 가역이고
$\mathcal{L} \cong f^*\mathcal{N}$라고 주장한다. 이것으로 증명이 끝난다.

\medskip\noindent
$\mathcal{O}_Y$-가군 $\mathcal{N} = f_*\mathcal{L}$는 Cohomology of Schemes,
명제 \ref{coherent-proposition-proper-pushforward-coherent}에 의해 연접이다.
따라서 이것이 가역 $\mathcal{O}_Y$-가군임을 보이려면 줄기들에서 확인하면
충분하다(Algebra, 보조정리 \ref{algebra-lemma-finite-projective}). Noether
국소환에서 그 완비화로 가는 사상은 충실히 평탄이므로, 완비화
$(f_*\mathcal{L})_y^\wedge$가 자유임을 확인하면 충분하다(Algebra, 절
\ref{algebra-section-completion-noetherian}과 보조정리
\ref{algebra-lemma-finite-projective-descends}를 보라). 이를 위해
Cohomology of Schemes, 보조정리
\ref{coherent-lemma-formal-functions-stalk}에 서술된 형식적 함수 정리를
사용한다. $f_*\mathcal{O}_X = \mathcal{O}_Y$이고 따라서
$(f_*\mathcal{O}_X)_y^\wedge \cong \mathcal{O}_{Y, y}^\wedge$이므로, 각 $n$에
대하여 $\mathcal{L}|_{X_n} \cong \mathcal{O}_{X_n}$임을 보이면 충분하다
% P05-EN-CH37-0343: frozen source omits the closing parenthesis after “compatibly for varying n”.
($n$이 변할 때 서로 양립하게). 보조정리
\ref{more-morphisms-lemma-picard-group-first-order-thickening}에 의해 완전열
$$
H^1(X_y, \mathfrak m_y^n\mathcal{O}_X/\mathfrak m_y^{n + 1}\mathcal{O}_X)
\to \Pic(X_{n + 1}) \to \Pic(X_n)
$$
을 얻는다. 표기는 형식적 함수 정리에서와 같다. 어떤 정수
$r_n \geq 0$에 대하여 전사
$$
\mathcal{O}_{X_y}^{\oplus r_n} \cong
\mathfrak m_y^n/\mathfrak m_y^{n + 1} \otimes_{\kappa(y)} \mathcal{O}_{X_y}
\longrightarrow
\mathfrak m_y^n\mathcal{O}_X/\mathfrak m_y^{n + 1}\mathcal{O}_X
$$
가 있음에 유의하자. $\dim(X_y) \leq 1$이므로 이 전사는 1차 코호몰로지
군들 위의 전사를 유도한다(Cohomology, 명제
\ref{cohomology-proposition-vanishing-Noetherian}에서 따르는 차수
$\geq 2$인 코호몰로지의 소멸을 사용한다). 따라서 완전열의 $H^1$은
0이고, 전이 사상 $\Pic(X_{n + 1}) \to \Pic(X_n)$은 원한 대로 단사이다.

\medskip\noindent
아직 $f^*\mathcal{N} \cong \mathcal{L}$임을 보여야 한다. 같은 방법으로
증명되며 세부사항은 생략한다.
\end{proof}









\section{아핀 층화}
\label{more-morphisms-section-affine-stratifications}

\noindent
이 내용은 \cite{RV}에서 가져왔다. 본 절을 읽기 전에 Topology, 절
\ref{topology-section-stratifications}에서 층화에 관한 내용을 조금 읽어 두라.

\medskip\noindent
$X$가 스킴이면 $X$의 층화는 보통 $X$의 바탕 위상공간의 층화를 뜻한다.
각 층은 국소 닫힌 부분집합이다. Schemes, 주석
\ref{schemes-remark-reduced-induced-locally-closed}을 사용하여 이 층들을
$X$의 축약 국소 닫힌 부분스킴으로 보겠다.

\begin{definition}
\label{more-morphisms-definition-affine-stratification}
$X$가 스킴이라 하자. {\it 아핀 층화}란 각 층 $X_i$가 아핀이고 포함 사상
$X_i \to X$가 아핀인 국소 유한 층화
$X = \coprod_{i \in I} X_i$를 말한다.
\end{definition}

\noindent
포함 사상 $X_i \to X$가 준콤팩트라는 조건 아래에서 층화
$X = \coprod X_i$가 국소 유한이라는 조건은 각 층이 $X$의 국소 구성가능
부분집합이라는 조건과 동치이다. Properties, 보조정리
\ref{properties-lemma-stratification-locally-finite-constructible}를 보라.

\medskip\noindent
% P05-EN-CH37-0344: frozen source says “independent on” instead of “independent of”.
$X_i \to X$가 아핀 사상이라는 조건은 국소 닫힌 부분집합 $X_i$에 부여한
스킴 구조에 무관하다. 보조정리 \ref{more-morphisms-lemma-thicken-property-morphisms}를
보라. 더욱이 $X$가 분리이거나, 더 일반적으로 아핀 대각을 가지고,
$X = \coprod X_i$가 아핀 층들로 이루어진 국소 유한 층화이면 사상
$X_i \to X$는 아핀이다. Morphisms, 보조정리
\ref{morphisms-lemma-affine-permanence}를 보라. 따라서 관심 있는 대부분의
경우 포함 사상 $X_i \to X$의 아핀성 조건을 따로 고려하지 않아도 된다.

\medskip\noindent
층화의 부분순서 지표집합 $I$가 유한인 경우에 흔히 관심이 있다.
부분순서집합 $I$의 {\it 길이}란 $I$의 원소들의 사슬
$i_0 < i_1 < \ldots < i_p$의 길이 $p$들의 상한임을 상기하자.

\begin{lemma}
\label{more-morphisms-lemma-improve-stratification}
$X$가 스킴이고 $X = \coprod_{i \in I} X_i$가 유한 아핀 층화라 하자.
$n$을 $I$의 길이라 하면, 지표집합이 $\{0, \ldots, n\}$인 아핀 층화가
존재한다.
\end{lemma}

\begin{proof}
모든 $i \in I$에 대하여 $X_i$의 폐포가
$\bigcup_{j \leq i} X_j$에 포함되게 하는 $I$ 위의 부분순서가 있음을
상기하자. $I' \subset I$를 $I$의 극대 지표들의 집합이라 하자.
$i \in I'$이면 다른 층들의 폐포의 합집합이 $X_i$의 여집합이므로
$X_i$는 $X$에서 열린집합이다. $U = \bigcup_{i \in I'} X_i$를 $X$의
열린 부분스킴으로 보아 스킴으로서
$U_{red} = \coprod_{i \in I'} X_i$가 되게 하자. 그러면 Schemes,
보조정리 \ref{schemes-lemma-disjoint-union-affines}와 보조정리
\ref{more-morphisms-lemma-thickening-affine-scheme}에 의해 $U$는 아핀 스킴이다.
각 $i \in I'$에 대해 $X_i \to X$가 아핀이므로, 보조정리
\ref{more-morphisms-lemma-thicken-property-morphisms}를 사용한 같은 논증에 의해
$U \to X$도 아핀이다. 여집합 $Z = X \setminus U$에 유도된 축약 스킴
구조를 부여하면 아핀 층화
$Z = \bigcup_{i \in I \setminus I'} X_i$를 갖는다. 여기서는 스킴의 사상
$T \to Z$가 아핀일 필요충분조건이 합성 $T \to X$가 아핀인 것임을
사용한다. 이는 Morphisms, 보조정리
\ref{morphisms-lemma-closed-immersion-affine},
\ref{morphisms-lemma-composition-affine} 및
\ref{morphisms-lemma-affine-permanence}에서 따른다. 부분순서집합
$I \setminus I'$의 길이는 $I$의 길이보다 정확히 하나 작다. 따라서
귀납법에 의해 $Z$는 지표집합이
% P05-EN-CH37-0345: frozen source gives {1,...,n} although the displayed strata are Z_0,...,Z_{n-1}.
$\{1, \ldots, n\}$인 아핀 층화
$Z = Z_0 \amalg \ldots \amalg Z_{n - 1}$를 갖는다. $Z_n = U$로 놓으면
$X$의 원하는 층화를 얻는다.
\end{proof}

\noindent
스킴 $X$가 유한 아핀 층화를 가지면 물론 $X$는 준콤팩트이다. 이로부터
$X$가 준분리이기도 하다는 사실은 조금 덜 자명하다.

\begin{lemma}
\label{more-morphisms-lemma-qc-affine-stratification}
$X$가 스킴이라 하자. 다음은 서로 동치이다.
\begin{enumerate}
\item $X$는 유한 아핀 층화를 갖는다.
\item $X$는 준콤팩트이고 준분리이다.
\end{enumerate}
\end{lemma}

\begin{proof}
$X = \bigcup X_i$가 유한 아핀 층화라 하자. 각 $X_i$가 아핀이고 따라서
준콤팩트이므로 $X$는 준콤팩트이다. $U, V \subset X$가 아핀 열린집합이라
하자. $X_i \to X$가 아핀 사상이므로 $U \cap X_i$와 $V \cap X_i$는
$X_i$의 아핀 열린집합이다. 따라서 $U \cap V \cap X_i$는 아핀 스킴
$X_i$의 아핀 열린집합이다(예를 들어 Schemes, 보조정리
\ref{schemes-lemma-characterize-separated}를 보라). 그러므로
$U \cap V = \coprod U \cap V \cap X_i$는 아핀 층들의 유한 합집합이므로
준콤팩트이다. Schemes, 보조정리
\ref{schemes-lemma-characterize-quasi-separated}에 의해 $X$가 준분리라고
결론 내린다.

\medskip\noindent
$X$가 준콤팩트이고 준분리라고 가정하자. Cohomology of Schemes, 보조정리
\ref{coherent-lemma-induction-principle}의 귀납 원리를 사용하여 $X$가 유한
아핀 층화를 갖는다는 명제를 증명할 수 있다. $X$가 공집합이면 공집합인
아핀 층화를 갖는다. $X$가 공집합이 아닌 아핀 스킴이면 하나의 층으로
이루어진 아핀 층화를 갖는다. 다음으로 $X = U \cup V$이고, $U$가
준콤팩트 열린집합이며 $V$가 아핀 열린집합이고,
% P05-EN-CH37-0346: frozen source says “a finite affine stratifications”.
$U = \bigcup_{i \in I} U_i$와 $U \cap V = \coprod_{j \in J} W_j$가
유한 아핀 층화들이라고 가정하자. $Z = X \setminus V$ 및
$Z' = X \setminus U$로 나타내자. $Z$는 $U$에서 닫혀 있고 $Z'$은
$V$에서 닫혀 있다. $U_i \cap Z$와
$U_i \cap W_j = U_i \times_U W_j$는 $U$ 위에서 아핀인 아핀 스킴이다.
(힌트: $U_i \times_U W_j \to W_j$는 $U_i \to U$의 밑변환이므로
아핀이고, 따라서 $U_i \cap W_j$가 아핀이며, 이에 따라
$U_i \cap W_j \to U_i$가 아핀이고, 결국 $U_i \cap W_j \to U$가
아핀임을 사용하라.) 따라서
$$
U = \coprod\nolimits_{i \in I} (U_i \cap Z) \amalg
\coprod\nolimits_{(i, j) \in I \times J} (U_i \cap W_j)
$$
는 $I \amalg I \times J$ 위의 부분순서를
$i' \leq (i, j) \Leftrightarrow i' \leq i$ 및
$(i', j') \leq (i, j) \Leftrightarrow i' \leq i$이고 $j' \leq j$인 것으로
정한 유한 아핀 층화이다. 또한
$(U_i \cap Z) \times_X V = \emptyset$이고
$(U_i \cap W_j) \times_X V = U_i \cap W_j$는 아핀이다. 사상의 아핀성은
치역 위에서 국소적으로 확인할 수 있고(Morphisms, 보조정리
\ref{morphisms-lemma-characterize-affine}), $U$와 $V$ 양쪽 위에서
아핀성을 가지므로 사상 $U_i \cap Z \to X$와
$U_i \cap W_j \to X$는 아핀이다. 증명을 끝내기 위해 위의 층화에
하나의 층 $Z'$을 더하고, 그 지표를 부분순서집합의 극소 원소로 추가한다.
\end{proof}

\begin{definition}
\label{more-morphisms-definition-affine-stratification-number}
$X$가 공집합이 아닌 준콤팩트 준분리 스킴이라 하자.
{\it 아핀 층화수}란 다음 동치 조건들을 만족하는 가장 작은 정수
$n \geq 0$이다.
\begin{enumerate}
\item $I$의 길이가 $n$인 유한 아핀 층화
$X = \coprod_{i \in I} X_i$가 존재한다.
\item 지표집합이 $\{0, \ldots, n\}$인 아핀 층화
$X = X_0 \amalg X_1 \amalg \ldots \amalg X_n$이 존재한다.
\end{enumerate}
\end{definition}

\noindent
조건들의 동치는 보조정리 \ref{more-morphisms-lemma-improve-stratification}에 의해 성립한다.
유한 아핀 층화의 존재는 보조정리
\ref{more-morphisms-lemma-qc-affine-stratification}에서 증명되었다.

\begin{lemma}
\label{more-morphisms-lemma-affine-stratification-number-bound}
$X$가 $n + 1$개의 아핀 열린집합으로 이루어진 열린 덮개를 갖는 분리
스킴이라 하자. 그러면 $X$의 아핀 층화수는 $n$ 이하이다.
\end{lemma}

\begin{proof}
$X = U_0 \cup \ldots \cup U_n$이 아핀 열린 덮개라 하자. 다음과 같이 놓는다.
$$
X_i = (U_i \cup \ldots \cup U_n) \setminus
(U_{i + 1} \cup \ldots \cup U_n)
$$
그러면 $X_i$는 $U_i$의 닫힌 부분스킴으로서 아핀이다. Morphisms,
보조정리 \ref{morphisms-lemma-affine-permanence}에 의해 사상
$X_i \to X$는 아핀이다. 마지막으로
% P05-EN-CH37-0347: frozen source omits a union operator before X_0.
$\overline{X_i} \subset X_i \cup X_{i - 1} \cup \ldots X_0$이다.
\end{proof}

\begin{lemma}
\label{more-morphisms-lemma-affine-stratification-number-bound-Noetherian}
$X$가 차원이 $\infty > d \geq 0$인 Noether 스킴이라 하자. 그러면
$X$의 아핀 층화수는 $d$ 이하이다.
\end{lemma}

\begin{proof}
$d$에 관한 귀납법을 쓴다. $d = 0$이면 $X$는 아핀이다. Properties,
보조정리 \ref{properties-lemma-locally-Noetherian-dimension-0}을 보라.
$d > 0$이라고 가정하자. $\eta_1, \ldots, \eta_n$을 $X$의 기약 성분들의
일반점들이라 하자(Properties, 보조정리
\ref{properties-lemma-Noetherian-irreducible-components}).
$\eta_1, \ldots, \eta_n$을 포함하는 아핀 열린집합들로 $X$를 덮을 수 있다.
Properties, 보조정리
\ref{properties-lemma-point-and-maximal-points-affine}를 보라. $X$가
준콤팩트이므로 모든 $j = 1, \ldots, m$에 대하여
$\eta_1, \ldots, \eta_n \in U_j$인 유한 아핀 열린 덮개
$X = \bigcup_{j = 1, \ldots, m} U_j$를 찾을 수 있다. $\eta_1, \ldots,
\eta_n$을 포함하는 아핀 열린집합 $U \subset U_1 \cap \ldots \cap U_m$을
택한다(이미 인용한 보조정리에 의해 가능하다). 모든 $j$에 대하여
$U \to U_j$가 아핀이므로 사상 $U \to X$는 아핀이다. Morphisms,
보조정리 \ref{morphisms-lemma-characterize-affine}를 보라.
$Z = X \setminus U$라 하자. 구성에 의해 $\dim(Z) < \dim(X)$이다.
귀납 가정에 의해 $n \leq \dim(Z)$인 $Z$의 아핀 층화
% P05-EN-CH37-0348: frozen source reuses n for the stratification length after n already counts irreducible components.
$Z = \bigcup_{i \in \{0, \ldots, n\}} Z_i$를 찾을 수 있다.
$U = X_{n + 1}$ 및 $i \leq n$에 대하여 $X_i = Z_i$로 놓으면 결론을 얻는다.
\end{proof}

\begin{proposition}
\label{more-morphisms-proposition-vanishing-affine-stratification-number}
$X$가 아핀 층화수가 $n$인 공집합이 아닌 준콤팩트 준분리 스킴이라 하자.
그러면 모든 준연접 $\mathcal{O}_X$-가군 $\mathcal{F}$와 $p > n$에 대하여
$H^p(X, \mathcal{F}) = 0$이다.
\end{proposition}

\begin{proof}
아핀 층화수 $n$에 관한 귀납법으로 증명하겠다. $n = 0$이면 $X$는
아핀이고, 결과는 Cohomology of Schemes, 보조정리
\ref{coherent-lemma-quasi-coherent-affine-cohomology-zero}이다.
$n > 0$이라고 가정하자. 정의 \ref{more-morphisms-definition-affine-stratification-number}에
의해 아핀 스킴 $U$와 아핀 열린 몰입 $j : U \to X$가 존재하여 여집합
$Z$의 아핀 층화수가 $n - 1$이다. $U$와 $j$가 아핀이므로 $p > 0$에 대하여
$H^p(X, j_*(\mathcal{F}|_U)) = 0$이다. Cohomology of Schemes, 보조정리
\ref{coherent-lemma-relative-affine-cohomology}와
\ref{coherent-lemma-relative-affine-vanishing}을 보라. 사상
$\mathcal{F} \to j_*(\mathcal{F}|_U)$의 핵과 여핵을 각각
$\mathcal{K}$와 $\mathcal{Q}$로 나타내자. 그러면 준연접
$\mathcal{O}_X$-가군들의 완전열
$$
0 \to \mathcal{K} \to \mathcal{F} \to j_*(\mathcal{F}|_U)
\to \mathcal{Q} \to 0
$$
을 얻는다(Schemes, 절 \ref{schemes-section-quasi-coherent}를 보라).
이 완전열을 짧은 완전열들로 나누고 코호몰로지 긴 완전열을 사용하는
표준 논증에 의해, $p \geq n$에 대하여
$H^p(X, \mathcal{K}) = 0$과 $H^p(X, \mathcal{Q}) = 0$임을 증명하면
충분하다. $\mathcal{F} \to j_*(\mathcal{F}|_U)$를 $U$로 제한하면 동형이므로
$\mathcal{K}$와 $\mathcal{Q}$는 $Z$ 위에 받침을 가진다. Properties,
보조정리 \ref{properties-lemma-quasi-coherent-colimit-finite-type}에 의해
이 가군들을 유한형 준연접 부분가군들의 여과 여극한으로 쓸 수 있다.
$X$ 위 층들의 코호몰로지가 여과 여극한과 가환한다는 사실(Cohomology,
보조정리 \ref{cohomology-lemma-quasi-separated-cohomology-colimit})을
사용하면, $\mathcal{G}$가 받침이 $Z$에 포함되는 유한형 준연접 가군일 때
$p \geq n$에 대하여 $H^p(X, \mathcal{G}) = 0$임을 보이면 충분하다.
$Z' \subset X$를 $\mathcal{G} \oplus \mathcal{O}_Z$의 스킴론적 받침이라
하자. $\mathcal{G}$를 $Z'$ 위의 준연접 가군으로 보아도 되고 그렇게 한다.
Morphisms, 절 \ref{morphisms-section-support}를 보라. 그러면 $Z'$과 $Z$의
바탕 위상공간은 같으므로 아핀 층화수도 같고, 그 값은 $n - 1$이다.
따라서 귀납 가정에 의해 $p \geq n$에 대하여
$H^p(X, \mathcal{G}) = H^p(Z', \mathcal{G})$는 0이다. 여기서 등식은
Cohomology of Schemes, 보조정리
\ref{coherent-lemma-relative-affine-cohomology}에 의한다.
\end{proof}

\begin{example}
\label{more-morphisms-example-affine-stratification-number}
$k$가 체이고 $X = \mathbf{P}^n_k$가 $k$ 위의 $n$차원 사영공간이라 하자.
Constructions, 보조정리
\ref{constructions-lemma-standard-covering-projective-space}에 의해 보조정리
\ref{more-morphisms-lemma-affine-stratification-number-bound}를 여기에 적용할 수 있다.
따라서 $\mathbf{P}^n_k$의 아핀 층화수는 $n$ 이하이다. 한편
$\mathbf{P}^n_k$ 위의 어떤 준연접 가군은 차수 $n$에서 0이 아닌
코호몰로지를 갖는다. Cohomology of Schemes, 보조정리
\ref{coherent-lemma-cohomology-projective-space-over-ring}를 보라.
명제 \ref{more-morphisms-proposition-vanishing-affine-stratification-number}를 사용하면
$\mathbf{P}^n_k$의 아핀 층화수가 $n$과 같다고 결론 내린다.
\end{example}








\section{보편 열린 사상}
\label{more-morphisms-section-universally-open}

\noindent
보편 열린 사상에 관한 몇 가지 내용이다.

\begin{lemma}
\label{more-morphisms-lemma-test-universally-open}
$f : X \to S$가 스킴의 사상이라 하자. 다음은 서로 동치이다.
\begin{enumerate}
\item $f$는 보편 열린 사상이다.
\item 국소 유한 표시인 모든 사상 $S' \to S$에 대하여 밑변환
$X_{S'} \to S'$는 열린 사상이다.
\item 모든 $n$에 대하여 사상
$\mathbf{A}^n \times X \to \mathbf{A}^n \times S$는 열린 사상이다.
\end{enumerate}
\end{lemma}

\begin{proof}
(1)이 (2)를 함의하고 (2)가 (3)을 함의함은 명백하다. (3)이 (1)을
함의함을 증명하자. 어떤 스킴 사상 $g : T \to S$에 대하여 밑변환
$X_T \to T$가 열린 사상이 아니라고 가정하자. 그러면
$f(U) \subset V$ 및 $g(W) \subset V$이고 $U \times_V W \to W$가 열린
사상이 아닌 아핀 열린집합 $V \subset S$, $U \subset X$, $W \subset T$를
찾을 수 있다. 이로부터
$\mathbf{A}^n \times U \to \mathbf{A}^n \times V$가 열린 사상이 아님을
보이면 $\mathbf{A}^n \times X \to \mathbf{A}^n \times S$도 열린 사상이
아니므로 증명이 끝난다. 따라서 다음 문단에서 증명하는 결과로 환원된다.

\medskip\noindent
$A \to B$가 환 사상이고, 어떤 $A$-대수 $A'$에 대하여
$A' \to B' = A' \otimes_A B$가 스펙트럼의 열린 사상을 유도하지
않는다고 하자. 주 열린집합들이 $\Spec(B')$의 위상의 기저를 이루므로,
어떤 $g \in B'$에 대하여 $D(g)$의 $\Spec(A')$에서의 상이 열리지 않음을
알 수 있다. 어떤 $n$, $a'_i \in A'$, $b_i \in B$에 대하여
$g = \sum_{i = 1, \ldots, n} a'_i \otimes b_i$
로 쓰자. $B[x_1, \ldots, x_n]$의 원소
$h = \sum_{i = 1, \ldots, n} x_i b_i$를 생각하자. 모순을 얻기 위해
$D(h)$가 사상
$$
\Spec(B[x_1, \ldots, x_n]) \longrightarrow \Spec(A[x_1, \ldots, x_n])
$$
아래에서 열린 부분집합으로 간다고 가정하자. 그러면 $D(h)$는 준콤팩트
열린집합 $U \subset \Spec(A[x_1, \ldots, x_n])$ 위로 전사적으로 간다.
$x_i$를 $a'_i$로 보내는 $A$-대수 준동형
$A[x_1, \ldots, x_n] \to A'$를 생각하자. 이는 $h$를 $g$로 보내는
$B$-대수 준동형 $B[x_1, \ldots, x_n] \to B'$도 유도한다. 도식
$$
\xymatrix{
\Spec(B[x_1, \ldots, x_n]) \ar[d] &
\Spec(B') \ar[l] \ar[d] \\
\Spec(A[x_1, \ldots, x_n]) &
\Spec(A') \ar[l]
}
$$
이 올곱이므로 $D(g)$의 $\Spec(A')$에서의 상은 $U$의 $\Spec(A')$에서의
역상과 같고, 따라서 열린집합이다. 이는 원하는 모순이다.
\end{proof}

\begin{lemma}
\label{more-morphisms-lemma-quasi-finite-Noetherian-universally-open}
$f : X \to Y$가 스킴의 사상이라 하자. 다음을 가정하자.
\begin{enumerate}
\item $f$는 국소 준유한이다.
\item $Y$는 기하적으로 단일가지이고 국소 Noether이다.
\item $X$의 모든 기약 성분은 $Y$의 어떤 기약 성분을 지배한다.
\end{enumerate}
그러면 $f$는 보편 열린 사상이다.
\end{lemma}

\begin{proof}
모든 $n$에 대하여 보조정리 \ref{more-morphisms-lemma-number-of-branches-and-smooth}와
Properties, 보조정리 \ref{properties-lemma-number-of-branches-1}에 의해
스킴 $\mathbf{A}^n \times Y$는 기하적으로 단일가지이다. 따라서 모든
$n$에 대하여 사상
$\mathbf{A}^n \times X \to \mathbf{A}^n \times Y$는 보조정리의 가정들을
만족한다. 보조정리 \ref{more-morphisms-lemma-test-universally-open}에 의해 $f$가 열린
사상임을 증명하면 충분하다. Morphisms, 보조정리
\ref{morphisms-lemma-locally-finite-presentation-universally-open}에 의해
일반화가 $f$를 따라 올라감을 보이면 충분하다. $y' \leadsto y$가 $Y$의
점들의 특수화이고 $x \in X$가 $y$로 가는 점이라고 하자. 보조정리
\ref{more-morphisms-lemma-etale-makes-quasi-finite-finite-at-point}에서와 같이 도식
$$
\xymatrix{
u \ar[d] & U \ar[d] \ar[r] & X \ar[d] \\
v & V \ar[r] & Y
}
$$
을 택한다. 여기서 $(V, v) \to (Y, y)$는 기본 \'etale 근방이고,
$U \to V$는 유한이며, $u$는 $v$로 가는 $U$의 유일한 점이고,
$U \subset V \times_Y X$는 열린 부분이며, $v \mapsto y$ 및
$u \mapsto x$이다. $u$를 지나는 $U$의 기약 성분 $E$를 택하자(적어도
하나는 존재한다). $U \to X$가 \'etale이므로 $E$는 $X$의 한 기약
성분으로 가고, 그 성분은 다시 가정에 의해 $Y$의 한 기약 성분을 지배한다.
$U \to V$가 유한이고 따라서 닫힌 사상이므로, $E$의 상 $E' \subset V$는
$v$를 지나고 $Y$의 한 기약 성분을 지배하는 기약 닫힌 부분집합이다.
$V \to Y$가 \'etale이므로 $E'$은 $v$를 지나는 $V$의 한 기약 성분이어야
한다. $Y$가 기하적으로 단일가지이므로 $E'$은 $v$를 지나는 $V$의 유일한
기약 성분이다(보조정리 \ref{more-morphisms-lemma-nr-branches}). $V$가 국소 Noether이므로
$V$를 축소한 뒤 $E' = V$라고 가정해도 된다(집합의 등식).

\medskip\noindent
$V \to Y$가 \'etale이므로 상이 $y' \leadsto y$인 특수화
$v' \leadsto v$를 찾을 수 있다. 위의 논의에 의해 $v'$로 가는
$u' \in U$를 찾을 수 있다. $u$가 $v$로 가는 $U$의 유일한 점이고
$U \to V$가 닫힌 사상이므로 $u' \leadsto u$이다. 마지막으로 $u'$의 상
$x' \in X$는 $x$로 특수화되고 $y'$로 가는 점이므로 증명이 끝난다.
\end{proof}

\begin{lemma}
\label{more-morphisms-lemma-characterize-universally-open-finite}
$A \to B$가 환 사상이라 하자. $B$가 $A$-가군으로서
$b_1, \ldots, b_d \in B$에 의해 생성된다고 하자.
$h = \sum x_ib_i \in B[x_1, \ldots, x_d]$로 놓자. 그러면
$\Spec(B) \to \Spec(A)$가 보편 열린 사상일 필요충분조건은 $D(h)$의
$\Spec(A[x_1, \ldots, x_d])$에서의 상이 열린집합인 것이다.
\end{lemma}

\begin{proof}
$\Spec(B) \to \Spec(A)$가 보편 열린 사상이면 물론 $D(h)$의 상은
열린집합이다. 역으로 $D(h)$의 상 $U$가 열린집합이라고 가정하자.
$A \to A'$이 환 사상이라 하자. 임의의 주 열린집합
$D(g) \subset \Spec(A' \otimes_A B)$의 $\Spec(A')$에서의 상이 열린집합임을
보이면 충분하다. 어떤 $a'_i \in A'$에 대하여
$g = \sum_{i = 1, \ldots, d} a'_i \otimes b_i$로 쓸 수 있다.
% P05-EN-CH37-0349: frozen source switches from d variables to an undefined n in this proof and diagram.
$x_i$를 $a'_i$로 보내는 $A$-대수 준동형
$A[x_1, \ldots, x_n] \to A'$를 생각하자. 이는 $h$를 $g$로 보내는
$B$-대수 준동형 $B[x_1, \ldots, x_n] \to A' \otimes_A B$도 유도한다.
% P05-EN-CH37-0350: frozen source uses undefined B' in the following diagram.
도식
$$
\xymatrix{
\Spec(B[x_1, \ldots, x_n]) \ar[d] &
\Spec(B') \ar[l] \ar[d] \\
\Spec(A[x_1, \ldots, x_n]) &
\Spec(A') \ar[l]
}
$$
이 올곱이므로 $D(g)$의 $\Spec(A')$에서의 상은 $U$의 $\Spec(A')$에서의
역상과 같고, 따라서 열린집합이다.
\end{proof}

\begin{lemma}
\label{more-morphisms-lemma-descend-quasi-finite-universally-open}
$S = \lim S_i$가 전이 사상들이 아핀인 스킴들의 유향계의 극한이라 하자.
$0 \in I$라 하고 $f_0 : X_0 \to Y_0$가 $S_0$ 위의 스킴 사상이라 하자.
$S_0$, $X_0$, $Y_0$가 준콤팩트이고 준분리라고 가정하자.
$f_i : X_i \to Y_i$를 $f_0$의 $S_i$로의 밑변환, $f : X \to Y$를
$f_0$의 $S$로의 밑변환이라 하자. 다음을 가정하자.
\begin{enumerate}
\item $f$는 국소 준유한이고 보편 열린 사상이다.
\item $f_0$는 국소 유한 표시이다.
\end{enumerate}
그러면 $f_i$가 국소 준유한이고 보편 열린 사상이 되는 $i \geq 0$가
존재한다.
\end{lemma}

\begin{proof}
Limits, 보조정리 \ref{limits-lemma-descend-quasi-finite}에 의해 $0$을
증가시킨 뒤 $f_0$가 국소 준유한이라고 가정해도 된다. $x \in X$라 하자.
준유한 사상의 \'etale 국소화에 의해 도식
$$
\xymatrix{
X \ar[d] & U \ar[l] \ar[d] \\
Y & V \ar[l]
}
$$
을 찾을 수 있다. 여기서 $V \to Y$는 \'etale이고,
$U \subset X_V$는 열린 부분이며, $U \to V$는 유한이고, $x$는
$U \to X$의 상에 속한다. 보조정리
\ref{more-morphisms-lemma-etale-makes-quasi-finite-finite-at-point}를 보라. $V$를 축소한
뒤 $V$와 $U$가 아핀이라고 가정해도 된다. $X$가 준콤팩트이므로 이러한
$V$와 $U$들의 유한 분리합을 취하여 $U \to X$가 전사인 위와 같은
도식을 만들 수 있다. Limits, 보조정리
\ref{limits-lemma-descend-finite-presentation},
\ref{limits-lemma-descend-opens},
\ref{limits-lemma-descend-surjective},
\ref{limits-lemma-descend-finite-finite-presentation},
\ref{limits-lemma-descend-etale} 및
\ref{limits-lemma-limit-affine}에 의해, 필요하면 $0$을 증가시킨 뒤 도식
$$
\xymatrix{
X_0 \ar[d] & U_0 \ar[l] \ar[d] \\
Y_0 & V_0 \ar[l]
}
$$
을 가진다고 가정해도 된다. 여기서 $V_0$는 아핀이고,
$V_0 \to Y_0$는 \'etale이며, $U_0 \subset (X_0)_{V_0}$는 열린 부분이고,
$U_0 \to V_0$는 유한이며, $U_0 \to X_0$는 전사이다.
% P05-EN-CH37-0351: frozen source omits the dummy subject before “follows”.
$V_i \to Y_i$가 \'etale이고 따라서 보편 열린 사상이므로, 충분히 큰 $i$에
대하여 $U_i \to V_i$가 보편 열린 사상임을 증명하면 충분하다. 따라서
다음 문단에서 논하는 경우로 환원된다.

\medskip\noindent
$A = \colim A_i$가 환들의 여과 여극한이라 하자. $A_0 \to B_0$가 환
사상이라 하고 $B = A \otimes_{A_0} B_0$ 및
$B_i = A_i \otimes_{A_0} B_0$로 놓자. $A_0 \to B_0$가 유한이고 유한
표시이며 $A \to B$가 보편 열린 사상이라고 가정하자. 충분히 큰 $i$에
대하여 $A_i \to B_i$가 보편 열린 사상임을 보여야 한다. $B_0$를
$A_0$-가군으로 생성하는 $b_{0, 1}, \ldots, b_{0, d} \in B_0$를 택하고,
$B_0[x_1, \ldots, x_d]$에서
$h_0 = \sum_{j = 1, \ldots, d} x_jb_{0, j}$로 놓자. $h_0$의
$B[x_1, \ldots, x_d]$ 및 $B_i[x_1, \ldots, x_d]$에서의 상을 각각
$h$ 및 $h_i$로 나타내자. $A \to B$가 보편 열린 사상이므로 $D(h)$의
$\Spec(A[x_1, \ldots, x_d])$에서의 상 $U$는 열린집합이다. 물론 $U$는
아핀 스킴의 상이므로 준콤팩트이다. 충분히 큰 $i$에 대하여
$\Spec(A[x_1, \ldots, x_d])$에서의 역상이 $U$인 준콤팩트 열린집합
$U_i \subset \Spec(A_i[x_1, \ldots, x_d])$가 존재한다. Limits, 보조정리
\ref{limits-lemma-descend-opens}를 보라. $i$를 증가시킨 뒤 $D(h_i)$가
$U_i$ 안으로 간다고 가정해도 된다. 이는 $D(h_i)$에서 $U_i$의 당김을
생각하여 같은 보조정리를 적용하면 따른다. 마지막으로 $i$를 더 증가시키면
이미 사용한 Limits, 보조정리 \ref{limits-lemma-descend-surjective}에 의해
스킴의 사상 $D(h_i) \to U_i$는 전사가 된다. 보조정리
\ref{more-morphisms-lemma-characterize-universally-open-finite}에 의해
$A_i \to B_i$가 보편 열린 사상이라고 결론 내린다.
\end{proof}

\begin{lemma}
\label{more-morphisms-lemma-count-geometric-fibres}
$f : X \to Y$가 국소 준유한 사상이라 하자. 그러면
\begin{enumerate}
\item 보조정리
\ref{more-morphisms-lemma-base-change-fibres-nr-geometrically-irreducible-components}와
\ref{more-morphisms-lemma-base-change-fibres-nr-geometrically-connected-components}의
함수 $n_{X/Y}$는 서로 일치한다.
\item $X$가 준콤팩트이면 $n_{X/Y}$는 최댓값 $d < \infty$를 갖는다.
\end{enumerate}
\end{lemma}

\begin{proof}
함수들의 일치는 국소 준유한 사상의 (기하학적) 올들이 이산이라는
사실에서 곧바로 따른다. Morphisms, 보조정리
\ref{morphisms-lemma-locally-quasi-finite-fibres}를 보라. 유계성은
Morphisms, 보조정리 \ref{morphisms-lemma-characterize-universally-bounded}와
\ref{morphisms-lemma-locally-quasi-finite-qc-source-universally-bounded}에서
따른다.
\end{proof}

\begin{lemma}
\label{more-morphisms-lemma-count-geometric-fibres-universally-open}
$f : X \to Y$가 분리이고 국소 준유한이며 보편 열린 스킴 사상이라 하자.
$n_{X/Y}$를 보조정리 \ref{more-morphisms-lemma-count-geometric-fibres}에서와 같이 두자.
어떤 $y \in Y$ 및 $d \geq 0$에 대하여 $n_{X/Y}(y) \geq d$이면,
$y$의 한 열린 근방에서 $n_{X/Y} \geq d$이다.
\end{lemma}

\begin{proof}
문제는 $Y$ 위에서 국소적이므로 $Y$가 아핀이라고 가정해도 된다.
$K$가 잉여체 $\kappa(y)$의 대수적 폐포라 하자. 가정은 $(X_y)_K$가
$\geq d$개의 연결 성분을 갖는다는 것이다. 그러면 적당한 준콤팩트
열린집합 $X' \subset X$에 대하여 스킴 $(X'_y)_K$가 $\geq d$개의 연결
성분을 갖는다. 세부사항은 생략한다. $X$를 $X'$으로 바꾼 뒤 $X$가
준콤팩트라고 가정해도 된다. 그러면 $f$는 준유한이다.
$x_1, \ldots, x_n$을 $y$ 위에 놓이는 $X$의 점들이라 하자. 보조정리
\ref{more-morphisms-lemma-etale-splits-off-quasi-finite-part-technical-variant}를 적용하여
\'etale 근방 $(U, u) \to (Y, y)$와 분해
$$
U \times_Y X =
W \amalg
\ \coprod\nolimits_{i = 1, \ldots, n}
\ \coprod\nolimits_{j = 1, \ldots, m_i}
V_{i, j}
$$
를 얻는다. 표기는 인용한 곳에서와 같다. 이 상황에서
$n_{X/Y}(y) = \sum_i m_i$임에 유의하자. 일부 세부사항은 생략한다.
$f$가 보편 열린 사상이므로 모든 $i, j$에 대하여 $V_{i, j} \to U$는
열린 사상이다. 따라서 $U$를 축소한 뒤 모든 $i, j$에 대하여
$V_{i, j} \to U$가 전사라고 가정해도 된다. 이로부터
$n_{U \times_Y X/U} \geq \sum_i m_i = n_{X/Y}(y) \geq d$임이 증명된다.
$n_{X/Y}$의 구성은 밑변환과 양립하므로 증명이 끝난다.
\end{proof}

\begin{lemma}
\label{more-morphisms-lemma-large-open}
$f : X \to Y$가 분리이고 국소 준유한이며 보편 열린 스킴 사상이라 하자.
$n_{X/Y}$를 보조정리 \ref{more-morphisms-lemma-count-geometric-fibres}에서와 같이 두자.
$n_{X/Y}$가 최댓값 $d < \infty$를 가지면 집합
$$
Y_d = \{y \in Y \mid n_{X/Y}(y) = d\}
$$
는 $Y$에서 열려 있고 사상 $f^{-1}(Y_d) \to Y_d$는 유한이다.
\end{lemma}

\begin{proof}
$Y_d$의 열림은 보조정리
\ref{more-morphisms-lemma-count-geometric-fibres-universally-open}에서 곧바로 따른다.
$Y_d$ 위에서 유한임을 증명하기 위해 그 보조정리의 증명 논증을 반복한다.
즉 $y \in Y_d$라 하자. 그러면 $y$ 위에 놓이는 $X$의 점들은 많아야
$d$개이다. $x_1, \ldots, x_n$이 $y$ 위에 놓이는 $X$의 점들이라 하자. 보조정리
\ref{more-morphisms-lemma-etale-splits-off-quasi-finite-part-technical-variant}를 적용하여
\'etale 근방 $(U, u) \to (Y, y)$와 분해
$$
U \times_Y X =
W \amalg
\ \coprod\nolimits_{i = 1, \ldots, n}
\ \coprod\nolimits_{j = 1, \ldots, m_i}
V_{i, j}
$$
를 얻는다. 표기는 인용한 곳에서와 같다. 이 상황에서
$d = n_{X/Y}(y) = \sum_i m_i$임에 유의하자. 일부 세부사항은 생략한다.
$f$가 보편 열린 사상이므로 모든 $i, j$에 대하여 $V_{i, j} \to U$는
열린 사상이다. 따라서 $U$를 축소한 뒤 모든 $i, j$에 대하여
$V_{i, j} \to U$가 전사이고,
% P05-EN-CH37-0352: frozen source says U maps into W, although W is a summand and the argument needs U to map into Y_d.
$U$가 $W$ 안으로 간다고 가정해도 된다. 이로부터
$n_{U \times_Y X/U} \geq \sum_i m_i = d$임이 증명된다. $n_{X/Y}$의
구성은 밑변환과 양립하므로 $n_{U \times_Y X/U} = d$임을 안다. 이는
$W$가 공집합이어야 함을 뜻하며, 따라서 $U \times_Y X \to U$가
유한이라고 결론 내린다. Descent, 보조정리
\ref{descent-lemma-descending-property-finite}에 의해 이는 $X \to Y$가
열린 사상 $U \to Y$의 상 위에서 유한임을 함의한다. 다시 말해 원한 대로
$f$는 $y$의 한 열린 근방 위에서 유한이다.
\end{proof}






\section{가중치}
\label{more-morphisms-section-weightings}

\noindent
본 절의 내용은 \cite[Exposee XVII, 6.2.4]{SGA4}에서 가져왔다.

\medskip\noindent
$\pi : U \to V$가 올들이 유한인 스킴의 국소 준유한 사상이라 하자.
함수 $w : U \to \mathbf{Z}$가 주어지면 함수
$$
\textstyle{\int}_\pi w : V \longrightarrow \mathbf{Z},\quad
v \longmapsto
\sum\nolimits_{u \in U,\ \pi(u) = v} w(u) [\kappa(u) : \kappa(v)]_s
$$
를 정의한다. 체 확대들은 유한이고(Morphisms, 보조정리
\ref{morphisms-lemma-residue-field-quasi-finite}),
$[\kappa' : \kappa]_s$는 분리 차수이며(Fields, 정의
\ref{fields-definition-insep-degree}), $\pi$의 올들이 유한이라고
가정했으므로 합도 유한임에 유의하자. 점 $v \in V$에서
$\int_\pi w$의 값을 계산하는 또 다른 방법은 다음과 같다. 대수적으로
닫힌 체 $k$와 상이 $v$인 사상 $\overline{v} : \Spec(k) \to V$를 택한다.
그러면
$$
(\textstyle{\int}_\pi w)(v) =
\sum\nolimits_{\overline{u} \in U_{\overline{v}}} w(\overline{u})
$$
이다. 여기서 물론 $w(\overline{u})$는 사상 $U_{\overline{v}} \to U$ 아래에서
점 $\overline{u}$의 상 $u$에서의 $w$의 값을 뜻한다.
$\overline{u} \in U_{\overline{v}}$를
$\pi \circ \overline{u} = \overline{v}$를 만족하는 사상
$\overline{u} : \Spec(k) \to U$로 볼 수 있음에 유의하자. 실제로
$U \to V$가 올들이 유한인 국소 준유한 사상이므로, 스킴
$U_{\overline{v}}$는 모든 소 아이디얼이 잉여체 $k$를 갖는 극대
아이디얼인 $k$ 위의 유한차원 대수의 스펙트럼이다.
% P05-EN-CH37-0353: frozen source has a malformed “finite dimension ... and all of whose” construction.
등식이 성립함을 보려면, 주어진 $u$ 위에 놓이는 사상 $\overline{u}$의
수가 Fields, 보조정리 \ref{fields-lemma-separable-degree}에 의해
$[\kappa(u) : \kappa(v)]_s$와 같음에 유의하면 된다.

\begin{lemma}
\label{more-morphisms-lemma-weighting-check-after-etale-base-change}
올곱 도식
$$
\xymatrix{
U \ar[d]_\pi & U' \ar[l]^h \ar[d]^{\pi'} \\
V & V' \ar[l]_g
}
$$
이 주어졌다고 하자. $\pi$가 올들이 유한인 국소 준유한 사상이고
$w : U \to \mathbf{Z}$가 함수이면
$(\int_\pi w) \circ g = \int_{\pi'} (w \circ h)$이다.
\end{lemma}

\begin{proof}
이는 위에서 준 $\int_\pi w$의 두 번째 기술에서 곧바로 따른다. 정의에서
직접 증명하려면 다음 사실을 사용한다. $E/F$가 유한 체 확대이고
$F'/F$가 또 다른 체 확대일 때, $F'$ 위의 유한 체들의 곱으로
$(E \otimes_F F')_{red} = \prod E'_i$라고 쓰면
$$
[E : F]_s = \sum [E'_i : F']_s
$$
이다. 이 등식을 증명하기 위해 대수적으로 닫힌 체 확대 $\Omega/F'$를
택하고 다음을 관찰한다.
\begin{align*}
[E : F]_s
& =
|\Mor_F(E, \Omega)| \\
& =
|\Mor_{F'}(E \otimes_F F', \Omega)| \\
& =
|\Mor_{F'}((E \otimes_F F')_{red}, \Omega)| \\
& =
\sum |\Mor_{F'}(E'_i, \Omega)| \\
& =
\sum [E'_i : F']_s
\end{align*}
여기서는 Fields, 보조정리 \ref{fields-lemma-separable-degree}를 사용했다.
\end{proof}

\begin{definition}
\label{more-morphisms-definition-weighting}
$f : X \to Y$를 국소 준유한 사상이라 하자. $f$의
{\it 가중치(weighting)} 또는 {\it pond\'eration}이란 임의의 도식
$$
\xymatrix{
X \ar[d]_f & U \ar[l]^h \ar[d]^\pi \\
Y & V \ar[l]_g
}
$$
에서 $V \to Y$가 에탈이고, $U \subset X_V$가 열린 부분이며,
$U \to V$가 유한일 때 함수 $\int_\pi (w \circ h)$가 국소 상수인
사상 $w : X \to \mathbf{Z}$이다.
\end{definition}

\noindent
물론 $w = 0$으로 두면 임의의 국소 준유한 사상 $f$의 가중치를
얻지만, 그다지 흥미로운 것은 아니다. 여러 목적에서 가장 흥미로운
것은 $w : X \to \mathbf{Z}_{> 0}$인 {\it 양의} 가중치임이 드러날 것이다.

\begin{lemma}
\label{more-morphisms-lemma-weighting-base-change}
$f : X \to Y$를 국소 준유한 사상이라 하고,
$w : X \to \mathbf{Z}$를 그 가중치라 하자. 사상 $Y' \to Y$에 의한
$f$의 밑변환을 $f' : X' \to Y'$라 하자. 그러면 사영 $X' \to X$와
$w$의 합성 $w' : X' \to \mathbf{Z}$는 $f'$의 가중치이다.
\end{lemma}

\begin{proof}
사상 $f'$에 대하여 정의 \ref{more-morphisms-definition-weighting}에 나오는 것과 같은 도식
$$
\xymatrix{
X' \ar[d]_{f'} & U' \ar[l]^{h'} \ar[d]^{\pi'} \\
Y' & V' \ar[l]_{g'}
}
$$
을 생각하자. 임의의 $v' \in V'$에 대하여
$\int_{\pi'} (w' \circ h')$가 $v'$의 어떤 열린 근방에서 상수임을
보여야 한다. 보조정리
\ref{more-morphisms-lemma-weighting-check-after-etale-base-change}와 에탈 사상이
열린 사상이라는 사실에 따라, $V'$를 $v'$의 임의의 에탈 근방으로
바꾸어도 된다. $V'$를 $v'$의 에탈 근방으로 바꾼 뒤
$U' = U'_1 \amalg \ldots \amalg U'_n$이고, 각 $U'_i$에는 $v'$ 위에
놓이는 유일한 점 $u'_i$가 있으며
$\kappa(u'_i)/\kappa(v')$가 순수 비분리라고 가정할 수 있다.
보조정리 \ref{more-morphisms-lemma-etale-splits-off-quasi-finite-part-technical-variant}를
보라. 따라서 $\int_{U'_i \to V'} w'|_{U'_i}$가 $v'$의 한 근방에서
상수임을 보이면 충분함이 분명하다. 이로써 다음 문단에서 다루는
경우로 환원된다.

\medskip\noindent
$v' \in V'$이고, $U'$에는 $v'$ 위에 놓이면서
$\kappa(u')/\kappa(v')$가 순수 비분리인 유일한 점 $u'$가 있다고
하자. $u'$와 $v'$의 상을 각각 $x \in X$와 $y \in Y$로 쓰자.
위에서 사용한 보조정리에 따라 에탈 근방
$(V, v) \to (Y, y)$와 열린 부분 $U \subset X_V$를 다음과 같이
잡을 수 있다. 즉, $\pi : U \to V$는 유한이고, $v$ 위에 놓이면서
사영 $h : U \to X$에 의해 $x \in X$로 가는 유일한 점 $u \in U$가
있으며, 더욱이 $\kappa(u)/\kappa(v)$는 순수 비분리이다.
사상
$$
U'' = U \times_X U' \longrightarrow V \times_Y V' = V''
$$
을 생각하자. $u$와 $u'$가 모두 $x \in X$로 가므로 $(u,u')$로
가는 점 $u'' \in U''$가 존재한다. $u''$의 상을 $v'' \in V''$로
쓰자. $V', v'$를 $V'', v''$로 바꾼 뒤에는 합성
$V' \to Y' \to Y$가 에탈 근방의 사상
$(V',v') \to (V,v)$을 통해 분해되고, 유도된 사상
$X'_{V'} = X_{V'} \to X_V$가 $u'$를 $u$로 보낸다고 가정할 수
있다. 밑변환 $X'_{V'} = X_{V'}$ 안에는 두 열린 부분스킴, 즉
$U'$와 $U \subset X_V$의 역상 $U_{V'}$가 있다. 구성상 둘 다
$v'$ 위의 유일한 점을 포함하며, 그 점은 어느 경우에도 $u'$이다.
따라서 $V'$를 축소한 뒤 이 두 열린 부분이 같다고 가정할 수 있다.
실제로 $U' \setminus (U' \cap U_{V'})$와
$U_{V'} \setminus (U' \cap U_{V'})$의 $V'$에서의 상은 닫혀 있고,
그 어느 상도 $v'$를 포함하지 않는다. 따라서 $U'=U_{V'}$이고
보조정리 \ref{more-morphisms-lemma-weighting-check-after-etale-base-change}에 나오는
데카르트 도식을 얻는다. 가정에 의해
$\int_\pi (w \circ h)$가 국소 상수이므로 결론이 따른다.
\end{proof}

\begin{lemma}
\label{more-morphisms-lemma-weighting-open}
$f : X \to Y$를 국소 준유한 사상이라 하고
$w : X \to \mathbf{Z}$를 $f$의 가중치라 하자. $X' \subset X$가
열린 부분이면 $w|_{X'}$는 $f|_{X'} : X' \to Y$의 가중치이다.
\end{lemma}

\begin{proof}
정의에서 바로 따른다.
\end{proof}

\begin{lemma}
\label{more-morphisms-lemma-weighting-composition}
$f : X \to Y$와 $g : Y \to Z$를 국소 준유한 사상이라 하자.
$w_f : X \to \mathbf{Z}$를 $f$의 가중치라 하고
$w_g : Y \to \mathbf{Z}$를 $g$의 가중치라 하자. 그러면 함수
$$
X \longrightarrow \mathbf{Z},\quad
x \longmapsto w_f(x) w_g(f(x))
$$
는 $g \circ f$의 가중치이다.
\end{lemma}

\begin{proof}
$x \in X$에 대하여
$w_{g \circ f}(x) = w_f(x) w_g(f(x))$로 두자. 도식
$$
\xymatrix{
X \ar[d]_{g \circ f} & U \ar[l] \ar[d]^\pi \\
Z & W \ar[l]
}
$$
에서 $W \to Z$가 에탈이고, $U \subset X_W$가 열린 부분이며,
$U \to W$가 유한이라고 하자.
$\int_\pi w_{g \circ f}|_U$가 국소 상수임을 보여야 한다.
점 $w \in W$를 잡자. 보조정리
\ref{more-morphisms-lemma-weighting-check-after-etale-base-change}와 에탈 사상이
열린 사상이라는 사실에 따라, $(W,w)$를 에탈 근방으로 바꾼 뒤
$\int_\pi w_{g \circ f}|_U$가 상수임을 보이면 충분하다.
$(W,w)$를 에탈 근방으로 바꾼 뒤
$U = U_1 \amalg \ldots \amalg U_n$이고, 각 $U_i$에는 $w$ 위에
놓이면서 $\kappa(u_i)/\kappa(w)$가 순수 비분리인 유일한 점 $u_i$가
있다고 가정할 수 있다. 보조정리
\ref{more-morphisms-lemma-etale-splits-off-quasi-finite-part-technical-variant}를 보라.
따라서 $\int_{U_i \to W} w_{g \circ f}|_{U_i}$가 $w$의 어떤 에탈
근방에서 상수임을 보이면 충분함이 분명하다. 이로써 다음 문단에서
다루는 경우로 환원된다.

\medskip\noindent
$w \in W$이고, $w$ 위에 놓이면서
$\kappa(u)/\kappa(w)$가 순수 비분리인 유일한 점 $u \in U$가 있다고
하자. 점 $v=f(u)\in Y$를 생각하자. $(W,w)$를 기본 에탈 근방으로
바꾼 뒤에는 $v$의 열린 근방 $V \subset Y_W$가 존재하여
$V \to W$가 유한이라고 가정할 수 있다. 보조정리
\ref{more-morphisms-lemma-etale-makes-quasi-finite-finite-at-point}를 보라.
여기서 $f_W : X_W \to Y_W$는 $f$의 $W$로의 밑변환이며,
$f_W^{-1}(V)\cap U$는 $u$의 열린 근방이다. 따라서 $W$를
Zariski 축소한 뒤에는 $f_W(U)\subset V$라고 가정할 수 있다.
이로써 사상
$$
U \xrightarrow{a} V \xrightarrow{b} W
$$
을 얻는다. $V \to W$가 유한이므로 분리이고, 따라서
$U \to V$는 유한이다. $w_f$와 $w_g$가 각각 $f$와 $g$의
가중치이므로 $\int_a w_f|_U$는 $V$ 위에서 국소 상수이고
$\int_b w_g|_V$는 $W$ 위에서 국소 상수이다. 그러므로 $W$를
한 번 더 축소한 뒤에는 이 함수들이 각각 값 $n$과 $m$을 갖는
상수 함수라고 가정할 수 있다. 따라서 원하는 대로
$\int_\pi w_{g \circ f}|_U = \int_{b \circ a} w_{g \circ f}|_U$는
값 $nm$을 갖는 상수 함수이다.
\end{proof}

\begin{lemma}
\label{more-morphisms-lemma-weighting-universally-open}
$f : X \to Y$를 국소 준유한 사상이라 하고,
$w : X \to \mathbf{Z}$를 가중치라 하자. 모든 $x \in X$에 대하여
$w(x)>0$이면 $f$는 보편 열린 사상이다.
\end{lemma}

\begin{proof}
이 성질은 밑변환에 의해 보존되므로(보조정리
\ref{more-morphisms-lemma-weighting-base-change}), $f$가 열린 사상임을 보이면
충분하다. 또한 $X$를 $X$의 임의의 열린 부분으로 바꿀 수 있으므로
$f(X)$가 열려 있음을 보이면 충분하다. $y\in f(X)$라 하고
$f(x)=y$인 $x\in X$를 잡자. $f(X)$가 $y$의 열린 근방을 포함함을
보이면 충분하며, 이를 위해 $Y$를 $y$의 에탈 근방으로 바꾸어도
된다. 준유한 사상의 에탈 국소화에 따라(절
\ref{more-morphisms-section-etale-localization}), $x$의 열린 근방
$U\subset X$가 존재하여 $\pi=f|_U:U\to Y$가 유한이라고 가정할
수 있다. 그러면 $\int_\pi w|_U$는 국소 상수이고 $y$에서 양의
값을 갖는다. 따라서 $\pi(U)$는 $y$의 열린 근방을 포함하며
증명이 끝난다.
\end{proof}

\begin{lemma}
\label{more-morphisms-lemma-weighting-flat-quasi-finite}
$f : X \to Y$를 스킴의 사상이라 하자. $f$가 국소 준유한이고,
국소 유한 표시이며, 평탄하다고 가정하자. 그러면 $y\in Y$ 위에
놓이는 $x\in X$를
$$
w(x) =
\text{length}_{\mathcal{O}_{X, x}}
(\mathcal{O}_{X, x}/\mathfrak m_y \mathcal{O}_{X, x})
[\kappa(x) : \kappa(y)]_i
$$
로 보내는 규칙으로 주어지는 $f$의 양의 가중치
$w:X\to\mathbf{Z}_{>0}$가 존재한다. 여기서
$[\kappa' : \kappa]_i$는 비분리 차수이다(Fields, 정의
\ref{fields-definition-insep-degree}).
\end{lemma}

\begin{proof}
정의 \ref{more-morphisms-definition-weighting}에 나오는 도식을 생각하자.
$u\in U$의 $X,Y,V$에서의 상을 각각 $x,y,v$라 하자. 다음 두
등식이 성립함을 주장한다.
$$
\text{length}_{\mathcal{O}_{X, x}}
(\mathcal{O}_{X, x}/\mathfrak m_y \mathcal{O}_{X, x}) =
\text{length}_{\mathcal{O}_{U, u}}
(\mathcal{O}_{U, u}/\mathfrak m_v \mathcal{O}_{U, u})
$$
그리고
$$
[\kappa(x) : \kappa(y)]_i =
[\kappa(u) : \kappa(v)]_i
$$
첫 번째 등식은 $\mathcal{O}_{X, x} \to \mathcal{O}_{U, u}$가
$\mathfrak m_y \mathcal{O}_{U, u} = \mathfrak m_v \mathcal{O}_{U, u}$이고
$\mathfrak m_x \mathcal{O}_{U, u} = \mathfrak m_u$인 평탄 국소
준동형이라는 사실과 Algebra, 보조정리
\ref{algebra-lemma-pullback-module}에서 따른다. 여기서 위 등식들은
$\mathcal{O}_{Y, y} \to \mathcal{O}_{V, v}$와
$\mathcal{O}_{X, x} \to \mathcal{O}_{U, u}$가 비분기이기 때문에
성립한다. 두 번째 등식은 $\kappa(v)/\kappa(y)$가 유한 분리
확대이고 $\kappa(u)$가
$\kappa(x) \otimes_{\kappa(y)} \kappa(v)$의 한 인자이므로 비분리
차수가 변하지 않는다는 사실에서 따른다. 이로써 보조정리에 나온
함수의 구성은 에탈 밑변환과 가환함을 알 수 있다. 따라서 문제는
다음 문단의 논의로 환원된다.

\medskip\noindent
$f$가 유한 표시인 유한 평탄 사상이라고 가정하자.
$\int_f w$가 $Y$ 위에서 국소 상수임을 보여야 한다. 실제로
$f$는 유한 국소 자유이고(Morphisms, 보조정리
\ref{morphisms-lemma-finite-flat}), $\int_f w$가 $f$의 차수와
같음을 보이겠다. 이 차수는 $Y$ 위의 국소 상수 함수이다.
실제로 $y\in Y$에 대하여
\begin{align*}
(\textstyle{\int}_f w)(y)
& =
\sum\nolimits_{f(x) = y}
\text{length}_{\mathcal{O}_{X, x}}
(\mathcal{O}_{X, x}/\mathfrak m_y \mathcal{O}_{X, x})
[\kappa(x) : \kappa(y)]_i
[\kappa(x) : \kappa(y)]_s \\
& =
\sum\nolimits_{f(x) = y}
\text{length}_{\mathcal{O}_{X, x}}
(\mathcal{O}_{X, x}/\mathfrak m_y \mathcal{O}_{X, x})
[\kappa(x) : \kappa(y)] \\
& =
\text{length}_{\mathcal{O}_{Y, y}}((f_*\mathcal{O}_X)_y/
\mathfrak m_y (f_*\mathcal{O}_X)_y)
\end{align*}
% P05-EN-CH37-0356
마지막 등식은 Algebra, 보조정리
\ref{algebra-lemma-pushdown-module}에서 따른다. 마지막 수는
원하는 대로 $y$에서 $f_*\mathcal{O}_X$의 계수이다.
\end{proof}

\begin{lemma}
\label{more-morphisms-lemma-weighting-quasi-finite-Noetherian}
$f : X \to Y$를 스킴의 사상이라 하자. 다음을 가정하자.
\begin{enumerate}
\item $f$는 국소 준유한이다.
\item $Y$는 기하적으로 단일가지이고 국소 Noether이다.
\end{enumerate}
그러면 $y \in Y$ 위에 놓이는 $x \in X$를
$\mathcal{O}_{Y, y}^{sh}$ 위에서 $\mathcal{O}_{X, x}^{sh}$의
``일반 분리 차수''로
보내는 규칙으로 주어지는 가중치 $w : X \to \mathbf{Z}_{\geq 0}$가 존재한다.
\end{lemma}

\begin{proof}
Algebra, 보조정리
\ref{algebra-lemma-quasi-finite-strict-henselization}
에 따라 $\mathcal{O}_{Y, y}^{sh} \to \mathcal{O}_{X, x}^{sh}$는
유한이다. $Y$가 기하적으로 단일가지이므로
$\mathcal{O}_{Y, y}^{sh}$에는 유일한 최소 소아이디얼
$\mathfrak p$가 존재한다. More on Algebra, 보조정리
\ref{more-algebra-lemma-geometrically-unibranch}를 보라.
$$
(\kappa(\mathfrak p) \otimes_{\mathcal{O}_{Y, y}^{sh}}
\mathcal{O}_{X, x}^{sh})_{red} =
\prod K_i
$$
를 체들의 유한 곱으로 쓰자.
$w(x) = \sum [K_i : \kappa(\mathfrak p)]_s$로 둔다.

\medskip\noindent
이 정의는 에탈 국소화에 영향을 받지 않음이 분명하므로, $w$가
가중치임을 보이는 문제는 $f$가 유한 사상일 때 $\int_f w$가
국소 상수임을 보이는 문제로 환원된다. $Y$의 일반점 $\eta$에서
$\int_f w$의 값은 바로 $\eta$ 위에서 $X \to Y$의 기하적 올
$X_{\overline{\eta}}$의 점 개수이다. 더구나 $Y$가 단일가지이므로
$Y$의 점 $y$는 유일한 일반점 $\eta$의 특수화이다. 따라서
$(\int_f w)(y)$가 $X_{\overline{\eta}}$의 점 개수와 같음을 보이면
충분하다. $y$의 아핀 근방으로 옮긴 뒤 $X \to Y$가 유한 환 사상
$A \to B$로 주어진다고 가정할 수 있다.
$\mathcal{O}_{Y, y}^{sh}$가 대수적으로 닫힌 체 $k$로 가는 사상
$\kappa(y) \to k$를 사용하여 구성되었다고 하자. 그러면
$$
\mathcal{O}_{Y, y}^{sh} \otimes_A B =
\prod\nolimits_{f(x) = y}
\prod\nolimits_{\varphi \in \Mor_{\kappa(y)}(\kappa(x), k)}
\mathcal{O}_{X, x}^{sh}
$$
가 Algebra, 보조정리 \ref{algebra-lemma-finite-over-henselian}와
위에서 사용한 보조정리에 의해 성립한다.
$\mathcal{O}_{Y, y}^{sh}$의 최소 소아이디얼 $\mathfrak p$는
$\eta$에 대응하는 $A$의 소아이디얼로 간다. 따라서 기하적 올의
점 개수는 체 확대에 의해 변하지 않으므로 원하는 등식이 성립한다.
\end{proof}

\begin{lemma}
\label{more-morphisms-lemma-weighting-specialization}
$f : X \to Y$를 스킴의 국소 준유한 사상이라 하고,
$w : X \to \mathbf{Z}$를 $f$의 가중치라 하자.
$y \leadsto y'$를 $Y$의 점들의 특수화라 하자.
$f(x) = y$인 모든 $x \in X$에 대하여 $w(x) = 0$이라고 가정하자.
그러면 $f(x') = y'$인 모든 $x' \in X$에 대하여 $w(x') = 0$이다.
\end{lemma}

\begin{proof}
$y'$로 가는 점 $x' \in X$를 잡자. 보조정리
\ref{more-morphisms-lemma-etale-makes-quasi-finite-finite-at-point}에 따라 도식
$$
\xymatrix{
U \ar[r] \ar[dr]_\pi & X_V \ar[d] \ar[r] & X \ar[d]^f \\
& V \ar[r]^\varphi & Y
}
$$
을 다음과 같이 잡을 수 있다. 여기서 $V \to Y$는 에탈이고,
$v' \in V$는 $y'$로 가며, $U \subset Y_V$는 열린 부분이고,
$U \to V$는 유한이며, $v'$로 가는 유일한 점 $u' \in U$가
존재하고 이 점은 다시 $X$의 $x'$로 간다.
% P05-EN-CH37-0354
% P05-EN-CH37-0355
$V \to Y$가 열린 사상이므로 $v'$로 특수화되고 $y$로 가는 점
$v \in V$가 존재함에 유의하자. 정의 \ref{more-morphisms-definition-weighting}에
따라 함수 $\int_\pi (w \circ h)$는 국소 상수이다. 여기서
$h : U \to X$는 위쪽 수평 화살표들의 합성이다.
$w$에 대한 가정에 의해 $\int_\pi (w \circ h)$는 $v$에서
영이므로 $v'$에서도 영이다. 한편 $v'$에서 그 값은
$w(x')[\kappa(u'):\kappa(v')]_s$이므로 결론이 따른다.
\end{proof}




\section{가중치에 관한 추가 결과}
\label{more-morphisms-section-more-weightings}

\noindent
가중치의 기본 성질 몇 가지를 더 증명한다. 처음에는 가중치가
매우 제멋대로일 수 있는 듯 보이지만, 실제로는 정의
\ref{more-morphisms-definition-weighting}에서 부과한 조건이 상당히 강함이 드러난다.

\begin{lemma}
\label{more-morphisms-lemma-jumps-w}
$f : X \to Y$를 국소 준유한 사상이라 하고,
$w : X \to \mathbf{Z}$를 $f$의 가중치라 하자. 그러면 함수 $w$의
준위집합들은 $X$에서 국소 구성가능하다.
\end{lemma}

\begin{proof}
아래 증명에서는 보조정리 \ref{more-morphisms-lemma-weighting-open}과
\ref{more-morphisms-lemma-weighting-base-change}를 더 언급하지 않고 사용하겠다.
또한 스킴과 위상공간의 구성가능 부분집합에 관한 초등 성질들을
사용한다. Topology, 절 \ref{topology-section-constructible}와
Properties, 절 \ref{properties-section-constructible}를 보라.
이를 사용하면 이 문제는 $X$와 $Y$ 위에서 국소적임을 알 수 있다.
세부사항은 생략한다.
% P05-EN-CH37-0357
따라서 $X$와 $Y$가 아핀이라고 가정할 수 있다. $w$의 준위집합들의
$X' = Y' \times_Y X$로의 역상이 국소 구성가능하도록 하는 유한 표시
전사 사상 $Y' \to Y$를 찾을 수 있으면, Morphisms, 정리
\ref{morphisms-theorem-chevalley}에 의해 결론이 따른다.

\medskip\noindent
$X$와 $Y$가 아핀이라고 가정하자. $T \to Y$가 유한인 몰입
$X \to T$를 잡을 수 있다. 보조정리
\ref{more-morphisms-lemma-quasi-finite-separated-pass-through-finite}를 보라.
Morphisms, 보조정리 \ref{morphisms-lemma-massage-finite}에 따라
$Y$를 $Y$ 위의 전사 유한 국소 자유 스킴 $Y'$로 바꾸고,
$X$를 $Y' \times_Y X$로 바꾸며, $T$를 $Y' \times_Y T$를 닫힌
부분스킴으로 포함하는 $Y'$ 위의 유한 국소 자유 스킴으로 바꾼 뒤에는
$T$가 $Y$ 위에서 유한 국소 자유이고, $Y$로 동형사상되는 닫힌
부분스킴들 $T_i$를 포함하며, 집합론적으로
$T = \bigcup_{i = 1, \ldots, n} T_i$라고 가정할 수 있다.
$T_i \subset T$는 구성가능 닫힌 부분집합이다. 실제로 이는 스킴의
유한 표시 사상 $Y \to T$의 상이다. 따라서
$I \subset \{1, \ldots, n\}$에 대하여 교집합
$\bigcap_{i \in I} T_i$는 $T$의 구성가능 닫힌 부분집합이고,
그 상은 $Y$의 구성가능 닫힌 부분집합이다.

\medskip\noindent
공집합이 아닌 부분들로 이루어진 서로소 합 분해
$\{1, \ldots, n\} = I_1 \amalg \ldots \amalg I_r$에 대하여,
$T_y = \{x_1, \ldots, x_r\}$가 정확히 $r$개의 점으로 이루어지고
$x_j \in T_i \Leftrightarrow i \in I_j$를 만족하는 점 $y \in Y$들의
부분집합 $Y_{I_1, \ldots, I_r} \subset Y$를 생각하자. 위의 관찰에
따라 이들은 $Y$의 구성가능 분할을 이룬다. $Y$ 위의 유한 표시
아핀 스킴 $Y'$가 존재하여 $Y' \to Y$의 상이 정확히
$Y_{I_1, \ldots, I_r}$가 된다. Algebra, 보조정리
\ref{algebra-lemma-constructible-is-image}를 보라. 따라서 어떤 서로소
합 분해 $\{1, \ldots, n\} = I_1 \amalg \ldots \amalg I_r$에 대하여
$Y = Y_{I_1, \ldots, I_r}$라고 가정할 수 있다. 이 경우
$T(j) = \bigcap_{i \in I_j} T_i$로 두면
$T = T(1) \amalg \ldots \amalg T(r)$은 $T$를 서로소인 닫힌
부분집합들, 따라서 열린 부분집합들로 분해한다. 국소 닫힌 부분스킴
$X$와 교차하면 열린닫힌 부분들로의 유사한 분해
$X = X(1) \amalg \ldots \amalg X(r)$을 얻는다. 사상
$X(j) \to Y$는 몰입이다.
% P05-EN-CH37-0358
$w$가 가중치이므로 $w|_{X(j)}$는 국소 상수이다\footnote{실제로
$f : X \to Y$가 몰입이고 $w$가 $f$의 가중치이면, $w$가 영이 아닌
자취에 제한한 $f$는 열린 사상이다.} 따라서 결론이 따른다.
\end{proof}

\begin{lemma}
\label{more-morphisms-lemma-weighting-on-dense-set}
$f : X \to Y$를 스킴의 국소 준유한 사상이라 하고,
$w : X \to \mathbf{Z}$를 $f$의 가중치라 하자.
$E \subset X$를 조밀한 부분집합이라 하자. 모든 $x \in E$에 대하여
$w(x) = 0$이라고 가정하자. 그러면 $w = 0$이다.
\end{lemma}

\begin{proof}
먼저 보조정리 \ref{more-morphisms-lemma-weighting-specialization}를 읽기를 권한다.
$x \in X$의 상을 $y \in Y$라 하자. $w(x) = 0$임을 보이겠다.
보조정리 \ref{more-morphisms-lemma-etale-makes-quasi-finite-finite-at-point}에 따라 도식
$$
\xymatrix{
U \ar[r] \ar[dr]_\pi & X_V \ar[d] \ar[r] & X \ar[d]^f \\
& V \ar[r]^\varphi & Y
}
$$
을 다음과 같이 잡을 수 있다. 여기서 $V \to Y$는 에탈이고,
$v \in V$는 $y$로 가며, $U \subset Y_V$는 열린 부분이고,
$U \to V$는 유한이며, $v$로 가는 유일한 점 $u \in U$가 존재하고
이 점은 다시 $X$의 $x$로 간다.
% P05-EN-CH37-0359
% P05-EN-CH37-0360
$V$, 따라서 $U$가 아핀이라고 가정해도 좋다. 정의
\ref{more-morphisms-definition-weighting}에 따라 함수 $\int_\pi (w \circ h)$는
국소 상수이다. 여기서 $h : U \to X$는 위쪽 수평 화살표들의 합성이다.
보조정리 \ref{more-morphisms-lemma-jumps-w}에 따라 $w \circ h$가 $0$인 점들의
집합 $U_0$는 $U$의 구성가능 부분집합이고, $h^{-1}(E)$를 포함하므로
$U$에서 조밀하다. 여기서는 $h$가 열린 사상임을 사용했다.
Topology, 보조정리 \ref{topology-lemma-dense-in-constructible}에 따라
$U_0$는 조밀한 열린 부분을 포함한다. 다시 말해 $w \circ h$가
영이 아닌 점들의 집합은 어디에서도 조밀하지 않은 닫힌 부분집합
$Z \subset U$에 포함된다. 그러면 Morphisms, 보조정리
\ref{morphisms-lemma-image-nowhere-dense-finite}에 따라 $\pi(Z)$는
$V$의 어디에서도 조밀하지 않은 닫힌 부분집합이다. 따라서
$\int_\pi (w \circ h)$는 조밀한 열린 부분에서 영이고, 그러므로
모든 곳에서 영이다. 한편 $v$에서 그 값은
$w(x)[\kappa(u):\kappa(v)]_s$이므로 결론이 따른다.
\end{proof}

\begin{lemma}
\label{more-morphisms-lemma-jumps-int-w}
$f : X \to Y$를 유한 표시인 국소 준유한 사상이라 하고,
$w : X \to \mathbf{Z}$를 $f$의 가중치라 하자. 그러면 함수
$\int_f w$의 준위집합들은 $Y$에서 국소 구성가능하다.
\end{lemma}

\begin{proof}
보조정리 \ref{more-morphisms-lemma-weighting-check-after-etale-base-change}에 따라
함수 $\int_f w$의 구성은 임의의 밑변환과 가환하고, 보조정리
\ref{more-morphisms-lemma-weighting-base-change}에 따라 밑변환한 뒤에도 가중치를
얻는다.
% P05-EN-CH37-0361
따라서 전사이고 유한 표시인 $Y' \to Y$를 찾을 수 있으면
$Y'$로 밑변환한 뒤 결과를 증명하는 것으로 충분하다. Morphisms, 정리
\ref{morphisms-theorem-chevalley}를 보라.

\medskip\noindent
이 문제는 $Y$ 위에서 국소적이므로 $Y$가 아핀이라고 가정할 수 있다.
그러면 $f$가 유한 표시이므로 $X$는 준콤팩트이고 준분리이다.
$X = U \cup V$가 준콤팩트 열린 부분들에 의한 덮개라고 하자. 그러면
$$
\textstyle{\int}_f w =
\textstyle{\int}_{f|_U} w|_U +
\textstyle{\int}_{f|_V} w|_V -
\textstyle{\int}_{f|_{U \cap V}} w|_{U \cap V}
$$
가 성립한다. 따라서 $w|_U$, $w|_V$, $w|_{U \cap V}$에 대한 결과를
알면 $w$에 대한 결과도 안다. 귀납 원리(Cohomology of Schemes,
보조정리 \ref{coherent-lemma-induction-principle})에 따라 $X$가
아핀인 경우에 보조정리를 증명하면 충분하다.

\medskip\noindent
$X$와 $Y$가 아핀이라고 가정하자. $T \to Y$가 유한인 열린 몰입
$X \to T$를 잡을 수 있다. 보조정리
\ref{more-morphisms-lemma-quasi-finite-separated-pass-through-finite}를 보라.
여전히 적당한 $Y' \to Y$로 밑변환할 수 있으므로 Morphisms,
보조정리 \ref{morphisms-lemma-massage-finite}를 사용하여 사상
$T \to Y$가 유도하는 모든 잉여체 확대가 자명한 경우로 환원할 수
있다. 그러면 당연히 $X \to Y$가 유도하는 확대도 자명하다.
% P05-EN-CH37-0362
이 상황에서 $\int_f w$는 단순히 각 올에서 $w$의 값들을 더하는
함수이다. $w$의 준위집합들은 $X$에서 국소 구성가능하다(보조정리
\ref{more-morphisms-lemma-jumps-w}). Properties, 보조정리
\ref{properties-lemma-stratification-locally-finite-constructible}에
따라 함수 $w$는 유한 개의 값만을 갖는다. 따라서 Morphisms, 정리
\ref{morphisms-theorem-chevalley}와 정수의 합에 관한 몇 가지 초등
논증으로 결론을 얻는다.
\end{proof}

\begin{lemma}
\label{more-morphisms-lemma-semicontinuous-w}
$f : X \to Y$를 국소 준유한 사상이라 하고,
$w : X \to \mathbf{Z}_{> 0}$를 $f$의 양의 가중치라 하자.
그러면 $w$는 상반연속이다.
\end{lemma}

\begin{proof}
$x \in X$의 상을 $y \in Y$라 하자. 에탈 근방
$(V, v) \to (Y, y)$와 열린 부분 $U \subset X_V$를 다음과 같이
잡자. 즉, $\pi : U \to V$는 유한이고, $v$로 가면서
$\kappa(u)/\kappa(v)$가 순수 비분리인 유일한 점 $u \in U$가
존재한다. 보조정리
\ref{more-morphisms-lemma-etale-makes-quasi-finite-finite-multiple-points-var}를 보라.
그러면 $(\int_\pi w|_U)(v) = w(u)$이다. 정의
\ref{more-morphisms-definition-weighting}에 따라 $V$를 $v$의 근방으로 바꾼 뒤
모든 $u' \in U$에 대하여
$w|_U(u') \leq w|_U(u) = w(x)$라고 할 수 있다.
% P05-EN-CH37-0363
실제로 $w|_U(u')$는 $(\int_\pi w|_U)(\pi(u'))$의 식에서 한
합항으로 나타난다. 에탈 사상 $U \to X$가 열려 있으므로 이것으로
보조정리가 증명된다.
\end{proof}

\begin{lemma}
\label{more-morphisms-lemma-semicontinuous-int-w}
$f : X \to Y$를 올들이 유한인 분리 국소 준유한 사상이라 하고,
$w : X \to \mathbf{Z}_{> 0}$를 $f$의 양의 가중치라 하자.
그러면 $\int_f w$는 하반연속이다.
\end{lemma}

\begin{proof}
$y \in Y$라 하자. $y$ 위에 놓이는 점들을
$x_1, \ldots, x_r \in X$라 하자. 보조정리
\ref{more-morphisms-lemma-etale-splits-off-quasi-finite-part-technical-variant}를
적용하여 에탈 근방 $(U, u) \to (Y, y)$와 분해
$$
U \times_Y X =
W \amalg
\ \coprod\nolimits_{i = 1, \ldots, n}
\ \coprod\nolimits_{j = 1, \ldots, m_i}
V_{i, j}
$$
를 얻는다. 여기서 분해는 앞서 인용한 결과에서와 같다.
$(\int_f w)(y) = \sum w(v_{i, j})$이고
$w(v_{i, j}) = w(x_i)$임에 유의하자. 정의에 따라
$\int_{V_{i, j} \to U} w|_{V_{i, j}}$는 국소 상수이므로 $U$를
축소한 뒤 이 함수들의 값이 $w(v_{i, j})$인 상수라고 가정할 수 있다.
따라서
$$
\textstyle{\int}_{U \times_Y X \to U} w|_{U \times_Y X} =
\textstyle{\int}_{W \to U} w|_W +
\sum \textstyle{\int}_{V_{i, j} \to U} w|_{V_{i, j}} =
\textstyle{\int}_{W \to U} w|_W + (\int_f w)(y)
$$
를 얻는다. 이는 $\geq (\int_f w)(y)$이고, $U \to Y$가 열려 있으며
적분의 구성이 밑변환과 가환하므로 결론이 따른다(보조정리
\ref{more-morphisms-lemma-weighting-check-after-etale-base-change}).
\end{proof}

\begin{lemma}
\label{more-morphisms-lemma-max-int-w}
$f : X \to Y$를 $X$가 준콤팩트인 국소 준유한 사상이라 하고,
$w : X \to \mathbf{Z}$를 $f$의 가중치라 하자. 그러면
$\int_f w$는 최댓값을 갖는다.
\end{lemma}

\begin{proof}
보조정리 \ref{more-morphisms-lemma-jumps-w}와 Properties, 보조정리
\ref{properties-lemma-stratification-locally-finite-constructible}에
따라 $w$는 $X$ 위에서 유한 개의 값만을 갖는다. Morphisms, 보조정리
\ref{morphisms-lemma-locally-quasi-finite-qc-source-universally-bounded}에
따라 $X \to Y$의 기하적 올들의 크기는 유계이다. 따라서
$\int_f w$는 유계이다.
\end{proof}

\begin{lemma}
\label{more-morphisms-lemma-max-int-finite}
$f : X \to Y$를 분리 국소 준유한 사상이라 하고,
$w : X \to \mathbf{Z}_{> 0}$를 $f$의 양의 가중치라 하자.
$\int_w f$가 최댓값 $d$를 갖는다고 가정하고,
$(\int_f w)(y) = d$인 점 $y$들의 열린집합을 $Y_d \subset Y$라 하자.
그러면 사상 $f^{-1}(Y_d) \to Y_d$는 유한이다.
% P05-EN-CH37-0364
\end{lemma}

\begin{proof}
보조정리 \ref{more-morphisms-lemma-semicontinuous-int-w}에 따라 $Y_d$는 열려 있다.
$y \in Y_d$라 하고, $y$ 위에 놓이는 $X$의 점들을
$x_1, \ldots, x_n$이라 하자. 보조정리
\ref{more-morphisms-lemma-etale-splits-off-quasi-finite-part-technical-variant}를
적용하여 에탈 근방 $(U, u) \to (Y, y)$와 분해
$$
U \times_Y X =
W \amalg
\ \coprod\nolimits_{i = 1, \ldots, n}
\ \coprod\nolimits_{j = 1, \ldots, m_i}
V_{i, j}
$$
를 얻는다. 여기서 분해는 앞서 인용한 결과에서와 같다.
$d = \sum w(v_{i, j})$이고 $w(v_{i, j}) = w(x_i)$임에 유의하자.
정의에 따라 $\int_{V_{i, j} \to U} w|_{V_{i, j}}$는 국소 상수이므로
$U$를 축소한 뒤 이 함수들의 값이 $w(v_{i, j})$인 상수라고 가정할
수 있다. 따라서
$$
\textstyle{\int}_{U \times_Y X \to U} w|_{U \times_Y X} =
\textstyle{\int}_{W \to U} w|_W +
\sum \textstyle{\int}_{V_{i, j} \to U} w|_{V_{i, j}} =
\textstyle{\int}_{W \to U} w|_W + (\int_f w)(y)
$$
를 얻는다. 이는 $\geq (\int_f w)(y) = d$이므로 $W$는 공집합이어야
한다.
% P05-EN-CH37-0365
따라서 $U \times_Y X \to U$는 유한이다. Descent, 보조정리
\ref{descent-lemma-descending-property-finite}에 따라
$X \to Y$는 열린 사상 $U \to Y$의 상 위에서 유한이다. 다시 말해
원하는 대로 $f$는 $y$의 어떤 열린 근방 위에서 유한이다.
\end{proof}

\begin{lemma}
\label{more-morphisms-lemma-open-and-closed-in-finite}
$A \to B$를 유한이고 유한 표시인 환 사상이라 하자. 유한 표시
환 사상 $A \to A_{univ}$와 멱등원
$e_{univ} \in B \otimes_A A_{univ}$가 존재하여, 임의의 환 사상
$A \to A'$와 멱등원 $e \in B \otimes_A A'$에 대하여 $e_{univ}$를
$e$로 보내는 환 사상 $A_{univ} \to A'$가 존재한다.
\end{lemma}

\begin{proof}
$B$를 $A$-가군으로 생성하는 $b_1, \ldots, b_n \in B$를 잡자.
각 $i$에 대하여 $B$에서 $P_i(b_i) = 0$인 일계수 다항식
$P_i \in A[x]$를 잡자. Algebra, 보조정리
\ref{algebra-lemma-finite-is-integral}를 보라. 그러면 $B$는 유한 자유
$A$-대수 $B' = A[x_1, \ldots, x_n]/(P_1(x_1), \ldots, P_n(x_n))$의
몫이다. 전사 $B' \to B$의 핵을 $J \subset B'$라 하자. $B$가
유한 생성 $A$-대수이므로 $J =(f_1, \ldots, f_m)$은 유한 생성이다.
Algebra, 보조정리 \ref{algebra-lemma-compose-finite-type}를 보라.
% P05-EN-CH37-0366
$B'$의 $A$-기저 $b'_1, \ldots, b'_N$을 잡자. 대수
$$
A_{univ} = A[z_1, \ldots, z_N, y_1, \ldots, y_m]/I
$$
를 생각하자. 여기서 $I$는 식
$$
(\sum z_j b'_j)^2 - \sum z_j b'_j + \sum y_k f_k
$$
을 $B'[z_1, \ldots, z_N, y_1, \ldots, y_m]$의 기저 원소
$b'_1, \ldots, b'_N$에 관해 전개했을 때
$A[z_1, \ldots, z_n, y_1, \ldots, y_m]$에 놓이는 계수들이 생성하는
아이디얼이다.
% P05-EN-CH37-0367
구성에 따라 원소 $\sum z_j b'_j$는 대수
$B \otimes_A A_{univ}$의 멱등원 $e_{univ}$로 간다. 더구나
$e \in B \otimes_A A'$가 멱등원이면, $e$를
$B' \otimes_A A'$의 $\sum b'_j \otimes a'_j$ 꼴 원소로 올릴 수
있고, 다음 식이
$$
(\sum b'_j \otimes a'_j)^2 - \sum b'_j \otimes a'_j + \sum f_k \otimes a''_k
$$
$B' \otimes_A A'$에서 영이 되도록 하는 $a''_k \in A'$를 찾을 수
있다. 따라서 $z_j$를 $a'_j$로, $y_k$를 $a''_k$로 보내며
$e_{univ}$를 $e$로 보내는 $A$-대수 사상 $A_{univ} \to A$를 얻는다.
% P05-EN-CH37-0368
이것으로 증명이 끝난다.
\end{proof}

\begin{lemma}
\label{more-morphisms-lemma-open-and-closed-in-quasi-finite}
$X \to Y$를 준유한이고 유한 표시인 아핀 스킴의 사상이라 하자.
유한 표시 사상 $Y_{univ} \to Y$와 열린 부분스킴
$U_{univ} \subset Y_{univ} \times_Y X$가 존재하여
$U_{univ} \to Y_{univ}$가 유한이고 다음 성질을 만족한다. 임의의
아핀 스킴의 사상 $Y' \to Y$와 $U' \to Y'$가 유한인 열린 부분스킴
$U' \subset Y' \times_Y X$가 주어지면, $U_{univ}$의 역상이 $U'$인
사상 $Y' \to Y_{univ}$가 존재한다.
\end{lemma}

\begin{proof}
유한형 사상이 준유한일 필요충분조건은 그 상대 차원이 $0$인
것임을 상기하자. Morphisms, 보조정리
\ref{morphisms-lemma-locally-quasi-finite-rel-dimension-0}를 보라.
보조정리 \ref{more-morphisms-lemma-Noetherian-approximation-dimension-d}를 $d = 0$으로
적용하여 $X$와 $Y$가 Noether인 경우로 환원한다. $X' \to Y$가
유한인 열린 몰입 $X \to X'$를 잡을 수 있다. Algebra, 보조정리
\ref{algebra-lemma-quasi-finite-open-integral-closure}를 보라.
$Y' \to Y$와 $U'$가 (2)에서와 같이 주어졌다면
% P05-EN-CH37-0369
$$
U' \to Y' \times_Y X \to Y' \times_Y X'
$$
은 $Y'$ 위의 유한 스킴들 사이의 열린 몰입이고, 따라서 닫힌
몰입이기도 하다. 그러므로 $U'$는
$$
\Gamma(Y', \mathcal{O}_{Y'})
\otimes_{\Gamma(Y, \mathcal{O}_Y)}
\Gamma(X', \mathcal{O}_{X'})
$$
의 한 멱등원에 대응하며, 그에 대응하는 열린닫힌 부분집합은 열린
부분 $Y' \times_Y X$에 포함된다. 보조정리
\ref{more-morphisms-lemma-open-and-closed-in-finite}에서 환 사상
$\Gamma(Y, \mathcal{O}_Y) \to \Gamma(X', \mathcal{O}_{X'})$에 대해
구성한 쌍을 $Y'_{univ} \to Y$와 멱등원
$$
e'_{univ} \in
\Gamma(Y_{univ}, \mathcal{O}_{Y_{univ}})
\otimes_{\Gamma(Y, \mathcal{O}_Y)}
\Gamma(X', \mathcal{O}_{X'})
$$
이라 하자.
% P05-EN-CH37-0370
여기서는 $Y$가 Noether라는 사실을 사용하여 $X'$가 $Y$ 위에서
유한 표시임을 안다. 대응하는 열린닫힌 부분스킴을
$U'_{univ} \subset Y'_{univ} \times_Y X'$라 하자. 그러면
$$
U'_{univ} \setminus Y'_{univ} \times_Y X
$$
는 $U'_{univ}$의 닫힌 부분집합이고, 따라서 그 상
$T \subset Y'_{univ}$도 닫혀 있다.
$Y_{univ} = Y'_{univ} \setminus T$로 두고 $U_{univ}$를
$U'_{univ}$의 $Y_{univ} \times_Y X$로의 제한으로 두면 보조정리가
성립함을 알 수 있다.
\end{proof}

\begin{lemma}
\label{more-morphisms-lemma-descend-weighting}
$Y = \lim Y_i$를 아핀 스킴들의 유향 극한이라 하자. $0 \in I$라
하고, $f_0 : X_0 \to Y_0$를 준유한이고 유한 표시인 아핀 스킴의
사상이라 하자. $f : X  \to Y$와 $i \geq 0$에 대한
$f_i : X_i \to Y_i$를 $f_0$의 밑변환들이라 하자.
$w : X \to \mathbf{Z}$가 $f$의 가중치이면, 충분히 큰 $i$에 대하여
$X$로 당긴 것이 $w$인 $f_i$의 가중치 $w_i : X_i \to \mathbf{Z}$가
존재한다.
\end{lemma}

\begin{proof}
보조정리 \ref{more-morphisms-lemma-jumps-w}에 따라 $w$의 준위집합들은 $X$의
구성가능 부분집합 $E_k$들이다. Properties, 보조정리
\ref{properties-lemma-stratification-locally-finite-constructible}에
따라 이는 함수 $w$가 유한 개의 값만을 가짐을 뜻한다. 따라서 어떤
$i$가 존재하여 모든 $k$에 대해 $E_k$가 $X_i$의 구성가능 부분집합
$E_{i, k}$로 내려가고, 더 나아가 $X_i = \coprod E_{i, k}$라고
가정할 수 있다.
% P05-EN-CH37-0371
이는 $X$의 위상공간이 스펙트럴 전이 사상들로 이루어진 유향계의
스펙트럴 공간 $X_i$들의 위상공간 범주에서의 극한이기 때문이다.
Limits, 절 \ref{limits-section-descent}와 Topology, 절
\ref{topology-section-limits-spectral}을 보라.
준위집합들이 구성가능 집합 $E_{i, k}$가 되도록
$w_i : X_i \to \mathbf{Z}$를 정의한다.

\medskip\noindent
보조정리 \ref{more-morphisms-lemma-open-and-closed-in-quasi-finite}에서와 같은
$Y_{i, univ} \to Y_i$와
$U_{i, univ} \subset Y_{i, univ} \times_{Y_i} X_i$를 잡자.
이 구성의 보편 성질에 따라 $w_i$가 가중치임을 보이려면
$$
\tau_i = \textstyle{\int}_{U_{i, univ} \to Y_{i, univ}} w_i|_{U_{i, univ}}
$$
가 $Y_{i, univ}$ 위에서 국소 상수임을 보이면 충분하다. 보조정리
\ref{more-morphisms-lemma-jumps-int-w}에 따라 이 함수의 준위집합들은 구성가능하지만
함수는 아직 국소 상수가 아닐 수 있다.
$Y_{univ} = Y_{i, univ} \times_{Y_i} Y$로 두고
$U_{univ} \subset Y_{univ} \times_Y X$를 $U_{i, univ}$의 역상이라
하자. $w$를 $Y_{univ} \times_Y X$로 당긴 것은
$Y_{univ} \times_Y X \to Y_{univ}$의 가중치이므로(보조정리
\ref{more-morphisms-lemma-weighting-base-change}), 실제로
$$
\tau = \textstyle{\int}_{U_{univ} \to Y_{univ}} w_i|_{U_{univ}}
$$
는 $Y_{univ}$ 위에서 국소 상수이다.
% P05-EN-CH37-0372
따라서 $\tau$의 준위집합들은 열린닫힌 집합이다. 마지막으로
$Y_{univ} = \lim_{i' \geq i} Y_{i', univ}$이고, $\tau$의
준위집합들은 마찬가지로 정의된 $\tau_{i'}$의 준위집합들의
역극한이다. 그러므로 위 참고문헌들에 따라 충분히 큰 $i'$에
대하여 $\tau_{i'}$의 준위집합들도 열려 있다. 그러한 지표 $i'$에
대하여 원하는 대로 $w_{i'}$가 $f_{i'}$의 가중치라고 결론 내린다.
\end{proof}







\section{가중치와 아핀 층화수}
\label{more-morphisms-section-bounds-asn}

\noindent
이 절에서는 특정한 종류의 아핀 스킴 덮개를 갖는 스킴의 아핀 층화수에
대한 상계를 제시한다.

\begin{lemma}
\label{more-morphisms-lemma-affineness-of-large-open}
$f : X \to Y$를 준유한이고 유한 표시인 아핀 스킴의 사상이라 하자.
$w : X \to \mathbf{Z}_{> 0}$를 $f$의 양의 가중치라 하고,
$d < \infty$를 $\int_f w$의 최댓값이라 하자. $Y$의 열린 부분
$$
Y_d = \{y \in Y \mid (\textstyle{\int}_f w)(y) = d \}
$$
는 아핀이다.
\end{lemma}

\begin{proof}
보조정리 \ref{more-morphisms-lemma-max-int-w}에 따라 $\int_f w$는 최댓값을 갖는다.
보조정리 \ref{more-morphisms-lemma-semicontinuous-int-w}에 따라 집합 $Y_d$는
열려 있다. 따라서 보조정리의 명제는 잘 정의되어 있다.

\medskip\noindent
Noether 경우로의 환원이다. 이 문단은 건너뛰어도 좋다. 유한형 사상이
준유한일 필요충분조건은 상대 차원이 $0$인 것임을 상기하자.
Morphisms, 보조정리
\ref{morphisms-lemma-locally-quasi-finite-rel-dimension-0}를 보라.
보조정리 \ref{more-morphisms-lemma-Noetherian-approximation-dimension-d}를 $d = 0$으로
적용하면 아핀 Noether 스킴의 준유한 사상 $f_0 : X_0 \to Y_0$와
사상 $Y \to Y_0$를 잡아 $f$가 $f_0$의 밑변환이 되게 할 수 있다.
그러면 $Y = \lim Y_i$를 $Y_0$ 위의 유한형 아핀 스킴들의 유향
극한으로 쓸 수 있다. Algebra, 보조정리
\ref{algebra-lemma-ring-colimit-fp}를 보라. 보조정리
\ref{more-morphisms-lemma-descend-weighting}에 따라 어떤 $i$가 존재하여 가중치
$w$가 $f_0$의 밑변환 $f_i : X_i \to Y_i$의 가중치 $w_i$로
내려간다. 이제 $f_i, w_i$에 대해 보조정리가 성립하면 $f$에
대해서도 성립한다. 실제로 $\int_f w$의 구성은 밑변환과 가환한다.
보조정리 \ref{more-morphisms-lemma-weighting-check-after-etale-base-change}를 보라.

\medskip\noindent
$X$와 $Y$가 Noether라고 가정하자. 사상 $g : Y' \to Y$에 의한
$f$의 밑변환을 $X' \to Y'$라 하자. $\int_f w$의 구성, 따라서
열린 부분 $Y_d$의 구성은 밑변환과 가환한다. $g$가 유한 전사이면
$Y'_d \to Y_d$도 유한 전사이다. 이 경우 $Y_d$가 아핀임을 보이는
것은 $Y'_d$가 아핀임을 보이는 것과 동치이다. Cohomology of Schemes,
보조정리
\ref{coherent-lemma-image-affine-finite-morphism-affine-Noetherian}를 보라.

\medskip\noindent
$T$가 $Y$ 위에서 유한인 몰입 $X \to T$를 잡을 수 있다. 보조정리
\ref{more-morphisms-lemma-quasi-finite-separated-pass-through-finite}를 보라.
유한 사상 $T \to Y$에 Morphisms, 보조정리
\ref{morphisms-lemma-massage-finite}를 적용하자. 이 보조정리에 따라
유한 전사 사상 $Y' \to Y$가 존재하여 $Y' \times_Y T$는 특별한 꼴을
갖는 $Y'$ 위의 유한 스킴 $T'$의 닫힌 부분스킴이 된다. 첫 문단의
논의에 따라 $Y$를 $Y'$로, $T$를 $T'$로, $X$를
$Y' \times_Y X$로 바꾸어도 된다. 따라서 반드시 열린 몰입이나 닫힌
몰입일 필요는 없는 몰입 $X \to T$와 닫힌 부분스킴들
$T_i \subset T$, $i = 1, \ldots, n$이 존재하고 다음을 만족한다고
가정할 수 있다.
\begin{enumerate}
\item $T \to Y$는 유한이고 국소 자유이다.
\item $T_i \to Y$는 동형사상이다.
\item 집합론적으로 $T = \bigcup_{i = 1, \ldots, n} T_i$이다.
\end{enumerate}
$Y$의 기약 성분들을 $Y$의 정역 닫힌 부분스킴으로 보고, 그 서로소
합을 $Y' = \coprod Y_k$라 하자. 그러면 $Y' \to Y$로 한 번 더
밑변환해도 된다. 여기서는 $Y$가 Noether라는 사실을 사용했다.
따라서 더 나아가 $Y$가 정역이고 Noether라고 가정할 수 있다.

\medskip\noindent
중복을 제거하여 $i \not = j$이면 $T_i \not = T_j$라고 가정할 수도
있고, 실제로 그렇게 하자. $Y$, 따라서 모든 $T_i$가 정역이므로,
$T_i$와 $T_j$가 교차한다면 그 교집합은 $Y$의 진 닫힌 부분집합으로
가는 닫힌 부분집합이다.

\medskip\noindent
$V_i = X \cap T_i$는 국소 닫힌 부분집합이면서 $X$의 닫힌
부분스킴이기도 하므로 아핀이다. $\eta \in Y$와 $\eta_i \in T_i$를
일반점들이라 하자. $\eta \not \in Y_d$이면
$Y_d = \emptyset$이므로 끝난다. $\eta \in Y_d$라고 가정하자.
$\eta_i \in V_i$인 지표 $i$들의 부분집합을
$I \in \{1, \ldots, n\}$이라 쓰자.
% P05-EN-CH37-0373
$i \in I$에 대하여 국소 닫힌 부분집합 $V_i \subset T_i$는 기약
공간 $T_i$의 일반점을 포함하므로 열려 있다. 한편 $f$가 열린
사상이므로(보조정리 \ref{more-morphisms-lemma-weighting-universally-open}),
임의의 $x \in X$에 대하여 어떤 $i \in I$와 특수화
$\eta_i \leadsto x$를 찾을 수 있다. 따라서 $x \in T_i$이고,
그러므로 $x \in V_i$이다. 다시 말해 집합론적으로
$X = \bigcup_{i \in I} V_i$이다.
$Y_d = \bigcap_{i \in I} \Im(V_i \to Y)$임을 주장한다. 그러면
$Y$의 아핀 열린 부분들 $\Im(V_i \to Y)$의 교집합이 아핀이므로
증명이 끝난다.

\medskip\noindent
$y \in Y$에 대하여 $X$에서
$f^{-1}(\{y\}) = \{x_1, \ldots, x_r\}$라 하자. 각 $i \in I$마다
$\eta_i \leadsto x_{j(i)}$를 만족하는
$j(i) \in \{1, \ldots, x_r\}$는 많아야 하나이다.
% P05-EN-CH37-0374
실제로 $j(i)$가 존재하여 $j$와 같을 필요충분조건은
$x_j \in V_i$인 것이다. $i \in I$에 대하여 $j = j(i)$가
존재한다면, $V_i \to Y$는 $x_j \mapsto y$의 한 근방에서
동형사상이다. 따라서 원천과 표적을 $x_j \mapsto y$의 근방들로
바꾼 뒤 $\bigcup_{i \in I,\ j(i) = j} V_i \to Y$는 유한이다.
그러므로 가중치의 정의에 따라
$w(x_j) = \sum_{i \in I,\ j(i) = j} w(\eta_i)$이다. 따라서
$$
(\textstyle{\int}_f w)(\eta) =
\sum\nolimits_{i \in I} w(\eta_i) \geq
\sum\nolimits_{j(i)\text{ exists}} w(\eta_i) =
\sum\nolimits_j w(x_j) = (\textstyle{\int}_f w)(y)
$$
를 얻는다. 등호가 성립할 필요충분조건은 $y$가
$\bigcap_{i \in I} \Im(V_i \to Y)$에 포함되는 것이며, 이것이
보이려던 바이다.
\end{proof}

\begin{proposition}
\label{more-morphisms-proposition-asn-weighting}
$f : X \to Y$를 스킴의 전사 준유한 사상이라 하고,
$w : X \to \mathbf{Z}_{> 0}$를 $f$의 양의 가중치라 하자.
$X$는 아핀이고 $Y$는 분리이며\footnote{$Y$의 대각선이
아핀이기만 해도 충분하다.} 공집합이 아니라고 가정하자. 그러면
$Y$의 아핀 층화수는 $\int_f w$가 갖는 서로 다른 값의 개수에서
$1$을 뺀 수 이하이다.
\end{proposition}

\begin{proof}
$Y$가 분리이므로 사상 $X \to Y$는 아핀임에 유의하자(Morphisms,
보조정리 \ref{morphisms-lemma-affine-permanence}). 보조정리
\ref{more-morphisms-lemma-max-int-w}에 따라 함수 $\int_f w$는 최댓값 $d$를 갖는다.
$d$에 관한 귀납법을 사용하겠다. $Y$의 열린 부분스킴
$Y_d = \{y \in Y \mid (\int_f w)(y) = d\}$를 생각하자.
$f^{-1}(Y_d) \to Y_d$가 유한임을 상기하자. 보조정리
\ref{more-morphisms-lemma-max-int-finite}를 보라. 보조정리
\ref{more-morphisms-lemma-affineness-of-large-open}에 따라 모든 아핀 열린 부분
$W \subset Y$에 대하여 $Y_d \cap W$는 아핀이다. 여기서는
$W \times_Y X$가 $X$ 위에서 아핀이므로 아핀이라는 사실을 사용했다.
따라서 $Y_d \to Y$는 스킴의 아핀 사상이다. 그러므로
$f^{-1}(Y_d) = Y_d \times_Y X$는 $X$ 위에서 아핀이므로 아핀
스킴이다. 이제 $f^{-1}(Y_d) \to Y_d$는 전사이므로 Limits, 보조정리
\ref{limits-lemma-affine}에 따라 $Y_d$는 아핀이다.
$Y_d = Y$이면 $Y$의 아핀 층화수는 $0$이고 명제가 성립한다.
% P05-EN-CH37-0375
그렇지 않다면 $Y' = Y \setminus Y_d$를 $Y$의 축약 닫힌
부분스킴으로 보자. $X' = Y' \times_Y X$로 둔다. $X'$는 $X$에서
닫혀 있으므로 아핀이고, $Y'$는 $Y$에서 닫혀 있으므로 분리이다.
사상 $f' : X' \to Y'$는 전사이고, $w$는 $f'$의 가중치 $w'$를
유도하며(보조정리 \ref{more-morphisms-lemma-weighting-base-change}),
$\int_{f'} w'$는 $\int_f w$의 $Y'$로의 제한이다. 귀납 가정에 따라
$Y'$는 길이가 $\leq$ $\int_{f'} w'$의 서로 다른 값의 개수에서
$1$을 뺀 수인 아핀 층화를 가지므로 증명이 끝난다.
\end{proof}









\section{완전 분해 사상}
\label{more-morphisms-section-completely-decomposed}

\noindent
Nishnevich는 오늘날 Nishnevich 위상이라 부르는 것을 정의하기 위해
에탈 사상의 완전 분해 족이라는 개념을 연구했다. 예를 들어
\cite{Nishnevich}를 보라.

\begin{definition}
\label{more-morphisms-definition-cd-morphism}
스킴의 사상 $f : X \to Y$가 {\it 완전 분해}되었다는 것은\footnote{이는
비표준 용어일 수 있다.} 모든 점 $y \in Y$에 대하여
$f(x) = y$이고 체 확대 $\kappa(x)/\kappa(y)$가 자명한 점
$x \in X$가 존재한다는 뜻이다. 표적이 고정된 스킴의 사상족
$\{f_i : X_i \to Y\}_{i \in I}$가 {\it 완전 분해}되었다는 것은
$\coprod f_i : \coprod Y_i \to X$가 완전 분해되었다는 뜻이다.
% P05-EN-CH37-0376
\end{definition}

\noindent
몇 가지 기본 보조정리로 시작한다.

\begin{lemma}
\label{more-morphisms-lemma-composition-cd}
스킴의 완전 분해 사상 두 개의 합성은 완전 분해된다.
$\{f_i : X_i \to Y\}_{i \in I}$가 완전 분해되고, 각 $i$에 대하여
완전 분해 족 $\{X_{ij} \to X_i\}_{j \in J_i}$가 주어졌다면
족 $\{X_{ij} \to Y\}_{i \in I, j \in J_i}$는 완전 분해된다.
\end{lemma}

\begin{proof}
생략한다.
\end{proof}

\begin{lemma}
\label{more-morphisms-lemma-base-change-cd}
스킴의 완전 분해 사상의 밑변환은 완전 분해된다.
$\{f_i : X_i \to Y\}_{i \in I}$가 완전 분해되고
$Y' \to Y$가 스킴의 사상이면
$\{X_i \times_Y Y' \to Y'\}_{i \in I}$는 완전 분해된다.
\end{lemma}

\begin{proof}
$f : X \to Y$와 $g : Y' \to Y$를 스킴의 사상이라 하자.
$y' \in Y'$의 $Y$에서의 상을 $y = g(y')$라 하자.
$f(x) = y$이고 $\kappa(x) = \kappa(y)$인 점 $x \in X$가 있으면,
$Y'$에서 $y'$로 가고 $X$에서 $x$로 가며
$\kappa(x') = \kappa(y')$인 유일한 점
$x' \in X' = X \times_Y Y'$가 존재한다. Schemes, 보조정리
\ref{schemes-lemma-points-fibre-product}를 보라. 이 사실에서
보조정리가 쉽게 따르므로 세부사항은 생략한다.
\end{proof}

\begin{lemma}
\label{more-morphisms-lemma-decompose}
\begin{reference}
\cite[Lemma 2.1.2]{EHIK}
\end{reference}
$f : X \to Y$를 스킴의 사상이라 하자. $f$가 완전 분해되고
$f$가 국소 유한 표시이며, $Y$가 준콤팩트이고 준분리라고 가정하자. 그러면
$n \geq 0$과 사상들 $Z_i \to Y$, $i = 1, \ldots, n$이 존재하여
다음 성질들을 만족한다.
\begin{enumerate}
\item $\coprod Z_i \to Y$는 전사이다.
\item 모든 $i$에 대하여 $Z_i \to Y$는 몰입이다.
\item 모든 $i$에 대하여 $Z_i \to Y$는 유한 표시이다.
\item 모든 $i$에 대하여 밑변환 $X \times_Y Z_i \to Z_i$는
단면을 갖는다.
\end{enumerate}
\end{lemma}

\begin{proof}
$y \in Y$라 하자. 가정에 따라 $Y$ 위의 사상
$\sigma : \Spec(\kappa(y)) \to X$가 존재한다.
$\Spec(\kappa(y))$를 $Y$ 위의 아핀 스킴 $Z$들의 유향 극한으로
쓸 수 있는데, 여기서 $Z \to Y$는 유한 표시 몰입이다. 구체적으로
아핀 열린 부분 $y \in \Spec(A) \subset Y$를 잡고 $y$가 $A$의
소아이디얼 $\mathfrak p$에 대응한다고 하자. 그러면
$\kappa(\mathfrak p)$는 환들 $(A/I)_f$의 여과 여극한이다. 여기서
$I \subset \mathfrak p$는 유한 생성 아이디얼이고
$f \in A$, $f \not \in \mathfrak p$이다. 사상들
$Z = \Spec((A/I)_f) \to Y$는 유한 표시 몰입이며, $Y$가 준분리이므로
$Z \to Y$는 준콤팩트이다. Schemes, 보조정리
\ref{schemes-lemma-quasi-compact-permanence}를 보라. Limits, 명제
\ref{limits-proposition-characterize-locally-finite-presentation}에
따라 그러한 어떤 $Z$에 대하여 $Y$ 위의 사상
$\sigma' : Z \to X$가 존재하고, $\kappa(\mathfrak p)$의 스펙트럼
위에서 $\sigma$와 일치한다. $\sigma'$가 $Y$ 위의 사상이므로
사영 $X \times_Y Z \to Z$의 단면을 얻는다.
% P05-EN-CH37-0377

\medskip\noindent
따라서 $Y$는 $X \times_Y Z \to Z$가 단면을 갖는 유한 표시 몰입
$Z \to Y$들의 상들의 합집합이다. $Z \to Y$의 상은 구성가능하고
(Morphisms, 보조정리 \ref{morphisms-lemma-chevalley}), $Y$는
구성가능 위상에서 콤팩트이므로(Properties, 보조정리
\ref{properties-lemma-quasi-compact-quasi-separated-spectral} 및
Topology, 보조정리
\ref{topology-lemma-constructible-hausdorff-quasi-compact}), 그중 유한
개만으로 충분하다.
\end{proof}

\begin{lemma}
\label{more-morphisms-lemma-descend-cd}
$S = \lim_{\lambda \in \Lambda} S_\lambda$를 전이 사상들이 아핀인
스킴의 유향계의 극한이라 하자. $0 \in \Lambda$라 하고,
$f_0 : X_0 \to Y_0$를 $S_0$ 위의 스킴의 사상이라 하자.
$\lambda \geq 0$에 대하여 $f_\lambda : X_\lambda \to Y_\lambda$를
$f_0$의 $S_\lambda$로의 밑변환이라 하고, $f : X \to Y$를
$f_0$의 $S$로의 밑변환이라 하자. 다음을 가정하자.
\begin{enumerate}
\item $f$는 완전 분해된다.
\item $Y_0$는 준콤팩트이고 준분리이다.
\item $f_0$는 국소 유한 표시이다.
\end{enumerate}
그러면 $f_\lambda$가 완전 분해되는 $\lambda \geq 0$가 존재한다.
% P05-EN-CH37-0378
\end{lemma}

\begin{proof}
$Y_0$가 준콤팩트이고 준분리이며 $Y$가 $Y_0$ 위에서 아핀이므로
그 스킴도 준콤팩트이고 준분리이다. 보조정리 \ref{more-morphisms-lemma-decompose}에서와
같은 $n \geq 0$과 $Z_i \to Y$, $i = 1, \ldots, n$을 잡자.
$Z_i$의 선택에 의해 존재하는 $Y$ 위의 사상들을
$\sigma_i : Z_i \to X$로 쓰자. $0 \in \Lambda$를 더 큰 지표로
바꾼 뒤, $S$로 밑변환하면 $Z_i \to Y$가 되는 유한 표시 사상
$Z_{i, 0} \to Y_0$들이 존재한다고 가정할 수 있다. Limits, 보조정리
\ref{limits-lemma-descend-finite-presentation}를 보라. Limits, 보조정리
\ref{limits-lemma-descend-immersion}에 따라 필요하면 $0$을 더 키운
뒤 $Z_{i, 0} \to Y_0$가 몰입이라고 가정할 수 있다.
$\coprod Z_i \to Y$가 전사이므로, 필요하면 $0$을 더 키운 뒤
$\coprod Z_{i, 0} \to Y_0$가 전사라고 가정할 수 있다. Limits,
보조정리 \ref{limits-lemma-descend-surjective}를 보라. 다음에 유의하자.
$Z_i = \lim_{\lambda \geq 0} Z_{i, \lambda}$
여기서 $Z_{i, \lambda} = Y_\lambda \times_{Y_0} Z_{i, 0}$이다.
합성
$$
Z_i \xrightarrow{\sigma_i} X \to X_0
$$
들을 $Y_0$ 위의 사상으로 보자. $f_0$가 국소 유한 표시이므로
Limits, 명제
\ref{limits-proposition-characterize-locally-finite-presentation}에 따라
어떤 $\lambda \geq 0$와 $Y_0$ 위의 사상
$\sigma'_{i, \lambda} : Z_{i, \lambda} \to X_0$들이 존재하여,
$Z_i \to Z_{i, \lambda}$와 먼저 합성하면 위에 표시한 화살표가
되게 할 수 있다. 물론 그러면 $\sigma'_{i, \lambda}$는
$Y_\lambda$ 위의 사상 $\sigma_{i, \lambda} : Z_{i, \lambda} \to
X_\lambda = X_0 \times_{Y_0} Y_\lambda$를 결정한다.
$\coprod Z_{i, \lambda} \to Y_\lambda$가 전사이므로
$X_\lambda \to Y_\lambda$는 완전 분해된다.
\end{proof}




\section{가역 가군들의 풍부한 족}
\label{more-morphisms-section-families-ample-invertible-modules}

\noindent
Morphisms, 절
\ref{morphisms-section-families-ample-invertible-modules}의 논의를
계속한다.

\begin{lemma}
\label{more-morphisms-lemma-ample-family-ample-relative}
$f : X \to Y$를 스킴의 사상이라 하자. 다음을 가정하자.
\begin{enumerate}
\item $Y$는 가역 가군들의 풍부한 족을 갖는다.
\item $X$ 위에 $f$-풍부한 가역 가군이 존재한다.
\end{enumerate}
그러면 $X$는 가역 가군들의 풍부한 족을 갖는다.
\end{lemma}

\begin{proof}
$\mathcal{L}$을 $X$ 위의 $f$-풍부한 가역 가군이라 하자. 특히 이는
$f$가 준콤팩트임을 뜻한다. Morphisms, 정의
\ref{morphisms-definition-relatively-ample}를 보라. Morphisms, 정의
\ref{morphisms-definition-family-ample-invertible-modules}에 따라
$Y$가 준콤팩트이므로 $X$도 준콤팩트이다. 따라서 $X$ 자체가
Morphisms, 정의
\ref{morphisms-definition-family-ample-invertible-modules}의 첫 번째
조건을 만족한다. $x \in X$의 상을 $y \in Y$라 하자. 가정 (2)에
따라 가역 $\mathcal{O}_Y$-가군 $\mathcal{N}$과 단면
$t \in \Gamma(Y, \mathcal{N})$을 잡아, $t$가 영이 아닌 자취
$Y_t$가 아핀이 되게 할 수 있다.
% P05-EN-CH37-0379
그러면 $\mathcal{L}$은 $f^{-1}(Y_t) = X_{f^*t}$ 위에서 풍부하므로
$x$를 포함하고 $(X_{f^*t})_s$가 아핀이 되게 하는 단면
$s \in \Gamma(X_{f^*t}, \mathcal{L})$을 찾을 수 있다. Properties,
보조정리 \ref{properties-lemma-invert-s-sections}에 따라 어떤
$n \geq 0$에 대하여 곱 $(f^*t)^n s$는 단면
$s' \in \Gamma(X, f^*\mathcal{N}^{\otimes n} \otimes \mathcal{L})$으로
확장된다. 마지막으로
$f^*\mathcal{N}^{\otimes n + 1} \otimes \mathcal{L}$의 단면
$s'' = f^* ts'$는 $X \setminus X_{f^*t}$의 모든 점에서 영이다.
따라서 원하는 대로 $X_{s''} = (X_{f^*t})_s$는 아핀이다.
\end{proof}

\begin{lemma}
\label{more-morphisms-lemma-resolution-property-goes-up-affine}
$f : X \to Y$를 스킴의 아핀 사상 또는 준아핀 사상이라 하자.
$Y$가 가역 가군들의 풍부한 족을 가지면 $X$도 그러하다.
\end{lemma}

\begin{proof}
Morphisms, 보조정리 \ref{morphisms-lemma-quasi-affine-O-ample}에 따라
이는 보조정리 \ref{more-morphisms-lemma-ample-family-ample-relative}의 특수한
경우이다.
\end{proof}






\section{블로업과 가역 가군들의 풍부한 족}
\label{more-morphisms-section-ample-family-by-blowing-up}

\noindent
\cite{Gross-thesis}의 한 결과를 증명한다.

\begin{lemma}
\label{more-morphisms-lemma-get-ample-family}
$X$를 스킴이라 하자. $X$ 위의 유효 Cartier 인자
$D_1, \ldots, D_m$과 가역 가군
$\mathcal{L}_1, \ldots, \mathcal{L}_m$이 주어져
$\bigcap D_i = \emptyset$이고
$\mathcal{L}_i|_{X \setminus D_i}$가 풍부하다고 하자. 그러면
$X$는 가역 가군들의 풍부한 족을 갖는다.
\end{lemma}

\begin{proof}
$x \in X$라 하자. $x \not \in D_i$인 지표
$i \in \{1, \ldots, m\}$을 잡고 $U_i = X \setminus D_i$로 두자.
$\mathcal{L}_i|_{U_i}$가 풍부하므로 어떤 $n \geq 1$과 단면
$s \in \Gamma(U_i, \mathcal{L}_i^{\otimes n})$을 잡아, $s$가 영이
아닌 자취 $(U_i)_s$가 아핀이 되게 할 수 있다(Properties, 정의
\ref{properties-definition-ample}).
% P05-EN-CH37-0380
$U_i$는 표준 단면 $1 \in \mathcal{O}_X(D_i)$가 영이 아닌
자취이므로, Properties, 보조정리
\ref{properties-lemma-invert-s-sections}에 따라 어떤 $N \geq 0$가
존재하여 $s$가 단면
$$
s' \in \Gamma(X, \mathcal{L}_i^{\otimes n} \otimes_{\mathcal{O}_X}
\mathcal{O}_X(N D_i))
$$
으로 확장되게 할 수 있다. $N$을 $N + 1$로 바꾼 뒤에는 $s'$가
$D_i$의 모든 점에서 영이고, 따라서 $X_{s'} = (U_i)_s$가
아핀임을 알 수 있다. 이로써 $X$가 가역 가군들의 풍부한 족을
가짐이 증명된다. Morphisms, 정의
\ref{morphisms-definition-family-ample-invertible-modules}를 보라.
\end{proof}

\begin{lemma}
\label{more-morphisms-lemma-blow-up-ample-family}
\begin{reference}
\cite[Proposition 1.3.1]{Gross-thesis}
\end{reference}
$X$를 기약 성분이 유한 개인 준콤팩트 준분리 스킴이라 하자.
준콤팩트 조밀 열린 부분 $U \subset X$와 $U$-허용 블로업
$X' \to X$가 존재하여 스킴 $X'$가 가역 가군들의 풍부한 족을 갖는다.
\end{lemma}

\begin{proof}
$\eta_1, \ldots, \eta_n \in X$를 $X$의 기약 성분들의 일반점이라
하자. Properties, 보조정리
\ref{properties-lemma-point-and-maximal-points-affine}와 $X$의
준콤팩트성에 따라 유한 아핀 열린 덮개
$X = U_1 \cup \ldots \cup U_m$을 잡아 각 $U_i$가
$\eta_1, \ldots, \eta_n$을 모두 포함하게 할 수 있다. 특히 준콤팩트
열린 부분 $U = U_1 \cap \ldots \cap U_m$은 $X$에서 조밀하다.
$Z_i = V(\mathcal{I}_i)$일 때 $U_i = X \setminus Z_i$가 되게 하는
유한형 준연접 아이디얼 층
$\mathcal{I}_i \subset \mathcal{O}_X$를 잡자. Properties, 보조정리
\ref{properties-lemma-quasi-coherent-finite-type-ideals}를 보라.
$$
f : X' \longrightarrow X
$$
를 아이디얼 층 $\mathcal{I} = \mathcal{I}_1 \cdots \mathcal{I}_m$에
관한 $X$의 블로업이라 하자. $V(\mathcal{I})$는 집합론적으로
$U$와 서로소인 $Z_i$들의 합집합이므로 $f$는 $U$-허용 블로업이다.
또한 $f$는 사영적 사상이고 $\mathcal{O}_{X'}(1)$은 $f$-상대적으로
풍부하다. Divisors, 보조정리
\ref{divisors-lemma-blowing-up-projective}를 보라. Divisors, 보조정리
\ref{divisors-lemma-blowing-up-two-ideals}에 따라 각 $i$에 대하여
사상 $f'$은 $X' \to X'_i \to X$로 분해된다.
% P05-EN-CH37-0381
여기서 $X'_i \to X$는 $\mathcal{I}_i$에 관한 블로업이고,
$X' \to X'_i$는 또 하나의 블로업이다. 구체적으로 후자는
$\mathcal{I}_i$를 제외한 아이디얼들 $\mathcal{I}_j$의 곱을 당긴
것에 관한 블로업이다. 따라서
$D_i = f^{-1}(Z_i) \subset X'$는 유효 Cartier 인자이다. Divisors,
보조정리 \ref{divisors-lemma-blow-up-pullback-effective-Cartier}와
\ref{divisors-lemma-blowing-up-gives-effective-Cartier-divisor}를 보라.
$X' \setminus D_i = f^{-1}(U_i)$이다. $\mathcal{O}_{X'}(1)$이
$f$-풍부하므로 $\mathcal{O}_{X'}(1)$을 $X' \setminus D_i$에
제한한 것은 풍부하다.
보조정리 \ref{more-morphisms-lemma-get-ample-family}에 따라 $X'$는 가역 가군들의
풍부한 족을 갖는다.
\end{proof}

\begin{proposition}
\label{more-morphisms-proposition-envelope-with-resolution-property}
$X$를 준콤팩트 준분리 스킴이라 하자. 유한 표시이고 고유하며 완전
분해된 사상 $f : Y \to X$가 존재하여(정의
\ref{more-morphisms-definition-cd-morphism}), 스킴 $Y$는 가역 가군들의 풍부한
족을 갖는다.
\end{proposition}

\begin{proof}
Limits, 명제 \ref{limits-proposition-approximate}에 따라 아핀 사상
$X \to X_0$가 존재하며, 여기서 $X_0$는 $\mathbf{Z}$ 위의 유한형
스킴이다. 아래에서는 원하는 모든 성질을 갖는 사상
$Y_0 \to X_0$를 구성한다. 그러면
$Y = X \times_{X_0} Y_0$로 두고 $f$를 사영
$f : Y \to X$로 두어 결론을 얻는다. 이를 보려면 $f$가 유한
표시이고(Morphisms, 보조정리
\ref{morphisms-lemma-base-change-finite-presentation}), $f$가
고유하며(Morphisms, 보조정리
\ref{morphisms-lemma-base-change-proper}), $f$가 완전
분해된다는 사실(보조정리 \ref{more-morphisms-lemma-base-change-cd})에 유의하면 된다.
마지막으로 $Y \to Y_0$는 $X \to X_0$의 밑변환이므로 아핀이고,
따라서 보조정리 \ref{more-morphisms-lemma-resolution-property-goes-up-affine}에
따라 $Y$는 가역 가군들의 풍부한 족을 갖는다. 이로써 다음 문단에서
다루는 경우로 환원된다.

\medskip\noindent
$X$가 $\mathbf{Z}$ 위의 유한형이라고 가정하자. 특히
$\dim(X) < \infty$이다. $\dim(X)$에 관한 귀납법을 사용하겠다.
$\dim(X) = 0$이면 $X$는 아핀이고 분해 성질을 갖는다.
% P05-EN-CH37-0382
일반적으로 조밀 열린 부분 $U \subset X$와 $U$-허용 블로업
$X' \to X$가 존재하여 $X'$가 가역 가군들의 풍부한 족을 갖는다.
보조정리 \ref{more-morphisms-lemma-blow-up-ample-family}를 보라.
$f : X' \to X$가 $U$ 위에서 동형사상이므로 $U$의 모든 점은 같은
잉여체를 갖는 $X'$의 점으로 올라간다. $Z = X \setminus U$에
유도된 축약 스킴 구조를 주자. $U$가 $X$에서 조밀하므로 위에서
본 대로 $\dim(Z) < \dim(X)$이다. 귀납 가정에 따라 고유이고 완전
분해된 사상 $W \to Z$를 잡아 $W$가 가역 가군들의 풍부한 족을
갖게 할 수 있다. 그러면 $Y = X' \amalg W \to X$가 원하는
사상이다.
\end{proof}




\section{닫힌 몰입에 대한 확장 판정법}
\label{more-morphisms-section-extensive-criterion-closed-immersions}

\noindent
이 절에서는 스킴의 사상이 닫힌 몰입일 판정법을 제시한다.

\begin{lemma}
\label{more-morphisms-lemma-affine-extensive-criterion}
아핀 스킴의 사상 $f : X \to Y$가 닫힌 몰입일 필요충분조건은,
모든 단사 환 사상 $A \to B$와 가환 사각형
$$
\xymatrix{
\Spec(B) \ar[d] \ar[r] & X \ar[d]^f \\
\Spec(A) \ar[r] \ar@{..>}[ur] & Y
}
$$
에 대하여 두 삼각형을 가환하게 하는 올림 $\Spec(A) \to X$가
존재하는 것이다.
\end{lemma}

\begin{proof}
사상 $f$가 환 사상 $\phi : R \to S$로 주어진다고 하자. 그러면
$f$가 닫힌 몰입일 필요충분조건은 $\phi$가 전사인 것이다.

\medskip\noindent
먼저 $\phi$가 전사라고 가정하자. $\psi : A \to B$를 단사 환
사상이라 하고, 가환 도식
$$
\xymatrix{
R \ar[r]^\alpha \ar[d]^\phi & A \ar[d]^\psi \\
S \ar[r]^\beta \ar@{..>}[ur] & B
}
$$
이 주어졌다고 하자. $\phi(r) = s$인 $r \in R$에 대하여
$s \mapsto \alpha(r)$로 두어 올림 $S \to A$를 정의한다.
$\psi$가 단사이고 사각형이 가환하므로 이는 잘 정의된다. 환의
스펙트럼을 취하는 것은 가환환과 아핀 스킴 사이의 반동치를
정의하므로 $f$에 대해 원하는 올림 성질이 성립한다.

\medskip\noindent
다음으로 $\phi$가 모든 단사 환 사상 $\psi: A \to B$에 대한
올림을 갖는다고 가정하자. $\phi(R)$는 $S$의 부분환이므로 가환
사각형
$$
\xymatrix{
R \ar[r] \ar[d]^\phi  & \phi(R) \ar[d] \\
S \ar@{=}[r] \ar@{..>}[ur] & S
}
$$
을 얻고, 여기에는 올림 $S \to \phi(R)$가 존재한다. 따라서 포함
$\phi(R) \to S$는 동형사상이어야 한다. 그러므로 $\phi$는
전사이고 증명이 끝난다.
\end{proof}

\begin{lemma}
\label{more-morphisms-lemma-scheme-affine-if-aff-is-split-mono}
$X$를 스킴이라 하자. Schemes, 보조정리
\ref{schemes-lemma-morphism-into-affine}의 표준 사상
$X \to \Spec(\Gamma(X, \mathcal{O}_X))$가 수축을 가지면 $X$는
아핀 스킴이다.
\end{lemma}

\begin{proof}
$S = \Spec(\Gamma(X, \mathcal{O}_X))$로 쓰고, 보조정리의 사상을
$f : X \to S$라 하자. $s : S \to X$를 수축이라 하면
$\text{id}_X = sf$이다. 따라서 $f s f = \text{id}_S f$이다.
$f$는 동형사상
$\Gamma(S, \mathcal{O}_S) \to \Gamma(X, \mathcal{O}_X)$를
유도하므로, 이는 $fs$와 $\text{id}_S$가
$\Gamma(S, \mathcal{O}_S)$ 위에서 같은 사상을 유도함을 뜻한다.
$S$가 아핀이므로 $f s = \text{id}_S$이다. 따라서 $f$는
동형사상이고, 보이려던 대로 $X$는 아핀 스킴이다.
\end{proof}

\begin{lemma}
\label{more-morphisms-lemma-aff-is-injective-on-sheaves}
$X$를 스킴이라 하고,
$f : X \to S = \Spec(\Gamma(X, \mathcal{O}_X))$를 Schemes, 보조정리
\ref{schemes-lemma-morphism-into-affine}의 표준 사상이라 하자.
$f^\sharp : \mathcal{O}_S \to f_*\mathcal{O}_X$의 핵에 포함되는 가장
큰 준연접 $\mathcal{O}_S$-가군은 영이다. $X$가 준콤팩트이면
$f^\sharp$는 단사이다. 특히 $X$가 준콤팩트이면 $f$는 우세 사상이다.
\end{lemma}

\begin{proof}
$f^\sharp$의 핵에 포함되는 가장 큰 준연접 $\mathcal{O}_S$-가군에
대응하는 부분가군을 $M \subset \Gamma(S, \mathcal{O}_S)$이라 하자.
그러면 임의의 원소 $a \in M$는 $f^\sharp$에 의해 영으로 간다.
그러나 $f^\sharp(a)$는
$$
\Gamma(S, f_*\mathcal{O}_X) = \Gamma(X, \mathcal{O}_X) =
\Gamma(S, \mathcal{O}_S)
$$
에서 $a$ 자체에 대응하는 원소이다. 따라서 $a = 0$이고 $M = 0$이다.
이로써 첫 번째 주장이 증명된다. 이는 사상 $f : X \to S$가
스킴론적으로 전사인 것과 동치임에 유의하자.

\medskip\noindent
$X$가 준콤팩트이면 Morphisms, 보조정리
\ref{morphisms-lemma-quasi-compact-scheme-theoretic-image}에 따라
$\Ker(f^\sharp)$는 준연접이다. 따라서 $\Ker(f^\sharp) = 0$이고
$f^\sharp$는 단사이다. 이 경우 Morphisms, 보조정리
\ref{morphisms-lemma-quasi-compact-scheme-theoretic-image}의 (4)에
따라 $f$는 우세 사상이다.
\end{proof}

\begin{lemma}
\label{more-morphisms-lemma-extensive-criterion}
$f: X \to Y$를 스킴의 준콤팩트 사상이라 하자. 그러면 $f$가 닫힌
몰입일 필요충분조건은 모든 단사 환 사상 $A \to B$와 가환 사각형
$$
\xymatrix{
\Spec(B) \ar[d] \ar[r] & X \ar[d]^f \\
\Spec(A) \ar[r] \ar@{..>}[ur] & Y
}
$$
에 대하여 도식을 가환하게 하는 올림 $\Spec A \to X$가 존재하는
것이다.
\end{lemma}

\begin{proof}
$f$가 닫힌 몰입이라고 가정하자. $A \to B$를 단사 환 사상이라
하고 가환 사각형
$$
\xymatrix{
\Spec(B) \ar[d] \ar[r]  & X \ar[d]^f \\
\Spec(A) \ar[r] \ar@{..>}[ur] & Y
}
$$
을 생각하자. 그러면 $\Spec(A) \times_Y X \to \Spec(A)$는 닫힌
몰입이므로 아이디얼 $I \subset A$와 가환 도식
$$
\xymatrix{
\Spec(B) \ar[d] \ar[r] & \Spec(A/I) \ar[r] \ar[d] & X \ar[d]^f \\
\Spec(A) \ar[r] \ar@{..>}[ur] & \Spec(A) \ar[r] & Y
}
$$
을 얻는다. 보조정리 \ref{more-morphisms-lemma-affine-extensive-criterion}에 따라
올림을 얻는다.

\medskip\noindent
이제 $f$가 보조정리에 명시된 올림 성질을 갖는다고 가정하자.
$f$가 닫힌 몰입임을 보이는 것은 $Y$ 위에서 국소적이므로 $Y$가
아핀이라고 가정해도 좋고, 실제로 그렇게 하자. 특히 $Y$는
준콤팩트이고, 따라서 $X$도 준콤팩트이다. 그러므로 유한 아핀 열린
덮개 $X = U_1 \cup \ldots \cup U_n$가 존재한다. 사상
$$
\pi : U = \coprod U_i \longrightarrow X
$$
의 원천은 아핀이고 유도된 환 사상
$\Gamma(X, \mathcal{O}_X) \to \Gamma(U, \mathcal{O}_U)$는 단사이다.
가정에 따라 도식
$$
\xymatrix{
U \ar[r]^\pi \ar[d] & X \ar[d]^f \\
\Spec(\Gamma(X, \mathcal{O}_X)) \ar[r]^-{f'} \ar@{..>}[ur]^h & Y
}
$$
에 올림이 존재한다. 여기서 $f'$은 환 사상
$\Gamma(Y, \mathcal{O}_Y) \to \Gamma(X, \mathcal{O}_X)$에 대응하는
아핀 스킴의 사상이다. $\pi$가 전사상이라는 사실에서 사상 $h$가
표준 사상 $X \to \Spec(\Gamma(X, \mathcal{O}_X))$의 수축임이
따른다. 세부사항은 생략한다. 따라서 보조정리
\ref{more-morphisms-lemma-scheme-affine-if-aff-is-split-mono}에 따라 $X$는 아핀이다.
보조정리 \ref{more-morphisms-lemma-affine-extensive-criterion}에 따라 $f$는 닫힌
몰입이다.
\end{proof}
